% =====================================================================
% SGA 1 中文翻译重排
% 《平展覆盖与基本群》(Révêtements étales et groupe fondamental)
% Séminaire de Géométrie Algébrique du Bois Marie 1960-61
% 由 A. Grothendieck 主持的讨论班, 附 M. Raynaud 女士的两个讲
% 源: Lecture Notes in Mathematics 224, Springer-Verlag, 1971
%     (arXiv:math/0206203 重排版)
% 编译方式: xelatex -shell-escape (两遍以上; 术语索引用 xindy, 记号索引用 makeindex)
% =====================================================================
\documentclass[11pt, fontsetup=font-setup-open.tex, titlesetup=titles-setup.tex, geometry=a4]{AJbook}

\usepackage{mycommand}
\usepackage{myarrows}
\usepackage{stmaryrd}
\usepackage{tikz-cd}
\usepackage{imakeidx}
\usepackage[bookmarksnumbered=true, colorlinks=true,
	linkcolor=MidnightBlue, citecolor=BrickRed, urlcolor=MidnightBlue,
	bookmarksnumbered=true, pdftitle={平展覆盖与基本群 (SGA 1)}]{hyperref}

% ---- 章编号: 按原书用罗马数字 (Exposé I-XIII, 无 VII) ----
\renewcommand{\thechapter}{\Roman{chapter}}
% \thesection 默认 = \thechapter.\arabic{section}, 定理编号随之
% \thetheorem = \thesection.\arabic{theorem}, 如 IX.4.7, 与原书 "IX 4.7" 对应

% ---- 章标题: "第 I 讲 ..." ----
\renewcommand{\AJchapterttl}[1]{第\enspace #1\enspace 讲}

% ---- 天眉: 章显示 "第 IX 讲", 节显示 "§IX.4" ----
\AtBeginDocument{
	\renewcommand{\chaptermark}[1]{\markboth{第\thechapter 讲\quad #1}{}}
	\renewcommand{\sectionmark}[1]{\markright{\S\thesection \quad #1}}
}

% ---- 索引: 术语索引 + 记号索引 (与原文 Index terminologique / Index des notations 对应) ----
\makeindex[
	columns=2,
	name=term,
	program=truexindy,
	intoc=true,
	options=-M texindy -I xelatex -C utf8,
	title={术语索引}]
% 记号索引: SGA 的记号以数学符号为主, xindy 的符号排序模块无法解析
% (如 $\pi_1(S,a)$ 等条目), 故记号索引沿用 makeindex; 术语索引用 xindy.
\makeindex[
	columns=3,
	name=notations,
	intoc=true,
	title={记号索引}]

% ---- 参考文献: 使用 biblatex + biber, 各讲内部独立编号 (refsection) ----
\addbibresource{references.bib}

% MinerU 通用: 压缩列表间距
\providecommand{\tightlist}{%
  \setlength{\itemsep}{0pt}\setlength{\parskip}{0pt}}

\begin{document}

% 超链接颜色按模板: 内部链接 MidnightBlue, 文献引用 BrickRed, 网址 MidnightBlue
\hypersetup{colorlinks=true, linkcolor=MidnightBlue, citecolor=BrickRed, urlcolor=MidnightBlue}

% ===================== 封面 =====================
\begin{titlepage}
\centering
\vspace*{2.0cm}
{\Huge\bfseries 平展覆盖与基本群\par}
\vspace{0.9cm}
{\large\bfseries Révêtements étales et groupe fondamental\par}
\vspace{1.6cm}
{\large Séminaire de Géométrie Algébrique du Bois Marie, 1960--61\par}
\vspace{0.7cm}
{\large 由 A. Grothendieck 主持\par}
\vspace{0.4cm}
{\large 附 M. Raynaud 女士的两个讲\par}
\vspace{3.2cm}
{\large 中文翻译重排版\par}
\vfill
{\small 依据 Lecture Notes in Mathematics 224, Springer-Verlag (1971)\par}
{\small 及 arXiv:math/0206203 重排版 (法国数学会 (SMF) 版) 翻译排版\par}
\vspace{0.5em}
{\normalsize 译者 翻译及重排\par}
\end{titlepage}

\frontmatter
% =====================================================================
% 前言 (Préface), 导论 (Introduction), 说明 (Avertissement)
% =====================================================================

\chapter*{前言 (Préface)}

本文是《平展覆盖与基本群》(Révêtements Étales et Groupe Fondamental)，Lecture Notes in Mathematics 224，Springer-Verlag，Berlin-Heidelberg-New York，1971，Alexander Grothendieck 等人一书的重排注释版。

LaTeX2e 的排版由一批志愿者完成，他们参与了一个由 Bas Edixhoven 主持的项目，有关详情见网址 http://www.math.leidenuniv.nl/\textasciitilde edix/. 版式由法国数学会 (Société mathématique de France) 最终完成.

这是对原文略作订正的一个版本. 它由法国数学会 (http://smf.emath.fr/) 在``数学文献''(Documents Mathématiques) 丛书中出版. Michel Raynaud 作了一些更新说明. 这些说明以方括号 [ ] 界定，并用记号 (MR) 标出 (第 211，220，256，292 页的注记以及第 67 页的注 7). 为统一记号，点 $x$ 的剩余域 (residue field) 记为 $k(x)$，局部环 $A$ 的剩余域记为 $k(A)$.

另有一个电子版本，意在复现原文.

两个版本共用一个源文件，该文件位于 arXiv.org 电子预印本服务器 http://arxiv.org/ 上. 两个版本之间的差异在源文件中均有记载.

旧页码已并入页边 (数字 $n$ 表示第 $n$ 页的开头).

\cleardoublepage

\chapter*{导论 (Introduction)}

在本导论的第一部分，我们对本卷的内容作一些说明；在第二部分，则对``Bois-Marie 代数几何讨论班''整体作说明，本卷构成其第一卷.

\section*{1.}

本卷在``Kronecker 式''的观点下，呈现了代数几何 (Géométrie Algébrique) 中基本群 (fundamental group) 理论的基础，该观点使人能够以同等的地位处理通常意义下的代数簇 (algebraic variety) 的情形，以及例如一个数域 (number field) 的整数环 (ring of integers) 的情形. 这一观点只有在概形 (scheme) 的语言中才能得到令人满意的表述，我们将自由地使用这种语言，以及 J. Dieudonné 与 A. Grothendieck 的《代数几何原理》(Éléments de Géométrie Algébrique，下文引用为 EGA) 前三章所阐述的主要结果. 研读本卷``Bois-Marie 代数几何讨论班''不需要代数几何的其他知识，因此对一位希望熟悉这些技巧的读者而言，它可作为现代代数几何技术的一个入门.

本书的讲 (Exposé) I 至 XI 是讨论班口头发言油印笔记的逐字复制，几乎未加改动，这些笔记曾由高等科学研究院 (Institut des Hautes Études Scientifiques)\footnote{以及随后的讨论班笔记. 由于这种分发方式长此以往被证明是不切实际且不充分的，今后所有的``Bois-Marie 代数几何讨论班''都将像本卷一样以书籍形式出版.} 分发. 我们只限于在原始文本中加入一些脚注，订正若干排印错误，并作了术语上的调整，其中``简单态射''(simple morphism) 一词已被``光滑态射''(smooth morphism) 所取代，后者不会引起同样的混淆.

讲 I 至 IV 呈现平展态射 (etale morphism) 与光滑态射的局部概念；它们几乎不用《原理》(Éléments) 第一章所阐述的概形语言\footnote{如今在《原理》第四章 §§17 和 18 中可以找到更完整的研究.}. 讲 V 呈现概形的基本群的公理化描述，即使在经典情形 (该概形化为一个域的谱 (spectrum)) 下也是有用的，在那里人们找到通常的 Galois 理论的一个非常便利的重述. 讲 VI 和 VIII 呈现下降 (descent) 理论，近年来下降理论在代数几何中获得了日益增长的重要性，并且它可能为解析几何 (Géométrie Analytique) 和拓扑学 (Topologie) 提供类似的服务. 应当指出，讲 VII 当时未曾写出，其内容已被并入 J. Giraud 的一篇著作 (Méthode de la Descente，Bull. Soc. Math. France，Mémoire 2，1964，viii + 150 p.). 在讲 IX 中，我们更专门地研究平展态射的下降，为基本群的 Van Kampen 型定理得到一种系统的方法，这些定理在这里表现为下降定理的简单翻译. 本质上，这是计算连通概形 $X$ 的基本群的一种方法，它给定了满射且紧合的 (proper) 态射，比方说 $X' \to X$，用 $X'$ 的连通分支的基本群、纤维积 $X' \times_X X'$，$X' \times_X X' \times_X X'$，以及上述概形之间的典范单纯 (simplicial) 态射在这些群之间诱导的同态来表达. 讲 X 给出基本群的特殊化 (specialization) 理论，针对紧合且光滑的态射，其最引人注目的结果在于 (几乎完全地) 确定特征 $p > 0$ 中光滑代数曲线的基本群，这得益于特征零中已知的超验方法的结果. 讲 XI 给出一些例子和补充，特别以同调的形式阐明了 Kummer 覆盖与 Artin-Schreier 覆盖的理论. 关于文本的其他评注，参见多印本版中的《说明》，它紧接本导论之后.

自从 1961 年本讨论班成文以来，M. Artin 和我本人合作发展了平展拓扑 (etale topology) 的语言以及相应的上同调理论，该理论在代数几何讨论班 SGA 4《概形的平展上同调》部分中阐述，即将在与本卷相同的丛书中出版. 这种语言以及它目前所拥有的成果，为基本群的研究提供了一个特别灵活的工具，使人们能够更好地理解并超越此处阐述的某些结果. 因此理应从这个观点出发彻底重做基本群理论 (事实上，所有关键结果目前都已出现在上引文献 (loc. cit.) 中). 这正是《原理》中拟用于基本群的那一章所计划的，它本应还包含若干无法在此占有一席之地的其他发展 (它们依据奇点解消 (resolution of singularities) 技术)：用该环的一个适当奇点解消来计算完备局部环的``局部基本群'' (local fundamental group)，无真性 (properness) 假设下基本群的局部与整体 Künneth 公式 (参见讲 XIII)，M. Artin 关于优的 (excellent) Hensel 局部环与其完备化的局部基本群之比较的结果 (SGA 4 XIX). 还应当指出，有必要发展一个意象 (topos) 的基本群理论，它将同时涵盖通常的拓扑理论、其半单纯 (semi-simplicial) 版本、本卷讲 V 中发展的``投射有限'' (profinite) 变体，以及 SGA 3 X 7 中稍更一般的投射离散 (pro-discrete) 变体 (它适用于非正规且非单支的概形的情形). 在从这一角度对理论进行全面改造之前，Raynaud 女士的讲 XIII 利用 SGA 4 的语言和结果，旨在表明在若干典型问题中可以得到的益处，特别是把讲 X 的某些结果推广到非紧合的相对概形. 其中特别给出了特征任意的非完备代数曲线的``素于 $p$''(prime to p) 基本群的结构 (这是我 1959 年宣布过的，但其证明至今未发表).

尽管有如此多的缺陷和不完善之处 (另一些人会说：正因为这些缺陷和不完善之处)，我认为本卷对于那些希望熟悉基本群理论的读者，以及作为一部参考著作，都是有用的，直至一部免于我刚列举的各种批评的文本写出并问世.

\section*{2.}

本卷构成``Bois-Marie 代数几何讨论班''的第 1 卷，其后续各卷预计在与本卷相同的丛书中出版. 讨论班的目标，与 J. Dieudonné 和 A. Grothendieck 的论著《代数几何原理》并驾齐驱，是按照后一部著作中的观点为代数几何奠定基础. 讨论班所有各卷的标准参考文献由《代数几何原理》第一、二、三章构成 (引用为 EGA I，II，III)，并假定读者已具备这些章所蕴含的交换代数 (commutative algebra) 与同调代数 (homological algebra) 的装备\footnote{关于这一点，参见 EGA I 的导论.}. 此外，讨论班的每一卷都将根据需要自由引用同一讨论班先前的各卷，或《原理》中已出版或即将出版的其他各章.

讨论班的每一部分都围绕一个主题，该主题在相应一卷或多卷的标题中给出；口头讨论班通常持续一个学年，有时更长. 讨论班每一部分内部的各讲之间通常具有密切的逻辑依存关系；相反，讨论班的不同部分在很大程度上彼此逻辑独立. 这样，``群概形'' (Schémas en Groupes) 部分几乎完全独立于在时间上先于它的讨论班的两个部分；相反，它频繁引用 EGA IV 的结果. 下面是讨论班即将出版的各部分的清单 (下文引用为 SGA 1 至 SGA 7)：

\begin{itemize}\tightlist
\item SGA 1. 平展覆盖与基本群 (Révêtements étales et groupe fondamental)，1960 和 1961.
\item SGA 2. 凝聚层的局部上同调与局部和整体 Lefschetz 定理 (Cohomologie locale des faisceaux cohérents et théorèmes de Lefschetz locaux et globaux)，1961/62.
\item SGA 3. 群概形 (Schémas en groupes)，1963 和 1964 (3 卷)，与 M. Demazure 合作.
\item SGA 4. 意象理论与概形的平展上同调 (Théorie des topos et cohomologie étale des schémas)，1963/64 (3 卷)，与 M. Artin 和 J.L. Verdier 合作.
\item SGA 5. $\ell$-adic 上同调与 $L$ 函数 (Cohomologie $\ell$-adique et fonctions $L$)，1964 和 1965 (2 卷).
\item SGA 6. 相交理论与 Riemann-Roch 定理 (Théorie des intersections et théorème de Riemann-Roch)，1966/67 (2 卷)，与 P. Berthelot 和 L. Illusie 合作.
\item SGA 7. 代数几何中的局部单值性群 (Groupes de monodromie locale en géométrie algébrique).
\end{itemize}

这些局部讨论班中有三个是与他人合作主持的，他们将作为合著者出现在相应各卷的封面上. 至于讨论班的其他活跃参与者，其作用 (无论在写作方面还是在数学定稿工作方面) 逐年增长，每位参与者的名字都列在他以演讲人或撰稿人身份负责的各讲之首，而某一特定卷中参与者的名单则在该卷的扉页上标明.

应当对讨论班与《原理》之间的关系作一些说明. 《原理》原则上旨在全面阐述在代数几何中被认为最基础的概念和技术，随着这些概念和技术本身在逻辑一致性与审美要求的自然作用之下显现出来. 从这个角度看，把讨论班视为《原理》的初版是自然的，它迟早几乎完全会被后者吸收. 这一进程在几年前就已经在一定程度上开始了，因为本卷 SGA 1 的讲 I 至 IV 已完全被 EGA IV 吸收，而讲 VI 至 VIII 本应在几年内被 EGA VI 吸收. 然而，随着《原理》和讨论班中建设的推进，以及整体规模日益明确，最初的宗旨 (按照它，讨论班不过是一个初步而临时的版本) 显得越来越不现实，原因 (除其他外) 在于深谋远虑的大自然为人的寿命设定的限度. 考虑到讨论班各部分的写作通常都很精心，也许只有在后来的写作进展能够在付出相当深刻的改动代价之下带来非常实质性的改进时，才有必要在《原理》(或接替它的论著) 中重写这样的部分. 当前的情形对 SGA 1 讨论班就是如此，如上所述，对 SGA 2 亦然 (得益于 Raynaud 女士近期的结果). 相反，目前没有任何迹象表明，上述 SGA 3 至 SGA 7 中任何一个部分在不久的将来会如此.

``Bois-Marie 代数几何讨论班''内部的参考文献是这样给出的. 对讨论班 SGA 1 至 SGA 7 某一部分的内部引用采用 III 9.7 的格式，其中数字 III 表示讲的编号，它出现在该讲的每一页的页首，而 9.7 是该讲内部陈述 (或定义、注记等) 的编号. 必要时可以使用更长的十进制数字，例如用 9.7.1，9.7.2 来指称命题 9.7 证明中的各个步骤. 引用 III 9 表示讲 III 的第 9 节. 对某一讲内部的引用省略讲的编号. 对讨论班其他部分的引用使用同样的记号，但前面冠以 SGA 相应部分的名称，如 SGA 1 III 9.7. 同样，引用 EGA IV 11.5.7 意指：《代数几何原理》第四章，陈述 (或定义等) 11.5.7；这里第一个阿拉伯数字仍表示节的编号. 除了在全部 SGA 中通行的这些约定外，与某一讲相关的文献 (bibliography) 通常汇集在该讲的末尾，讲内部按照惯例用方括号中的数字如 [3] 来引用.

最后，为读者方便起见，只要看起来有必要，我们将在 SGA 各卷的末尾附上一个记号索引 (index of notations) 和一个术语索引 (terminological index)，后者在适当处包含所用法语术语的英文翻译.

我要在本导论之后附上一段数学之外的评论. 1969 年 11 月，我得知高等科学研究院 (IHES) -- 我自其成立以来基本上是那里的教授 -- 三年来一直接受国防部 (Ministère des Armées) 的资助. 早在我作为初出茅庐的研究者时，我就发现大多数科学家愿意以这种或那种形式与军事机构合作，他们对此几乎毫无顾忌，这令我极为遗憾. 我当时的动机本质上属于道德性质，因此很难被认真对待. 今天，鉴于军事机构的泛滥及其所掌握的大规模毁灭手段对人类构成毁灭的危险，这些动机获得了新的力量与维度. 我已经在别处以更详细的方式说明过这些问题，它们比任何科学 (包括数学) 的进展都重要得多；例如可就此查阅 G. Edwards 在《Survivre》杂志第 1 期 (1970 年 8 月) 上的文章，它概括了我在别处就这些问题所作的更详尽的阐述. 就这样，我发现自己在一个机构里工作了三年，而它在我毫不知情的情况下参与了在我看来既不道德又危险的资助方式\footnote{不言而喻，我刚才表达的意见仅代表我个人的责任，而不代表出版本卷的 Springer 出版社的责任.}. 由于目前在 IHES 的同事中只有我持这种意见，这使我争取从 IHES 预算中取消军费资助的努力归于失败，于是我在 1970 年 9 月 30 日作出了势在必行的决定，离开了 IHES，并且只要该机构继续接受此类资助，我就暂停与该机构的一切科学合作. 我已请求 IHES 所长 Motchane 先生，使 IHES 从 1970 年 10 月 1 日起不再传播我署名的数学文本，或属于 Bois-Marie 代数几何讨论班的文本. 如上所述，本讨论班的传播将由 Julius Springer 出版社在 Lecture Notes 丛书中承担. 我在此高兴地感谢 Springer 出版社和 K. Peters 先生，感谢他们为出版本书所提供的有效而殷勤的帮助，特别是承担了新补入旧讨论班的各讲以及补全不全讨论班的各讲所需的照相胶印排版.

我还要感谢 J.P. Delale 先生，他承担了编制记号索引和术语索引这一吃力不讨好的工作.

\begin{flushright}
Massy，1970 年 8 月.
\end{flushright}

\cleardoublepage

\chapter*{说明 (Avertissement)}

已写成的每个讲都是若干连续口头发言内容的结晶. 似乎没有必要说明其具体日期.

讲 VII 在讲 VIII 中被反复引用，但它未曾由演讲人写出；在口头报告中，演讲人只限于从严格实用的观点勾勒一般范畴中下降的语言，而没有深入这种语言所引发的逻辑困难. 事实证明，对这种语言的正确阐述将超出本笔记的范围，仅就其篇幅而言也是如此. 关于下降理论的规范阐述，我请读者参阅 Jean Giraud 一篇正在准备中的文章. 在它问世之前\footnote{现已问世：J. Giraud，Méthodes de la descente，Mémoire no 2 de la Société mathématiques de France，1964.}，我认为细心的读者不难凭借自己的力量补足讲 VIII 中的``幻影引用'' (ghost references).

另有一些口头发言，位于讲 XI 之后，文本的某些地方曾提到它们，它们也未曾写出，本应构成一个讲 XII 和一个讲 XIII 的内容. 这些口头发言的前几个，在带幂零元 (nilpotent elements) 的概形与解析空间 (analytic space) 的框架内 (如 Cartan 讨论班 1960/61 所引入的)，重述了与一个在完备赋值域 (complete valued field) $k$ 上局部有限型的预概形 (prescheme) 相关联的解析空间的构造、$k$ 为复数域情形的 GAGA 型定理，以及它们在超验途径定义的基本群与本笔记所研究的基本群的比较上的应用 (比较 A. Grothendieck，Fondements de la Géométrie Algébrique，Séminaire Bourbaki no 190，第 10 页，1959 年 12 月). 后几个口头发言勾勒了本文所发展的方法向允许缓和分歧 (tame ramification) 的覆盖的研究、以及去掉有限个点后完备曲线的基本群结构的研究的推广 (比较上引文献 (loc. cit.)，no 182，第 27 页，定理 14). 这些发言没有引入任何本质上新的思想，因此在《代数几何原理》相应各章出版之前，似乎没有必要对它们作规范的书面整理\footnote{它们已包含在本卷中 Raynaud 女士的讲 XII 里，其证明与口头讨论班中所阐述的原始证明不同 (参见导论).}.

相反，针对基本群和 Picard 群的 Lefschetz 型定理，无论从局部还是整体观点看，都已构成 1962 年一个独立讨论班的对象，它已完全写出，可供使用者查阅\footnote{凝聚层的平展上同调与局部和整体 Lefschetz 定理 (SGA 2)，已由 North Holland Pub. Cie. 出版.}. 应当指出，无论是本讨论班还是 1962 年讨论班所发展的结果，都将在预概形平展上同调中若干关键结果的发表中得到本质性的使用，这些关键结果将构成一个 (由 M. Artin 和我本人主持的) 1963/64 年讨论班的对象，目前正在准备之中\footnote{概形的平展上同调 (引用为 SGA 4)，即将在同一丛书中出版.}.

讲 I 至 IV 本质上是局部性的、非常初等的，将被《代数几何原理》第四章完全吸收，该章的第一部分正在付印，大概将于 64 年底出版. 尽管如此，它们对于一位希望在深入系统论著的奥秘之前了解光滑、平展或平坦态射基本性质的读者仍是有用的. 至于其他各讲，它们将被《原理》第八章吸收\footnote{事实上，由于《原理》最初计划的改变，基本群的研究被推迟到比刚才所指出的章更靠后的某一章. 比较位于本《说明》之前的导论.}，而该章在数年之内很难考虑出版.

\begin{flushright}
Bures，1963 年 6 月.
\end{flushright}

\cleardoublepage


\mainmatter
\begin{refsection}% =====================================================================
% 第一讲 (Exposé I) 平展态射 (Morphismes étales)
% A. Grothendieck
% 原书页码: pp. 1-24
% =====================================================================
\chapter{平展态射}

\section{微分演算的概念}

设 $X$ 是 $Y$ 上的\emph{预概形} (prescheme)\index[term]{yugaixing@预概形}，设 $\Delta_{X/Y}$ 或 $\Delta$ 是对角态射 (diagonal morphism) $X \to X \times_Y X$. 它是\emph{浸入} (immersion) 态射，因而是从 $X$ 到 $X \times_Y X$ 的某个开集 $V$ 中的闭浸入态射. 设 $\mathcal{I}_X$ 是 $V$ 中对应于该对角线的闭子预概形的理想 (N.B. 若想作内蕴的处理，而不假设 $X$ 在 $Y$ 上分离 (separated) —— 这个假设会是荒谬的 (canularesque) —— 就应当考虑 $\mathcal{O}_{X \times X}$ 在 $X$ 中的集合论原像，并用 $\mathcal{I}_X$ 表示后者中的增广理想 (idéal d'augmentation)...). \emph{层} (sheaf) $\mathcal{I}_X / \mathcal{I}_X^2$ 可以看作 $X$ 上的\emph{拟凝聚层} (quasi-coherent sheaf)，记为 $\varOmega_{X/Y}^{1}$. 若 $X \to Y$ 是有限型 (of finite type) 的，则它是有限型的. 它对基扩张 (base extension) $Y' \to Y$ 表现良好. 还引入层 $\mathcal{O}_{X \times_Y X} / \mathcal{I}_X^{n+1} = \mathcal{P}_{X/Y}^{n}$，它们是 $X$ 上的环层，使 $X$ 成为一个预概形，可记为 $\Delta_{X/Y}^{n}$ 并称为 $X/Y$ 的第 $n$ 个\emph{无穷小邻域} (infinitesimal neighborhood)\index[term]{wuqiongxiaolinyu@无穷小邻域}. 关于它的这套推论 (sorites) 完全平凡，尽管相当长\footnote{参见 EGA IV 16.3.}；谨慎起见，最好还是等到我们对\emph{光滑态射} (smooth morphism) 有可用的结论时再谈它.

\section{拟有限态射}

\begin{proposition}\label{prop:I.2.1}
设 $A \to B$ 是局部同态 (local homomorphism) (N.B. 现在环都是诺特 (noetherian) 的)，$m$ 是 $A$ 的极大理想 (maximal ideal). 下列条件等价：
\begin{enumerate}[(i)]
\item $B/mB$ 在 $k = A/m$ 上是有限维的；
\item $mB$ 是定义理想 (ideal of definition)，且 $B/\mathfrak{r}(B) = k(B)$ 是 $k = k(A)$ 的扩张；
\item $B$ 的完备化 (completion) $\widehat{B}$ 在 $A$ 的完备化 $\widehat{A}$ 上有限.
\end{enumerate}
\end{proposition}

这时称 $B$ 在 $A$ 上拟有限 (quasi-finite). 态射 $f \colon X \to Y$ 称为在 $x$ 处拟有限 (或 $Y$-预概形\index[term]{niyouxiantaishe@拟有限态射} $f$ 称为在 $x$ 处拟有限)，若 $\mathcal{O}_x$ 在 $\mathcal{O}_{f(x)}$ 上拟有限. 这也等价于说 $x$ 在其纤维 (fiber) $f^{-1}(x)$ 中是孤立的 (isolated). 态射称为拟有限的，若它在每一点处都是拟有限的\footnote{在 EGA II 6.2.3 中还进一步假设 $f$ 是有限型的.}.

\begin{corollary}\label{cor:I.2.2}
若 $A$ 完备，则拟有限等价于有限.
\end{corollary}

我们本可以给出\emph{拟有限态射}的通常那套推论 (i)，(ii)，(iii)，(iv)，(v)，但在这里似乎没有必要.

\section{非分歧态射}

\begin{proposition}\label{prop:I.3.1}
设 $f \colon X \to Y$ 是有限型态射，$x \in X$，$y = f(x)$. 下列条件等价：
\begin{enumerate}[(i)]
\item $\mathcal{O}_x / m_y \mathcal{O}_x$ 是 $k(y)$ 的有限可分扩张 (finite separable extension)；
\item $\varOmega_{X/Y}^{1}$ 在 $x$ 处为零；
\item 对角态射 $\Delta_{X/Y}$ 在 $x$ 的邻域中是开浸入.
\end{enumerate}
\end{proposition}

\begin{proof}
对于蕴含关系 (i)$\Rightarrow$(ii)，由 Nakayama 引理立即归结到 $Y = \mathrm{Spec}(k)$，$X = \mathrm{Spec}(k')$ 的情形，在那里这是众所周知的，而且从可分的定义看甚至是平凡的；(ii)$\Rightarrow$(iii) 依据一个用 Krull 定理的开浸入的巧妙而容易的刻画；(iii)$\Rightarrow$(i) 是因为我们仍归结到 $Y = \mathrm{Spec}(k)$ 且对角态射处处是开浸入的情形. 于是必须证明 $X$ 是 $k$ 上的有限可分环代数，为此归结到 $k$ 代数闭 (algebraically closed) 的情形. 但此时 $X$ 的每个闭点都是孤立的 (因为它等于由 $x$ 定义的态射 $X \to X \times_k X$ 之下对角线的原像)，故 $X$ 有限. 然后可以假设 $X$ 化为一个点，其环为 $A$，从而 $A \otimes_k A \to A$ 是同构，故 $A = k$，证毕.
\end{proof}

\begin{definition}\label{def:I.3.2}
\begin{enumerate}[a)]
\item 这时称 $f$ 在 $x$ 处是非分歧的 (unramified)\index[term]{feifenqitaishe@非分歧态射} (unramified)，或称 $X$ 在 $Y$ 上于 $x$ 处非分歧，或非分歧.
\item 设 $A \to B$ 是局部同态，称它是非分歧的，或称 $B$ 是 $A$ 上的非分歧局部代数\index[term]{feifenqidaishu@非分歧代数} (net local algebra)，若 $B/\mathfrak{r}(A)B$ 是 $A/\mathfrak{r}(A)$ 的有限可分扩张，即若 $\mathfrak{r}(A)B = \mathfrak{r}(B)$ 且 $k(B)$ 是 $k(A)$ 的可分扩张\footnote{参见 \ref{rem:III.1.2} 中的更正 (retraction).}.
\end{enumerate}
\end{definition}

\begin{remarkstar}
$B$ 在 $A$ 上非分歧这一事实从 $A$ 与 $B$ 的完备化上即可看出. 非分歧蕴含拟有限.
\end{remarkstar}

\begin{corollary}\label{cor:I.3.3}
$f$ 为非分歧的那些点构成的集合是开的.
\end{corollary}

\begin{corollary}\label{cor:I.3.4}
设 $X'$，$X$ 是 $Y$ 上的两个有限型预概形，$g \colon X' \to X$ 是 $Y$-态射. 若 $X$ 在 $Y$ 上非分歧，则图像态射 (graph morphism) $\Gamma_g \colon X' \to X \times_Y X$ 是开浸入.
\end{corollary}

\begin{proof}
事实上，它就是对角态射 $X \to X \times_Y X$ 在
\[
g \times_Y \mathrm{id}_{X'} \colon X' \times_Y X \longrightarrow X \times_Y X
\]
之下的原像.
\end{proof}

还可以引入 $\varOmega_{X/Y}^{1}$ 的零化理想 $\mathfrak{d}_{X/Y}$，称为 $X/Y$ 的差理想 (different ideal)\index[term]{panbielixiang@判别理想} 或微分歧理想；它定义 $X$ 的一个闭子预概形，从集合观点看，它就是 $X/Y$ 分歧 (ramified)，即非非分歧的那些点构成的集合.

\begin{proposition}\label{prop:I.3.5}
\begin{enumerate}[(i)]
\item 浸入是非分歧的.\index[term]{jinru@浸入}
\item 两个\emph{非分歧态射}的复合是非分歧的.
\item 非分歧态射的基扩张 (base extension) 是非分歧的.
\end{enumerate}
\end{proposition}

\begin{proof}
可以从 (ii) 或 (iii) 直接看出 (第二个在我看来更有趣). 当然也可以把它们细化成逐点的陈述；这只是表面上更一般 (除定义 b) 的框架外)，而且令人厌烦.
\end{proof}

照例得到推论：

\begin{corollary}\label{cor:I.3.6}
\begin{enumerate}[(iv)]
\item 两个非分歧态射的笛卡儿积 (Cartesian product) 是非分歧的.
\item 若 $g \circ f$ 非分歧，则 $f$ 非分歧.
\item 若 $f$ 非分歧，则 $f_{\mathrm{r\acute{e}d}}$ 非分歧.
\end{enumerate}
\end{corollary}

\begin{proposition}\label{prop:I.3.7}
设 $A \to B$ 是局部同态，假设剩余扩张 (residue field extension) $k(B)/k(A)$ 平凡，或 $k(A)$ 代数闭. 为使 $B/A$ 非分歧，必要且充分条件是 $\widehat{B}$ (作为 $A$-代数) 是 $\widehat{A}$ 的商.
\end{proposition}

\begin{remarkstar}
在不假设剩余扩张平凡的情形，可以对 $A$ 作一个适当的有限平坦扩张，以破坏所说的扩张，从而归结到这一情形.

举例：设 $A$ 是曲线的一个普通二重点 (ordinary double point) 的局部环，$B$ 是其正规化上一个点的局部环：那么 $A \subset B$，$B$ 在 $A$ 上非分歧且剩余扩张平凡，而 $\widehat{A} \to \widehat{B}$ 是满射但非单射. 因此我们要强化非分歧性的概念.
\end{remarkstar}

\section{平展态射。平展覆盖}

我们将承认关于平坦 (flat) 态射所需的一切；这些事实将在稍后予以证明，如果有必要的话\footnote{这里原书注 (4) 缺失.}.

\begin{definition}\label{def:I.4.1}
\begin{enumerate}[a)]
\item 设 $f \colon X \to Y$ 是有限型态射. 称 $f$ 在 $x$ 处平展 (etale)，若 $f$ 在 $x$ 处平坦且在 $x$ 处非分歧. 称 $f$ 平展，若它在所有点处平展. 称 $X$ 在 $Y$ 上于 $x$ 处平展，或它是 $x$ 处的平展 $Y$-预概形，等等.
\item 设 $f \colon A \to B$ 是局部同态. 称 $f$ 平展，或称 $B$ 在 $A$ 上平展，若 $B$ 在 $A$ 上平坦且非分歧\footnote{原书注 (5) 缺失.}.
\end{enumerate}
\end{definition}

\begin{proposition}\label{prop:I.4.2}
为使 $B/A$ 平展，必要且充分条件是 $\widehat{B}/\widehat{A}$ 平展.
\end{proposition}

\begin{proof}
事实上，对``非分歧''和``平坦''分别成立即可.
\end{proof}

\begin{corollary}\label{cor:I.4.3}
设 $f \colon X \to Y$ 是有限型的，$x \in X$. $f$ 在 $x$ 处平展这一事实只依赖于局部同态 $\mathcal{O}_{f(x)} \to \mathcal{O}_x$，甚至只依赖于对完备化所作的相应同态.
\end{corollary}

\begin{corollary}\label{cor:I.4.4}
假设剩余扩张 $k(A) \to k(B)$ 平凡，或 $k(A)$ 代数闭. 则 $B/A$ 平展当且仅当 $\widehat{A} \to \widehat{B}$ 是同构.
\end{corollary}

\begin{proof}
把平坦性与 \ref{prop:I.3.7} 结合起来即得.
\end{proof}

\begin{proposition}\label{prop:I.4.5}
设 $f \colon X \to Y$ 是有限型态射. 则 $f$ 为平展的那些点构成的集合是开的.
\end{proposition}

\begin{proof}
事实上，对``非分歧''和``平坦''分别成立即可.
\end{proof}

本命题说明，在研究处处平展的有限型态射时，可以放弃``逐点''的陈述.

\begin{proposition}\label{prop:I.4.6}
\begin{enumerate}[(i)]
\item 开浸入是平展的.
\item 两个平展态射的复合是平展的.
\item 基扩张 (base extension).
\end{enumerate}
\end{proposition}

\begin{proof}
事实上，(i) 是平凡的，而对 (ii) 和 (iii)，只需注意这对``非分歧''和``平坦''成立. 老实说，对局部同态 (无有限性条件) 也有相应的陈述，它们在各种情形都应当进入大全 (multiplodoque)，首先是非分歧的情形.
\end{proof}

\begin{corollary}\label{cor:I.4.7}
两个平展态射的笛卡儿积也是平展的.
\end{corollary}

\begin{corollary}\label{cor:I.4.8}
设 $X$ 和 $X'$ 在 $Y$ 上有限型，$g \colon X \to X'$ 是 $Y$-态射. 若 $X'$ 在 $Y$ 上非分歧且 $X$ 在 $Y$ 上平展，则 $g$ 平展.
\end{corollary}

\begin{proof}
事实上，$g$ 是图像态射 $\Gamma_g \colon X \to X \times_Y X'$ (由 \ref{cor:I.3.4} 它是开浸入) 与投影态射的复合，而后者平展，因为它由平展态射 $X \to Y$ 经``基扩张'' $X' \to Y$ 得到.
\end{proof}

\begin{definition}\label{def:I.4.9}
称 $Y$ 的平展覆盖 (etale covering) (相应地：非分歧覆盖) 为 $Y$ 上的一个 $Y$-概形 $X$，它在 $Y$ 上有限且平展 (相应地：非分歧).
\end{definition}

第一个条件意味着 $X$ 由 $Y$ 上的一个凝聚 (coherent) 代数层 $\mathcal{B}$ 定义. 第二个条件则意味着 $\mathcal{B}$ 在 $Y$ 上局部自由 (locally free) (相应地：什么也不意味着)，并且此外对每个 $y \in Y$，纤维 (fiber) $\mathcal{B}(y) = \mathcal{B}_y \otimes_{\mathcal{O}_y} k(y)$ 是 $k(y)$ 上的可分代数 (= 有限个有限可分扩张的直积).

\begin{proposition}\label{prop:I.4.10}
设 $X$ 是 $Y$ 的 $n$ 次平坦覆盖 (这个术语的定义本应放在 \ref{def:I.4.9} 中)，由凝聚局部自由代数层 $\mathcal{B}$ 定义. 以众所周知的方式定义迹同态 (trace homomorphism) $\mathcal{B} \to \mathcal{A}$ (它是 $\mathcal{A}$-模同态，其中 $\mathcal{A} = \mathcal{O}_Y$). 为使 $X$ 平展，必要且充分条件是相应的双线性型 (bilinear form) $\mathrm{tr}_{\mathcal{B}/\mathcal{A}}\, xy$ 定义 $\mathcal{B}$ 到其对偶 $\check{\mathcal{B}}$ 的同构，或等价地，判别式 (discriminant) 截面
\[
d_{X/Y} = d_{\mathcal{B}/\mathcal{A}} \in \Gamma\bigl(Y, \textstyle{\bigwedge^{n}} \check{\mathcal{B}} \otimes_{\mathcal{A}} \textstyle{\bigwedge^{n}} \check{\mathcal{B}}\bigr)
\]
可逆，或最后，由该截面定义的\emph{判别理想} (discriminant ideal) 是单位理想.
\end{proposition}

\begin{proof}
事实上，归结到 $Y = \mathrm{Spec}(k)$ 的情形，此时这是众所周知的可分性判则，通过 $k$ 的代数闭包 (algebraic closure) 即可平凡地看出.
\end{proof}

\begin{remarkstar}
下面，当不预先假设 $X$ 在 $Y$ 上平坦而只作正规性 (normality) 假设时，我们有一个不那么平凡的陈述.
\end{remarkstar}

\section{平展态射的基本性质}

\begin{theorem}\label{thm:I.5.1}
设 $f \colon X \to Y$ 是有限型态射. 为使 $f$ 是开浸入，必要且充分条件是它是平展且根式 (radicial) 的态射.
\end{theorem}

\begin{proof}
回忆根式意为：单射，且剩余扩张是根式的 (还可以回忆这等价于：该态射对任何基扩张仍是单射). 必要性是平凡的，剩下充分性. 我们给出两个不同的证明，第一个更短，第二个更初等.
\end{proof}

\begin{enumerate}[1)]
\item 平坦态射是开的，因此可以假设 (把 $Y$ 换成 $f(X)$) $f$ 是满的同胚 (homeomorphism). 对任何基扩张，$f$ 仍将平坦、根式、满射，因而仍是同胚，从而是闭的. 故 $f$ \emph{紧合} (proper). 因此 $f$ 有限 (参考文献：Chevalley 定理)，由凝聚代数层 $\mathcal{B}$ 定义. $\mathcal{B}$ 局部自由，且由假设它在处处秩 (rank) 为 1，故 $X = Y$，证毕.
\item 可以假设 $Y$ 和 $X$ 仿射 (affine). 我们更容易归结到证明：若 $Y = \mathrm{Spec}(A)$，$A$ 局部，且 $f^{-1}(y)$ 非空 ($y$ 是 $Y$ 的闭点)，则 $X = Y$ (事实上，这将蕴含每个 $y \in f(X)$ 有一个开邻域 $U$ 使得 $X|U = U$). 设 $X = \mathrm{Spec}(B)$，我们要证明 $A = B$. 为此归结到把 $A$ 换成 $\widehat{A}$，$B$ 换成 $B \otimes_A \widehat{A}$ 的相应断言 (考虑到 $\widehat{A}$ 在 $A$ 上忠实平坦 (faithfully flat)). 故可假设 $A$ 完备. 设 $x$ 是 $y$ 上的点，由推论 \ref{cor:I.2.2}，$\mathcal{O}_x$ 在 $A$ 上有限，从而 (在 $A$ 上平坦且根式) 与 $A$ 相同. 于是 $X = Y \amalg X'$ (不交并). 由于 $X$ 在 $Y$ 上根式，$X'$ 为空. 证毕.
\end{enumerate}

\begin{corollary}\label{cor:I.5.2}
设 $f \colon X \to Y$ 是闭浸入且平展的态射. 若 $X$ 连通，则 $f$ 是 $X$ 到 $Y$ 的一个连通分支上的同构.
\end{corollary}

\begin{proof}
事实上，$f$ 也是开浸入. 由此得出：
\end{proof}

\begin{corollary}\label{cor:I.5.3}
设 $X$ 是 $Y$ 上的非分歧概形，$Y$ 连通. 则 $X$ 在 $Y$ 上的每个截面 (section) 都是 $Y$ 到 $X$ 的一个连通分支上的同构. 因此在所有这些截面的集合与 $X$ 的满足``投影 $X_i \to Y$ 是同构''(或由 \ref{thm:I.5.1} 等价地：满射且根式) 的连通分支 $X_i$ 的集合之间存在一一对应. 特别地，一个截面由其在一个点的值即可确定.
\end{corollary}

\begin{proof}
只有第一个断言需要证明；由 \ref{cor:I.5.2}，只需注意一个截面是闭浸入 (因为 $X$ 在 $Y$ 上分离) 且由 \ref{cor:I.4.8} 它是平展的.
\end{proof}

\begin{corollary}\label{cor:I.5.4}
设 $X$ 和 $Y$ 是 $S$ 上的两个预概形，$X$ 在 $S$ 上非分歧且分离，$Y$ 连通. 设 $f$，$g$ 是 $Y$ 到 $X$ 的两个 $S$-态射，$y$ 是 $Y$ 的一个点，假设 $f(y) = g(y) = x$，且由 $f$ 和 $g$ 定义的剩余同态 $k(x) \to k(y)$ 相同 (<<$f$ 和 $g$ 在 $y$ 处几何上重合>>). 则 $f$ 和 $g$ 相同.
\end{corollary}

\begin{proof}
由 \ref{cor:I.5.3} 归结到 $Y = S$ 的情形，把 $X$ 换成 $X \times_S Y$.
\end{proof}

下面是 \ref{cor:I.5.3} 的一个特别重要的变体.

\begin{theorem}\label{thm:I.5.5}
设 $S$ 是预概形，$X$ 和 $Y$ 是两个 $S$-预概形，$S_0$ 是 $S$ 的一个与 $S$ 有相同底空间的闭子预概形，$X_0 = X \times_S S_0$ 和 $Y_0 = Y \times_S S_0$ 是 $X$ 和 $Y$ 在 $S_0$ 上的``限制''. 假设 $X$ 在 $S$ 上平展. 则自然映射
\[
\operatorname{Hom}_{S}(Y, X) \longrightarrow \operatorname{Hom}_{S_0}(Y_0, X_0)
\]
是双射.
\end{theorem}

\begin{proof}
仍归结到 $Y = S$ 的情形，于是由 \ref{cor:I.5.3} 中给出的 $X/Y$ 的截面的``拓扑''描述即得.
\end{proof}

\begin{scholie}
这个结果包含态射的存在性与唯一性断言. 它还可以表述为 (当 $X$ 和 $Y$ 都在 $S$ 上平展时)：从平展 $S$-概形的范畴到平展 $S_0$-概形的范畴的函子 $X \mapsto X_0$ 是全忠实 (fully faithful) 的，即建立起第一个范畴与第二个范畴的一个全子范畴 (full subcategory) 之间的等价. 我们将在下面看到，它甚至是第一个与第二个范畴之间的等价 (这将是一个平展 $S$-概形的存在性定理).
\end{scholie}

下面是 \ref{thm:I.5.5} 的表面上更一般的形式，它常常是方便的：

\begin{corollary}[«提升延拓定理» (Theorem of extension of liftings)]\label{cor:I.5.6}
考虑态射的交换图
\[
\begin{tikzcd}[row sep=large, column sep=large]
X & Y_0 \arrow[l] \arrow[d, hook] \\
S & Y \arrow[l] \arrow[u, dashed, "\exists!"]
\end{tikzcd}
\]
其中 $X \to S$ 平展，$Y_0 \to Y$ 是双射闭浸入. 则可以找到唯一的态射 $Y \to X$，使相应的两个三角形交换.
\end{corollary}

\begin{proof}
事实上，把 $S$ 换成 $Y$，$X$ 换成 $X \times_S Y$，归结到 $Y = S$ 的情形，此时它是 \ref{thm:I.5.5} 对 $Y = S$ 的特殊情形.
\end{proof}

再指出 \ref{thm:I.5.1} 的如下直接推论 (我们当初未把它列为推论，是为了不打断 \ref{thm:I.5.1} 之后所发展的思路)：

\begin{proposition}\label{prop:I.5.7}
设 $X$，$X'$ 是两个在 $Y$ 上有限型且平坦的预概形，$g \colon X \to X'$ 是 $Y$-态射. 为使 $g$ 是开浸入 (相应地：同构)，必要且充分条件是对每个 $y \in Y$，在纤维上诱导的态射
\[
g \otimes_Y k(y) \colon X \otimes_Y k(y) \longrightarrow X' \otimes_Y k(y)
\]
是开浸入 (相应地：同构).
\end{proposition}

\begin{proof}
只需证明充分性；由于这对满射概念成立，归结到开浸入的情形. 由 \ref{thm:I.5.1}，必须验证 $g$ 根式 (这是平凡的) 且平展 (这由下面的推论 \ref{cor:I.5.9} 得出).
\end{proof}

\begin{corollary}\label{cor:I.5.8}
(本应移到第 3 节) 设 $X$ 和 $X'$ 是两个 $Y$-预概形，$g \colon X \to X'$ 是 $Y$-态射，$x$ 是 $X$ 的点，$y$ 是它在 $Y$ 上的投影. 为使 $g$ 在 $x$ 处拟有限 (相应地：非分歧)，必要且充分条件是 $g \otimes_Y k(y)$ 亦然.
\end{corollary}

\begin{proof}
事实上，为了确认在 $x$ 处确有拟有限 (相应地：非分歧) 态射而需要考察的 $k(g(x))$ 上的两个代数，对 $g$ 和 $g \otimes_Y k(y)$ 是相同的.
\end{proof}

\begin{corollary}\label{cor:I.5.9}
沿用 \ref{cor:I.5.8} 的记号，假设 $X$ 和 $X'$ 在 $Y$ 上平坦且有限型. 为使 $g$ 在 $x$ 处平坦 (相应地：平展)，必要且充分条件是 $g \otimes_Y k(y)$ 平坦 (相应地：平展).
\end{corollary}

\begin{proof}
对于``平坦''，这个陈述只是备忘，它是平坦性的基本判则之一\footnote{参见 EGA IV 11.8.2 (原书注 (6)).}. 对于平展，结合 \ref{cor:I.5.8} 即得.
\end{proof}

\section{完备局部环的平展扩张的应用}

本节是关于形式预概形 (formal prescheme) 的结果的一个特殊情形，那些结果应当进入大全 (multiplodoque). 不过在这里我们更省事，即不需要 §7 中平展态射的显式局部刻画 (它使用主定理 (Main Theorem)). 这也许是保留本节 (即使在大全中也保留) 在此位置的充分理由.

\begin{theorem}\label{thm:I.6.1}
设 $A$ 是完备局部环 (当然是诺特的)，剩余域为 $k$. 对每个 $A$-代数 $B$，设 $R(B) = B \otimes_A k$ 看作 $k$-代数，它函子性地依赖于 $B$. 则 $R$ 建立 $A$ 上有限且平展的 $A$-代数的范畴与 $k$ 上有限阶可分代数的范畴之间的等价.
\end{theorem}

\begin{proof}
首先，所述函子全忠实，这由下述更一般的结论得出：
\end{proof}

\begin{corollary}\label{cor:I.6.2}
设 $B$，$B'$ 是两个在 $A$ 上有限的 $A$-代数. 若 $B$ 在 $A$ 上平展，则典范映射
\[
\operatorname{Hom}_{A\text{-alg}}(B, B') \longrightarrow \operatorname{Hom}_{k\text{-alg}}\bigl(R(B), R(B')\bigr)
\]
是双射.
\end{corollary}

\begin{proof}
归结到 $A$ 为 Artin 环的情形 (把 $A$ 换成 $A/\mathfrak{m}^n$)，此时它是 \ref{thm:I.5.5} 的特殊情形.
\end{proof}

剩下要证明：对每个有限可分 $k$-代数 (为什么不说：平展，这更短) $L$，存在平展于 $A$ 上的 $B$ 使得 $R(B)$ 同构于 $L$. 可以假设 $L$ 是 $k$ 的可分扩张，作为这样的扩张它有一个生成元 (generator) $x$，即同构于一个代数 $k[t]/Fk[t]$，其中 $F \in k[t]$ 是首一多项式 (monic polynomial). 把 $F$ 提升为首一多项式 $F_1 \in A[t]$，并取 $B = A[t]/F_1 A[t]$.

\section{非分歧态射与平展态射的局部构造}

\begin{proposition}\label{prop:I.7.1}
设 $A$ 是诺特环，$B$ 是 $A$ 上的有限代数，$u$ 是 $B$ 在 $A$ 上的生成元，$F \in A[t]$ 满足 $F(u) = 0$ (不假设 $F$ 首一)，$u' = F'(u)$ (其中 $F'$ 是导多项式 (derived polynomial))，$q$ 是 $B$ 的一个不含 $u'$ 的素理想，$p$ 是它在 $A$ 上的迹. 则 $B_{\mathfrak{q}}$ 在 $A_{\mathfrak{p}}$ 上非分歧.
\end{proposition}

\begin{proof}
换言之，令 $Y = \mathrm{Spec}(A)$，$X = \mathrm{Spec}(B)$，$X_{u'} = \mathrm{Spec}(B_{u'})$，则 $X_{u'}$ 在 $Y$ 上非分歧. 本陈述由下面更精确的结论得出：
\end{proof}

\begin{corollary}\label{cor:I.7.2}
$B/A$ 的差理想包含 $u'B$，并且若自然同态 $A[t]/FA[t] \to B$ (把 $t$ 映到 $u$) 是同构，则与之相等.
\end{corollary}

\begin{proof}
设 $J$ 是同态 $C = A[t] \to B$ 的核，它包含 $FA[t]$，并且在 \ref{cor:I.7.2} 所考虑的第二种情形中与之相等. 由于它是满射，$\varOmega_{B/A}^{1}$ 等同于 $\varOmega_{C/A}^{1}$ 除以由 $J\varOmega_{C/A}^{1}$ 和 $d(J)$ 生成的子模 (本应在 §1 中详细说明同态 $d$ 的定义，以及多项式代数的 $\varOmega^{1}$ 的计算). 借助基 $dt$ 把 $\varOmega_{C/A}^{1}$ 等同于 $C$，得到 $B / B \cdot J'$，故差理想由 $J$ 中诸 $G$ 的导数在 $B$ 中的像的集合 $J'$ 生成 (只需取生成 $J$ 的 $G$). 由于 $F \in J$ (相应地：$F$ 生成 $J$)，证毕. (N.B. 应当把 \ref{cor:I.7.2} 作为命题、把 \ref{prop:I.7.1} 作为推论). 得到：
\end{proof}

\begin{corollary}\label{cor:I.7.3}
在 \ref{prop:I.7.1} 的条件下，假设 $F$ 首一且 $A[t]/FA[t] \to B$ 是同构，为使 $B_{\mathfrak{q}}$ 在 $A_{\mathfrak{p}}$ 上平展，必要且充分条件是 $q$ 不含 $u'$.
\end{corollary}

\begin{proof}
事实上，由于 $B$ 在 $A$ 上平坦，平展等价于非分歧，并可应用 \ref{cor:I.7.2}.
\end{proof}

\begin{corollary}\label{cor:I.7.4}
在 \ref{cor:I.7.3} 的条件下，为使 $B$ 在 $A$ 上平展，必要且充分条件是 $u'$ 可逆，或等价地，$F$，$F'$ 在 $A[t]$ 中生成的理想是单位理想.
\end{corollary}

\begin{proof}
最后一个判则由第一个判则和 Nakayama 引理 (在 $B$ 中) 得出.
\end{proof}

具有推论 \ref{cor:I.7.4} 所述性质的首一多项式 $F \in A[t]$ 称为\emph{可分多项式} (separable polynomial)\index[term]{kefenduoxiangshi@可分多项式} (若 $F$ 非首一，至少应要求其首项系数可逆；当 $A$ 是域时，就得到通常的定义).

\begin{corollary}\label{cor:I.7.5}
设 $B$ 是局部环 $A$ 上的有限代数. 假设 $k(A)$ 无限，或 $B$ 局部. 设 $n$ 是 $L = B \otimes_A k(A)$ 在 $k(A) = k$ 上的秩 (rank). 为使 $B$ 在 $A$ 上非分歧 (相应地：平展)，必要且充分条件是 $B$ 同构于 $A[t]/FA[t]$ 的商 (相应地：同构于 $A[t]/FA[t]$)，其中 $F$ 是首一可分多项式，可假设 (相应地：必然) 次数为 $n$.
\end{corollary}

\begin{proof}
只需证明必要性. 假设 $B$ 在 $A$ 上非分歧，故 $L$ 在 $k$ 上可分，由所作假设 $L/k$ 有生成元 $\xi$，于是 $\xi^i$ ($0 \leqslant i < n$) 构成 $L$ 在 $k$ 上的基. 设 $u \in B$ 提升 $\xi$，由 Nakayama 引理，$u^i$ ($0 \leqslant i < n$) 生成 $A$-模 $B$ (相应地：构成其基)，特别可以找到首一多项式 $F \in A[t]$ 使得 $F(u) = 0$，且 $B$ 同构于 $A[t]/FA[t]$ 的商 (相应地：同构于 $A[t]/FA[t]$). 最后，由应用于 $L/k$ 的 \ref{cor:I.7.4}，$F$ 和 $F'$ 模 $\mathfrak{m}_A[t]$ 生成 $A[t]$，从而 (由 $A[t]/FA[t]$ 中的 Nakayama 引理) $F$ 和 $F'$ 生成 $A[t]$，证毕.
\end{proof}

\begin{theorem}\label{thm:I.7.6}
设 $A$ 是局部环，$A \to \mathcal{O}$ 是局部同态，使得 $\mathcal{O}$ 同构于 $A$ 上有限型代数的局部化 (localized algebra). 假设 $\mathcal{O}$ 在 $A$ 上非分歧. 则可以找到 $A$ 上的整 (integral) 代数 $B$，$B$ 的一个极大理想 $n$，$B$ 在 $A$ 上的生成元 $u$，首一多项式 $F \in A[t]$，使得 $\mathfrak{r} \not\ni F'(u)$ 且 $\mathcal{O}$ (作为 $A$-代数) 同构于 $B_{\mathfrak{n}}$. 若 $\mathcal{O}$ 在 $A$ 上平展，可取 $B = A[t]/FA[t]$.
\end{theorem}

(当然，这里的条件也是充分的...)

先指出令人愉快的推论：

\begin{corollary}\label{cor:I.7.7}
为使 $\mathcal{O}$ 在 $A$ 上非分歧，必要且充分条件是 $\mathcal{O}$ 同构于一个在 $A$ 上平展的类似代数的商.
\end{corollary}

\begin{proof}
事实上，取 $\mathcal{O}' = B'_{\mathfrak{n}'}$，其中 $B' = A[t]/FA[t]$，$\mathfrak{n}'$ 是 $n$ 在 $B'$ 中的原像.
\end{proof}

\begin{corollary}\label{cor:I.7.8}
设 $f \colon X \to Y$ 是有限型态射，$x \in X$. 为使 $f$ 在 $x$ 处非分歧，必要且充分条件是存在 $x$ 的开邻域 $U$ 使得 $f|U$ 分解为 $U \to X' \to Y$，其中第一个箭头是闭浸入，第二个是平展态射.
\end{corollary}

\begin{proof}
这是 \ref{cor:I.7.7} 的简单翻译.
\end{proof}

我们来证明 \ref{thm:I.7.6} 的行话是如何从主要陈述得出的：事实上，由 \ref{cor:I.7.7} 存在满同态 $\mathcal{O}' \to \mathcal{O}$，其中 $\mathcal{O}$ 具有所要求的性质；但由于 $\mathcal{O}'$ 和 $\mathcal{O}$ 都在 $A$ 上平展，由 \ref{cor:I.4.8} 态射 $\mathcal{O}' \to \mathcal{O}$ 平展，从而是同构.

\paragraph*{\ref{thm:I.7.6} 的证明.} 它沿用 Chevalley 讨论班的一个证明. 由主定理 (Main Theorem)，有 $\mathcal{O} = B_{\mathfrak{n}}$，其中 $B$ 是 $A$ 上的有限代数，$n$ 是它的一个极大理想. 于是 $B/\mathfrak{n} = k(\mathcal{O})$ 是 $k$ 的可分扩张，因而是单生成的；若 $\mathfrak{n}_i$ ($1 \leqslant i \leqslant r$) 是 $B$ 的异于 $n$ 的极大理想，则存在 $B$ 的元素 $u$，属于所有 $\mathfrak{n}_i$，且其在 $B/\mathfrak{n}$ 中的像是该环的生成元. 而 $B/\mathfrak{n} = B_{\mathfrak{n}}/\mathfrak{n}B_{\mathfrak{n}} = B_{\mathfrak{n}}/mB_{\mathfrak{n}}$ (其中 $m$ 是 $A$ 的极大理想). 暂且承认下述引理：

\begin{lemma}\label{lem:I.7.9}
设 $A$ 是局部环，$B$ 是 $A$ 上的有限代数，$n$ 是 $B$ 的极大理想，$u$ 是 $B$ 的元素，其在 $B_{\mathfrak{n}}/mB_{\mathfrak{n}}$ 中的像作为 $k = A/\mathfrak{m}$ 上的代数生成该环，且 $u$ 属于 $B$ 的所有异于 $n$ 的极大理想. 设 $B' = B[u]$，$\mathfrak{n}' = \mathfrak{n}B'$. 则典范同态 $B'_{\mathfrak{n}'} \to B_{\mathfrak{n}}$ 是同构.
\end{lemma}

\begin{lemma}\label{lem:I.7.10}
(本应作为 \ref{prop:I.7.1} 的推论，排在它所蕴含的 \ref{cor:I.7.5} 之前.) 设 $B$ 是 $A$ 上由元素 $u$ 生成的有限代数，$n$ 是 $B$ 的极大理想，使得 $B_{\mathfrak{n}}$ 在 $A$ 上非分歧. 则存在首一多项式 $F \in A[t]$ 使得 $F(u) = 0$ 且 $F'(u) \notin \mathfrak{n}$.
\end{lemma}

\begin{proof}
事实上，设 $n$ 是 $k$-代数 $L = B \otimes_A k$ 的秩，由 Nakayama 引理存在 $A[t]$ 中的 $n$ 次首一多项式，使得 $F(u) = 0$. 设 $f$ 是由 $F$ 模 $m$ 约化得到的多项式，则 $L$ 与 $k[t]/fk[t]$ 是 $k$-同构的，故由 \ref{cor:I.7.3}，$f'(\xi)$ 不含于 $L$ 中对应于 $n$ 的极大理想 ($\xi$ 表示 $t$ 在 $L$ 中的像，即 $u$ 在 $L$ 中的像). 由于 $f'(\xi)$ 是 $F'(u)$ 的像，证毕.
\end{proof}

定理 \ref{thm:I.7.6} 现在由 \ref{lem:I.7.9} 与 \ref{lem:I.7.10} 的结合得出. 剩下要证明 \ref{lem:I.7.9}. 令 $S' = B' - \mathfrak{n}'$，故 $B'S'^{-1} = B'_{\mathfrak{n}'}$，并有交换图
\[
\begin{tikzcd}[row sep=large, column sep=large]
B \arrow[r] & BS'^{-1} \arrow[r] & BS^{-1} = B_{\mathfrak{n}} \\
B' \arrow[u] \arrow[r] & B'S'^{-1} \arrow[u] \arrow[r,"\sim"] & B'_{\mathfrak{n}'}
\end{tikzcd}
\]
同样令 $S = B - \mathfrak{n}$，故 $BS^{-1} = B_{\mathfrak{n}}$，从而有自然同态 $BS'^{-1} \to BS^{-1} = B_{\mathfrak{n}}$，我们证明它是同构，即 $S$ 的元素在 $BS'^{-1}$ 中可逆，即后者的每个极大理想 $p$ 都不与 $S$ 相交，即在 $B$ 上诱导 $n$. 事实上，由于 $BS'^{-1}$ 在 $B'S'^{-1} = B'_{\mathfrak{n}'}$ 上有限，$p$ 诱导 $B'_{\mathfrak{n}'}$ 的惟一极大理想 $\mathfrak{n}'B'_{\mathfrak{n}'}$，因而诱导 $B'$ 的极大理想 $\mathfrak{n}'$；由于 $B$ 在 $B'$ 上有限，由 $p$ 在 $B$ 上诱导的理想 $q$ 在 $\mathfrak{n}'$ 之上，必为极大理想，且不含 $u$，故等于 $n$ (这里用到了 $u$ 属于 $B$ 的所有异于 $n$ 的极大理想). 现在证明 $BS'^{-1}$ 等于 $B'S'^{-1}$：由于它在后者上有限，由 Nakayama 引理归结为证模 $\mathfrak{n}'BS'^{-1}$ 的相等，更不妨证明模 $mBS'^{-1}$ 的相等；而 $BS'^{-1}/mBS'^{-1} = B_{\mathfrak{n}}/mB_{\mathfrak{n}}$ 由 $u$ 在 $k$ 上生成 (这里用到 $u$ 的另一性质)，故 $B'$ (从而 $B'S'^{-1}$) 在其中的像是全部 (作为包含 $k$ 和 $u$ 的像的子环).

\begin{remarkstar}
定理 \ref{thm:I.7.6} 应当可以推广到 $\mathcal{O}$ 仅为半局部 (semi-local) 环的情形，以同时涵盖 \ref{cor:I.7.5}：我们假设 $\theta/m\mathcal{O}$ 是单生成的 $k$-代数；于是可以找到 $u \in B$，其在 $B/mB$ 中的像是生成元，且属于 $B$ 的所有不由 $\mathcal{O}$ 诱导的极大理想. 引理 \ref{lem:I.7.9} 和 \ref{lem:I.7.10} 应当不难适应. 更一般地，...
\end{remarkstar}

\section{平展概形的无穷小提升。形式概形的应用}

\begin{proposition}\label{prop:I.8.1}
设 $Y$ 是预概形，$Y_0$ 是子预概形，$X_0$ 是平展 $Y_0$-概形，$x$ 是 $X_0$ 的点. 则存在平展 $Y$-概形 $X$，$x$ 在 $X_0$ 中的邻域 $U_0$，以及 $Y_0$-同构
\[
U_0 \stackrel{\sim}{\longrightarrow} X \times_Y Y_0.
\]
\end{proposition}

\begin{proof}
事实上，设 $y$ 是 $x$ 在 $Y_0$ 中的投影，把 \ref{thm:I.7.6} 应用于 $Y_0$ 和 $X_0$ 中点 $y$ 与 $x$ 的局部环的平展局部同态 $A_0 \to B_0$：得到同构
\[
B_0 = C_{0,\mathfrak{n}_0}, \qquad C_0 = A_0[t]/F_0 A_0[t],
\]
其中 $F_0$ 是首一多项式，$\mathfrak{n}_0$ 是 $C_0$ 的一个不含 $F'_0(t)$ 在 $C_0$ 中的类的极大理想. 设 $A$ 是 $y$ 在 $Y$ 中的局部环，$F$ 是 $A[t]$ 中的首一多项式，经满同态 $A \to A_0$ 给出 $F_0$ (提升 $F_0$ 的系数)，最后设 $C = A[t]/FA[t]$，$n$ 是 $C$ 中由自然满同态 $C \to C \otimes_A A_0 = C_0$ 下 $n$ 的原像构成的极大理想. 令
\[
B = C_{\mathfrak{n}}.
\]
由构造和 \ref{prop:I.7.1} 立即可知 $B$ 在 $A$ 上平展，且有同构 $B \otimes_A A_0 = A_0$ (原文如此，应为 $B_0$). 由第一章 (Chap. I) 知道，存在有限型 $Y$-概形 $X$ 和 $X$ 中位于 $y$ 之上的点 $z$，使得 $\mathcal{O}_z$ 与 $C$ 是 $A$-同构的；由于后者在 $A = \mathcal{O}_y$ 上平展，可以 (取 $X$ 足够小) 假设 $X$ 在 $Y$ 上平展. 设 $X'_0 = X \times_Y Y_0$，则 $z$ 在 $X'_0$ 中的局部环等同于 $\theta_z \otimes_A A_0 = B \otimes_A A_0$，故同构于 $B_0$. 这个同构由一个把 $x$ 在 $X_0$ 中的邻域 $U_0$ 映到 $z$ 在 $X'_0$ 中的邻域上的同构 (上引文献 (loc. cit.)) 定义，取 $X$ 足够小可以假设该邻域等于 $X'_0$. 证毕.
\end{proof}

\begin{corollary}\label{cor:I.8.2}
对平展覆盖有类似的陈述，需假设剩余域 $k(y)$ 无限.
\end{corollary}

\begin{proof}
证明相同，以 \ref{cor:I.7.5} 代替 \ref{thm:I.7.6}.
\end{proof}

\begin{theorem}\label{thm:I.8.3}
\ref{thm:I.5.5} 中所考虑的函子是范畴的等价.
\end{theorem}

\begin{proof}
由定理 \ref{thm:I.5.5}，剩下要证明每个平展 $S_0$-概形 $X_0$ 都同构于某个 $X \times_S S_0$，其中 $X$ 是平展 $S$-概形. $X$ 的底空间 (underlying topological space) 必然与 $X_0$ 的相同，且 $X_0$ 还等同于 $X$ 的闭子预概形. 因此问题等价于：在 $X_0$ 的底空间 $|X_0|$ 上找 $f_0^{*}(\mathcal{O}_S)$ 的代数层 $\mathcal{O}_X$ (其中 $f_0$ 是投影 $X_0 \to S_0$，这里看作底空间之间的连续映射)，使 $|X_0|$ 成为平展 $S$-预概形 $X$，并有代数同态 $\mathcal{O}_X \to \mathcal{O}_{X_0}$，与标量层上的同态 $f_0^{*}(\mathcal{O}_S) \to f_0^{*}(\mathcal{O}_{S_0})$ 相容，诱导同构 $\mathcal{O}_X \otimes_{f_0^{*}(\mathcal{O}_S)} f_0^{*}(\mathcal{O}_{S_0}) \stackrel{\sim}{\longrightarrow} \mathcal{O}_{X_0}$ (于是 $X$ 将是沿 $X_0$ 约化的平展 $S$-预概形，因 $X_0$ 在 $S_0$ 上分离故 $X$ 在 $S$ 上分离，且 $X$ 回答了问题). 此外，若 $(U_i)$ 是 $X_0$ 的开覆盖，且在每个 $U_i$ 上都找到了问题的解，则由唯一性定理 \ref{thm:I.5.5}，这些解可以粘合 (i.e. 定义它们的代数层连同它们的增广同态可以粘合)，并且立即看到这样在 $S$ 之上构造的环空间 (ringed space) 是带同构 $X \times_S S_0 \to X_0$ 的平展 $S$-预概形 $X$. 因此只需局部求解，这由 \ref{prop:I.8.1} 保证.
\end{proof}

\begin{corollary}\label{cor:I.8.4}
设 $S$ 是局部诺特形式预概形，带定义理想 (ideal of definition) $J$，设 $S_0 = (|S|, \mathcal{O}_S/\mathcal{J})$ 是相应的普通预概形. 则函子 $\mathfrak{X} \mapsto \mathfrak{X} \times_S S_0$ 从 $S$ 的平展覆盖的范畴到 $S_0$ 的平展覆盖的范畴是范畴的等价.
\end{corollary}

当然，形式预概形 $S$ 的平展覆盖是指 $S$ 的覆盖，i.e. $S$ 上的形式预概形，由凝聚代数层 $\mathcal{B}$ 定义，使得 $\mathcal{B}$ 局部自由，且 $\mathcal{B}$ 的剩余纤维 $\mathcal{B}_s \otimes_{\mathcal{O}_s} k(s)$ 是 $k(s)$ 上的可分代数. 若以 $S_n$ 记普通预概形 $\bigl(|S|, \mathcal{O}_S/J^{n+1}\bigr)$，则 $S$ 上凝聚代数层 $\mathcal{B}$ 的给出等价于诸 $S_n$ 上一族凝聚代数层 $\mathcal{B}_n$ 的给出，它们带有传递的同态系统 $\mathcal{B}_m \to \mathcal{B}_n$ ($m \geqslant n$)，定义同构 $\mathcal{B}_m \otimes_{\mathcal{O}_{S_m}} \mathcal{O}_{S_n} \stackrel{\sim}{\longrightarrow} \mathcal{B}_n$. 立即可知 $\mathcal{B}$ 局部自由当且仅当诸 $S_n$ 上的 $\mathcal{B}_n$ 局部自由，且可分性条件成立当且仅当它对 $\mathcal{B}_0$ 成立，或对所有的 $\mathcal{B}_n$ 成立. 这样，$\mathcal{B}$ 在 $S$ 上平展当且仅当诸 $S_n$ 上的 $\mathcal{B}_n$ 平展. 据此，\ref{cor:I.8.4} 立即由 \ref{thm:I.8.3} 得出.

\begin{remarkstar}
在 \ref{cor:I.8.4} 中没有必要局限于覆盖的情形. 不过目前只用到这一情形.
\end{remarkstar}

\section{永久性质}

设 $A \to B$ 是局部且平展的同态，这里我们考察若干``对 $A$ 成立的某种性质蕴含对 $B$ 的同一性质成立，或反之''的情形. 其中相当一部分命题只是 $B$ 在 $A$ 上拟有限且平坦这一简单事实的推论，我们只限于``回忆''其中的一些：$A$ 与 $B$ 有相同的 Krull 维数 (Krull dimension)，相同的深度 (depth) (在仍通用的术语中即 Serre 的``上同调余维数'' (cohomological codimension)). 例如由此得出 $A$ 是 Cohen-Macaulay 的当且仅当 $B$ 是. 此外，对 $B$ 的每个在 $A$ 上诱导 $p$ 的素理想 $q$，$B_{\mathfrak{q}}$ 仍在 $A_{\mathfrak{p}}$ 上拟有限且平坦，只要假设 $B$ 是 $A$ 上有限型代数的局部化 (这由``有限型态射为拟有限 (相应地：平坦) 的点集是开的''这一事实得出)；而且 $A$ 的每个素理想 $p$ 都由 $B$ 的某个素理想 $q$ 诱导 (因为 $B$ 在 $A$ 上忠实平坦). 例如由此得出 $p$ 与 $q$ 有相同的秩 (rank)；还得出 $A$ 无嵌入素理想 (embedded prime ideal) 当且仅当 $B$ 无嵌入素理想.

因此我们只限于对平展态射更特殊的情形.

\begin{proposition}\label{prop:I.9.1}
设 $A \to B$ 是局部平展同态. 为使 $A$ 正则 (regular)，必要且充分条件是 $B$ 正则.
\end{proposition}

\begin{proof}
事实上，设 $k$ 是 $A$ 的剩余域，$L$ 是 $B$ 的剩余域. 由于 $B$ 在 $A$ 上平坦且 $L = B \otimes_A k$，i.e. $n = mB$ (其中 $m$，$n$ 是 $A$，$B$ 的极大理想)，$B$ 上的 $m$-adic 滤链 (filtration) 与其 $n$-adic 滤链相同，且有
\[
\operatorname{gr}^{*}(B) = \operatorname{gr}^{*}(A) \otimes_k L.
\]
由此推出 $\mathrm{gr}^{*}(B)$ 是 $L$ 上的多项式代数当且仅当 $\mathrm{gr}^{*}(A)$ 是 $k$ 上的多项式代数. 证毕. (N.B. 没有用到 $L/k$ 可分这一事实).
\end{proof}

\begin{corollary}\label{cor:I.9.2}
设 $f \colon X \to Y$ 是平展态射. 若 $Y$ 正则\index[term]{pingzhantaishe@平展态射}，则 $X$ 正则，逆命题在 $f$ 满射时成立.
\end{corollary}

% 原书编号笔误: 下面这一条在原文中与上面的推论同为 ``\ref{cor:I.9.2}''
\begingroup
\addtocounter{theorem}{-1}%
\renewcommand{\thetheorem}{\thesection.2}%
\begin{proposition}\label{prop:I.9.2red}
设 $f \colon X \to Y$ 是平展态射. 若 $Y$ \emph{约化} (reduced)，则 $X$ 亦然，逆命题在 $f$ 满射时成立.
\end{proposition}

\endgroup

这等价于：

\begin{corollary}\label{cor:I.9.3}
设 $f \colon A \to B$ 是局部平展同态，$B$ 同构于 $A$ 上有限型代数的局部化. 为使 $A$ 约化，必要且充分条件是 $B$ 约化.
\end{corollary}

\begin{proof}
必要性是平凡的，因为 $A \to B$ 单射 ($B$ 在 $A$ 上忠实平坦). 充分性：设 $\mathfrak{p}_i$ 是 $A$ 的极小素理想，由假设自然映射 $A \to \prod_i A/\mathfrak{p}_i$ 单射，用平坦 $A$-模 $B$ 作张量，得到 $B \to \prod_i B/\mathfrak{p}_i B$ 单射，归结为证明诸 $B/\mathfrak{p}_i B$ 约化. 由于 $B/\mathfrak{p}_i B$ 在 $A/\mathfrak{p}_i$ 上平展，归结到 $A$ 整 (integral domain) 的情形. 设 $K$ 是其分式域 (field of fractions)，则 $A \to K$ 单射，故 ($B$ 为 $A$-平坦) $B \to B \otimes_A K$ 亦单射，归结为证明后一个环约化. 而 $B$ 是 $A$ 上有限型代数的局部化，它是 $Y = \mathrm{Spec}(A)$ 上有限型平展概形 $X = \mathrm{Spec}(C)$ 中某点 $x$ 的局部环，故 $B \otimes_A K$ 是 $X \otimes_A K$ 的环 $C \otimes_A K$ 的局部化 (关于某个适当的乘法稳定集 (multiplicatively closed set)). 由于 $X \otimes_A K$ 在 $K$ 上平展，其环是有限个域的乘积 ($K$ 的可分扩张)，$B \otimes_A K$ 亦然，证毕.
\end{proof}

\begin{corollary}\label{cor:I.9.4}
设 $f \colon A \to B$ 是局部平展同态，假设 $A$ 解析约化 (analytically reduced) (i.e. $A$ 的完备化 $\widehat{A}$ 无幂零元 (nilpotent element)). 则 $B$ 解析约化，从而是约化的.
\end{corollary}

\begin{proof}
事实上，$\widehat{B}$ 在 $\widehat{A}$ 上有限且平展，应用 \ref{cor:I.9.3} 即得.
\end{proof}

\begin{theorem}\label{thm:I.9.5}
设 $f \colon A \to B$ 是局部同态，$B$ 同构于 $A$ 上有限型代数的局部化. 则：
\begin{enumerate}[(i)]
\item 若 $f$ 平展，则 $A$ 正规 (normal) 当且仅当 $B$ 正规.
\item 若 $A$ 正规，则 $f$ 平展当且仅当 $f$ 单射且非分歧 (此时由 (i)，$B$ 正规).
\end{enumerate}
\end{theorem}

\begin{proof}
我们对 (i) 给出两个不同的证明，第一个用拟有限平坦态射的某些性质 (本节开头所回忆的)，不用 \ref{thm:I.7.6} (从而不用主定理)；第二个证明则相反. 最后，对 (ii) 看来无论如何都需要主定理.
\end{proof}

\paragraph*{第一个证明.} 使用下述诺特局部环 $A$ (维数 $\neq 0$) 正规的必要充分条件.

\begin{remarkstar}
\textbf{\emph{Serre 判则} (Serre criterion).}\index[term]{Serre panze@Serre 判则} (i) 对 $A$ 的每个秩 1 的素理想 $p$，$A_{\mathfrak{p}}$ 正规 (等价地：正则)；(ii) 对 $A$ 的每个秩 $\geqslant 2$ 的素理想 $p$，深度 $A_{\mathfrak{p}} \geqslant 2$\footnote{参见 EGA IV 5.8.6.}.
\end{remarkstar}

这里我们承认这个判则，它本应出现在关于平坦态射的那一节. 它的主要优点是它不预先假设 $A$ 约化，更不假设 $A$ 整. 这里已经可以假设 $\dim A = \dim B \neq 0$.

由本节开头的回忆，$A$ 中秩 1 (相应地：秩 $\geqslant 2$) 的素理想 $p$ 恰是 $B$ 中秩 1 (相应地：秩 $\geqslant 2$) 的素理想 $q$ 在 $A$ 上的迹. 最后，若 $p$ 与 $q$ 对应，则 $B_{\mathfrak{q}}$ 在 $A_{\mathfrak{p}}$ 上平展，故与 $A_{\mathfrak{p}}$ 有相同的深度，且正则当且仅当 $A_{\mathfrak{p}}$ 正则 (\cref{prop:I.9.1}). 应用 Serre 判则，得到 $A$ 正规当且仅当 $B$ 正规.

\paragraph*{第二个证明.} 假设 $B$ 正规，设 $L$ 是其分式域，$K$ 是 $A$ 的分式域 ($A$ 因 $B$ 整而整). 在 \ref{cor:I.9.3} 的证明中看到 $B \otimes_A K$ 是有限个域的直积，由于它含于 $L$ 故是域，又因它包含 $B$ 故等于 $L$. $K$ 中在 $A$ 上整的元素在 $B$ 上整，因 $B$ 正规故属于 $B$，又因 $B \cap K = A$ (由 $B$ 在 $A$ 上忠实平坦得出) 故属于 $A$.

现在假设 $A$ 正规，证明 $B$ 正规. 由 \ref{thm:I.7.6} 有 $B = B'_{\mathfrak{n}}$，其中 $B' = A[t]/FA[t]$，$F$ 和 $n$ 如 \ref{thm:I.7.6} 所述. 于是 $\mathcal{L} = \mathcal{B} \otimes_A K$ 是 $B' \otimes_A K = K[t]/FK[t]$ 的局部化，并且是域的直积；这些 $K$ 的有限可分扩张 (每当我们把 Artin 环，这里是 $B'_K$，关于一个乘法稳定集局部化时) 是 $B'_K$ 的直因子，因而对应于 $K[t]$ 中的分解 $F = F_1 F_2$，而 $L$ 中对应于 $t$ 的生成元已被 $F_1$ 零化. 再者，因 $A$ 正规，$F_i \in A[t]$ (假设它们首一). 注意 $B \to L = B \otimes_A K$ 单射 ($A \to K$ 单射且 $B$ 在 $A$ 上平坦)，故已有 $F_1(u) = 0$，其中 $u$ 是 $t$ 在 $L$ 中的类. 假设 $F$ 的次数取最小，则 $F_2 = 1$ (N.B. 有 $F'(u) = F'_1(u)F_2(u) + F_1(u)F'_2(u) = F'_1(u)F_2(u)$ 因为 $F_1(u) = 0$，故由 $F'(u) \neq 0$ 得 $F'_1(u) \neq 0$.)

因此有
\begin{equation}
L = B \otimes_A K = K[t]/FK[t] \tag{($*$)}
\end{equation}
$F$ 因而是 $K[t]$ 中的可分多项式 (但显然不必在 $A[t]$ 中). (N.B. 目前我们本质上只证明了：在 \ref{thm:I.7.6} 中可以选取 $F$ 和 $n$，使得 —— 用这里的记号 —— $B' \to B'_{\mathfrak{n}} = B$ 单射；这用到了 $A$ 的正规性；我不知道无正规性假设时是否仍然成立).

现在回忆一个众所周知的引理，取自 Serre 去年的课程：

\begin{lemma}\label{lem:I.9.6}
设 $K$ 是环，$F \in K[t]$ 是首一可分多项式，$L = K[t]/FK[t]$，$u$ 是 $t$ 在 $L$ 中的类 (从而 $F'(u)$ 是 $L$ 的可逆元). 则有以下公式 (其中 $n = \deg F$)：
\[
\mathrm{tr}_{L/K}\, u^{i}/F'(u) = 0 \quad \text{若 } 0 \leqslant i < n - 1,
\]
\[
\mathrm{tr}_{L/K}\, u^{n-1}/F'(u) = 1.
\]
\end{lemma}

\begin{corollary}\label{cor:I.9.7}
矩阵 $\bigl(u^{j} \cdot u^{i}/F'(u)\bigr)_{0 \leqslant i, j \leqslant n-1}$ 的行列式等于 $(-1)^{n(n-1)/2}$，因而在 $K$ 的任何子环 $A$ 中可逆.
\end{corollary}

\begin{corollary}\label{cor:I.9.8}
设 $A$ 是 $K$ 的子环，$V$ 是 $L$ 中由 $u^i$ ($0 \leqslant i \leqslant n-1$) 生成的 $A$-模，$V'$ 是 $L$ 中由满足``对一切 $y \in V$ (i.e. 对形如 $u^i$，$0 \leqslant i \leqslant n-1$ 的 $y$) 有 $\mathrm{tr}_{L/K}(xy) \in A$''的 $x \in L$ 构成的子 $A$-模. 则 $V'$ 是以 $u^{i}/F'(u)$ ($0 \leqslant i \leqslant n-1$) 为基的 $A$-模.
\end{corollary}

\begin{corollary}\label{cor:I.9.9}
假设 $K$ 是整正规环 $A$ 的分式域，$F$ 的系数在 $A$ 中. 则用 \ref{cor:I.9.8} 的记号，$V'$ 包含 $A$ 在 $L$ 中的正规闭包 (normal closure) $A'$，后者因而包含在 $A[u]/F'(u)$ 中，从而更包含在 $A[u][F'(u)^{-1}]$ 中.
\end{corollary}

\begin{proof}
把这个最后的推论应用于我们证明中得到的情形：由于 $F'(u)$ 在包含 $A[u]$ 的 $B$ 中可逆，$B$ 包含 $A'$. 由主定理 (或由 $B = A[u]_{\mathfrak{n}}$ 这一事实)，$B$ 是 $A'$ 的局部化代数. 由于 $A'$ 正规，$B$ 亦然.
\end{proof}

\paragraph*{(ii) 的证明.} 像前面的证明那样，证明在 \ref{thm:I.7.6} 中可以选取 $F$ 使得 ($*$) 仍成立. 惟一的先行障碍是，不再假设 $B$ 在 $A$ 上平坦，就不能再断言 $B \to L$ 单射，因此推理原则上只适用于 $B$ 在该同态下的像 $B_1$. 由此立即得出 $B_1$ 在 $A$ 上平坦 (作为 $A$ 上自由代数的局部化). 由 \ref{cor:I.4.8}，态射 $B \to B_1$ 平展，从而是同构，证明完毕.

(从写作的角度看，应当把最后两个证明的顺序对调，并把引理及其推论的形式计算放到单独的一节.)

\begin{corollary}\label{cor:I.9.10}
设 $f \colon X \to Y$ 是平展态射. 若 $Y$ 正规，则 $X$ 正规，逆命题在 $f$ 满射时成立.
\end{corollary}

\begin{corollary}\label{cor:I.9.11}
设 $f \colon X \to Y$ 是支配 (dominant) 态射，$Y$ 正规且 $X$ 连通. 若 $f$ 非分歧，则 $f$ 平展，故 $X$ 正规，从而 (连通) 不可约 (irreducible).
\end{corollary}

\begin{proof}
设 $U$ 是 $f$ 平展的点集，它是开的，只需证明它也是闭的且非空. $U$ 包含 $Y$ 的生成点 (generic point) 的原像 (因为对域上的代数，非分歧 = 平展)，故 ($X$ 支配 $Y$) 非空. 若 $x$ 属于 $U$ 的闭包，则它属于 $U$ 的某个不可约分支 $U_i$ 的闭包，因而属于 $X$ 的与 $U_i$ 相交的某个不可约分支 $X_i = \overline{U}_i$，从而支配 $Y$ (因为 $U$ 的每个分支都在 $Y$ 上平坦，因而支配 $Y$). 于是，若 $y$ 是 $x$ 在 $Y$ 上的投影，则 $\mathcal{O}_y \to \mathcal{O}_x$ 单射 (考虑到 $\mathcal{O}_y$ 整). 由于 $\mathcal{O}_y$ 正规且 $\mathcal{O}_y \to \mathcal{O}_x$ 非分歧，借助 \ref{thm:I.9.5}(ii) 即得结论.
\end{proof}

\begin{corollary}\label{cor:I.9.12}
设 $f \colon X \to Y$ 是有限型支配态射，$Y$ 正规且 $X$ 不可约. 则 $f$ 平展的点集恰为 $\varOmega_{X/Y}^{1}$ 的支集 (support) 的补集，i.e. 由差理想 $\mathfrak{d}_{X/Y}$ 定义的 $X$ 的闭子预概形的补集.
\end{corollary}

(这正是 §4 的注记中所提到的``不那么平凡的''陈述.)

\begin{remarkstar}
我们切勿以为不可约概形的连通平展覆盖本身不可约，当不假设基正规时尤其如此. 这个问题将在 §11 中研究.
\end{remarkstar}

\section{正规概形的平展覆盖}

\begin{proposition}\label{prop:I.10.1}
设 $X$ 是正规连通概形 $Y$ (其函数域为 $K$) 上的平展分离预概形. 则 $X$ 的连通分支 $X_i$ 是整的，其函数域 $K_i$ 是 $K$ 的有限可分扩张，且 $X_i$ 等同于 $X$ 在 $K_i$ 中的正规化 (normalized) 的一个非空开部分 (故 $X$ 等同于 $Y$ 在 $R(X) = L = \prod K_i$ 中的正规化的一个开稠密部分).
\end{proposition}

\begin{proof}
由 \ref{cor:I.9.10}，$X$ 正规，更不用说其局部环整，故 $X$ 的连通分支不可约. 由于 $X_i$ 正规，且在 $Y$ 上有限且支配，由主定理的一个 (事实上几乎平凡) 特殊情形得出，$X_i$ 是 $X$ 在 $K_i$ 中的正规化的一个开集.
\end{proof}

\begin{corollary}\label{cor:I.10.2}
在 \ref{prop:I.10.1} 的条件下，$X$ 在 $Y$ 上有限 (i.e. 是 $Y$ 的平展覆盖) 当且仅当 $X$ 同构于 $Y$ 在 $L = R(X)$ (即 $X$ 上的有理函数 (rational functions) 环) 中的正规化 $Y'$.
\end{corollary}

\begin{proof}
事实上，我们知道这个正规化在 $Y$ 上有限 ($Y$ 正规且 $R/K$ 可分)，反之若 $X$ 在 $Y$ 上有限则它在 $Y'$ 上有限，故其在 $Y'$ 中的像是闭的，另一方面它又是稠密的.
\end{proof}

$K$ 上的有限阶代数 $L$ 称为在 $X$ 上非分歧 (或简称在 $K$ 上非分歧，若 $X$ 隐含) 的，若 $L$ 是 $K$ 上的可分代数，i.e. 是有限个可分扩张 $K_i$ 的直积，且 $Y$ 在 $L$ 中的正规化 $Y'$ ($Y$ 在诸 $K_i$ 中的正规化的不交并) 在 $Y$ 上非分歧 ($=$ 由 \ref{cor:I.9.11} 平展). 因此：

\begin{corollary}\label{cor:I.10.3}
对每个在 $Y$ 上有限、其每个不可约分支都支配 $Y$ 的 $X$，设 $R(X)$ 是 $X$ 上的有理函数环 ($X$ 的不可约分支的生成点的局部环的乘积)，于是 $X \mapsto R(X)$ 是取值于 $K = R(Y)$ 上有限阶代数的函子. 该函子建立 $Y$ 的连通平展覆盖的范畴与 $K$ 的在 $Y$ 上非分歧的扩张 $L$ 的范畴之间的等价.
\end{corollary}

逆函子是正规化函子.

假设 $Y$ 仿射，故由分式域为 $K$ 的正规环 $A$ 定义\index[term]{zhengguigaixing@正规概形}. 设 $L$ 是 $K$ 的有限扩张，是域的直积，则按定义 $Y$ 在 $L$ 中的正规化 $Y'$ 同构于 $\mathrm{Spec}(A')$，其中 $A'$ 是 $A$ 在 $L$ 中的正规化. 说 $L$ 在 $Y$ 上非分歧，意味着 $A'$ 在 $A$ 上非分歧 (或：平展). 若 $A$ 局部，则等价于说局部环 $A'_{\mathfrak{n}}$ ($n$ 遍历 $A'$ 的有限个极大理想，即诱导 $A$ 的极大理想 $m$ 的素理想) 在局部环 $A$ 上非分歧 ($=$ 平展). 最后还注意到，判别式判则 (\cref{prop:I.4.10}) 在这一情形也可应用 (更一般地，该判则的一个变体应当这样陈述，无需预先的平坦性条件，当 $X$ 支配 $Y$，$Y$ 仍假设局部整时：$A \to B$ 和 $B \to B \otimes_A K = L$ 单射 —— 则 $\mathrm{tr}_{L/K}$ 有定义 —— 且 $\mathrm{tr}_{L/K}(xy)$ 诱导基本双线性型 $B \times B \to A$，i.e. 存在 $x_i \in B$ ($1 \leqslant i \leqslant n$，$n$ 为 $L$ 在 $K$ 上的秩) 使得对所有 $i, j$ 有 $\mathrm{tr}(x_i x_j) \in A$ 且 $\det\bigl(\mathrm{tr}(x_i x_j)\bigr)$ 在 $A$ 中可逆).

推论组 (\cref{prop:I.4.6}) 立即蕴含经典框架中非分歧性的推论组：

\begin{proposition}\label{prop:I.10.4}
设 $Y$ 是正规整预概形，函数域为 $K$. (i) $K$ 在 $Y$ 上非分歧. (ii) 若 $L$ 是 $K$ 的在 $Y$ 上非分歧的扩张，$Y'$ 是函数域为 $L$ 且支配 $Y$ 的正规预概形 (例如 $Y$ 在 $L$ 中的正规化)，$M$ 是 $L$ 的在 $Y'$ 上非分歧的扩张，则 $M/K$ 在 $X$ 上非分歧 (非分歧性的传递性 (transitivity)). (iii) 设 $Y'$ 是支配 $Y$ 的正规整预概形，函数域为 $K'/K$；若 $L$ 是 $K$ 的在 $Y$ 上非分歧的扩张，则 $L \otimes_K K'$ 是 $K'$ 的在 $Y'$ 上非分歧的扩张 (平移性质 (translation property)).
\end{proposition}

此外：

\begin{corollary}\label{cor:I.10.5}
在 (iii) 的条件下，若 $Y = \mathrm{Spec}(A)$，$Y' = \mathrm{Spec}(A')$，则 $Y'$ 在 $L' = L \otimes_K K'$ 中的正规化 $\overline{A'}$ 等同于 $\overline{A} \otimes_A A'$，其中 $\overline{A}$ 是 $A$ 在 $L$ 中的正规化.
\end{corollary}

通常人们 (那些厌恶考虑非整环、即使它是域的直积的人) 把平移性质表述为如下 (更弱) 的形式：

\begin{corollary}\label{cor:I.10.6}
在 (iii) 的条件下，设 $L_1$ 是 $L/K$ (在 $Y$ 上非分歧) 与 $K'/K$ 的复合扩张 (composite extension). 则 $L_1/K'$ 在 $Y'$ 上非分歧. 在 $Y = \mathrm{Spec}(A)$，$Y' = \mathrm{Spec}(A')$ 的情形，还成立
\[
\overline{A'} = A[\overline{A}, A']
\]
i.e. $A'$ 在 $L_1$ 中的正规化环 $\overline{A'}$ 是由 $A'$ 和 $A$ 在 $L$ 中的正规化 $\overline{A}$ 生成的 $A$-代数.
\end{corollary}

后一事实在无非分歧性假设时是假的，即使在数域 (number field) 的复合扩张的情形也是如此...

为了结束本节，我们给出与平展覆盖概念的直观图像相对应的解释：在 $Y$ 的所考虑点 $y$ 之上应当有``最大数目''的点，特别地，在 $y$ 之上不应当有``多个重合的点''. 为了以所希望的全部普遍性证明这个方向的结果，我们在这里承认下面的命题 \ref{prop:I.10.7} (其证明将在大全第四章 §15 中给出，它使用 Chevalley 的可构造集 (constructible sets) 技术\index[term]{kegoubufen@可构部分}，以及一点点下降 (descent) 理论...).

有限型态射 $f \colon X \to Y$ 称为泛开 (universally open)\index[term]{wanyoukaitaishe@万有开态射} 的，若对任何基扩张 $Y' \to Y$ (其中 $Y'$ 局部诺特)，态射 $f' \colon X' = X \times_Y Y' \to Y'$ 是开的，i.e. 把开集变成开集. 还可以只限于 $Y'$ 在 $Y$ 上有限型的情形 (甚至 $Y'$ 取形如 $Y[t_1, \dots, t_r]$，其中 $t_i$ 是未定元). 泛开态射更是开的 (逆命题不真)，另一方面若 $f$ 开，且 $X$ 与 $Y$ 不可约，则 $f$ 的所有纤维的所有分支都有相同的维数 (即生成纤维 $f^{-1}(z)$ 的维数，$z$ 为 $Y$ 的生成点). 最后，若 $Y$ 正规，则最后一个条件已经蕴含 $f$ 泛开 (Chevalley 定理). 例如由此得出：若 $f \colon X \to Y$ 是拟有限态射，$Y$ 正规不可约，则 $f$ 泛开 (或：开) 当且仅当 $X$ 的每个不可约分支都支配 $Y$. 还回忆：(有限型的) 平坦态射既然开，也就泛开. 在这些准备之后，``回忆''下述

\begin{proposition}\label{prop:I.10.7}
设 $f \colon X \to Y$ 是拟有限分离泛开态射. 对每个 $y \in Y$，设 $n(y)$ 是纤维 $f^{-1}(y)$ 的``点的几何个数''\index[term]{jihedianshu@几何点数} (geometric number of points)，等于对 $x \in f^{-1}(y)$，剩余扩张 $k(x)/k(y)$ 的可分次数 (separable degree) 之和. 则 $Y$ 上的函数 $y \mapsto n(y)$ 上半连续 (upper semicontinuous). 为使它在点 $y$ 的邻域上为常数 (i.e. 使 $n(y) = n(z_i)$，其中 $z_i$ 是 $Y$ 的包含 $y$ 的不可约分支的生成点)，必须存在 $y$ 的邻域 $U$ 使得 $X|U$ 在 $U$ 上有限\footnote{参见 EGA IV 13.1.4 (原书注 (8)).}.
\end{proposition}

\begin{corollary}\label{cor:I.10.8}
若 $y \mapsto n(y)$ 为常数且 $Y$ \emph{几何单支} (geometrically unibranched)\index[term]{jihedanzhi@几何单支}\footnote{参见 §11 的末尾 (原书注 (9)).}，则 $X$ 的不可约分支两两不交.
\end{corollary}

\begin{proposition}\label{prop:I.10.9}
设 $f \colon X \to Y$ 是平展分离态射. 用 \ref{prop:I.10.7} 的记号，函数 $y \mapsto n(y)$ 上半连续. 为使它在点 $y$ 的邻域上为常数 (i.e. 使 $n(y) = n(z_i)$，其中 $z_i$ 是 $Y$ 的包含 $y$ 的不可约分支的生成点)，必要且充分条件是存在 $y$ 的开邻域 $U$ 使得 $X|U$ 在 $U$ 上有限，i.e. 是 $U$ 的平展覆盖.
\end{proposition}

\begin{corollary}\label{cor:I.10.10}
为使平展分离态射 $f \colon X \to Y$ ($Y$ 连通) 有限 (i.e. 使 $X$ 成为 $Y$ 的平展覆盖)，必要且充分条件是 $f$ 的所有纤维有相同的点的几何个数.
\end{corollary}

在 \ref{prop:I.10.7} 及其推论中，对 $Y$ 没有正规性假设. 若作这样的假设，就得到更强的陈述 (通常被取作覆盖非分歧性的定义)：

\begin{theorem}\label{thm:I.10.11}
设 $f \colon X \to Y$ 是拟有限分离态射. 假设 $Y$ 不可约，$X$ 的每个分支都支配 $Y$，$X$ 约化 (i.e. $\mathcal{O}_X$ 无幂零元). 设 $n$ 是 $X$ 在 $Y$ 上的次数 (degree) ($X$ 的不可约分支 $X_i$ 的域 $K_i$ 在 $Y$ 的域 $K$ 上的次数之和). 设 $y$ 是 $Y$ 的正规点. 则 $X$ 在 $y$ 之上的点的几何个数 $n(y)$ 小于等于 $n$，等号成立当且仅当存在 $y$ 的开邻域 $U$ 使得 $X|U$ 是 $U$ 的平展覆盖.
\end{theorem}

\begin{proof}
``仅当''是平凡的，证明``当''. 设 $z$ 是 $Y$ 的生成点，有 $n(z) =$ (诸 $K_i/K$ 的可分次数之和) $\leqslant n$，由 \ref{prop:I.10.7} 有 $n(y) \leqslant n(z) \leqslant n$，等号蕴含 $X|U$ 对 $y$ 的某个适当邻域 $U$ 在 $U$ 上有限. 故可假设 $X$ 在 $Y$ 上有限，且 $Y$ 上的函数 $n(y')$ 为常数. 最后由 \ref{cor:I.10.8}，$X$ 是其不可约分支的不交并，为证明它在 $y$ 处非分歧，归结到 $X$ 不可约从而整的情形. 最后可以假设 $Y = \mathrm{Spec}(\mathcal{O}_y)$. 于是定理归结为下述经典陈述：
\end{proof}

\begin{corollary}\label{cor:I.10.12}
设 $A$ 是 (如往常一样诺特的) 正规局部环，函数域为 $K$，$L$ 是 $K$ 的 $n$ 次有限扩张，可分次数为 $n_s$，$B$ 是 $L$ 的包含在 $L$ 中的子环，在 $A$ 上有限，分式域为 $L$，$m$ 是 $A$ 的极大理想，$n'$ 是 $B/\mathfrak{m}B$ 在 $A/\mathfrak{m}A = k$ 上的可分次数 (= 该环的剩余扩张的可分次数之和). 有 $n' \leqslant n_s$，从而 $n' \leqslant n$. 最后一个不等式取等号当且仅当 $B$ 在 $A$ 上非分歧 ($=$ 平展).
\end{corollary}

\begin{proof}
剩下只需证明 $n' = n$ 蕴含 $B$ 在 $A$ 上平展. 回忆 $k$ 无限时的证明：只需证明 $R = B/\mathfrak{m}B$ 在 $k$ 上可分；若不然，则 (由已知引理) 存在 $R$ 的元素 $a$，其在 $k$ 上的最小多项式次数 $> n'$. 该元素来自 $B$ 的元素 $x$，后者在 $K$ 上 (作为 $L$ 的元素) 的最小多项式次数 $\leqslant n$；另一方面，由于 $A$ 正规，该最小多项式的系数在 $A$ 中，模 $m$ 约化给出 $k[t]$ 中的首一多项式 $F$，次数 $\leqslant n = n'$，使得 $F(a) = 0$，矛盾.
\end{proof}

在一般情形 ($k$ 可以有限)，回到几何语言，考虑 $Y' = \mathrm{Spec}(A[t])$，它在 $Y$ 上忠实平坦，以及 $Y'$ 在 $y$ 上的纤维 $\mathrm{Spec}(k[t])$ 的生成点 $y'$. 则 $X$ 在 $Y$ 上的 $y$ 处非分歧，当且仅当 $X' = X \times_Y Y' = \mathrm{Spec}(B[t])$ 在 $Y'$ 上的 $y'$ 处非分歧，这立即可见. 另一方面，由 $y'$ 的选取，其剩余域是 $k(t)$，因而是无限的. 由于 $y'$ 是 $Y'$ 的正规点，归结到前面的情形.

\section{一些补充}

我们已经说过，整概形的连通平展覆盖未必整. 下面是这一事实的两个例子.

\begin{enumerate}[a)]
\item 设 $C$ 是带普通二重点 $x$ 的代数曲线，$C'$ 是其正规化，$a$ 和 $b$ 是 $C'$ 中位于 $x$ 之上的两个点. 设 $C'_i$ ($i = 1, 2$) 是 $C'$ 的两份拷贝，$a_i$ 和 $b_i$ 是 $C'_i$ 中分别对应于 $a$ 和 $b$ 的点. 在和曲线 $C'_1 \amalg C'_2$ 中，把 $a_1$ 与 $b_2$ 等同，把 $a_2$ 与 $b_1$ 等同 (把具体的等同过程留给读者；它将在大全第六章中说明，而在代数闭域上的曲线的情形，已在 Serre 关于代数曲线的书中处理). 得到一个连通但可约的曲线 $C''$，它是 $C$ 的 2 次平展覆盖. 读者将验证：一般地，$C$ 的连通``Galois 型''平展覆盖 $C''$，其原像 $C'' \times_C C'$ 是 $C'$ 的平凡覆盖 (i.e. 同构于若干份 $C'$ 的拷贝的和)，是 $n$ 次``循环''的，且对每个整数 $n > 0$，可以构造 $C$ 的一个 $n$ 次连通循环平展覆盖. 在稍后将发展的基本群语言中，这意味着 $\pi_1(C)$ 除以由 $\pi_1(C') \to \pi_1(C)$ (投影诱导的同态) 的像生成的闭不变子群所得的商，同构于 $\mathbf{Z}$ 的紧化 (compactification). 更精确地说，应当能够证明 $C$ 的基本群同构于 $C'$ 的基本群与 $Z$ 的紧化的 (拓扑) 自由积 (free product). 注意，正是这类问题催生了概形的``下降理论''.
\item 设 $A$ 是整完备局部环，我们知道其正规化 $A'$ 在 $A$ 上有限 (Nagata)，故它是完备半局部环，又因它整故是局部的. 假设它所定义的剩余扩张 $L/k$ 非根式 (在相反的情形，称 $A$ 几何单支，参见下文). 例如环 $\mathbf{R}[\![s, t]\!]/(s^2 + t^2)\mathbf{R}[\![s, t]\!]$ 就是如此，其中 $\mathbf{R}$ 是实数域. 设 $k'$ 是 $k$ 的有限 伽罗瓦扩张，使 $L \otimes_k k'$ 分解；设 $B$ 是对应于剩余扩张 $k'$ 的 $A$ 上的有限平展代数 (回忆 $B$ 本质上唯一). 则 $B$ 上的 $B' = A' \otimes_A B$ 有剩余代数 $L \otimes_k k'$，它不是局部的，故 $B'$ 不是局部环，从而 (完备) 有零因子 (zero divisors). 由于它包含在 $B$ 的全分式环 (total ring of fractions) 中 (因为它在 $A'$ 上自由，故在 $A'$ 上无挠 (torsion-free)，故在 $A$ 上无挠，故包含在 $B' \otimes_A K = B'_{(K)} = A'_{(K)} \otimes_K B_{(K)} = B_{(K)}$ 中，因为 $A'_{(K)} = K$)，由此推出 $B$ 不整. 对环 $\mathbf{R}[s, t]/(s^2 + t^2)\mathbf{R}[s, t]$，取 $k'/k = \mathbf{C}/\mathbf{R}$，得到的 $B$ 是平面中两条相交直线在交点处的局部环.
\end{enumerate}

还注意到，若存在 $Y$ (整) 的连通平展覆盖 $X$ 不是不可约的，则 $X$ 的每个不可约分支都给出支配 $Y$ 的、在 $Y$ 上非分歧但不在 $Y$ 上平展的覆盖 $X'$ 的例子. 在例子 a) 的情形，由此得到 $C'$ 在 $C$ 上非分歧，但在 $a$ 和 $b$ 两点处不平展 (这也可以直接通过检查 $x$ 和 $a$ 的局部环的完备化看出：从``形式''的观点，$C'$ 在 $a$ 点等同于 $C$ 在 $x$ 点的闭子概形，即经过 $x$ 的 $C$ 的两条``分支''之一).

在 a) 和 b) 中，我们看到 \ref{thm:I.9.5} (i) 和 (ii) 的结论不成立，这直接与``$Y$ 的一个点'爆裂'成正规化中的不同点''这一事实相关 (在 b) 中，剩余扩张非根式这一事实应当从几何上作如此解释). 精确地说，称整局部环 $A$ 几何单支 (geometrically unibranched)，若其正规化只有一个极大理想，且相应的剩余扩张是根式的；整预概形的点 $y$ 称为几何单支的，若其局部环几何单支. 例子：正规点，曲线的普通尖点 (cusp)，等等. 似乎若 $Y$ 有一个非单支的点，则 $Y$ 总存在一个连通非不可约的平展覆盖；至少在 $Y$ 是完备局部环的谱的情形，这就是我们在 b) 中所证明的. 反之可以证明，若 $Y$ 的所有点都几何单支，则每个在 $Y$ 上非分歧且支配 $Y$ 的连通 $Y$-预概形都平展且不可约. 证明沿用 \ref{thm:I.9.5} 的证明，使用下述定理 \ref{thm:I.8.3} 的推广，它将在稍后用下降技术证明(...)

设 $Y' \to Y$ 是有限、根式、满射的态射 (i.e. 可以称之为``泛同胚 (universal homeomorphism)''). 考虑从 $Y$-预概形到 $Y'$-预概形的函子 $X \mapsto X \times_Y Y' = X'$. 该函子诱导平展 $Y$-概形的范畴与平展 $Y'$-概形的范畴之间的等价. 例如，可以把这一结果应用于 $Y'$ 是 $Y$ 的正规化、$Y$ 假设为单支 (且 $Y'$ 在 $Y$ 上有限，这在实践中所遇到的所有情形都成立) 的情形，或应用于``夹在'' $Y$ 与其正规化之间的 $Y''$ (它不再需要在 $Y$ 上有限) 的情形.


\end{refsection}
\begin{refsection}% =====================================================================
% 第二讲 (Exposé II) 光滑态射: 一般理论、微分性质
% (Morphismes lisses: généralités, propriétés différentielles)
% A. Grothendieck
% 原书页码: pp. 25-48
% =====================================================================
\chapter{光滑态射：一般理论、微分性质}

对讲 I 的引用以 I 表示. 回忆环都是诺特 (noetherian) 的，预概形都是局部诺特的.

\section{一般理论}

设 $Y$ 是预概形，$t_1, \dots, t_n$ 是未定元，令
\begin{equation}
Y[t_1, \dots, t_n] = Y \otimes_{\mathbf{Z}} \mathbf{Z}[t_1, \dots, t_n]. \tag{1.1}
\label{eq:II.1.1}
\end{equation}
因此 $Y[t_1, \dots, t_n]$ 是 $Y$-概形，在 $Y$ 上仿射 (affine)，由拟凝聚 (quasi-coherent) 代数层 $\mathcal{O}_Y[t_1, \dots, t_n]$ 定义. 该预概形在 $Y$ 上的一个截面的给出，等价于 $\mathcal{O}_Y$ 的 $n$ 个截面的给出 (对应于诸 $t_i$ 在相应同态下的像). 若 $Y'$ 在 $Y$ 之上，则有
\begin{equation}
Y[t_1, \dots, t_n] \times_Y Y' = Y'[t_1, \dots, t_n], \tag{1.2}
\label{eq:II.1.2}
\end{equation}
(这意味着 $Y'$ 到 $Y[t_1, \dots, t_n]$ 的一个 $Y$-态射的给出等价于 $\mathcal{O}_{Y'}$ 的 $n$ 个截面的给出)，另一方面，由 $\mathbf{Z}$ 上多项式环的类似公式，有
\begin{equation}
\bigl(Y[t_1, \dots, t_n]\bigr)[t_{n+1}, \dots, t_m] = Y[t_1, \dots, t_m]. \tag{1.3}
\label{eq:II.1.3}
\end{equation}
公式 \eqref{eq:II.1.2} 蕴含 $Y[t_1, \dots, t_n]$ 随 $Y$ 函子性地变化.

$Y[t_1, \dots, t_n]$ 在 $Y$ 上有限型且平坦 (flat).

\begin{definition}\label{def:II.1.1}
设 $f \colon X \to Y$ 是态射，使 $X$ 成为 $Y$-预概形. 称 $f$ 在 $x \in X$ 处光滑 (smooth)\index[term]{guanghuataishe@光滑态射}\footnote{旧术语：$f$ 在 $x$ 处简单 (simple)，或 $x$ 是 $f$ 的简单点. 这一术语在多种语境 (单代数 (simple algebra)，单群 (simple group)) 中会引起混淆，因而不得不放弃.}，或称 $X$ 在 $Y$ 上于 $x$ 处光滑，若存在整数 $n \geqslant 0$，$x$ 的开邻域 $U$，以及 $U$ 到 $Y[t_1, \dots, t_n]$ 中的平展 $Y$-态射. 称 $f$ (相应地：$X$) 光滑，若它在 $X$ 的所有点处光滑. 环 $A$ 上的代数 $B$ 称为在 $B$ 的素理想 $p$ 处光滑，若 $\mathrm{Spec}(B)$ 在 $\mathrm{Spec}(A)$ 上于 $p$ 点光滑；$B$ 称为在 $A$ 上光滑，若它在 $A$ 上于 $B$ 的每个素理想 $p$ 处光滑. 最后，局部环的局部同态 $A \to B$ 称为光滑的 (或称 $B$ 在 $A$ 上光滑)\footnote{此时最好按 EGA IV 18.6.1 的说法，称 $B$ 在 $A$ 上``本质上光滑'' (essentially smooth)\index[term]{benzhiguanghua@本质光滑}.}，若 $B$ 是 $A$ 上光滑的有限型代数 $B_1$ 的局部化.
\end{definition}

注意 $X$ 在 $Y$ 上的光滑性概念对 $X$ 和 $Y$ 都是局部的；若 $X$ 在 $Y$ 上光滑，则它在 $Y$ 上局部有限型.

% 原书编号笔误: 下面这条命题在原文中与上面的定义同为 ``1.1''
\begingroup
\addtocounter{theorem}{-1}%
\renewcommand{\thetheorem}{\thesection.1}%
\begin{proposition}\label{prop:II.1.1}
$X$ 中 $f$ 光滑的点 $x$ 的集合是开的.
\end{proposition}

\endgroup

由定义即得，平凡.

\begin{corollary}\label{cor:II.1.2}
若 $B$ 在 $A$ 上于 $p$ 处光滑，则它对 $B$ 的每个含于 $p$ 的素理想 $q$ 在 $A$ 上光滑.
\end{corollary}

\ref{prop:II.1.1} 还蕴含 \ref{prop:II.1.1} 中最后两个定义在其共同的存在域内一致.

\begin{proposition}\label{prop:II.1.3}
(i) 平展态射，特别是开浸入、恒等态射，是光滑的. (ii) 光滑态射的基扩张给出光滑态射. (iii) 两个光滑态射的复合是光滑的.
\end{proposition}

\begin{proof}
(i) 由定义即得，平凡，更精确地有：
\end{proof}

\begin{corollary}\label{cor:II.1.4}
平展 = 拟有限 + 光滑.
\end{corollary}

\begin{proof}
(ii) 立即由平展态射 (\ref{prop:II.1.3}) 和投影 $Y[t_1, \dots, t_n] \to Y$ (参见 \eqref{eq:II.1.2}) 的相应事实得出. 对 (iii)，由``平展''(\ref{prop:II.1.3}) 和 $Y[t_1, \dots, t_n] \to Y$ 型投影 (参见 \eqref{eq:II.1.3}) 分别成立，以及 (ii) 中所引的两个事实，形式地得出：设 $Y$ 在 $Z$ 上光滑且 $X$ 在 $Y$ 上光滑，证明 $X$ 在 $Z$ 上光滑；可以假设 $Y$ 在 $Z[t_1, \dots, t_n]$ 上平展且 $X$ 在 $Y[s_1, \dots, s_m]$ 上平展，于是第一个假设蕴含 $Y[s_1, \dots, s_m]$ 在 $Z[t_1, \dots, t_n][s_1, \dots, s_m] = Z[t_1, \dots, s_m]$ 上平展，故 $X$ 在 $Z[t_1, \dots, s_m]$ 上平展，证毕.
\end{proof}

\begin{remark}\label{rem:II.1.5}
定义 \ref{prop:II.1.1} 中的整数 $n$ 是确定好的，因为立即看出它就是 $x$ 的局部环在其纤维 $f^{-1}\bigl(f(x)\bigr)$ 中的维数. 称它为 $X$ 在 $Y$ 上的``相对维数'' (relative dimension)\index[term]{xiangduiweishu@相对维数}. 它对态射的复合具有可加性.
\end{remark}

\section{光滑性的若干判则}

\begin{theorem}\label{thm:II.2.1}
设 $f \colon X \to Y$ 是局部有限型态射，$x \in X$，$y = f(x)$. 为使 $f$ 在 $x$ 处光滑，必要且充分条件是 (a) $f$ 在 $x$ 处平坦，且 (b) $f^{-1}(y)$ 在 $k(y)$ 上于 $x$ 处光滑.
\end{theorem}

\begin{proof}
两个平坦态射的复合平坦，且 $Y[t_1, \dots, t_n] \to Y$ 是平坦态射，故光滑蕴含平坦；结合 \ref{prop:II.1.3} (ii) 即证明必要性. 设 (a) 和 (b) 成立，设 $V$ 是 $y$ 的环为 $A$ 的仿射邻域，$U$ 是 $x$ 在 $V$ 之上的环为 $B$ 的仿射邻域. 取 $U$ 足够小，由 (b) 可以假设存在平展的 $k(y)$-态射
\[
g \colon U|f^{-1}(y) \longrightarrow \mathrm{Spec}\, k[t_1, \dots, t_n] \quad (k = k(y)),
\]
由 $U|f^{-1}(y)$ 的结构层的 $n$ 个截面 $g_i$ 定义. 容易看到可以假设诸 $g_i$ (先验地是 $B \otimes_A k = BS^{-1}$ 的元素，其中 $S = A - \mathfrak{p}$，$\mathfrak{p}$ 是对应于 $y$ 的 $A$ 的素理想) 来自 $U$ 的结构层的截面，因此 $g$ 由一个仍记为 $g$ 的态射
\[
g \colon U \longrightarrow Y[t_1, \dots, t_n]
\]
诱导 (必要时把诸 $g_i$ 同乘以 $k$ 的同一个非零元素). 而由 (a)，$U$ 在 $Y$ 上平坦，$Y[t_1, \dots, t_n]$ 亦然，另一方面 $g$ 在 $y$ 之上的纤维之间诱导平展态射，故由 (\ref{cor:I.5.8})，$g$ 在 $x$ 处平展，证毕.
\end{proof}

\begin{corollary}\label{cor:II.2.2}
设 $S$ 是预概形，$f \colon X \to Y$ 是有限型 $S$-态射，$Y$ 在 $S$ 上有限型且平坦，$x \in X$，$s$ 是 $x$ 在 $S$ 上的投影. 为使 $f$ 在 $x$ 处光滑，必要且充分条件是 $X$ 在 $S$ 上于 $x$ 处平坦 (或：光滑)，且 $f$ 在 $s$ 的纤维上诱导的态射 $f_s \colon X_s \to Y_s$ 在 $x$ 处光滑.
\end{corollary}

\begin{proof}
只需证明充分性，它由判则 \ref{thm:II.2.1} 连同平坦性判则 (\ref{cor:I.5.9}) 得出.
\end{proof}

为了陈述下面的结果，``回忆''一个局部有限型态射 $f \colon X \to Y$ 称为在点 $x \in X$ 处等维 (equidimensional) 的，若 (令 $y = f(x)$) 可以找到 $x$ 的开邻域 $U$，其每个分支都支配 (dominate) $Y$ 的某个分支，使得对每个 $y' \in Y$，$f^{-1}(y') \cap U$ 的不可约分支都有同一个维数，且该维数与 $y'$ 无关. 在这个条件中，只需取 $y'$ 为 $Y$ 的经过 $y$ 的不可约分支的生成点，以及点 $y$. 例如若 $X$ 和 $Y$ 整且 $f$ 支配，则条件意味着经过 $x$ 的 $f^{-1}(y)$ 的分支有``正确''的维数，i.e. 生成纤维的维数 (回忆它们总是 $\geqslant$ 生成纤维的维数). 若 $f$ 在 $x$ 处等维，其纤维在 $x$ 处的维数为 $n$，且 $g \colon U \to Y' = Y[t_1, \dots, t_n]$ 是 $x$ 的邻域 $U$ 的 $Y$-态射，在 $y$ 的纤维上诱导在 $x$ 处拟有限 (或等价地，若 $g$ 在 $x$ 处拟有限) 的态射，则能证明 $U$ 的每个经过 $x$ 的不可约分支都支配 $Y'$ 的某个不可约分支. 此外，由``正规化引理'' (normalization lemma)，这样的 $g$ 总存在 (反过来，若存在 $x$ 的开邻域 $U$ 到形如 $Y' = Y[t_1, \dots, t_n]$ 的 $Y$-概形中的拟有限 $Y$-态射 $g$，使得 $U$ 的每个经过 $x$ 的分支都支配 $Y'$ 的某个分支，则 $f$ 在 $x$ 处等维). 有了这些准备：

\begin{proposition}\label{prop:II.2.3}
设 $f \colon X \to Y$ 是局部有限型态射，$x$ 是 $X$ 的点，$y = f(x)$，假设 $\mathcal{O}_y$ 正规 (normal). 为使 $f$ 在 $x$ 处光滑，必要且充分条件是 $f$ 在 $x$ 处等维，且 $f^{-1}(y)$ 在 $k(y)$ 上于 $x$ 处光滑.
\end{proposition}

\begin{proof}
由定义立即可见光滑态射是等维的\index[term]{dengweitaishe@等维态射} (N.B. 有限型平坦态射未必在 $x$ 处等维，即使其纤维在 $x$ 处不可约). 证明逆命题. 由于 $f^{-1}(y)$ 在 $k(y)$ 上于 $x$ 处光滑，可以假设 (必要时把 $X$ 换成 $x$ 的适当邻域) 存在 $Y$-态射
\[
g \colon X \longrightarrow Y[t_1, \dots, t_n] = Y',
\]
在 $y$ 的纤维上诱导平展态射，从而在 $x$ 处拟有限. 故 $g$ 非分歧，且 ($f$ 在 $x$ 处等维) $X$ 的经过 $x$ 的不可约分支各支配 $Y'$ 的某个分支，更不用说由 $g$ 推出的同态 $\mathcal{O}_{y'} \to \mathcal{O}_x$ (其中 $y' = g(x)$) 单射. 该同态还非分歧，且 $\mathcal{O}_{y'}$ 正规，因为它是环 $\mathcal{O}_y[t_1, \dots, t_n]$ 的局部化，而后者因 $\mathcal{O}_y$ 正规而正规. 故同态 $\mathcal{O}_{y'} \to \mathcal{O}_x$ 平展 (\ref{thm:I.9.5} (ii)).
\end{proof}

\begin{remark}\label{rem:II.2.4}
前面的陈述在把``$\mathcal{O}_y$ 正规''换成更弱的假设``$Y$ 在 $y$ 处几何单支 (geometrically unibranched)''(参见 I 11) 时仍成立 —— 因为 (\ref{thm:I.9.5}) 在该假设下成立. 借此机会同时指出：若整局部环 $A$ 的剩余域代数闭，则解析整 (analytically integral) (i.e. $\widehat{A}$ 整) 蕴含几何单支，反过来在``好环''的任何范畴中都成立，精确地说，在由通常运算封闭、且正规局部环的完备化仍正规的环的范畴中成立 (这一条件，由 Zariski 的``解析正规性定理'' (analytic normality theorem)，在仿射代数及其局部化的范畴中满足)\footnote{参见 EGA IV 7.8.}. 最后在本语境下``回忆''下述 Hironaka 的结果\footnote{参见 EGA IV 5.12.10.}，它有时能保证 $f^{-1}(y)$ 是约化 (reduced) 概形，i.e. 也就是许多代数几何学家滥用为 $f$ 在 $x$ 之上的纤维 (非分歧数) 的东西 (即 $f^{-1}(y)_{\mathrm{r\acute{e}d}}$)：
\end{remark}

\begin{proposition}\label{prop:II.2.5}
设 $f \colon X \to Y$ 是约化预概形的有限型支配态射，$y$ 是 $Y$ 的点，使得 $\mathcal{O}_y$ 正则 (regular). 假设 $f^{-1}(y)$ 的所有分支都重数 (multiplicity) 为 1 (参见下面的定义)，且 $f^{-1}(y)_{\mathrm{r\acute{e}d}}$ 正规. 则 $f^{-1}(y)$ 约化从而正规，$X$ 在 $f^{-1}(y)$ 的所有点处正规，最后 $X$ 在 $f^{-1}(y)$ 的所有点处平坦.
\end{proposition}

称 $f^{-1}(y)$ 的分支 $Z$ 重数为 1，若 $x$ 表示 $Z$ 的生成点，有 (i) $\dim \mathcal{O}_x = \dim \mathcal{O}_y$ (i.e. $Z$ 不是``过剩分支'' (excess component)，即不是``维数太大''的)；(ii) $\mathcal{O}_x$ 的极大理想由 $\mathcal{O}_y$ 的极大理想生成 (后者按 $x$ 的选取先验地生成 $\mathcal{O}_x$ 的一个定义理想).

结合 \ref{prop:II.2.3} 或 \ref{thm:II.2.1}，得到：

\begin{corollary}\label{cor:II.2.6}
设 $f \colon X \to Y$ 是约化预概形的有限型支配态射，$y$ 是 $Y$ 的点，使得 $\mathcal{O}_y$ 正则. 为使 $f$ 在 $X$ 中位于 $y$ 之上的点处光滑，必要且充分条件是 $f^{-1}(y)$ 的分支重数为 1，且 $f^{-1}(y)_{\mathrm{r\acute{e}d}}$ 在 $k(y)$ 上光滑.
\end{corollary}

在过去，当 $Y$ 是离散赋值环 (discrete valuation ring) $A$ 的谱时，人们尤其考虑这种情形，并通常以诸如此类的措辞称之：``若 $X$ 关于所给赋值的约化 (reduction) 是好的''... 此外，那时 $X$ 表示 $\mathbf{P}_K^n$ 的一个闭子概形 ($K$ 是 $A$ 的分式域)，而由于缺乏适当的语言，``定义在 $A$ 上''(而不仅在 $K$ 上) 的对象的更内蕴的作用几乎不显现.

\section{永久性质}

\begin{proposition}\label{prop:II.3.1}
设 $f \colon X \to Y$ 是态射，$x \in X$，$y = f(x)$. 假设 $f$ 在 $x$ 处光滑. 为使 $\mathcal{O}_x$ 约化 (相应地：正则，相应地：正规)，必要且充分条件是 $\mathcal{O}_y$ 约化 (相应地：正则，相应地：正规).
\end{proposition}

\begin{proof}
当 $X$ 形如 $Y[t_1, \dots, t_n]$ 时该陈述确为已知，且对平展态射在 (I，§9) 中已证明；一般情形借助定义 \ref{prop:II.1.1} 立即得出.
\end{proof}

我们不在此详述其他永久性质，它们已经由平坦性本身得出，或由``$X$ 在一个形如 $Y[t_1, \dots, t_n]$ 的 $Y$-预概形之上局部拟有限且平坦''(或者，如我们所说，``$X$ 在 $Y$ 之上 Cohen-Macaulay'') 这一事实得出. 只指出由后一事实得出
\begin{equation}
\dim \mathcal{O}_x = \dim \mathcal{O}_y + n - d, \qquad \operatorname{prof} \mathcal{O}_x = \operatorname{prof} \mathcal{O}_y + n - d, \tag{3.1}
\label{eq:II.3.1}
\end{equation}
其中 $n$ 是 $f$ 的纤维在 $x$ 处的维数，$d$ 是 $k(x)$ 在 $k(y)$ 上的超越次数 (transcendence degree)，于是 (令 coprof = dim $-$ prof)\footnote{对这些公式参见 EGA IV 6.1 和 6.3.}
\begin{equation}
\operatorname{coprof} \mathcal{O}_x = \operatorname{coprof} \mathcal{O}_y. \tag{3.2}
\label{eq:II.3.2}
\end{equation}
例如由此得出 $\mathcal{O}_x$ 是 Cohen-Macaulay 的 (相应地：无嵌入分支) 当且仅当 $\mathcal{O}_y$ 亦然.

\section{光滑态射的微分性质}

为简化起见，我们主要限于一阶微分演算，只对高阶作些快速说明 (那里的结果同样简单).

关于 $Y$-预概形 $X$ 的一次微分 (1-différentielle) 层 $\varOmega_{X/Y}^{1}$ 的定义，参见 (I §1). 假设 $X$ 和 $Y$ 是 $S$-预概形，结构态射 $f \colon X \to Y$ 是 $S$-态射. 则 $f$ 定义一个 (与 $f$ 相容的) 模层 (Module) 同态
\begin{equation}
f^{*} \colon \varOmega_{Y/S}^{1} \longrightarrow \varOmega_{X/S}^{1}, \tag{4.1}
\label{eq:II.4.1}
\end{equation}
换言之，$\varOmega_{X/S}^{1}$ 对 $S$-预概形 $X$ 是反变的. 此外，\eqref{eq:II.4.1} 等价于 $X$ 上的模层同态
\begin{equation}
f^{*}\bigl(\varOmega_{Y/S}^{1}\bigr) \longrightarrow \varOmega_{X/S}^{1}, \tag{4.1bis}
\label{eq:II.4.1bis}
\end{equation}
也姑且记为 $f^{*}$，它嵌入一个模层同态的典范正合序列
\begin{equation}
f^{*}\bigl(\varOmega_{Y/S}^{1}\bigr) \longrightarrow \varOmega_{X/S}^{1} \longrightarrow \varOmega_{X/Y}^{1} \longrightarrow 0. \tag{4.2}
\label{eq:II.4.2}
\end{equation}
所有这些同态都由``局部性质 (归结到仿射情形) 且与算子 $d$ 交换''这一条件定义. \eqref{eq:II.4.2} 的正合性是经典而平凡的，在仿射情形翻译成 (对应于 $A$-代数的同态 $B \to C$) 的正合序列：
\begin{equation}
\varOmega_{B/A}^{1} \otimes_B C \longrightarrow \varOmega_{C/A}^{1} \longrightarrow \varOmega_{C/B}^{1} \longrightarrow 0. \tag{4.2bis}
\label{eq:II.4.2bis}
\end{equation}

\begin{lemma}\label{lem:II.4.1}
设 $f \colon X \to Y$ 是 $S$-预概形的态射. 若 $f$ 非分歧 (相应地：平展)，则 $f^{*}\bigl(\varOmega_{Y/S}^{1}\bigr) \to \varOmega_{X/S}^{1}$ 满射 (相应地：同构). 逆命题在``非分歧''的情形、且 $f$ 假设为局部有限型时成立.
\end{lemma}

\begin{proof}
非分歧的情形由正合序列 \eqref{eq:II.4.2} 和 (\ref{prop:I.3.1}) 得出，但也可以像平展的情形那样直接看. 考虑图
\[
\begin{tikzcd}[row sep=large, column sep=large]
X \arrow[r,"\Delta_{X/Y}"] \arrow[d] & X \times_Y X \arrow[r] & X \times_S X \\
Y \arrow[r,"\Delta_{Y/S}"] & Y \times_S Y
\end{tikzcd}
\]
(原文图式如此)，其中 $X \times_Y X$ 等同于 $Y$ 与 $X \times_S X$ 在 $Y \times_S Y$ 上的纤维积. 由于 $f$ 非分歧，$X \to X \times_Y X$ 是开浸入，故复合浸入 $\Delta_{X/S}$ (由后者与 $X \times_Y X \to X \times_S X$ 复合而得) 的``余法'' (conormal) 层同构于浸入 $X \times_Y X \to X \times_S X$ 的余法层在 $X$ 上的原像. 另一方面，$X \to Y$ 平展从而平坦，故 $X \times_S X \to Y \times_S Y$ 平坦，因此浸入 $X \times_Y X \to X \times_S X$ 的余法层同构于浸入 $Y \to Y \times_S Y$ 的余法层的原像，i.e. $\varOmega_{Y/S}^{1}$ 的原像. 结论由此得出.
\end{proof}

\begin{lemma}\label{lem:II.4.2}
设 $X = Y[t_1, \dots, t_n]$，$Y$ 是 $S$-预概形. 则典范同态序列
\[
0 \longrightarrow f^{*}\bigl(\varOmega_{Y/S}^{1}\bigr) \longrightarrow \varOmega_{X/S}^{1} \longrightarrow \varOmega_{X/Y}^{1} \longrightarrow 0
\]
正合，且 $\varOmega_{X/Y}^{1}$ 以 $d_{X/Y} t_i$ 为基自由.
\end{lemma}

\begin{proof}
(纯仿射的) 验证是直接的. (N.B. 已经知道 \eqref{eq:II.4.2} 的正合性).
\end{proof}

结合这两个陈述和定义 \ref{prop:II.1.1}，得到：

\begin{theorem}\label{thm:II.4.3}
设 $f \colon X \to Y$ 是 $S$-预概形的光滑态射，则：
\begin{enumerate}[(i)]
\item 典范同态序列
\[
0 \longrightarrow f^{*}\bigl(\varOmega_{Y/S}^{1}\bigr) \longrightarrow \varOmega_{X/S}^{1} \longrightarrow \varOmega_{X/Y}^{1} \longrightarrow 0
\]
正合.
\item $\varOmega_{X/Y}^{1}$ 局部自由，其在 $x$ 处的秩 (rank) $n$ 等于 $f$ 在 $x$ 处的相对维数.
\end{enumerate}
\end{theorem}

\begin{corollary}\label{cor:II.4.4}
同态 $f^{*}\bigl(\varOmega_{Y/S}^{1}\bigr) \to \varOmega_{X/S}^{1}$ 单射，其在 $\varOmega_{X/S}^{1}$ 中的像局部是直因子 (direct factor).
\end{corollary}

设 $u \colon F \to G$ 是预概形 $X$ 上的模层同态，称它在 $x \in X$ 处泛单射 (universally injective)，若 $\mathcal{O}_x$-模同态 $F_x \to G_x$ 单射，且与任何 $\mathcal{O}_x$-代数 (或等价地，与任何 $\mathcal{O}_x$-模) 张量后仍单射. 例如，只需存在 $x$ 的开邻域 $U$ 使得 $u$ 诱导 $F|U$ 到 $G|U$ 的某个直因子上的同构；当 $F$ 和 $G$ 在 $x$ 的邻域中自由 (且有限型) 时，这一条件也是必要的，精确地说此时下列条件等价：
\begin{enumerate}[(i)]
\item $u$ 在 $x$ 处单射且 Coker $u$ 在 $x$ 处自由；
\item 存在 $x$ 的开邻域 $U$ 使得 $u$ 诱导 $F|U$ 到 $G|U$ 的某个直因子上的同构；
\item $u$ 在 $x$ 处泛单射；
\item $u$ 在``限制纤维'' $F_x \otimes k(x) \to G_x \otimes k(x)$ 上诱导的同态单射；
\item 转置同态 $\check{G} \to \check{F}$ 在点 $x$ (或等价地，在 $x$ 的邻域) 满射.
\end{enumerate}
(循环证明：(iv)$\Rightarrow$(v) 由 Nakayama 引理得出，另一方面 (v)$\Rightarrow$(i)，因为局部自由商层必为直因子). 几何上，所考虑的情形意味着 $u$ 对应于截面层为 $F$ 的向量丛 (vector bundle) 到由 $G$ 定义的类似向量丛的一个子丛上的同构. 当然，仅凭 $F \to G$ 单射并不够.

\begin{corollary}\label{cor:II.4.5}
设 $f \colon X \to Y$ 是 $S$-预概形的局部有限型态射，$x \in X$，$y = f(x)$，$s$ 是 $x$ 和 $y$ 在 $S$ 上的投影. 假设 $Y$ 在 $y$ 处 $S$-光滑. 下列条件等价：
\begin{enumerate}[(i)]
\item $f$ 在 $x$ 处光滑.
\item $X$ 在 $S$ 上于 $x$ 处光滑，且 $f^{*}\bigl(\varOmega_{Y/S}^{1}\bigr) \to \varOmega_{X/S}^{1}$ 在 $x$ 处泛单射，i.e. 它在 $x$ 处是单射且其余核 $\varOmega_{X/Y}^{1}$ 在 $x$ 处自由.
\end{enumerate}
\end{corollary}

\begin{proof}
必要性由 \ref{prop:II.1.3} (iii) 和 \ref{thm:II.4.3} (i) (ii) 得出，证明充分性. 由于 $dg$ ($g \in \mathcal{O}_x$) 在 $x$ 处生成模 $\varOmega_{X/Y}^{1}$，可以找到 $g_i$ ($1 \leqslant i \leqslant n$) 使得 $dg_i$ 在 $\bigl(\varOmega_{X/Y}^{1}\bigr)_x$ 中的像构成该模的基. 取 $X$ 足够小，可以假设诸 $g_i$ 来自 $\mathcal{O}_X$ 的截面，从而定义 $Y$-态射 $g \colon X \to Y' = Y[t_1, \dots, t_n]$. 利用假设和引理 \ref{lem:II.4.2}，容易看到相应的同态 $g^{*}\bigl(\varOmega_{Y'/S}^{1}\bigr) \to \varOmega_{X/S}^{1}$ 在 $x$ 处双射，这归结为证明：
\end{proof}

\begin{corollary}\label{cor:II.4.6}
设 $f \colon X \to Y$ 是光滑 $S$-预概形的态射. 为使 $f$ 在 $x \in X$ 处平展，必要且充分条件是 $f^{*}\bigl(\varOmega_{Y/S}^{1}\bigr) \to \varOmega_{X/S}^{1}$ 在 $x$ 处是同构.
\end{corollary}

\begin{proof}
由 \ref{lem:II.4.1} 知其必要，且该条件由同一引理蕴含 $f$ 在 $x$ 处非分歧. 由 \ref{cor:II.4.6} 归结到 $S = \mathrm{Spec}(k)$ 的情形. 由于 $Y$ 在 $k$ 上光滑，它正则，从而是正规的，由 (\ref{thm:I.9.5} (ii)) 归结为证明 $\mathcal{O}_y \to \mathcal{O}_x$ 单射，或 $\mathcal{O}_y$ 与 $\mathcal{O}_x$ 维数相同. 而这两个维数分别是 $\varOmega_{Y/k}^{1}$ 和 $\varOmega_{X/k}^{1}$ 在 $y$ 处 (相应地：$x$ 处) 的秩，由假设它们相等.
\end{proof}

\begin{remark}\label{rem:II.4.7}
设 $X$ 和 $Y$ 在 $S$ 上光滑，$f \colon X \to Y$ 的光滑性判则 \ref{cor:II.4.5} (ii) 还可以说成：对每个 $x \in X$，$f$ 在 $x$ 处 (相对基 $S$ 的) 切映射 (tangent map)，i.e. $f^{*}\bigl(\varOmega_{Y/S}^{1}\bigr)$ 与 $\varOmega_{X/S}^{1}$ 在 $x$ 处的限制纤维这两个有限维 $k(x)$-向量空间之间同态的转置，是满射. 这对于使用解析空间 (analytic space) 的人们尤其是一个很熟悉的假设. 他们通常所作的非奇异 (non-singularity) 假设 (它意味着 $X$ 和 $Y$``在 $\mathbf{C}$ 上光滑''，参见 §5) 似乎只是出于代数簇或解析空间的奇异点至今仍使许多几何学家感到恐惧.
\end{remark}

指出 \ref{cor:II.4.6} 的下列特殊情形：

\begin{corollary}\label{cor:II.4.8}
设 $X$ 是 $S$-预概形，$g \colon X \to S[t_1, \dots, t_n]$ 是由 $\mathcal{O}_X$ 的截面 $g_i$ ($1 \leqslant i \leqslant n$) 定义的 $S$-态射，$x$ 是 $X$ 的点，使得 $X$ 在 $S$ 上于 $x$ 处光滑. 为使 $g$ 在 $x$ 处平展，必要且充分条件是 $dg_i$ ($1 \leqslant i \leqslant n$) 在 $x$ 处构成 $\varOmega_{X/S}^{1}$ 的基 (或等价地，它们在 $\varOmega_{X/S}^{1}(x) = \bigl(\varOmega_{X/S}^{1}\bigr)_x \otimes_{\mathcal{O}_x} k(x)$ 中的像构成该 $k(x)$-向量空间的基).
\end{corollary}

\begin{proof}
设 $X$ 是预概形，$Y$ 是由凝聚 (coherent) 理想层 $\mathcal{J}$ 定义的 $X$ 的闭子预概形. 于是 $\mathcal{J}/\mathcal{J}^2$ 可以看作 $Y$ 上的凝聚层 ($Y$ 在 $X$ 中的余法层). 若现在 $X$ 是 $S$-预概形，则有 $Y$ 上的拟凝聚层典范正合序列
\begin{equation}
\mathcal{J}/\mathcal{J}^2 \xrightarrow{d} \varOmega_{X/S}^{1} \otimes_{\mathcal{O}_X} \mathcal{O}_Y \longrightarrow \varOmega_{Y/S}^{1} \longrightarrow 0, \tag{4.3}
\label{eq:II.4.3}
\end{equation}
其右边部分正是 \eqref{eq:II.4.2} (其中 $X$ 与 $Y$ 的角色互换，考虑到 $\varOmega_{Y/X}^{1} = 0$)，而同态 $\mathcal{J}/\mathcal{J}^2 \to \varOmega_{X/S}^{1} \otimes_{\mathcal{O}_X} \mathcal{O}_Y$ 由 (一般非线性) 同态 $g \mapsto dg$ 取商得出. \eqref{eq:II.4.3} 的正合性是经典而平凡的，在仿射情形由下列正合序列解释 (对应于 $A$-代数满同态 $B \to C$，核为 $J$)：
\begin{equation}
J/J^2 \longrightarrow \varOmega_{B/A}^{1} \otimes_B C \longrightarrow \varOmega_{C/A}^{1} \longrightarrow 0 \quad (C = B/J). \tag{4.3bis}
\label{eq:II.4.3bis}
\end{equation}
(这个正合序列在 (\ref{prop:II.4.9}) 的证明中其实已被隐含使用!)
\end{proof}

\begin{proposition}\label{prop:II.4.9}
设 $X$ 是 $S$-预概形，$Y$ 是由 $X$ 上的凝聚理想层 $\mathcal{J}$ 定义的 $X$ 的闭子预概形，$x$ 是 $X$ 的点，$g_i$ ($1 \leqslant i \leqslant n$) 是 $\mathcal{O}_X$ 的截面，定义 $S$-态射
\[
g \colon X \longrightarrow S[t_1, \dots, t_n] = X',
\]
最后 $p$ 是整数，$0 \leqslant p \leqslant n$. 假设 $X$ 在 $S$ 上于 $x$ 处光滑. 下列条件等价：
\begin{enumerate}[(i)]
\item 存在 $x$ 在 $X$ 中的开邻域 $X_1$，使得 $g|X_1$ 平展，且 $Y_1 = Y \cap X_1$ ($Y$ 在 $X_1$ 上的迹) 是 $X' = S[t_1, \dots, t_n]$ 的闭子预概形 $Y' = S[t_{p+1}, \dots, t_n]$ 的原像 (i.e. $g_i$ ($1 \leqslant i \leqslant p$) 生成 $\mathcal{J}|X_1$)；
\item $Y$ 在 $S$ 上于 $x$ 处光滑，$g_i$ ($1 \leqslant i \leqslant p$) 定义 $\mathcal{J}_x$ 的元素，$dg_i(x)$ ($1 \leqslant i \leqslant n$) 构成 $k(x)$ 上的基 $\varOmega_{X/S}^{1}(x)$，$dg'_i(x)$ ($p+1 \leqslant i \leqslant n$) 构成 $k(x)$ 上的基 $\varOmega_{Y/S}^{1}(x)$ (其中 $g'_i$ 表示 $g_i$ 到 $Y$ 的限制；微分是关于 $S$ 取的)；
\item $g_i$ ($1 \leqslant i \leqslant p$) 定义 $\mathcal{J}_x$ 的生成元组，且 $dg_i(x)$ ($1 \leqslant i \leqslant n$) 构成 $k(x)$ 上的基 $\varOmega_{X/S}^{1}(x)$；
\item $Y$ 在 $S$ 上于 $x$ 处光滑，$g_i$ ($1 \leqslant i \leqslant p$) 构成 $\mathcal{J}_x$ 的极小生成元组，$dg'_i(x)$ ($p+1 \leqslant i \leqslant n$) 构成 $k(x)$ 上的基 $\varOmega_{Y/S}^{1}(x)$.
\end{enumerate}
此外，在这些条件下，$\mathcal{J}/\mathcal{J}^2$ 是 $Y$ 上在 $x$ 处自由的模层，以诸 $g_i$ ($1 \leqslant i \leqslant p$) 的类为 $x$ 处的基，且典范同态 $\mathcal{J}/\mathcal{J}^2 \to \varOmega_{X/S}^{1} \otimes \mathcal{O}_Y$ 在 $x$ 处泛单射.
\end{proposition}

\begin{remarkstar}
这意味着 $p$ 由其他条件确定，或作为 $\mathcal{J}/\mathcal{J}^2$ 在 $Y$ 上于 $x$ 处的自由模层的秩，或作为 $\mathcal{J}_x$ 在 $X$ 上的生成元的最小个数，或最后由``$Y$ 相对 $S$ 在 $x$ 处的相对维数为 $n - p$''这一事实确定.
\end{remarkstar}

\paragraph*{证明.} 先假设 (i) 成立. 则由 (\ref{prop:II.4.9} (iii))，$Y_1$ 在 $Y'$ 上平展，故按定义它在 $S$ 上于 $x$ 处光滑 (相对维数 $n - p$)，$Y$ 亦然. 于是由 (\cref{cor:II.4.8}) 得出 $dg_i$ ($1 \leqslant i \leqslant n$) 构成 $\varOmega_{X/S}^{1}$ 在 $x$ 处的基，$dg'_i$ ($p+1 \leqslant i \leqslant n$) 构成 $\varOmega_{Y/S}^{1}$ 在 $x$ 处的基，由此由正合序列 \eqref{eq:II.4.3} 得出 $g_i$ ($1 \leqslant i \leqslant p$) 在 $\mathcal{J}/\mathcal{J}^2$ (看作 $Y$ 上的模层) 中于 $x$ 处线性无关；由于 $g_i$ ($1 \leqslant i \leqslant p$) 生成 $\mathcal{J}_x$，故 $g_i$ 模 $\mathcal{J}_x^2$ 构成 $\mathcal{J}/\mathcal{J}^2$ 在 $x$ 处的基. 这既蕴含 $g_i$ ($1 \leqslant i \leqslant p$) 构成 $\mathcal{J}_x$ 的极小生成元组，又蕴含 \eqref{eq:II.4.3} 的同态 $\mathcal{J}/\mathcal{J}^2 \to \varOmega_{X/S}^{1} \otimes \mathcal{O}_Y$ 在 $x$ 处泛单射 (因为它把一个 $x$ 处自由模层的基映到 $x$ 处自由模层的基的一部分 —— N.B. 这是 $Y$-模层). 这证明 (i) 蕴含 (ii)，(iii)，(iv) 以及命题 \ref{prop:II.4.9} 的最后断言.

(iii) 由推论 \ref{cor:II.4.8} 蕴含 (i).

(ii) 蕴含 (i). 事实上，(ii) 中的第一个假设意味着 (把 $X$ 换成 $x$ 在 $X$ 中的开邻域) $g$ 诱导态射 $h \colon Y \to Y'$. 由 \ref{cor:II.4.8}，(ii) 的其他两个假设意味着 $g$ 在 $x$ 处平展，$h$ 在 $x$ 处平展. 设 $Y''$ 是 $Y'$ 在 $g$ 下的原像. 于是 $Y$ 是 $Y''$ 的闭子预概形，而由 (\ref{prop:II.4.9} (iii))，$Y''$ 在 $Y'$ 上于 $x$ 处平展 (因为 $g$ 在 $x$ 处平展). 故浸入态射 $Y \to Y''$ 本身平展 (\ref{cor:II.4.8})，从而是开浸入 (\ref{cor:I.5.8} 或 \ref{cor:I.5.2})，于是再把 $X$ 换成 $x$ 的适当开邻域 $X_1$，得到 (i).

以上建立了条件 (i)，(ii)，(iii) 的等价性，以及它们蕴含 (iv)；剩下证明 (iv)$\Rightarrow$(ii)，一旦知道``$Y$ 在 $S$ 上于 $x$ 处光滑蕴含 $\mathcal{J}/\mathcal{J}^2$ 在 $Y$ 上于 $x$ 处自由，且同态 $\mathcal{J}/\mathcal{J}^2 \to \varOmega_{X/S}^{1} \otimes \mathcal{O}_Y$ 在 $x$ 处泛单射''(考虑到 $\varOmega_{X/S}^{1}$ 在 $X$ 上于 $x$ 处自由) 就立即得出. 最后这一点包含在

\begin{theorem}\label{thm:II.4.10}
设 $X$ 是光滑 $S$-预概形，$Y$ 是由 $X$ 上的凝聚理想层 $\mathcal{J}$ 定义的 $X$ 的闭子预概形，$x$ 是 $X$ 的点. 下列条件等价：
\begin{enumerate}[(i)]
\item $Y$ 在 $S$ 上于 $x$ 处光滑.
\item 存在 $x$ 在 $X$ 中的开邻域 $X_1$ 和平展 $S$-态射
\[
g \colon X_1 \longrightarrow X' = S[t_1, \dots, t_n],
\]
使得 $Y_1 = Y \cap X_1$ ($Y$ 在 $X_1$ 上的迹) 是 $X_1$ 中由 $g$ 从 $X' = S[t_1, \dots, t_n]$ 的闭子预概形 $Y' = S[t_{p+1}, \dots, t_n]$ 取原像得到的子预概形，对某个适当的 $p$.
\item 存在 $\mathcal{J}_x$ 的生成元 $g_i$ ($1 \leqslant i \leqslant p$)，使得 $dg_i$ 构成 $\varOmega_{X/S}^{1}$ 在 $x$ 处的基的一部分 (或等价地，使 $dg_i(x)$ 在 $\varOmega_{X/S}^{1}$ 中在 $k(x)$ 上线性无关).
\item 层 $\mathcal{J}/\mathcal{J}^2$ 在 $Y$ 上于 $x$ 处自由，且典范同态
\[
d \colon \mathcal{J}/\mathcal{J}^2 \longrightarrow \varOmega_{X/S}^{1} \otimes \mathcal{O}_Y
\]
在 $x$ 处泛单射；或：典范同态序列
\[
0 \longrightarrow \mathcal{J}/\mathcal{J}^2 \longrightarrow \varOmega_{X/S}^{1} \otimes \mathcal{O}_Y \longrightarrow \varOmega_{Y/S}^{1} \longrightarrow 0
\]
在 $x$ 处正合，且 $\varOmega_{Y/S}^{1}$ 在 $x$ 处局部自由.
\end{enumerate}
\end{theorem}

\paragraph*{证明.} 由前面的讨论已知 (ii) 蕴含 (i)，(iii)，(iv). 证明 (i)$\Rightarrow$(ii) (这将同时完成 \ref{prop:II.4.9} 的证明). 由定理 \ref{thm:II.4.10} (ii)，正合序列 \eqref{eq:II.4.3} 中最后两项是 $Y$ 上的自由模层. 于是，由于 $dg$ ($g \in \mathcal{O}_X$) 在 $\varOmega_{X/S}^{1} \otimes_{\mathcal{O}_X} \mathcal{O}_Y$ 中的像在 $x$ 处生成该模，故它们在 $\varOmega_{Y/S}^{1}$ 中的像在 $x$ 处生成后者，可以找到 $g_i$ ($p+1 \leqslant i \leqslant n$) $\in \mathcal{O}_X$ 使得 $dg'_i$ 构成 $\varOmega_{Y/S}^{1}$ 的基，然后 (由 \eqref{eq:II.4.3} 的正合性) 用形如 $dg_i$ ($1 \leqslant i \leqslant p$，诸 $g_i \in \mathcal{J}_x$) 的元素把 $dg_i$ ($p+1 \leqslant i \leqslant n$) 的组补成中项的基. 诸 $g_i$ 来自 $X$ 中 $x$ 的邻域上的 $\mathcal{O}_X$ 的截面，可以假设该邻域等于 $X$. 于是处于 \ref{cor:II.4.8} (ii) 的条件下，而我们已经建立了这蕴含条件 \ref{cor:II.4.8} (i)，由此得 \ref{thm:II.4.10} (ii).

\ref{thm:II.4.10} 中 (iii)$\Rightarrow$(ii) 由 \ref{cor:II.4.8} 中 (iii)$\Rightarrow$(i) 立即得出. 故 (i)，(ii)，(iii) 等价，且蕴含 (iv). 最后，(iv) 蕴含 (iii) 是平凡的，考虑到模 $\mathcal{J}_x^2$ 构成 $\mathcal{J}_x$ 的基的 $g_i \in \mathcal{J}_x$ 生成 $\mathcal{J}_x$ (Nakayama 引理).

此外，前面的证明还显示：

\begin{corollary}\label{cor:II.4.11}
设 $X$ 是 $S$-预概形，$Y$ 是由 $X$ 上的凝聚理想层 $\mathcal{J}$ 定义的闭子预概形，$x$ 是 $Y$ 的点. 假设 $X$ 和 $Y$ 在 $S$ 上于 $x$ 处光滑. 设 $g_i$ 是 $\mathcal{J}$ 的截面 ($1 \leqslant i \leqslant p$). 下列条件等价：
\begin{enumerate}[(i)]
\item $g_i$ 生成 $\mathcal{J}_x$，且 $dg_i(x)$ 在 $\varOmega_{X/S}^{1}(x)$ 中在 $k(x)$ 上线性无关；
\item $g_i$ 模 $\mathcal{J}^2$ 构成 $\mathcal{J}/\mathcal{J}^2$ 在 $x$ 处的基；
\item $g_i$ 构成 $\mathcal{J}_x$ 的极小生成元组；
\item 可以找到 $X$ 的邻域 $X_1$ 上 $\mathcal{O}_X$ 的另外一些截面 $g_i$ ($p+1 \leqslant i \leqslant n$)，与前面的截面一起定义平展态射 $X_1 \to X' = S[t_1, \dots, t_n]$，使得 $Y_1 = Y \cap X_1$ 是 $g$ 从 $X' = S[t_1, \dots, t_n]$ 的闭子预概形 $Y' = S[t_{p+1}, \dots, t_n]$ 取原像得到的.
\end{enumerate}
\end{corollary}

特别地：

\begin{corollary}\label{cor:II.4.12}
设 $X$ 是 $S$-预概形，$F$ 是 $\mathcal{O}_X$ 的截面，$Y$ 是 $F$ 的零点子概形 (由凝聚理想 $F \cdot \mathcal{O}_X$ 定义)，$x$ 是 $Y$ 的点. 假设 $X$ 在 $S$ 上于 $x$ 处简单. 为使 $Y$ 在 $S$ 上于 $x$ 处光滑，必要且充分条件是 $F$ 在 $x$ 的邻域内为零，或 $dF(x) \neq 0$ (其中 $dF(x)$ 表示 $dF$ 在 $k(x)$ 上的向量空间 $\varOmega_{X/S}^{1}(x)$ 中的像).
\end{corollary}

\begin{proof}
由 \ref{thm:II.4.10} 判则 (iii) 知其充分. 它也是必要的，因为 $\mathcal{J}$ 由一个元素生成，首先必须 $\mathcal{J}/\mathcal{J}^2$ 在点 $x$ 处自由且秩 $\leqslant 1$. 若此秩为 0，i.e. $\mathcal{J}/\mathcal{J}^2 = 0$ 在 $x$ 处，则由 Nakayama 引理推出 $\mathcal{J} = 0$ 在 $x$ 处，i.e. $F$ 在 $x$ 的邻域内为零. 若此秩为 1，则 $F$ 构成 $\mathcal{J}$ 在 $x$ 处的极小生成元组，由 (\cref{cor:II.4.11}) 中 (i) 与 (iii) 的等价性即得结论.
\end{proof}

\begin{corollary}\label{cor:II.4.13}
设 $Y$ 是局部有限型 $S$-预概形，$S'$ 是平坦 $S$-预概形，$Y' = Y \times_S S'$，$x'$ 是 $Y'$ 的点，$x$ 是它在 $Y$ 中的典范像. 为使 $Y$ 在 $S$ 上于 $x$ 处光滑，必要且充分条件是 $Y'$ 在 $S'$ 上于 $x'$ 处光滑. 特别地，若 $S' \to S$ 平坦且满射，则 $Y$ 在 $S$ 上光滑当且仅当 $Y'$ 在 $S'$ 上光滑.
\end{corollary}

\begin{proof}
只需证明充分性 (必要性已在 \ref{prop:II.1.3} (ii) 中看到). 可以假设 (把 $Y$ 换成 $x$ 的适当邻域，$Y'$ 换成它的原像) $Y$ 是仿射 $S$ 上的仿射有限型概形，故 $Y$ 同构于 $S[t_1, \dots, t_n]$ 的闭子概形. 于是 $Y'$ 等同于 $X' = X \times_S S'$ 的闭子预概形. 由于 $X$ 在 $S$ 上光滑，$X'$ 在 $S'$ 上光滑，可以应用光滑性判则 \ref{thm:II.4.10}. 这里判则 (iv) 立即给出结果.
\end{proof}

\begin{remark}\label{rem:II.4.14}
定理 \ref{thm:II.4.10} 的判则 (iii) 值得称为 \emph{Jacobi 光滑性判则} (Jacobian criterion of smoothness)\index[term]{Jacobi guanghuaxingpanze@Jacobi 光滑性判则}. 它使得理论上能够识别一个给定的 $S$-预概形 $Y$ 是否在 $S$ 上于 $Y$ 的点 $x$ 处光滑，因为 $Y$ 的邻域总同构于光滑 $S$-预概形 $X$ 的子概形，例如 $X = S[t_1, \dots, t_n]$. 事实上，通常正是在 $X = S[t_1, \dots, t_n]$，$S = \mathrm{Spec}(A)$ 时陈述 Jacobi 判则 (当然，在 Zariski 所考虑的经典情形，$A$ 是域). 把与 $A[t_1, \dots, t_n]$ 的理想 $J$ 及包含它的素理想有关的陈述留给读者. 注意到目前看来 (尤其自从 Nagata 用非微分方法推广了 Zariski 的``代数概形的正则点集是开的''这一定理以来)，Jacobi 判则几乎只有以我们这里所给出的形式才有意义 (i.e. 只使用相对微分，而不使用绝对微分，i.e. 相对绝对常数环 $\mathbf{Z}$ 的微分). 和常常见到的情况一样，在这里考虑微分比考虑导子更方便. 最后注意，若 $Y$ 在 $S$ 上于 $x$ 处光滑，相对维数为 $n$，则存在 $Y$ 中 $x$ 的开邻域同构于 $X = S[t_1, \dots, t_n]$ 的子概形，其中 $n = m + 1$，这由定义和 (\ref{thm:II.4.10}) 得出.
\end{remark}

设 $A$ 是诺特环，$x_i$ ($1 \leqslant i \leqslant n$) 是 $A$ 的元素，$J$ 是由诸 $x_i$ 生成的理想. 称诸 $x_i$ 构成 $J$ 的正则生成元组 (regular system of generators)，若典范满同态
\[
(A/J)[t_1, \dots, t_n] \longrightarrow \mathrm{gr}^{J}(A)
\]
(由诸 $x_i$ 定义，其中右端表示与 $A$ 按 $J$ 的幂滤链所得的相伴分次环) 是同构. 这一条件等价于下列两个条件：
\begin{enumerate}[(i)]
\item 典范满同态
\[
S_{A/J}(J/J^2) \longrightarrow \mathrm{gr}^{J}(A)
\]
(其中左端表示 $A/J$-模 $J/J^2$ 的对称代数) 是同构，且
\item $J/J^2$ 自由，以诸 $x_i$ 模 $J^2$ 的类为基.
\end{enumerate}
以这种形式看出，若 $J \neq A$，则诸 $x_i$ 构成 $J$ 的极小生成元组，且 $J$ 的任何其他极小生成元组都是正则生成元组 (N.B. ``极小'' 取严格意义：元素个数最少，它仅当 $A$ 局部时才等价于``按包含关系极小''的意义)；另一方面，若 $J = A$，则 $J$ 的任何生成元组都是正则的.

理想的正则生成元组的正则性对任意乘法稳定集 (multiplicatively closed set) 的局部化稳定，另一方面立即看出，为使 $(x_i)$ 是 $J$ 的极小生成元组，只需对包含 $J$ 的每个极大理想 $m$，$x_i$ 在 $A_{\mathfrak{m}}$ 中定义 $JA_{\mathfrak{m}}$ 的正则生成元组. 因此归结到 $A$ 是以 $m$ 为极大理想的局部环、诸 $x_i$ 属于 $m$ 的情形. 此时诸 $x_i$ 构成 $J$ 的正则生成元组，当且仅当它们构成 Serre 意义下的 $A$-序列 (A-sequence)\footnote{现在我们宁愿说 ``$A$-正则序列'' (A-regular sequence)，参见 EGA 0$_{\mathrm{IV}}$ 15.1.7 和 15.1.11.}，i.e. 若对每个满足 $1 \leqslant i \leqslant n$ 的 $i$，$x_i$ 在 $A/(x_1, \dots, x_{i-1})A$ 中非零因子 (non-diviseur de 0).

最后，当 $A$ 是环 $B$ 上的代数，且 $A/J$ 作为 $B$-代数同构于 $B$ (从而 $J$ 是 $B$-代数同态 $A \to B$ 的核) 时，诸 $x_i$ 构成 $J$ 的正则生成元组，当且仅当典范同态
\[
B[\![t_1, \dots, t_n]\!] \longrightarrow \widehat{A}
\]
(由诸 $x_i$ 定义，其中右端表示 $A$ 按 $J$ 的幂所定义的拓扑的分离完备化 $\varprojlim A/J^{n+1}$) 是同构 (它总是满射).

所有这些事实都是众所周知的，大概几乎完全见于 Serre 的交换代数课程 (由 Gabriel 整理) 中 (那里，在 $A$ 是局部环的情形，还可以找到 $A$-序列的其他刻画).

设 $J$ 是诺特环 $A$ 中的理想. 称 $J$ 为\emph{正则理想} (regular ideal)，若对 $A$ 的每个素理想 $\mathfrak{p}$，$JA_{\mathfrak{p}}$ 有正则生成元组. 显然只需对 $\mathfrak{p} \supset J$ 验证，且可以只限于 $p$ 极大. 更一般地，设 $\mathcal{J}$ 是局部诺特预概形 $X$ 上的理想层，称 $\mathcal{J}$ 为正则理想层，若对每个 $x \in X$，$\mathcal{J}_x$ 是 $\mathcal{O}_x$ 的有正则生成元组的理想. 这等价于下列两个条件的合取：
\begin{enumerate}[a)]
\item 典范满同态
\[
S_{\mathcal{O}_X/\mathcal{J}}(\mathcal{J}/\mathcal{J}^2) \longrightarrow \mathrm{gr}^{\mathcal{J}}(\mathcal{O}_X)
\]
是同构，且
\item $\mathcal{O}_X/\mathcal{J}$-模层 $\mathcal{J}/\mathcal{J}^2$ 局部自由.
\end{enumerate}
此时还说由 $\mathcal{J}$ 定义的 $X$ 的子预概形 $Y$ (即 $\mathcal{O}_Y$ 以 0 延拓后同构于 $\mathcal{O}_X/\mathcal{J}$ 者) 在 $X$ 中正则浸入 (regularly immersed)，并同样 (以显然的方式) 定义正则浸入态射 (morphisme d'immersion régulière) (相应地：在点 $x$ 处正则) 的概念，即把 $Y$ (相应地：$x$ 的适当邻域) 等同于 $X$ 的某个开集中的正则浸入闭子预概形的浸入态射. (不应当说：正则子预概形，因为那将意味着 $Y$ 的局部环正则). 最后，$\mathcal{J}$ 的截面 $x_i$ 称为正则生成元组，若对每个 $x \in X$，$\mathcal{O}_x$ 中相应元素构成 $\mathcal{J}_x$ 的正则生成元组 (该术语与为环的理想生成元引入的术语相容). 这也意味着由诸 $x_i$ 定义的典范满同态
\[
\mathcal{O}_Y[t_1, \dots, t_n] \longrightarrow \mathrm{gr}^{\mathcal{J}}(\mathcal{O}_X)
\]
是同构. 若事先知道理想层 $\mathcal{J}$ 是正则的\index[term]{zhengzelixiang@正则理想}，这也简单意味着在 $Y$ 的每个点 $x$ 处，诸 $x_i$ 定义 $\mathcal{J}/\mathcal{J}^2$ 在 $\mathcal{O}_{Y,x}$ 上的基. (N.B. 若 $Y$ 为空，该条件为空.) 于是，为使 $\mathcal{J}$ 有正则生成元组\index[term]{zhengzeshengchengyuanxi@正则生成元系}，必要且充分条件是 $\mathcal{J}$ 正则，且 $\mathcal{O}_Y$-模层 $\mathcal{J}/\mathcal{J}^2$ 整体自由 (而不只是局部自由)，i.e. 典范同态 $S_{\mathcal{O}_Y}(\mathcal{J}/\mathcal{J}^2) \to \mathrm{gr}^{\mathcal{J}}(\mathcal{O}_X)$ 满射，且 $\mathcal{O}_Y$-模层 $\mathcal{J}/\mathcal{J}^2$ 整体自由.

增广环 (augmented ring) 称为正则的\index[term]{zhengzezengguanghuan@正则增广环}，若增广理想 (idéal d'augmentation) 正则. 于是，若 $A$ 是局部环，看作增广到其剩余域 $k$，则 $A$ 是正则局部环当且仅当它是\emph{正则增广环}.

(老实说，似乎本不必先从环开始做这套初步推论 (sorites)，而最好立即从层开始. 若想要诺特情形的某种东西，这里所采纳的定义 —— 先验地比 Serre 的 $A$-序列定义更宽松 —— 似乎对微分演算的需要更可取. 当然，要做得好的话，还应当在非诺特框架下至少发展一部分光滑态射理论\footnote{如前言所说，这现在已实现，参见 EGA IV 17，18.}，大概从 Jacobi 判则出发，以便尽可能得到光滑态射和\emph{平展态射} (i.e. 光滑且拟有限) 的一切本质形式性质；只有逆命题才用到诺特假设.)

在这些冗长的术语准备之后，一个小定理：

\begin{theorem}\label{thm:II.4.15}
设 $X$ 是局部有限型 $S$-预概形，$Y$ 是由 $X$ 上的凝聚理想层 $\mathcal{J}$ 定义的 $X$ 的闭子预概形，$x$ 是 $X$ 的点. 现在假设 $Y$ 在 $S$ 上于 $x$ 处光滑 (对 $X$ 不作假设). 则下列条件等价：
\begin{enumerate}[(i)]
\item $X$ 在 $S$ 上于 $x$ 处光滑；
\item 浸入 $i \colon Y \to X$ 在 $x$ 处正则，i.e. $\mathcal{J}_x$ 是 $\mathcal{O}_x$ 的正则理想.
\end{enumerate}
\end{theorem}

\begin{corollary}\label{cor:II.4.16}
假设 $Y$ 在 $S$ 上光滑. 为使 $X$ 在 $Y$ 的邻域中 (i.e. 在 $Y$ 的点处) 在 $S$ 上光滑，必要且充分条件是 $Y$ 正则嵌入 $X$，i.e. 浸入 $i \colon Y \to X$ 正则.
\end{corollary}

\paragraph*{证明.} (i) 蕴含 (ii). 应用 \ref{thm:II.4.10} 判则 (ii)，由于 $g \colon X_1 \to X$ 平坦，为证明 $X'$ 的子预概形 $Y'$ 在 $g$ 下的原像正则嵌入，归结为证明 $Y' = S[t_{p+1}, \dots, t_n]$ 正则嵌入 $S[t_1, \dots, t_n]$，这是平凡的 ($t_i$ ($1 \leqslant i \leqslant p$) 构成定义 $X'$ 中 $Y'$ 的理想的正则生成元组).

(ii) 蕴含 (i). 设 $g_i$ ($1 \leqslant i \leqslant p$) 是 $\mathcal{J}_x$ 的正则生成元组，$g_i$ ($p+1 \leqslant i \leqslant n$) 是 $\mathcal{O}_{X,x}$ 的元素，使得它们在 $\mathcal{O}_{Y,x}$ 中的像 $g'_i$ 定义平展态射
\[
Y_1 \longrightarrow Y' = S[t_{p+1}, \dots, t_n]
\]
(把 $Y$ 的邻域 $Y_1$ 映到 $Y'$). 诸 $g_i$ ($1 \leqslant i \leqslant n$) 来自 $x$ 的邻域 $X_1$ 上 $\mathcal{O}_X$ 的 (同名) 截面，且可以假设 $X_1 = X$，$Y_1 = Y$. 由此得到态射
\[
g \colon X \longrightarrow X' = S[t_1, \dots, t_n],
\]
一切归结为证明该态射在 $x$ 处平展. 取 $X_1$ 足够小，可以假设 $g_i$ ($1 \leqslant i \leqslant p$) 在整个 $X$ 上构成 $\mathcal{J}$ 的正则生成元组. 特别地，它们生成 $\mathcal{J}$，故 $X$ 的子预概形 $Y$ 等同于 $g$ 从 $X'$ 的子预概形 $Y'$ 取原像得到的子预概形. 设 $x' = g(x)$，则 $X' \to X$ 在 $x'$ 处的纤维等同于 $Y \to Y'$ 在 $x$ 处的纤维，故在 $k(x')$ 上平展，于是 $g$ 在 $x$ 处非分歧，剩下证明 $g$ 在 $x$ 处平坦. 而按 $\mathcal{J}'_x$ 的幂滤链所得的 $\mathcal{O}_{X',x'}$ 的相伴分次环在各次数都是 $\mathcal{O}_{Y',x'}$ 上的自由环，另一方面按 $\mathcal{J}_x = \mathcal{J}'_x \mathcal{O}_{X,x}$ 的幂滤链所得的 $\mathcal{O}_{X,x}$ 的相伴分次环 (在典范同态下) 同构于前者的 $\mathcal{O}_{Y,x}$ 张量 (因为两个环都是常数环 $\mathcal{O}_{Y',x'}$ (相应地：$\mathcal{O}_{Y,x}$) 上 $n - p$ 个未定元的多项式环)，最后在 $\mathcal{O}_{X',x'}/\mathcal{J}'_{x'} = \mathcal{O}_{Y',x'}$ 上 $\mathcal{O}_{X,x}/\mathcal{J}_x = \mathcal{O}_{Y,x}$ 平坦.

由平坦性的一般判则 (适用于诺特局部环的局部同态 $A' \to A$，$A'$ 带有理想 $J' \neq A'$，使得相伴分次环在 $A'/J'$ 上各维数自由) 推出 $X$ 在 $X'$ 上于 $x$ 处平坦，证毕.

\begin{corollary}\label{cor:II.4.17}
设 $X$ 是在 $Y$ 上局部有限型的预概形，$i$ 是 $X$ 在 $Y$ 上的截面，$y$ 是 $Y$ 的点，$x = i(y)$，$\mathcal{J}$ 是由子预概形 $i(Y)$ 定义的 $X$ 上的理想层 (为简化陈述我们假设它闭，当 $X$ 是概形时此条件成立).

下列条件等价：
\begin{enumerate}[(i)]
\item $X$ 在 $Y$ 上于 $x$ 处光滑；
\item $i$ 在 $y$ 处是正则浸入；
\item $\mathcal{O}_x$ 按 $\mathcal{J}_x$ 的幂所定义的拓扑的完备化 $\mathcal{O}_y$-代数是形式幂级数代数 $\mathcal{O}_y[\![t_1, \dots, t_n]\!]$；
\item[(iii bis)] 存在 $y$ 的开邻域 $U$，使得 $\mathcal{O}_Y$ 上的代数层 $\varprojlim i^{*}\bigl(\mathcal{O}_X/\mathcal{J}^{n+1}\bigr)$ 同构于 $U$ 之上的形如 $\mathcal{O}_Y[\![t_1, \dots, t_n]\!]$ 的层；
\item 存在 $y$ 的开邻域 $U$，$x$ 的开邻域 $V$，以及 $Y$-态射 $g \colon V \to U[t_1, \dots, t_n]$，使得 $g$ 平展，$i$ 诱导 $V$ 在 $U$ 上的截面，经 $g$ 变为 $U[t_1, \dots, t_n]$ 在 $U$ 上的零截面.
\end{enumerate}
\end{corollary}

\begin{proof}
(i) 与 (ii) 的等价性是定理 \ref{thm:II.4.15} 在 $Y = S$ 时的特殊情形；(ii) 与 (iii) 的等价性 (以及实质上 (ii) 与 (iii bis) 的等价性) 已在``回忆''中指出；至于 (i) 与 (iv) 的等价性，容易由定理 \ref{thm:II.4.15} (其条件 (i) 与 (ii) 的等价性) 推出.
\end{proof}

\begin{corollary}\label{cor:II.4.18}
设 $X$ 是在 $S$ 上光滑的预概形. 则对角态射 (diagonal morphism)
\[
\Delta_{X/S} \colon X \longrightarrow X \times_S X
\]
是正则浸入，或如人们所说的，$X$ 在 $S$ 上``微分光滑''\index[term]{weifenguanghua@微分光滑} (differentially smooth).
\end{corollary}

\begin{proof}
事实上，这是推论 \ref{cor:II.4.16} 的特殊情形，因为 $X$ 和 $X \times_S X$ 都在 $S$ 上光滑.
\end{proof}

% 原书编号笔误: 下面这条在原文中与上面的推论同为 ``\ref{cor:II.4.18}''
\begingroup
\addtocounter{theorem}{-1}%
\renewcommand{\thetheorem}{\thesection.18}%
\begin{remark}\label{rem:II.4.18bis}
回忆 (I 1)：若 $X$ 是 $S$ 上的预概形，引入 $X$ 上的拟凝聚代数层 $\mathcal{P}_{X/S}^{n} = \mathcal{O}_{X \times_S X}/\mathcal{I}_X^{n+1}$ (其中 $\mathcal{I}_X$ 表示在 $X \times_S X$ 中定义对角线的理想层)，借助第一投影 $\mathrm{pr}_1 \colon X \times_S X \to X$ 把它看作 $\mathcal{O}_X$-代数层. 诸 $\mathcal{P}_{X/S}^{n}$ 构成 $X$ 上的代数层投射系 (projective system)，其投射极限记为 $\hat{\mathcal{P}}_{X/S}^{\infty}$，它就是 $X \times_S X$ 沿对角线的形式完备化 (formal completion) 的结构层 (现在假设 $X$ 在 $S$ 上局部有限型，故诸 $\mathcal{P}_{X/S}^{n}$ 凝聚). 说 $X$ 在 $S$ 上\emph{微分光滑} (i.e. 对角态射 $\Delta_{X/S}$ 是正则浸入) 也意味着，作为增广到 $\mathcal{O}_X$ 的增广代数层，$\mathcal{P}_{X/S}^{\infty}$ 正则，i.e. $\varOmega_{X/S}^{1}$ 局部自由，且典范满同态
\[
S_{\mathcal{O}_X}\bigl(\varOmega_{X/S}^{1}\bigr) \longrightarrow \mathrm{gr}_{*}\bigl(\mathcal{P}_{X/S}^{\infty}\bigr)
\]
是同构，或最后，$X$ 的每个点都有一个开邻域，在其上增广代数层 $\mathcal{P}_{X/S}^{\infty}$ 同构于层 $\mathcal{O}_X[\![t_1, \dots, t_n]\!]$.

设 $s$ 是 $X$ 在 $S$ 上的截面，$\mathcal{J}$ 是它所定义的 $X$ 上的理想层 (为简化假设 $s(S)$ 闭)，则有 $\mathcal{O}_X$-增广代数层的典范同构：
\begin{equation}
s^{*}\bigl(\mathcal{P}_{X/S}^{n}\bigr) = \mathcal{O}_X/\mathcal{J}^{n+1}, \qquad s^{*}\bigl(\mathcal{P}_{X/S}^{\infty}\bigr) = \varprojlim_n \mathcal{O}_X/\mathcal{J}^{n+1}. \tag{4.4}
\label{eq:II.4.4}
\end{equation}
这些同构对基变换在显然的意义下是函子性的，并且 (考虑到这一点) 重新给出 $S$ 上代数层 $\mathcal{P}_{X/S}^{n}$ 的一个刻画. 例如若 $S = \mathrm{Spec}(k)$，$k$ 是域，则 $X$ 在 $S$ 上的截面 $s$ 的给出等价于 $X$ 在 $k$ 上有理点 $x$ 的给出，而前面的公式意味着有 $k$-代数同构
\begin{equation}
\mathcal{P}_{X/S}^{n}(x) = \mathcal{O}_x/\mathfrak{m}_x^{n+1}, \tag{4.5}
\label{eq:II.4.5}
\end{equation}
这证明``$X$ 相对 $S$ 的 $n$ 阶\emph{主部层} (faisceau des parties principales d'ordre n)\index[term]{zhubuceng@主部层}''这一名称之合理，该名称是给 $\mathcal{P}_{X/S}^{n}$ 的. 还由 \eqref{eq:II.4.4} 看到，若 $X$ 在 $s(S)$ 的所有点处微分光滑，则 $X$ 在 $s(S)$ 的所有点处光滑 (推论 \ref{cor:II.4.17})，逆命题也成立 (推论 \ref{cor:II.4.18}). 结合 \ref{cor:II.4.13}，容易由此推出：若 $X$ 是局部有限型 $S$-预概形，则 $X$ 在 $S$ 上光滑当且仅当它在 $S$ 上平坦且微分光滑. (N.B. 平坦性假设是本质的，取 $X$ 为 $S$ 的闭子预概形即可看出.)

再作为回忆指出，借助投影 $\mathrm{pr}_2 \colon X \times_S X \to X$ 在 $\mathcal{P}_{X/S}^{n}$ 上得到第二种代数结构，它由第一种结构经 $X \times_S X$ 的对称自同构诱导的环层 $\mathcal{P}_{X/S}^{n}$ 的典范对合 (involution) 得出. 用 $d_{X/S}^{n}$ 或简称 $d^n$ 记环层同态
\begin{equation}
d_{X/S}^{n} \colon \mathcal{O}_X \longrightarrow \mathcal{P}_{X/S}^{n}, \tag{4.6}
\label{eq:II.4.6}
\end{equation}
它对应于这第二种代数结构. 考虑到同构 \eqref{eq:II.4.4}，该同态把 $\mathcal{O}_X$ 的截面 $f$ 变成 $\mathcal{P}_{X/S}^{n}$ 的截面 $d^n(f)$，后者经 $X$ 在 $S$ 上的截面 $s$ 取原像后等同于 $f$ 在 $\Gamma(X, \mathcal{O}_X/\mathcal{J}^{n+1})$ 中的典范像. 这证明给 $d^n f$ 的``$f$ 的 $n$ 阶主部组 (système des parties principales d'ordre n)''名称之合理，特别是在 $S = \mathrm{Spec}(k)$ 的情形，如公式 \eqref{eq:II.4.5} 所考虑的.

最后注意到，同态 \eqref{eq:II.4.6} 可以看作 (相对常数预概形 $S$ 的) $\mathcal{O}_X$ 上的万有 $\leqslant n$ 阶\emph{微分算子} (opérateur différentiel d'ordre $\leqslant n$)\index[term]{weifensuanzi@微分算子}\footnote{关于本段的全部内容，可参阅 EGA IV 16.8 至 16.12.}，这里约定把 $\mathcal{O}_X$ 到模层 $F$ 的 $\leqslant n$ 阶微分算子称为可分解为
\[
D \colon \mathcal{O}_X \xrightarrow{d^n} \mathcal{P}_{X/S}^{n} \xrightarrow{u} F
\]
的层同态 $D$，其中 $u$ 是 $\mathcal{O}_X$-模层同态，且由 $D$ 唯一确定. 这个定义与直观的递归定义一致 ($D$ 是 $\leqslant n$ 阶微分算子，若对 $X$ 的开集 $U$ 上 $\mathcal{O}_X$ 的每个截面 $g$，$f \mapsto D(fg) - D(f)$ 是 $U$ 上的 $\leqslant n-1$ 阶微分算子). 由此推出，若 $X$ 在 $S$ 上微分光滑，则所有阶微分算子构成的环层具有微分流形 (differentiable manifold) 上微分演算中众所周知的简单结构，特别在局部有 $\mathcal{O}_X$-模层基，由可交换算子 $\delta/\delta x_i$ ($1 \leqslant i \leqslant n$) 的分次幂 (divided powers) 构成.

若 $S$ 是 $\mathbf{Q}$-代数层 ($\mathbf{Q}$ 为有理数域)，只需取诸 $\delta/\delta x_i$ 的通常多项式. 此时且仅在此时，极为特殊地，为使 $X$ 在 $S$ 上微分光滑，只需 $\varOmega_{X/S}^{1}$ 局部自由.
\end{remark}
\endgroup

\begin{remark}\label{rem:II.4.19}
本节引入的``正则浸入''、``正则理想''等术语遇到了相当激烈而普遍的反对 (Chevalley，Serre). 有人提议``Cohen-Macaulay 理想''、``Macaulay 理想''或``macaulay 理想''，这从道义上也将迫使我们采用``Cohen-Macaulay 浸入''或``Macaulay 浸入''. 然而这一术语与大全未来稿中已经使用的另一术语冲突，那里一个 (有限型) 态射称为在一点``Cohen-Macaulay''，若它在这一点平坦，且经过该点的纤维在该点的局部环是 Cohen-Macaulay 环. 在找到满意的方案之前，我们暂且保留本节引入的术语\footnote{这就是 EGA 0$_{\mathrm{IV}}$ 15.1.7 所采纳的.}.
\end{remark}

\section{基域的情形}

\begin{proposition}\label{prop:II.5.1}
设 $k$ 是域，$X$ 是 $k$ 上的有限型预概形，$x$ 是 $X$ 的点，$n$ 是 $X$ 在 $x$ 处的维数，$f \colon X \to \mathrm{Spec}\, k[t_1, \dots, t_n] = Y$ 是由 $\Gamma(X, \mathcal{O}_X)$ 的元素 $f_i$ 定义的态射. 下列条件等价 (且蕴含 $X$ 在 $k$ 上于 $x$ 处光滑，从而由 \ref{prop:II.3.1} 在 $x$ 处正则)：
\begin{enumerate}[(i)]
\item $f$ 在 $x$ 处平展.
\item $df_i$ 构成 $\varOmega_{X/k}^{1}$ 在 $x$ 处的基.
\item $df_i$ 在 $x$ 处生成 $\varOmega_{X/k}^{1}$.
\end{enumerate}
\end{proposition}

\begin{proof}
由于 (i) 蕴含 $X$ 在 $k$ 上于 $x$ 处光滑，(i)$\Rightarrow$(ii) 是 \ref{cor:II.4.8} 的特殊情形，(ii)$\Rightarrow$(iii) 平凡，剩下证明 (iii)$\Rightarrow$(i). 若 (iii) 成立，由引理 \ref{lem:II.4.1}，$f$ 在 $x$ 处非分歧，故 (把 $x$ 换成开邻域) 拟有限，从而由维数理由支配. 由于 $Y$ 正则，由 (\ref{thm:I.9.5} (ii)) 或 (\ref{cor:I.9.11}) 得出 $f$ 平展.
\end{proof}

\begin{corollary}\label{cor:II.5.2}
在 \ref{prop:II.5.1} 的预备条件下，假设 $k(x)$ 是 $k$ 的有限可分扩张，且 $f_i$ ($1 \leqslant i \leqslant n$) 定义 $\mathfrak{m}_x$ 的元素. 则前面的条件等价于
\begin{enumerate}[(i)]
\item[(iv)] $f_i$ 构成 $\mathfrak{m}_x$ 的生成元组 (或：$f_i$ 模 $\mathfrak{m}_x^2$ 构成 $k(x)$ 上的基 $\mathfrak{m}_x/\mathfrak{m}_x^2$).
\end{enumerate}
\end{corollary}

\begin{proof}
事实上，(iv)$\Rightarrow$(iii) 由正合序列
\begin{equation}
\mathfrak{m}_x/\mathfrak{m}_x^2 \longrightarrow \varOmega_{\mathcal{O}_x/k}^{1} \longrightarrow \varOmega_{k(x)/k}^{1} \longrightarrow 0 \tag{5.1}
\label{eq:II.5.1}
\end{equation}
得出，并考虑到 $k(x)$ 在 $k$ 上平展故 $\varOmega_{k(x)/k}^{1} = 0$. 另一方面 (ii) 蕴含 (iv)，因为 $X$ 和 $\mathrm{Spec}\bigl(k(x)\bigr)$ 都在 $k$ 上于 $x$ 处光滑，由 \ref{thm:II.4.10} (iv) 可以在上述正合序列左边加一个 0.
\end{proof}

\begin{corollary}\label{cor:II.5.3}
设 $x$ 是 (在 $k$ 上有限型的) $X$ 的点. 若 $X$ 在 $k$ 上于 $x$ 处光滑，则 $\mathcal{O}_x$ 正则，逆命题在 $k(x)$ 是 $k$ 的有限可分扩张时成立.
\end{corollary}

\begin{proof}
事实上，逆命题由 \ref{cor:II.5.2} 得出，取 $\mathfrak{m}_x$ 的正则生成元组 ($f_i$) (N.B. 除 \ref{cor:II.5.2} 外，也可以援引定理 \ref{thm:II.4.15}). 于是：
\end{proof}

\begin{proposition}\label{prop:II.5.4}
设 $X$ 是 $k$ 上的有限型预概形. 若 $X$ 在 $k$ 上光滑，则它正则，逆命题在 $k$ 完美 (perfect) 时成立.
\end{proposition}

\begin{proof}
对逆命题，注意由 \ref{cor:II.5.3}，$X$ 在 $k$ 上于每个闭点光滑，从而处处光滑，因为光滑的点集是开的.
\end{proof}

\begin{theorem}\label{thm:II.5.5}
设 $X$ 是 $k$ 上的有限型预概形，$x$ 是 $X$ 的点，$n$ 是 $X$ 在 $x$ 处的维数，$k'$ 是 $k$ 的完美扩张. 下列条件等价：
\begin{enumerate}[(i)]
\item $X$ 在 $k$ 上于 $x$ 处光滑.
\item $\varOmega_{X/k}^{1}$ 在 $x$ 处是秩为 $n$ 的自由模.
\item[(ii bis)] $\varOmega_{X/k}^{1}$ 在 $x$ 处由 $n$ 个元素生成.
\item $X$ 在 $k$ 上于 $x$ 处微分光滑.
\item 存在 $x$ 的开邻域 $U$，使得 $U \otimes_k k'$ 正则 (i.e. 其各点的局部环正则).
\end{enumerate}
\end{theorem}

\begin{proof}
由 \ref{thm:II.4.3} (ii) 有 (i)$\Rightarrow$(ii)，(ii)$\Rightarrow$(ii bis) 平凡，(ii bis)$\Rightarrow$(i) 由 \ref{prop:II.5.1} 得出. 由于 $X$ 在 $k$ 上平坦，由 \ref{cor:II.4.18} 有 (i)$\Leftrightarrow$(iii). 有 (i)$\Rightarrow$(iv)，因为光滑对基扩张不变且蕴含正则，而 (iv)$\Rightarrow$(i) 因为由命题 \ref{prop:II.5.4}，$U \otimes_k k'$ 在 $k'$ 上简单，故由 \ref{cor:II.4.13}，$U$ 在 $k$ 上简单.
\end{proof}

对 $x$ 取 $X$ 的生成点 (假设 $X$ 不可约)，得到：

\begin{corollary}\label{cor:II.5.6}
设 $K$ 是域 $k$ 上有限型代数的局部化所成的 Artin 局部环 (例如 $K$ 可以是 $k$ 的有限型扩张)，设 $n$ 是其剩余域在 $K$ 上的超越次数. 下列条件等价：
\begin{enumerate}[(i)]
\item $K$ 是 $k$ 的纯超越扩张 (purely transcendental extension) $k(t_1, \dots, t_n)$ 的有限可分扩张.
\item $\varOmega_{K/k}^{1}$ 是秩为 $n$ 的自由 $K$-模.
\item[(ii bis)] $\varOmega_{K/k}^{1}$ 是可由 $n$ 个元素生成的 $K$-模.
\item $K \otimes_k K$ 按增广理想 $K \otimes_k K \to K$ 的幂所定义的拓扑的完备化 $\mathcal{O}'$ 是``正则''增广 $K$-代数，i.e. 同构于 $K$ 上的形式幂级数代数 (N.B. 若 $K$ 是域，这等价于说 $\mathcal{O}'$ 是正则局部环).
\item $K$ 是 $k$ 的可分扩张.
\end{enumerate}
\end{corollary}

\begin{proof}
事实上，总可以把 $K$ 看作 $k$ 上不可约有限型概形 $X$ 的生成点的局部环，而所考虑的条件正是 \ref{thm:II.5.5} 中同名条件，在 (iv) 中对 $k'$ 取 $k$ 的代数闭包. 只有``$K$ 在 $k$ 上可分 $\Rightarrow$ $X$ 在 $k$ 上于 $x$ 处光滑''这一蕴含需要证明. 借助 \ref{cor:II.4.13} 立即归结到基域为 $k'$ 从而代数闭的情形，此时存在 $X$ 在 $k$ 上有理点 $a$. 但此时由 \ref{prop:II.5.4}，$X$ 在 $k$ 上于 $a$ 处光滑，更不用说它在 $k$ 上于生成点 $x$ 处光滑，证毕\footnote{注意对 $k$ 所作的假设是不可缺少的，如下例所示：$k = \mathbf{R}$，$X = \mathrm{Spec}(k)$，$Y = \mathrm{Spec}\bigl(k[t]/(t^2+1)\bigr)$. 则 $Y$ 是 $X$ 的平展覆盖但不平凡，尽管 $k$ 在 $k(t)$ 中代数闭.}.
\end{proof}

注意，当 $K$ 是 $k$ 的有限型扩张时，(i)，(ii)，(ii bis)，(iv) 的等价性是众所周知的，但我们没有使用这些已知等价性中的任何一个. 当然，命题 \ref{prop:II.5.1} 包含作为特殊情形：元素组 $x_i$ ($1 \leqslant i \leqslant n$) 是 $K$ 在 $k$ 上的``可分超越基'' (separating transcendence base)，当且仅当 $dx_i$ 构成 $K$-模 $\varOmega_{K/k}^{1}$ 的基，这事实别处也都知道.

\begin{corollary}\label{cor:II.5.7}
设 $X$ 是域 $k$ 上的有限型预概形. 为使 $X$ 在 $k$ 上光滑，必要且充分条件是 $\varOmega_{X/k}^{1}$ 局部自由，且 $X$ 的不可约分支的生成点的局部环是 $k$ 的可分扩张 (最后一个条件在 $k$ 完美且 $X$ 约化时自动满足).
\end{corollary}

\begin{proof}
可以假设 $X$ 连通，设 $n$ 是假设局部自由的 $\varOmega_{X/k}^{1}$ 的秩. 由假设和 \ref{cor:II.5.6}，它也是 $X$ 的生成点的局部环所定义的 $k$ 的扩张的超越次数，故 $X$ 的所有不可约分支维数为 $n$. 于是由 \ref{thm:II.5.5} 得出结论.
\end{proof}

注意，若 $K$ 是 $k$ 的有限 (不必可分的) 扩张，则 $\varOmega_{K/k}^{1}$ 是自由 $k$-模，故令 $X = \mathrm{Spec}(K)$，$\varOmega_{X/k}^{1}$ 是局部自由层，且 $X$ 约化，而 $X$ 未必在 $k$ 上光滑. 再把标量扩张到 $k$ 的代数闭包，得到类似的例子，其中 $k$ 代数闭，但 $X$ 却不再约化.

\begin{corollary}\label{cor:II.5.8}
设 $X$ 是域 $k$ 上的有限型预概形，$x$ 是 $X$ 的点，$n$ 是 $X$ 在 $x$ 处的维数，$p$ 是 $\mathcal{O}_x$ 的维数，i.e. $x$ 在 $X$ 中的闭包 $Y$ 的余维数；故 $n - p$ 是 $k(x)$ 在 $k$ 上的超越次数. 设 $f_i$ ($1 \leqslant i \leqslant n$) 是 $\mathcal{O}_x$ 的元素，对 $1 \leqslant i \leqslant p$ 有 $f_i \in \mathfrak{m}_x$. 下列条件等价：
\begin{enumerate}[(i)]
\item 由诸 $f_i$ 定义的 $x$ 处的态射芽
\[
X \longrightarrow \mathrm{Spec}\bigl(k[t_1, \dots, t_n]\bigr)
\]
在 $x$ 处平展.
\item $f_i$ ($1 \leqslant i \leqslant p$) 生成 $\mathfrak{m}_x$，i.e. 构成 $\mathcal{O}_x$ 的正则参数系 (regular system of parameters)，且 $f_j$ ($p+1 \leqslant j \leqslant n$) 在 $k(x)$ 中的类构成可分超越基 (i.e. $d\overline{f}_j$ ($p+1 \leqslant j \leqslant n$) 构成 $\varOmega_{k(x)/k}^{1}$ 的基，或生成 $\varOmega_{k(x)/k}^{1}$).
\end{enumerate}
\end{corollary}

\begin{proof}
假设 (i) 成立. 由此得出 $df_i(x)$ 构成 $\varOmega_{X/k}^{1}(x)$ 的基 (\cref{cor:II.4.8})，故它们在 $\varOmega_{k(x)/k}^{1}$ 中的像 $d\overline{f}_i(x)$ 生成该 $k$-向量空间. 由于对 $1 \leqslant i \leqslant p$，$\overline{f}_i$ 为零，只需取 $p+1 \leqslant i \leqslant n$ 的 $d\overline{f}_i(x)$. 由于 $k(x)$ 在 $k$ 上的超越次数为 $n - p$，由推论 \ref{cor:II.5.8} 判则 (iii) (应用于 $K = k(x)$) 得出 $Y$ 在 $k$ 上于其生成点 $x$ 处光滑，且 $d\overline{f}_i(x)$ ($p+1 \leqslant i \leqslant n$) 构成 $\varOmega_{k(x)/k}^{1}$ 在 $k(x)$ 上的基. 于是 \ref{prop:II.4.9} 的条件 (ii) 成立，从而条件 (iii) 也成立，特别地 $f_i$ ($1 \leqslant i \leqslant p$) 构成 $\mathfrak{m}_x$ 的生成元组. 由于 $\mathcal{O}_x$ 的维数为 $p$，它们构成 $x$ 处的正则参数系. 这证明 (ii).
\end{proof}

假设 (ii) 成立. 由正合序列 \eqref{eq:II.5.1}，由此得出 $df_i(x)$ 生成 $\varOmega_{X/k}^{1}$，从而由命题 \ref{prop:II.5.1} 得 (i).

\begin{corollary}\label{cor:II.5.9}
设 $X$ 是域 $k$ 上的有限型预概形，$x$ 是 $X$ 的点，$n$ 是 $X$ 在 $x$ 处的维数，$p$ 是 $\mathcal{O}_x$ 的维数，i.e. $x$ 在 $X$ 中的闭包 $Y$ 的余维数，故 $n - p$ 是 $k(x)$ 在 $k$ 上的超越次数. 下列条件等价：
\begin{enumerate}[(i)]
\item $\mathcal{O}_x$ 正则且 $k(x)$ 是 $k$ 的可分扩张.
\item $X$ 在 $k$ 上于 $x$ 处光滑，且典范同态
\[
\mathfrak{m}_x/\mathfrak{m}_x^2 \longrightarrow \varOmega_{\mathcal{O}_x/k}^{1} \otimes_{\mathcal{O}_x} k(x) = \varOmega_{X/k}^{1}(x)
\]
单射.
\item 存在 $f_i \in \mathcal{O}_x$ ($1 \leqslant i \leqslant n$)，对 $1 \leqslant i \leqslant p$ 有 $f_i \in \mathfrak{m}_x$，使得由诸 $f_i$ 定义的 $X$ 到 $\mathrm{Spec}\bigl(k[t_1, \dots, t_n]\bigr)$ 的 $x$ 处态射芽在 $x$ 处平展 (i.e. 由 \ref{prop:II.5.1}，使得 $df_i(x)$ 生成 $\varOmega_{X/k}^{1}(x)$).
\item 存在 $f_i \in \mathcal{O}_x$ ($1 \leqslant i \leqslant n$)，使得 $f_i$ ($1 \leqslant i \leqslant p$) 生成 $\mathfrak{m}_x$，且 $df_j(x)$ ($p+1 \leqslant j \leqslant n$) 在 $k(x)$ 上生成 $\varOmega_{k(x)/k}^{1}$.
\end{enumerate}
\end{corollary}

\begin{proof}
(iii) 与 (iv) 的等价性由推论 \ref{cor:II.5.8} 得出，且由 \ref{prop:II.4.9}，这些条件还等价于``$X$ 在 $k$ 上于 $x$ 处光滑且 \ref{thm:II.4.10} 的条件 (ii) 成立''. 故它们等价于``$X$ 在 $k$ 上于 $x$ 处光滑且 \ref{thm:II.4.10} 的条件 (iv) 成立''，从而等价于 \ref{cor:II.4.8} (ii). 或等价于``$X$ 在 $k$ 上于 $x$ 处光滑且 \ref{thm:II.4.10} 的条件 (i) 成立''，在这里它简单意味着 $k(x)$ 在 $k$ 上可分. 这蕴含 \ref{cor:II.5.9} (i)，剩下证明 \ref{cor:II.5.9} (i) 蕴含它，i.e. 证明：
\end{proof}

\begin{corollary}\label{cor:II.5.10}
设 $x$ 是域 $k$ 上有限型预概形的点，使得 $k(x)$ 在 $k$ 上可分. 为使 $X$ 在 $k$ 上于 $x$ 处光滑，必要且充分条件是它在 $x$ 处正则 (i.e. 局部环 $\mathcal{O}_x$ 正则).
\end{corollary}

\begin{proof}
事实上，若如此，显然可以找到 $f_i \in \mathcal{O}_x$ ($1 \leqslant i \leqslant n$) 满足条件 \ref{cor:II.5.9} (iv).
\end{proof}

\begin{remarkstar}
\textbf{勘误 (Errata).} 在本节 \ref{cor:II.5.6} 的证明中，我们使用了``代数闭域上非空的有限型约化概形至少有一个正则 (从而光滑) 点''这一事实，该事实通常用微分方法证明 (借助 Zariski 的``$X$ 的正则点集是开的''这一定理). 若要避免循环论证，必须证明：若 $K/k$ 是有限型可分扩张，且 $f_i \in K$ 使得 $d_{K/k} f_i$ ($1 \leqslant i \leqslant n$) 构成 $\varOmega_{K/k}^{1}$ 的基，则 $n$ 是 $K$ 在 $k$ 上的超越次数，i.e. 诸 $f_i$ 代数无关 (algebraically independent). 用 Mac Lane 判则证明这一事实是众所周知的，参见 Bourbaki，Algèbre，Chap. V，§9，th. 2：取次数最小的多项式 $g \in k[t_1, \dots, t_n]$ 使得 $g(f_1, \dots, f_n) = 0$，于是有
\[
\sum_i \frac{dg}{dt_i}(f_1, \dots, f_n)\, df_i = 0,
\]
从而 (因为 $df_i$ 构成 $\varOmega_{K/k}^{1}$ 的基) $dg/dt_i$ 在 $(f_1, \dots, f_n)$ 处取值为零，故由 $g$ 的极小性它们为零，于是若 $k$ 特征为 0 则 $g = 0$，若 $k$ 特征 $p \neq 0$ 则 $g = h(t_1^p, \dots, t_n^p)$. 用 Mac Lane 判则，多项式 $h \in k[t_1, \dots, t_n]$ 也零化 $(f_1, \dots, f_n)$，从而仍由 $g$ 的极小性得 $g = 0$.
\end{remarkstar}

\end{refsection}
\begin{refsection}% =====================================================================
% 第三讲 (Exposé III) 光滑态射: 延拓性质
% (Morphismes lisses: propriétés de prolongement)
% A. Grothendieck
% 原书页码: pp. 49-73
% =====================================================================
\chapter{光滑态射：延拓性质}

\section{形式光滑同态}

在讲 II 中，我们只限于考虑有限型同态，因此在局部环的局部同态 $A \to B$ 中，只限于 $B$ 同构于 $A$ 上有限型代数的局部化的情形. 这个情形对多种应用 (特别是形式几何 (formal geometry) 或解析几何 (analytic geometry)) 是不够的. 例如，形式幂级数环 $B = A[\![t_1, \dots, t_n]\!]$ 具有 (从形式几何的观点看) $A$ 上光滑代数的性质. 在解析几何中，乘积 $Y \times \mathbf{C}^n$ 的一个点 $(y, z)$ 的局部环亦然，它被看作 $y$ 的局部环上的代数；而且该代数的完备化同构于基环 $\mathcal{O}_x$ 的完备化上的 $n$ 个未定元的形式幂级数代数. 正是这一点导致我们给出下面的定义.

\begin{definition}\label{def:III.1.1}
设 $u \colon A \to B$ 是局部环 (回忆它们是诺特的) 的局部同态. 假设 $k(B)$ 在 $k(A)$ 上有限. 称 $u$ 是形式光滑 (formally smooth) 同态，或称代数 $B$ 在 $A$ 上形式光滑，若存在 $A$ 上的有限局部 $A$-代数 $A'$，在 $A$ 上自由，使得半局部环 $\overline{B} \otimes_{\overline{A}} A' = B'$ 的局部分量作为 $A'$-代数同构于 $A'$ 上的形式幂级数代数\footnote{关于由下面 \ref{thm:III.2.1} 启发的更一般、更概念化的定义，参见 EGA 0$_{\mathrm{IV}}$ 19.3.1.}.
\end{definition}

(用 $\overline{A}$，$\overline{B}$ 记 $A$，$B$ 的完备化.) 由于 $B'$ 在 $\overline{B}$ 上有限且自由，它确是一个半局部环，是完备局部环的直积，其中每个局部环仍是 $\overline{B}$ 上的自由模，故维数与 $\overline{B}$ 相同，从而与 $B$ 相同. 由此得出，\ref{def:III.1.1} 中所考虑的形式幂级数环中变元 $t_i$ 的个数等于 $\dim B - \dim A = \dim \overline{B} - \dim \overline{A}$，特别地与所选择的局部分量无关. 立即看出，它也就是环 $B \otimes k = B/\mathfrak{m}B$ 的维数，其中 $k = A/\mathfrak{m}$ 是 $A$ 的剩余域；称它为 $B$ 相对 $A$ 的\emph{相对维数} (relative dimension).

\begin{remark}\label{rem:III.1.2}
显然，定义 \ref{def:III.1.1} 只依赖于由 $A \to B$ 推出的完备化上的同态 $\overline{A} \to \overline{B}$，这在一定程度上证明了该术语的合理性. 我们在此对定义 \ref{def:I.3.2} b) 和 \ref{def:I.4.1} b) 感到遗憾，它们有误导之嫌，并宁愿在那些定义所考虑的情形说``\emph{形式非分歧}\index[term]{xingshifeifenqi@形式非分歧}'' (formally unramified) 和``\emph{形式平展}\index[term]{xingshipingzhan@形式平展}'' (formally etale)，把``非分歧''和``平展''的术语保留给 $B$ 是 $A$ 上有限型代数的局部化的情形\footnote{或更好地说``本质上非分歧'' (essentially unramified)，相应地``本质上平展'' (essentially etale)，比较 EGA IV 18.6.1.}. 读者立即验证``形式平展''等价于``形式光滑且拟有限''. 最后指出，对剩余扩张 $k(B)/k(A)$ (这里假设有限) 不作任何预先假设，也存在``形式光滑''的合理定义，它尤其涵盖这样的局部同态 $A \to B$：$B$ 在 $A$ 上平坦且 $B/\mathfrak{m}B$ 是 $A/\mathfrak{m} = k$ 的可分扩张 (不必有限型)；例如，Cohen $p$-环在 $p$-adic 整数环上形式光滑. 在这个一般情形，同态的提升性质 (比较 \ref{thm:III.2.1}) 应当成为定义. 对我们所考虑的应用，定义 \ref{def:III.1.1} 所处理的情形已经足够；在下面的``形式光滑''中，我们隐含``剩余扩张有限''这一条件.
\end{remark}

\begin{lemma}\label{lem:III.1.3}
若 $B$ 在 $A$ 上形式光滑，则 $B$ 在 $A$ 上平坦.
\end{lemma}

\begin{proof}
由于平坦性对完备化不变，可以假设 $A$ 和 $B$ 完备. 由于平坦性对基环的平坦 (从而忠实平坦) 局部扩张不变，由定义 \ref{def:III.1.1} 归结到 $B$ 是 $A$ 上的形式幂级数代数的情形. 但此时作为 $A$-模，$B$ 同构于若干同构于 $A$ 的 $A$-模的乘积，故 (基环 $A$ 诺特) 作为平坦 $A$-模的乘积是 $A$-平坦的.
\end{proof}

设我们处于 \ref{def:III.1.1} 的条件下. 由于 $B'$ 的局部分量在 $A'$ 上的剩余扩张平凡，由此得出 $L \otimes_k k'$ 是 $k'$-代数，是 Artin 的，其局部分量的剩余扩张平凡 (其中 $L$，$k$，$k'$ 是 $A$，$B$，$A'$ 的剩余域). 这个有限自由扩张 $A'$ 满足 \ref{def:III.1.1} 所述条件的必要条件也是充分的，这由下面的 \ref{prop:III.1.4} (i) 和 \ref{prop:III.1.5} 立即得出.

\begin{proposition}\label{prop:III.1.4}
设 $A \to B$ 是剩余扩张有限的局部环局部同态，$A'$ 是 $A$ 上的有限局部 $A$-代数，于是 $B' = B \otimes_A A'$ 在 $B$ 上有限，从而也是诺特半局部环. (i) 若 $B$ 在 $A$ 上形式光滑，则 $B$ 在其极大理想处的局部化在 $A'$ 上形式光滑. (ii) 若 $A'$ 在 $A$ 上自由，则逆命题成立.
\end{proposition}

\begin{proof}
立即归结到 $A$，$B$ 完备的情形.
\end{proof}

\begin{enumerate}[(i)]
\item 设 $A''$ 是 $A$ 上的有限自由局部扩张，使得 $B'' = B \otimes_A A''$ 的局部分量是 $A''$ 上的形式幂级数代数. 作标量扩张 $A'' \to A'' \otimes_A A' \to A'''$，其中 $A'''$ 是 $A'' \otimes_A A'$ 的一个局部分量，看到 $B'' \otimes_{A''} A''' = B \otimes_A A'''$ 的局部分量是 $A'''$ 上的形式幂级数代数. 另一方面还有 $B \otimes_A A''' = (B \otimes_A A') \otimes_{A'} A''' = B' \otimes_{A'} A'''$；又由于 $A''$ 在 $A$ 上自由，$A'' \otimes_A A'$ 在 $A'$ 上自由，作为其直因子的 $A'''$ 亦然，这证明 $B'$ 在 $A'$ 上形式光滑.
\item 设 $A''$ 是 $A'$ 上的有限自由局部代数，使得 $B' \otimes_{A'} A'' = B \otimes_A A''$ 的局部分量是 $A''$ 上的形式幂级数代数. 由于 $A'$ 在 $A$ 上自由，$A''$ 也自由，故 $B$ 在 $A$ 上形式光滑.
\end{enumerate}

\begin{proposition}\label{prop:III.1.5}
设 $A \to B$ 是剩余扩张平凡的局部环局部同态. 为使 $B$ 在 $A$ 上形式光滑，必要且充分条件是 $\overline{B}$ 同构于 $\overline{A}$ 上的形式幂级数代数.
\end{proposition}

\begin{proof}
只需证明必要性，且可以假设 $A$ 和 $B$ 完备. 设 $m$ ($n$) 是 $A$ ($B$) 的极大理想，$t_1, \dots, t_n$ 是 $n$ 的元素，定义向量空间
\[
(\mathfrak{n}/\mathfrak{n}^2)/\operatorname{Im}(\mathfrak{m}/\mathfrak{m}^2) = \mathfrak{n}/(\mathfrak{n}^2 + \mathfrak{m}B)
\]
的基. 这些元素因此定义局部 $A$-代数同态
\[
B_1 = A[\![t_1, \dots, t_n]\!] \longrightarrow B,
\]
证明它是同构. 只需证明对 $m$ 的每个幂 $m^q$，模 $\mathfrak{m}^q$ 约化后得到同构 ($B_1$ 和 $B$ 是相应约化环按变动的 $q$ 的投射极限). 由于 $B$ 和 $B_1$ 是 $A$-\emph{平坦模}，故按 $m$-adic 滤链所得的相伴分次环在 $k = A/\mathfrak{m}$ 上由 $\mathrm{gr}(A)$ 张量得到，归结为证明 $B_1/\mathfrak{m}B_1 \to B/\mathfrak{m}B$ 是同构. 考虑到 \ref{lem:III.1.3}，归结到 $A$ 是域 $k$ 的情形. 另一方面，若 $A'$ 是 $A$ 上的有限自由局部代数，使得 $B \otimes_A A'$ 是 $A'$ 上的形式幂级数代数 (N.B. 该代数是局部的，因为 $B$ 在 $A$ 上的剩余扩张平凡)，为证明 $B_1 \to B$ 是同构，只需证明 $B_1 \otimes_A A' \to B \otimes_A A'$ 是同构. 这归结到 $B$ 已经是形式幂级数代数的情形 (本应先作这个约化，再归结到基域是域的情形). 但此时 $B$ 是带代表元域 (field of representatives) $k$ 的正则局部环，而众所周知 (并且由考察 $B_1$ 和 $B$ 上 $\mathfrak{n}_1$-adic 和 $n$-adic 滤链的相伴分次环立即看出) $B_1 \to B$ 是同构，证明完毕.
\end{proof}

\begin{corollary}\label{cor:III.1.6}
若 $B$ 在 $A$ 上形式光滑，则存在 $A$ 上的有限局部 $A$-代数 $A'$，使得 $\overline{B} \otimes_{\overline{A}} \overline{A'} = \overline{(B \otimes_A A')}$ 的局部分量同构于 $\overline{A'}$ 上的形式幂级数代数.
\end{corollary}

\begin{proof}
事实上，若 $L/k$ 是 $B/A$ 的剩余扩张，考虑扩张 $k'/k$，使得 $k'$-代数 $L \otimes_k k'$ 中的剩余扩张平凡. 设 $A'$ 是 $A$ 上的有限自由代数，使得 $A'/\mathfrak{m}A' = k'$ (知道它存在：例如逐步归结到 $k'/k$ 单生成的情形，此时把 $k'$ 在 $k$ 上的生成元的最小多项式的系数提升到 $A$). 它是局部的. 则 $B \otimes_A A'$ 在其极大理想处对 $A'$ 的剩余扩张 $k'$ 的剩余扩张平凡，用 \ref{prop:III.1.5} 完成证明.
\end{proof}

\begin{corollary}\label{cor:III.1.7}
设 $A \to B$ 是局部环的局部同态. 为使 $B$ 在 $A$ 上形式光滑，必要且充分条件是 $B$ 在 $A$ 上平坦且 $B/\mathfrak{m}B$ 在 $k = A/\mathfrak{m}$ 上形式光滑.
\end{corollary}

\begin{proof}
对 $A$ 作适当的有限自由局部扩张 $A'$，用 \ref{prop:III.1.4} (ii)，归结到 $B/A$ 的剩余扩张平凡的情形. 此外，由 \ref{prop:III.1.4} (i) 和 \ref{lem:III.1.3} 已知所述条件是必要的. 对充分性，只需注意 \ref{prop:III.1.5} 的证明在所作的假设下 (假设 $A$ 和 $B$ 完备，这是允许的) 证明 $B$ 是 $A$ 上的形式幂级数代数.
\end{proof}

\begin{remark}\label{rem:III.1.8}
对形式光滑同态，不难发展讲 II 中所研究的光滑态射的一切性质的类似物. 对微分性质，这需要修改 Kähler 微分的通常定义 (参见 I 1)，用完备张量积代替通常张量积. 我们只在此提及这些深渊，因为前面的内容对我们的目的已经足够.
\end{remark}

还要在形式光滑性与讲 II 中发展 (且我们至今尚未使用) 的光滑性概念之间建立联系：

\begin{proposition}\label{prop:III.1.9}
设 $A \to B$ 是局部同态，$B$ 是 $A$ 上有限型代数的局部化. 为使 $B$ 在 $A$ 上光滑，必要且充分条件是它在 $A$ 上形式光滑.
\end{proposition}

\begin{proof}
用 \ref{cor:III.1.7} 和 \ref{thm:II.2.1}，归结到 $A$ 是域的情形.
\end{proof}

用 \ref{prop:III.1.4} (ii) 和 \ref{cor:II.4.13}，$k$ 的适当扩张 $k'$ 把我们归结到 $B/k$ 的剩余扩张平凡的情形. 由 \ref{prop:III.1.5} (相应地：\ref{cor:II.5.2})，$B$ 在 $k$ 上光滑 (相应地：在 $k$ 上形式光滑) 当且仅当 $B$ 是正则局部环 (相应地：其完备化是 $k$ 上的形式幂级数代数). 而这两个条件众所周知是等价的 (剩余扩张平凡时).

\section{形式光滑同态的特征提升性质}

\begin{theorem}\label{thm:III.2.1}
设 $A \to B$ 是定义有限剩余扩张的局部环局部同态. 下列条件等价：
\begin{enumerate}[(i)]
\item $B$ 在 $A$ 上形式光滑；
\item 对每个局部同态 $A \to C$ (其中 $C$ 是完备局部环)，$C$ 的每个含于 $\mathfrak{r}(C)$ 的理想 $J$，以及每个局部 $A$-同态 $B \to C/J$，存在提升它的 (必然局部的) $A$-同态 $B \to C$；
\item 对每个 $A$-代数 $C$ (不必是诺特环)，$C$ 的每个幂零理想 $J$，以及每个连续的 (i.e. 在 $\mathfrak{r}(B)$ 的某个幂上为零的) $A$-同态 $B \to C/J$，存在提升它的 (也必然连续的) $A$-同态 $B \to C$；
\item 同 (ii) 和 (iii) 的陈述，但 $C$ 是在 $A$ 上有限的 Artin 局部环.
\item[(iv bis)] 同 (iv)，但 $J$ 的平方为零.
\end{enumerate}
\end{theorem}

\begin{remarkstar}
对本讲的其余部分，我们只用蕴含 (iv)$\Rightarrow$(i) 或 (iv bis)$\Rightarrow$(i)；直接蕴含 (i)$\Rightarrow$(ii) 将在下一节用另一方法、在 $B$ 是 $A$ 上有限型代数的局部化的情形证明. 回忆在 Cohen 定理的``好''理论\footnote{参见 EGA 0$_{\mathrm{IV}}$ 19.3，19.8.}中，性质 (ii) 或 (iii) 成为形式光滑同态的定义，而 \ref{def:III.1.1} 成为仅在剩余扩张有限的情形成立的一个特征性质. 注意 (ii) 和 (iii) 中任何一个都不比另一个更一般. 可以给出 (等价的) 性质，同时涵盖它们两者：引入线性拓扑化的分离完备环 $C$，$C$ 的闭的拓扑幂零理想，以及连续同态 $A \to C$ (从而把 $C$ 变成拓扑 $A$-代数)；我们把这个修改留给读者.
\end{remarkstar}

 \paragraph*{\ref{thm:III.2.1} 的证明.} 证明 (i)$\Rightarrow$(iii)$\Rightarrow$(ii)，然后 (iv)$\Rightarrow$(i). 由于 (ii)$\Rightarrow$(iv) 平凡，而 (iv) 与 (iv bis) 的等价性通过对满足 $J^n = 0$ 的整数 $n$ 作直接归纳看出，这将完成证明.

(i)$\Rightarrow$(iii). 直接归纳归结到 $J^2 = 0$ 的情形. 由于 $C$ 在 $A$ 上有限，存在 $A$ 的极大理想的幂 $\mathfrak{m}^q$ 零化 $C$. 除以 $\mathfrak{m}^q$，并注意由 \ref{prop:III.1.4} (i)，$B/\mathfrak{m}^q B$ 仍在 $A/\mathfrak{m}^q$ 上形式光滑，可以假设 $A$ Artin. 由 \ref{lem:III.1.3}，$B$ 在 $A$ 上平坦，故 $A$ Artin 时 $B$ 在 $A$ 上自由. 因此存在提升所给同态 $u \colon B \to C/J$ 的 $A$-模同态
\[
w \colon B \longrightarrow C.
\]
令
\[
f(x, y) = w(xy) - w(x)w(y) \quad (x, y \in B),
\]
则 $f(x, y) \in J$，故 $f$ 是 $B \times B$ 到 $J$ 的 $A$-双线性映射. 为使存在提升 $u$ 的代数同态 $v \colon B \to C$，必要且充分条件是存在 $A$-线性映射 $g \colon B \to J$，使得 $v = w + g$ 是代数同态，这写成
\[
\begin{array}{c}
g(1) = 1 - w(1), \\
g(xy) - u(x)g(y) - u(y)g(x) = -f(x, y) \qquad (\text{对 } x, y \in B).
\end{array}
\]
这是 $\mathrm{Hom}_A(B, J)$ 中的线性方程组，右端在 $J$ 中，故它有解当且仅当在 $\mathrm{Hom}_{A}(B, J) \otimes_A A'$ 中、右端在 $J' = J \otimes_A A'$ 中的相应方程组有解，其中 $A'$ 是 $A$ 上的忠实平坦代数. 而设 $A'$ 是 $A$ 上的有限自由局部代数，使得 $B' = B \otimes_A A'$ 是 $A'$ 上的形式幂级数代数 (N.B. 在证明中可以假设 $A$ 和 $B$ 完备，立即看出). 由于 $A'$ 在 $A$ 上有限自由，有
\[
\operatorname{Hom}_{A}(B, J) \otimes_A A' = \operatorname{Hom}_{A'}(B', J'),
\]
并看到在 $\mathrm{Hom}_{A'}(B', J')$ 中得到的方程组，正是确定那些经``修正'' $A'$-模同态 $g' \colon B' \to J'$ 提升由 $u$ 标量扩张得到的 $A'$-模同态 $w' \colon B' \to C'$、从而提升 $u' \colon B' \to C'/J'$ 的 $A'$-代数同态 $B' \to C' = C \otimes_A A'$ 的方程组. (注意 $B$ 作为 $A'$-模生成 $B'$.) 于是归结到证明 $B$ 是 $A$ 上的形式幂级数代数 $B = A[\![t_1, \dots, t_n]\!]$ 时 (iii) 成立. 此时把 $t_i$ 在 $C/J$ 中的像任意提升为 $C$ 的元素 $z_i$. 由于 $z_i$ 模 $J$ 是幂零的 ($u \colon B \to C/J$ 连续)，$z_i$ 亦然 (因为 $J$ 幂零)，故诸 $z_i$ 定义 $B$ 到离散拓扑 $C$ 中的连续拓扑 $A$-代数同态，显然提升 $u$，证毕.

(iii)$\Rightarrow$(ii). 设 $n$ 是 $C$ 的极大理想，对每个整数 $q > 0$，设
\[
C_q = C/\mathfrak{n}^q, \qquad J_q = (J + \mathfrak{n}^q)/\mathfrak{n}^q,
\]
则 $C_q/J_q$ 等同于 $C/J$ 的商代数，另一方面复合同态 $u_q \colon B \to C/J \to C_q/J_q$ 是 $B$ 到离散环 $C_q/J_q$ 的连续同态，且 $J_q$ 是 $C_q$ 中的幂零理想. 于是逐步构造 $A$-同态
\[
v_q \colon B \longrightarrow C_q
\]
使得 (a) $v_q$ 提升 $u_q$，(b) $v_q$ 提升 $v_{q-1}$. 归纳的可行性容易验证，因为 $u_q \colon B \to C/(J + \mathfrak{n}^q)$ 和 $v_{q-1} \colon B \to C/\mathfrak{n}^{q-1}$ 定义同一个同态
\[
B \longrightarrow C/\bigl((J + \mathfrak{n}^q) + \mathfrak{n}^{q-1}\bigr) = C/(J + \mathfrak{n}^{q-1}) = C_{q-1}/J_{q-1},
\]
即 $u_{q-1}$，它们定义一个同态
\[
B \longrightarrow C/J'_q \quad \text{其中} \quad J'_q = (J + \mathfrak{n}^q) \cap \mathfrak{n}^{q-1} \supset \mathfrak{n}^q
\]
(它们都由它约化得到). 归结为把 $B \to C/J'_q$ 提升为 $B$ 到 $C_q$ 的某个含于 $J_q$ 中的理想 $J'_q/\mathfrak{n}^q$ 之商的同态，因此该理想幂零，由假设 (iii) 这是可能的.

做完之后，诸 $v_q$ 定义 $B$ 到 $C_q$ 的投射极限 $C$ 的同态. 由于 $J$ 闭，$J$ 是诸 $J_q$ 的投射极限，故 $v$ 提升 $u$，证毕.

(iv)$\Rightarrow$(i). 首先立即看出，若 (iv) 成立，则对 $A$ 上的有限局部代数 $A'$，它对 $B \otimes_A A'$ 的局部分量在 $A'$ 上仍成立. 取 $A'$ 在 $A$ 上自由且使 $B'$ 在 $A'$ 上的剩余扩张平凡，由 \ref{prop:III.1.4} (ii) 归结到 $B$ 在 $A$ 上的剩余扩张平凡的情形. 此时我们证明一个稍更精确的结果：

\begin{corollary}\label{cor:III.2.2}
在 \ref{thm:III.2.1} 的条件下，再假设 $B$ 在 $A$ 上的剩余扩张平凡. 则 \ref{thm:III.2.1} 的等价条件还等价于下列两个条件 (在 (v) 中假设 $A$ 和 $B$ 完备)：
\begin{enumerate}[(iv ter)]
\item[(iv ter)] 同 (iv)，但 $A$ 上有限的 Artin 局部环 $C$ 限于具有平凡的剩余扩张 (且，如果愿意，理想 $J$ 的平方为零).
\item[(v)] 存在局部 $A$-同态 (其中 $n = \dim \mathfrak{n}/(\mathfrak{n}^2 + \mathfrak{m}B)$)
\[
u \colon B \longrightarrow B_1 = A[\![t_1, \dots, t_n]\!]
\]
诱导同构
\[
\mathfrak{n}/(\mathfrak{n}^2 + \mathfrak{m}B) \stackrel{\sim}{\longrightarrow} \mathfrak{n}_1/(\mathfrak{n}_1^2 + \mathfrak{m}B_1),
\]
其中 $n$ ($n_1$) 是 $B$ ($B_1$) 的极大理想，$m$ 是 $A$ 的极大理想.
\end{enumerate}
\end{corollary}

\paragraph*{证明.} 由于 (iv bis) 显然蕴含 (iv ter) —— 撇开平方为零理想的玩笑不谈 —— 只需证明 (iv ter)$\Rightarrow$(v)$\Rightarrow$(i).

(iv ter)$\Rightarrow$(v). 选择 $\mathfrak{n}/(\mathfrak{n}^2 + \mathfrak{m}B)$ 的基 $a_1, \dots, a_n$，它因此定义局部 $A$-代数同态
\[
B \longrightarrow B_1/(\mathfrak{n}_1^2 + \mathfrak{m}B_1) = k[t_1, \dots, t_n]/(t_1, \dots, t_n)^2,
\]
由 (iv ter)，可以把它逐步提升为 $B$ 到 $B_1/\mathfrak{n}_1^2$，$B_1/\mathfrak{n}_1^3$ 等的 $A$-代数同态，取投射极限得到具有所述性质的同态 $B \to B_1$.

(v)$\Rightarrow$(i). 在交换图中 (两行正合，两端垂直箭头满射)
\[
\begin{tikzcd}[row sep=large, column sep=large]
\mathfrak{n}/(\mathfrak{n}^2 + \mathfrak{m}B) \arrow[r] \arrow[d] & \mathfrak{n}_1/(\mathfrak{n}_1^2 + \mathfrak{m}B_1) \arrow[d] \\
B/\mathfrak{m}B \arrow[r] & B_1/\mathfrak{m}B_1
\end{tikzcd}
\]
中间箭头满射，由此 ($B$ 完备) 得出 $B \to B_1$ 满射. 设 $x_i$ ($1 \leqslant i \leqslant n$) 是 $B$ 中提升 $t_i$ 的元素. 它们因此定义 $A$-代数同态 $B_1 \to B$，它同 $u$ 一样是满射，且按构造与 $u$ 的复合是恒等. 故 $B_1 \to B$ 也单射，从而是同构. 于是得到

\begin{corollary}\label{cor:III.2.3}
在 \ref{cor:III.2.2} (v) 的条件下，$u$ 必然是同构.
\end{corollary}

\begin{proof}
这完成 $B$ 在 $A$ 上形式光滑的证明. 同时还重新得到 \ref{prop:III.1.5} (但这没什么了不起).
\end{proof}

\section{态射到光滑 \texorpdfstring{$S$-概形}{S-概形}中的局部无穷小延拓}

\begin{theorem}\label{thm:III.3.1}
设 $f \colon X \to Y$ 是局部有限型态射. 下列条件等价：
\begin{enumerate}[(i)]
\item $f$ 光滑.
\item 对 $Y$ 上的每个预概形 $Y'$，每个与 $Y'$ 有相同底空间的 $Y'$ 的闭子预概形 $Y'_0$，每个 $Y$-态射 $g_0 \colon Y'_0 \to X$ 以及每个 $z \in Y'_0$，存在 $z$ 在 $Y'$ 中的开邻域 $U$ 以及 $g_0|Y'_0 \cap U$ 到 $Y$-态射 $U \to X$ 的延拓 $g$.
\item[(ii bis)] 对 (ii) 中的 $Y'$，$Y'_0$ 和 $z$，令 $X' = X \times_Y Y'$，$X'_0 = X \times_Y Y'_0$，$X'_0$ 在 $Y'_0$ 上的每个截面延拓为 $X'$ 在 $z$ 的开邻域 $U$ 上的截面.
\item 对每个 $Y$-概形 $Y'$ (为某个 $\mathcal{O}_y$ ($y \in Y$) 上有限的 Artin 局部环的谱)，$Y'$ 的每个非空闭子预概形 $Y'_0$，以及每个 $Y$-态射 $g_0 \colon Y'_0 \to X$，存在延拓 $g_0$ 的 $Y$-态射 $g \colon Y' \to X$.
\item[(iii bis)] 对 (iii) 中的每个 $Y'$，令 $X' = X \times_Y Y'$，$X'_0 = X \times_Y Y'_0$，$X'_0$ 在 $Y'_0$ 上的每个截面延拓为 $X'$ 在 $Y'$ 上的截面.
\end{enumerate}
\end{theorem}

\paragraph*{证明.} (ii) 与 (ii bis) 的等价性，以及 (iii) 与 (iii bis) 的等价性，是平凡的，(ii)$\Rightarrow$(iii) 亦然. 剩下证明 (i)$\Rightarrow$(ii) 和 (iii)$\Rightarrow$(i).

(i)$\Rightarrow$(ii). 设 $x = g_0(z)$. 把 $X$ 换成 $x$ 的适当开邻域，$Y'$ 换成 $g_0$ 在 $x$ 的邻域的原像上的诱导预概形，可以假设 $X$ 在 $Y[t_1, \dots, t_n]$ 上平展. 考虑复合 $Y$-态射 $Y'_0 \to X \to Y[t_1, \dots, t_n]$，它由层 $\mathcal{O}_{Y'_0}$ 的 $n$ 个截面定义，这些截面可以延拓到 $z$ 的邻域成为 $\mathcal{O}_{Y'}$ 的截面，故可以假设该态射已延拓为 $Y$-态射 $Y' \to Y[t_1, \dots, t_n]$. 由 (\ref{cor:I.5.6})，存在唯一提升前者的 $Y$-态射 $g \colon Y' \to X$，它同时延拓 $g_0$，证毕.

(iii)$\Rightarrow$(i). 由于 $f$ 光滑的点集是开的，只需证明它包含每个在纤维中闭的 $x \in X$. 设 $y = f(x)$，则 $\mathcal{O}_x$ 是 $\mathcal{O}_y$ 上的代数，是有限型代数的局部化，剩余扩张有限. 另一方面，假设 (iii) 蕴含：$\mathcal{O}_x$ 到 $A/J$ (其中 $A$ 是 $\mathcal{O}_y$ 上的有限 Artin 局部代数，$J$ 是含于其根式中的理想) 的每个同态，可以提升为 $\mathcal{O}_x$ 到 $A$ 的同态 (考虑到 $\mathrm{Spec}(B)$ ($B$ 局部环) 到 $X$ 的态射由 $\mathcal{O}_x$ ($x \in X$) 到 $B$ 的局部同态一一确定). 由 \ref{thm:III.2.1}，由此得出 $\mathcal{O}_x$ 在 $\mathcal{O}_y$ 上形式光滑，从而由 \ref{prop:III.1.9} 在 $\mathcal{O}_y$ 上光滑.

\begin{corollary}\label{cor:III.3.2}
设 $f \colon X \to Y$ 如 \ref{thm:III.3.1}. 下列条件等价：
\begin{enumerate}[(i)]
\item $f$ 平展；
 \item \ref{thm:III.3.1} 的条件 (ii) 成立，且 $g_0$ 到 $U$ 的延拓 $g$ 唯一；
 \item \ref{thm:III.3.1} 的条件 (iii) 成立，且 $g$ 唯一.
\end{enumerate}
\end{corollary}

\begin{proof}
只需注意，在上述 (i)$\Rightarrow$(ii) 的证明中，只有当 $n = 0$ (已知充分的条件) 时才能有唯一性 (当 $Y'_0$ 在 $z$ 的邻域内不恒等于 $Y'$ 时).
\end{proof}

\begin{corollary}\label{cor:III.3.3}
设 $X$ 是完备局部环 $A$ 上的局部有限型预概形，$y$ 是 $Y = \mathrm{Spec}(A)$ 的闭点，$x$ 是 $f^{-1}(y)$ 中在 $k(y)$ 上有理的点. 若 $X$ 在 $A$ 上于 $x$ 处光滑，则存在 $X$ 在 $Y$ 上的截面 $s$，``经过 $x$''，i.e. 使得 $s(y) = x$.
\end{corollary}

\begin{proof}
特别地，若 $X$ 在 $A$ 上光滑，则自然映射
\[
\Gamma(X/Y) \longrightarrow \Gamma(X \otimes_A k/k)
\]
(从 $X$ 在 $Y$ 上的截面到纤维 $f^{-1}(y) = X \otimes_A k$ 中 $k$ 上有理点的集合) 满射. 这个事实在 $A$ 是离散赋值环且 $X$ 在 $A$ 上紧合 (事实上：射影) 时尤为著名，此时 $X$ 在 $Y$ 上的截面 (i.e.``$X$ 的在 $A$ 中取值的点'') 也等同于有理截面，i.e. $X \otimes_A K = X_K$ (它是 $K$ 上紧合且简单的概形) 中 $K$ 上有理的点 (i.e. $X$ 的 $K$-有理点)，其中 $K$ 是 $A$ 的分式域.
\end{proof}

\section{光滑 \texorpdfstring{$S$-概形}{S-概形}的局部无穷小延拓}

\begin{theorem}\label{thm:III.4.1}
设 $Y$ 是局部诺特预概形，$Y_0$ 是与其有相同底空间的闭子预概形，$X_0$ 是光滑的 $Y_0$-预概形，$x$ 是 $X_0$ 的点. 则存在 $x$ 的开邻域 $U_0$，在 $Y$ 上光滑的预概形 $X$，以及 $Y_0$-同构
\[
h \colon U_0 \stackrel{\sim}{\longrightarrow} X \times_Y Y_0.
\]
此外，若 $(U'_0, X', h')$ 是该问题的另一个解，则``它在 $x$ 的邻域内同构于第一个解''.
\end{theorem}

\begin{proof}
``它''的含义留给读者. 可以注意到，对给定的 $U_0$，该问题的一个解相当于在 $U_0$ 上给出 $f_0^{-1}(\mathcal{O}_Y)$ (其中 $f_0$ 是结构态射 $U_0 \to Y_0$ 之下的连续映射) 的代数层 $\mathcal{B}$，带有与同态 $f^{-1}(\mathcal{O}_Y) \to f^{-1}(\mathcal{O}_{Y_0})$ 相容的环同态 $\mathcal{B} \to \mathcal{O}_{U_0}$，使得
\begin{enumerate}[a)]
\item 该同态诱导同构
\[
\mathcal{B} \otimes_{f^{-1}(\mathcal{O}_Y)} f^{-1}(\mathcal{O}_{Y_0}) \stackrel{\sim}{\longrightarrow} \mathcal{O}_{Y_0};
\]
\item $U_0$ 配备 $\mathcal{B}$ 后成为光滑的 $Y$-预概形.
\end{enumerate}
这样，局部唯一性断言的精确含义变得特别明显.
\end{proof}

\paragraph*{证明.} 已经可以假设 $X_0$ 在 $Y_0[t_1, \dots, t_n] = Y'_0$ 上平展. 而后者可以看作 $Y' = Y[t_1, \dots, t_n]$ 的与它有相同底空间的闭子预概形. 由 (\ref{thm:I.8.3})，存在 $Y'$ 上平展的 $X$ 以及 $Y'_0$-同构 $X \times_{Y'} Y'_0 \stackrel{\sim}{\longrightarrow} X'_0$. 存在性已经解决. 对唯一性，用光滑态射的性质 \ref{thm:III.3.1} (ii)，并考虑下述引理：

\begin{lemma}\label{lem:III.4.2}
设 $Y$ 是预概形，$Y_0$ 是由局部幂零理想层 $\mathcal{J}$ 定义的闭子预概形，$X$ 和 $X'$ 是 $Y$-预概形，$u \colon X \to X'$ 是 $Y$-态射. 假设 $X$ 在 $Y$ 上平坦. 为使 $u$ 是同构，必要且充分条件是 $u_0 \colon X \times_Y Y_0 \to X' \times_Y Y_0$ 是同构.
\end{lemma}

\begin{proof}
证明容易，归结到仿射情形并考察相伴分次环. 还注意到把``同构''换成``闭浸入''得到的类似陈述也成立，且无需平坦性假设.
\end{proof}

\begin{remark}\label{rem:III.4.3}
必须注意， \ref{thm:III.4.1} 中得到的局部延拓 $X$ 不是典范的，换言之，两个解之间的局部同构不唯一，i.e. 一般存在 $X$ 的非平凡 $Y$-自同构，在闭子预概形 $X_0 = X \times_Y Y_0$ 上诱导恒等. 正是因此，在构造光滑预概形的整体无穷小延拓时，应当预期存在上同调性质的\emph{障碍}\index[term]{zhangai@障碍} (obstruction)，它将在下面 (§6) 精确化.
\end{remark}

\section{态射的整体无穷小延拓}

设 $T$ 是拓扑空间，$G$ 是 $X$ 上的群层，$\mathcal{P}$ 是 $T$ 上 $\mathcal{G}$ 在其上作用的集合层 (为确定起见作用在右边). 称 $\mathcal{P}$ 在 $\mathcal{G}$ 下\emph{形式主齐性} (formally principal homogeneous)，若众所周知的集合层同态
\[
\mathcal{G} \times \mathcal{P} \longrightarrow \mathcal{P} \times \mathcal{P}
\]
(由 $\mathcal{G}$ 在 $\mathcal{P}$ 上的作用得出) 是同构. 这等价于说：对每个 $x \in T$，$\mathcal{P}_x$ 是空的，或是在通常群 $\mathcal{G}_x$ 下的主齐性空间；也等价于说：对 $T$ 的每个开集 $U$，$\mathcal{P}(U)$ 是空的，或是在通常群 $\mathcal{G}(U)$ 下的主齐性空间. 称 $\mathcal{P}$ 是 $\mathcal{G}$ 下的主齐性层 (principal homogeneous sheaf)，若它形式主齐性，且诸 $\mathcal{P}_x$ 非空 (换言之，若所有 $\mathcal{P}_x$ 都是 $\mathcal{G}_x$ 下的 (非空的) 主齐性空间)\footnote{看来最好采用 Giraud 论文中引入的更短更形象的术语``$G$ 下的挠子 (torseur under G)''.}. 回忆 $\mathcal{G}$ 下的主齐性层的 (同构) 类的集合等同于上同调集 $H^1(T, \mathcal{G})$，当 $\mathcal{G}$ 交换时它也是 $T$ 以 $\mathcal{G}$ 为系数的通常上同调群. 于是，对每个主齐性 $\mathcal{P}$，有特征类 (characteristic class) $c(\mathcal{P}) \in H^1(T, \mathcal{G})$，其平凡性是 $\mathcal{P}$ 平凡 (i.e. 同构于 $\mathcal{G}$，$\mathcal{G}$ 在其上通过右平移作用) 的必要且充分条件，也就是 $\mathcal{P}$ 有截面的必要且充分条件.

\begin{proposition}\label{prop:III.5.1}
设 $S$ 是预概形，$X$ 和 $Y$ 是 $S$ 上的预概形，$Y_0$ 是由 $Y$ 上平方为零的理想层 $\mathcal{J}$ 定义的 $Y$ 的闭子预概形. 设 $g_0$ 是 $Y_0$ 到 $X$ 的 $S$-态射，$\mathcal{P}(g_0)$ 是 $Y$ 上的层，其在开集 $U$ 上的截面是 $g_0|U \cap Y_0$ 到 $S$-态射 $g \colon U \to X$ 的延拓 $g$. 则 $\mathcal{P}(g_0)$ 是交换群层
\[
\mathcal{G} = \mathbf{Hom}_{\mathcal{O}_{Y_0}}\bigl(g_0^{*}(\varOmega_{X/S}^{1}), \mathcal{J}\bigr)
\]
下的 (自然地) 形式主齐性层.
\end{proposition}

\begin{proof}
令 $\mathcal{P} = \mathcal{P}(g_0)$. 必须对 $Y$ 的每个开集 $U$ 定义映射
\[
\mathcal{P}(U) \times \mathcal{G}(U) \longrightarrow \mathcal{P}(U),
\]
使得 (a) 对固定的 $g \in \mathcal{P}(U)$，映射 $s \mapsto gs$ 从 $\mathcal{G}(U)$ 到 $\mathcal{P}(U)$ 是双射；(b) $\mathcal{P}(U)$ 成为以 $\mathcal{G}(U)$ 为算子群的集合；(c) 前面的映射与 $V \subset U$ 的限制算子相容. (c) 的验证平凡，为简化可以假设 $U = Y$. (b) 的验证 (如所愿是局部性质的) 留给读者，因此对给定的 $g \in \mathcal{P}(Y)$，只定义从 $\mathcal{G}(Y)$ 到 $\mathcal{P}(Y)$ 的自然双射. 于是假设已给出 $S$-态射 $g \colon Y \to X$，求典范双射
\begin{equation}
\mathrm{Hom}_{\mathcal{O}_{Y_0}}\bigl(g_0^{*}(\varOmega_{X/S}^{1}), \mathcal{J}\bigr) \stackrel{\sim}{\longrightarrow} \mathcal{P}(Y), \tag{($*$)}
\end{equation}
其中 $\mathcal{P}(Y)$ 是与 $g$ 诱导同一个态射 $g_0 \colon Y_0 \to X$ 的 $Y$ 到 $X$ 的 $S$-态射 $g'$ 的集合. 这样的 $g'$ 的给出等价于 $S$-态射 $h \colon Y \to X \times_S X$ 的给出，满足 $\mathrm{pr}_1 \circ h = g$，且 $h \circ i = (g_0, g_0)$，其中 $\mathrm{pr}_1 \colon X \times_S X \to X$ 是第一投影，$i \colon Y_0 \to Y$ 是典范浸入，$(g_0, g_0) \colon Y_0 \to X \times_S X$ 是分量为 $g_0$，$g_0$ 的态射 $\Delta_{X/S} g_0$：
\[
\begin{tikzcd}[row sep=large, column sep=large]
X \times_S X \arrow[d,"\mathrm{pr}_1"] & Y_0 \arrow[l,"{h_0 = (g_0, g_0)}"] \\
X & Y \arrow[l,"h"]
\end{tikzcd}
\]
由于 $h_0$ 通过对角浸入 $\Delta_{X/S}$ 分解，且 $Y$ 在 $Y_0$ 的一阶无穷小邻域中 (i.e. $\mathcal{J}^2 = 0$)，所求的 $h$ 必然 (唯一地) 通过对角线的一阶无穷小邻域分解，后者 (作为 $X$-预概形，借助 $\mathrm{pr}_1$) 等同于代数层 $\mathcal{O}_X + \varOmega_{X/S}^{1}$ 的谱 $X'$ (其中第二项看作平方为零的理想)，对角态射 $X \to X'$ 对应于该代数层的典范增广. 令 $Y' = X' \times_X Y$，$Y'_0 = Y' \times_Y Y_0 = X' \times_X Y_0$，于是所求的 $h$ 与 $Y'$ 在 $Y$ 上的、延拓 $Y'_0$ 在 $Y_0$ 上的给定截面 $u_0$ 的截面 $u$ 一一对应. 此外可以把 $Y'$ 等同于 $Y$ 上的代数层 $\mathcal{A} = g^{*}(\mathcal{O}_X + \varOmega_{X/S}^{1}) = \mathcal{O}_Y + g^{*}(\varOmega_{X/S}^{1})$ 的谱，把 $Y'_0$ 等同于代数层 $\mathcal{A}_0 = \mathcal{A} \otimes_{\mathcal{O}_Y} \mathcal{O}_{Y_0} = \mathcal{O}_{Y_0} + g_0^{*}(\varOmega_{X/S}^{1})$ 的谱，则截面 $u_0$ 是由 $A_0$ 到 $\mathcal{O}_{Y_0}$ 的典范增广定义的截面. 故 $\mathcal{P}(Y)$ 等同于诱导典范增广 $\mathcal{A}_0 \to \mathcal{O}_{Y_0}$ 的代数同态 $\mathcal{A} \to \mathcal{O}_Y$ 的集合. 而代数同态 $\mathcal{A} \to \mathcal{O}_Y$ 与模层同态 $\mathcal{M} \to \mathcal{A}$ (为简化令 $\mathcal{M} = g^{*}(\varOmega_{X/S}^{1})$) 一一对应，且所关注的是那些诱导零同态 $\mathcal{M}_0 \to \mathcal{O}_{Y_0}$ (其中 $\mathcal{M}_0 = \mathcal{M} \otimes_{\mathcal{O}_Y} \mathcal{O}_{Y_0}$)，i.e. 把 $\mathcal{M}$ 映入增广理想 $\mathcal{J}$ 的同态. 于是得到集合
\[
\mathrm{Hom}_{\mathcal{O}_Y}(\mathcal{M}, \mathcal{J}) = \mathrm{Hom}_{\mathcal{O}_{Y_0}}(\mathcal{M}_0, \mathcal{J})
\]
(因为 $\mathcal{J}$ 被 $\mathcal{J}$ 零化). 这就是所求的典范双射 ($*$).
\end{proof}

\begin{corollary}\label{cor:III.5.2}
若 $X$ 在 $S$ 上光滑 (至少在 $g_0(Y_0)$ 的点处)，则 $\mathcal{P}$ 甚至是交换群层 $G$ 下的主齐性层，此时 $G$ 也可以写成
\[
\mathcal{G} = g_0^{*}(\mathfrak{g}_{X/S}) \otimes_{\mathcal{O}_{Y_0}} \mathcal{J},
\]
其中 $\mathfrak{g}_{X/S}$ 是 $\varOmega_{X/S}^{1}$ 的对偶层，i.e. $X$ 相对 $S$ 的切层 (tangent sheaf) (或导子层). (最后一个公式来自 $\varOmega_{X/S}^{1}$ 此时有限自由这一事实.)
\end{corollary}

\begin{proof}
特别地，与这个主齐性层相对应的是 $H^1(Y_0, \mathcal{G})$ 中的一个上同调类，其消没是存在延拓 $g_0$ 的 $S$-态射 $g$ 的必要且充分条件. 若存在这样的延拓，则一切可能延拓的集合是群 $H^0(Y_0, \mathcal{G})$ 下的齐性空间.
\end{proof}

在应用形式几何的方法时，情形通常如下：给出两个 $S$-预概形 $X$ 和 $Y$，$S$ 上的凝聚理想层 $\mathcal{I}$；用 $S_n$ 记由 $\mathcal{I}^{n+1}$ 定义的 $S$ 的闭子预概形，并令
\[
X_n = X \times_S S_n, \qquad Y_n = Y \times_S S_n.
\]
假设给出 $S_n$-态射
\[
g_n \colon Y_n \longrightarrow X_n
\]
(或等价地，$S$-态射 $Y_n \to X$，或 $S_{n+1}$-态射 $Y_n \to X_{n+1}$，因为这样的态射必然诱导 $Y_n \to X_n$)，并设法把它延拓为 $S_{n+1}$-态射
\[
g_{n+1} \colon Y_{n+1} \longrightarrow X_{n+1}.
\]
(若可以无限继续，就得到 $Y$ 和 $X$ 关于理想 $\mathcal{I}\mathcal{O}_Y$ 和 $\mathcal{I}\mathcal{O}_X$ 的形式完备化预概形 $\widehat{Y} \to \widehat{X}$ 的态射.) 可以应用 \ref{prop:III.5.1}，把 $(S, X, Y, Y_0, g_0)$ 换成 $(S_{n+1}, X_{n+1}, Y_{n+1}, Y_n, g_n)$，此时群层 $G$ 成为 $g_n^{*}\bigl(\varOmega_{X_{n+1}/S_{n+1}}^{1}\bigr)$ 到 $\mathcal{J} = \mathcal{I}^{n+1}\mathcal{O}_Y/\mathcal{I}^{n+2}\mathcal{O}_Y$ 的模层同态层. 由于 $\mathcal{J}$ 被 $\mathcal{I}\mathcal{O}_Y$ 零化，可以把 $g_n^{*}\bigl(\varOmega_{X_{n+1}/S_{n+1}}^{1}\bigr)$ 换成它在 $Y_0$ 上诱导的层，即 $h_0^{*}\bigl(\varOmega_{X/S}^{1}\bigr)$，其中 $h_0$ 是复合 $Y_0 \to Y_n \to X_{n+1}$，或复合 $Y_0 \to X_0 \to X_{n+1}$，其中 $g_0 \colon Y_0 \to X_0$ 由 $g_n$ 诱导. 由于 $\varOmega_{X_{n+1}/S_{n+1}}^{1}$ 在 $X_0 = X_{n+1} \times_{S_{n+1}} S_0$ 上的原像为 $\varOmega_{X_0/S_0}^{1}$，立即看到
\[
\mathcal{G} = \mathbf{Hom}_{\mathcal{O}_{Y_0}}\bigl(g_0^{*}(\varOmega_{X_0/S_0}^{1}), \mathcal{I}^{n+1}\mathcal{O}_Y/\mathcal{I}^{n+2}\mathcal{O}_Y\bigr).
\]
于是得到：

\begin{corollary}\label{cor:III.5.3}
设 $S$，$X$，$Y$，$\mathcal{I}$，$g_n$ 如上，$\mathcal{P}(g_n)$ 是 $Y$ 上的层，其在开集 $U$ 上的截面是把 $g_n$ 延拓为 $S_{n+1}$-态射 $Y_{n+1} \to X_{n+1}$ 的 $g_{n+1}$. 则 $\mathcal{P}(g_n)$ 是群层
\[
\mathcal{G} = \mathbf{Hom}_{\mathcal{O}_{Y_0}}\bigl(g_0^{*}(\varOmega_{X_0/S_0}^{1}), \mathrm{gr}_{\mathcal{I}\mathcal{O}_Y}^{n+1}(\mathcal{O}_Y)\bigr)
\]
下的形式主齐性层.
\end{corollary}

特别地：

\begin{corollary}\label{cor:III.5.4}
若此外 $X$ 在 $S$ 上光滑 (至少在 $g_0(Y_0)$ 的点处)，则 $\mathcal{P}(g_n)$ 甚至是主齐性层. 特别地，它定义 $H^1(Y_0, \mathcal{G})$ 中的一个障碍类，其消没是存在 $g_n$ 的整体延拓 $g_{n+1}$ 的必要且充分条件. 若存在这样的延拓，则一切整体延拓的集合是 $H^0(Y_0, \mathcal{G})$ 下的主齐性空间. 最后，在所考虑的情形，层 $\mathcal{G}$ 也可以写成
\[
\mathcal{G} = g_0^{*}(\mathfrak{g}_{X_0/S_0}) \otimes_{\mathcal{O}_{Y_0}} \mathrm{gr}_{\mathcal{I}\mathcal{O}_Y}^{n+1}(\mathcal{O}_Y).
\]
\end{corollary}

\begin{proof}
逐步进行，看到若所有 $H^1(Y_0, \mathcal{G}_n)$ 都消没 (其中 $\mathcal{G}_n = g_0^{*}\bigl(\mathfrak{g}_{X_0/S_0}\bigr) \otimes \mathrm{gr}_{\mathcal{I}\mathcal{O}_Y}^{n}(\mathcal{O}_Y)$)，则从任意的 $g_k$ 出发，可以把它依次延拓为 $g_{k+1}$，$\dots$. 特别地，若 $\mathcal{I}$ 幂零，就能找到 $g_k$ 到 $Y$ 的延拓 $g$. $H^1$ 消没的条件在 $Y_0$ 仿射时尤其成立. 于是得到：
\end{proof}

\begin{corollary}\label{cor:III.5.5}
在定理 \ref{thm:III.3.1} 的陈述中，得到一个与其余条件等价的必要且充分条件：假设 (ii) (或 (ii bis)) 中出现的 $Y'$ 仿射，并要求 $g_0$ 存在到整个 $Y'$ 的整体延拓 $g$.
\end{corollary}

注意 \ref{thm:III.3.1} 的证明不能直接给我们这个结果.

一个重要情形是 $Y$ 在 $S$ 上平坦，此时有
\[
\mathrm{gr}^{n}(\mathcal{O}_Y) = \mathrm{gr}^{n}(\mathcal{O}_S) \otimes_{\mathcal{O}_{S_0}} \mathcal{O}_{Y_0},
\]
且当诸 $\mathrm{gr}^{n}(\mathcal{O}_S)$ 在 $S$ 上局部自由时，得到
\[
\mathcal{G}_n = \mathbf{Hom}_{\mathcal{O}_{Y_0}}\bigl(g_0^{*}(\varOmega_{X_0/S_0}^{1}), \mathcal{O}_{Y_0}\bigr) \otimes_{\mathcal{O}_{S_0}} \mathrm{gr}^{n}(\mathcal{O}_S),
\]
或，若 $\varOmega_{X_0/S_0}^{1}$ 也局部自由 (例如 $X$ 在 $S$ 上光滑)
\[
\mathcal{G}_n = g_0^{*}\bigl(\mathfrak{g}_{X_0/S_0}\bigr) \otimes_{\mathcal{O}_{S_0}} \mathrm{gr}^{n}(\mathcal{O}_S).
\]
例如若 $S$ 仿射，环为 $A$，$\mathcal{I}$ 由 $A$ 的理想 $I$ 定义，得到
\[
H^{i}(Y_0, \mathcal{G}_n) = H^{i}(Y_0, \mathcal{G}_0) \otimes_A \mathrm{gr}_{I}^{n}(A)
\]
对每个 $i$ 成立 (事实上，问题在 $S_0$ 上是局部的，归结为与自由模层张量的情形). 此时，$H^1(Y_0, \mathcal{G}_0)$ 的消没蕴含 $g_n$ 的一切逐次延拓障碍都为零. 于是得到：

\begin{corollary}\label{cor:III.5.6}
设 $(S, X, Y, \mathcal{I}, g_n)$ 如上，再假设 $X$ 在 $S$ 上光滑，$Y$ 在 $S$ 上平坦，最后 $S$ 仿射，且诸 $\mathrm{gr}^{n}(\mathcal{O}_S) = \mathcal{I}^n/\mathcal{I}^{n+1}$ 局部自由. 则构造 $g_{n+1}$ 的障碍位于 $H^1(Y_0, \mathcal{G}_0) \otimes_A \mathrm{gr}_{I}^{n+1}(A)$ (其中 $A$ 是 $S$ 的环，$I$ 是定义 $\mathcal{I}$ 的 $A$ 的理想)，只要令
\[
\mathcal{G}_0 = g_0^{*}(\mathfrak{g}_{X_0/S_0}).
\]
若 $H^1(Y_0, \mathcal{G}_0) = 0$，则 $g_n$ 可以延拓为 $S$-态射 $\widehat{g} \colon \widehat{Y} \to \widehat{X}$.
\end{corollary}

当然，如果出发点不是通常的 $S$-预概形 $X$ 和 $Y$，而是 $\widehat{S}$-adic 形式预概形 $X$ 和 $Y$，这个结果原样仍然成立. 例如，它可以用来证明某些在完备局部环上紧合的形式概形 (例如) 实际上是代数的. 事实上，像引理 \ref{lem:III.4.2} 那样进行，得到：

\begin{corollary}\label{cor:III.5.7}
在 \ref{cor:III.5.6} 的条件下，若 $g_0$ 是同构，则 $\widehat{g}$ 亦然.
\end{corollary}

(N.B. 对闭浸入也有同样的结果.)

于是得到：

\begin{proposition}\label{prop:III.5.8}
设 $A$ 是极大理想为 $m$、剩余域为 $k$ 的完备局部环，$X$ 和 $Y$ 是 $A$ 上的两个 $m$-adic 形式预概形，在 $A$ 上平坦 (i.e. 对每个 $n$，$X_n$ 和 $Y_n$ 在 $A_n = A/\mathfrak{m}^{n+1}$ 上平坦)，假设 $X_0 = \mathfrak{X} \otimes_A k$ 在 $k$ 上光滑且 $H^1(X_0, \mathfrak{g}_{X_0/k}) = 0$. 则 $Y_0$ 到 $X_0$ 的每个 $k$-同构延拓为 $Y$ 到 $X$ 的 $A$-同构；若还有 $H^0(X_0, \mathfrak{g}_{X_0/k}) = 0$，则这个延拓唯一.
\end{proposition}

\begin{proof}
这特别给出 $A$ 上光滑形式预概形按给定预概形 $X_0$ 约化的唯一性结果 (在 $H^1(X_0, \mathfrak{g}_{X_0/k}) = 0$ 的条件下). 此外，若 $X$ 和 $Y$ 来自 $A$ 上紧合的普通概形 $X$ 和 $Y$，则由形式几何中层存在性定理 (参见 Bourbaki 讨论班报告，N$^{\circ}$182\footnote{证明参见 EGA III 5.4.1.})，知道 $A$-同构 $Y \xrightarrow{\sim} X$ 与形式完备化的 $A$-同构之间有一一对应，因此
\end{proof}

\begin{corollary}\label{cor:III.5.9}
前面的陈述 \ref{prop:III.5.8} 在把 $X$ 和 $Y$ 换成 $A$ 上紧合的普通 $A$-概形 $X$ 和 $Y$ 时仍然成立.
\end{corollary}

最后，当 $X$ 是 $A$ 上紧合的形式概形，且 $Y$ 形如 $\widehat{Y}$ (其中 $Y$ 是 $A$ 上紧合的普通概形) 时，命题 \ref{prop:III.5.8} 给出 $X$ 与 $\widehat{Y}$ 同构的充分条件，因而形式概形 $X$ 事实上``代数化'' (i.e. 同构于 $\widehat{X}$，其中 $X$ 是 $A$ 上紧合的普通概形，后者将被典范地确定). 例如当 $X_0 = \mathbf{P}_k^r$ 时 (或更一般地，当 $X_0$ 是 Severi-Brauer 概形，i.e. 在 $k$ 的代数闭包上同构于标准射影空间时) 就如此：$A$ 上以 $\mathbf{P}_k^r$ 为纤维的每个紧合平坦形式概形都可代数化，更精确地说同构于 $\mathbf{P}_A^r$ 的 $m$-adic 形式完备化. 特别地 (借助``存在性定理'')，$A$ 上以 $\mathbf{P}_k^r$ 为纤维的每个紧合普通概形同构于 $\mathbf{P}_A^r$ ($A$ 是完备局部环). 用下降理论可以证明，若 $A$ 不完备，则 $X$ 在作基的有限平展扩张 $A \to A'$ 后同构于 $\mathbf{P}^r$ (在这种形式下，结果对 Severi-Brauer 概形纤维仍然成立).

\section{光滑 \texorpdfstring{$S$-概形}{S-概形}的整体无穷小延拓}

在定理 \ref{thm:III.4.1} 的条件下，问题是要知道是否存在 $Y$ 上光滑的预概形 $X$，使得 $X \times_Y Y_0$ 是 $Y_0$-同构于 $X_0$ 的，已知这样的概形``在 $X_0$ 上局部存在''. 沿用逐步构造的方法，把 $Y$ 换成字母 $S$，假设给出由理想层 $\mathcal{I}$ (不再需要假设局部幂零) 定义的 $S$ 的闭子预概形 $S_0$，引入由 $\mathcal{I}^{n+1}$ 定义的 $S$ 的闭子预概形 $S_n$，并假设给出了在 $S_n$ 上光滑的子预概形 $X_n$. 目标是找到``按 $X_n$ 约化''的 $S_{n+1}$-预概形 $X_{n+1}$，i.e. 配备同构
\[
X_{n+1} \times_{S_{n+1}} S_n \stackrel{\sim}{\longleftarrow} X_n,
\]
且在 $S_{n+1}$ 上光滑 (或由 \ref{thm:II.2.1}，等价地在 $S_{n+1}$ 上平坦). 如 §4 所指出的，这样的给出相当于在 $f^{-1}(\mathcal{O}_{S_{n+1}})$ (其中 $f$ 是结构态射 $X_n \to S_n$ 之下的连续映射) 上给出代数层 $\mathcal{B}$，配备增广 $\mathcal{B} \to \mathcal{O}_{X_n}$，与增广 $f^{-1}(\mathcal{O}_{S_{n+1}}) \to f^{-1}(\mathcal{O}_{S_n})$ 相容，并满足 (我们不重写，只指出它们是 $X_n$ 的底空间上的局部性质) 两个条件 (a) 和 (b). 由 \ref{thm:III.4.1} 知道一个解局部存在. 它还是同构 (不唯一) 意义下唯一的，至少在局部如此. 先精确化这一点：

\begin{proposition}\label{prop:III.6.1}
设 $X_{n+1}$ 是在 $S_{n+1}$ 上按 $S_n$ 上的 $X_n$ 约化的. 则 ($X_n$ 的底空间上，或 $X_0$ 上) 在 $X_n$ 上诱导恒等的 $X_{n+1}$ 的 $S_{n+1}$-自同构层，典范同构于
\[
\mathcal{G} = \mathfrak{g}_{X_0/S_0} \otimes_{\mathcal{O}_{S_0}} \mathrm{gr}_{\mathcal{I}}^{n+1}(\mathcal{O}_S)
\]
(作为群层).
\end{proposition}

\begin{proof}
事实上，由 \ref{cor:III.5.4} 和 \ref{lem:III.4.2}，这个层是 $\mathcal{G}$ 下的主齐性层. 由于它带有特惠截面 (即 $X_{n+1}$ 的恒等自同构)，它因此作为集合层等同于 $G$. 必须验证这个等同与群结构相容. 这是容易的，而且它是 \ref{prop:III.5.1} 和 \ref{cor:III.5.3} 中主丛结构关于态射复合的相容性的更一般结果的特殊情形 (我们在此不陈述该结果，但它应当载入大全 (hyperplodoque)).
\end{proof}

特别地，$X_0$ 上 $X_{n+1}$ 的自同构芽层 (带所述结构) 是交换的. 由此得出，若 $X'_{n+1}$ 是该问题的另一个解，在 $X_0$ 的开集 $U$ 之上同构于 $X_{n+1}$，则由同构 $X_{n+1}|U \stackrel{\sim}{\longrightarrow} X'_{n+1}|U$ 通过结构转移得到的 $\mathbf{Aut}(X_{n+1})|U$ 到 $\mathbf{Aut}(X'_{n+1})|U$ 的同构不依赖于后者的选择. (它无非是 $\mathcal{G}$ 的恒等同构，当我们借助 \ref{prop:III.6.1} 把两个自同构层都等同于 $G$ 时.)

由 \ref{prop:III.6.1} 推出：

\begin{corollary}\label{cor:III.6.2}
设 $X_{n+1}$，$X'_{n+1}$ 在 $S_{n+1}$ 上光滑且``按 $X_n$ 约化''. 则在 $X_n$ 上诱导恒等的 $X_{n+1}$ 到 $X'_{n+1}$ 的 $S_{n+1}$-同构层 (在 $X_0$ 的底空间上) 自然地是 $\mathcal{G}$ 下的主齐性层.
\end{corollary}

\begin{proof}
这确实表示 $X_{n+1}$ 与 $X'_{n+1}$ 局部同构，且第一个的自同构芽层为 $\mathcal{G}$.
\end{proof}

现在注意，由 \ref{thm:III.4.1}，总可以找到 $X_n$ 的开覆盖 $(U_i)$ (可以假设它们仿射)，以及每个 $i$ 在 $S_{n+1}$ 上按 $U_i$ 约化的光滑概形 $X^i$. 为简化假设 $X_n$ 分离，故 $U_{ij} = U_i \cap U_j$ 仍是 $X_n$ 的仿射开集. 由于仿射开集的以拟凝聚层 $\mathcal{G}$ 为系数的 $H^1$ 为零，由推论 \ref{cor:III.6.2} 得出 $X^i|U_{ij}$ 同构于 $X^j|U_{ij}$；设
\[
f_{ji} \colon X^i|U_{ij} \xrightarrow{\sim} X^j|U_{ij}
\]
是这样一个同构. 它由 $\mathcal{G}$ 在 $U_{ij}$ 上的一个截面差一个差确定. 对每个三元指标组令：
\[
f_{ji}^{(k)} = f_{ji}|U_{ijk}, \quad \text{其中 } U_{ijk} = U_i \cap U_j \cap U_k.
\]
若成立
\begin{equation}
f_{kj}^{(i)} f_{ji}^{(k)} = f_{ki}^{(j)}, \tag{1}
\label{eq:III.1}
\end{equation}
则诸 $X^i$ 可以通过 $f_{ji}$``粘合''，因此它们定义所求问题的解 $X = X_{n+1}$. 更一般地，若可以把诸 $f_{ji}$ 修改为
\begin{equation}
f'_{ji} = f_{ji} g_{ji} \quad (g_{ji} \in \Gamma(U_{ij}, \mathcal{G})) \tag{2}
\label{eq:III.2}
\end{equation}
使得诸 $f'_{ji}$ 满足上面的传递条件，则这样的解存在. 解存在的这个充分条件也是必要的，回忆：这样的解 $X$ 必须在每个 $U_i$ 上同构于 $X^i$，这使我们能够选择同构
\[
f_i \colon X|U_i \stackrel{\sim}{\longrightarrow} X^i,
\]
并定义满足粘合条件的
\[
f'_{ji} = (f_j|U_{ij})(f_i|U_{ij})^{-1} \colon X^i|U_{ij} \xrightarrow{\sim} X^j|U_{ij}.
\]

令
\begin{equation}
f_{ijk} = (f_{ki}^{(j)})^{-1} f_{kj}^{(i)} f_{ji}^{(k)}, \tag{3}
\label{eq:III.3}
\end{equation}
这是 $X^i|U_{ijk}$ 的自同构，借助 \ref{prop:III.6.1} 把它等同于 $\mathcal{G}$ 的截面. 通过留给读者的小小形式计算，看到它是开覆盖 $\mathcal{U} = (U_i)$ 以 $\mathcal{G}$ 为系数的 2-上闭链 $f$. 同样的计算表明，在 (2) 的条件下，对 $f'_{ij}$ 的粘合条件 (1) 等价于公式
\begin{equation}
f = dg, \tag{4}
\label{eq:III.4}
\end{equation}
其中 $g = (g_{ij})$ 看作 $\mathcal{U}$ 以 $\mathcal{G}$ 为系数的 1-上链. 故问题有解的必要且充分条件是上闭链 (3) 定义的 $H^2(\mathcal{U}, \mathcal{G})$ 中的上同调类为零. 还注意到，由于 $\mathcal{U} = (U_i)$ 是概形 $X_0$ 的仿射覆盖，$H^2(\mathcal{U}, \mathcal{G})$ 等同于 $H^2(X_0, \mathcal{G})$. 显然，在 $H^2(X_0, \mathcal{G})$ 中这样得到的上同调类不依赖于所考虑的仿射覆盖. 称它为把 $X_n$ 延拓为 $S_{n+1}$ 上光滑概形 $X_{n+1}$ 的障碍类.

假设这个障碍为零. 则上面概述的推理表明，每个解 $X = X_{n+1}$ 同构于由同构 $f'_{ji}$ 粘合得到的解，可以假设 $f'_{ji}$ 形如 (2)，粘合条件正是 (3). 故允许的 $g$ 的集合是以 $\mathcal{U}$ 以 $\mathcal{G}$ 为系数的 1-上闭链群 $Z^1(\mathcal{U}, \mathcal{G})$ 为主齐性空间. 此外立即看到，两个上链 $g$ 和 $g'$ (满足 $dg = dg' = f$) 定义同构的解，当且仅当上闭链 $g - g'$ 形如 $dh$，其中 $h = (h_i) \in C^0(\mathcal{U}, \mathcal{G})$. 于是得到：

\begin{theorem}\label{thm:III.6.3}
设 $(S, \mathcal{I}, X_n)$ 如上，假设 $X_n$ 分离\footnote{这个条件实际上不必要，可以避免上面的上闭链计算. 参见 J. Giraud 的书《Cohomologie Non Abélienne》(将出版于 Springer Verlag 1971). 比较注记 \ref{rem:III.6.4}.}. 则可以典范地定义 $H^2(X_0, \mathcal{G})$ (其中 $\mathcal{G}$ 在 \ref{prop:III.6.1} 中定义) 中的障碍类，其消没是存在按 $X_n$ 约化的、在 $S_{n+1}$ 上光滑的概形 $X_{n+1}$ 的必要且充分条件. 若这个障碍为零，则按 $X_n$ 约化的 $S_{n+1}$-预概形 $X_{n+1}$ 的 (在 $X_n$ 上诱导恒等的) 同构类集合，自然地是 $H^1(X_0, \mathcal{G})$ 下的主齐性空间.
\end{theorem}

\begin{remark}\label{rem:III.6.4}
从 \ref{prop:III.6.1} 出发，这里所做的推理纯属形式，并且可以有利地搬到局部范畴甚至一般纤维范畴的框架中. 在一般语境下这样得到的``整体''对象存在性的障碍类 (对其可以局部找到一个对象，且任意两个对象总是``局部同构''，每个对象的自同构群交换) 包含作为特殊情形：不必交换的群层正合序列中的``第二个边缘同态'' (例如 Grothendieck 在 Kansas 或 Tôhoku 中所研究的). 这里所做的笨拙上闭链计算应被看作权宜之计，因为尚不存在令人满意的参考文献.
\end{remark}

\refstepcounter{theorem}\noindent\textbf{\thesection.\arabic{theorem}.}\quad 注意，在 \ref{thm:III.6.3} 中，一般没有 $H^1(X_0, \mathcal{G})$ 之下所考虑主齐性空间中的特惠元素. 这尤其表现为：得到 (在 $S$ 上局部化) $S_0$ 上结构群为 $R^1 f_{*}(\mathcal{G})$ 的主齐性层，它不必平凡，i.e. 定义 $H^1(S_0, R^1 f_{*}(\mathcal{G}))$ 中不必为零的上同调类. (当假设类 $d \in H^2(X_0, \mathcal{G})$ 不为零、但在``$S$ 之上局部为零''，i.e. 定义 $R^2 f_{*}(\mathcal{G})$ 的零截面，i.e. $H^0(S_0, R^2 f_{*}(\mathcal{G}))$ 中的零元素时，正是如此.)

\refstepcounter{theorem}\noindent\textbf{\thesection.\arabic{theorem}.}\quad 目前对本节所引入上同调类的一般代数机制、以及它与前一节的关系几乎一无所知，即使在最简单的特殊情形 (例如 Artin 环上的阿贝尔概形的情形) 也说不出的确切的东西\footnote{现在知道在这个情形障碍总是为零 [2003 年添加 (MR)：参见 F. Oort，Finite group schemes，local moduli for abelian varieties and lifting problems，Algebraic Geometry Oslo 1970，Wolters-Noordhoff，1972，p. 223--254].}. 希望有人会潜心研究这个问题，它似乎特别有趣. 它与代数结构的``模理论'' (moduli theory) 尤其密切相关.

\begin{corollary}\label{cor:III.6.7}
假设 $H^2(X_0, \mathcal{G}) = 0$，则 $X_{n+1}$ 存在，且若还有 $H^1(X_0, \mathcal{G}) = 0$ 则在同构意义下唯一.
\end{corollary}

\begin{proof}
特别地，逐步进行 (并注意仿射概形对拟凝聚层是无轮回 (acyclic) 的)，得出结论：
\end{proof}

\begin{corollary}\label{cor:III.6.8}
在定理 \ref{thm:III.4.1} 的条件下，若 $X_0$ 仿射，则存在按 $X_0$ 约化的、在 $Y$ 上光滑的 $X$，且这个 $X$ 在同构 (不唯一) 意义下唯一.
\end{corollary}

注意定理 \ref{thm:III.4.1} 的直接证明不能给我们这个结果.

\begin{corollary}\label{cor:III.6.9}
在 \ref{thm:III.6.3} 的条件下，假设 $S$ 仿射，环为 $A$，$\mathcal{I}$ 由 $A$ 的理想 $I$ 定义，最后诸 $\mathrm{gr}_{\mathcal{I}}^{n}(\mathcal{O}_S) = \mathcal{I}^n/\mathcal{I}^{n+1}$ 局部自由. 则 $H^i(X_0, \mathcal{G})$ 等同于 $H^i(X_0, \mathcal{G}_0) \otimes_A \mathrm{gr}_{I}^{n+1}(A)$，其中
\[
\mathcal{G}_0 = \mathfrak{g}_{X_0/S_0},
\]
故 $X_n$ 的延拓障碍类位于 $H^2(X_0, \mathcal{G}_0) \otimes_A \mathrm{gr}_{I}^{n+1}(A)$，且若它为零，解的 (同构) 类集合是 $H^1(X_0, \mathcal{G}_0) \otimes_A \mathrm{gr}_{I}^{n+1}(A)$ 下的主齐性空间.
\end{corollary}

\begin{proof}
特别地：
\end{proof}

\begin{corollary}\label{cor:III.6.10}
在 \ref{cor:III.6.9} 的条件下假设
\[
H^2(X_0, \mathfrak{g}_{X_0/S_0}) = 0,
\]
则在 $S$ 的形式 $\mathcal{I}$-adic 完备化 $\widehat{S}$ 上存在 $\mathcal{I}$-adic 形式概形 $X$，它``在 $S$ 上光滑''(i.e. 诸 $\mathfrak{X}_p$ 在诸 $S_p$ 上光滑) 且按 $X_n$ 约化，i.e. 配备同构
\[
\mathfrak{X} \times_S S_n \stackrel{\sim}{\longleftarrow} X_n.
\]
若还有 $H^1(X_0, \mathfrak{g}_{X_0/S_0}) = 0$，则这样的 $X$ 在同构意义下唯一.
\end{corollary}

\begin{proof}
事实上，逐步构造 $X_{n+1}$，$X_{n+2}$ 等，取诸 $X_i$ 的归纳极限得到 $X$. 唯一性断言已见于前一节.
\end{proof}

\section{在完备局部环 \texorpdfstring{$A$}{A} 上构造光滑形式概形与普通概形的应用}

前一节的结果有时能证明：在一个这样的环上，存在按 $k$ 上给定光滑概形 $X_0$ 约化的 $m$-adic 形式概形. 区分两种情形：

\begin{enumerate}[a)]
\item $A$ 是``等特征''的 (当 $k$ 特征为 0 时尤其如此). 此时知道 $A$ 存在代表元子域 (subfield of representatives)，i.e. 子域 $k'$ 使得 $A \to k$ 诱导同构 $k' \xrightarrow{\sim} k$. 则甚至存在按 $X_0$ 约化的 $A$ 上光滑普通概形，即 $X = X_0 \otimes_k A$，其中 $A$ 借助由 $k'$ 定义的同态 $k \to k' \to A$ 看作 $k$ 上的代数. 但必须注意这个构造不是``自然的''；容易相信 (已经可以在 $A = k[t]/(t^2)$，对偶数代数的情形) 另一个提升同态 $k \to A$ (此处由 $k$ 到自身的绝对导子定义) 定义 $A$ 上的 $X'$，它一般不同构于 $X$ (若 $H^1(X_0, \mathfrak{g}_{X_0/k}) \neq 0$). 研究 (对特征 0 或非完美特征 $p > 0$ 的 $k$) 这样得到的 $A$ 上光滑 $X$ 是哪些、两个同态 $k \to A$ 在什么条件下定义同构的 $A$-概形，将是很有趣的. 然而，$k'$ 的存在足以推出 $X_0$ 提升的第一个障碍 (位于 $H^2\bigl(X_0, \mathfrak{g}_{X_0/k}\bigr) \otimes_k \mathfrak{m}/\mathfrak{m}^2$) 必然为零. 当然，把 $X_0$ 提升为 $A/\mathfrak{m}^2$ 上光滑的 $X_1$ 后，构造 $X_2$ 的新障碍一般不为零：它依赖于 $H^1(X_0, \mathcal{G}_0) \otimes \mathfrak{m}/\mathfrak{m}^2$ 之下某个主齐性空间中一个可变元素，且位于 $H^2(X_0, \mathcal{G}_0) \otimes \mathfrak{m}^2/\mathfrak{m}^3$：应当详细研究这一情形\footnote{它大概由 Kodaira-Spencer 的括号运算描述 (参见 Séminaire Cartan，1960/61，Exp. 4).}.
\item $A$ 是不等特征的. 此时几乎一无所知 (除非碰巧 $H^2(X_0, \mathfrak{g}_{X_0/k}) = 0$，此时可以构造按 $k$ 约化的 $A$ 上 $m$-adic 光滑形式概形). 即使 $A = \mathbf{Z}/p^2\mathbf{Z}$ 且 $X_0$ 是维数 2 的``阿贝尔''概形，也不知道能否把它提升为 $A$ 上光滑的 $X = X_1$\footnote{现在这已被证明，参见第 67 页注 7.}；另一方面，没有例子说明某个 $X_0$ 已被证明不来自 $A$ 上光滑普通概形 $X$. (我觉得这应当存在，取 $X_0$ 为射影曲面)\footnote{这样的例子后来由 J.-P. Serre 构造 (Proc. Nat. Acad. Sc. USA 47 (1961)，No 1，p. 108--109)，至少在若干维数. D. Mumford 给出了一个例子 (未发表)，用代数曲面.}. 只指出，由 Cohen 定理，存在剩余域为 $k$ 的 Cohen $p$-环 $B$ 以及诱导剩余域上恒等同构的同态 $B \to A$；因此，最强的提升结果应取 $A$ 为 Cohen $p$-环：若在一个这样的环之上存在 (普通或形式的) 解，则在剩余域为 $k$ 的每个完备局部环之上也存在. 特别地，由于对 Cohen $p$-环，$\mathfrak{m}/\mathfrak{m}^2$ 典范等同于 $k$，看到对特征 $p > 0$ 的域 $k$ 上的每个光滑概形 $X_0$，存在 $H^2\bigl(X_0, \mathfrak{g}_{X_0/k}\bigr)$ 中的上同调类，即把 $X_0$ 提升为 Cohen $p$-环上光滑概形的第一个障碍；不知道它能否非零\footnote{它可以非零，如第 69 页注 10 所指出的.}.
\end{enumerate}

即使能逐步构造按 $X_0$ 约化的 $X_n$，一般也只能得到按 $X_0$ 约化的、$A$ 上光滑的形式概形 $X$. 当 $X_0$ 在 $A$ 上紧合时，还剩下问题：$X$ 是否事实上可代数化，以便得到按 $X_0$ 约化的、$A$ 上紧合且简单的普通概形. 唯一已知的判则 (在 Bourbaki 讨论班中提及，见于《原理》第三章 4.7.1) 如下：若 $X$ 在 $A$ 上紧合，且 $\mathcal{L}$ 是 $X$ 上的可逆层，使得在 $X_0$ 上诱导的层 $\mathcal{L}_0$ 丰沛 (ample) (i.e. 适当的张量幂 $\mathcal{L}_0^{\otimes n}$ ($n > 0$) 来自 $X_0$ 的射影浸入)，则存在 $A$ 上射影的概形 $X$ 和 $X$ 上的丰沛可逆层 $\mathcal{L}$，使得 $(X, \mathcal{L})$ 经 $m$-adic 完备化得到 $(\mathfrak{X}, \mathcal{L})$. 于是，给定 $X_0$ 上的局部自由层 $\mathcal{E}_0$ (对我们的目的选择丰沛可逆的)，要把它延拓为 $X$ 上的局部自由层 $\mathcal{E}$. 为此归结为在诸 $X_n$ 上逐步构造局部自由层 $\mathcal{E}_n$. 讨论完全类似于 §6 的讨论 (参见注记 \ref{rem:III.6.4})，基本角色由在 $\mathcal{E}_n$ 上诱导恒等的 $\mathcal{E}_{n+1}$ 的自同构层扮演：立即证明该层等同于
\begin{equation}
\mathcal{G} = \operatorname{Hom}_{\mathcal{O}_{X_0}}\bigl(\mathcal{E}_0, \mathcal{E}_0 \otimes \mathrm{gr}_{\mathcal{I}}^{n+1}(\mathcal{O}_X)\bigr) = \operatorname{Hom}_{\mathcal{O}_{X_0}}\bigl(\mathcal{E}_0, \mathcal{E}_0\bigr) \otimes \mathrm{gr}_{\mathcal{I}}^{n+1}(\mathcal{O}_X), \tag{($*$)}
\end{equation}
它仍是交换群层. 得到：

\begin{proposition}\label{prop:III.7.1}
设 $S$ 是配备拟凝聚理想层 $\mathcal{I}$ 的预概形，$X$ 是 $S$ 上的预概形，$S_n$ 是由 $\mathcal{I}^{n+1}$ 定义的 $S$ 的子预概形，$X_n = X \times_S S_n$ (对每个整数 $n$). 设 $\mathcal{E}_n$ 是 $X_n$ 上的局部自由层，要把它延拓为 $X_{n+1}$ 上的局部自由层 $\mathcal{E}_{n+1}$. 则 $\mathcal{E}_n$ 定义 $H^2(X_0, \mathcal{G})$ 中的典范障碍类 (其中 $G$ 是上式给出的拟凝聚层)，该类消没是存在延拓 $\mathcal{E}_n$ 的 $\mathcal{E}_{n+1}$ 的必要且充分条件. 若该类为零，则 (在 $\mathcal{E}_n$ 上诱导恒等的) 同构意义下的解类 $\mathcal{E}_{n+1}$ 的集合，是 $H^1(X_0, \mathcal{G})$ 下的主齐性空间.
\end{proposition}

这个命题给出通常的推论. 只指出，若 $X$ 在 $S$ 上平坦，则可以写
\[
\mathcal{G} = \operatorname{Hom}_{\mathcal{O}_{X_0}}\bigl(\mathcal{E}_0, \mathcal{E}_0\bigr) \otimes_{\mathcal{O}_{S_0}} \mathrm{gr}_{\mathcal{I}}^{n+1}(\mathcal{O}_S),
\]
从而，若 $S$ 仿射环为 $A$ 且诸 $\mathcal{I}^n/\mathcal{I}^{n+1}$ 局部自由，得充分条件
\[
H^2(X_0, \mathcal{G}_0) = 0, \quad \text{其中} \quad \mathcal{G}_0 = \mathbf{Hom}_{\mathcal{O}_{X_0}}\bigl(\mathcal{E}_0, \mathcal{E}_0\bigr),
\]
它保证 $\mathcal{E}_{n+1}$ 存在，从而逐步保证逐次延拓 $\mathcal{E}_m$ ($m = n, n+1$，等) 存在.

回到初始情形，得到：

\begin{proposition}\label{prop:III.7.2}
设 $A$ 是完备局部环，$X$ 是 $A$ 上紧合且平坦的形式概形，使得 $X_0$ 射影且 $H^2(X_0, \mathcal{O}_{X_0}) = 0$. 则存在 $A$ 上射影的概形 $X$，其 $m$-adic 形式完备化同构于 $X$.
\end{proposition}

与 \ref{cor:III.6.10} 结合，得到：

\begin{theorem}\label{thm:III.7.3}
设 $A$ 是剩余域为 $k$ 的完备局部环，$X_0$ 是 $k$ 上射影且光滑的概形，使得
\[
H^2(X_0, \mathfrak{g}_{X_0/k}) = H^2(X_0, \mathcal{O}_{X_0}) = 0.
\]
则存在按 $X_0$ 约化的、$A$ 上光滑且射影的概形 $X$.
\end{theorem}

更一般地，若给出按 $X_0$ 约化的、在 $A_n = A/\mathfrak{m}^{n+1}$ 上光滑的 $X_n$，则存在 $A$ 上光滑且紧合的 $X$ 以及同构 $X \otimes_A A_n = X_n$.

\begin{corollary}\label{cor:III.7.4}
$k$ 上的每条光滑紧合曲线都经约化来自 $A$ 上的光滑紧合曲线.
\end{corollary}

\begin{proof}
正是这个结果 (连同形式几何中层存在性定理) 将是经由超验途径研究 $X_0$ 的基本群的基本工具.
\end{proof}

\end{refsection}
\begin{refsection}% =====================================================================
% 第四讲 (Exposé IV) 平坦态射 (Morphismes plats)
% A. Grothendieck
% 原书页码: pp. 71-86
% =====================================================================
\chapter{平坦态射}

我们在这里主要给出在前面各讲中用到的\emph{平坦性} (flatness) 性质. 更详细的研究将在准备中的《代数几何原理》第四章\footnote{参见 EGA IV 11 和 12.}中给出，那里系统研究下列情形：$X$ 在局部诺特 $Y$ 上局部有限型，$\mathcal{F}$ 在 $X$ 上凝聚\index[term]{ningjuceng@凝聚层}且 $Y$-平坦，给出 $Y$ 的性质、$\mathcal{F}$ 的性质、以及 $\mathcal{F}$ 在 $X \to Y$ 的纤维上诱导的凝聚层的性质之间的关系 (特别是从维数、上同调维数、深度等观点). 尤其有系统获得 Seidenberg 型或 Bertini 型定理 (对超平面截面的) 的方法. 在这一语境下应用平坦性方法的关键结果如下 (将在下面证明)：若 $Y$ 整，$X$ 在 $Y$ 上有限型，$\mathcal{F}$ 在 $X$ 上凝聚，则存在 $Y$ 的非空开集 $U$，使得 $\mathcal{F}$ 在 $U$ 之上的 $X$ 的点处 $Y$-平坦. 平坦性进入代数几何的第二种方式，大概更为重要，是下降 (descent) 理论：参见例如 Grothendieck 在 Bourbaki 讨论班就这一主题的两个报告\footnote{以及，更详细的报告，下面的讲 VIII 和 IX.}. 平坦性因而似乎是代数几何中一个中心技术概念.

回忆平坦性和忠实平坦性概念是由 Serre 在 GAGA 中引入的\index[term]{pingtanxing@平坦性}. 下面第 1、2 节的阐述也可以见 Bourbaki 的《交换代数》(Alg. Comm.) (当然，如书名所示，它不限于交换基环的情形)\footnote{N. Bourbaki，Algèbre Commutative，Chap. I (Modules Plats)，Act. Sc. Ind 1290，Paris，Hermann (1961).}.

与前几讲相反，我们不假设所考虑的环必然是诺特的.

\section{平坦模的平凡推论}

环 $A$ 上的模 $M$ 称为平坦的 (flat) (若想指明 $A$ 则称 $A$-平坦)，若函子
\[
T_M \colon N \longmapsto M \otimes_A N
\]
(它总是右正合的) 正合，i.e. 把单同态变为单同态. 这等价于说第一个右导出函子，或所有右导出函子为零，i.e. 有
\[
\mathrm{Tor}_{1}^{A}(M, N) = 0 \qquad \text{对每个 } N,
\]
相应地有
\[
\operatorname{Tor}_{i}^{A}(M, N) = 0 \quad \text{对 } i > 0, \text{每个 } N.
\]
由于诸 $\mathrm{Tor}_i$ 与归纳极限可交换，只需对有限型 $N$ 验证这些条件，甚至 (取 $N$ 的以单生成模为商的合成序列) 只需有
\[
\mathrm{Tor}_{1}^{A}(M, N) = 0
\]
若 $N$ 单生成，i.e. 形如 $A/I$，$I$ 是 $A$ 的理想. 还注意
\[
\mathrm{Tor}_{1}^{A}(M, A/I) = 0 \quad \Longleftrightarrow \quad I \otimes_A M \to M = A \otimes_A M \text{ 单射},
\]
这从 Tor 的正合序列 (考虑到 $\mathrm{Tor}_{1}^{A}(M, A) = 0$) 看出. 故 $M$ 平坦等价于说：对每个理想 $I$，自然同态
\[
I \otimes_A M \longrightarrow IM
\]
是同构. 只需对有限型 $I$ 验证，更不用说只需验证函子 $M \otimes$ 在有限型模上正合.

每当我们有一个正合函子 $T$ 时，若把 $N$ 的子对象 $N'$ 的 $T(N')$ 等同于 $T(N)$ 的子对象，则对 $N$ 的两个子对象 $N'$，$N''$ 有
\[
\begin{array}{l}
T(N' \cap N'') = T(N') \cap T(N''), \\
T(N' + N'') = T(N') + T(N'').
\end{array}
\]

平坦模的直和、平坦模的直因子都是平坦的. 特别地，$A$ 平坦，自由模从而射影模平坦. 两个平坦模的张量积平坦，且若 $M$ 在 $A$ 上平坦，则对每个基变换 $A \to B$，$M \otimes_A B$ 在 $B$ 上平坦 (因为张量积的结合律和正合函子的复合仍正合). 若 $M$ 在 $B$ 上平坦，$B$ 在 $A$ 上平坦，则 $M$ 在 $A$ 上平坦 (同理).

Tor 的正合序列加上 Tor 的``交换性''，给出：

\begin{proposition}\label{prop:IV.1.1}
设 $0 \to M' \to M \to M'' \to 0$ 是 $A$-模的正合序列，$M''$ 平坦. 则
\begin{enumerate}[(i)]
\item 该序列与任何 $A$-模 $N$ 张量后仍正合; 
\item 为使 $M$ 平坦，必要且充分条件是 $M'$ 平坦.
\end{enumerate}
\end{proposition}

可以说，从张量积的行为的观点看，平坦模与自由模或射影模``一样好'' (而 \ref{prop:IV.1.1} 中的正合序列``一样好''，仿佛它分裂).

设 $S$ 是 $A$ 的乘法稳定子集，则 $S^{-1}A$ 在 $A$ 上平坦，因为 $S^{-1}A \otimes N = S^{-1}N$ 是 $N$ 的正合函子. 若 $M$ 是 $A$-平坦的，则 $S^{-1}M = S^{-1}A \otimes M$ 是 $S^{-1}A$-平坦的，逆命题在 $M \to S^{-1}M$ 是同构时成立，i.e. 若 $S$ 的元素在 $M$ 中是双射 (因为平坦性的传递性，$S^{-1}A$ 在 $A$ 上平坦). 更一般地，预概形态射 $X \to Y$ 以及要研究相对 $Y$ 的平坦性的 $X$ 上拟凝聚层 $\mathcal{F}$ 的情形，导致带两个环的情形：

\begin{proposition}\label{prop:IV.1.2}
设 $A \to B$ 是环同态，$M$ 是 $B$-模，$T$ 是 $B$ 的乘法稳定子集.
\begin{enumerate}[(i)]
\item 若 $M$ 是 $A$-平坦的，则 $T^{-1}M$ 是 $A$-平坦的 (从而对 $A$ 的任意映入 $T$ 的乘法稳定子集 $S$，也是 $S^{-1}A$-平坦的).
\item 反过来，若对 $B$ 的每个极大理想 $n$，$M_{\mathfrak{n}}$ 在 $A_{\mathfrak{n}}$ 上平坦 ($M_{\mathfrak{n}}$ 在 $A$ 上平坦，或等价地在 $A_{\mathfrak{m}}$ 上平坦，其中 $m$ 是 $n$ 的原像)，则 $M$ 是 $A$-平坦的.
\end{enumerate}
\end{proposition}

\begin{proof}
事实上，有对 $A$-模 $N$ 函子性的公式：
\[
T^{-1}M \otimes_A N = T^{-1}(M \otimes_A N),
\]
因为由 $\otimes$ 的结合律公式，两端函子性地同构于 $T^{-1}B \otimes_B M \otimes_B N_{(B)}$ (其中 $N_{(B)} = N \otimes_A B$). 立即得出，若 $M \otimes_A N$ 对 $N$ 正合，则 $T^{-1}\mathcal{M} \otimes_A \mathcal{N}$ 亦然 (作为两个正合函子的复合)，由此得 (i). 同样得出 (ii)，因为为验证 $B$-模序列的正合性，只需验证它在 $B$ 的所有极大理想处的局部化正合.
\end{proof}

\begin{proposition}\label{prop:IV.1.3}
(i) 设 $M$ 是 $A$-平坦模\index[term]{pingtanmo@平坦模}. 若 $x \in A$ 在 $A$ 中非零因子，则它在 $M$ 中非零因子. 特别地，若 $A$ 整，则 $M$ 无挠.
(ii) 假设 $A$ 整，且对 $A$ 的每个极大理想 $m$，$A_{\mathfrak{m}}$ 是主理想环 (例如 $A$ 是 Dedekind 环，甚至是主理想环). 为使 $A$-模 $M$ 平坦，必要且充分条件是它无挠.
\end{proposition}

\begin{proof}
(i) 由下述事实得出：$M$ 中的倍扩 $x$ 由 $A$ 中的倍扩 $x$ 与 $M$ 张量得到. 对 (ii)，借助 \ref{prop:IV.1.2} (ii) 可以先假设 $A$ 主; 必须证明若 $M$ 无挠，则对 $A$ 的每个理想 $I$，单射 $I \to A$ 与 $M$ 张量后仍是单射，这意味着 $I$ 的生成元 $x$ 在 $M$ 中非零因子，好.
\end{proof}

\section{忠实平坦模}

从一个范畴到另一个范畴的函子 $F$ 称为忠实的 (faithful)，若对每个 $X$，$Y$，映射 $\mathrm{Hom}(X, Y) \to \operatorname{Hom}\bigl(F(X), F(Y)\bigr)$ 单射. 若是加性范畴的加性函子，这等价于说 $F(u) = 0$ 蕴含 $u = 0$，且这蕴含 $F(X) = 0$ 蕴含 $X = 0$. 为使 $F$ 忠实且正合，必要且充分条件是：对 $\mathcal{C}$ 中态射的任何序列 $M' \to M \to M''$，变换后的序列 $F(M') \to F(M) \to F(M'')$ 正合当且仅当前一个正合. 或：$F$ 正合，且 $F(X) = 0$ 蕴含 $X = 0$. (N.B. 为能谈论正合性，必须假设所涉及的范畴是 Abel 的.) 假设有 $\mathcal{C}$ 的非零对象族 $(M_i)$，使得 $\mathcal{C}$ 的每个非零对象都有一个子对象，其商同构于某个 $M_i$. 则 $F$ 忠实且正合等价于：$F$ 正合，且对每个 $i$ 有 $F(M_i) \neq 0$. 若 $\mathcal{C}$ 是环 $A$ 上的模范畴，例如可取 $(M_i)$ 为诸 $A/\mathfrak{m}$ 组成的族，$m$ 遍历 $A$ 的极大理想. (事实上，每个非零模都有非零单生成子模，故同构于某个 $A/I$ ($I \neq A$)，后者由 Krull 定理有一个同构于某个 $A/\mathfrak{m}$ 的商.) 由这些平凡推论特别得出：

\begin{proposition}\label{prop:IV.2.1}
设 $M$ 是 $A$-模. 下列条件等价：
\begin{enumerate}[(i)]
\item 函子 $M \otimes_A$ 忠实且正合.
\item[(i bis)] $M$ 平坦，且 $M \otimes_A N = 0$ 蕴含 $N = 0$.
\item[(i ter)] $M$ 平坦，且对 $A$ 的每个极大理想 $m$，$M \otimes A/\mathfrak{m} \neq 0$.
\item 对任何同态序列 $N' \to N \to N''$，与 $M$ 张量后的序列正合当且仅当初始序列正合.
\end{enumerate}
\end{proposition}

此时称 $M$ 是 $A$-\emph{忠实平坦模} (faithfully flat module)\index[term]{zhongshipingtanmo@忠实平坦模}. 特别地，若 $M$ 忠实平坦，则 $N \to N'$ 是单同态 (满同态，同构) 当且仅当与 $M$ 张量后的同态亦然. 忠实平坦模是忠实的，因为 $M$ 中的倍扩 $f$ 由 $A$ 中的倍扩 $f$ 与 $M$ 张量得到.

像 §1 中那样看到通常的传递性性质：两个忠实平坦模的张量积忠实平坦; 若 $M$ 在 $A$ 上忠实平坦，则对任何基扩张 $A \to B$，$M \otimes_A B$ 在 $B$ 上忠实平坦; 若 $B$ 是在 $A$ 上忠实平坦的 $A$-代数且 $M$ 是忠实平坦 $B$-模，则 $M$ 是忠实平坦 $A$-模.

\begin{corollary}\label{cor:IV.2.2}
设 $A \to B$ 是局部环的局部同态，$M$ 是有限型 $B$-模. 为使 $M$ 在 $A$ 上忠实平坦，必要且充分条件是它在 $A$ 上平坦且非零.
\end{corollary}

\begin{proof}
由判则 (i ter) 和 Nakayama 引理得出. 特别地，为使 $B$ 是 $A$-平坦的，必要且充分条件是它是忠实 $A$-平坦的.
\end{proof}

\begin{proposition}\label{prop:IV.2.3}
设 $A \to B$ 是环同态，$M$ 是在 $A$ 上忠实平坦的 $B$-模. 对 $A$ 的每个素理想 $p$，存在 $B$ 的素理想 $q$ 诱导它.
\end{proposition}

\begin{proof}
除以 $p$，归结到 $\mathfrak{p} = 0$ 的情形. 在素理想 0 处局部化，归结到 $A$ 是域的情形. 但 $M$ 在 $A$ 上忠实平坦故非零，更不用说 $B \neq 0$，故 $B$ 有素理想，它只能诱导 $A$ 的惟一的素理想! 几何上可以说，在 $X = \operatorname{Spec}(B)$ 上存在相对 $A$``忠实平坦''的拟凝聚层 $\mathcal{F}$，蕴含 $X \to Y = \operatorname{Spec}(A)$ 满射.
\end{proof}

\begin{corollary}\label{cor:IV.2.4}
假设 $M$ 在 $A$ 上平坦，在 $B$ 上有限型，且 $\mathrm{Supp}\, M = \mathrm{Spec}(B)$ (i.e. 对 $B$ 的每个素理想 $q$ 有 $M_{\mathfrak{q}} \neq 0$). 则 $B$ 的包含 $pB$ 的极小素理想 $q$ 诱导 $p$.
\end{corollary}

\begin{proof}
仍归结到 $\mathfrak{p} = 0$ 的情形 (因为所有假设除以后都保持)，故 $A$ 整. 归结为下述陈述：
\end{proof}

\begin{corollary}\label{cor:IV.2.5}
$M$ 如上，$B$ 的每个极小素理想 $q$ 诱导 $A$ 的素理想 $p$，它是极小的.
\end{corollary}

\begin{proof}
事实上，在 $p$ 和 $q$ 处局部化，归结为证明：若 $A$ 和 $B$ 局部，同态 $A \to B$ 局部，$M$ 是在 $A$ 上平坦的非零 $B$-模，且 $\dim B = 0$，则 $\dim A = 0$. 由 \ref{prop:IV.2.3} 和 \ref{cor:IV.2.2}，结论是 $A$ 的每个素理想都由 $B$ 的素理想诱导，从而由 $B$ 的极大理想诱导，故是极大理想，证毕. 几何上，\ref{cor:IV.2.5} 意味着 $X = \operatorname{Spec}(B)$ 的每个不可约分支都支配 $Y = \operatorname{Spec}(A)$ 的某个不可约分支 (在 $X$ 上存在支集为 $X$、相对 $Y$ 平坦的有限型拟凝聚层的条件下).
\end{proof}

注意在 \ref{cor:IV.2.4} 中不必假设 $M$ 在 $A$ 上忠实平坦，但此时不保证存在包含 $pB$ 的素理想 (从而不保证存在其中极小的).

\begin{proposition}\label{prop:IV.2.6}
设 $i \colon A \to B$ 是环同态. 下列条件等价：
\begin{enumerate}[(i)]
\item $B$ 是忠实平坦 $A$-模.
\item $B$ 在 $A$ 上平坦，且 $\operatorname{Spec}(B) \to \operatorname{Spec}(A)$ 满射.
\item[(ii bis)] $B$ 在 $A$ 上平坦，且每个极大理想由 $B$ 的理想诱导.
\item $i$ 单射，且 Coker $i$ 是 $A$-平坦模.
\item 函子 $M_{(B)} = M \otimes_A B$ (关于 $A$-模 $M$) 正合，且典范函子同态 $M \to M_{(B)}$ 单射.
\item[(iv bis)] 对 $A$ 的每个理想 $I$，$I \otimes_A B \to IB$ 是同构，且 $IB$ 在 $A$ 中的原像等于 $I$.
\end{enumerate}
\end{proposition}

\begin{proof}
有 (i)$\Rightarrow$(ii) (由 \ref{prop:IV.2.3})，(ii)$\Rightarrow$(ii bis) 平凡，(ii bis)$\Rightarrow$(i) 由 \ref{prop:IV.2.1} 的判则 (i ter). 有 (iii)$\Rightarrow$(iv) (由 \ref{prop:IV.1.1})，(iv)$\Rightarrow$(iv bis) 平凡 (在第二个条件中取 $M = A/I$)，(iv bis)$\Rightarrow$(i) 由 §1 开头所见的理想平坦性判则和判则 \ref{prop:IV.2.1} (i ter). 最后，(iv)$\Rightarrow$(iii) 由 \ref{prop:IV.1.1} 的简单逆命题，(i)$\Rightarrow$(iv) 因为若 $N$ 是 $M \to M \otimes_A B = T(M)$ 的核，则 ($T$ 正合) $N \to T(N)$ 为零，从而 $T(N) = N \otimes_A B = 0$，故 $N = 0$，证毕.
\end{proof}

\section{与完备化的关系}

设 $A$ 是诺特环，$I$ 是 $A$ 中的理想，$\widehat{A}$ 是 $A$ 按 $I$-预adic 拓扑的分离完备化，对每个 $A$-模 $M$，$\widehat{M}$ 是它按 $I$-预adic 拓扑的完备化. 它是 $A$-模，于是有典范同态
\[
M \otimes_A \widehat{A} \longrightarrow \widehat{M}.
\]
当 $M$ 遍历有限型模时，函子 $M \mapsto \widehat{M}$ 正合，这容易由 Krull 定理得出：若 $N \subset M$，则 $N$ 的拓扑是 $M$ 的拓扑诱导的. 由于 $M \otimes_A \widehat{A}$ 右正合，容易由此推出 (把 $M$ 用 $L \to L' \to M$ 解出，其中 $L$ 和 $L'$ 有限自由) 上面的函子同态是同构 ($\widehat{M}$ 也右正合)，从而 $M \otimes_A \widehat{A}$ 也是 $M$ 的正合函子. 因此：

\begin{proposition}\label{prop:IV.3.1}
设 $A$ 是诺特环，$I$ 是 $A$ 的理想，则 $A$ (按 $I$-预adic 拓扑) 的分离完备化 $\widehat{A}$ 在 $A$ 上平坦.
\end{proposition}

\begin{corollary}\label{cor:IV.3.2}
为使 $\widehat{A}$ 在 $A$ 上忠实平坦，必要且充分条件是 $I$ 含于 $A$ 的根式 (radical) 中.
\end{corollary}

\begin{proof}
事实上，只需应用判则 \ref{prop:IV.2.1} (i ter).
\end{proof}

这些结果概括了从线性代数的观点所能说的、关于 $A$ 与 $\widehat{A}$ 之间关系的一切. 推论 \ref{cor:IV.3.2} 在 $A$ 是诺特局部环且 $I$ 含于极大理想 $m$ 中 (通常与之相等) 时最常用.

\section{与自由模的关系}

\begin{proposition}\label{prop:IV.4.1}
设 $A$ 是环，$I$ 是 $A$ 的理想，$M$ 是 $A$-模. 假设处于下列假设之一：
\begin{enumerate}[a)]
\item $I$ 幂零; 
\item $A$ 诺特，$I$ 在 $A$ 的根式中，且 $M$ 有限型.
\end{enumerate}
为使 $M$ 在 $A$ 上自由，必要且充分条件是 $M \otimes A/I$ 在 $A/I$ 上自由，且 $\mathrm{Tor}_{1}^{A}(M, A/I) = 0$.
\end{proposition}

\begin{proof}
必要性显然，证明充分性. 设 $(e_i)$ 是 $M$ 的元素族，其在 $M \otimes A/I = M/IM$ 中的像定义 $A/I$ 上的基 (在情形 (b) 中是有限族). 设 $L$ 是在同一个指标集上构造的 $A$-自由模，于是有同态 $L \to M$，其与 $A/I$ 的张量化 $T$ 诱导同构 $T(L) \stackrel{\sim}{\longrightarrow} T(M)$. 若 $Q$ 是 $L \to M$ 的余核，则有 $T(Q) = 0$，由 Nakayama 引理 (在条件 (a) 或 (b) 下成立) 得 $Q = 0$. 故 $L \to M$ 满射，设 $R$ 是其核，则有正合序列
\[
0 \longrightarrow R \longrightarrow L \longrightarrow M \longrightarrow 0,
\]
由此，由于 $\operatorname{Tor}_{1}^{A}(M, A/I) = 0$，得正合序列 $0 \to T(R) \to T(L) \to T(M) \to 0$，故 $T(R) = 0$，再由 Nakayama 引理 (考虑到在情形 (b) 中，$A$ 诺特故 $R$ 有限型) 得 $R = 0$.
\end{proof}

\begin{corollary}\label{cor:IV.4.2}
可以把条件 $\mathrm{Tor}_{1}^{A}(M, A/I) = 0$ 换成：典范满同态
\begin{equation}
\mathrm{gr}_{I}^{0}(M) \otimes_{A/I} \mathrm{gr}_{I}(A) \longrightarrow \mathrm{gr}_{I}(M) \tag{($*$)}
\end{equation}
是同构.
\end{corollary}

\begin{proof}
事实上，若 $M$ 自由，这当然成立. 所以必须证明：若 $M \otimes A/I$ 在 $A/I$ 上自由且关于 gr 的条件成立，则 $M$ 自由. 重做上面的证明，构造 $L \to M$，由假设该同态诱导相伴分次环的同构，故其核含于诸 $I^n L$ 的交，从而为零 (在 (a) 中平凡，在 (b) 中众所周知). 证毕.
\end{proof}

\begin{corollary}\label{cor:IV.4.3}
假设 $A/I$ 是域. 则 $M$ 的下列条件等价：
\begin{enumerate}[(i)]
\item $M$ 自由; 
\item $M$ 射影; 
\item $M$ 平坦; 
\item $\mathrm{Tor}_{1}^{A}(M, A/I) = 0$; 
\item 典范同态 (∗) 双射.
\end{enumerate}
\end{corollary}

\begin{proof}
事实上，在所考虑的情形，$M \otimes A/I$ 自动自由.
\end{proof}

前面的结果在下列两种情形下成立：
\begin{enumerate}[a)]
\item $M$ 是局部环 $A$ (其极大理想 $I$ 幂零，例如 Artin 局部环) 上的任意模.
\item $M$ 是诺特局部环上的有限型模.
\end{enumerate}

为备忘回忆：

\begin{corollary}\label{cor:IV.4.4}
假设 $A$ 是极大理想 $\mathfrak{m} = I$、剩余域 $k = A/I$、分式域 $K$ 的诺特整局部环. 设 $M$ 是 $A$ 上的有限型模. 则前面 (i)$\dots$(v) 的等价条件还等价于
\begin{enumerate}[(i)]
\item[((vi))] $M \otimes_A K$ 和 $M \otimes_A k$ 是维数相同的向量空间 (i.e. $M$ 在 $A$ 上的秩等于 $A$-模 $M$ 的生成元的最小个数).
\end{enumerate}
\end{corollary}

\begin{proof}
证明直接; 推广到 $A$ 只假设无幂零元的情形留给读者：此时必须要求 $M$ 在 $A$ 的诸极小素理想处的秩等于向量空间 $M \otimes_A k$ 的维数.
\end{proof}

\section{平坦性的局部判则}

\begin{proposition}\label{prop:IV.5.1}
设 $A$ 是带理想 $I$ 的环，$M$ 是 $A$-模. 假设
\[
\operatorname{Tor}_{1}^{A}(M, A/I^n) = 0 \quad \text{对 } n > 0,
\]
则典范满同态
\begin{equation}
\mathrm{gr}_{I}^{0}(M) \otimes_{A/I} \mathrm{gr}_{I}(A) \longrightarrow \mathrm{gr}_{I}(M) \tag{($*$)}
\end{equation}
是同构. 逆命题在 $I$ 幂零时成立.
\end{proposition}

\begin{proof}
假设意味着同态
\[
I^n \otimes_A M \longrightarrow I^n M
\]
是同构，由此立即得出同态
\[
I^n/I^{n+1} \otimes_A M \longrightarrow I^n M/I^{n+1} M
\]
是同构. 反过来，假设此条件成立且 $I$ 幂零，证明对每个 $n$ 有 $\operatorname{Tor}_{1}^{A}(M, A/I^n) = 0$. 对大的 $n$ 这成立，对 $n$ 作向下归纳，假设已对 $n+1$ 证明. 有交换图
\[
\begin{tikzcd}[row sep=large, column sep=large]
& M \otimes I^{n+1} \arrow[r] \arrow[d] & M \otimes I^n \arrow[r] \arrow[d] & M \otimes (I^n/I^{n+1}) \arrow[r] \arrow[d] & 0 \\
0 \arrow[r] & MI^{n+1} \arrow[r] & MI^n \arrow[r] & MI^n/MI^{n+1}
\end{tikzcd}
\]
其中两行正合. 由假设，最后一个垂直箭头是同构，归纳假设也意味着第一个垂直箭头是同构. 故中间的垂直箭头亦然，证明完毕.
\end{proof}

下面的命题是讨论班进行期间 Serre 提出的; 它使本节的证明大为简化.

\begin{proposition}\label{prop:IV.5.2}
设 $A \to B$ 是环同态，$M$ 是 $A$-模. 下列条件等价：
\begin{enumerate}[(i)]
\item 对每个 $B$-模 $N$，有 $\mathrm{Tor}_{1}^{A}(M, N) = 0$; 
\item $\mathrm{Tor}_{1}^{A}(M, B) = 0$，且 $M_{(B)} = M \otimes_A B$ 是 $B$-平坦的.
\end{enumerate}
\end{proposition}

\begin{proof}
有函子同构
\[
M \otimes_A N = (M \otimes_A B) \otimes_B N,
\]
它把左端 (看作 $M$ 的函子) 表示为两个函子 $M \mapsto M \otimes_A B$ 和 $P \mapsto P \otimes_B N$ 的复合. 由于前者把 $A$ 上的自由模变为 $B$ 上的自由模，从而把射影变为射影，有复合函子的谱序列
\[
\operatorname{Tor}_{n}^{A}(M, N) \Longleftarrow \operatorname{Tor}_{p}^{B}\bigl(\operatorname{Tor}_{q}^{A}(M, B), N\bigr),
\]
由此得低次项的正合序列
\[
0 \longleftarrow \operatorname{Tor}_{1}^{B}(M \otimes_A B, N) \longleftarrow \operatorname{Tor}_{1}^{A}(M, N) \longleftarrow \operatorname{Tor}_{1}^{A}(M, B) \otimes_A N.
\]
若 (i) 成立，由这个正合序列得出对每个 $N$ 有 $\operatorname{Tor}_{1}^{B}(M \otimes_A B, N) = 0$，i.e. $M \otimes_A B$ 是 $B$-平坦的，由此得 (ii). 若反过来 (ii) 成立，则在正合序列中围住 $\mathrm{Tor}_{1}^{A}(M, N)$ 的项为零，故有 (i).
\end{proof}

\begin{corollary}\label{cor:IV.5.3}
假设 $B = A/I$，则前面的条件等价于：
\begin{enumerate}[(i)]
\item[((iii))] 对每个被 $I$ 的幂零化的 $A$-模 $N$，$\mathrm{Tor}_{1}^{A}(M, N) = 0$.
\end{enumerate}
\end{corollary}

\begin{proof}
事实上，(i) 意味着 $N$ 被 $I$ 零化时如此. 把假设应用于诸 $I^n N/I^{n+1}N$，由此推出 (iii).
\end{proof}

\begin{corollary}\label{cor:IV.5.4}
在 \ref{cor:IV.5.3} 的条件下，所考虑的条件蕴含函子同态
\begin{equation}
\mathrm{gr}_{I}^{0}(M) \otimes_{A/I} \mathrm{gr}_{I}(A) \longrightarrow \mathrm{gr}_{I}(M) \tag{($*$)}
\end{equation}
是同构，且 $M \otimes_A A/I$ 在 $A/I$ 上平坦.
\end{corollary}

\begin{proof}
只需应用 (iii) 和 (\cref{prop:IV.5.1}). 用 \ref{prop:IV.5.1} 在 $I$ 幂零情形的逆命题，得到：
\end{proof}

\begin{corollary}\label{cor:IV.5.5}
设 $A$ 是带幂零理想 $I$ 的环，$M$ 是 $A$-模. 下列条件等价：
\begin{enumerate}[(i)]
\item $M$ 是 $A$-平坦的; 
\item $M \otimes_A A/I$ 是 $A/I$-平坦的，且 $\mathrm{Tor}_{1}^{A}(M, A/I) = 0$; 
\item $M \otimes_A A/I$ 是 $A/I$-平坦的，且相伴分次环上的典范同态 (∗) 是同构.
\end{enumerate}
\end{corollary}

\begin{proof}
事实上，它们分别是前面的条件 (iii) 和 (ii)，以及推论 \ref{cor:IV.5.4} 的条件.
\end{proof}

不再假设 $I$ 幂零，则先验地只有蕴含 (i)$\Rightarrow$(ii)$\Rightarrow$(iii). 另一方面，由于条件 (iii) 除以 $I$ 的幂后仍稳定，由 \ref{cor:IV.5.5} 看到它蕴含
\begin{enumerate}[(i)]
\item[((iv))] 对每个整数 $n$，$M \otimes A/I^n$ 在 $A/I^n$ 上平坦.
\end{enumerate}
要给出由此能推出 (i) (i.e. $M$ 是 $A$-平坦的) 的条件. 我说为此只需 $A$ 诺特且 $M$ 满足下列有限性条件：对 $A$ 上的每个有限型模 $N$，$M \otimes_A N$ 按 $I$-预adic 拓扑分离. (只需对 $A$ 中有限型理想 $N$ 验证.) 事实上，证明在这些条件下，若 $N' \to N$ 是有限型模的单同态，则 $\mathcal{M} \otimes_A \mathcal{N}' \to \mathcal{M} \otimes_A \mathcal{N}$ 是单同态. 只需证明其核含于诸 $\mathcal{I}^n(M \otimes_A N') = \operatorname{Im}(M \otimes_A \mathcal{I}^n N' \to M \otimes_A \mathcal{N}')$，或含于诸 $\operatorname{Im}(M \otimes_A V'_n \to M \otimes_A N') = \mathrm{Ker}(M \otimes_A N' \to M \otimes_A (N'/V'_n))$，其中 $V'_n$ 遍历 $N'$ (配备 $I$-adic 拓扑) 中 0 的可数个基本邻域组. 由 Krull 定理，$N'$ 的 $I$-adic 拓扑由 $N$ 的拓扑诱导，可以取 $V'_n = N' \cap I^n N$. 考虑交换图
\[
\begin{tikzcd}[row sep=large, column sep=large]
M \otimes_A N' \arrow[r] \arrow[d] & M \otimes_A (N'/V'_n) \arrow[d] \\
M \otimes_A N \arrow[r] & M \otimes_A (N/I^n N)
\end{tikzcd}
\]
由于 $N'/V'_n$ 和 $N/I^n N$ 被 $I^n$ 零化，第二个垂直同态等同于由单射 $N'/V'_n \to N/I^n N$ 与平坦 $(A/I^n)$-模 $M \otimes_A A/I^n$ 张量得到的同态，故单射. 因此 $M \otimes_A N' \to M \otimes_A N$ 的核含于 $M \otimes_A N' \to M \otimes_A (N'/V'_n)$ 的核，这正是要证的.

$M$ 所满足的``有限性''条件，在 $M$ 是诺特 $A$-代数 $B$ (满足 $IB$ 含于 $B$ 的根式) 上的有限型模时尤其成立：此时对 $A$ 上的每个有限型模 $N$，$M \otimes_A N$ 是 $B$ 上的有限型模，故由 Krull 定理对 $I$-adic 拓扑 (= 其 $(IB)$-adic 拓扑) 分离. 于是得到：

\begin{theorem}\label{thm:IV.5.6}
设 $A \to B$ 是诺特环的同态，$I$ 是 $A$ 的理想，使得 $IB$ 含于 $B$ 的根式，$M$ 是有限型 $B$-模. 下列条件等价：
\begin{enumerate}[(i)]
\item $M$ 是 $A$-平坦的; 
\item $M \otimes_A A/I$ 是 $A/I$-平坦的，且 $\mathrm{Tor}_{1}^{A}(M, A/I) = 0$; 
\item $M \otimes_A A/I$ 是 $A/I$-平坦的，且典范同态
\[
\mathrm{gr}_{I}^{0}(M) \otimes_{A/I} \mathrm{gr}_{I}(A) \longrightarrow \mathrm{gr}_{I}(M)
\]
是同构; 
\item 对每个整数 $n$，$M \otimes_A A/I^n$ 在 $A/I^n$ 上平坦.
\end{enumerate}
\end{theorem}

\begin{proof}
这个结果在 $A$，$B$ 是诺特局部环、$A \to B$ 是局部同态、$I$ 是 $A$ 中含于其极大理想的理想 (且 \ref{thm:IV.5.6} 可以立即归结到这个情形) 时最常用. 一个有趣的情形是 $A/I$ 是域，i.e. $I$ 极大，此时``$M \otimes_A (A/I)$ 在 $A/I$ 上平坦''的条件成为多余; 此外，由于诸 $A/I^n$ 是 Artin 局部环，条件 (iv) 意味着诸 $M \otimes_A (A/I^n)$ 在诸 $A/I^n$ 上自由.
\end{proof}

\begin{corollary}\label{cor:IV.5.7}
设 $A \to B$ 是诺特局部环的局部同态，$u \colon M' \to M$ 是有限型 $B$-模的同态，假设 $M$ 在 $A$ 上平坦. 则下列条件等价：
\begin{enumerate}[(i)]
\item $u$ 单射，且 Coker $u$ 在 $A$ 上平坦; 
\item $u \otimes_A k \colon M' \otimes_A k \to M \otimes_A k$ 单射
\end{enumerate}
\end{corollary}

(其中 $k$ 表示 $A$ 的剩余域.)

(i)$\Rightarrow$(ii) 由 \ref{prop:IV.1.1}，证明逆命题. 首先 $u$ 单射，因为只需在相伴分次环上验证，在那里由读者将写出的交换方块得出. 设 $M''$ 是其余核，则有正合序列
\[
0 \longrightarrow M' \longrightarrow M \longrightarrow M'' \longrightarrow 0,
\]
由此由 Tor 的正合序列，考虑到假设 (ii) 和 $\mathrm{Tor}_{1}^{A}(M, k) = 0$，有关系 $\mathrm{Tor}_{1}^{A}(M'', k) = 0$，故由定理 \ref{thm:IV.5.6}，$M''$ 在 $A$ 上平坦.

\begin{corollary}\label{cor:IV.5.8}
在 \ref{thm:IV.5.6} 的条件下，设 $J$ 是 $B$ 的含 $IB$ 且含于根式的理想. 设 $\widehat{A}$ 是 $A$ 的 $I$-adic 完备化，$\widehat{B}$ 和 $M$ 是 $B$ 和 $M$ 的 $J$-adic 完备化. 为使 $M$ 是 $A$-平坦的，必要且充分条件是 $\widehat{M}$ 是 $\widehat{A}$-平坦的.
\end{corollary}

\begin{proof}
(N.B. 充分性已经容易由 \ref{cor:IV.3.2} 得出.) 在情形 $(A, B, I, M)$ 和 $(\widehat{A}, \widehat{B}, I\widehat{A}, \widehat{M})$ 中使用 \ref{thm:IV.5.6} 的判则 (iii). 借助 \ref{cor:IV.3.2}，看到对两种情形得到的条件等价.
\end{proof}

\begin{corollary}\label{cor:IV.5.9}
设 $A \to B \to C$ 是诺特局部环的局部同态，$M$ 是有限型 $C$-模 (N.B. $C$ 只是为了能对 $M$ 施加有限性条件). 假设 $B$ 在 $A$ 上平坦. 设 $k$ 是 $A$ 的剩余域. 下列条件等价：
\begin{enumerate}[(i)]
\item $M$ 在 $B$ 上平坦; 
\item $M$ 在 $A$ 上平坦，且 $M \otimes_A k$ 在 $B \otimes_A k$ 上平坦.
\end{enumerate}
\end{corollary}

\begin{proof}
蕴含 (i)$\Rightarrow$(ii) 平凡，证明 (ii)$\Rightarrow$(i). 对 $(B, C, \mathfrak{m}B = I, M)$ 应用判则 (iii) of \ref{thm:IV.5.6}，由于 $M \otimes_B (B/I) = M \otimes_B (B \otimes_A k) = M \otimes_A k$，该判则的第一个条件恰意味着 $M \otimes_A k$ 在 $B \otimes_A k$ 上平坦，好. 判则的第二个条件成立，因为 $M$ 在 $A$ 上平坦且 $B$ 在 $A$ 上平坦，由张量积结合律公式. —— 当然，若以 \ref{cor:IV.5.5} 代替 \ref{thm:IV.5.6}，则在假设 $A$ 的理想 $m$ 幂零时，得到无诺特条件和有限性条件的类似陈述. ($m$ 取极大这一事实其实没有用到; 但某种意义上``$m$ 极大''的情形是``尽可能好''的.)
\end{proof}

\section{平坦态射与开集}

先回忆一些关于可构造集 (constructible sets) 的结果，它们在流传的 Dieudonné-Rosenlicht 概形讨论班笔记中已经证明\footnote{参见 EGA 0$_{\mathrm{III}}$ 9，EGA IV 1.8 和 1.10.}.

设 $X$ 是拓扑空间. 随 Chevalley，称 $X$ 的子集可构造，若它是有限个局部闭子集的并.

\begin{lemma}\label{lem:IV.6.1}
设 $X$ 是诺特拓扑空间，$Z$ 是 $X$ 的子集. 为使 $Z$ 可构造，必要且充分条件是：对 $X$ 的每个不可约闭子集 $Y$，$Z \cap Y$ 在 $Y$ 中不稠密，或包含空间 $Y$ 的一个非空开子集.
\end{lemma}

由此，用交换代数的众所周知的引理，得出：

 \begin{lemma}[Chevalley]\label{lem:IV.6.2}
设 $f \colon X \to Y$ 是预概形的有限型态射，$Y$ 诺特. 则 $f(X)$ 可构造.
\end{lemma}

\begin{lemma}\label{lem:IV.6.3}
设 $X$ 是诺特拓扑空间，其每个不可约闭子集都有生成点，$U$ 是 $X$ 的可构造子集，$x \in X$. 为使 $U$ 是 $x$ 的邻域，必要且充分条件是 $x$ 的每个推广 (generalization) $y$ (i.e. 每个满足 $x \in \overline{\{y\}}$ 的 $y \in X$) 都在 $U$ 中.
\end{lemma}

特别地：

\begin{corollary}\label{cor:IV.6.4}
设 $X$ 是诺特拓扑空间，其每个不可约闭子集都有生成点，$U$ 是 $X$ 的子集. 为使 $U$ 开，必要且充分条件是它满足下列两个条件：
\begin{enumerate}[a)]
\item $U$ 包含其每个点的所有推广; 
\item 若 $x \in U$，则 $U \cap \overline{\{x\}}$ 包含空间 $\overline{\{x\}}$ 的一个非空开子集.
\end{enumerate}
\end{corollary}

\begin{proof}
事实上，由 \ref{lem:IV.6.1}，$U$ 必然可构造，应用证明 $U$ 是其每个点的邻域的判则 \ref{lem:IV.6.3}.
\end{proof}

\begin{corollary}\label{cor:IV.6.5}
设 $f \colon X \to Y$ 是预概形的有限型态射，$Y$ 局部诺特，$x$ 是 $X$ 的点，$y = f(x)$. 为使 $f$ 把 $x$ 的每个邻域变成 $y$ 的邻域，必要且充分条件是：对 $y$ 的每个推广 $y'$，存在 $x$ 的推广 $x'$，使得 $f(x') = y'$.
\end{corollary}

\begin{proof}
显然可以假设 $X$ 和 $Y$ 仿射，从而诺特. 条件充分，因为只需证明 $f(X)$ 是 $y$ 的邻域，而由 \ref{lem:IV.6.1}，$f(X)$ 可构造，只需应用判则 \ref{lem:IV.6.3}. 条件必要，因为设 $Y' = \overline{\{y'\}}$，$F$ 是不含 $x$ 的 $f^{-1}(Y')$ 的不可约分支的并. 则 $X - F$ 是 $x$ 的开邻域，故其像是 $y$ 的邻域，更不用说包含 $y'$，故存在 $x'_1 \in X - F$ 使得 $f(x'_1) = y'$. 考虑 $f^{-1}(Y')$ 的包含 $x'_1$ 的不可约分支，它必然包含 $x$ (否则它含于 $F$)，设 $x'$ 是其生成点. 它是 $x$ 的推广，且 $f(x')$ 是 $f(x'_1) = y'$ 的推广，含于 $Y'$，故等于 $y'$，证毕.
\end{proof}

\begin{theorem}\label{thm:IV.6.6}
设 $f \colon X \to Y$ 是局部有限型态射，$Y$ 局部诺特，$F$ 是 $X$ 上支集为 $X$、相对 $Y$ 平坦的凝聚层. 则 $f$ 是开态射 (i.e. 把开集变为开集).
\end{theorem}

\begin{proof}
只需对每个点 $x \in X$ 证明判则 \ref{cor:IV.6.5}. 而 $x$ 的推广 $x'$ 对应于 $\mathcal{O}_x$ 的素理想，$y$ 的推广 $y'$ 对应于 $\mathcal{O}_y$ 的素理想，因此必须验证 $\mathcal{O}_y$ 的每个素理想都由 $\mathcal{O}_x$ 的素理想诱导. 而 $\mathcal{F}_x$ 是 $\mathcal{O}_x$-模，非零且 $\mathcal{O}_y$ 平坦，故由 \ref{cor:IV.2.2} 在 $\mathcal{O}_y$ 上忠实平坦. 于是可以应用 \ref{prop:IV.2.3}，证明完毕.
\end{proof}

\begin{remarkstar}
由于平坦性对基扩张保持，在 \ref{cor:IV.6.5} 的条件下 $f$ 甚至泛开 (universally open). 不过我并不知道，当 $Y$ 整且 $X$ 在 $Y$ 上有限型时，$f$ 是否在 $X$ 的每个分支 $X_i$ 上诱导开态射，或甚至只是等维 (equidimensional) 态射\index[term]{dengweitaishe@等维态射}\footnote{对第二个问题的答案是肯定的，对第一个问题即使 $f$ 平展也是否定的; 参见 EGA IV 12.1.1.5 和 EGA Err$_{\mathrm{IV}}$ 33.}，i.e. 其所有纤维的分支都有相同维数 (只知道 $X_i$ 支配 $Y$). 问题与下面有关：设 $A \to B$ 是诺特局部环的局部同态，使得 $B$ 在 $A$ 上平坦且 $mB$ 是 $B$ 的定义理想 (这还蕴含 $\dim B = \dim A$). 是否对 $B$ 的每个极小素理想 $\mathfrak{p}_i$ 都有 $\dim B/\mathfrak{p}_i = \dim B$? 只指出：把 \ref{thm:IV.6.6} 的平坦性假设换成``$f$ 泛开''这一假设，第一个问题的答案是肯定的.
\end{remarkstar}

\begin{lemma}\label{lem:IV.6.7}
设 $A$ 是诺特整环，$B$ 是 $A$ 上的有限型代数，$M$ 是有限型 $B$-模. 则存在 $A$ 的非零元素 $f$，使得 $M_f$ 是 $A_f$ 上的自由模 (更不用说平坦).
\end{lemma}

\begin{proof}
设 $K$ 是 $A$ 的分式域，则 $B \otimes_A K$ 是 $K$ 上的有限型代数，$M \otimes_A K$ 是后者的有限型模. 设 $n$ 是该模支集的维数，对 $n$ 归纳. 若 $n < 0$，i.e. $M \otimes_A K = 0$，则取 $M$ 在 $B$ 上的有限个生成元，看到存在零化这些生成元 (从而零化 $M$) 的 $f \in A$，故 $M_f = 0$，已经成功. 假设 $n \geqslant 0$. 知道 $B$-模 $M$ 有合成序列，其逐次商同构于模 $B/\mathfrak{p}_i$，诸 $\mathfrak{p}_i$ 是 $B$ 的素理想. 由于自由模的扩张是自由的，归结到 $M$ 本身形如 $B/\mathfrak{p}$，或等同于 $B$，其中 $B$ 是整 $A$-代数. 把 Noether 正规化引理应用于 $K$-代数 $B \otimes_A K$，容易看到存在 $A$ 的非零元素 $f$，使得 $B_f$ 在子环 $A_f[t_1, \dots, t_n]$ (其中 $t_i$ 是未定元) 上整. 故可以先假设 $B$ 在 $C = A[t_1, \dots, t_n]$ 上整，从而是有限型无挠 $C$-模. 设 $m$ 是其秩，则有 $C$-模的正合序列：
\[
0 \longrightarrow C^m \longrightarrow B \longrightarrow M' \longrightarrow 0,
\]
其中 $M'$ 是挠 $C$-模. 由此，$C \otimes_A K$-模 $M' \otimes_A K$ 的 Krull 维数严格小于 $C \otimes_A K$ 的维数 $n$. 由归纳假设，只要对 $A$ 的适当非零元素 $f$ 局部化，可以假设 $M'$ 是 $A$-自由模. 另一方面，$C^m$ 是 $A$-自由模. 故 $B$ 此时是 $A$-自由模，完成.
\end{proof}

\begin{lemma}\label{lem:IV.6.8}
设 $A$ 是诺特环，$B$ 是 $A$ 上的有限型代数，$M$ 是有限型 $B$-模，$p$ 是 $B$ 的素理想，$q$ 是它在 $A$ 上诱导的素理想. 假设 $M_{\mathfrak{p}}$ 在 $A_{\mathfrak{q}}$ 上平坦 (或在 $A$ 上平坦，都一样). 则存在 $g \in B - \mathfrak{p}$，使得
\begin{enumerate}[a)]
\item $(M/\mathfrak{q}M)_g$ 在 $A/\mathfrak{q}$ 上平坦; 
\item $\mathrm{Tor}_{1}^{A}(M, A/\mathfrak{q})_g = 0$.
\end{enumerate}
\end{lemma}

\begin{proof}
事实上，把 \ref{lem:IV.6.7} 应用于 $\bigl(A/(q), B/\mathfrak{q}B, M/\mathfrak{q}M\bigr)$，先看到存在 $f \in A - \mathfrak{q}$，使得 $(M/\mathfrak{q}M)_f$ 在 $A/\mathfrak{q}$ 上平坦. 另一方面，由于 $M_{\mathfrak{p}}$ 在 $A$ 上平坦，有 $\mathrm{Tor}_{1}^{A}(M, A/\mathfrak{q})_{\mathfrak{p}} = \mathrm{Tor}_{1}^{A}(M_{\mathfrak{p}}, A/\mathfrak{q}) = 0$，故因 $\operatorname{Tor}_{1}^{A}(M, A/\mathfrak{q})$ 是有限型 $B$-模，存在 $g \in B - \mathfrak{p}$ 使得 (b) 成立. 于是可以 (把 $g$ 换成 $gf$) 假设 (a) 同时成立，证明完毕.
\end{proof}

\begin{corollary}\label{cor:IV.6.9}
用 \ref{lem:IV.6.8} 的记号，对 $B$ 的每个含 $p$ 而不含 $g$ 的素理想 $p'$，$M_{\mathfrak{p}'}$ 在 $A$ 上平坦 (或等价地，在 $A_{\mathfrak{q}'}$ 上平坦，其中 $q'$ 是 $p'$ 诱导的 $A$ 的素理想).
\end{corollary}

\begin{proof}
只需对系统 $(A, B_{\mathfrak{q}'}, \mathfrak{q}, M_{\mathfrak{q}'})$ 应用判则 \ref{thm:IV.5.6} (ii)，使用 Tor 的局部化.
\end{proof}

\begin{theorem}\label{thm:IV.6.10}
设 $f \colon X \to Y$ 是有限型态射，$Y$ 局部诺特，$F$ 是 $X$ 上的凝聚层. 设 $U$ 是满足``$\mathcal{F}_x$ 在 $\mathcal{O}_{f(x)}$ 上平坦''的点 $x \in X$ 的集合. 则 $U$ 是开集.
\end{theorem}

\paragraph*{证明.} 可以假设 $X$ 和 $Y$ 仿射，环为 $B$ 和 $A$，故 $\mathcal{F}$ 由有限型 $B$-模 $M$ 定义. 应用判则 \ref{cor:IV.6.4}. 条件 (a) 由 \ref{prop:IV.1.2} (i) 平凡成立，剩下验证 \ref{cor:IV.6.4} 的条件 (b). 这正是引理 \ref{lem:IV.6.8} 和推论 \ref{cor:IV.6.9} 所做的.

在许多问题中，定理 \ref{thm:IV.6.10} 的下列较弱形式就足够了 (它已经由引理 \ref{lem:IV.6.7} 得出，因而不需要可构造集的技术，也不需要定理 \ref{thm:IV.5.6})：

\begin{corollary}\label{cor:IV.6.11}
在 \ref{thm:IV.6.10} 的条件下，若假设 $Y$ 整，则存在 $Y$ 中的非空开集 $V$，使得 $\mathcal{F}$ 在 $f^{-1}(V)$ 的所有点处相对 $Y$ 平坦.
\end{corollary}

\begin{proof}
事实上，开集 $U$ 包含 $Y$ 的生成点的纤维 (因为该点的局部环是域)，故它包含形如 $f^{-1}(V)$ 的开集，因为 $X$ 在 $Y$ 上有限型. 由 \ref{cor:IV.6.11} 还容易得出下列结果，其中 $Y$ 假设诺特 (但不必整)：存在 $Y$ 的局部闭子集组成的划分 $Y_i$，使得 (给 $Y_i$ 配备诱导约化结构) $\mathcal{F}$ 在每个 $X_i = X \times_Y Y_i$ 上诱导相对 $Y_i$ 平坦的层.
\end{proof}
\end{refsection}
\begin{refsection}% =====================================================================
% 第五讲 (Exposé V) 基本群: 一般理论
% (Le groupe fondamental: généralités)
% A. Grothendieck
% 原书页码: pp. 87-115
% =====================================================================
\chapter{基本群：一般理论}

\section*{0. 引言}

本讨论班是 1960 年讨论班的续篇. 我们用诸如 \ref{cor:I.9.7} 这样的记号引用后者，它表示：代数几何讨论班，第 I 讲，第 9.7 号. 1961 年各讲的编号紧随 1960 年的编号. 我们用诸如 (EGA I 8.7.3) 这样的记号引用 Dieudonné-Grothendieck 的《代数几何原理》(Éléments de Géométrie Algébrique).

本讲概括 (并稍加补充) 1960 年最后几个未成文的讲.

与 1961 年一样，一般限于局部诺特预概形，尽管这个限制常常无关紧要. 我们将在讲 VI 中承认忠实平坦下降理论，它概括在 Bourbaki 讨论班 N$^{\circ}$ 190 中. 若有必要，我们将在以后的讲中给出更详细的阐述\footnote{参见讲 VI 和讲 VIII.}，一旦读者有机会相信这一技术对基本群理论的有用性.

\section{带有限算子群的预概形。商预概形}

设 $X$ 是预概形，$G$ 是通过自同构作用在 $X$ 上的有限群 (为确定起见作用在右边). 若 $X$ 是环为 $A$ 的仿射概形，则 $G$ 通过左自同构作用在 $A$ 上.

对每个预概形 $Z$，$G$ 左作用在集合 $\mathrm{Hom}(X, Z)$ 上，于是可以考虑集合
\[
\mathrm{Hom}(X, Z)^{G}
\]
(被 $G$ 不变的态射). 它函子性地依赖于 $Z$，可以问这个函子是否``可表示'' (representable)，i.e. 同构于函子 $Z \mapsto \mathrm{Hom}(Y, Z)$.

这意味着可以找到预概形 $Y$ 和 $G$-不变态射
\[
p \colon X \longrightarrow Y,
\]
使得对每个 $Z$，相应的映射 $g \mapsto gp$
\[
\operatorname{Hom}(Y, Z) \longrightarrow \operatorname{Hom}(X, Z)^{G}
\]
双射. 此时称 $(Y, p)$ 是 $X$ 除以 $G$ 的商预概形 (quotient prescheme) (它在唯一同构意义下确定).

\begin{proposition}\label{prop:V.1.1}
设 $A$ 是有限群 $G$ 在其上左作用的环，$B = A^{G}$ 是 $A$ 的不变元子环，$X = \operatorname{Spec}(A)$，$Y = \operatorname{Spec}(B)$，$p \colon X \to Y$ 是 (显然 $G$-不变的) 典范态射. 则
\begin{enumerate}[(i)]
\item $A$ 在 $B$ 上整，i.e. $p$ 是整态射 (integral morphism)；
\item $p$ 满射，其纤维是 $G$ 的轨道，$Y$ 的拓扑是 $X$ 的拓扑的商；
\item 设 $x \in X$，$y = p(x)$，$G_x$ 是 $x$ 的稳定化子，则 $k(x)$ 是 $k(y)$ 的拟 Galois 代数扩张，且从 $G_x$ 到 $k(x)$ 的 $k(y)$-自同构群 $\operatorname{Gal}\bigl(k(x)/k(y)\bigr)$ 的典范映射满射；
\item $(Y, p)$ 是 $X$ 除以 $G$ 的商预概形.
\end{enumerate}
\end{proposition}

\begin{proof}
陈述 (i)，(ii)，(iii) 在交换代数中众所周知\footnote{参见 N. Bourbaki，Alg. Comm. Chap. 5，§ 1 和 § 2，th. 2.}，只作为备忘，只有关于拓扑的断言例外，它来自下列一般事实 (Cohen-Seidenberg 定理的简单推论)：整态射是闭的 (i.e. 把闭集变为闭集). 立即指出：
\end{proof}

\begin{corollary}\label{cor:V.1.2}
在以上条件下，自然同态
\[
\mathcal{O}_Y \longrightarrow p_{*}(\mathcal{O}_X)^{G}
\]
是同构.
\end{corollary}

\begin{proof}
这由公式
\[
(S^{-1}A)^{G} = S^{-1}(A^{G})
\]
立即得出，该公式对 $B = A^{G}$ 的每个乘法稳定子集 $S$ 成立 (该公式对模也可模块化，更一般地，对平坦的基变换 $A \to A'$ 成立)，把它应用于 $S$ 由 $B$ 的一个元素 $f$ 生成的情形.
\end{proof}

断言 (ii) 和推论 \ref{cor:V.1.2} 容易蕴含 (iv)；更一般地，有：

\begin{proposition}\label{prop:V.1.3}
设 $X$ 是带有限自同构群 $G$ 的预概形，$p \colon X \to Y$ 是 $G$-不变仿射态射，使得 $\mathcal{O}_Y \stackrel{\sim}{\longrightarrow} p_{*}(\mathcal{O}_X)^{G}$. 则 \ref{prop:V.1.1} 的结论 (i)，(ii)，(iii)，(iv) 仍然成立.
\end{proposition}

\begin{proof}
事实上，对 (i)，(ii)，(iii)，可以假设 $Y$ 从而 $X$ 仿射，若 $B$，$A$ 是其环，假设蕴含 $B = A^{G}$，只需应用 \ref{prop:V.1.1}. 对 (iv)，用 (ii) 和 $\mathcal{O}_Y = p_{*}(\mathcal{O}_X)^{G}$.
\end{proof}

\begin{corollary}\label{cor:V.1.4}
在 \ref{prop:V.1.3} 的条件下，对 $Y$ 的每个开集 $U$，$U$ 是 $X|U = p^{-1}(U)$ 除以 $G$ 的商.
\end{corollary}

\begin{proof}
事实上，由 $p$ 诱导的 $p^{-1}(U) \to U$ 满足与 $p$ 相同的假设.
\end{proof}

若现在 $X$ 是 $Z$-预概形且 $G$ 的作用是 $Z$-自同构，则由 (iv)，$Y$ 是 $Z$-预概形. 据此：

\begin{corollary}\label{cor:V.1.5}
为使 $X$ 在 $Z$ 上仿射 (相应地：分离)，必要且充分条件是 $Y$ 如此. 若 $X$ 在 $Z$ 上有限型，则它在 $Y$ 上有限；若还有 $Z$ 局部诺特，则 $Y$ 在 $Z$ 上有限型.
\end{corollary}

\begin{proof}
由于 $X$ 在 $Y$ 上仿射 (更不用说分离)，若 $Y$ 在 $Z$ 上仿射 (相应地：分离)，则 $X$ 亦然. 反过来，假设 $X$ 在 $Z$ 上仿射，证明 $Y$ 仿射：借助 \ref{cor:V.1.4} 可以假设 $Z$ 仿射，归结为证明若 $X$ 仿射则 $Y$ 仿射，这由 \ref{prop:V.1.1} 中 $Y$ 作为 $\operatorname{Spec}(A)^{G}$ 的显式确定得出. 同样，由于 $p \colon X \to Y$ 整故泛闭，且满射，推出若 $X$ 在 $Z$ 上分离则 $Y$ 分离 (一个有待阐明的引理!)；事实上，在交换图
\[
\begin{tikzcd}[row sep=large, column sep=large]
X \times_Z X \arrow[r,"p \times_Z p"] \arrow[d] & Y \times_Z Y \arrow[d] \\
X \arrow[r,"p"] & Y
\end{tikzcd}
\]
中，态射 $X \times_Z X \to Y \times_Z Y$ 闭，故把 $X \times_Z X$ 的 (闭的) 对角线变为 $Y \times_Z Y$ 的闭子集，由于 $p$ 满射它正是后者的对角线. —— 若 $X$ 在 $Z$ 上有限型，则更不用说在 $Y$ 上有限型，故在 $Y$ 上有限 (因为它已在 $Y$ 上整). 再假设 $Z$ 局部诺特，证明 $Y$ 在 $Z$ 上有限型. 借助 \ref{cor:V.1.4} 可以假设 $Z$ 仿射. 由于拓扑空间 $X$ 拟紧且 $p \colon X \to Y$ 满射，$Y$ 也拟紧，故是有限个仿射开集的并，由 \ref{cor:V.1.4} 归结到 $Y$ 仿射从而 $X$ 仿射的情形. 但此时 $X$ 的环 $A$ 是诺特环 $Z$ 的环 $C$ 上的有限型代数，且众所周知 $B = A^{G}$ 也是 $C$ 上的有限型代数 (因为 $A$ 将整，故在 $B$ 的 $C$ 上有限型的子代数 $B'$ 上有限，而 $B'$ 诺特，故 $B$ 也在 $B'$ 上有限，从而在 $C$ 上有限型).
\end{proof}

\begin{corollary}\label{cor:V.1.6}
为使 $X$ 仿射 (相应地：是概形)，必要且充分条件是 $Y$ 如此.
\end{corollary}

\begin{definition}\label{def:V.1.7}
设 $X$ 是有限群 $G$ 在其上右作用的预概形. 称 $G$ 的作用可允许 (admissible)，若存在具有 \ref{prop:V.1.3} 性质的态射 $p \colon X \to Y$ (这意味着 $X/G$ 存在且同构于 $Y$).
\end{definition}

\begin{proposition}\label{prop:V.1.8}
设 $X$ 是有限群 $G$ 在其上右作用的预概形. 为使 $G$ 的作用可允许，必要且充分条件是 $X$ 是 $G$-不变仿射开集的并，或等价地，$G$ 在 $X$ 中的每个轨道都含于一个仿射开集.
\end{proposition}

\begin{proof}
后一个条件显然由前一个蕴含，反过来它也蕴含前一个；因为设 $T$ 是轨道，$U$ 是含它的仿射开集，$U$ 在 $G$ 诸 $g$ 下的变换之交是含 $T$、含于仿射开集 $U$ 的 $G$-稳定开集 $U'$. 由于在 $U$ 中每个有限子集都有仿射开邻域基本组，存在含于 $U'$ 的 $T$ 的仿射开邻域 $V$. 它在 $G$ 诸 $g$ 下的变换是仿射的且含于分离的 $U'$，故其交 $U''$ 是 $G$-不变的仿射开集且含 $T$. —— 据此，\ref{prop:V.1.8} 中所考虑的条件必要，因为取覆盖 $Y$ 的仿射开集 $Y_i$ 的原像 $X_i$ 即可. 它充分，因为此时可以借助 \ref{prop:V.1.1} 构造商 $Y_i = X_i/G$；在每个 $Y_i$ 中，$X_i \cap X_j$ 的像是开集 $Y_{ij}$，由 \ref{cor:V.1.4} 等同于 $X_{ij}/G$，特别由此得到同构 $Y_{ij} \xrightarrow{\sim} Y_{ji}$，使得可以把诸 $Y_i$ 粘合起来构造 $Y$. Serre 宁愿直接构造 $X$ 除以 $G$ 的商拓扑空间 $Y$，在其上放层 $p_{*}(\mathcal{O}_X)^{G}$，并验证 $Y$ 成为预概形，且处于 \ref{prop:V.1.3} 的条件下.
\end{proof}

% 原书编号笔误: 下面这条推论在原文中与上面的定义同为 ``\ref{def:V.1.7}''
\begingroup
\addtocounter{theorem}{-1}%
\renewcommand{\thetheorem}{\thesection.7}%
\begin{corollary}\label{cor:V.1.7sub}
若作用在 $X$ 上的 $G$ 可允许，则 $G$ 的每个子群 $H$ 亦然 (故 $X/H$ 存在).
\end{corollary}

\endgroup

这也可以直接在 \ref{prop:V.1.3} 的情形下验证，注意总可以假设 $X$ 在某个 $Z$ 上仿射，且 $s \in G$ 的元素通过 $Z$-自同构作用 (例如取 $Z = Y$)；事实上：

\begin{corollary}\label{cor:V.1.8}
假设 $X$ 在 $Z$ 上仿射，且 $G$ 的作用是 $Z$-自同构. 则 $G$ 在 $X$ 上的作用可允许. 若 $X$ 由拟凝聚代数层 $\mathcal{A}$ 定义，则 $Y$ 由 $\mathcal{A}$ 的 $G$ 不变元层 $\mathcal{A}^{G}$ 定义.
\end{corollary}

\begin{proposition}\label{prop:V.1.9}
假设 $G$ 在 $X$ 上的作用可允许，且 $X/G = Y$ 是 $Z$ 上的预概形. 考虑基变换态射 $Z' \to Z$，令 $X' = X \times_Z Z'$，$Y' = Y \times_Z Z'$，于是 $G$ 通过结构的转移仍作用在 $X'$ 上，态射 $p' \colon X' \to Y'$ 不变. 若 $Z'$ 在 $Z$ 上平坦，则 $p'$ 仍满足 \ref{prop:V.1.3} 的假设，i.e. $\mathcal{O}_{Y'} \longrightarrow p'_{*}(\mathcal{O}_{X'})^{G}$ 是同构 ($p'$ 无论如何都是仿射的). 故 $G$ 在 $X'$ 上的作用可允许，且
\[
(X/G) \times_Z Z' \approx (X \times_Z Z')/G.
\]
\end{proposition}

\begin{proof}
显然可以假设 $Z = Y$，归结到 $Y$ 和 $Y'$ 仿射的情形. 必须证明：若 $B$ 是 $G$ 在 $A$ 中作用的不变元子环，$B'$ 是在 $B$ 上平坦的 $B$-代数，则 $B'$ 是 $A' = A \otimes_B B'$ 的不变元子代数. 这立即可见，因为正合序列
\[
0 \longrightarrow B \stackrel{i}{\longrightarrow} A \stackrel{j}{\longrightarrow} A^{(G)}
\]
(其中最后一项表示 $A$ 的幂，且 $j(x)$ 是系统 $s \cdot x - x$ ($s \in G$)) 与 $B'$ 张量后仍正合.
\end{proof}

注意平坦性假设对结果的成立是本质的；特别地，若 $Y'$ 是 $Y$ 的闭子预概形 (例如甚至是 $Y$ 的闭点)，$X'$ 是它在 $X$ 中的原像，则 $Y'$ 一般不等同于 $X'/G$. 我们将看到，若 $X$ 在 $Y$ 上平展，则情况仍然如此.

为结束本节，给出一个既方便又平凡的框架. 设 $Y$ 是预概形. 由于在预概形范畴中直和存在，可以对每个集合 $E$ 考虑一族 $(Y_i)_{i \in E}$ (全部等同于 $Y$) 的直和预概形，记为 $Y \times E$. 它由公式
\begin{equation}
\operatorname{Hom}(Y \times E, Z) = \operatorname{Hom}(E, \operatorname{Hom}(Y, Z)) \tag{($*$)}
\end{equation}
刻画，其中第二个 Hom 当然表示从集合 $E$ 到集合 $\mathrm{Hom}(Y, Z)$ 的映射的集合. 有典范态射
\[
Y \times E \longrightarrow Y,
\]
使 $Y \times E$ 成为 $Y$ 上的预概形. 由于纤维积与直和可交换 (在预概形范畴中)，若 $Y$ 是另一个 $Z$ 上的预概形，则对基变换 $Z' \to Z$ 有：
\[
(Y \times E) \times_Z Z' = (Y \times_Z Z') \times E
\]
(该公式在 $Z = Y$ 时尤其有用). 另一方面，由定义平凡地得出
\[
(Y \times E) \times F = Y \times (E \times F) = (Y \times E) \times_Y (Y \times F)
\]
(不过最后一个公式来自上面指出的交换性).

对固定的 $Y$，可以把 $Y \times E$ 看作 $E$ 的函子，取值于 $Y$ 上的预概形，由前面的公式该函子与有限积可交换 (这使每个通常群 $G$ 对应 $Y$ 上的群概形 $Y \times G$，若 $Y$ 有限则它有限，等等). 更一般地，这个函子``左正合''，但我们在此用不到. 它还与直和可交换，并且``右正合''，这从定义公式 (∗) 立即看出. 特别地，若有限群 $G$ 右作用在集合 $E$ 上，则它右作用在 $Y \times E$ 上，且有
\[
(Y \times E)/G = Y \times (E/G),
\]
其中第一个成员的商确实满足 \ref{prop:V.1.9} 的条件 (立即可见).

\section{分解群与惯性群。平展情形}

设 $G$ 是右作用在预概形 $X$ 上的有限群. 若 $x \in X$，称 $x$ 的\emph{分解群} (decomposition group)\index[term]{fenjiequn@分解群} 为 $x$ 的稳定化子 $G_d(x)$. 该群典范地 (左) 作用在剩余域 $k(x)$ 上，其中平凡作用的元素集称为 $x$ 的\emph{惯性群} (inertia group)\index[term]{guanxingqun@惯性群}，记为 $G_i(x)$.

假设 $G$ 在 $X$ 上的作用可允许，$Y$ 是 $Z$ 上的预概形. 固定 $z \in Z$，以及 $k(z)$ 的代数闭扩张 $\Omega$，其超越次数高于诸 $k(x)/k(z)$ 的超越次数，其中 $x$ 是 $z$ 之上的 $X$ 的点. 可以把 $\operatorname{Spec}(\Omega)$ 看作 $Z$-概形，$X$ 在 $\Omega$ 中取值的点对应于 $k(z)$-代数同态 $k(x) \to \Omega$，其中 $x$ 是 $z$ 之上的 $X$ 的点；由于 $\Omega$ 取得足够大，$z$ 之上的每个点 $x$ 都是 $X$ 在 $\Omega$ 中取值的一个点的所在点. 若 $X(\Omega)$ 和 $Y(\Omega)$ 分别表示 $X$ 和 $Y$ 在 $\Omega$ 中取值的点的集合，则有自然映射
\[
X(\Omega) \longrightarrow Y(\Omega),
\]
另一方面，$G$ 作用在 $X(\Omega)$ 上，前面的映射 $G$-不变. 据此，\ref{prop:V.1.1} 的结论 (ii) 和 (iii) 也解释为：前面的映射满射，且把 $Y(\Omega)$ 等同于商 $X(\Omega)/G$. 此外，若 $x$ 是 $a \in X(\Omega)$ 的所在点，则 $a$ 在 $G$ 中的稳定化子正是惯性群 $G_i(x)$. 这一切甚至在不假设 $\Omega$``足够大''时也成立，后一个假设只用于保证可以把 $z$ 之上 $X$ 的每个元素的惯性群刻画为``几何''稳定化子. 例如立即得出：

\begin{proposition}\label{prop:V.2.1}
作基扩张 $Z' \to Z$，由此 $X' = X \times_Z Z'$. 设 $x'$ 是 $X'$ 的点，$x$ 是它在 $X$ 中的像，则有 $G_i(x) = G_i(x')$.
\end{proposition}

\begin{proof}
只需在上述考虑中取 $\Omega$ 为 $k(z')$ 的足够大的扩张 (其中 $z$，$z'$ 是 $x$，$x'$ 在 $Z$，$Z'$ 中的像).
\end{proof}

\begin{proposition}\label{prop:V.2.2}
在 \ref{prop:V.2.2} 的条件下，假设 $Y$ 局部诺特，$X$ 在 $Y$ 上有限. 设 $H$ 是 $G$ 的子群，考虑 $X' = X/H$ (参见 \ref{def:V.1.7})，设 $x \in X$，$x'$ 是它在 $X'$ 中的像，$y$ 是它在 $Y$ 中的像.
\begin{enumerate}[(i)]
\item 若 $H \supset G_d(x)$，则同态 $\mathcal{O}_y \to \mathcal{O}_{x'}$ 在完备化上诱导同构.
\item 若 $H \supset G_i(x)$，则同态 $\mathcal{O}_y \to \mathcal{O}_{x'}$ 平展，i.e. $X'$ 在 $x'$ 处平展于 $Y$ 上.
\end{enumerate}
\end{proposition}

\begin{proof}
设 $Y_1 = \mathrm{Spec}(\widehat{\mathcal{O}_y})$，作基变换 $Y_1 \to Y$，得到在 $Y_1$ 上有限的 $X_1 = X \times_Y Y_1$，$G$ 在其上作用，商由 \ref{prop:V.1.9} 为 $Y_1$. 设 $y_1$ 是 $Y_1$ 在 $y$ 之上的惟一点，由于 $k(y) = k(y_1)$，得出 $X$ 在 $y$ 处的纤维同构于 $X_1$ 在 $y$ 处的纤维，故 $X_1$ 在 $x$ 之上有惟一点 $x_1$. 此外由 \ref{prop:V.1.9} 有 $X_1/H = X'_1 = X' \times_Y Y_1$，设 $x'_1$ 是 $x_1$ 在 $X'_1$ 中的像，它在 $x'$ 之上，且容易验证 ($X'$ 在 $Y$ 上有限型) 同态 $\theta_{x'} \to \theta_{x'_1}$ 在完备化上诱导同构. 故归结到 $Y$ 是完备局部环 $B$ 的谱的情形，从而 $X$ 是 $B$ 上有限环 $A$ 的谱，它由对应于 $Y$ 上 $X$ 的诸点 $x_i$ 的有限个局部环 $A_x$ 组成. 若 $A_0$ 对应于 $x = x_0$，则 $A$ 等同于环 $\mathrm{Hom}_{\mathcal{G}_d}(G, A_0)$，即满足 $f(st) = sf(t)$ ($s \in G_d$) 的函数 $f \colon G \to A_0$ 的环，$G$ 在这些函数上的作用由 $(uf)(t) = f(tu)$ 定义. 于是看到，若 $H$ 是 $G$ 的任意子群，则 $A^{H}$ 是满足
\[
f(stu) = sf(t) \quad \text{对 } s \in G_d, u \in H
\]
的函数 $f \colon G \to A_0$ 的环，故它是半局部环，其局部分量对应于 $G$ 中的双陪集 $G_d a H$，对应于 $a \in G$ 的双陪集 (借助映射 $f \mapsto f(a)$) 的是 $A_0$ 的子环 $A_0^{H(a)}$，其中 $H(a) = G_d \cap aHa^{-1}$. 此外，$A^{H}$ 中对应于 $x$ 的像 $x'$ 的局部分量，也是对应于单位元素的双陪集 $G_d H$ 的那个分量，故其局部环为 $A_0^{G_d \cap H}$. 因此若 $G_d \subset H$，则得到 $A_0^{G_d} = A^{G} = B$，这证明 (i). 为证明 (ii)，可以对 $A$ 作适当的有限平坦扩张，用 \ref{prop:V.2.1}，归结到剩余扩张 $k(x)/k(y)$ 平凡的情形. 但此时 $G_i(x) = G_d(x)$，归结到前一情形.
\end{proof}

\begin{corollary}\label{cor:V.2.3}
在 \ref{prop:V.2.2} 的条件下，假设 $G_i(x) = (e)$，则 $X$ 在 $x$ 处平展于 $Y$ 上. 故若对每个 $x \in X$ 有 $G_i(x) = (e)$，则 $X \to Y$ 是平展态射.
\end{corollary}

有部分逆命题：

\begin{corollary}\label{cor:V.2.4}
假设 $X$ 连通且群 $G$ 在 $X$ 上忠实. 为使 $p \colon X \to Y = X/G$ 平展，必要且充分条件是 $X$ 的诸点的惯性群约化为单位元. 若如此，则 $G$ 等同于 $Y$-概形 $X$ 的所有 $Y$-自同构的群.
\end{corollary}

\begin{proof}
考虑到 \ref{cor:V.2.3}，可以假设 $X$ 在 $Y$ 上平展. 但若 $s \in G$ 属于某个 $G_i(x)$，则由 \ref{cor:I.5.4} 得出 $s$ 在 $G$ 上平凡作用，故是单位元，因为 $G$ 忠实，这证明第一个断言. 设 $u$ 是 $X$ 的 $Y$-自同构，设 $x \in X$. 由命题 \ref{prop:V.1.3}，存在 $s \in G$，使得 $s(x) = u(x)$，且诱导与 $u$ 相同的剩余同态 $k(x)/k(x')$. 由上引文献 (loc. cit.) 有 $s = u$，证毕.
\end{proof}

\begin{remark}\label{rem:V.2.5}
在推论 \ref{cor:V.2.4} 中``$G$ 忠实''的假设当然不是多余的. ``$X$ 连通''的假设也不是，例如取 $X = Y \times E$，$E$ 是有限集，$G$ 是 $E$ 的置换群：$G$ 的作用带强惯性，然而 $(Y \times E)/G = Y \times (E/G) = Y$，且 $X$ 在 $Y$ 上平展. 取 $G$ 为严格小于 $E$ 的对称群、但传递作用在 $E$ 上的群，看到还有不来自 $G$ 的 $X$ 的 $Y$-自同构.
\end{remark}

无惯性作用的典型例子是 $Y \times G$，让 $G$ 借助它在因子 $G$ 上的右平移作用：右算子群为 $G$ 的 $Y$-预概形 $X$ 称为平凡的，若它同构于 $Y \times G$.

为建立带有限算子群的预概形与范畴中主丛 (principal bundle) 概念之间的联系 (这一联系对本讨论班其余部分虽不需要，但在其他语境中很重要)，下面的考虑是有用的. 固定基预概形 $Y$，处于 $Y$-预概形范畴中. 若 $G$ 是有限群，为简短令 $G_Y = Y \times G$，它是 $Y$ 上的有限群概形 (参见第 1 节)，若 $X$ 是 $Y$-预概形，有
\[
X \times_Y G_Y = X \times G
\]
(同一参考). $Y$-态射 $X \times_Y G_Y \to X$ 的给出因此等价于 $Y$-态射 $X \times G \to X$ 的给出，i.e. 对每个 $g \in G$ 给出 $Y$-态射 $T_g \colon X \to X$. 立即看到，为使诸 $T_g$ 在 $X$ 上定义右算子群 $G$ 的预概形结构 (i.e. $T_{gg'} = T_{g'}T_g$，$T_e = \mathrm{id}_x$)，必要且充分条件是相应的 $Y$-态射 $X \times_Y G_Y \to X$ 在 $X$ 上定义 $Y$-群概形 $Y$-算子群的 $Y$-预概形结构 (在范畴 $\mathcal{C}$ 中 $\mathcal{C}$-算子群对象的普遍意义下). 假设如此. 回忆若典范态射
\[
X \times_Y G_Y \longrightarrow X \times_Y X
\]
(其分量分别为 $\mathrm{pr}_1$ 和乘法态射 $\pi \colon X \times_Y G_Y \to X$) 是同构，则称 $X$ 在 $G_Y$ 下形式主齐性\footnote{现在宁愿说：$X$ 是 $G_Y$ 下的伪挠子 (pseudo-torseur sous $G_Y$).}. 在目前情形，把第一个成员等同于 $X \times G$，所考虑的态射是给每个 $g \in G$ 对应态射
\[
(\operatorname{id}_X, T_g) = (\operatorname{id}_X \times_Y T_g)\Delta_{X/Y} \colon X \longrightarrow X \times_Y X
\]
的那个，从而说 $X$ 在 $G_Y$ 下形式主齐性，也意味着 $X \times_Y X$ 同构于对角线在 $G \times G$ 的元素 $(e, g)$ 下的变换的直和 ($G \times G$ 以显然的方式作用在 $X \times_Y X$ 上)，其中 $e$ 是 $G$ 的单位元. 若不想区分左右、给出对多于两个等同于 $X$ 的因子乘积仍适用的公式，可以把条件表述为：典范态射
\[
X \times_G (G \times G) \longrightarrow X \times_Y X
\]
(它由给配对 $(g, g')$ 对应态射
\[
(T_g, T_{g'}) = (T_g \times_Y T_{g'})\Delta_{X/Y} \colon X \longrightarrow X \times_Y X
\]
并使 $G$ 通过对角同态左作用在 $G \times G$ 上：
\[
s(g, g') = (sg, sg')
\]
得到) 是同构.

主齐性空间 (principal homogeneous space) 的概念由形式主齐性空间的概念添加一个公理得出，该公理保证 $X$ 除以 $G_Y$ 的``商''存在，且正是范畴的右幺元对象，这里即 $Y$. 这个公理依语境可以变化，在``下降''的框架中通常最方便地明确为：要求算子对象经一个合适类型的合适基变换后成为``平凡''的，i.e. 同构于乘积 $X \times_Y G_Y$ (此处即 $X \times G$)(如此，实际上允许下降技术；参见 Grothendieck，Technique de descente et théorèmes d'existence en Géométrie Algébrique，Sém. Bourbaki N$^{\circ}$ 190，第 26 至 28 页)\footnote{关于平坦下降理论参见讲 VIII.}. 在这种思路下，这里指出 $G$ 群主齐性丛 (在上引文献 (loc. cit.) 的意义下) 的刻画：

\begin{proposition}\label{prop:V.2.6}
设 $Y$ 是局部诺特预概形，$X$ 是带右作用的有限算子群 $G$ 的 $Y$-预概形. 下列条件等价：
\begin{enumerate}[(i)]
\item $X$ 在 $Y$ 上有限，$Y = X/G$，$X$ 的诸点的惯性群约化为单位元；
\item 存在忠实平坦拟紧的基变换 $Y_1 \to Y$，使得 $X_1 = X \times_Y Y_1$ 是平凡的 $Y_1$-算子预概形，i.e. 同构于 $Y_1 \times G$；
\item[(ii bis)] 同 (ii)，但 $Y_1 \to Y$ 有限、平展、满射；
\item $X$ 在 $G_Y$ 下形式主齐性，且在 $Y$ 上忠实平坦且拟紧.
\end{enumerate}
\end{proposition}

\paragraph*{证明.} (i)$\Rightarrow$(ii bis). 取 $Y_1 = X$，注意 $X \to Y$ 确为有限、平展 (由 \ref{cor:V.2.3}) 且满射. 证明 $X_1$ 于是在 $Y_1$ 上平凡，这由下述推论得出：

\begin{corollary}\label{cor:V.2.7}
若 (i) 成立且 $X$ 在 $Y$ 上有截面，则 $X$ 是平凡算子空间.
\end{corollary}

\begin{proof}
事实上，这个截面可以定义 $G$-态射 $X \times G \to X$，由于 $G$ 在 $X$ 的纤维上传递它满射，由于 $G$ 无惯性作用它单射；最后，由 \ref{cor:I.5.3}，因为 $X$ 在 $Y$ 上平展，它是局部同构. 故它是同构.
\end{proof}

(ii bis) 平凡蕴含 (ii)，(ii) 蕴含 (i)，因为 (i) 的成分对基的忠实平坦拟紧扩张``不变''(关于``有限''参见上面引用的 Bourbaki 讨论班；对惯性群应用 \ref{prop:V.2.1}；对 $Y = X/G$，用 \ref{prop:V.1.9} 在忠实平坦基变换情形的逆命题，我们忘了把它明确化).

我们通过证明 (i)$\Rightarrow$(ii bis) 证明了 (i)$\Rightarrow$(iii). 最后 (iii)$\Rightarrow$(ii)，因为 (iii) 中的第一个假设恰意味着 $X$ 在作基变换 $Y_1 = X$ 后平凡；由此 (ii)，因为 $X$ 在 $Y$ 上忠实平坦且拟紧.

\begin{definition}\label{def:V.2.8}
满足 \ref{prop:V.2.6} 等价条件的带右算子群 $G$ 的 $Y$-预概形 $X$，称为 $Y$ 的以 $G$ 为 伽罗瓦群的主覆盖 (principal covering).
\end{definition}

\section{平展覆盖的自同构与态射}

\begin{proposition}\label{prop:V.3.1}
设 $X$ 在局部诺特 $Y$ 上平展、分离、有限型，$G$ 是通过 $Y$-自同构作用在 $X$ 上的有限群. 则 $G$ 的作用可允许，且商预概形 $X/G$ 在 $Y$ 上平展.
\end{proposition}

\begin{proof}
不假设 $X$ 在 $Y$ 上有限，但 $X$ 在 $Y$ 上拟射影，从而由 \ref{cor:V.1.8} 存在 $X/G$. 先证明：
\end{proof}

\begin{corollary}\label{cor:V.3.2}
态射 $X \to X/G$ 平展.
\end{corollary}

\begin{proof}
显然可以先假设 $G$ 在 $X$ 的连通分支集合上传递，然后考虑一个连通分支的稳定化子，甚至可以假设 $X$ 连通. 最后可以假设 $G$ 忠实作用. 但此时像 \ref{cor:V.2.4} 中那样看到 $G$ 无惯性作用，故由 \ref{cor:V.2.3} 得出 $X \to X/G$ 平展. 借助下述引理得出结论：
\end{proof}

 \begin{lemma}[对讲 I 的更正]\label{lem:V.3.3}
设 $X \to X' \to Y$ 是有限型态射，$x$ 是 $X$ 的点，$x'$ 和 $y$ 是它的像. 假设 $Y$ 局部诺特. 若所考虑的三个态射中有两个在标记点处平展，则第三个亦然.
\end{lemma}

\begin{proof}
只剩下 $X \to X'$ 和 $X \to Y$ 在 $x$ 处平展、要证明 $X' \to Y$ 在 $x'$ 处平展的情形 (这正是 \ref{prop:V.3.1} 所需). 作 $Y$ 的适当平坦扩张，归结到剩余扩张 $k(x)/k(y)$ 平凡的情形. 考虑同态 $\mathcal{O}_y \to \mathcal{O}_{x'} \to \mathcal{O}_x$ 及由完备化得到的同态，假设意味着 $\widehat{\theta_y} \to \widehat{\theta_x}$ 和 $\widehat{\mathcal{O}_{x'}} \to \widehat{\mathcal{O}_x}$ 是同构，由此立即得出 $\widehat{\mathcal{O}_y} \to \widehat{\mathcal{O}_{x'}}$ 是同构，这证明引理.
\end{proof}

\begin{corollary}\label{cor:V.3.4}
若 $X$ 在 $Y$ 上有限且平展，则 $X/G$ 在 $Y$ 上有限且平展.
\end{corollary}

\begin{proposition}\label{prop:V.3.5}
设 $X$，$X'$ 是 $Y$ 的两个平展覆盖. 则每个 $Y$-态射 $f \colon X \to X'$ 分解为平展满射态射 $X \to X''$ 与 $X'$ 的一个既开又闭子集 $X''$ 的典范浸入 $X'' \to X'$ 的乘积.
\end{proposition}

\begin{proof}
知道 (\ref{cor:I.4.8}) $f$ 平展，故开；另一方面 $X$ 在 $Y$ 上有限，故 $f$ 闭，因此 $f(X) = X''$ 是 $X'$ 的既开又闭子集. 证明完毕 (N.B. 只需 $X'$ 在 $Y$ 上非分歧，而不必是平展覆盖).
\end{proof}

\begin{corollary}\label{cor:V.3.6}
用前面的记号，$X \to X''$ 是预概形范畴中的严格满态射 (strict epimorphism)，$X'' \to X'$ 是预概形范畴中的单态射 (monomorphism) (甚至是严格单态射).
\end{corollary}

\begin{proof}
第一个断言按定义意味着态射序列
\[
X \times_{X''} X \xrightarrow[\mathrm{pr}_2]{\mathrm{pr}_1} X \longrightarrow X''
\]
正合，这由 $X \to X''$ 有限且忠实平坦的事实得出 (容易看出，参见 Grothendieck，loc. cit.). 对 $X'' \to X'$ 的对偶断言更平凡.
\end{proof}

推论 \ref{cor:V.3.6} 在下一节的基本群理论中对我们有用；(对那些不喜欢严格满态射概念的人) 可以把推论 \ref{cor:V.3.6} 换成读者按个人口味安排的变体. 只借此机会指出：分解 $f = f'f''$ (其中 $f''$ 是严格满态射，$f'$ 是单态射) 在 (任何范畴中) 唯一同构意义下必然唯一；然而，可以同时存在具有对偶性质的分解 $f = f_1 f_2$ ($f_2$ 是满态射，$f_1$ 是严格单态射，也在唯一同构意义下唯一)，它不同构于前一个：只需取例如拓扑向量空间 (若愿意，分离的) 的范畴，及使 $u(X)$ 不闭的 $u \colon X \to X'$.

\begin{proposition}\label{prop:V.3.7}
设 $Y$ 是连通局部诺特预概形，$y$ 是 $Y$ 的点，$\Omega$ 是 $k(y)$ 的代数闭扩张. 对 $Y$ 上的每个 $X$，用 $X(\Omega)$ 表示 $X$ 在 $\Omega$ 中取值的点的集合. 设 $X$，$X'$ 是 $Y$ 的平展覆盖，$u \colon X \to X'$ 是 $Y$-态射，使得相应的映射 $X(\Omega) \to X'(\Omega)$ 是同构. 则 $u$ 是同构.
\end{proposition}

\begin{proof}
立即归结到 $X'$ 连通的情形. 由于 $X \to X'$ 有限且平展，知道 $X \to X'$ 的纤维中点的几何个数恒定，且等于 1 当且仅当所考虑的态射是同构. 而这个数也是 $X(\Omega) \to X'(\Omega)$ 的一个纤维中的元素个数，由此得结论.
\end{proof}

\section{Galois 理论的公理条件}

设 $\mathcal{C}$ 是范畴，$F$ 是 $\mathcal{C}$ 到有限集范畴的共变函子. 假设下列条件满足：
\begin{itemize}\tightlist
\item[(G 1)] $\mathcal{C}$ 有终对象 (final object)\footnote{回忆 $\mathcal{C}$ 的对象 $e$ 称为终对象，若对 $\mathcal{C}$ 中的每个 $X$，$\mathrm{Hom}(X, e)$ 恰好有一个元素. 对偶地定义 $\mathcal{C}$ 的始对象 (initial object).}，且 $\mathcal{C}$ 中两个对象在第三个之上的纤维积存在 (这个公理也可以说成：$\mathcal{C}$ 中存在有限投射极限).
\item[(G 2)] $\mathcal{C}$ 中存在有限直和 (故也存在扮演空集角色的始对象 $\varnothing_{\mathcal{C}}$)，以及对象除以有限自同构群的商.
\item[(G 3)] 设 $u \colon X \to Y$ 是 $\mathcal{C}$ 中的态射，则 $u$ 分解为乘积 $X \xrightarrow{u'} Y' \xrightarrow{u''} Y$，其中 $u'$ 是严格满态射，$u''$ 是单态射，且是 $Y$ 的一个直因子上的同构.
\item[(G 4)] 函子 $F$ 左正合 (i.e. 把右幺元变为右幺元，且与纤维积可交换).
\item[(G 5)] $F$ 与有限直和可交换，把严格满态射变为满态射，且与除以有限自同构群的商可交换.
\item[(G 6)] 设 $u \colon X \to Y$ 是 $\mathcal{C}$ 中的态射，使得 $F(u)$ 是同构，则 $u$ 是同构.
\end{itemize}

我们的目标是从条件 (G 1) 至 (G 6) 构造一个投射有限 (projective limit of finite groups) 的拓扑群 $\pi$，以及 $\mathcal{C}$ 与有限集范畴 $\mathcal{C}(\pi)$\index[notations]{$\mathcal{C}(\pi)$} (其中 $\pi$ 连续作用，i.e. 作用的稳定化子是开子群，或等价地，已作用在所考虑集合上的一个离散商群存在) 的等价，所构造的等价把给定函子 $F$ 变为 $\mathcal{C}(\pi)$ 到有限集范畴的显然包含函子. 立即注意，用拓扑群 $\pi$ 构造的范畴 $\mathcal{C}(\pi)$ 和前面的包含函子确实满足条件 (G 1) 至 (G 6).

分几步进行.

\begin{enumerate}[a)]
\item 设 $u \colon X \to Y$ 在 $\mathcal{C}$ 中. 为使 $u$ 是单态射，必要且充分条件是 $F(u)$ 单射. (用 (G 1)，(G 4)，(G 6).)
\end{enumerate}

事实上，说 $u$ 是单态射意味着投影 $\mathrm{pr}_1 \colon X \times_Y X \to X$ 是同构.

\begin{enumerate}[a)]
\setcounter{enumi}{1}
\item 对象 $X$ 都是 $\mathcal{C}$ 的 Artin 对象.
\end{enumerate}

事实上，若 $X' \to X'' \to X$ 是单态射，使得 $F(X')$ 和 $F(X'')$ 在 $F(X)$ 中有相同的像，则由 a) $F(X') \to F(X'')$ 是同构，故由 (G 6)，$X' \to X''$ 是同构.

\begin{enumerate}[a)]
\setcounter{enumi}{2}
\item 函子 $F$ 严格 pro-可表示 (strictly pro-representable). (参见 Grothendieck，Technique de descente et théorèmes d'existence en Géométrie Algébrique，II，Séminaire Bourbaki 195，février 1960.)
\end{enumerate}

事实上，由上引文献 (loc. cit.) 命题 \ref{prop:V.3.1}，这由 b) 和 (G 4) 得出. 故可以找到 $I$ 上的有序滤过投射系：
\[
P = (P_i)_{i \in I}
\]
(看作 $\mathcal{C}$ 的 pro-对象)，以及函子同构
\begin{equation}
F(X) = \operatorname{Hom}_{\operatorname{Pro}(\mathcal{C})}(P, X) = \varinjlim_{i} \operatorname{Hom}_{\mathcal{C}}(P_i, X). \tag{($*$)}
\end{equation}
该同构由元素
\[
\varphi \in \varprojlim_{i} F(P_i) = F(P)
\]
实现. 还可以假设转移同态 $\varphi_{ji} \colon P_i \to P_j$ ($i \geqslant j$) 是满态射，且每个满态射 $P_i \to P'$ 等价于对适当的 $j \leqslant i$ 的满态射 $P_i \to P_j$ (这基本上唯一地确定投射系 $P$).

称 $X \in \mathcal{C}$ 连通 (connected)，若它不同构于两个不都同构于始对象 $\varnothing_{\mathcal{C}}$ 的对象的和.

\begin{enumerate}[a)]
\setcounter{enumi}{3}
\item 诸 $P_i$ 连通且不同构于 $\varnothing_{\mathcal{C}}$.
\end{enumerate}

若 $X$ 是左幺元，由应用于空族的直和的 (G 5)，$F(X) = \varnothing$，反之亦然 (由 (G 6)). 故若 $X'$ 是 $\mathcal{C}$ 中非左幺元的对象，i.e. $F(X') \neq \varnothing$，则不存在从 $X'$ 到 $X$ 的态射. 故若某个 $P_i$ 是左幺元，则 $i$ 是有序滤过指标集 $I$ 的最大元素，且公式 (∗) 将意味着 $F(X) = \mathrm{Hom}(P_i, X) =$ 对每个 $X$ 约化为一个元素的集合，这荒谬，因为 $F(\varnothing_{\mathcal{C}}) = \varnothing$. 故诸 $P_i$ 不同构于 $\varnothing_{\mathcal{C}}$.

假设 $P_i = A \amalg B$，由此由 (G 5) $F(P_i) = F(A) \amalg F(B)$，特别是 $F(P_i)$ 中对应于 (由 (∗)) 恒等同态 $P_i \to P_i$ 的元素 $a_i$ 在 $F(A) \amalg F(B)$ 中，例如在 $F(A)$ 中. 这意味着存在 $j \geqslant i$，使得 $\varphi_{ij} \colon P_j \to P_i$ 分解为 $P_j \to A \to P_i = A \amalg B$，其中第二个箭头是典范态射. 故 $F(P_j) \to F(P_i)$ 分解为 $F(P_j) \to F(A) \to F(P_i) = F(A) \amalg F(B)$，且由于由 (G 5)，$F(P_j) \longrightarrow F(P_i)$ 满射，得出 $F(B) = \varnothing$，故 $B$ 同构于 $\varnothing_{\mathcal{C}}$.

\begin{enumerate}[a)]
\setcounter{enumi}{4}
\item 每个态射 $u \colon X \to Y$ (其中 $X$ 不同构于 $\varnothing_{\mathcal{C}}$，$Y$ 连通) 都是严格满态射. 连通对象的每个自同态都是自同构.
\end{enumerate}

考虑 $u$ 的 (G 3) 分解，由于 $X \neq \varnothing_{\mathcal{C}}$，由 (G 6) 得出 $F(X) \neq \varnothing$，故 $F(Y') \neq \varnothing$，从而 $Y' \neq \varnothing_{\mathcal{C}}$，故 $Y$ 连通时 $Y'$ 等同于 $Y$，因此 $u$ 是严格满态射. 假设 $u$ 是连通对象 $X$ 的自同态，证明它是自同构. 事实上，可以假设 $X$ 不同构于 $\varnothing_{\mathcal{C}}$，故由前面所述 $u$ 是严格满态射，因此由 (G 5) $F(u)$ 是满态射，又因 $F(X)$ 是有限集，$F(u)$ 双射. 故由 (G 6)，$u$ 是自同构.

特别地，每个 $P_i$ 的自同态都是自同构.

\begin{enumerate}[a)]
\setcounter{enumi}{5}
\item 对 $P_i$ 的下列条件等价：(i) 自然单射 $\mathrm{Hom}(P_i, P_i) \to \mathrm{Hom}(P, P_i) \simeq F(P_i)$ 也满射，i.e. 对每个 $u \colon P \to P_i$，存在 $v \colon P_i \to P_i$ 使得 $u = v\varphi_i$ (其中 $\varphi_i$ 是典范同态 $P \to P_i$). (ii) 群 $\operatorname{Aut}(P_i)$ 传递作用在 $F(P_i)$ 上. (iii) 群 $\operatorname{Aut}(P_i)$ 简单传递作用在 $F(P_i)$ 上.
\end{enumerate}

事实上，把 $\mathrm{Hom}(P, P_i)$ 等同于 $F(P_i)$，(i) 中所考虑的映射正是 $v \mapsto F(v)(\varphi_i)$. 三个条件的等价性来自 $\mathrm{Hom}(P_i, P_i) = \operatorname{Aut}(P_i)$ 且前面的映射已经单射这一事实.

满足 f) 的等价条件 (i)，(ii)，(iii) 的 $P_i$ 称为 Galois 的 (Galois).

\begin{enumerate}[a)]
\setcounter{enumi}{6}
\item 对 $\mathcal{C}$ 中的每个 $X$，存在 Galois 的 $P_i$，使得每个 $u \in \mathrm{Hom}(P, X)$ 分解为 $P \xrightarrow{\varphi_i} P_i \longrightarrow X$.
\end{enumerate}

设 $J = \operatorname{Hom}(P, X) = F(X)$，它是有限集，故存在 $P_j$，使得每个 $u \colon P \to X$ 分解为 $P \xrightarrow{\varphi_j} P_j \longrightarrow X$，或等价地，使得自然态射
\[
P \longrightarrow X^{J} \quad (J = \operatorname{Hom}(P, X))
\]
分解为
\[
P \xrightarrow{\varphi_j} P_j \longrightarrow X^{J}.
\]
由 (G 3)，态射 $P_j \to X^{J}$ 分解为单态射与严格满态射的乘积，可以取为形如 $\varphi_{ij} \colon P_j \to P_i$. 故归结为证明 $P_i$ Galois. 设 $k \geqslant j$ 是使得每个态射 $P \to P_i$ 都经 $P \stackrel{\varphi_k}{\longrightarrow} P_k \longrightarrow P_i$ 分解的指标. 注意自然态射 $P_k \to X^{J}$ 仍分解为复合
\[
P_k \xrightarrow{\varphi_{ik}} P_i \xrightarrow{U} X^{J},
\]
其中第一个箭头由 e) 是严格满态射，第二个是单态射. 要证明：对给定的态射 $\psi \colon P_k \to P_i$，存在 $P_i$ 的自同态 $v$ 使得 $\psi = v\varphi_{ik}$. 但对每个 $u \in \mathrm{Hom}(P_i, X)$，考虑 $u\psi \in \mathrm{Hom}(P_k, X)$，它形如 $u'\varphi_{ik}$，其中 $u' \in \mathrm{Hom}(P_i, X)$ 完全确定. 由 $\psi$ 如此定义的 $J$ 到 $J$ 的映射 $u \mapsto u'$ 还单射，因为由 e) $\psi$ 是满态射；由于 $J$ 有限它双射. 故 $J$ 到 $J$ 的双射映射 $u \mapsto u'$ 定义同构 $\alpha \colon X^{J} \xrightarrow{\sim} X^{J}$，使下面的图交换. 由态射分解为单态射与严格满态射的乘积的唯一性，得出 (由于 $\psi$ 由 e) 也是严格满态射) 可以找到使图交换的态射 $v \colon P_i \to P_i$，证毕.

特别得出，Galois 的 $P_i$ 在诸 $(P_j)$ 中构成共尾系. 由于对 Galois 对象 $P_i$ 有
\[
\operatorname{Hom}(P, P_i) = \operatorname{Hom}(P_i, P_i) = \operatorname{Aut}(P_i),
\]
取极限得：
\[
\operatorname{Hom}(P, P) = \varprojlim_{i} \operatorname{Hom}(P, P_i) = \varprojlim_{i} \operatorname{Hom}(P_i, P_i) = \varprojlim_{i} \operatorname{Aut}(P_i),
\]
其中投射极限在 Galois 的 $P_i$ 上取. 此外，借助等同 $\mathrm{Hom}(P, P_i) = F(P_i)$，并考虑到 $F$ 把满态射变为满态射，看到前面投射系的转移同态满射. 由这一切得出结论：

\begin{enumerate}[a)]
\setcounter{enumi}{7}
\item 有
\[
\operatorname{Hom}(P, P) = \operatorname{Aut}(P) = \varprojlim_{i} F(P_i) = \varprojlim_{i} \operatorname{Aut}(P_i),
\]
其中投射极限在 Galois 的 $P_i$ 上取.
\end{enumerate}

特别地，$\operatorname{Aut}(P)$ 表现为有限群投射系 (转移同态满射) 的投射极限，给它配备离散拓扑的投射极限拓扑. 记 $\pi$ 并称为\index[notations]{$\pi$} ($\mathcal{C}$ 配备 $F$ 的) \emph{基本群} (fundamental group)\index[term]{jibenqun@基本群} 的是 $\operatorname{Aut}(P)$ 的对偶群 (opposite group). 该群因此右作用在 $P$ 上，它是右作用在 Galois 的 $P_i$ 上的有限群 $\pi_i$ 的投射极限，$\pi_i$ 是 $\operatorname{Aut}(P_i)$ 的对偶群.

考虑到函子同构
\[
F(X) = \operatorname{Hom}(P, X)
\]
和 $\pi$ 的定义，看到 $\pi$ 左作用在 $F(X)$ 上，且由 g) 连续作用 (因为用 g) 的记号，实际上是 $\pi_i$ 作用在 $F(X)$ 上). 对 $\mathcal{C}$ 中的每个态射 $u \colon X \to Y$，态射 $F(u) \colon F(X) \to F(Y)$ 与 $\pi$ 的作用相容，这是平凡的. 因此可以把 $F$ 看作共变函子
\[
F \colon \mathcal{C} \longrightarrow \mathcal{C}(\pi),
\]
其中 $\mathcal{C}(\pi)$ 是 $\pi$ 在其上连续左作用的有限集的范畴.

现在定义反方向的函子：
\[
G \colon \mathcal{C} \longleftarrow \mathcal{C}(\pi)
\]
由公式
\[
G(E) = P \times_\pi E,
\]
其中 $P \times_\pi E$ 定义为由下式概括的万有问题的解
\[
\operatorname{Hom}_{\mathcal{C}}(P \times_\pi E, X) \xrightarrow{\sim} \operatorname{Hom}_\pi(E, \operatorname{Hom}(P, X))
\]
(其中第二个成员中的 $\mathrm{Hom}(P, X) = F(X)$ 看作 $\pi$ 在其上左作用的集合). 必须证明对象 $P \times_\pi E$ 的存在性.

\begin{enumerate}[a)]
\setcounter{enumi}{8}
\item 设 $Q$ 是有限群 $G$ 在其上右作用的 $\mathcal{C}$ 的对象，$E$ 是 $G$ 在其上左作用的有限集. 则 $Q \times_G E$ 存在，且典范映射
\[
F(Q) \times_G E \longrightarrow F(Q \times_G E)
\]
是同构.
\end{enumerate}

由于由 (G 2) $\mathcal{C}$ 中存在有限直和，且 $F$ 由 (G 5) 与其可交换，立即归结到 $G$ 传递作用在 $E$ 上的情形，因为若 $E_j$ 是 $G$ 在 $E$ 中的轨道，则有
\[
Q \times_G E = \coprod_{j} Q \times_G E_j.
\]
设 $a \in E$，$H$ 是其稳定化子，由定义立即看到 $Q \times_G E$ 等同于 $Q/H$. 由此得到存在性 (借助 (G 2))，以及对 $F$ 的交换性 (借助 (G 5)).

\begin{enumerate}[a)]
\setcounter{enumi}{9}
\item 设 $E$ 是 $\mathcal{C}(\pi)$ 的对象，$P_i$ 是 $\pi_i$ 已作用在 $E$ 上的 Galois 对象. 则 $P_i \times_{\pi_i} E$ 存在，且有典范同构
\[
E \stackrel{\sim}{\longrightarrow} F(P_i \times_{\pi_i} E).
\]
若 $j \geqslant i$ 使得 $P_j$ Galois，则典范同态 $P_j \times_{\pi_j} E \to P_i \times_{\pi_i} E$ 是同构.
\end{enumerate}

第一个断言是 i) 的特殊情形，考虑到 $\pi_i$ 简单传递作用在配备标记点 $\varphi_i$ 的 $F(P_i)$ 上，由此得同构 $F(P_i) \times_{\pi_i} E \simeq E$. 对第二个断言例如用 (G 6).

对每个 $j$，设 $\mathcal{C}_j$ 是 $\mathcal{C}$ 的由满足``$\mathrm{Hom}(P_j, X) \to \mathrm{Hom}(P, X) \simeq F(X)$ 双射''的 $X$ 组成的全子范畴. 由 g) 知道 $\mathcal{C}$ 是诸 $\mathcal{C}_j$ 的滤过并. 故对 $X \in \mathcal{C}_j$ 有
\[
\begin{array}{r}
\mathrm{Hom}_\pi(E, \mathrm{Hom}(P, X)) \simeq \mathrm{Hom}_\pi(E, \mathrm{Hom}(P_j, X)) \simeq \mathrm{Hom}_{\pi_j}(E, \mathrm{Hom}(P_j, X)) \\
\simeq \mathrm{Hom}(P_j \times_{\pi_j} E, X),
\end{array}
\]
考虑到 j) 中最后一个断言，得到对 $\mathcal{C}_j$ 的对象 $X$ 函子性的同构：
\[
\mathrm{Hom}_\pi(E, \mathrm{Hom}(P, X)) \simeq \mathrm{Hom}(P_i \times_{\pi_i} E, X).
\]
由于这对每个 $j$ 成立，且这些函子同构随 $j$ 变动而互相诱导，得出结论：

\begin{enumerate}[a)]
\setcounter{enumi}{10}
\item 在 j) 的条件下，由典范态射复合而成的态射
\[
E \longrightarrow \operatorname{Hom}(P_i, P_i \times_{\pi_i} E) \longrightarrow \operatorname{Hom}(P, P_i \times_{\pi_i} E)
\]
使 $P_i \times_{\pi_i} E$ 成为定义 $P \times_\pi E$ 的万有问题的解，i.e. 后者存在且有同构
\[
P \times_\pi E \stackrel{\sim}{\longrightarrow} P_i \times_{\pi_i} E.
\]
\end{enumerate}

这完成函子 $G(E)$ 的构造. 另一方面有函子同态
\[
\alpha \colon \mathrm{id}_{\mathcal{C}(\pi)} \longrightarrow FG,
\]
i.e. 对 $\mathcal{C}(\pi)$ 的对象 $E$ 函子性的同态：
\[
\alpha(E) \colon E \longrightarrow FG(E) = F(P \times_\pi E),
\]
即典范态射的复合
\[
E \longrightarrow F(P) \times_\pi E \longrightarrow F(P \times_\pi E)
\]
(其中第一个来自标记点 $\varphi \in F(P)$). 结合 j) 和 k)，得到：

\begin{enumerate}[a)]
\setcounter{enumi}{11}
\item 同态 $\alpha$ 是同构.
\end{enumerate}

类似地定义函子同态
\[
\beta \colon GF \longrightarrow \mathrm{id}_{\mathcal{C}},
\]
i.e. 对 $\mathcal{C}$ 的对象 $X$ 函子性的同态：
\[
\beta(X) \colon P \times_\pi F(X) \longrightarrow X,
\]
作为与 $\pi$-同态
\[
F(X) \longrightarrow \operatorname{Hom}(P, X)
\]
(它是典范同构 $\mathrm{Hom}(P, X) \stackrel{\sim}{\longrightarrow} F(X)$ 的逆) 相伴的同态.

\begin{enumerate}[a)]
\setcounter{enumi}{12}
\item 复合
\[
\begin{array}{c}
F(X) \xrightarrow{\alpha(F(X))} FGF(X) \xrightarrow{F(\beta(X))} F(X), \\
G(E) \xrightarrow{G(\alpha(E))} GFG(E) \xrightarrow{\beta(G(E))} G(E)
\end{array}
\]
是恒等同构.
\end{enumerate}

(此处原书有一句无法辨认的旁注.)

考虑到 l) 得出：

\begin{enumerate}[a)]
\setcounter{enumi}{13}
\item 同态 $\beta$ 是同构.
\end{enumerate}

这样得到所承诺的结果：

\begin{theorem}\label{thm:V.4.1}
设 $\mathcal{C}$ 是满足本节开头条件 (G 1)，(G 2)，(G 3) 的范畴，$F$ 是 $\mathcal{C}$ 到有限集范畴的共变函子，满足条件 (G 4)，(G 5)，(G 6). 则前面的典范构造定义互为拟逆的\emph{范畴等价} $F \colon \mathcal{C} \to \mathcal{C}(\pi)$ 和 $G \colon \mathcal{C}(\pi) \to \mathcal{C}$. 精确地说，存在 $\mathcal{C}$ 的 pro-对象 $P$ 和函子同构 $F(X) \longleftarrow \mathrm{Hom}(P, X)$，$\pi$ 是 $P$ 的自同构群的对偶群，以合适的方式拓扑化，使得 $\pi$ 连续作用在集合 $\mathrm{Hom}(P, X) \simeq F(X)$ 上. 最后，$G$ 由 $G(E) \simeq P \times_\pi E$ 给出.
\end{theorem}

\begin{remark}\label{rem:V.4.2}
若把 (G 2) 和 (G 5) 分别换成：
\begin{itemize}\tightlist
\item[(G' 2)] $\mathcal{C}$ 中存在有限归纳极限；
\item[(G' 5)] 函子 $F$ 右正合 (i.e. 与有限归纳极限可交换)，
\end{itemize}
则条件 (G 1) 至 (G 6) 的陈述变得更简单更友好. 这些条件表面上强于 (G 2) 和 (G 5)，但由结构定理 \ref{thm:V.4.1} 立即得出它们被 (G 1) 至 (G 6) 蕴含. 然而注意，在我们关心的情形，(G 2) 和 (G 5) 的验证似乎确实比 (G' 2) 和 (G' 5) 简单. 我不知道在条件 (G 3) 中，``$u''$ 是 $Y$ 的一个直因子上的同构''这一事实是否可以被省略.
\end{remark}

\section{Galois 范畴}

\begin{definition}\label{def:V.5.1}
称范畴 $\mathcal{C}$ 为 \emph{Galois 范畴} (Galois category)\index[term]{Galois fanchou@Galois 范畴}，若它等价于某个范畴 $\mathcal{C}(\pi)$，其中 $\pi$ 是紧致群，有限群的投射极限 (i.e. 完全不连通).
\end{definition}

$\mathcal{C}(\pi)$ 的定义参见第 4 节开头. 由定理 \ref{thm:V.4.1}，$\mathcal{C}$ 是 Galois 范畴当且仅当它满足条件 (G 1) 至 (G 3)，且存在 $\mathcal{C}$ 到有限集范畴的、满足条件 (G 4) 至 (G 6) 的函子 $F$ (i.e. 在一般术语中，正合且保守的函子). 这样的函子称为 Galois 范畴 $\mathcal{C}$ 的\emph{基本函子} (fundamental functor)\index[term]{jibenhanzi@基本函子}\footnote{看来最好采用更形象的术语``\emph{纤维函子}'' (fiber functor).}；它被我们记为 $P_F$ 的 pro-对象 pro-表示；使得相伴函子 $F$ 基本 (fundamental) 的 pro-对象 $P$ 称为\emph{基本 pro-对象} (fundamental pro-object)\index[term]{jiben pro-duixiang@基本 pro-对象}. 这样，$\mathcal{C}$ 上的基本函子范畴与基本 pro-对象范畴反等价；若 $F$ 和 $P$ 对应，则群 $\mathrm{Aut}\, F$ 同构于群 $\mathrm{Aut}\, P$ 的对偶，故前一节记为 $\pi$ 的群正是 $\mathrm{Aut}\, P$. 回忆在前一节中我们从给定的基本函子 $F$ 构造了 $\mathcal{C}$ 与 $\mathcal{C}(\pi)$ (其中 $\pi = \operatorname{Aut}(F)$) 的等价，它把函子 $F$ 变为 $\mathcal{C}(\pi)$ 到有限集范畴的典范函子. 在这个典型情形 $\mathcal{C} = \mathcal{C}(\pi)$，$F =$ 典范函子中，与 $F$ 相伴的基本 pro-对象正是 $\pi$ 的离散商 $\pi_i$ 的投射系.

明确化 $\mathcal{C}(\pi)$ 的 pro-对象范畴可能是有用的. 得到：

\begin{proposition}\label{prop:V.5.2}
范畴 Pro-$\mathcal{C}(\pi)$ 典范等价于拓扑群 $\pi$ 的紧致完全不连通算子空间范畴 $\mathcal{C}'(\pi)$.
\end{proposition}

\begin{proof}
由于后者把 $\mathcal{C}(\pi)$ 作为全子范畴 (对应于紧致离散算子空间) 包含，且其中存在投射极限，无论如何有典范函子 $g \colon \mathrm{Pro-}\mathcal{C}(\pi) \to \mathcal{C}'(\pi)$，它给投射系 $Q = (Q_i)$ 对应 $\mathcal{C}'(\pi)$ 的对象 $X = \varprojlim_i Q_i$. 为定义反方向的函子，只需定义 $\mathcal{C}'(\pi)$ 到 $\mathcal{C}$-Ens 左正合函子范畴的反变函子，对每个 $X \in \mathcal{C}'(\pi)$ 取函子 $h(X)(E) = \operatorname{Hom}(X, E)$ (Hom 在 $\mathcal{C}'(\pi)$ 中取). 由定义立即看出函子 $h$ 和 $g$ 互为伴随，且 $hg$ 典范同构于 Pro-$\mathcal{C}(\pi)$ 的恒等函子. 剩下证明 (为建立 $g$ 和 $h$ 互为拟逆) $\mathcal{C}'(\pi)$ 的每个对象同构于形如 $g(Q)$ 的对象，其中 $Q \in \mathrm{Pro-}\mathcal{C}(\pi)$；换言之：每个以拓扑群 $\pi$ 为算子群、紧致且完全不连通的算子空间 $X$ 都同构于有限离散算子空间的投射极限. 由于 $X$ 是它的有限离散商 (作为无算子拓扑空间) 的投射极限，归结为证明在这些商的集合中存在 $\pi$-不变的共尾系. 为此只需证明：对这样的商 $X'$，$X'$ 在 $\pi$ 的作用下的变换集有限 (于是取这些变换的上确界，它是包含 $X'$ 的不变商). 或等价地，存在 $\pi$ 的开不变子群 $\pi'$，使该子群的元素不变 $X'$. 而 $X'$ 对应于 $X$ 的有限开集 $X_i$ 划分. 由 $\pi$ 的连续性和紧致性，存在单位元邻域 $V$，使得 $s \in V$ 蕴含对每个 $i$ 有 $s \cdot X_i \subset X_i$，故 $s$ 不变 $X'$. 而知道 $\pi$ 的开不变子群构成单位元的基本邻域组，证明完毕.
\end{proof}

注意还可以更简单地看到 Ind-$\mathcal{C}(\pi)$ 范畴典范等价于 $\pi$ 在其上连续作用的集合范畴. 我们在此用不到.

\begin{proposition}\label{prop:V.5.3}
设 $\mathcal{C}$ 是 Galois 范畴，$F$ 是 $\mathcal{C}$ 上的基本函子，$P = (P_i)$ 是按通常方式正规化的相伴 pro-对象. 设 $X \in \mathcal{C}$；为使 $X$ 连通，必要且充分条件是 $\pi$ 传递作用在 $E = F(X)$ 上.
\end{proposition}

\begin{proof}
归结到典型情形 $\mathcal{C} = \mathcal{C}(\pi)$，$F$ = 典范函子，在那里这是平凡的.
\end{proof}

\begin{corollary}\label{cor:V.5.4}
$X$ 的下列条件等价：(i) $X$ 连通且 $\neq \varnothing_{\mathcal{C}}$；(ii) 群 $\pi$ 在 $E = F(X)$ 上传递，且 $F(X) \neq \varnothing$；(iii) $X$ 同构于某个 $P_i$.
\end{corollary}

\begin{proof}
(i) 与 (iii) 的等价性也已经容易由第 4 节的 e) 得出.
\end{proof}

\begin{proposition}\label{prop:V.5.5}
设 $Q = (Q_i)_{i \in I}$ 是 $\mathcal{C}$ 的按通常方式正规化的 pro-对象，$G$ 是相应的 $\mathcal{C}$ 到 (Ens) 的函子 $G(X) = \operatorname{Hom}(Q, X)$. 下列条件等价：
\begin{enumerate}[(i)]
\item $G$ 与有限直和可交换；
\item $G$ 与两个对象的和可交换；
\item 诸 $Q_i$ 连通且 $\neq \varnothing$；
\item $Q$ 同构于 $\pi/H$，其中 $H$ 是 $\pi$ 的闭子群；
\item 函子 $G$ 同构于由 $\pi$ 的闭子群 $H$ 定义的函子 $E \mapsto E^{H}$ ($H$ 的不变元集).
\end{enumerate}
\end{proposition}

N.B. 在 (iv) 和 (v) 的陈述中，假设选定了基本函子，从而把 $\mathcal{C}$ 等同于范畴 $\mathcal{C}(\pi)$.

\paragraph*{证明.} 可以假设 $\mathcal{C} = \mathcal{C}(\pi)$. 蕴含 (i)$\Rightarrow$(ii) 平凡，(ii)$\Rightarrow$(iii) 像第 4 节的 d) 那样证明. 证明 (iii)$\Rightarrow$(iv). 事实上，知道 $\lim Q_i$ 作为非空有限集的投射极限非空. 设 $a$ 是 $\lim Q_i$ 的点，它定义算子空间同态
\[
\pi \longrightarrow Q,
\]
它满射，因为对每个 $i$，复合 $\pi \to Q \to Q_i$ 满射 (由 \ref{prop:V.5.3}，$\pi$ 在 $Q_i$ 中传递). 故若 $H$ 是 $a$ 的稳定化子群，得到同构 $\pi/H \xrightarrow{\sim} Q$. 蕴含 (iv)$\Rightarrow$(v) 和 (v)$\Rightarrow$(i) 又是平凡的.

\begin{proposition}\label{prop:V.5.6}
设 $\mathcal{C}$ 是 Galois 范畴，$P$ 是 $\mathcal{C}$ 的基本 pro-对象，$F$ 是相伴的基本函子. 设 $P' = (P'_i)_{i \in I}$ 是 $\mathcal{C}$ 的 pro-对象，化为标准形式，$F'$ 是相伴函子 $F'(X) = \operatorname{Hom}(P', X)$ ($\mathcal{C}$ 到 Ens). 下列条件等价：
\begin{enumerate}[(i)]
\item $P' \simeq P$，或 $F' \simeq F$；
\item $P'$ 基本，或 $F'$ 基本；
\item $F'$ 把两个对象的和变为和，且 $X \neq \varnothing_{\mathcal{C}}$ 蕴含 $F(X) \neq \varnothing$；
\item $\mathcal{C}$ 的连通且 $\neq \varnothing_{\mathcal{C}}$ 的对象恰是那些同构于某个 $P'_i$ 的对象.
\end{enumerate}
\end{proposition}

\begin{proof}
平凡地有 (i)$\Rightarrow$(iii) 和 (i)$\Rightarrow$(ii)，此外由 \ref{cor:V.5.4} (应用于 $P'$ 以代替 $P$) 有 (ii)$\Rightarrow$(iv). 再者，(iii) 或 (iv) 由 \ref{prop:V.5.5} 蕴含 $P'$ 形如 $\pi/H$，其中 $H$ 是 $\pi$ 的闭子群. 在情形 (iii)，对 $\pi$ 的每个开不变子群 $\pi'$ 存在 $\pi$-同态 $P' = \pi/H \to \pi/\pi'$，故 $H \subset \pi'$，从而 $H = (e)$，由此得 (i)，证毕.
\end{proof}

\begin{corollary}\label{cor:V.5.7}
设 $\mathcal{C}$ 是 Galois 范畴. 基本 pro-对象同构，基本函子同构.
\end{corollary}

\begin{proof}
换言之，基本函子范畴是连通\emph{群胚} (groupoid) $\Gamma$，可称为 Galois 范畴 $\mathcal{C}$ 的\emph{基本群胚} (fundamental groupoid)\index[term]{jibenqunpei@基本群胚}. 若 $\mathcal{C} = \mathcal{C}(\pi)$，则基本群胚的对象自同构群同构于 $\pi$，该同构在内自同构意义下完全确定. (N.B. 称所有态射都是同构的范畴为群胚，连通群胚即所有对象都同构的群胚.) $\mathcal{C}$ 的基本 pro-对象构成一个连通群胚，等价于基本群胚的对偶. 若 $F$，$F'$ 是两个基本函子，相伴的基本 pro-对象为 $P$，$P'$，则 $\mathrm{Hom}(F, F') = \mathrm{Isom}(F, F')$ 有时记为 $\pi_{F', F}$\index[notations]{$\pi_{F',F}$}\index[notations]{$\pi_{F',F}$}，扮演``从 $F$ 到 $F'$ 的\emph{道路类}集合''的角色. 特别地，$\pi_{F, F} = \pi_F$ 正是前一节构造的 $\mathcal{C}$ 在 $F$ 处的基本群. 至于与 $F$ 相伴的 pro-对象 $P$，它扮演 $\mathcal{C}$ 的终对象 $e_{\mathcal{C}}$ 在 $F$ 处的\emph{万有覆盖} (universal covering)\index[term]{wanyoufugai@万有覆盖} 的角色.
\end{proof}

给出 $\mathcal{C}$ (在等价意义下) 用其基本群胚 $\Gamma$ 的描述 (不经过选择其对象 $F$) 可能更方便. 对 $\mathcal{C}$ 的每个对象 $X$ 相伴基本群胚上的函子 $E_X$，定义为
\[
E_X(F) = F(X),
\]
取值于 Ens. (在拓扑中这样的函子称为群胚上的``\emph{局部系}'' (local system)\index[term]{jubuxi@局部系}.) $F(X) = E_X(F)$ 可称为 $X$ 在 $F$ 处的纤维 (fiber)，函子 $E_X$ 称为与 $X$ 相伴的纤维函子\index[term]{xianweihanzi@纤维函子}. 函子 $E_X$ 有下列性质：对每个 $F$，$E_X(F)$ 是拓扑群 $\pi_F = \operatorname{Aut}(F)$ 在其上连续作用的有限集. 对给定的从基本群胚到 Ens 的共变函子 $\xi$，前面的条件等价于对任意固定的 $F$ 的同样条件. 据此：

\begin{proposition}\label{prop:V.5.8}
函子 $X \mapsto E_X$ 是 $\mathcal{C}$ 与满足上面指出的条件的、$\mathcal{C}$ 的基本群胚 $\Gamma$ 到 Ens 的共变函子范畴之间的等价.
\end{proposition}

\begin{proof}
事实上，设 $F_0$ 是基本群胚的对象，$\pi_0 = \mathrm{Aut}(F_0)$，则函子 $\xi \mapsto \xi(F_0)$ 是 \ref{prop:V.5.8} 所考虑的第二个范畴与范畴 $\mathcal{C}(\pi_0)$ 之间的等价，立即看出. 另一方面，后者与函子 $X \mapsto E_X$ 的复合是自然等价 $\mathcal{C} \to \mathcal{C}(\pi_0)$. 由此得出函子 $X \mapsto E_X$ 本身是等价.
\end{proof}

\begin{corollary}\label{cor:V.5.9}
范畴 Pro-$\mathcal{C}$ 典范等价于基本群胚 $\Gamma$ 到拓扑空间范畴的共变函子 $\xi$ 的范畴，这些函子满足条件：对 $\Gamma$ 的每个对象 $F$，$\xi(F)$ 是以拓扑群 $\pi_F$ 为算子群的紧致完全不连通算子空间.
\end{corollary}

\begin{proof}
这里再次可以对 $\xi$ 验证这个条件，只需对一个 $F$ 验证. 证明与 \ref{prop:V.5.8} 相同.
\end{proof}

\begin{remark}\label{rem:V.5.10}
设 $(F_s)_{s \in S}$ 是基本群胚 $\Gamma$ 的对象族. 对 $s, s' \in S$ 令
\[
\operatorname{Hom}(s, s') = \operatorname{Hom}(F_s, F_{s'}),
\]
于是 $S$ 本身成为连通群胚，映射 $s \mapsto F_s$ 是 $S$ 到 $\Gamma$ 的全忠实函子，记为 $f$. 考虑 $\mathcal{C}$ 到函子范畴 $\mathrm{Hom}(S, \mathrm{Ens})$ 的函子 $X \mapsto E_X \circ f$，得到 \ref{prop:V.5.8} (和 \ref{cor:V.5.9}) 以 $\Gamma$ 换成 $S$ 的变体. 当 $S$ 约化为一点时，如此得到的陈述归结为定理 \ref{thm:V.4.1}，当 $S$ 是 $\Gamma$ 的对象集时正是 \ref{prop:V.5.8} 本身.
\end{remark}

我们将用 \ref{cor:V.5.9} 定义 $\mathcal{C}$ 的典范 pro-对象. 为此考虑 $\Gamma$ 到拓扑空间 (甚至拓扑群) 范畴的函子：
\[
f \colon F \longmapsto \operatorname{Aut}(F) = \pi_F.
\]
该函子满足 \ref{prop:V.5.8} 所考虑的条件，算子空间 $f(F)$ (在 $\pi_F$ 下) 正是 $\pi_F$ 本身，看作在自身内自同构作用下的算子空间. 故函子 $f$ 对应 $\mathcal{C}$ 的一个在唯一同构意义下确定的 pro-对象，它甚至是 $\mathcal{C}$ 的 pro-群，称为 $\mathcal{C}$ 的\emph{基本 pro-群} (fundamental pro-group)\index[term]{jiben pro-qun@基本 pro-群}，扮演\emph{基本群局部系}的角色. 因此它是 $\mathcal{C}$ 的 pro-群 $\Pi$，由下述条件定义：有对 $F$ 函子性的同构
\[
F(\Pi) \simeq \pi_F.
\]
若 $X$ 是 $\mathcal{C}$ 的任意 pro-对象，有典范态射
\[
\Pi \times X \longrightarrow X,
\]
它使 $X$ 成为 Pro-$\mathcal{C}$ 中以 $G$ 为左算子群的算子对象. 只需注意对变动的 $F$，有典范映射
\[
\Pi(F) \times X(F) \longrightarrow X(F),
\]
i.e.
\[
\operatorname{Aut}(F) \times E_X(F) \longrightarrow E_X(F), \quad \text{或} \quad \pi_F \times F(X) \longrightarrow F(X),
\]
它关于 $F$ 函子性. 它也关于 $X$ 函子性，故对 pro-对象的每个态射 $X \to Y$，相应的图交换.

\begin{remark}\label{rem:V.5.11}
务必不要把基本 pro-对象 $P$ (它不配备群结构，且连通) 与基本 pro-群 (它是 pro-群，一般不必连通) 混淆. 精确地说，$G$ 连通当且仅当 $\pi_F$ 通过内自同构作用在自身上传递，i.e. 当且仅当 $\pi$ 约化为单位元，或等价地 $\mathcal{C}$ 等价于有限集范畴. 另一个本质区别是 $G$ 在唯一同构意义下确定，而 $P$ 只在非唯一同构意义下确定.
\end{remark}

设 $E$ 是有限集，考虑群胚 $\Gamma$ 上取常值 $E$ 的函子：由 \ref{prop:V.5.8} 它定义 $\mathcal{C}$ 的对象，记为 $E_{\mathcal{C}}$，也可以解释为 $\mathcal{C}$ 的终对象 $e_{\mathcal{C}}$ 的 $E$ 份拷贝的和. 可以把 $E_{\mathcal{C}}$ 看作有限集范畴到范畴 $\mathcal{C}$ 的函子，该函子正合，故把有限群变为 $\mathcal{C}$ 的群，等等. 故若 $X$ 是有限群 $G$ 在其上右作用的 $\mathcal{C}$ 的对象，看到可以把 $X$ 看作具有右 $\mathcal{C}$-算子群 $G_{\mathcal{C}}$ 的 $\mathcal{C}$ 对象. 因此，按 $\mathcal{C}$-算子群对象一般术语的延伸，若 $X$ 在 $G_{\mathcal{C}}$ 下形式主齐性，i.e. 若由 $G_{\mathcal{C}}$ 在 $X$ 上的右作用得出的典范态射
\[
X \times G_{\mathcal{C}} \longrightarrow X \times X
\]
是同构，则称 $X$ 在 $G$ 下形式主齐性. 称 $X$ 在 $G$ 下主齐性\index[term]{zhuqixingduixiang@主齐性对象}，若它在 $G_{\mathcal{C}}$ 下主齐性，i.e. 若它形式主齐性，且此外商 $X/G = X/G_{\mathcal{C}}$ 是 $e_{\mathcal{C}}$. 若选定基本函子，从而给出 $\mathcal{C}$ 与某个范畴 $\mathcal{C}(\pi)$ 的等价，$X$ 对应一个 $\pi$ 在其上连续左作用的集合，即 $E = F(X)$. 让 $G$ 右作用在 $X$ 上，等价于让 $G$ 右作用在集合 $E$ 上，使 $G$ 的作用与 $\pi$ 的作用交换. 立即看到 $X$ 在 $G$ 下主齐性当且仅当集合 $E$ 是 $G$ 下的主齐性空间，i.e. 当且仅当 $G$ 在其上简单传递作用. (此外，$X$ 形式主齐性当且仅当 $E$ 主齐性或空.) 与 \ref{prop:V.5.3} 比较，看到若 $X$ 在 $G$ 下主齐性且连通，则给定同态 $G \to \operatorname{Aut}(X)$ 的对偶群是同构；且为使 $\mathcal{C}$ 的对象 $X$ 连通且在 $\mathrm{Aut}(X)$ 的对偶群下主齐性，必要且充分条件 (用第 4 节的记号) 是它同构于某个 Galois 的 $P_i$. 在典型情形 $\mathcal{C} = \mathcal{C}(\pi)$，这意味着 $X$ 同构于 $\pi$ 除以一个不变子群的商.

仍假设给定基本函子 $F$. 则给出在有限群 $G$ 下主齐性、带 $a \in F(X)$ 的 $X$，等价于给出从 $\pi$ 到群 $G$ 的同态. 事实上，这样的同态给出集合 $E = G$，让 $\pi$ 借助给定同态 $\pi \to G$ 和 $G$ 的左平移在其上左作用，让 $G$ 通过右平移右作用，$E$ 的标记点 $a$ 是 $G$ 的单位元. 由前面所述，由此以基本唯一的方式得到具有上述性质的每个三元组 $(X, G, a)$，因为带标记点的 $G$ 下主齐性集合等同于 $G$ 本身. 这样，有限群范畴到 Ens 的函子 $G \mapsto \operatorname{Hom}(\pi, G)$ 有直接的几何解释，该函子借助 $\pi$ pro-可表示，对它的考察因此给出与 $F$ 相伴的群 $\pi$ 的另一个构造.

\section{从一个 Galois 范畴到另一个的正合函子}

\begin{proposition}\label{prop:V.6.1}
设 $\mathcal{C}$，$\mathcal{C}'$ 是两个 Galois 范畴，$H \colon \mathcal{C} \to \mathcal{C}'$ 是共变函子，$F'$ 是 $\mathcal{C}'$ 上的基本函子，$F = F' \circ H$. 下列条件等价：
\begin{enumerate}[(i)]
\item $H$ 正合，i.e. 左正合且右正合；
\item $H$ 左正合，把有限和变为有限和，且把满态射变为满态射 (或：把 $\neq \varnothing_{\mathcal{C}}$ 的对象变为 $\neq \varnothing_{\mathcal{C}'}$ 的对象)；
\item $F$ 是 $\mathcal{C}$ 上的基本函子.
\end{enumerate}
\end{proposition}

\begin{proof}
蕴含 (i)$\Rightarrow$(ii) 是范畴论的一般事实. 此外 (ii) 的第一种形式蕴含第二种，这注意到：若 $X$ 是 $\mathcal{C}$ 的对象，则 $X \neq \varnothing_{\mathcal{C}}$ 当且仅当态射 $X \to e_{\mathcal{C}}$ 是满态射；注意 $F$ 假设左正合，把 $e_{\mathcal{C}}$ 变为 $e_{\mathcal{C}'}$. (ii) 的第二种形式蕴含 (iii)，因为 $F$ 左正合从而 pro-可表示，可以套用判则 \ref{prop:V.5.6} (iii). 最后 (iii) 蕴含 (i)，由 $F$ 正合且``保守'' (i.e. 满足第 4 节的公理 (G 6)) 的事实得出.
\end{proof}

设 $\Gamma$ 是 $\mathcal{C}$ 的基本群胚，$\Gamma'$ 是 $\mathcal{C}'$ 的基本群胚. 故若 $H$ 正合，则 $F' \mapsto F' \circ H$ 是从群胚 $\Gamma'$ 到群胚 $\Gamma$ 的函子，记为 ${}^{t}H$：
\[
{}^{t}H(F')(X) = F'(H(X)),
\]
它也可以用第 6 节引入的记号 $F(X) = E_X(F)$ 写成
\[
E_{H(X)}(F') = E_X({}^{t}H(F')).
\]
后一个公式 (考虑到 \ref{prop:V.5.8} 或 \ref{thm:V.4.1}) 表明正合函子 $H$ 在知道相应函子 ${}^{t}H$ 时 (在唯一同构意义下) 被确定. 固定 $F'$，令 $F = {}^{t}H(F')$，则 ${}^{t}H$ 定义同态 $\pi_{F'} = \operatorname{Aut}(F') \to \pi_F = \operatorname{Aut}(F)$：
\[
{}^{t}H \colon \pi_{F'} \longrightarrow \pi_F \quad (F = {}^{t}H(F') = F' \circ H).
\]
此外，上面的公式 (考虑到 \ref{prop:V.5.8}) 表明这个同态具有下列性质：对每个 $\pi_F$ 在其上连续作用的有限集 $E$，群 $\pi_{F'}$ 借助前面的同态 $\pi_{F'} \to \pi_F$ 也连续作用. 把它应用于 $\pi_F$ 除以开不变子群的商，看到前面的条件也意味着所考虑的同态连续. 反过来，给出 $\Gamma$ 的对象 $F$，$\Gamma'$ 的对象 $F'$ 和连续同态
\[
u \colon \pi_{F'} \longrightarrow \pi_F,
\]
它就给出从 $\mathcal{C}(\pi)$ 到 $\mathcal{C}(\pi')$ 的函子，显然正合，故由 \ref{thm:V.4.1} 它给出从 $\mathcal{C}$ 到 $\mathcal{C}'$ 的正合函子 $H$，使得 ${}^{t}H \colon \pi_{F'} \to \pi_F$ 恰为 $u$. 也可以不从群同态出发，而从函子
\[
U \colon \Gamma' \longrightarrow \Gamma
\]
出发，它满足：对每个 $F' \in \Gamma'$ (或一个 $F' \in \Gamma'$，都一样)，相应同态 $\pi_{F'} \to \pi_F$ 连续：这样的函子同构于形如 ${}^{t}H$ 的函子，其中 $H \colon \mathcal{C} \to \mathcal{C}'$ 是在唯一同构意义下确定的正合函子. 因此：

\begin{corollary}\label{cor:V.6.2}
为使 Galois 范畴的函子 $H \colon \mathcal{C} \to \mathcal{C}'$ 正合，必要且充分条件是存在等价 $\mathcal{C}(\pi) \to \mathcal{C}$ 和 $\mathcal{C}' \to \mathcal{C}(\pi')$，把函子 $H$ 变为与拓扑群同态 $\pi' \to \pi$ 相伴的函子 $\mathcal{C}(\pi) \to \mathcal{C}(\pi')$.
\end{corollary}

\begin{corollary}\label{cor:V.6.3}
设 $\mathcal{C}$，$\mathcal{C}'$ 是两个 Galois 范畴，$\Gamma$，$\Gamma'$ 是它们的基本群胚. 则从 $\mathcal{C}$ 到 $\mathcal{C}'$ 的正合函子范畴等价于具有下列性质的函子 $U \colon \Gamma' \to \Gamma$ 的范畴：对 $\Gamma'$ 中的每个 $F'$ (或 $\Gamma'$ 中的一个 $F'$，都一样)，令 $\mathcal{F} = \mathcal{U}(\mathcal{F}')$，由 $U$ 定义的同态
\[
\pi_{F'} = \operatorname{Aut}(F') \longrightarrow \pi_F = \operatorname{Aut}(F)
\]
连续.
\end{corollary}

\begin{proof}
考虑 $\mathcal{C}$ 的基本 pro-群 $\Pi$，正合函子 $H$ 把它变为 $\mathcal{C}'$ 的 pro-群 $H(\Pi)$，定义同态
\[
\Pi' \longrightarrow H(\Pi)
\]
(其中 $\Pi'$ 是 $\mathcal{C}'$ 的基本 pro-群)，其条件为：对 $\Gamma'$ 的每个对象 $F'$，相应的同态
\[
F'(\Pi') = \pi_{F'} \longrightarrow F'(H(\Pi)) = \pi_F \quad \left(\text{其中 } F = F' \circ H = {}^{t}H(F')\right)
\]
是自然同态
\[
\operatorname{Aut}(F') \longrightarrow \operatorname{Aut}(F' \circ H).
\]
由于后者关于 $F'$ 函子性，由 \ref{prop:V.5.8} 它确实定义 $\mathcal{C}'$ 的 pro-对象、事实上 pro-群的同态. 称该同态为与函子 $H$ 相伴的同态.
\end{proof}

现在设 $H'$ 是从 Galois 范畴 $\mathcal{C}'$ 到 Galois 范畴 $\mathcal{C}''$ 的第二个正合函子. 平凡地有
\[
{}^{t}(H'H) = {}^{t}H\,{}^{t}H'
\]
(N.B. 这里是函子的恒等，而不只是典范同构). 基本 pro-群的相伴同态有类似的传递性.

现在用相应同态
\[
u \colon \pi_{F'} \longrightarrow \pi_F \quad (\text{其中 } F' = F' \circ H)
\]
解释正合函子 $H$ 的性质.

引入 (配备基本函子 $F$ 的) Galois 范畴 $\mathcal{C}$ 的\emph{带点对象} (pointed object) 概念是方便的\index[term]{daidianduixiang@带点对象}：它按定义是配备 $F(X)$ 中元素 $a$ 的 $\mathcal{C}$ 的对象 $X$. 它因此解释为 $\pi_F$ 在其上连续左作用的有限集，配备一点 $a$. 故由 \ref{prop:V.5.3}，$\mathcal{C}$ 的连通带点对象等同于 $\pi_F$ 的开子群. 若 $U$ 和 $V$ 是两个这样的子群，对应于 $\mathcal{C}$ 的连通带点对象 $X$，$Y$，则存在从 $X$ 到 $Y$ 的带点同态当且仅当 $U \subset V$，且此时该同态唯一. 当然，函子 $H$ 把带点对象变为带点对象 (因为 $F = F' \circ H$). 另一方面注意，像 $\pi_F$ 这样的群的闭子群是包含它的开子群的交；因此，若 $M$，$N$ 是两个闭子群，则 $M \subset N$ 当且仅当每个包含 $N$ 的开子群也包含 $M$. 借助这些评注，容易证明下列结果：

\begin{proposition}\label{prop:V.6.4}
设 $X$ 是 $\mathcal{C}$ 的连通带点对象，对应于 $\pi_F$ 的开子群 $U$. 为使 $U$ 包含 $u(\pi_{F'})$ (相应地：$u(\pi_{F'})$ 生成的闭不变子群)，必要且充分条件是 $H(X)$ 有带点截面 (相应地：完全分裂).
\end{proposition}

称对象 $X$ (在终对象之上的) 截面 (section) 为 Galois 范畴 $\mathcal{C}$ 的对象 $X$ 的从终对象 $e_{\mathcal{C}}$ 出发的态射，它等价于给出 $\pi_F$ 不变的 $F(X)$ 的元素 $a$；若 $X$ 带点，当截面与 $X$ 和 $e_{\mathcal{C}}$ 上的带点结构相容，i.e. 当 $a$ 恰是 $F(X)$ 的标记对象时，称有带点截面 (pointed section). 这样的截面因此唯一，且存在当且仅当 $F(X)$ 的标记对象被 $\pi_F$ 不变. 最后，若 Galois 范畴的对象同构于若干终对象的和，i.e. 若 $\pi_F$ 在 $F(X)$ 中平凡作用 —— 当 $X$ 带点时，这显然是比存在带点截面更强的条件 —— 则称该对象\emph{完全分裂} (completely decomposed)\index[term]{wanquanfenlie@完全分裂}. 命题 \ref{prop:V.6.4} 由定义和前面的评注平凡地得出.

\begin{corollary}\label{cor:V.6.5}
为使 $u$ 平凡，必要且充分条件是：对 $\mathcal{C}$ 的每个对象 $X$，$H(X)$ 完全分裂.
\end{corollary}

\begin{proposition}\label{prop:V.6.6}
设 $X'$ 是 $\mathcal{C}'$ 的连通带点对象，对应于 $\pi_{F'}$ 的开子群 $U'$. 为使 $U'$ 包含 Ker $u$，必要且充分条件是存在 $\mathcal{C}$ 的连通带点对象 $X$ 和从 $H(X)$ 的连通带点分量 $X'_0$ 到 $X'$ 的带点同态 (故 $X$ 作为带点对象同构于 $H(X)$ 的原像的带点中性连通分量的商). 若 $u$ 满射，前面的条件也等价于：$X'$ 同构于某个 $H(X)$，其中 $X$ 是 $\mathcal{C}$ 的带点对象.
\end{proposition}

(Galois 范畴 $\mathcal{C}$ 的带点对象 $X$ 的中性连通分量 (neutral connected component)，是 $X$ 的惟一的连通带点子对象；由 \ref{prop:V.5.3}，它对应于 $\pi_F$ 在 $F(X)$ 的标记点上的轨道.) 由于``$U'$ 包含 Ker $u$''不依赖于 $X'$ 所选择的带点结构 (因为另一个带点结构就是把 $U$ 换成共轭子群 $\dot{\mathrm{a}} U$)，看到：

\begin{corollary}\label{cor:V.6.7}
为使 $U'$ 包含 Ker $u$，必要且充分条件是存在 $\mathcal{C}$ 的对象 $X$ (可以假设连通) 和从 $H(X)$ 的一个连通分量到 $X'$ 的态射. 若 $u$ 满射，这也意味着 $X'$ 同构于形如 $H(X)$ 的对象.
\end{corollary}

\begin{corollary}\label{cor:V.6.8}
为使 $u$ 单射，必要且充分条件是：对 $\mathcal{C}'$ 的每个对象 $X'$，存在 $\mathcal{C}$ 的对象 $X$ 和从 $H(X)$ 的一个连通分量到 $X'$ 的同态.
\end{corollary}

\begin{proposition}\label{prop:V.6.9}
下列条件等价：
\begin{enumerate}[(i)]
\item 同态 $u \colon \pi_{F'} \to \pi_F$ 满射；
\item 对 $\mathcal{C}$ 的每个连通对象 $X$，$H(X)$ 连通；
\item 函子 $H$ 全忠实.
\end{enumerate}
\end{proposition}

\begin{proof}
最后一个事实意味着对 $\mathcal{C}$ 的两个对象 $X$，$Y$，自然映射
\[
\operatorname{Hom}(X, Y) \longrightarrow \operatorname{Hom}(H(X), H(Y))
\]
双射.
\end{proof}

\begin{corollary}\label{cor:V.6.10}
为使 $u$ 是同构，必要且充分条件是 $H$ 是范畴等价，或等价地下列两个条件成立：
\begin{enumerate}[a)]
\item 对 $\mathcal{C}$ 的每个连通对象 $X$，$H(X)$ 连通；
\item $\mathcal{C}'$ 的每个对象同构于形如 $H(X)$ 的对象.
\end{enumerate}
\end{corollary}

\begin{proposition}\label{prop:V.6.11}
设 $H \colon \mathcal{C} \to \mathcal{C}'$ 和 $H' \colon \mathcal{C}' \to \mathcal{C}''$ 是 Galois 范畴之间的正合函子，$F''$ 是 $\mathcal{C}''$ 上的基本函子，令 $F' = F''H'$ 和 $F = F'H$，考虑相伴同态
\[
u' \colon \pi_{F''} \longrightarrow \pi_{F'} \qquad u \colon \pi_{F'} \longrightarrow \pi_F.
\]
为使 Ker $u \subset$ Im $u'$ (i.e. 为使 $uu'$ 是平凡同态)，必要且充分条件是：对 $\mathcal{C}$ 的每个对象 $X$，$H'(H(X))$ 完全分裂. 为使 Ker $u \supset$ Im $u'$，必要且充分条件是：对 $\mathcal{C}'$ 的每个使得 $H'(X')$ 有带点截面的连通带点对象 $X'$，存在 $\mathcal{C}$ 的对象 $X$ 和从 $H(X)$ 的一个连通分量到 $X'$ 的同态.
\end{proposition}

\begin{proof}
第一个断言由 \ref{prop:V.6.4} 的最后一个断言得出. 第二个断言由 \ref{prop:V.6.4} 和 \ref{prop:V.6.6} 的结合得出.
\end{proof}

\begin{remark}\label{rem:V.6.12}
一般地，在 \ref{cor:V.6.8} 的条件下，$X'$ 同构于形如 $H(X)$ 的对象并不成立. 可以证明，为使 $\mathcal{C}'$ 的每个连通 (从而每个) 对象同构于形如 $H(X)$ 的对象，必要且充分条件是 $u$ 是 $\pi_{F'}$ 到 $\pi_F$ 的一个直因子子群上的同构. 在实践中，用 $\mathcal{C}'$ 到 $\mathcal{C}$ 的适当正合函子直接构造 $u$ 的右逆 $\pi_F \to \pi_{F'}$.
\end{remark}

\begin{proposition}\label{prop:V.6.13}
设 $\mathcal{C}$ 是配备基本函子 $F$ 的 Galois 范畴，$S$ 是 $\mathcal{C}$ 的连通对象，$\mathcal{C}'$ 是 $\mathcal{C}$ 在 $S$ 之上的对象范畴. 则 $\mathcal{C}'$ 是 Galois 范畴，且从 $\mathcal{C}$ 到 $\mathcal{C}'$ 的函子 $X \mapsto H(X) = X \times S$ 正合. 设 $a \in F(S)$，设 $F'$ 是从 $\mathcal{C}'$ 到有限集范畴的函子，定义为
\[
F'(X') = a \text{ 在 } F(X') \longrightarrow F(S) \text{ 下的原像}.
\]
则有同构 $F \cong F' \circ H$，且相应同态
\[
u \colon \pi_{F'} \longrightarrow \pi_F
\]
是 $\pi_{F'}$ 到 $\pi_F$ 的稳定化子 $F(X)$ 的标记元素 $a$ 的开子群 $U$ 上的同构.
\end{proposition}

\begin{proof}
证明留给读者.
\end{proof}

\section{预概形的情形}

设 $S$ 是局部诺特连通预概形，设
\[
a \colon \operatorname{Spec}(\Omega) \longrightarrow S
\]
是 $S$ 的几何点 (geometric point)，取值于代数闭域 $\Omega$. 令
\[
\mathcal{C} = S \text{ 的平展覆盖的范畴}
\]
对 $\mathcal{C}$ 的对象 $X$，i.e. $S$ 的平展覆盖 $X$，令
\[
F(X) = X \text{ 在 } a \text{ 之上的几何点的集合}.
\]
于是 $F$ 成为 $\mathcal{C}$ 上取值于有限集范畴的函子. 性质 (G 1) 至 (G 6) 满足：对 (G 1)，它包含在 \ref{prop:I.4.6} 的平凡推论中；(G 2) 由 \ref{cor:V.3.4} 得出；(G 3) 由 \ref{prop:V.3.5} 得出；(G 4) 由定义平凡；(G 5) 由 \ref{prop:V.3.5} 和第 2 节开头得出；最后 (G 6) 在 \ref{prop:V.3.7} 中证明. 故可以应用第 4、5、6 节的结果. 这特别使我们能够定义表示 $F$ 的 $\mathcal{C}$ 的 pro-对象 $P$，称为 $S$ 在点 $a$ 的万有覆盖 (universal covering)\index[term]{wanyoufugai (yugaixing)@万有覆盖 (预概形)}，以及拓扑群 $\pi = \operatorname{Aut}(F) = \operatorname{Aut}^{0}P$，称为 $S$ 在 $a$ 处的基本群 (fundamental group)，记为 $\pi_1(S, a)$\index[notations]{$\pi_1(S, a)$}\index[notations]{$\pi_1(S, a)$}. 函子 $F$ 于是定义 $\mathcal{C}$ 与 $\pi = \pi_1(S, a)$ 在其上连续作用的有限集范畴的等价. 这个等价因此允许把覆盖上的通常投射极限和有限归纳极限运算 (积、纤维积、和、取商等) 解释为 $\mathcal{C}(\pi)$ 中的类似运算，i.e. 解释为 $\pi$ 在其上作用的有限集上的显然运算. 还注意，由于平展覆盖的拓扑连通分量也是平展覆盖，$\mathcal{C}$ 的对象 $X$ 在 $\mathcal{C}$ 中连通当且仅当 $X$ 拓扑连通；由 \ref{prop:V.5.3} 这意味着 $\pi_1$ 在 $F(X)$ 中传递作用. 注意，为使 $\mathcal{C}$ 的对象 $X$ 在 $S$ 上忠实平坦且拟紧 (它已在 $S$ 上平坦且拟紧)，必要且充分条件是 $X \to S$ 满射，i.e. 是 $\mathcal{C}$ 中的满态射，或 $X \neq \varnothing$. 于是由判则 \ref{prop:V.2.6} (iii) 得出结论：$X$ 是 $S$ 的以 $G$ 为 伽罗瓦群的主覆盖，当且仅当它是范畴 $\mathcal{C}$ 中 $G$ 下的主齐性空间 (如第 5 节所定义的).

若 $a'$ 是 $S$ 的另一个几何点 (对应于代数闭域 $\Omega'$，它可以不同于 $\Omega$，甚至可以有不同的特征)，它定义 $\mathcal{C}$ 到有限集范畴的纤维函子 $F' = F_{a'}$，它仍正合，故同构于 $F = F_a$. 因此，对变动的 $a$，基本群 $\pi_1(S; a)$ 彼此同构. 若用 $\pi_1(S; a, a')$ 记相伴纤维函子的同构\index[notations]{$\pi_1(S; a, a')$} (或等价地，同态) $F_a \to F_{a'}$ 的集合，则得到一个群胚，其对象集是 $S$ 的几何点集，基本群是所述群胚的对象的自同构群. 集合 $\pi_1(S; a', a)$ 可称为从 $a$ 到 $a'$ 的道路类集合\index[term]{daolulei@道路类}. 这些类因此以显然的方式复合. 最后，可以定义 $\mathcal{C}$ 的 pro-群 $\Pi_1^{S}$\index[notations]{$\Pi_1^S$}，可称为 $S$ 的基本 pro-群或 $S$ 上的基本群局部系\index[term]{jibenqunjubuxi@基本群局部系}，由下列条件在唯一同构意义下定义：有对 $S$ 的几何点 $a$ 函子性的同构
\[
F_a(\Pi_1^{S}) = \pi_1(S; a)
\]
(参见评注 \ref{rem:V.5.10}). 特别地，若 $s$ 是 $S$ 的普通点，则 $G$ 在 $s$ 处的纤维是 $k(s)$ 上的 pro-群，是 $k(s)$ 上有限平展群的投射极限；可以把这个 pro-群称为 $S$ 在普通点 $s$ 处的基本群，记为 $\pi_1(S, s)$. 按定义，$S$ 在 $k(s)$ 的代数闭扩张 $\Omega$ 中取值的这些点，正是 $\pi_1(S; a)$ 的元素，其中 $a$ 是由该扩张定义的 $S$ 的几何点. 特别地，(取 $S$ 为域的谱) 每个域 $k$ 典范地且函子性地相伴一个 $k$ 上的 pro-群，可记为 $\pi_1(k)$，它是 $k$ 上有限平展群的投射极限，其在 $k$ 的代数闭扩张 $\Omega$ 中的点等同于 $\overline{\{k\}}/k$ 的拓扑 伽罗瓦群 (其中 $\overline{k}$ 是 $k$ 在 $\Omega$ 中的 Galois 闭包) 的元素 (参见 \ref{prop:V.8.1}). 这个群 $\pi_1(k)$ 似乎尚未引起代数学家的注意.

现在设
\[
f \colon S' \longrightarrow S
\]
是从一个连通局部诺特预概形到另一个的态射，$a'$ 是 $S'$ 的几何点，$a = f(a')$ 是它在 $S$ 中的直接像. 则``原像''函子诱导从 $S$ 的平展覆盖范畴 $\mathcal{C}(S)$ 到 $S'$ 的平展覆盖范畴 $\mathcal{C}(S')$ 的函子：
\[
f^{\bullet} \colon \mathcal{C}(S) \longrightarrow \mathcal{C}(S').
\]
此外有函子同构
\[
F_a \cong F_{a'} \circ f^{\bullet},
\]
从而 $f^{\bullet}$ 是正合函子，可以对其应用第 6 节的结果. 特别有典范同态
\[
u = \pi_1(f; a')\colon \pi_1(S', a') \longrightarrow \pi_1(S; a) \quad (a' = f(a)),
\]
它可以重构原像函子，即作为算子群限制运算. 函子 $f^{\bullet}$ 的性质因此用相伴群同态的性质简单表达，如第 6 节所明确的. 若特别地 $S'$ 是 $S$ 的平展覆盖，则同态 $u$ 是 $\pi_1(S', a')$ 到定义 $S$ 的连通带点平展覆盖 $S'$ 的 $\pi_1(S, a)$ 的开子群 (i.e. $a' \in F_a(S')$ 在 $\pi_1(S; A)$ 中的稳定化子 $U$) 上的同构.

若想对变动的几何点 $a'$ 解释同态 $\pi_1(f; a')$，则应当 (按照第 6 节所说的) 考虑 pro-群同态
\[
\Pi_1(f) \colon \Pi_1^{S'} \longrightarrow f^{\bullet}(\Pi_1^{S}),
\]
并取几何纤维的相应同态.

\section{正规基预概形的情形}

\begin{proposition}\label{prop:V.8.1}
设 $S$ 是域 $k$ 的谱，$\Omega$ 是 $k$ 的代数闭扩张，定义 $S$ 在 $\Omega$ 中取值的几何点 $a$. 设 $\overline{k}$ 是 $k$ 在 $\Omega$ 中的可分闭包. 则存在 $\pi_1(S, a)$ 到 $\overline{k}/k$ 的拓扑 伽罗瓦群的典范同构.
\end{proposition}

\begin{proof}
设 $k'$ 是 $k$ 在 $\Omega$ 中的代数闭包，它对应 $S$ 在 $k'$ 中取值的几何点 $b$. 函子的自然同态 $F_b \to F_a$ 显然是同构，因为 $k$ 的有限可分扩张的 $k$-同态在 $\Omega$ 中取值必然在 $\overline{k}$ 中，更不用说在 $k'$ 中. 另一方面，$k$ 的 $k'/k$ 的自同构群 $\pi'$ 以显然的方式作用在 $F_b$ 上，由此得同态 $\pi' \to \operatorname{Aut}(F_b) \stackrel{\sim}{\longrightarrow} \operatorname{Aut}(F_a) = \pi_1(S; a)$. 另一方面，众所周知从 $\pi'$ 到 $\overline{k}/k$ 的自同构群 $\pi$ 的自然同态是同构. 由此得到典范同态 $\pi \to \pi_1(S; a)$，剩下证明它是同构. 事实上，该同态单射，因为核的元素是 $\overline{k}/k$ 的自同构，它在每个有限可分子扩张上诱导恒等，故平凡. 该同态满射，因为若 $X$ 是 $S$ 的连通平展覆盖，从而由 $k$ 的有限可分扩张 $L/k$ 定义，则 $\pi$ 在 $L$ 到 $k'$ 的 $k$-同态集合上传递，众所周知.
\end{proof}

\begin{proposition}\label{prop:V.8.2}
设 $S$ 是连通局部诺特正规预概形，$K = k(s)$ 是它的函数域 = 其在生成点 $s$ 处的剩余域，$\Omega$ 是 $K$ 的代数闭扩张，定义 $S' = \mathrm{Spec}(K)$ 的几何点 $a'$ 和 $S$ 的几何点 $a$. 则同态 $\pi_1(S'; a') \longrightarrow \pi_1(S; a)$ 满射. 当把第一个群等同于 $K$ 在 $\Omega$ 中的可分闭包 $\overline{K}$ 的 伽罗瓦群时 (参见 \ref{prop:V.8.1})，前面同态的核在 Galois 理论下对应于 $\overline{K}/K$ 的在 $S$ 上非分歧的 $K$ 的有限扩张组成的子扩张.
\end{proposition}

\begin{proof}
第一个断言意味着 $S$ 的连通平展覆盖 $X$ 在 $S'$ 上的原像连通，i.e. $X$ 整，这正是 (\ref{prop:I.10.1}). 前面同态的核解释为 $\overline{K}/K$ 的在集合 $F_a(X)$ 上诱导恒等的自同构 (可以假设 $S$ 的平展覆盖 $X$ 连通). 但这意味着该自同构在 $\overline{K}/K$ 的在 $S$ 上非分歧的有限子扩张上诱导恒等，这证明最后一个断言.
\end{proof}

\begin{remarkstar}
借助这种把正规预概形 $S$ 的基本群解释为通常 Galois 理论的方式，这个情形的定义很久以前就已知晓.
\end{remarkstar}

\section{非连通预概形的情形：多 Galois 范畴}

设 $S$ 是局部诺特预概形，$(S_i)_{i \in I}$ 是它的连通分量. 则 $S$ 的平展覆盖范畴 $\mathcal{C}(S)$ 等价于诸 $\mathcal{C}(S_i)$ 的乘积范畴，因此可以用诸 $S_i$ 的基本群 (在每个 $S_i$ 中选定一个几何点后) 描述. 在平展态射下降理论的应用中，有时为每个 $S_i$ 选择几何点很不方便. 更方便的是借助 \ref{prop:V.5.8} 的自然推广，把 $\mathcal{C}(S)$ 解释为 $S$ 的几何点群胚上的函子范畴，该群胚看作 $S$ 的诸连通分量相应群胚的和；所考虑的函子是取值于有限集范畴的函子，满足与 \ref{prop:V.5.8} 中所援引的类似的连续性性质. 在实践中，会有 $S$ 的几何点族 $(a_t)_{t \in E}$，使得 $S$ 的每个连通分量 $S_i$ 都含有一个这样的点，于是像 \ref{rem:V.5.10} 中那样，可以用底集为 $E$ 的类似群胚代替 $S$ 的所有几何点的群胚. 当然，这些考虑应当纳入等价于 $\mathcal{C}(\pi)$ 型范畴的乘积范畴的范畴的一般定义中，可称之为\emph{多 Galois 范畴} (multigaloisian categories)\index[term]{duo Galois fanchou@多 Galois 范畴}. 我们把细节留给读者.

\end{refsection}
\begin{refsection}% =====================================================================
% 第六讲 (Exposé VI) 纤维范畴与下降
% (Catégories fibrées et descente)
% A. Grothendieck
% 原书页码: pp. 116-152
% =====================================================================
\chapter{纤维范畴与下降}

% 与原书一致: 0. Introduction 编号为 0
\setcounter{section}{-1}

\section{引言}

与前一讲引言所宣布的相反，事实证明，不事先在一般范畴中足够仔细地发展下降的语言，就不可能对预概形范畴进行下降，即使在特殊情形下也不行.

"下降"的概念为对象的"粘合"(gluing) 乃至范畴的"粘合"的所有过程提供了普遍框架. 粘合最经典的情形与拓扑空间 $X$ 及用开集 $X_i$ 覆盖 $X$ 有关；若对每个 $i$ 在 $X_i$ 之上给一个纤维空间 (fibre space) $E_i$，并对每个指标对 $(i, j)$ 给一个同构 $f_{ji} \colon E_i|X_{ij} \to E_j|X_{ij}$ (其中 $X_{ij} = X_i \cap X_j$)，满足众所周知的传递条件 (简写为 $f_{kj} f_{ji} = f_{ki}$)，就知道在 $X$ 上存在纤维空间 $E$，它在唯一同构意义下由下列条件确定：存在同构 $f_i \colon E|X_i \xrightarrow{\sim} E_i$，满足关系 $f_{ji} = f_j f_i^{-1}$ (用通常的滥用记号). 设 $X'$ 是诸 $X_i$ 的和空间，它是 $X$ 上的纤维空间 (i.e. 配有连续映射 $X' \to X$). 可以把诸 $E_i$ 的给出更简洁地解释为 $X'$ 上的纤维空间 $E'$，把诸 $f_{ji}$ 解释为 $E'$ 在 $X'' = X' \times_X X'$ 上的两个逆像 $E''_1$ 和 $E''_2$ (按两个典范投影) 之间的同构，于是粘合条件可以写成三重纤维积 $X''' = X' \times_X X' \times_X X'$ 上纤维空间 $E'''_1$ 和 $E'''_3$ 之间的同构的恒等式 (其中 $E'''_i$ 表示 $E'$ 按第 $i$ 个典范投影在 $X'''$ 上的逆像). 从 $E'$ 和 $f$ 构造 $E$ 是"下降"过程的典型情形. 此外，从 $X$ 上的纤维空间 $E$ 出发，若存在 $X$ 的开覆盖 $(X_i)$ 使得诸 $E|X_i$ 同构于 $F \times X_i$，或等价地使 $E$ 在 $X' = \coprod_i X_i$ 上的逆像 $E'$ 同构于 $X' \times F$，则说 $X$ "局部平凡" (locally trivial)，纤维为 $F$.

这样，对象的"粘合"概念与性质的"局部化"概念，都与"基变换" $X' \to X$ 的某类研究有关. 在代数几何 (algebraic geometry) 中，还必须考虑许多其他类型的基变换，特别是忠实平坦态射 $X' \to X$，它们应被看作关于 $X$ 之上"预概形或其他对象"的"局部化"过程. 这类局部化与简单的拓扑局部化 (它不过是其特例) 使用得同样多. 在解析几何 (analytic geometry) 中 (程度较轻) 也是如此.

大多数证明归结为验证，因而省略或仅略述；在需要时我们说明证明中出现的较不显然的图式.

\section{宇宙、范畴、范畴等价}

为避免某些逻辑困难，这里采用"宇宙" (universe) 的概念，它是足够"大"的集合，使通常的集合论运算不越出它；一个"宇宙公理"保证每个对象都位于某个宇宙中. 细节参见 C. Chevalley 与报告人合著的一本正在准备的书\footnote{参见讲 VIII.}. 这样，记号 Ens 不表示所有集合的范畴 (这个观念无意义)，而表示位于给定宇宙中的集合的范畴 (我们在记号中不指明它). 同样，Cat 表示位于该宇宙中的范畴的范畴，Cat 的对象 $X$ 到对象 $Y$ 的"态射"按定义是 $X$ 到 $Y$ 的函子.

若 $\mathcal{C}$ 是范畴，用 $\mathrm{Ob}(\mathcal{C})$ 记 $\mathcal{C}$ 的对象集，用 Fl(C) 记 $\mathcal{C}$ 的箭集 (或 $\mathcal{C}$ 的态射集). 因此写 $X \in \mathrm{Ob}(\mathcal{C})$ 以避免通常的滥用记号 $X \in \mathcal{C}$. 若 $\mathcal{C}$ 和 $\mathcal{C}'$ 是两个范畴，$\mathcal{C}$ 到 $\mathcal{C}'$ 的函子总指通常所说的 $\mathcal{C}$ 到 $\mathcal{C}'$ 的共变函子；它的给出蕴含到达范畴和出发范畴 $\mathcal{C}$，$\mathcal{C}'$ 的给出. $\mathcal{C}$ 到 $\mathcal{C}'$ 的函子构成一个集合，记为 Hom$(\mathcal{C}, \mathcal{C}')$，它是一个记为 Hom$(\mathcal{C}, \mathcal{C}')$ 的范畴的对象集. 按定义，$\mathcal{C}$ 到 $\mathcal{C}'$ 的反变函子是 $\mathcal{C}$ 的对偶范畴 $\mathcal{C}^{\circ}$ 到 $\mathcal{C}'$ 的函子.

我们承认函子 $F \colon \mathcal{S} \to \mathcal{C}$ 的投射极限和归纳极限的概念，特别是这一概念最常见的特殊情形：笛卡儿积与纤维积，直和与余直和 (amalgamated sum) 的对偶概念，以及这些运算的通常形式性质.

例如，在刚才引入的范畴 Cat 中，投射极限 (相对于位于所选定宇宙中的范畴 $\mathcal{I}$) 存在；诸 $\mathcal{C}_i$ 的投射极限范畴 $\mathcal{C}$ 的对象集 (相应地：箭集) 通过取诸范畴 $\mathcal{C}_i$ 的对象集 (相应地：箭集) 的投射极限得到. 最熟知的情形是范畴族的产品. 我们将在下文不断使用两个范畴在第三个之上的纤维积.

关于范畴和函子的一切，在等待已经提到的正在准备的书的同时，参见 \cite{exp06-1} (即使在所论及的一般理论方面也必然很不完备).

借此机会明确范畴等价的概念，它在 \cite{exp06-1} 中阐述得不能令人满意. 函子 $F \colon \mathcal{C} \to \mathcal{C}'$ 称为忠实 (faithful) (相应地：全忠实 (fully faithful))，若对 $\mathcal{C}$ 的每对对象 $S$，$T$，映射 $u \mapsto F(u)$ (从 $\mathrm{Hom}(S, T)$ 到 $\mathrm{Hom}(F(S), F(T))$) 单射 (相应地：双射). 称 $F$ 是范畴等价\index[term]{fanchoudengjia@范畴等价}，若 $F$ 全忠实，且此外 $\mathcal{C}'$ 的每个对象 $S'$ 同构于形如 $F(S)$ 的对象. 证明这等价于：存在 $F$ 的拟逆 (quasi-inverse) 函子 $G \colon \mathcal{C}' \to \mathcal{C}$，i.e. 使 $GF$ 同构于 $\mathrm{id}_{\mathcal{C}}$. 若如此，则 $G$ 的给出和同构 $\varphi \colon GF \to \mathrm{id}_{\mathcal{C}'}$ 等价于：对每个 $S' \in \mathrm{Ob}(\mathcal{C}')$，给出由 $\mathcal{C}$ 的对象 $S$ 和同构 $u \colon F(S) \to S'$ 组成的对 $(S, u)$. (用这些记号，存在唯一函子 $\mathcal{C}' \to \mathcal{C}$，它以所给映射 $S' \mapsto G(S')$ 为对象映射，且使映射 $S' \mapsto \varphi(S')$ 成为函子同态 $FG \to \mathrm{id}_{\mathcal{C}'}$.) 最后，若 $G$ 是 $F$ 的拟逆函子，且选择同构 $\varphi \colon FG \xrightarrow{\sim} \mathrm{id}_{\mathcal{C}}$ 和 $\psi \colon GF \stackrel{\sim}{\longrightarrow} \mathrm{id}_{\mathcal{C}}$，则 \cite[I.1.2]{exp06-1} 中所陈述的关于 $\varphi$，$\psi$ 的两个相容性条件实际上彼此等价，且对任何选定的同构 $\varphi$，存在唯一同构 $\psi$ 使所述条件满足.

\section{一个范畴之上的范畴}

设 $\mathcal{E}$ 是 Univ 中的范畴，它是 Cat 的对象，可以考虑 $\mathcal{E}$ 之上的 Cat 对象的范畴 $\mathbf{Cat}_{/\mathcal{E}}$\index[notations]{$\mathbf{Cat}_{/\mathcal{E}}$}. 这个范畴的对象是函子
\[
p \colon \mathcal{F} \longrightarrow \mathcal{E}.
\]
也说配备这样函子的范畴 $\mathcal{F}$ 是 $\mathcal{E}$ 之上的范畴\index[term]{fanchouzhishangdefanchou@范畴之上的范畴}，或 $\mathcal{E}$-范畴. 把 $\mathcal{E}$ 之上的范畴 $\mathcal{F}$ 到 $\mathcal{E}$ 之上的范畴 $\mathcal{G}$ 的 $\mathcal{E}$-函子称为函子
\[
f \colon \mathcal{F} \longrightarrow \mathcal{G}
\]
使得
\[
q f = p,
\]
其中 $p$ 和 $q$ 分别是 $\mathcal{F}$ 和 $\mathcal{G}$ 的投影函子. 因此 $\mathcal{F}$ 到 $\mathcal{G}$ 的 $\mathcal{E}$-函子 $f$ 的集合与 $\mathbf{Cat}_{/\mathcal{E}}$ 中以 $\mathcal{F}$ 为源、$\mathcal{G}$ 为目标的箭的集合一一对应，尽管没有恒等 (因为像上面那样的 $f$ 的给出不把 $\mathcal{F}$ 和 $\mathcal{G}$ 确定为 $\mathcal{E}$ 之上的范畴)；但当然，如在任何其他范畴 $\mathcal{C}_{/S}$ 中那样，通常滥用语言，把 $\mathcal{E}$-函子 (按上面明确的含义) 等同于范畴 Cat$_{/\mathcal{E}}$ 中的箭.

用
\[
\operatorname{Hom}_{\mathcal{E}}(\mathcal{F}, \mathcal{G})
\]
记 $\mathcal{F}$ 到 $\mathcal{G}$ 的 $\mathcal{E}$-函子的集合. 当然，$\mathcal{E}$-函子的复合是 $\mathcal{E}$-函子 (所考虑的复合按定义对应于 $\mathbf{Cat}_{/\mathcal{E}}$ 中箭的复合).

现在考虑两个 $\mathcal{E}$-函子
\[
f, g \colon \mathcal{F} \longrightarrow \mathcal{G}
\]
和函子同态
\[
u \colon f \longrightarrow g.
\]
称 $u$ 为 $\mathcal{E}$-同态或"$\mathcal{E}$-函子同态"，若对每个 $\xi \in \mathrm{Ob}(\mathcal{F})$，有
\[
q(u(\xi)) = \mathrm{id}_{p(\xi)},
\]
用语言说：令 $S = p(\xi) = qf(\xi) = qg(\xi) \in \mathrm{Ob}(\mathcal{E})$，$\mathcal{G}$ 中的态射
\[
u(\xi) \colon f(\xi) \longrightarrow g(\xi)
\]
是 $\mathrm{id}_S$-态射. (一般地，对 $\mathcal{E}$ 中的每个态射 $\alpha \colon T \to S$ 和 $\mathcal{E}$ 之上的每个范畴 $\mathcal{G}$，若 $q(v) = \alpha$ (其中 $q$ 记投影函子 $\mathcal{G} \to \mathcal{E}$)，则 $\mathcal{G}$ 中的态射 $v$ 称为 $\alpha$-态射.) 若有第三个 $\mathcal{E}$-函子 $h \colon \mathcal{F} \to \mathcal{G}$ 和 $\mathcal{E}$-同态 $v \colon g \to h$，则 $vu$ 也是 $\mathcal{E}$-同态. 这样，$\mathcal{F}$ 到 $\mathcal{G}$ 的 $\mathcal{E}$-函子及它们的 $\mathcal{E}$-同态构成 $\mathcal{F}$ 到 $\mathcal{G}$ 的所有函子的范畴 Hom$(\mathcal{F}, \mathcal{G})$ 的一个子范畴，称它为 $\mathcal{F}$ 到 $\mathcal{G}$ 的 $\mathcal{E}$-函子范畴，记为
\[
\mathbf{Hom}_{\mathcal{E}/-}(\mathcal{F}, \mathcal{G}).
\]
它也是函子对
\[
R, S \colon \operatorname{Hom}(\mathcal{F}, \mathcal{G}) \Longrightarrow \operatorname{Hom}(\mathcal{F}, \mathcal{E})
\]
的核子范畴，其中 $R$ 是由 Hom$(\mathcal{F}, \mathcal{E})$ 的对象 $p$ 定义的常函子，$S$ 是由 $q \colon \mathcal{G} \to \mathcal{E}$ 定义的函子 $f \mapsto q \circ f$.

为了结束这些一般理论，还需定义范畴 Hom$_{\mathcal{E}/-}(\mathcal{F}, \mathcal{G})$ 之间由 $\mathcal{E}$-函子复合给出的自然配对. 换言之，想定义"复合函子"：
\begin{equation}
\mathbf{Hom}_{\mathcal{E}/-}(\mathcal{F}, \mathcal{G}) \times \mathbf{Hom}_{\mathcal{E}/-}(\mathcal{G}, \mathcal{H}) \longrightarrow \mathbf{Hom}_{\mathcal{E}/-}(\mathcal{F}, \mathcal{H}) \tag{i}
\label{eq:VI.i}
\end{equation}
(当 $\mathcal{F}, \mathcal{G}, \mathcal{H}$ 是 $\mathcal{E}$ 之上的三个范畴时)，使该函子对对象给出 $\mathcal{E}$-函子 $f \colon \mathcal{F} \to \mathcal{G}$ 和 $g \colon \mathcal{G} \to \mathcal{H}$ 的复合映射 $(f, g) \mapsto gf$. 为此，回忆定义典范函子：
\begin{equation}
\operatorname{Hom}(\mathcal{F}, \mathcal{G}) \times \operatorname{Hom}(\mathcal{G}, \mathcal{H}) \longrightarrow \operatorname{Hom}(\mathcal{F}, \mathcal{H}) \tag{ii}
\label{eq:VI.ii}
\end{equation}
它对对象正是函子复合映射 $(f, g) \mapsto gf$，并把箭 $(u, v)$ (其中
\[
u \colon f \longrightarrow f', \quad v \colon g \longrightarrow g'
\]
分别是 Hom$(\mathcal{F}, \mathcal{G})$ 和 Hom$(\mathcal{G}, \mathcal{H})$ 中的箭) 变为箭
\[
v * u \colon gf \longrightarrow g'f'
\]
由关系定义：
\[
v * u(\xi) = v(f'(\xi)) \cdot g(u(\xi)) = g'(u(\xi)) \cdot v(f(\xi))
\]
众所周知由此确实得到 $gf$ 到 $g'f'$ 的同态，且 (对变动的 $f, g$ 和 $u, v$) 得到函子 (ii)，i.e. 有
\begin{enumerate}[(I)]
\item $\mathrm{id}_g * \mathrm{id}_f = \mathrm{id}_{gf}$，
\item $(v' * u') \circ (v * u) = (v' \circ v) * (u' \circ u)$.
\end{enumerate}
还回忆典范配对 (ii) 的结合律公式，它一方面由函子复合的结合性 $(hg)f = h(gf)$ 表达，另一方面由公式
\begin{equation}
(w * v) * u = w * (v * u) \tag{III}
\label{eq:VI.III}
\end{equation}
表达 (其中 $u \colon f \to f'$，$v \colon g \to g'$ 如上，且还给出函子 $h, h' \colon \mathcal{H} \to \mathcal{H}$ 的同态 $w \colon h \to h'$). 现在说，当 $\mathcal{F}, \mathcal{G}$ 是 $\mathcal{E}$-范畴时，典范复合函子 (ii) 诱导函子 (i). 由于已知两个 $\mathcal{E}$-函子的复合是 $\mathcal{E}$-函子，这等于说：当 $u \colon f \to f'$ 和 $v \colon g \to g'$ 是 $\mathcal{E}$-函子同态时，$v * u \colon gf \to g'f'$ 也是 $\mathcal{E}$-函子同态. 这确实由定义平凡得出. 由于配对 (i) 由配对 (ii) 诱导，它们满足同样的结合律性质，也用公式 $(hg)f = h(gf)$ 和 $(w * v) * u = w * (v * u)$ 表达 (对 $\mathcal{E}$-函子和 $\mathcal{E}$-函子同态).

为完成公式表 (I)，(II)，(III)，还回忆公式：
\begin{equation}
v * \mathrm{id}_{\mathcal{F}} = v \quad \text{与} \quad \mathrm{id}_{\mathcal{G}} * u = u, \tag{IV}
\label{eq:VI.IV}
\end{equation}
其中为简化，当 $u$ 或 $v$ 分别是 $f$ 或 $g$ 的恒等自同构时，用 $v * f$ 或 $u * g$ 代替 $v * u$.

由配对 (i) 的定义得出，Hom$_{\mathcal{E}/-}(\mathcal{F}, \mathcal{G})$ 是 $\mathcal{F}, \mathcal{G}$ 的函子，从范畴 $\mathbf{Cat}_{/\mathcal{E}}^{\circ} \times \mathbf{Cat}_{/\mathcal{E}}$ 到 Cat. 若确实 $g \colon \mathcal{G} \to \mathcal{G}_1$ 是 $\mathcal{E}$-函子，i.e. Hom$_{\mathcal{E}/-}(\mathcal{G}, \mathcal{G}_1)$ 的对象，则在 (i) 中取 $\mathcal{H} = \mathcal{G}_1$，它对应函子
\[
g_{*} \colon \operatorname{Hom}_{\mathcal{E}/-}(\mathcal{F}, \mathcal{G}) \longrightarrow \operatorname{Hom}_{\mathcal{E}/-}(\mathcal{F}, \mathcal{G}_1).
\]
类似地，对 $\mathcal{E}$-函子 $f \colon \mathcal{F}_1 \to \mathcal{F}$，定义函子
\[
f^{*} \colon \operatorname{Hom}_{\mathcal{E}/-}(\mathcal{F}, \mathcal{G}) \longrightarrow \operatorname{Hom}_{\mathcal{E}/-}(\mathcal{F}_1, \mathcal{G}).
\]
为简洁，这些函子也用记号 $f \mapsto g \circ f$ (相应地：$g \mapsto g \circ f$) 表示 (它们事实上只表示对象集上的相应映射). 由上面指出的结合律性质得出，如此确实得到如所宣布的函子 $\mathbf{Cat}_{/\mathcal{E}}^{\circ} \times \mathbf{Cat}_{/\mathcal{E}} \to \mathbf{Cat}$.

\section{\texorpdfstring{$\mathcal{E}$}{E} 之上范畴中的基变换}

由于在 Cat 中投射极限 (相对于 Univ 的元素的范畴 $\mathcal{I}$) 存在，在 $\mathbf{Cat}_{/\mathcal{E}}$ 中也如此，特别笛卡儿积存在，它解释为 Cat 中的纤维积. 按照一般记号，若 $\mathcal{F}$ 和 $\mathcal{G}$ 是 $\mathcal{E}$ 之上的范畴，用
\[
\mathcal{F} \times_{\mathcal{E}} \mathcal{G}
\]
记它们在 $\mathbf{Cat}_{/\mathcal{E}}$ 中的积，i.e. 它们在 Cat 中 $\mathcal{E}$ 之上的纤维积，作为 $\mathcal{E}$ 之上的范畴. 因此 $\mathcal{F} \times_{\mathcal{E}} \mathcal{G}$\index[notations]{$\mathcal{F} \times_{\mathcal{E}} \mathcal{G}$} 配备两个 $\mathcal{E}$-函子 $\mathrm{pr}_1$ 和 $\mathrm{pr}_2$，它们对 $\mathcal{E}$ 之上的每个范畴 $\mathcal{H}$ 定义双射
\[
\operatorname{Hom}_{\mathcal{E}}(\mathcal{H}, \mathcal{F} \times_{\mathcal{E}} \mathcal{G}) \stackrel{\sim}{\longrightarrow} \operatorname{Hom}_{\mathcal{E}}(\mathcal{H}, \mathcal{F}) \times \operatorname{Hom}_{\mathcal{E}}(\mathcal{H}, \mathcal{G}).
\]
这个双射还来自范畴同构
\[
\mathbf{Hom}_{\mathcal{E}/-}(\mathcal{H}, \mathcal{F} \times_{\mathcal{E}} \mathcal{G}) \stackrel{\sim}{\longrightarrow} \mathbf{Hom}_{\mathcal{E}/-}(\mathcal{H}, \mathcal{F}) \times \mathbf{Hom}_{\mathcal{E}/-}(\mathcal{H}, \mathcal{G})
\]
(取两个成员的对象集)，其中所写的函子是以 $h \mapsto \mathrm{pr}_1 \circ h$ 和 $h \mapsto \mathrm{pr}_2 \circ h$ 为分量的函子. 验证确实得到同构留给读者 (在范畴的投射极限的更一般情形也成立，而不限于纤维积).

回忆 (如第 1 号中所说)
\[
\operatorname{Ob}(\mathcal{F} \times_{\mathcal{E}} \mathcal{G}) = \operatorname{Ob}(\mathcal{F}) \times_{\operatorname{Ob}(\mathcal{E})} \operatorname{Ob}(\mathcal{G})
\]
\[
\operatorname{Fl}(\mathcal{F} \times_{\mathcal{E}} \mathcal{G}) = \operatorname{Fl}(\mathcal{F}) \times_{\operatorname{Fl}(\mathcal{E})} \operatorname{Fl}(\mathcal{G}),
\]
箭的复合逐分量进行.

在下面，考虑函子
\[
\lambda \colon \mathcal{E}' \longrightarrow \mathcal{E}
\]
并对 $\mathcal{E}$ 之上的每个范畴 $\mathcal{F}$，借助 $\mathrm{pr}_2$ 把 $\mathcal{F} \times_{\mathcal{E}} \mathcal{E}'$ 看作 $\mathcal{E}'$ 之上的范畴；换言之，把"纤维积"运算解释为"基变换"运算，函子 $\lambda \colon \mathcal{E}' \to \mathcal{E}$ 称为"基变换函子". 按照众所周知的一般事实，得到函子，称为 $\lambda$ 的基变换函子：
\[
\lambda^{*} \colon \mathbf{Cat}_{\mathcal{E}}\index[notations]{$\lambda^{*}$} \longrightarrow \mathbf{Cat}_{\mathcal{E}'},
\]
("基限制"函子的伴随，后者给 $\mathcal{E}'$ 之上的范畴 $\mathcal{F}'$ (投影函子为 $p'$) 对应由函子 $p = \lambda p'$ 看作 $\mathcal{E}$ 之上范畴的 $\mathcal{F}'$). 如范畴中基变换函子的一般情形所熟知，基变换函子"与投射极限可交换"，特别"把" $\mathcal{E}$ 上的纤维积"变为" $\mathcal{E}'$ 上的纤维积.

设 $\mathcal{F}$ 和 $\mathcal{G}$ 是 $\mathcal{E}$ 之上的两个范畴；定义典范同构：
\begin{equation}
\mathbf{Hom}_{\mathcal{E}'/-}(\mathcal{F}', \mathcal{G}') \xrightarrow{\sim} \mathbf{Hom}_{\mathcal{E}/-}(\mathcal{F} \times_{\mathcal{E}} \mathcal{E}', \mathcal{G}) \tag{i}
\label{eq:VI.i2}
\end{equation}
其中 $\mathcal{F}' = \mathcal{F} \times_{\mathcal{E}} \mathcal{E}'$，$\mathcal{G}' = \mathcal{G} \times_{\mathcal{E}} \mathcal{E}'$.

为此考虑函子
\[
\mathrm{pr}_1 \colon \mathcal{G}' = \mathcal{G} \times_{\mathcal{E}} \mathcal{E}' \longrightarrow \mathcal{G},
\]
并用
\[
F \longmapsto \mathrm{pr}_1 \circ F
\]
定义 (i)，它先验地表示函子
\[
\operatorname{Hom}(\mathcal{F}', \mathcal{G}') \longrightarrow \operatorname{Hom}(\mathcal{F}', \mathcal{G})
\]
只需验证后者对子范畴诱导函子 (i)，且它是同构. (ii) 诱导双射
\[
\operatorname{Hom}_{\mathcal{E}'/-}(\mathcal{F}', \mathcal{G}') \stackrel{\sim}{\longrightarrow} \operatorname{Hom}_{\mathcal{E}/-}(\mathcal{F} \times_{\mathcal{E}} \mathcal{E}', \mathcal{G})
\]
是基变换函子的特征性质. 剩下证明：若 $F$，$G$ 是 $\mathcal{F}' \to \mathcal{G}'$ 的 $\mathcal{E}'$-函子，则映射
\[
u \longmapsto \mathrm{pr}_1 \circ u
\]
诱导双射
\[
\operatorname{Hom}_{\mathcal{E}'}(F, G) \stackrel{\sim}{\longrightarrow} \operatorname{Hom}_{\mathcal{E}}(\mathrm{pr}_1 \circ F, \mathrm{pr}_1 \circ G).
\]
这个事实的验证是直接的，留给读者.

由这个同构 (i) 和前一节末尾，得出
\[
\mathbf{Hom}_{\mathcal{E}'/-}(\mathcal{F} \times_{\mathcal{E}} \mathcal{E}', \mathcal{G} \times_{\mathcal{E}} \mathcal{E}')
\]
可以看作 $\mathcal{E}', \mathcal{F}, \mathcal{G}$ 的函子，从范畴 $\mathbf{Cat}_{/\mathcal{E}}^{\circ} \times \mathbf{Cat}_{/\mathcal{E}}^{\circ} \times \mathbf{Cat}_{/\mathcal{E}}$ 到 Cat，同构于由表达式 Hom$_{\mathcal{E}/-}(\mathcal{F} \times_{\mathcal{E}} \mathcal{E}', \mathcal{G})$ 定义的函子. 特别地，对固定的 $\mathcal{F}, \mathcal{G}$，得到 $\mathcal{E}'$ 的函子 $\mathbf{Hom}_{\mathcal{E}'/-}(\mathcal{F}', \mathcal{G}') = \mathbf{Hom}_{\mathcal{E}/-}(\mathcal{F} \times_{\mathcal{E}} \mathcal{E}', \mathcal{G} \times_{\mathcal{E}} \mathcal{E}')$，特别投影 $\mathcal{E}$-函子 $\lambda \colon \mathcal{E}' \to \mathcal{E}$ 定义态射，i.e. 函子
\[
\lambda_{\mathcal{F},\mathcal{G}}^{*} \colon \operatorname{Hom}_{\mathcal{E}/-}(\mathcal{F}, \mathcal{G}) \longrightarrow \operatorname{Hom}_{\mathcal{E}'/-}(\mathcal{F}', \mathcal{G}')
\]
我们把它明确化. 对两个成员的对象集，它是映射
\[
f \longmapsto f' = f \times_{\mathcal{E}} \mathcal{E}'
\]
它表达 $\mathcal{F} \times_{\mathcal{E}} \mathcal{E}'$ 对 $\mathcal{E}$ 上的对象 $\mathcal{F}$ 的函子性依赖. 另一方面，考虑两个 $\mathcal{E}$-函子
\[
f, g \colon \mathcal{F} \longrightarrow \mathcal{G}
\]
和 $\mathcal{E}$-函子同态
\[
u \colon f \longrightarrow g,
\]
明确相应的 $\mathcal{E}'$-函子同态：
\[
u' \colon f' \longrightarrow g'.
\]
对每个
\[
\xi' = (\xi, S') \in \operatorname{Ob}(\mathcal{F}')
\]
(其中
\[
\xi \in \operatorname{Ob}(\mathcal{F}), \quad S' \in \operatorname{Ob}(\mathcal{E}'), \quad p(\xi) = \lambda(S') = S
\])
态射
\[
u'(\xi') \colon f'(\xi') = (f(\xi), S') \longrightarrow g'(\xi') = (g(\xi), S') \quad \text{在 } \mathcal{G}' \text{ 中}
\]
由公式定义：
\[
u'(\xi') = (u(\xi), \mathrm{id}_{S'})
\]
(这确实是 $\mathcal{G}'$ 中的 $S'$-态射，因为 $q(u(\xi)) = \lambda(\mathrm{id}_{S'}) = \mathrm{id}_S$).

现在考虑任意 $\mathcal{E}$-函子
\[
\lambda' \colon \mathcal{E}'' \longrightarrow \mathcal{E}'
\]
和相应函子
\[
\mathbf{Hom}_{\mathcal{E}'/-}(\mathcal{F} \times_{\mathcal{E}} \mathcal{E}', \mathcal{G} \times_{\mathcal{E}} \mathcal{E}') \longrightarrow \mathbf{Hom}_{\mathcal{E}''/-}(\mathcal{F} \times_{\mathcal{E}} \mathcal{E}'', \mathcal{G} \times_{\mathcal{E}} \mathcal{E}'')
\]
我说这个函子正是用前述方法从 $\mathcal{E}'$ 上的 $\mathcal{F}'$ 和 $\mathcal{G}'$ 出发、把 $\mathcal{E}''$ 看作 $\mathcal{E}'$ 之上的范畴得到的函子，考虑到"基变换传递性"同构：
\[
\mathcal{F}' \times_{\mathcal{E}'} \mathcal{E}'' \stackrel{\sim}{\longrightarrow} \mathcal{F}'' = \mathcal{F} \times_{\mathcal{E}} \mathcal{E}'' \quad \text{与} \quad \mathcal{G}' \times_{\mathcal{E}'} \mathcal{E}'' \stackrel{\sim}{\longrightarrow} \mathcal{G}'' = \mathcal{G} \times_{\mathcal{E}} \mathcal{E}'',
\]
它蕴含典范同构
\[
\mathbf{Hom}_{\mathcal{E}''/-}(\mathcal{F}' \times_{\mathcal{E}'} \mathcal{E}'', \mathcal{G}' \times_{\mathcal{E}'} \mathcal{E}'') \stackrel{\sim}{\longrightarrow} \mathbf{Hom}_{\mathcal{E}''/-}(\mathcal{F} \times_{\mathcal{E}} \mathcal{E}'', \mathcal{G} \times_{\mathcal{E}} \mathcal{E}'')
\]
这个相容性的验证是直接的，留给读者.

刚才定义的函子与前一节定义的配对相容，精确地说，若 $\mathcal{F}, \mathcal{G}, \mathcal{H}$ 是 $\mathcal{E}$ 之上的范畴且令
\[
\mathcal{F}' = \mathcal{F} \times_{\mathcal{E}} \mathcal{E}', \quad \mathcal{G}' = \mathcal{G} \times_{\mathcal{E}} \mathcal{E}', \quad \mathcal{H}' = \mathcal{H} \times_{\mathcal{E}} \mathcal{E}',
\]
则在函子图式中有交换性：
\[
\begin{tikzcd}[row sep=large, column sep=large]
\mathbf{Hom}_{\mathcal{E}/-}(\mathcal{F}, \mathcal{G}) \times \mathbf{Hom}_{\mathcal{E}/-}(\mathcal{G}, \mathcal{H}) \arrow[r] \arrow[d,"{\lambda_{\mathcal{F},\mathcal{G}}^{*} \times \lambda_{\mathcal{G},\mathcal{H}}^{*}}"] & \mathbf{Hom}_{\mathcal{E}/-}(\mathcal{F}, \mathcal{H}) \arrow[d,"{\lambda_{\mathcal{F},\mathcal{H}}^{*}}"] \\
\mathbf{Hom}_{\mathcal{E}'/-}(\mathcal{F}', \mathcal{G}') \times \mathbf{Hom}_{\mathcal{E}'/-}(\mathcal{G}', \mathcal{H}') \arrow[r] & \mathbf{Hom}_{\mathcal{E}'/-}(\mathcal{F}', \mathcal{H}')
\end{tikzcd}
\]
其中水平箭是前一节定义的复合函子. 这个交换性用公式表达：
\[
(gf)' = g'f'
\]
(对 $f \in \mathrm{Hom}_{\mathcal{E}}(\mathcal{F}, \mathcal{G})$，$g \in \mathrm{Hom}_{\mathcal{E}}(\mathcal{G}, \mathcal{H})$，该公式简单表达基变换的函子性)，以及公式
\[
(v * u)' = v' * u'
\]
(当 $u \colon f \to f_1$ 是 Hom$_{\mathcal{E}/-}(\mathcal{F}, \mathcal{G})$ 的箭，$v \colon g \to g_1$ 是 Hom$_{\mathcal{E}/-}(\mathcal{G}, \mathcal{H})$ 的箭时). 这个公式的验证由定义容易得出.

下面，我们主要关心 Hom$_{\mathcal{E}}(\mathcal{F}, \mathcal{G})$ (及其些值得注意的子范畴) 当 $\mathcal{F} = \mathcal{E}$ 的情形，为此引入特殊记号：
\[
\boldsymbol{\Gamma}(\mathcal{G}/\mathcal{E}) = \mathbf{Hom}_{\mathcal{E}}(\mathcal{E}, \mathcal{G}), \quad \Gamma(\mathcal{G}/\mathcal{E}) = \mathrm{Ob}\index[notations]{$\Gamma(\mathcal{G}/\mathcal{E})$}(\boldsymbol{\Gamma}(\mathcal{G}/\mathcal{E})) = \mathrm{Hom}_{\mathcal{E}}(\mathcal{E}, \mathcal{G}).
\]

\begin{remarkstar}
当 $\mathcal{E}$ 是单点范畴，i.e. $\mathrm{Ob}(\mathcal{E})$ 和 $\mathrm{Fl}(\mathcal{E})$ 约化为一个元素，这也就意味着 $\mathcal{E}$ 是 Cat 的终对象，则 $\mathcal{E}$ 之上范畴的给出等价于一个范畴的给出 (因为将有 $\mathcal{F}$ 到 $\mathcal{E}$ 的唯一函子). 更精确地说，$\mathbf{Cat}_{/\mathcal{E}}$ 此时同构于 Cat. 此外，范畴 Hom$_{\mathcal{C}/-}(\mathcal{F}, \mathcal{G})$ 此时正是 $\mathbf{Hom}(\mathcal{F}, \mathcal{G})$. 回忆基本公式
\[
\operatorname{Hom}(\mathcal{H}, \operatorname{Hom}(\mathcal{F}, \mathcal{G})) \stackrel{\sim}{\longrightarrow} \operatorname{Hom}(\mathcal{F} \times \mathcal{H}, \mathcal{G})
\]
(对其中出现的三个自变量函子性的同构)，允许公理化地、以范畴 Cat 的內在术语解释 Hom$(\mathcal{F}, \mathcal{G})$，使得 Hom 范畴的已知公式表看起来是"Hom 对象"存在的范畴 (如 Cat) 中有效公式表的特殊情形. Hom$_{\mathcal{E}/-}(\mathcal{F}, \mathcal{G})$ 有类似的解释，当重新假设 $\mathcal{E}$ 任意时，由公式
\[
\operatorname{Hom}(\mathcal{H}, \mathbf{Hom}_{\mathcal{E}/-}(\mathcal{F}, \mathcal{G})) \stackrel{\sim}{\longrightarrow} \operatorname{Hom}_{\mathcal{E}}(\mathcal{F} \times \mathcal{H}, \mathcal{G})
\]
给出 (对三个自变量函子性的同构). 这样，第 2、3 号中阐述的形式性质是更一般结果的特殊情形，后者在"Hom 对象" $_{\mathcal{E}/-}(\mathcal{F}, \mathcal{G})$ (当 $\mathcal{F}, \mathcal{G}$ 是第三个范畴 $\mathcal{E}$ 之上的两个对象时) 存在的范畴中成立.
\end{remarkstar}

\section{纤维范畴；\texorpdfstring{$\mathcal{E}$}{E}-范畴的等价}

设 $\mathcal{F}$ 是 $\mathcal{E}$ 之上的范畴，$S \in \mathrm{Ob}(\mathcal{E})$. 称 $\mathcal{F}$ 在 $S$ 处的\emph{纤维范畴} (fiber category)\index[term]{xianweifanchou (xianwei)@纤维范畴 (纤维)} 为 $\mathcal{F}$ 的子范畴 $\mathcal{F}_S$，它是 $\mathcal{E}$ 的由 $S$ 定义的单点子范畴的原像. 故 $\mathcal{F}_S$ 的对象是 $p(\xi) = S$ 的 $\mathcal{F}$ 的对象 $\xi$，其态射是 $p(u) = \mathrm{id}_S$ 的 $\mathcal{F}$ 的态射 $u$，i.e. $\mathcal{F}$ 中的 $S$-态射. 当然，$\mathcal{F}_S$ 典范同构于纤维积 $\mathcal{F} \times_{\mathcal{E}} \{S\}$，其中 $\{S\}$ 表示 $\mathcal{E}$ 的由 $S$ 定义的单点子范畴，配备其到 $\mathcal{E}$ 的包含函子. (考虑到基变换的传递性) 由此得出，若作基变换 $\lambda \colon \mathcal{E}' \to \mathcal{E}$，则对每个 $S' \in \mathrm{Ob}(\mathcal{E}')$，投影 $\mathrm{pr}_1 \colon \mathcal{F}' = \mathcal{F} \times_{\mathcal{E}} \mathcal{E}' \to \mathcal{F}$ 诱导同构
\[
\mathcal{F}'_{S'} \longrightarrow \mathcal{F}_S \quad (\text{其中 } S = \lambda(S')).
\]

\begin{proposition}\label{prop:VI.4.1}
设 $f \colon \mathcal{F} \to \mathcal{G}$ 是 $\mathcal{E}$-函子. 若 $f$ 全忠实，则对任何基变换 $\mathcal{E}' \to \mathcal{E}$，相应函子 $f' \colon \mathcal{F}' = \mathcal{F} \times_{\mathcal{E}} \mathcal{E}' \to \mathcal{G}' = \mathcal{G} \times_{\mathcal{E}} \mathcal{E}'$ 全忠实.
\end{proposition}

验证是直接的；更一般地，可以证明全忠实函子的任何投射极限 (这里 $f$ 和 $\mathcal{E}$，$\mathcal{E}'$ 中的恒等函子) 是全忠实函子.

注意与 \ref{prop:VI.4.1} 类似、把"全忠实"换成"范畴等价"的断言是假的，对 $\mathcal{G} = \mathcal{E}$ 已如此. 然而：

\begin{proposition}\label{prop:VI.4.2}
设 $f \colon \mathcal{F} \to \mathcal{G}$ 是 $\mathcal{E}$-函子. 下列条件等价：
\begin{enumerate}[(i)]
\item 存在 $\mathcal{E}$-函子 $g \colon \mathcal{G} \to \mathcal{F}$ 和 $\mathcal{E}$-同构
\[
gf \stackrel{\sim}{\longrightarrow} \mathrm{id}_{\mathcal{F}}, \quad fg \stackrel{\sim}{\longrightarrow} \mathrm{id}_{\mathcal{G}}.
\]
\item 对 $\mathcal{E}$ 之上的每个范畴 $\mathcal{E}'$，函子
\[
f' = f \times_{\mathcal{E}} \mathcal{E}' \colon \mathcal{F}' = \mathcal{F} \times_{\mathcal{E}} \mathcal{E}' \longrightarrow \mathcal{G}' = \mathcal{G} \times_{\mathcal{E}} \mathcal{E}'
\]
是范畴等价.
\item $f$ 是范畴等价，且对每个 $S \in \mathrm{Ob}(\mathcal{E})$，$f$ 诱导的函子 $f_S \colon \mathcal{F}_S \to \mathcal{G}_S$ 是范畴等价.
\item[(iii bis)] $f$ 全忠实，且对每个 $S \in \mathrm{Ob}(\mathcal{E})$ 和每个 $\eta \in \mathrm{Ob}(\mathcal{G}_S)$，存在 $\xi \in \mathrm{Ob}(\mathcal{F}_S)$ 和 $S$-同构 $u \colon f(\xi) \to \eta$.
\end{enumerate}
\end{proposition}

\paragraph*{证明.} 显然 (i) 蕴含 $f$ 是范畴等价 (该概念用同样条件定义，但其中的函子同构不必是 $\mathcal{E}$-态射). 另一方面，由前一节的函子性得出，条件 (i) 在基变换 $\mathcal{E}' \to \mathcal{E}$ 后保持. 由此得 (i)$\Rightarrow$(ii). 显然 (ii)$\Rightarrow$(iii)，因为只需作 $\mathcal{E}' = \mathcal{E}$ 和 $\mathcal{E}' = \{S\}$. (iii)$\Rightarrow$(iii bis) 更平凡，剩下证明 (iii bis)$\Rightarrow$(i). 为此，对每个 $\eta \in \mathrm{Ob}(\mathcal{G})$ 选择 $g(\eta) \in \mathrm{Ob}(\mathcal{F})$ 和同构 $u(\eta) \colon f(g(\eta)) \to \eta$，使得 $q(u(\eta)) = \mathrm{id}_S$，其中 $S = q(\eta)$. 由第二个条件 (iii bis) 这是可能的. 由于 $f$ 全忠实，众所周知且立即看出，$g$ 可以唯一地看作 $\mathcal{G}$ 到 $\mathcal{F}$ 的函子，使诸 $u(\eta)$ 定义函子同态 (从而是同构) $u \colon fg \stackrel{\sim}{\longrightarrow} \mathrm{id}_{\mathcal{G}}$. 此外，由构造 $g$ 是 $\mathcal{E}$-函子，$u$ 是 $\mathcal{E}$-同态. 与前述数据对应的是函子同构 $v \colon gf \to \mathrm{id}_{\mathcal{F}}$，由条件 $f * v = u * f$ 定义，立即看出它也是 $\mathcal{E}$-态射，证毕.

\begin{definition}\label{def:VI.4.3}
若前面的条件满足，则称 $f$ 是 $\mathcal{E}$ 之上的范畴等价，或 $\mathcal{E}$-等价.
\end{definition}

\begin{corollary}\label{cor:VI.4.4}
假设投影函子 $p \colon \mathcal{F} \to \mathcal{E}$ 是可运输函子 (transportable functor)，i.e. 对 $\mathcal{E}$ 中的每个同构 $\alpha \colon T \to S$ 和 $\mathcal{F}_T$ 的每个对象 $\xi$，存在 $\mathcal{F}$ 中以 $\xi$ 为源、满足 $p(u) = \alpha$ 的同构 $u$. 则 $\mathcal{F}$ 到 $\mathcal{G}$ 的每个是范畴等价的 $\mathcal{E}$-函子 $f \colon \mathcal{F} \to \mathcal{G}$ 是 $\mathcal{E}$-等价.
\end{corollary}

由判则 (iii bis) 得出.

\begin{corollary}\label{cor:VI.4.5}
设 $f \colon \mathcal{F} \to \mathcal{G}$ 是 $\mathcal{E}$-等价. 则对 $\mathcal{E}$ 之上的每个范畴 $\mathcal{H}$，相应函子
\[
\operatorname{Hom}_{\mathcal{E}/-}(\mathcal{G}, \mathcal{H}) \longrightarrow \operatorname{Hom}_{\mathcal{E}/-}(\mathcal{F}, \mathcal{H})
\]
\[
\mathbf{Hom}_{\mathcal{E}/-}(\mathcal{H}, \mathcal{F}) \longrightarrow \mathbf{Hom}_{\mathcal{E}/-}(\mathcal{H}, \mathcal{G})
\]
(参见第 2 号) 是范畴等价.
\end{corollary}

由判则 (i) 按通常推理得出.

\section{笛卡儿态射、逆像、笛卡儿函子}

设 $\mathcal{F}$ 是 $\mathcal{E}$ 之上的范畴，投影函子为 $p$.

\begin{definition}\label{def:VI.5.1}
考虑 $\mathcal{F}$ 中的态射
\[
\alpha \colon \eta \longrightarrow \xi
\]
并设
\[
S = p(\xi), \quad T = p(\eta), \quad f = p(\alpha).
\]
称 $\alpha$ 是\emph{笛卡儿态射} (cartesian morphism)\index[term]{dikaertaishe@笛卡儿态射}，若对每个 $\eta' \in \mathrm{Ob}(\mathcal{F}_T)$ 和每个 $f$-态射 $u \colon \eta' \to \xi$，存在唯一 $T$-态射 $\overline{u} \colon \eta' \to \eta$，使得 $u = \alpha \circ \overline{u}$.

因此，这意味着对每个 $\eta' \in \mathrm{Ob}(\mathcal{F}_T)$，映射 $v \mapsto \alpha \circ v$
\begin{equation}
\operatorname{Hom}_T(\eta', \eta) \longrightarrow \operatorname{Hom}_f(\eta', \xi) \tag{i}
\label{eq:VI.i3}
\end{equation}
双射. 这也意味着对 $(\eta, \alpha)$ 表示第二个成员关于 $\eta'$ 的函子 $\mathcal{F}_T^{\circ} \to$ Ens. 若对 $\mathcal{E}$ 中给定的态射 $f \colon T \to S$ 和给定的 $\xi \in \mathrm{Ob}(\mathcal{F}_S)$，存在这样的对 $(\eta, \alpha)$，i.e. 存在 $\mathcal{F}$ 中以 $\xi$ 为靶、满足 $p(\alpha) = f$ 的笛卡儿态射 $\alpha$，则 $\eta$ 在 $\mathcal{F}_T$ 中在唯一同构意义下确定. 此时称 $\xi$ 经 $f$ 的逆像 (inverse image) 存在，$\mathcal{F}_T$ 中配备 $f$-笛卡儿态射 $\alpha \colon \eta \to \xi$ 的对象 $\eta$ 称为 $\xi$ 经 $f$ 的逆像. 常假设在每次存在时选定了这样的逆像 ($\mathcal{F}$ 固定)；逆像用 $f_{\mathcal{F}}^{*}(\xi)$，或简写 $f^{*}(\xi)$，或当这些记号不引起混淆时用 $\xi \times_S T$ 表示；典范态射 $\alpha \colon \eta \to \xi$ 在下面用 $\alpha_f(\xi)$ 表示. 若对每个 $\xi \in \mathrm{Ob}(\mathcal{F}_S)$，$\xi$ 经 $f$ 的逆像存在，也说 $\mathcal{F}$ 中经 $f$ 的逆像函子存在，$f^{*}(\xi)$ 于是成为 $\xi$ 的共变函子，从 $\mathcal{F}_S$ 到 $\mathcal{F}_T$. 这来自 (i) 中第二个成员对 $\xi$ 的共变依赖，i.e. 精确地说表示 $\mathcal{F}_T^{\circ} \times \mathcal{F}_S$ 到 Ens 的函子. 对 $f^{*}(\xi)$ 的这种函子性依赖如下明确：考虑 $f$-笛卡儿态射
\[
\alpha \colon \eta \longrightarrow \xi, \quad \alpha' \colon \eta' \longrightarrow \xi'
\]
和 $S$-态射 $\lambda \colon \xi \to \xi'$，则存在且只存在一个 $T$-态射 $\mu \colon \eta \to \eta'$，使得
\[
\alpha'\mu = \lambda\alpha
\]
(由 $\alpha'$ 笛卡儿得出). 还注意如下立即事实：考虑 $\mathcal{F}$ 中的交换图式
\[
\begin{tikzcd}[row sep=large, column sep=large]
\xi \arrow[d,"\lambda"] & \eta \arrow[l,"\alpha"] \arrow[d,"\mu"] \\
\xi' & \eta' \arrow[l,"\alpha'"]
\end{tikzcd}
\]
其中 $\alpha$ 和 $\alpha'$ 是 $f$-态射，$\lambda$ 是 $S$-同构，$\mu$ 是 $T$-同构. 为使 $\alpha$ 笛卡儿，必要且充分条件是 $\alpha'$ 笛卡儿.
\end{definition}

\begin{definition}\label{def:VI.5.2}
$\mathcal{E}$-函子 $F \colon \mathcal{F} \to \mathcal{G}$ 称为\emph{笛卡儿函子} (cartesian functor)\index[term]{dikaerhanzi@笛卡儿函子}，若它把笛卡儿态射变为笛卡儿态射. 用 $\mathbf{Hom}_{\mathrm{cart}}(\mathcal{F}, \mathcal{G})$\index[notations]{$\mathbf{Hom}_{\mathrm{cart}}(\mathcal{F}, \mathcal{G})$}\index[notations]{$\mathbf{Hom}_{\mathrm{cart}}(\mathcal{F}, \mathcal{G})$} 记 $\mathbf{Hom}_{\mathcal{E}/-}(\mathcal{F}, \mathcal{G})$ 的由笛卡儿函子组成的全子范畴.
\end{definition}

例如，借助恒等函子把 $\mathcal{E}$ 看作 $\mathcal{E}$ 之上的范畴，$\mathcal{E}$ 的每个态射都是笛卡儿态射，因此 $\mathcal{E}$ 到 $\mathcal{F}$ 的笛卡儿函子是把 $\mathcal{E}$ 的每个态射变为笛卡儿态射的截面函子 $F \colon \mathcal{E} \to \mathcal{F}$；这样的函子称为 $\mathcal{F}$ 在 $\mathcal{E}$ 上的笛卡儿截面 (cartesian section).

\begin{proposition}\label{prop:VI.5.3}
(i) 是 $\mathcal{E}$-等价的函子 $F \colon \mathcal{F} \to \mathcal{G}$ 是笛卡儿函子. (ii) 设 $F$，$G$ 是两个同构的 $\mathcal{E}$-函子 $\mathcal{F} \to \mathcal{G}$. 若其中之一笛卡儿，则另一个亦然. (iii) 两个笛卡儿函子 $\mathcal{F} \to \mathcal{G}$ 和 $\mathcal{G} \to \mathcal{H}$ 的复合是笛卡儿函子.
\end{proposition}

断言 (iii) 由定义平凡，(ii) 由 \ref{def:VI.5.2} 前的评注得出，(i) 容易由定义和判则 \ref{prop:VI.4.2} (iii) 得出；更精确地说，$\mathcal{F}$ 中的态射 $\alpha$ 笛卡儿当且仅当 $F(\alpha)$ 笛卡儿.

\begin{corollary}\label{cor:VI.5.4}
设 $F \colon \mathcal{F} \to \mathcal{G}$ 是 $\mathcal{E}$-等价. 则对 $\mathcal{E}$ 之上的每个范畴 $\mathcal{H}$，相应函子 $G \mapsto G \circ F$ 和 $G \mapsto F \circ G$ 诱导范畴等价：
\[
\begin{array}{rl}
& \mathbf{Hom}_{\mathrm{cart}}(\mathcal{G}, \mathcal{H}) \xrightarrow{\approx} \mathbf{Hom}_{\mathrm{cart}}(\mathcal{F}, \mathcal{H}) \\
& \mathbf{Hom}_{\mathrm{cart}}(\mathcal{H}, \mathcal{F}) \xrightarrow{\approx} \mathbf{Hom}_{\mathrm{cart}}(\mathcal{H}, \mathcal{G})
\end{array}
\]
\end{corollary}

由 \ref{prop:VI.4.2} 判则 (i) 和 \ref{prop:VI.5.3} (i) (ii) (iii) 按通常方式推出. 可以精确化：$\mathcal{E}$-函子 $G \colon \mathcal{G} \to \mathcal{H}$ 笛卡儿当且仅当 $G \circ F$ 笛卡儿，同样 $\mathcal{E}$-函子 $G \colon \mathcal{H} \to \mathcal{F}$ 笛卡儿当且仅当 $F \circ G$ 笛卡儿.

由 \ref{cor:VI.5.4} (iii) 得出，若考虑 $\mathbf{Cat}_{/\mathcal{E}}$ 的子范畴 $\mathbf{Cat}_{/\mathcal{E}}^{\mathrm{cart}}$ (其对象与 $\mathbf{Cat}_{/\mathcal{E}}$ 的相同，态射是笛卡儿函子)，则如第 2 号那样有配对：
\[
\operatorname{Hom}_{\mathrm{cart}}(\mathcal{F}, \mathcal{G}) \times \operatorname{Hom}_{\mathrm{cart}}(\mathcal{G}, \mathcal{H}) \longrightarrow \operatorname{Hom}_{\mathrm{cart}}(\mathcal{F}, \mathcal{H})
\]
由第 2 号的配对诱导，使 $\mathbf{Hom}_{\mathrm{cart}}(\mathcal{F}, \mathcal{G})$ 可以看作 $\mathcal{F}, \mathcal{G}$ 的函子，从范畴 $(\mathbf{Cat}_{/\mathcal{E}}^{\mathrm{cart}})^{\circ} \times \mathbf{Cat}_{/\mathcal{E}}^{\mathrm{cart}}$ 到 Cat. 我们主要对 $\mathcal{F} = \mathcal{G}$ 的情形需要这个评注.

\begin{definition}\label{def:VI.5.5}
设 $\mathcal{F}$ 是 $\mathcal{E}$ 之上的范畴. 用
\[
\varprojlim \mathcal{F}/\mathcal{E}\index[notations]{$\varprojlim \mathcal{F}/\mathcal{E}$}
\]
记 $\mathcal{E}$ 到 $\mathcal{F}$ 的笛卡儿 $\mathcal{E}$-函子的范畴，i.e. $\mathcal{F}$ 在 $\mathcal{E}$ 上的笛卡儿截面的范畴.
\end{definition}

由前面所述，Lim $\mathcal{F}/\mathcal{E}$ 是 $\mathcal{F}$ 的函子，从范畴 Cat$^{\mathrm{cart}}_{/\mathcal{E}}$ 到 Cat.

下面将看到这个运算 Lim 与范畴投射极限概念的关系，以及大量例子.

\section{纤维范畴与预纤维范畴。其中的积与基变换}

\begin{definition}\label{def:VI.6.1}
$\mathcal{E}$ 之上的范畴 $\mathcal{F}$ 称为纤维范畴 (fibered category)\index[term]{xianweifanchou@纤维范畴} (此时说函子 $\mathcal{F} \to \mathcal{E}$ 是纤维的 (fibrant))，若它满足下列两个公理：
\begin{itemize}\tightlist
\item[Fib$_1$] 对 $\mathcal{E}$ 中的每个态射 $f \colon T \to S$，$\mathcal{F}$ 中经 $f$ 的逆像函子存在.
\item[Fib$_{II}$] 两个笛卡儿态射的复合是笛卡儿态射.
\end{itemize}
$\mathcal{E}$ 之上满足条件 Fib$_1$ 的范畴 $\mathcal{F}$ 称为 $\mathcal{E}$ 上的\emph{预纤维范畴} (prefibered category)\index[term]{yuxianweifanchou@预纤维范畴}.

若 $\mathcal{F}$ 是 $\mathcal{E}$ 上的纤维范畴 (相应地：预纤维范畴)，$\mathcal{F}$ 的子范畴 $\mathcal{G}$ 称为纤维子范畴 (相应地：预纤维子范畴)，若它是 $\mathcal{E}$ 上的纤维范畴 (相应地：预纤维范畴)，且包含函子笛卡儿. 例如若 $\mathcal{G}$ 是 $\mathcal{F}$ 的全子范畴，看到这意味着：对 $\mathcal{E}$ 中的每个态射 $f \colon T \to S$ 和每个 $\xi \in \mathrm{Ob}(\mathcal{G}_S)$，$f_{\mathcal{F}}^{*}(\xi)$ 与 $\mathcal{G}_T$ 的某个对象 $T$-同构. 另一个有趣的情形是：$\mathcal{F}$ 是 $\mathcal{E}$ 上的纤维范畴，考虑 $\mathcal{F}$ 的子范畴 $\mathcal{G}$，其对象相同，态射是 $\mathcal{F}$ 的笛卡儿态射；特别 $\mathcal{G}_S$ 的态射是 $\mathcal{F}_S$ 的同构. 立即看到这确实是 $\mathcal{F}$ 的纤维子范畴，因为在与 $\mathcal{F}$ 中 $f$-笛卡儿态射 $\alpha$ 相关的双射
\[
\operatorname{Hom}_T(\eta', \eta) \xrightarrow{\sim} \operatorname{Hom}_f(\eta', \xi)
\]
中，与第一成员的 $T$-同构对应的是第二成员的笛卡儿态射. 按定义，$\mathcal{E}$ 到 $\mathcal{F}$ 的笛卡儿截面于是与 $\mathcal{E}$ 到 $\mathcal{G}$ 的任意 $\mathcal{E}$-函子一一对应 (但注意自然函子
\[
\mathbf{Hom}_{\mathcal{E}/-}(\mathcal{E}, \mathcal{G}) \longrightarrow \mathbf{Hom}_{\mathrm{cart}}(\mathcal{E}, \mathcal{F}) = \underline{\mathrm{Lim}}(\mathcal{F}/\mathcal{E})
\]
忠实，但一般不严格全忠实，i.e. 不是同构).
\end{definition}

\begin{remarkstar}
设 $\mathcal{F}$ 是 $\mathcal{E}$ 之上的范畴. 下列条件等价：(i) $\mathcal{F}$ 的所有态射笛卡儿；(ii) $\mathcal{F}$ 是 $\mathcal{E}$ 上的纤维范畴，且诸 $\mathcal{F}_S$ 是群胚 (groupoid)\index[term]{qunpei@群胚} (i.e. $\mathcal{F}_S$ 中每个态射是同构). 此时称 $\mathcal{F}$ 是 $\mathcal{E}$ 上的\emph{群胚纤维范畴} (fibered category in groupoids)\index[term]{qunpeixianweifanchou@群胚纤维范畴}. 在"模理论"中遇到的正是它们. 若 $\mathcal{E}$ 是群胚，证明条件 (i) 和 (ii) 还等价于：(iii) $\mathcal{F}$ 是群胚，且投影函子 $p \colon \mathcal{F} \to \mathcal{E}$ 可运输 (参见 \ref{cor:VI.4.4}). 例如，若 $\mathcal{E}$ 和 $\mathcal{F}$ 是使 $\mathrm{Ob}(\mathcal{E})$ 和 $\mathrm{Ob}(\mathcal{F})$ 约化为一点的群胚，从而 $\mathcal{E}$ 和 $\mathcal{F}$ 在同构意义下由群 $E$ 和 $F$ 定义，且函子 $p \colon \mathcal{F} \to \mathcal{E}$ 由群同态 $p \colon F \to E$ 定义，则 $\mathcal{F}$ 在 $\mathcal{E}$ 上纤维当且仅当 $p$ 满射，i.e. 当且仅当 $p$ 定义群 $E$ 以群 $G = \operatorname{Ker} p$ 为核的扩张.
\end{remarkstar}

\begin{proposition}\label{prop:VI.6.2}
设 $F \colon \mathcal{F} \to \mathcal{G}$ 是 $\mathcal{E}$-等价. 为使 $\mathcal{F}$ 是 $\mathcal{E}$ 上的纤维范畴 (相应地：预纤维范畴)，必要且充分条件是 $\mathcal{G}$ 如此.
\end{proposition}

由定义和上面指出的评注 (即 $\mathcal{F}$ 中的态射 $\alpha$ 笛卡儿当且仅当 $F(\alpha)$ 笛卡儿) 容易得出.

\begin{proposition}\label{prop:VI.6.3}
设 $\mathcal{F}_1$，$\mathcal{F}_2$ 是 $\mathcal{E}$ 之上的两个范畴，$\alpha = (\alpha_1, \alpha_2)$ 是 $\mathcal{F} = \mathcal{F}_1 \times_{\mathcal{E}} \mathcal{F}_2$ 中的态射. 为使 $\alpha$ 笛卡儿，必要且充分条件是它的两个分量笛卡儿.
\end{proposition}

确实设 $\xi_i$ 是 $\alpha_i$ 的靶，$\eta_i$ 是源，$f \colon T \to S$ 是使 $\alpha_1$ 和 $\alpha_2$ 成为 $f$-态射的 $\mathcal{E}$ 的态射. 对每个 $\eta' = (\eta'_1, \eta'_2) \in \mathcal{F}_T$，有交换图式 (其中垂直箭是双射). 故若两个水平箭之一是双射，则另一个亦然. 这已表明：若 $\alpha_1$，$\alpha_2$ 笛卡儿 (从而第二个水平箭双射)，则 $\alpha$ 笛卡儿. 反过来，在上图式中取 $\eta'_i = \eta_i$，由此得 Hom$_T(\eta'_i, \eta_i) \neq \emptyset$，先对 $i = 2$，这证明 $\alpha_1$ 笛卡儿，再对 $i = 1$，证明 $\alpha_2$ 笛卡儿.

\begin{corollary}\label{cor:VI.6.4}
设 $\mathcal{F} = \mathcal{F}_1 \times_{\mathcal{E}} \mathcal{F}_2$，$F = (F_1, F_2)$ 是 $\mathcal{E}$-函子 $\mathcal{G} \to \mathcal{F}$. 为使 $F$ 笛卡儿，必要且充分条件是 $F_1$ 和 $F_2$ 笛卡儿. 由此得到范畴同构
\[
\mathbf{Hom}_{\mathrm{cart}}(\mathcal{G}, \mathcal{F}_1 \times_{\mathcal{E}} \mathcal{F}_2) \xrightarrow{\sim} \mathbf{Hom}_{\mathrm{cart}}(\mathcal{G}, \mathcal{F}_1) \times \mathbf{Hom}_{\mathrm{cart}}(\mathcal{G}, \mathcal{F}_2)
\]
特别 (取 $\mathcal{G} = \mathcal{E}$) 得到范畴同构
\[
\varprojlim (\mathcal{F}_1 \times_{\mathcal{E}} \mathcal{F}_2/\mathcal{E}) \xrightarrow{\sim} \varprojlim (\mathcal{F}_1/\mathcal{E}) \times \varprojlim (\mathcal{F}_2/\mathcal{E})
\]
\end{corollary}

\begin{corollary}\label{cor:VI.6.5}
设 $\mathcal{F}_1$ 和 $\mathcal{F}_2$ 是 $\mathcal{E}$ 上的两个纤维范畴 (相应地：预纤维范畴)，则它们的纤维积 $\mathcal{F} = \mathcal{F}_1 \times_{\mathcal{E}} \mathcal{F}_2$ 是 $\mathcal{E}$ 上的纤维范畴 (相应地：预纤维范畴).
\end{corollary}

这些结果还推广到 $\mathcal{E}$ 之上任意族范畴的纤维积的情形.

\begin{proposition}\label{prop:VI.6.6}
设 $\mathcal{F}$ 是 $\mathcal{E}$ 之上的范畴，投影函子为 $p$，$\lambda \colon \mathcal{E}' \to \mathcal{E}$ 是函子，借助投影函子 $p' = p \times_{\mathcal{E}} \mathrm{id}_{\mathcal{E}'}$ 把 $\mathcal{F}' = \mathcal{F} \times_{\mathcal{E}} \mathcal{E}'$ 看作 $\mathcal{E}'$ 之上的范畴. 设 $\alpha'$ 是 $\mathcal{F}'$ 的态射，为使 $\alpha'$ 是笛卡儿态射，必要且充分条件是它在 $\mathcal{F}$ 中的像 $\alpha$ 笛卡儿.
\end{proposition}

证明是直接的，留给读者.

\begin{corollary}\label{cor:VI.6.7}
对 $\mathcal{E}$ 之上范畴的每个笛卡儿函子 $F \colon \mathcal{F} \to \mathcal{G}$，函子 $F' = F \times_{\mathcal{E}} \mathcal{E}'$ (从 $\mathcal{F}' = \mathcal{F} \times_{\mathcal{E}} \mathcal{E}'$ 到 $\mathcal{G}' = \mathcal{G} \times_{\mathcal{E}} \mathcal{E}'$) 笛卡儿.
\end{corollary}

因此，第 3 号所考虑的函子 $\mathbf{Hom}_{\mathcal{E}}(\mathcal{F}, \mathcal{G}) \to \mathbf{Hom}_{\mathcal{E}'}(\mathcal{F}', \mathcal{G}')$ 诱导函子
\[
\operatorname{Hom}_{\mathrm{cart}}(\mathcal{F}, \mathcal{G}) \longrightarrow \operatorname{Hom}_{\mathrm{cart}}(\mathcal{F}', \mathcal{G}');
\]
换言之，对固定的 $\mathcal{F}, \mathcal{G}$，可以把
\[
\mathbf{Hom}_{\mathrm{cart}}(\mathcal{F} \times_{\mathcal{E}} \mathcal{E}', \mathcal{G} \times_{\mathcal{E}} \mathcal{E}')
\]
看作 $\mathcal{E}'$ 的函子，从范畴 $\mathbf{Cat}_{/\mathcal{E}}^{\circ}$ 到 Cat. 若也让 $\mathcal{F}, \mathcal{G}$ 变动，得到从范畴 $\mathbf{Cat}_{/\mathcal{E}}^{\circ} \times (\mathbf{Cat}_{/\mathcal{E}}^{\mathrm{cart}})^{\circ} \times \mathbf{Cat}_{/\mathcal{E}}^{\mathrm{cart}}$ 到 Cat 的函子. 考虑到第 3 号所考虑的同构
\[
\operatorname{Hom}_{\mathcal{E}'}(\mathcal{F}', \mathcal{G}') \stackrel{\sim}{\longrightarrow} \operatorname{Hom}_{\mathcal{E}}(\mathcal{F} \times_{\mathcal{E}} \mathcal{E}', \mathcal{G}),
\]
则 $\mathcal{F}'$ 到 $\mathcal{G}'$ 的笛卡儿 $\mathcal{E}'$-函子与把"第一投影是 $\mathcal{F}$ 的笛卡儿态射"的每个态射变为 $\mathcal{G}$ 的笛卡儿态射的 $\mathcal{E}$-函子 $\mathcal{F} \times_{\mathcal{E}} \mathcal{E}' \to \mathcal{G}$ 对应. 取 $\mathcal{F} = \mathcal{E}$，(更换记号后) 得到：

\begin{corollary}\label{cor:VI.6.8}
Lim$(\mathcal{F}'/\mathcal{E}')$ 同构于 $\mathbf{Hom}_{\mathcal{E}/-}(\mathcal{E}', \mathcal{F})$ 的由把任意态射变为笛卡儿态射的 $\mathcal{E}$-函子 $\mathcal{E}' \to \mathcal{F}$ 组成的全子范畴. 特别，若 $\mathcal{F}$ 是纤维范畴且 $\widetilde{\mathcal{F}}$ 是以 $\mathcal{F}$ 的笛卡儿态射为态射的 $\mathcal{F}$ 的子范畴，则双射
\[
\mathrm{Ob} \varprojlim (\mathcal{F}'/\mathcal{E}') \xrightarrow{\sim} \mathrm{Hom}_{\mathcal{E}/-}(\mathcal{E}', \widetilde{\mathcal{F}}).
\]
这精确了 $\underline{\mathrm{Lim}}(\mathcal{F} \times_{\mathcal{E}} \mathcal{E}'/\mathcal{E}')$ 应被看作 $\mathcal{E}'$ 和 $\mathcal{F}$ 的函子的方式，从范畴 Cat$^{\circ}_{/\mathcal{E}} \times \mathbf{Cat}^{\mathrm{cart}}_{/\mathcal{E}}$ 到 Cat. 后面在 $\mathcal{F}$ 限制为纤维范畴时将看到对 $\mathcal{E}'$ 的更完全的函子依赖.
\end{corollary}

\begin{corollary}\label{cor:VI.6.9}
设 $\mathcal{F}$ 是 $\mathcal{E}$ 上的纤维范畴 (相应地：预纤维范畴)，则 $\mathcal{F}' = \mathcal{F} \times_{\mathcal{E}} \mathcal{E}'$ 是 $\mathcal{E}'$ 上的纤维范畴 (相应地：预纤维范畴).
\end{corollary}

\begin{proposition}\label{prop:VI.6.10}
设 $\mathcal{F}, \mathcal{G}$ 是 $\mathcal{E}$ 上的预纤维范畴，$F$ 是 $\mathcal{F}$ 到 $\mathcal{G}$ 的笛卡儿 $\mathcal{E}$-函子. 为使 $F$ 忠实 (相应地：全忠实，相应地：是 $\mathcal{E}$-等价)，必要且充分条件是：对每个 $S \in \mathrm{Ob}(\mathcal{E})$，诱导函子 $F_S \colon \mathcal{F}_S \to \mathcal{G}_S$ 忠实 (相应地：全忠实，相应地：是等价).
\end{proposition}

由定义立即证明.

为结束本节，给出纤维范畴的若干性质，使用公理 Fib$_{II}$.

\begin{proposition}\label{prop:VI.6.11}
设 $\mathcal{F}$ 是 $\mathcal{E}$ 上的预纤维范畴. 为使 $\mathcal{F}$ 纤维，必要且充分条件是它满足下列条件：
\begin{itemize}\tightlist
\item[Fib$'_{II}$：] 设 $\alpha \colon \eta \to \xi$ 是 $\mathcal{F}$ 中 $\mathcal{E}$ 的态射 $f \colon T \to S$ 之上的笛卡儿态射. 对 $\mathcal{E}$ 中的每个态射 $g \colon U \to T$ 和每个 $\zeta \in \mathrm{Ob}(\mathcal{F}_U)$，映射 $u \mapsto \alpha \circ u$
\[
\mathrm{Hom}_g(\zeta, \eta) \longrightarrow \mathrm{Hom}_{fg}(\zeta, \xi)
\]
双射.
\end{itemize}
\end{proposition}

换言之，在 $\mathcal{E}$ 上的纤维范畴中，笛卡儿图式由先验更强的性质刻画 (在下列陈述中取 $g = \mathrm{id}_T$ 即得定义中的性质).

\begin{corollary}\label{cor:VI.6.12}
设 $\mathcal{F}$ 是 $\mathcal{E}$ 之上的范畴，$\alpha$ 是 $\mathcal{F}$ 中的态射. 为使 $\alpha$ 是同构，必须 $p(\alpha) = f$ 是同构且 $\alpha$ 笛卡儿；若 $\mathcal{F}$ 在 $\mathcal{E}$ 上纤维，则逆命题成立.
\end{corollary}

事实上，若 $\alpha$ 是同构则 $f = p(\alpha)$ 显然亦然；对每个 $\eta' \in \mathrm{Ob}(\mathcal{F}_T)$，映射 $u \mapsto \alpha \circ u$
\[
\operatorname{Hom}(\eta', \eta) \longrightarrow \operatorname{Hom}(\eta', \xi)
\]
双射；由于 $f$ 是同构，立即看到第一成员的 $T$-态射当且仅当其像为第二成员的 $f$-态射，故得到双射
\[
\operatorname{Hom}_T(\eta', \eta) \longrightarrow \operatorname{Hom}_f(\eta', \xi)
\]
证明第一个断言. 反过来，假设 $f$ 是同构且 $\alpha$ 满足 Fib$'_{II}$ 中陈述的条件 (当 $\mathcal{F}$ 在 $\mathcal{E}$ 上纤维时，这意味着 $\alpha$ 笛卡儿)，立即看到对每个 $\zeta \in \mathrm{Ob} \mathcal{F}$，映射 $u \mapsto \alpha \circ u$ (从 Hom$(\zeta, \eta)$ 到 Hom$(\zeta, \xi)$) 双射，故 $\alpha$ 是同构.

\begin{corollary}\label{cor:VI.6.13}
设 $\alpha \colon \eta \to \xi$ 和 $\beta \colon \zeta \to \eta$ 是纤维范畴 $\mathcal{F}$ 中的两个可复合态射. 若 $\alpha$ 笛卡儿，则 $\beta$ 笛卡儿当且仅当 $\alpha\beta$ 笛卡儿.
\end{corollary}

用 \ref{prop:VI.6.11} 的强化形式的笛卡儿态射定义.

\section{\texorpdfstring{$\mathcal{E}$}{E} 上的劈开范畴}

\begin{definition}\label{def:VI.7.1}
设 $\mathcal{F}$ 是 $\mathcal{E}$ 之上的范畴. 称 $\mathcal{F}$ 在 $\mathcal{E}$ 上的\emph{劈开} (cleavage)\index[term]{pikai@劈开} 为给每个 $f \in \mathrm{Fl}(\mathcal{E})$ 配 $\mathcal{F}$ 中经 $f$ 的逆像函子 $f^{*}$\index[notations]{$f^{*}$} 的函数. 若 $f = \mathrm{id}_S$ 蕴含 $f^{*} = \mathrm{id}_{\mathcal{F}_S}$，则称该劈开规范化 (normalized)\index[term]{guifanhuapikai@规范化劈开}. 配备劈开 (相应地：\emph{规范化劈开}) 的 $\mathcal{E}$ 之上的范畴 $\mathcal{F}$ 称为\emph{劈开范畴} (cleaved category)\index[term]{pikaifanchou@劈开范畴} (相应地：规范化劈开范畴).
\end{definition}

显然，$\mathcal{F}$ 有劈开当且仅当 $\mathcal{F}$ 是 $\mathcal{E}$ 上的预纤维范畴，且此时 $\mathcal{F}$ 有规范化劈开. $\mathcal{F}$ 上的劈开集与 $\mathrm{Fl}(\mathcal{F})$ 的满足下列条件的子集 $K$ 一一对应：
\begin{enumerate}[a)]
\item $K$ 的诸 $\alpha$ 是笛卡儿态射.
\item 对 $\mathcal{E}$ 中的每个态射 $f \colon T \to S$ 和每个 $\xi \in \mathrm{Ob}(\mathcal{F}_S)$，存在 $K$ 中以 $\xi$ 为靶的唯一 $f$-态射.
\end{enumerate}
为使 $K$ 定义的劈开规范化，必要且充分条件是 $K$ 还满足条件
\begin{enumerate}[a)]
\setcounter{enumi}{2}
\item $\mathcal{F}$ 中的恒等态射属于 $K$.
\end{enumerate}
$K$ 的态射可称为所考虑劈开的"\emph{运输态射}" (transport morphisms)\index[term]{yunshutaishe@运输态射}.

$\mathcal{E}$ 上劈开范畴的同构概念是清楚的. 更一般地，可以把 $\mathcal{E}$-劈开范畴的态射定义为把运输态射变为运输态射的 $\mathcal{E}$-范畴函子 $\mathcal{F} \to \mathcal{G}$. (它们特别是笛卡儿函子). 这样，$\mathcal{E}$ 上的劈开范畴是某个范畴的对象，即 $\mathcal{E}$ 上的劈开范畴范畴. 读者将明确积的存在性，它与下列事实有关：若 $\mathcal{E}$ 之上的范畴 $\mathcal{F}$ 是配备各自劈开的 $\mathcal{E}$ 之上范畴 $\mathcal{F}_i$ 的积，则 $\mathcal{F}$ 配备相应的自然劈开. 也把劈开范畴中的基变换概念留给读者明确.

用 $\alpha_f(\xi)$ 记典范态射\index[notations]{$\alpha_f(\xi)$}
\[
\alpha_f(\xi) \colon f^{*}(\xi) \longrightarrow \xi.
\]
已说过，它对 $\xi$ 函子性，i.e. 有函子同态
\[
\alpha_f \colon i_T f^{*} \longrightarrow i_S,
\]
其中对每个 $S \in \mathrm{Ob}(\mathcal{E})$，$i_S$ 记包含函子
\[
i_S \colon \mathcal{F}_S \longrightarrow \mathcal{F}.
\]
现在考虑 $\mathcal{E}$ 中的态射
\[
f \colon T \longrightarrow S \quad \text{和} \quad g \colon U \longrightarrow T,
\]
设 $\xi \in \mathrm{Ob}(\mathcal{F}_S)$，则存在唯一 $U$-态射
\[
c_{f,g}(\xi) \colon g^{*} f^{*}(\xi) \longrightarrow (fg)^{*}(\xi)
\]
使图式交换 (由 $(fg)^{*}(\xi)$ 的定义)：
\begin{equation}
\begin{tikzcd}[row sep=large, column sep=large]
f^{*}(\xi) \arrow[d,"{\alpha_f(\xi)}"] & g^{*}(f^{*}(\xi)) \arrow[l,"{\alpha_g(f^{*}(\xi))}"] \arrow[d,"{c_{f,g}(\xi)}"] \\
\xi & (fg)^{*}(\xi) \arrow[l,"{\alpha_{fg}(\xi)}"]
\end{tikzcd} \tag{172}
\label{eq:VI.172}
\end{equation}
对变动的 $\xi$，这个同态函子性，i.e. 有函子 $\mathcal{F}_S \to \mathcal{F}_U$ 的同态
\[
c_{f,g} \colon g^{*} f^{*} \longrightarrow (fg)^{*}
\]
立即注意：

\begin{proposition}\label{prop:VI.7.2}
为使 $\mathcal{E}$ 上的劈开范畴 $\mathcal{F}$ 纤维，必要且充分条件是诸 $c_{f,g}$ 是同构.
\end{proposition}

取 $f$ 为同构、$g$ 为其逆，并考虑同构 $c_{f,g}$ 和 $c_{g,f}$，由此得出：

\begin{corollary}\label{cor:VI.7.3}
若 $\mathcal{F}$ 是 $\mathcal{E}$ 上的纤维劈开范畴，则对 $\mathcal{E}$ 中的每个同构 $f \colon T \to S$，$f^{*}$ 是范畴等价 $\mathcal{F}_S \to \mathcal{F}_T$.
\end{corollary}

\begin{proposition}\label{prop:VI.7.4}
设 $\mathcal{F}$ 是 $\mathcal{E}$ 上的劈开范畴. 有
\[
\begin{array}{ll}
A) & \left\{ \begin{array}{l}
c_{f,\mathrm{id}_T}(\xi) = \alpha_{\mathrm{id}_T}(f^{*}(\xi)) \\
c_{\mathrm{id}_S,f}(\xi) = f^{*}(\alpha_{\mathrm{id}_S}(\xi))
\end{array} \right. \\
B) & c_{f,gh}(\xi) \cdot c_{g,h}(f^{*}(\xi)) = c_{fg,h}(\xi) \cdot h^{*}(c_{f,g}(\xi))
\end{array}
\]
(在这些公式中，$f, g, h$ 表示态射
\begin{equation}
V \longrightarrow U \longrightarrow T \longrightarrow S \tag{173}
\label{eq:VI.173}
\end{equation}
$\xi$ 是 $\mathcal{F}_S$ 的对象.)
\end{proposition}

在规范化劈开的情形，第一和第二个关系取更简单的形式
\[
A') \qquad c_{f,\mathrm{id}_T} = \mathrm{id}_{f^{*}}, \quad c_{\mathrm{id}_S,f} = \mathrm{id}_{f^{*}}.
\]
至于第三个，它由图式的交换性可视化
\begin{equation}
\begin{tikzcd}[row sep=large, column sep=large]
h^{*} g^{*} f^{*}(\xi) \arrow[r,"{c_{g,h}(f^{*}(\xi))}"] \arrow[d,"{h^{*}(c_{f,g}(\xi))}"] & (gh)^{*}(f^{*}(\xi)) \arrow[d,"{c_{f,gh}(\xi)}"] \\
h^{*}(fg)^{*}(\xi) \arrow[r,"{c_{fg,h}(\xi)}"] & (fgh)^{*}(\xi)
\end{tikzcd} \tag{D}
\label{eq:VI.D}
\end{equation}
在纤维范畴的情形 (诸 $c_{f,g}$ 是同构)，这个交换性可以直观地表达为：依次使用形如 $c_{f,g}$ 的同构不导致"矛盾的等同". 还可以不写自变量 $\xi$ 地写这个公式，用函子同态的卷积积：
\[
c_{fg,h} \circ (h^{*} * c_{f,g}) = c_{f,gh} \circ (c_{g,h} * f^{*}).
\]
\ref{prop:VI.7.4} 前两个公式的证明平凡，略述第三个的证明. 为此，除方形 (D) 外，考虑同态方形
\begin{equation}
\begin{tikzcd}[row sep=large, column sep=large]
g^{*} f^{*}(\xi) \arrow[r,"{\alpha_g(f^{*}(\xi))}"] \arrow[d,"{c_{f,g}(\xi)}"] & f^{*}(\xi) \arrow[d,"{\alpha_f(\xi)}"] \\
(fg)^{*}(\xi) \arrow[r,"{\alpha_{fg}(\xi)}"] & \xi
\end{tikzcd} \tag{(D')}
\end{equation}
它按 $c_{f,g}(\xi)$ 的定义交换. 考虑把 (D) 的顶点与 (D') 的相应顶点用形如 $\alpha$ 的同态连接得到的图式：
\[
\begin{array}{l}
\alpha_h(g^{*} f^{*}(\xi)), \alpha_{gh}(f^{*}(\xi)) \\
\alpha_h((fg)^{*}(\xi)), \alpha_{fgh}(\xi).
\end{array}
\]
如此得到的立方体的四个侧面也交换：对左侧面，来自 (D) 的左列由 (D') 的左列经 $h$ 作用得到，且 $\alpha_h$ 是函子同态；对其余三个，不过是 (D) 其余三边的运算 $c$ 的定义. 故立方体中除顶面外的五个面交换. 由此得出由 (D) 定义的两个 $(fgh)$-态射 $h^{*} g^{*} f^{*}(\xi) \to (fgh)^{*}(\xi)$ 与 $\alpha_{fgh}(\xi) \colon (fgh)^{*}(\xi) \to \xi$ 的复合相同，故由 $(fgh)^{*}$ 的定义它们相等.

下面限于 $\mathcal{E}$ 上的规范化劈开范畴. 这样的范畴产生下列对象：
\begin{enumerate}[a)]
\item 从 $\mathrm{Ob}(\mathcal{E})$ 到 Cat 的映射 $S \mapsto \mathcal{F}_S$.
\item 给每个 $f \in \mathrm{Fl}(\mathcal{E})$ (源为 $T$，靶为 $S$) 配函子 $f^{*} \colon \mathcal{F}_S \to \mathcal{F}_T$ 的映射 $f \mapsto f^{*}$.
\item 给 $\mathcal{E}$ 的每对箭 $(f, g)$ 配函子同态 $c_{f,g} \colon g^{*} f^{*} \to (fg)^{*}$ 的映射 $(f, g) \mapsto c_{f,g}$.
\end{enumerate}
此外，这些数据满足上面给出的公式 $A')$ 和 B) 所表达的条件. (N.B. 若不限于规范化劈开的情形，就必须引入额外对象，即给 $\mathcal{E}$ 的每个对象 $S$ 配函子同态 $\alpha_S \colon (\mathrm{id}_S)^{*} \to \mathrm{id}_{\mathcal{F}_S}$ 的函数 $S \mapsto \alpha_S$；条件 $A')$ 将被条件 A) 代替.)

下面说明如何用前述对象 (在唯一同构意义下) 重建 $\mathcal{E}$ 上的规范化劈开范畴 $\mathcal{F}$.

\section{由伪函子 \texorpdfstring{$\mathcal{E}^{\circ} \to \mathbf{Cat}$}{E^o→Cat} 定义的劈开范畴}

为简短，称从 $\mathcal{E}^{\circ}$ 到 Cat 的\emph{伪函子} (pseudo-functor)\index[term]{weihanzi@伪函子} (应当说：规范化伪函子) 为像上面 a)，b)，c) 那样、满足条件 $A')$ 和 B) 的数据集. 在前一节，我们把 $\mathcal{E}$ 上的规范化劈开范畴与伪函子 $\mathcal{E}^{\circ} \to \mathbf{Cat}$ 对应；这里指出相反的构造. 大多数细节、以及这些构造确实"互逆"的验证留给读者. 精确地说，应当把伪函子 $\mathcal{E}^{\circ} \to$ Cat 看作某个新范畴的对象，并证明我们的构造在它和前一节定义的 $\mathcal{E}$ 之上劈开范畴范畴之间提供互为拟逆的等价.

令
\[
\mathcal{F}_{\circ} = \coprod_{S \in \operatorname{Ob}(\mathcal{E})} \operatorname{Ob}(\mathcal{F}(S)),
\]
集合 $\mathrm{Ob}(\mathcal{F}(S))$ 的和集 (N.B. 这里用 $\mathcal{F}(S)$ 而不用 $\mathcal{F}_S$ 记\index[notations]{$\mathcal{F}_S$}给定伪函子在 $\mathcal{E}$ 的对象 $S$ 处的值，以避免后面的记号混淆). 故有显然映射：
\[
p_{\circ} \colon \mathcal{F}_{\circ} \longrightarrow \mathrm{Ob} \mathcal{E}.
\]
设
\[
\overline{\xi} = (S, \xi), \quad \overline{\eta} = (T, \eta) \qquad (\text{其中 } \xi \in \operatorname{Ob} \mathcal{F}(S), \eta \in \operatorname{Ob} \mathcal{F}(T))
\]
是 $\mathcal{F}_{\circ}$ 的两个元素，$f \in \operatorname{Hom}(T, S)$，令
\[
h_f(\overline{\eta}, \overline{\xi}) = \operatorname{Hom}_{\mathcal{F}(T)}(\eta, f^{*}(\xi)).
\]
若还有 $\mathcal{E}$ 中的态射 $g \colon U \to T$ 和 $\zeta \in \operatorname{Ob} \mathcal{F}(U)$，定义映射，记为 $(u, v) \mapsto u \circ v$
\[
h_f(\overline{\eta}, \overline{\xi}) \times h_g(\overline{\zeta}, \overline{\eta}) \longrightarrow h_{fg}(\overline{\zeta}, \overline{\xi}),
\]
i.e. 映射
\[
\operatorname{Hom}_{\mathcal{F}(T)}(\eta, f^{*}(\xi)) \times \operatorname{Hom}_{\mathcal{F}(U)}(\zeta, g^{*}(\eta)) \longrightarrow \operatorname{Hom}_{\mathcal{F}(U)}(\zeta, (fg)^{*}(\xi)),
\]
由公式
\[
u \circ v = c_{f,g}(\xi) \cdot g^{*}(u) \cdot v,
\]
i.e. $u \circ v$ 是序列
\[
\zeta \xrightarrow{v} g^{*}(\eta) \xrightarrow{g^{*}(u)} g^{*} f^{*}(\xi) \xrightarrow{c_{f,g}(\xi)} (fg)^{*}(\xi)
\]
的复合. 另一方面令
\[
h(\overline{\eta}, \overline{\xi}) = \coprod_{f \in \operatorname{Hom}(T, S)} h_f(\overline{\eta}, \overline{\xi}),
\]
前述配对定义配对
\[
h(\overline{\eta}, \overline{\xi}) \times h(\overline{\zeta}, \overline{\eta}) \longrightarrow h(\overline{\zeta}, \overline{\xi}),
\]
而 $h(\overline{\eta}, \overline{\xi})$ 的定义蕴含显然映射：
\[
p_{\overline{\eta},\overline{\xi}} \colon h(\overline{\eta}, \overline{\xi}) \longrightarrow \operatorname{Hom}(T, S).
\]
这样说，验证下列各点：
\begin{enumerate}[1)]
\item $h(\overline{\eta}, \overline{\xi})$ 的元素之间的复合结合.
\item 对 $\mathcal{F}_{\circ}$ 中的每个 $\overline{\xi} = (\xi, S)$，考虑 $h_{\mathrm{id}_S}(\overline{\xi}, \overline{\xi}) = \operatorname{Hom}_{\mathcal{F}(S)}(\mathrm{id}_S^{*}(\xi), \xi) = \operatorname{Hom}_{\mathcal{F}(S)}(\xi, \xi)$ 的恒等元素，及其在 $h(\overline{\xi}, \overline{\xi})$ 中的像. 该对象是 $h(\overline{\eta}, \overline{\xi})$ 元素之间复合的左幺元和右幺元.
\end{enumerate}
这已表明得到范畴 $\mathcal{F}$，令
\[
\mathrm{Ob} \mathcal{F} = \mathcal{F}_{\circ}, \quad \mathrm{Fl} \mathcal{F} = \coprod_{\overline{\xi}, \overline{\eta} \in \mathcal{F}_{\circ}} h(\overline{\eta}, \overline{\xi}).
\]
(N.B. 不能简单取 $\mathrm{Fl} \mathcal{F}$ 为诸集合 $h(\overline{\eta}, \overline{\xi})$ 的并，因为它们不必不交). 此外：
\begin{enumerate}[1)]
\setcounter{enumi}{2}
\item 映射 $p_{\circ} \colon \mathrm{Ob} \mathcal{F} \to \mathrm{Ob} \mathcal{E}$ 和 $p_1 = (p_{\overline{\eta},\overline{\xi}}) \colon \mathrm{Fl} \mathcal{F} \to \mathrm{Fl} \mathcal{E}$ 定义函子 $p \colon \mathcal{F} \to \mathcal{E}$. 这样，$\mathcal{F}$ 成为 $\mathcal{E}$ 之上的范畴，且显然映射 $h_f(\overline{\eta}, \overline{\xi}) \to \mathrm{Hom}(\overline{\eta}, \overline{\xi})$ 诱导双射
\end{enumerate}
\[
h_f(\overline{\eta}, \overline{\xi}) \stackrel{\sim}{\longrightarrow} \operatorname{Hom}_f(\overline{\eta}, \overline{\xi}).
\]
\begin{enumerate}[1)]
\setcounter{enumi}{3}
\item 显然映射
\end{enumerate}
\[
\operatorname{Ob} \mathcal{F}(S) \longrightarrow \mathcal{F}_{\circ} = \operatorname{Ob} \mathcal{F}, \quad \operatorname{Fl} \mathcal{F}(S) \longrightarrow \operatorname{Fl} \mathcal{F},
\]
(其中第二个由显然映射
\[
\mathrm{Hom}_{\mathcal{F}(S)}(\xi, \xi') = h_{\mathrm{id}_S}(\overline{\xi}, \overline{\xi}') \longrightarrow \mathrm{Hom}(\overline{\xi}, \overline{\xi}')
\]
定义) 定义同构
\[
i_S \colon \mathcal{F}(S) \stackrel{\sim}{\longrightarrow} \mathcal{F}_S.
\]
\begin{enumerate}[1)]
\setcounter{enumi}{4}
\item 对 $\mathcal{F}$ 的每个对象 $\overline{\xi} = (S, \xi)$ 和 $\mathcal{E}$ 的每个态射 $f \colon T \to S$，考虑 $\mathcal{F}_T$ 的元素 $\overline{\eta} = (T, \eta)$ (其中 $\eta = f^{*}(\xi)$) 和 $\mathrm{Hom}(\overline{\eta}, \overline{\xi})$ 的元素 $\alpha_f(\xi)$，它是 $\mathrm{id}_{f^{*}(\xi)}$ 在映射 Hom$_{\mathcal{F}(T)}(f^{*}(\xi), f^{*}(\xi)) = h_f(\overline{\eta}, \overline{\xi}) \to \mathrm{Hom}_f(\overline{\eta}, \overline{\xi})$ 下的像. 该元素笛卡儿，且 $f = \mathrm{id}_S$ 时是 $\overline{\xi}$ 的恒等态射；换言之，诸 $\alpha_f(\xi)$ 定义 $\mathcal{F}$ 在 $\mathcal{E}$ 上的规范化劈开. 此外，由构造，在函子图式中交换
\end{enumerate}
\[
\begin{tikzcd}[row sep=large, column sep=large]
\mathcal{F}(T) \arrow[r,"{f^{*}_{\mathcal{F}}}"] \arrow[d,"{i_T}"] & \mathcal{F}(S) \arrow[d,"{i_S}"] \\
\mathcal{F}_T \arrow[r,"{f^{*}_{\mathcal{F}}}"] & \mathcal{F}_S
\end{tikzcd}
\]
其中 $f^{*}_{\mathcal{F}}$ 是所考虑的劈开在 $\mathcal{F}$ 上相应的经 $f$ 的逆像函子. 最后：6) 与伪函子一起给出的同态 $c_{f,g}$ 经同构 $i_S$ 变为与 $\mathcal{F}$ 的劈开相关的函子同态 $c_{f,g}$.

只给出 1) 的验证 (如果可能，它比其他验证较不平凡). 只需证明形如 $h_f(\overline{\eta}, \overline{\xi})$ 的集合的元素之间复合的结合性. 因此考虑 $\mathcal{E}$ 中的态射
\[
S \xleftarrow{f} T \xleftarrow{g} U \xleftarrow{h} V
\]
和对象
\[
\xi, \eta, \zeta, \tau
\]
分别在 $\mathcal{F}(S), \mathcal{F}(T), \mathcal{F}(U), \mathcal{F}(V)$ 中，最后考虑元素
\[
\begin{array}{rl}
& u \in h_f(\overline{\eta}, \overline{\xi}) = \mathrm{Hom}_{\mathcal{F}(T)}(\eta, f^{*}(\xi)) \\
& v \in h_g(\overline{\zeta}, \overline{\eta}) = \mathrm{Hom}_{\mathcal{F}(U)}(\zeta, g^{*}(\eta)) \\
& w \in h_h(\overline{\tau}, \overline{\zeta}) = \mathrm{Hom}_{\mathcal{F}(V)}(\tau, h^{*}(\zeta)).
\end{array}
\]
要证明公式
\[
(u \circ v) \circ w = u \circ (v \circ w),
\]
它是 $\mathrm{Hom}_{\mathcal{F}(V)}(\tau, (fgh)^{*}(\xi))$ 中的等式. 按定义，该等式的两个成员按下面图式的上轮廓和下轮廓复合得到：
\[
\tau \xrightarrow{w} h^{*}(\zeta) \xrightarrow{h^{*}(v)} h^{*} g^{*}(\eta) \xrightarrow{h^{*} g^{*}(u)} h^{*} g^{*} f^{*}(\xi) \xrightarrow{h^{*}(c_{f,g}(\xi))} h^{*}(fg)^{*}(\xi) \xrightarrow{c_{f,gh}(\xi)} (fgh)^{*}(\xi)
\]
与
\[
\tau \xrightarrow{w} h^{*}(\zeta) \xrightarrow{c_{g,h}(\eta)} (gh)^{*}(\eta) \xrightarrow{(gh)^{*}(u)} (gh)^{*} f^{*}(\xi) \xrightarrow{c_{f,gh}(\xi)} (fgh)^{*}(\xi)
\]
(中间方形交换是因为 $c_{g,h}$ 是函子同态，右边方形按伪函子的条件 B) 交换). 由此得所述结果.

当然，还需在伪函子已来自 $\mathcal{E}$ 上的规范化劈开范畴 $\mathcal{F}'$ 的情形明确如何得到 $\mathcal{F}'$ 与 $\mathcal{F}$ 之间的自然同构. 细节留给读者.

也把劈开范畴 $\mathcal{F}$ (在 $\mathcal{E}$ 上) 经基变换函子 $\mathcal{E}' \to \mathcal{E}$ 的逆像概念用伪函子解释留给读者.

\section{例：由函子 \texorpdfstring{$\mathcal{E}^{\circ} \to \mathbf{Cat}$}{E^o→Cat} 定义的劈开范畴；\texorpdfstring{$\mathcal{E}$}{E} 上的分裂范畴}

假设有函子
\[
\varphi \colon \mathcal{E}^{\circ} \longrightarrow \mathbf{Cat},
\]
它定义伪函子，令
\[
\mathcal{F}(S) = \varphi(S), \quad f^{*} = \varphi(f), \quad c_{f,g} = \mathrm{id}_{(fg)^{*}}.
\]
因此前一节的构造给出 $\mathcal{E}$ 上的劈开范畴 $\mathcal{F}$，称为与函子 $\varphi$ 相关的劈开范畴. 为使 $\mathcal{E}$ 上的劈开范畴同构于由函子 $\varphi \colon \mathcal{E}^{\circ} \to$ Cat 定义的劈开范畴，显然必要且充分条件是它满足条件：
\[
(fg)^{*} = g^{*} f^{*}, \quad c_{f,g} = \mathrm{id}_{(fg)^{*}}.
\]
用运输态射集 $K$ 的术语，这也简单意味着两个运输态射的复合是运输态射. $\mathcal{E}$ 上的范畴 $\mathcal{F}$ 的满足前述条件的劈开称为 $\mathcal{F}$ 在 $\mathcal{E}$ 上的\emph{分裂} (splitting)\index[term]{fenlie@分裂}，配备分裂的 $\mathcal{E}$ 之上的范畴 $\mathcal{F}$ 称为 $\mathcal{E}$ 上的\emph{分裂范畴} (split category)\index[term]{fenliefanchou@分裂范畴}. 因此它是劈开范畴概念的特殊情形. $\mathcal{E}$ 上的分裂范畴范畴于是等价于 Hom$(\mathcal{E}^{\circ}, \mathrm{Cat})$. 注意 $\mathcal{E}$ 上的分裂范畴更是 $\mathcal{E}$ 上的劈开范畴.

若 $\mathcal{F}$ 是 $\mathcal{E}$ 上的纤维范畴，$\mathcal{F}$ 上不一定存在分裂. 例如假设 $\mathrm{Ob} \mathcal{E}$ 和 $\mathrm{Ob} \mathcal{F}$ 约化为一个元素，且该对象的自同态集是群 $E$ (相应地：$F$)，从而投影函子 $p$ 由群同态 $p \colon F \to E$ 给出，由于 $p$ 纤维故满射. 立即验证 $\mathcal{F}$ 在 $\mathcal{E}$ 上的劈开集与映射 $s \colon E \to F$ (满足 $ps = \mathrm{id}_E$，i.e. 满射同态 $p \colon F \to E$ 的核 $G$ 的陪集类的"代表元系") 的集合一一对应. 劈开是分裂当且仅当 $s$ 是群同态. 说存在分裂，意味着 $E$ 以 $G$ 为核的群扩张 $F$ 平凡，当 $G$ 交换时这由 $H^2(E, G)$ 中某个上同调类的消没表达 (其中 $G$ 看作 $E$ 作用其上的群).

然而假设 $\mathcal{F}$ 是 $\mathcal{E}$ 上的纤维范畴，使诸 $\mathcal{F}_S$ 是\emph{刚性范畴} (rigid category)\index[term]{gangxingfanchou@刚性范畴}，i.e. $\mathcal{F}_S$ 的每个对象的自同构群约化为恒等. 则容易证明 $\mathcal{F}$ 在 $\mathcal{E}$ 上有分裂. 事实上，先看到分裂的存在性问题不受把 $\mathcal{F}$ 换成 $\mathcal{E}$-等价范畴的影响，这把我们归结到诸 $\mathcal{F}_S$ 是刚性且约化 (reduced)\index[term]{yuehua@约化} (i.e. $\mathcal{F}_S$ 中两个同构对象相同) 范畴的情形. 但若 $G$ 是刚性约化范畴，则两个函子 $H \to G$ 的任何同构 (其中 $H$ 是任意范畴) 是恒等. 由此得出，若 $\mathcal{F}$ 是 $\mathcal{E}$ 上的纤维范畴，使纤维范畴刚性且约化，则 $\mathcal{F}$ 在 $\mathcal{E}$ 上有唯一劈开，它必然是一个分裂. 故 $\mathcal{F}$ 同构于由函子 $\varphi \colon \mathcal{E}^{\circ} \longrightarrow \mathbf{Cat}$ 定义的范畴，其中诸 $\varphi(S)$ 是刚性且离散的范畴，且函子 $\varphi$ 在唯一同构意义下确定.

\section{余纤维范畴、双纤维范畴}

考虑 $\mathcal{E}$ 之上的范畴 $\mathcal{F}$，投影函子为
\[
p \colon \mathcal{F} \longrightarrow \mathcal{E},
\]
它借助投影函子
\[
p^{\circ} \colon \mathcal{F}^{\circ} \longrightarrow \mathcal{E}^{\circ}
\]
定义 $\mathcal{E}^{\circ}$ 之上的范畴 $\mathcal{F}^{\circ}$.

$\mathcal{F}$ 中的态射 $\alpha \colon \eta \to \xi$ 称为余笛卡儿 (co-cartesian)\index[term]{yudikaertaishe@余笛卡儿态射}，若它是 $\mathcal{E}^{\circ}$ 之上的 $\mathcal{F}^{\circ}$ 的笛卡儿态射. 明确化，看到这意味着：对 $\mathcal{F}_S$ 的每个对象 $\xi'$，映射 $u \mapsto u \circ \alpha$
\[
\mathrm{Hom}_S(\xi, \xi') \longrightarrow \mathrm{Hom}_f(\eta, \xi')
\]
双射. 此时也说 $(\xi, \alpha)$ 是 $\eta$ 经 $f$ 在 $\mathcal{E}$ 之上的范畴 $\mathcal{F}$ 中的正像 (direct image). 若它对每个 $\eta \in \mathcal{F}_T$ 存在，说经 $f$ 的正像函子存在，选定后记为 $f_{*}^{\mathcal{F}}$ 或 $f_{*}$\index[notations]{$f_{*}$}. 它由 $\mathcal{F}_T^{\circ} \times \mathcal{F}_S$ 上的双函子同构定义：
\[
\mathrm{Hom}_S(f_{*}(\eta), \xi) \stackrel{\sim}{\longrightarrow} \mathrm{Hom}_f(\eta, \xi).
\]
故若 $f_{*}$ 存在，为使 $f^{*}$ 存在，必要且充分条件是 $f_{*}$ 有伴随函子，i.e. 存在函子 $f^{*} \colon \mathcal{F}_S \to \mathcal{F}_T$ 和双函子同构
\[
\operatorname{Hom}_S(f_{*}(\eta), \xi) \xrightarrow{\sim} \operatorname{Hom}_T(\eta, f^{*}(\xi)).
\]
设 $g \colon U \to T$ 是 $\mathcal{E}$ 中的另一个态射，假设经 $f$，$g$ 和 $fg$ 的逆像和正像存在. 考虑函子同态
\[
\begin{array}{l}
c^{f,g} \colon f_{*} g_{*} \longleftarrow (fg)_{*} \\
c_{f,g} \colon g^{*} f^{*} \longrightarrow (fg)^{*}.
\end{array}
\]
看到若把 $f_{*} g_{*}$ 和 $g^{*} f^{*}$ 看作一对伴随函子，同样 $(fg)_{*}$ 和 $(fg)^{*}$，则前面两个同态彼此伴随. 故其中一个同构当且仅当另一个同构. 特别：

\begin{proposition}\label{prop:VI.10.1}
假设 $\mathcal{E}$ 之上的范畴 $\mathcal{F}$ 预纤维且余预纤维. 为使它纤维，必要且充分条件是它余纤维.
\end{proposition}

当然，若 $\mathcal{F}^{\circ}$ 在 $\mathcal{E}$ 上预纤维 (相应地：纤维)，则说 $\mathcal{F}$ 在 $\mathcal{E}$ 上余预纤维 (co-prefibered) (相应地：余纤维 (co-fibered)\index[term]{yuxianweifanchou@余纤维范畴}). 若 $\mathcal{F}$ 在 $\mathcal{E}$ 上同时纤维和余纤维，则说 $\mathcal{F}$ 在 $\mathcal{E}$ 上双纤维 (bi-fibered)\index[term]{shuangxianweifanchou@双纤维范畴}.

\section{各种例子}

\begin{enumerate}[a)]
\item \textbf{$\mathcal{E}$ 的箭范畴.} 设 $\mathcal{E}$ 是范畴. 用 $\Delta^{1}$ 记与全序二元集 $[0, 1]$ 相关的范畴；它因此有两个对象 0 和 1，除两个恒等态射外，还有一条从 0 到 1 的箭 $(0, 1)$. 设
\[
\mathbf{Fl}(\mathcal{E}) = \mathbf{Hom}(\Delta^{1}, \mathcal{E}),
\]
称它为 $\mathcal{E}$ 的箭范畴. $\Delta^{1}$ 的对象 1 定义典范函子，称为\emph{靶函子} (target functor)\index[term]{bahanzi@靶函子}
\[
\mathbf{Fl}(\mathcal{E}) \longrightarrow \mathcal{E}
\]
($\Delta^{1}$ 的对象 0 定义的函子称为源函子). 对 $\mathcal{E}$ 的每个对象 $S$，纤维范畴 $\mathbf{Fl}(\mathcal{E})_S$ 典范同构于 $S$ 之上的 $\mathcal{E}$ 的对象范畴 $\mathcal{E}_{/S}$.

考虑 $\mathcal{E}$ 中的态射 $f \colon T \to S$，它对应典范函子
\[
f_{*} \colon \mathcal{E}_{/T} = \mathcal{F}_T \longrightarrow \mathcal{E}_{/S} = \mathcal{F}_S
\]
和函子同构
\[
\operatorname{Hom}_S(f_{*}(\eta), \xi) \xrightarrow{\sim} \operatorname{Hom}_f(\eta, \xi)
\]
它因此使 $f_{*}$ 成为 $\mathcal{F}$ 中经 $f$ 的正像函子. 这里还有
\[
(\mathrm{id}_S)_{*} = \mathrm{id}_{\mathcal{F}_S}, \quad (fg)_{*} = f_{*} g_{*}, \quad c^{f,g} = \mathrm{id}_{(fg)_{*}},
\]
i.e. $\mathcal{F}$ 配备 $\mathcal{E}$ 上的余分裂 (co-splitting). 更不用说，$\mathcal{F}$ 在 $\mathcal{E}$ 上余纤维. 注意 $\mathcal{F}$ 中态射集与 $\mathcal{E}$ 中的交换方形图式一一对应. 按定义，当所述方形在 $\mathcal{E}$ 中笛卡儿，i.e. 使 $Y$ 成为 $X$ 和 $T$ 在 $S$ 上的纤维积时，该态射笛卡儿. 故经 $f$ 的逆像函子 $f^{*}$ 存在，当且仅当对 $S$ 上的每个对象 $X$，纤维积 $X \times_S T$ 存在. 由 \ref{prop:VI.10.1} 得出，若 $\mathcal{E}$ 中第三个之上的两个对象的积总是存在，i.e. 若 $\mathcal{F}$ 在 $\mathcal{E}$ 上预纤维，则 $\mathcal{F}$ 甚至在 $\mathcal{E}$ 上纤维.

\item \textbf{变动空间上的预层或层范畴.} 设 $\mathcal{E} = \mathbf{Top}$ 是拓扑空间范畴. 若 $T$ 是拓扑空间，用 $\mathcal{U}(T)$ 记 $T$ 的开集范畴，其中态射是包含映射. 若 $\mathcal{C}$ 是范畴，从 $\mathcal{U}(T)^{\circ}$ 到 $\mathcal{C}$ 的函子称为 $T$ 上取值于 $\mathcal{C}$ 的\emph{预层} (presheaf)\index[term]{yuceng@预层}，若它满足这里不再重复的左正合条件，则称为层 (sheaf)\index[term]{ceng@层}. $T$ 上取值于 $\mathcal{C}$ 的预层范畴 $\mathcal{P}(T)$ 按定义是范畴 Hom$(\mathcal{U}(T)^{\circ}, \mathcal{C})$，$T$ 上取值于 $\mathcal{C}$ 的层范畴 $\mathcal{F}(T)$ 是 Hom$(\mathcal{U}(T)^{\circ}, \mathcal{C})$ 的由是层的对象组成的全子范畴. 若 $f \colon T \to S$ 是 $\mathcal{E}$ 中的态射，i.e. 拓扑空间的连续映射，由增映射 $U \mapsto f^{-1}(U)$ 对应函子 $\mathcal{U}(S) \to \mathcal{U}(T)$，由此得函子
\[
f_{*} \colon \operatorname{Hom}(\mathcal{U}(T)^{\circ}, \mathcal{C}) \longrightarrow \operatorname{Hom}(\mathcal{U}(S)^{\circ}, \mathcal{C})
\]
称为经 $f$ 的预层正像函子. 立即看到层的正像是层，故函子 $f_{*} \colon \mathcal{P}(T) \to \mathcal{P}(S)$ 诱导函子，仍记为 $f_{*} \colon \mathcal{F}(T) \to \mathcal{F}(S)$. 平凡地验证 (用函子复合的结合性) 对第二个连续映射 $g \colon U \to T$ 有恒等
\[
(gf)_{*} = g_{*} f_{*}, \quad \text{同样 } (\mathrm{id}_S)_{*} = \mathrm{id}_{\mathcal{P}(S)}.
\]
这样，得到函子
\[
S \longmapsto \mathcal{P}(S)
\]
(相应地：
\[
S \longmapsto \mathcal{F}(S)
\])
从 $\mathcal{E}$ 到 Cat. 事实上，我们关心相应函子
\[
S \longmapsto \mathcal{P}(S)^{\circ}, \quad \text{相应地: } S \longmapsto \mathcal{F}(S)^{\circ}.
\]
它定义拓扑空间范畴上的余纤维 (甚至余分裂) 范畴，称为取值于 $\mathcal{C}$ 的预层 (相应地：层) \emph{余纤维范畴} (暗指：在变动空间之上). 明确第 8 号的构造，看到 $T$ 上的预层 $B$ 到 $S$ 上的预层 $A$ 的态射是对 $(f, u)$，由 $T$ 到 $S$ 的连续映射和 $\mathcal{P}(S)$ 中的态射 $u \colon A \to f_{*}(B)$ 组成. 这个描述对层的态射同样有效，$\mathcal{F}$ 是 $\mathcal{P}$ 的全子范畴.

在最重要的情形，$\mathcal{E}$ 之上的范畴 $\mathcal{P}$ 和 $\mathcal{F}$ 也是纤维范畴，i.e. 对每个连续映射，正像函子 $\mathcal{P}(T) \to \mathcal{P}(S)$ 和 $\mathcal{F}(T) \to \mathcal{F}(S)$ 有伴随函子，记为 $f^{*}$，称为经连续映射 $f$ 的预层逆像函子 (相应地：层逆像函子). 例如 $\mathcal{C} = \mathbf{Ens}$ 时该函子存在. 可以证明函子 $f^{*} \colon \mathcal{P}(S) \to \mathcal{P}(T)$ 只要 $\mathcal{C}$ 中存在归纳极限 (相对于所考虑宇宙中的图式) 就存在. 对 $\mathcal{F}$ 的问题较不易；注意 (即使 $\mathcal{C} = \mathbf{Ens}$) 是层的预层的逆像一般不是层，换言之，层逆像函子不同构于预层逆像函子诱导的函子 (尽管共用记号 $f^{*}$). 这样，$\mathcal{F}$ 是 $\mathcal{P}$ 的余纤维子范畴，但不是纤维子范畴，i.e. 包含函子 $\mathcal{F} \to \mathcal{P}$ 不是纤维的.

余纤维范畴 $\mathcal{P}$ 可以由更一般的余纤维 (或更确切地说纤维) 范畴推出，后者如下得到. 对每个范畴 $\mathcal{U}$ (在固定宇宙中)，令
\[
\mathcal{P}(\mathcal{U}) = \operatorname{Hom}(\mathcal{U}, \mathcal{C})
\]
并注意 $\mathcal{U} \mapsto \mathcal{P}(\mathcal{U})$ 自然地是 Cat 到 Cat 的 $U$ 的反变函子. 它因此定义 $\mathcal{E} = \mathbf{Cat}$ 之上的分裂范畴，记为 $\mathbf{Cat}_{//\mathcal{C}}$. 该范畴的对象是范畴 $\mathcal{U}$ 和函子 $p \colon \mathcal{U} \to \mathcal{C}$ 组成的对 $(\mathcal{U}, p)$，$(\mathcal{U}, p)$ 到 $(\mathcal{V}, q)$ 的态射本质上是函子 $f \colon \mathcal{U} \to \mathcal{V}$ 和函子同态 $u \colon p \to qf$ 组成的对 $(f, u)$. 把 $\mathbf{Cat}_{//\mathcal{C}}$ 中态射的复合明确留给读者. 投影函子
\[
\mathcal{F} = \mathrm{Cat}_{//\mathcal{C}} \longrightarrow \mathcal{E} = \mathrm{Cat}
\]
给对 $(\mathcal{U}, p)$ 配对象 $\mathcal{U}$；在 $\mathcal{U}$ 处的纤维范畴是范畴 Hom$(\mathcal{U}, \mathcal{C})$ (在同构意义下). 当 $\mathcal{C}$ 中存在归纳极限时，容易证明 Cat 之上的纤维范畴 $\mathbf{Cat}_{//\mathcal{C}}$ 在 Cat 上也是余纤维的，i.e. 可以定义函子 $p \colon \mathcal{U} \to \mathcal{C}$ 经函子 $f \colon \mathcal{U} \to \mathcal{V}$ 的正像概念. 预层范畴由上述纤维范畴经基变换
\[
\mathbf{Top}^{\circ} \longrightarrow \mathbf{Cat}
\]
(前面定义的函子 $S \mapsto \mathcal{U}(S)$) 推出，得到 $\mathbf{Top}^{\circ}$ 之上的纤维范畴，取对偶范畴得到 Top 之上的预层余纤维范畴 $\mathcal{P}$. 函子的逆像概念对应预层的正像概念，函子的正像概念对应预层的逆像概念.

\item \textbf{算子对象之上的算子对象.} 设 $\mathcal{F}$ 是 $\mathcal{E}$ 之上的范畴，$S$ 是 $\mathcal{E}$ 的对象，有限群 $G$ 在其上 (为确定起见) 左作用. 这个算子对象可以解释为对应函子 $\lambda \colon \mathcal{E}' \to \mathcal{E}$ (其中 $\mathcal{E}'$ 是以 $G$ 为自同态群的单对象范畴) 到 $\mathcal{E}$，并因此由基变换定义 $\mathcal{E}'$ 之上的范畴 $\mathcal{F}'$，当 $\mathcal{F}$ 在 $\mathcal{E}$ 上纤维 (相应地：余纤维) 时 $\mathcal{F}'$ 在 $\mathcal{E}'$ 上纤维 (相应地：余纤维). $\mathcal{E}'$ 在 $\mathcal{F}'$ 上的截面 (必然笛卡儿，因为 $\mathcal{E}'$ 是群胚，且由 \ref{cor:VI.6.12}，$\mathcal{F}'$ 中的每个同构笛卡儿) 也可以解释为 $\lambda$ 之上的 $\mathcal{E}$-函子 $\mathcal{E}' \to \mathcal{F}$，或解释为算子对象 $S_{\circ}$ "之上"的 $\mathcal{F}$ 中的算子对象 $\xi$.

\item \textbf{互为拟逆的伴随函子对；自对偶.} 当基范畴 $\mathcal{E}$ 约化为两个对象 $a$，$b$，且除恒等箭外还有互为逆的两个同构 $f \colon a \to b$ 和 $g \colon b \to a$ (i.e. $\mathcal{E}$ 是带两个对象的连通刚性群胚) 时，$\mathcal{E}$ 上的规范化劈开范畴本质上是两个范畴 $\mathcal{F}_a$ 和 $\mathcal{F}_b$ 及一对伴随函子 $G \colon \mathcal{F}_a \to \mathcal{F}_b$ 和 $F \colon \mathcal{F}_b \to \mathcal{F}_a$ 组成的系统，它们互为范畴等价 (故互为拟逆). 取 $\mathcal{F}_a$ 和 $\mathcal{F}_b$ 为 $\mathcal{F}$ 的纤维范畴，取 $F$ 和 $G$ 为函子 $f^{*}$ 和 $g^{*}$，两个同构
\[
u \colon FG \stackrel{\sim}{\longrightarrow} \mathrm{id}_{\mathcal{F}_a} \qquad v \colon GF \stackrel{\sim}{\longrightarrow} \mathrm{id}_{\mathcal{F}_b}
\]
是 $c_{g,f}$ 和 $c_{f,g}$. $u$ 与 $v$ 之间的两个通常相容性条件正是 $fgf$ 和 $gfg$ 的复合的 \ref{prop:VI.7.4} B) 条件. 容易证明这些条件足以蕴含确实有伪函子 $\mathcal{E}^{\circ} \to \mathbf{Cat}$.

有趣的情形是
\[
\mathcal{F}_b = \mathcal{F}_a^{\circ}, \quad G = F^{\circ}, \quad v = u^{\circ}.
\]
称范畴 $\mathcal{C}$ 中的\emph{自对偶} (self-duality)\index[term]{ziduiou@自对偶} 为函子 $D \colon \mathcal{C} \to \mathcal{C}^{\circ}$ 和同构 $u \colon DD^{\circ} \xrightarrow{\sim} \mathrm{id}_{\mathcal{C}}$ 的给出，使 $u$ 与同构 $u^{\circ} \colon D^{\circ}D \stackrel{\sim}{\longrightarrow} \mathrm{id}_{\mathcal{C}^{\circ}}$ 使 $(D, D^{\circ})$ 成为一对伴随函子 (必然互为拟逆). 该条件写成：
\[
D(u(x)) = u(D(x)) \quad \text{对每个 } x \in \mathrm{Ob}(\mathcal{C}).
\]
\item \textbf{离散范畴 $\mathcal{E}$ 之上的范畴.} 若 $\mathcal{E}$ 的每条箭都是恒等箭，则称 $\mathcal{E}$ 是\emph{离散范畴} (discrete category)\index[term]{lisanfanchou@离散范畴}，此时 $\mathcal{E}$ 在唯一同构意义下由集合 $I = \mathrm{Ob}(\mathcal{E})$ 确定. $\mathcal{E}$ 之上范畴 $\mathcal{F}$ 的给出于是 (在唯一同构意义下) 等价于一族范畴 $\mathcal{F}_i$ $(i \in I)$ 的给出，即纤维范畴. $\mathcal{E}$ 之上的每个范畴 $\mathcal{F}$ 纤维，每个 $\mathcal{E}$-函子 $\mathcal{F} \to \mathcal{G}$ 笛卡儿，有典范同构
\[
\mathbf{Hom}_{\mathcal{E}/-}(\mathcal{F}, \mathcal{G}) \stackrel{\sim}{\longrightarrow} \prod_{i} \mathbf{Hom}(\mathcal{F}_i, \mathcal{G}_i).
\]
特别得到
\[
\Gamma(\mathcal{F}/\mathcal{E}) = \varprojlim \mathcal{F}/\mathcal{E} \stackrel{\sim}{\longrightarrow} \prod_{i} \mathcal{F}_i.
\]
\item 假设 $\mathcal{E}$ 恰有两个对象 $S$ 和 $T$，且除恒等态射外还有态射 $f \colon T \to S$. 则 $\mathcal{E}$ 之上的范畴 $\mathcal{F}$ 在唯一 $\mathcal{E}$-同构意义下由两个范畴 $\mathcal{F}_S$ 和 $\mathcal{F}_T$ 及取值于 Ens 的 $\mathcal{F}_T^{\circ} \times \mathcal{F}_S$ 上的双函子 $H(\eta, \xi)$ 定义. 事实上，若 $\mathcal{F}$ 是 $\mathcal{E}$ 之上的范畴，给它配两个纤维范畴 $\mathcal{F}_S$ 和 $\mathcal{F}_T$ 及双函子 $H(\eta, \xi) = \mathrm{Hom}_f(\eta, \xi)$. 反向构造的明确留给读者. 为使所考虑的范畴纤维 (或预纤维，这在此相同)，必要且充分条件是函子 $H$ 关于自变量 $\xi$ 可表示. 为使它余纤维，必要且充分条件是 $H$ 关于自变量 $\eta$ 可表示.
\item 设 $\mathcal{F} = \mathcal{C} \times \mathcal{E}$，借助 $\mathrm{pr}_2$ 把它看作 $\mathcal{E}$ 之上的范畴. 则 $\mathcal{F}$ 在 $\mathcal{E}$ 上纤维且余纤维，甚至配备由 $\mathcal{E}$ (相应地：$\mathcal{E}^{\circ}$) 上取常值 $\mathcal{C}$ 的函子对应的典范分裂和余分裂. 有
\[
\Gamma(\mathcal{F}/\mathcal{E}) \simeq \operatorname{Hom}(\mathcal{E}, \mathcal{C})
\]
且 Lim $\mathcal{F}/\mathcal{E}$ 对应于由把任意态射变为同构的函子 $F \colon \mathcal{E} \to \mathcal{C}$ 组成的全子范畴.
\end{enumerate}

\section{劈开范畴上的函子}

设 $\mathcal{F}$ 是 $\mathcal{E}$ 上的规范化劈开范畴. 对 $\mathcal{E}$ 的每个对象 $S$，用
\[
i_S \colon \mathcal{F}_S \longrightarrow \mathcal{F}
\]
记包含函子. 对 $\mathcal{E}$ 中的每个态射 $f \colon T \to S$，有函子同态
\[
\alpha_f \colon i_T f^{*} \longrightarrow i_S,
\]
其中 $f^{*}$ 是劈开定义的 $f$ 的基变换函子 $\mathcal{F}_S \to \mathcal{F}_T$. 现在设
\[
F \colon \mathcal{F} \longrightarrow \mathcal{C}
\]
是 $\mathcal{F}$ 到范畴 $\mathcal{C}$ 的函子，对每个 $S \in \mathrm{Ob}(\mathcal{E})$ 令
\[
F_S = F \circ i_S \colon \mathcal{F}_S \longrightarrow \mathcal{C},
\]
对 $\mathcal{E}$ 中的每个 $f \colon T \to S$，
\[
\varphi_f = F * \alpha_f \colon F_T f^{*} \longrightarrow F_S.
\]
这样，对每个函子 $F \colon \mathcal{F} \to \mathcal{C}$，配了一族函子 $F_S \colon \mathcal{F}_S \to \mathcal{C}$ 和一族函子同态 $\varphi_f \colon F_T f^{*} \to F_S$. 这些族满足下列条件：
\begin{enumerate}[a)]
\item $\varphi_{\mathrm{id}_S} = \mathrm{id}_{F_S}$.
\item 对 $\mathcal{E}$ 中的两个态射 $f \colon T \to S$ 和 $g \colon U \to T$，在函子同态方形中有交换性：
\end{enumerate}
\[
\begin{tikzcd}[row sep=large, column sep=large]
F_U g^{*} f^{*} \arrow[r,"{F_U * c_{f,g}}"] \arrow[d,"{\varphi_g * f^{*}}"] & F_U (fg)^{*} \arrow[d,"{\varphi_{fg}}"] \\
F_T f^{*} \arrow[r,"{\varphi_f}"] & F_S.
\end{tikzcd}
\]
第一个关系平凡，第二个关系通过把函子 $F$ 应用于交换图式
\[
\begin{tikzcd}[row sep=large, column sep=large]
g^{*} f^{*}(\xi) \arrow[r,"{c_{f,g}(\xi)}"] \arrow[d,"{\alpha_g(f^{*}(\xi))}"] & (fg)^{*}(\xi) \arrow[d,"{\alpha_{fg}(\xi)}"] \\
f^{*}(\xi) \arrow[r,"{\alpha_f(\xi)}"] & \xi
\end{tikzcd}
\]
(对 $\mathcal{F}_S$ 中变动的对象 $\xi$) 得到.

若 $G$ 是第二个函子 $\mathcal{F} \to \mathcal{C}$，产生函子 $G_S \colon \mathcal{F}_S \to \mathcal{C}$ 和函子同态 $\psi_f \colon G_T f \to G_S$，且 $u \colon F \to G$ 是函子同态，则它对应函子同态 $u * i_S$：
\[
u_S \colon F_S \longrightarrow G_S,
\]
并立即看到对 $\mathcal{E}$ 中的每个态射 $f \colon T \to S$，在方形中交换
\begin{equation}
\begin{tikzcd}[row sep=large, column sep=large]
F_T f^{*} \arrow[r,"{\phi_f}"] \arrow[d,"{u_T * f^{*}}"] & F_S \arrow[d,"{u_S}"] \\
G_T f^{*} \arrow[r,"{\psi_f}"] & G_S
\end{tikzcd} \tag{c}
\label{eq:VI.c}
\end{equation}

\begin{proposition}\label{prop:VI.12.1}
设 $\mathcal{H}(\mathcal{F}, \mathcal{C})$ 是对象为"函子 $F_S \colon \mathcal{F}_S \to \mathcal{C}$ 的族 $(F_S)$ ($S \in \mathrm{Ob}(\mathcal{E})$) 与函子同态 $F_T f^{*} \to F_S$ 的族 $(\varphi_f)$ ($f \in \mathrm{Fl}(\mathcal{E})$) 组成的对、满足条件 a) 和 b)"、态射为满足前面写出的交换条件 c) 的同态族 $(u_S)$ ($S \in \mathrm{Ob}(\mathcal{E})$) 的范畴 (态射的复合通过函子同态 $\mathcal{F}_S \to \mathcal{C}$ 的复合进行). 则上面明确的两个法则定义 Hom$(\mathcal{F}, \mathcal{C})$ 与范畴 $\mathcal{H}(\mathcal{F}, \mathcal{C})$ 的同构 $K$.
\end{proposition}

显然得到第一个范畴到第二个的函子. 该函子全忠实，因为对给定的 $F$，$G$，映射
\[
\mathrm{Hom}(F, G) \longrightarrow \mathrm{Hom}(K(F), K(G))
\]
平凡单射；为证满射，只需注意交换条件 c) 表达映射 $u(\xi) = u_S(\xi) \colon F(S) = F_S(\xi) \to G(\xi) = G_S(\xi)$ 对 $\mathcal{F}$ 中形如 $\alpha_f(\xi)$ 的同态的函子性，另一方面在每个纤维范畴上 (i.e. 对 $\mathcal{F}$ 中是 $T$-态射 ($T \in \mathrm{Ob}(\mathcal{E})$) 的态射) 有函子性，故对 $\mathcal{F}$ 中的每个态射有函子性，因为 $f$-态射 ($f \colon T \to S$ 是 $\mathcal{E}$ 中的态射) 唯一地是 $\alpha_f(\xi)$ 与 $T$-态射的复合. 剩下证明函子 $K$ 对对象双射. 前面的论证已表明 $K$ 对对象单射，剩下证明满射，i.e. 若从满足 a) 和 b) 的系统 $(F_S), (\varphi_f)$ 出发，用
\[
F(\xi) = F_S(\xi) \quad \text{对 } \xi \in \mathrm{Ob} \mathcal{F}_S \subset \mathrm{Ob} \mathcal{F}
\]
定义映射 $\mathrm{Ob} \mathcal{F} \to \mathrm{Ob} \mathcal{C}$，用
\[
F(\alpha_f(\xi) u') = \varphi_f(\xi) F_T(u')
\]
(对 $\mathcal{E}$ 中的每个态射 $f \colon T \to S$，$\mathcal{F}_S$ 的每个对象 $\xi$ 和以 $f^{*}(\xi)$ 为靶的每个 $T$-态射 $u'$) 定义映射 $\mathrm{Fl}(\mathcal{F}) \to \mathrm{Fl}(\mathcal{C})$，则得到 $\mathcal{F}$ 到 $\mathcal{C}$ 的函子 $F$. 事实上，关系 $F(\mathrm{id}_{\xi}) = \mathrm{id}_{F(\xi)}$ 平凡，剩下在有一个 $f$-态射 $u \colon \eta \to \xi$ 和一个 $g$-态射 $v \colon \zeta \to \nu$ (其中 $f \colon T \to S$，$g \colon U \to T$ 是 $\mathcal{E}$ 的态射) 时证明乘性 $F(uv) = F(u)F(v)$. 令 $w = uv$，有
\[
u = \alpha_f(\xi) u', v = \alpha_g(\eta) v', w = \alpha_{fg}(\xi) w'
\]
其中
\[
w' = c_{f,g}(\xi) g^{*}(u') v' \quad (\text{参见第 8 号}).
\]
用这些记号，要证明下面图式外轮廓的交换性：
\[
\begin{tikzcd}[row sep=large, column sep=large]
F_U(\zeta) \arrow[rrr, phantom, "F(w')"] \arrow[r,"{F_U(v')}"] \arrow[d,"{\varphi_g(\eta)}"] & F_U g^{*}(\eta) \arrow[r,"{F_U g^{*}(u')}"] \arrow[d,"{\varphi_g(f^{*}(\xi))}"] & F_U g^{*} f^{*}(\xi) \arrow[r,"{F_U(c_{f,g}(\xi))}"] \arrow[d,"{\varphi_{fg}(\xi)}"] & F_U (fg)^{*}(\xi) \\
F_T(\eta) \arrow[rr, phantom, "F(u)"] \arrow[r,"{F_T(u')}"] & F_T f^{*}(\xi) \arrow[r,"{\varphi_f(\xi)}"] & F_S(\xi)
\end{tikzcd}
\]
左边三角形按 $F(v)$ 的定义交换，中间方形按 $u' \colon \xi \to f(\eta)$ 经函子同态 $\alpha_g$ 得出的定义交换，右边方形按条件 b) 交换. 结论由此得出.

现在假设 $\mathcal{C}$ 也是 $\mathcal{E}$ 上的规范化劈开范畴，此后称它为 $\mathcal{G}$，且关心 $\mathcal{F}$ 到 $\mathcal{G}$ 的 $\mathcal{E}$-函子. 若 $F$ 是这样一个函子，它对纤维范畴诱导函子
\[
F_S \colon \mathcal{F}_S \longrightarrow \mathcal{G}_S.
\]
另一方面，对 $\mathcal{E}$ 中的每个态射 $f \colon T \to S$ 和 $\mathcal{F}_S$ 的每个对象 $\xi$，$f$-态射 $F(\alpha_f(\xi))$ 唯一地经 $T$-态射
\[
\varphi_f(\xi) \colon F_T(f_{\mathcal{F}}^{*}(\xi)) \longrightarrow f_{\mathcal{G}}^{*}(F_S(\xi))
\]
分解 (其中下标 $\mathcal{F}$ 或 $\mathcal{G}$ 指明取逆像函子的劈开范畴)，由此得 $\mathcal{F}_S$ 到 $\mathcal{G}_T$ 的函子同态：
\[
\varphi_f \colon F_T f_{\mathcal{F}}^{*} \longrightarrow f_{\mathcal{G}}^{*} F_S.
\]
两个系统 $(F_S)$ 和 $(\varphi_f)$ 满足下列条件：
a') $\varphi_{\mathrm{id}_S} = \mathrm{id}_{F_S}$.
b') 对 $\mathcal{E}$ 中的两个态射 $f \colon T \to S$ 和 $g \colon U \to T$，在下列函子同态图式中有交换性：
\[
\begin{tikzcd}[row sep=large, column sep=large]
F_U g_{\mathcal{F}}^{*} f_{\mathcal{F}}^{*} \arrow[r,"{F_U * c^{\mathcal{F}}_{f,g}}"] \arrow[d,"{\varphi_g f_{\mathcal{F}}^{*}}"] & F_U (fg)_{\mathcal{F}}^{*} \\
g_{\mathcal{G}}^{*} F_T f_{\mathcal{F}}^{*} \arrow[d,"{g_{\mathcal{G}}^{*} * \varphi_f}"] & \\
g_{\mathcal{G}}^{*} f_{\mathcal{G}}^{*} F_S \arrow[r,"{c^{\mathcal{G}}_{f,g} * F_S}"] & (fg)_{\mathcal{G}}^{*} F_S.
\end{tikzcd}
\]
验证留给读者，\cref{prop:VI.12.1} 的类似陈述的表述和证明也留给读者，它蕴含 $\mathcal{F}$ 到 $\mathcal{G}$ 的 $\mathcal{E}$-函子集与满足上面条件 a') 和 b') 的系统 $(F_S), (\varphi_f)$ 集之间的一一对应. 当然，在这个对应中，笛卡儿函子由"诸同态 $\varphi_f$ 是同构"这一性质刻画.

\begin{remarkstar}
当然，大多数时候直接在纤维范畴上推理而不使用显式劈开是有利的，这尤其使简单概念 $\mathcal{E}$-函子或笛卡儿 $\mathcal{E}$-函子不需要上面那样笨重的解释. 正是为了避免难以忍受的笨拙、并得到更内蕴的陈述，我们不得不放弃像 \cite{exp06-2} 那样从劈开范畴 (loc. cit. 中称为"纤维范畴") 的概念出发，后者退居其次，让位于纤维范畴概念. 而且很可能，与目前仍占优势、与旧思维习惯相关的用法相反，在万有问题中最终将被证明更便利的，不是把重点放在某个假定一次选定好的解上，而是让所有解处于平等地位.
\end{remarkstar}

\section{文献}\printbibliography[heading=none]


\end{refsection}
\begin{refsection}% =====================================================================
% 第八讲 (Exposé VIII) 忠实平坦下降
% (Descente fidèlement plate)
% M. Raynaud
% 原书页码: pp. 153-176
% =====================================================================
\setcounter{chapter}{7}
\chapter{忠实平坦下降}

\section{拟凝聚模的下降}

设 $\mathbf{Sch}$ 是预概形范畴\index[notations]{$\mathbf{Sch}$}. 如 VI 11 b) 那样处理，我们发现由预概形 $X$ 和 $X$ 上的模 $F$ 组成的对 $(X, F)$ 的范畴 (其态射如 loc. cit. 那样借助模经环化空间态射的直像概念定义) 可以看作 $\mathbf{Sch}$ 之上的纤维范畴，相对 $\mathbf{Sch}$ 中态射 $f \colon X \to Y$ 的基变换函子是模经 $f$ 的逆像函子. (注意前述纤维范畴在 $X \in \mathrm{Ob}(\mathbf{Sch})$ 处的纤维范畴是 $X$ 上模范畴的对偶范畴). 由于拟凝聚模的逆像拟凝聚，看到对 $(X, F)$ 中 $F$ 拟凝聚的配对组成的范畴的全子范畴是前述纤维范畴的纤维子范畴. (相反，若对 $f$ 不作假设，拟凝聚模的直像一般不拟凝聚). 我们简称为预概形上拟凝聚模的纤维范畴.

另一方面回忆，环化空间态射 $f \colon X \to Y$ 称为忠实平坦，若它平坦 (i.e. 对每个 $x \in X$，$\mathcal{O}_{X,x}$ 是 $\mathcal{O}_{Y,f(x)}$ 上的平坦模，(IV))，且满射. 称 $f$ 是\emph{拟紧态射}\index[term]{nijintaishe@拟紧态射}，若任何拟紧部分的原像是拟紧的；当 $f$ 是预概形态射时，这也意味着 $Y$ 的仿射开集的原像是 $X$ 的有限个仿射开集的并.

\begin{theorem}\label{thm:VIII.1.1}
设 $\mathcal{F}$ 是预概形上拟凝聚模的纤维范畴. 设 $g \colon S' \to S$ 是预概形态射，忠实平坦且拟紧. 则 $g$ 是有效 $\mathcal{F}$-下降态射\index[term]{youxiaoxiajiangtaishe@有效下降态射}.
\end{theorem}

回忆\footnote{这里承认下降的一般理论，它在 J. Giraud 的文章 (引自 Avertissement 的脚注 (5)) 中有详细阐述，此后我们把它引用为 \cite{exp08-D}. 也可参见 \cite{exp08-2} 的简明阐述.} 这意味着两件事：

\begin{corollary}\label{cor:VIII.1.2}
(\emph{模同态的下降})\index[term]{motongtaidexiajiang@模同态的下降} 设 $g \colon S' \to S$ 是预概形态射，忠实平坦且拟紧，$F$ 和 $G$ 是 $S$ 上的两个拟凝聚模，$F'$ 和 $G'$ 是它们在 $S'$ 上的逆像，最后 $F''$ 和 $G''$ 是它们在 $S'' = S' \times_S S'$ 上的逆像. 考虑由经 $g$，$p_1$，$p_2$ 的基变换函子 (其中 $p_1, p_2 \colon S' \times_S S' \rightrightarrows S'$ 是两个投影) 定义的集合映射图式：
\[
\mathrm{Hom}_S(F, G) \longrightarrow \mathrm{Hom}_{S'}(F', G') \rightrightarrows \mathrm{Hom}_{S''}(F'', G'').
\]
该图式正合，i.e. 定义第一个集合到第二个到第三个的两个映射的一致集的双射.
\end{corollary}

换言之，经 $g$ 的基变换函子 $F \mapsto F'$ 定义从 $S$ 上拟凝聚模范畴到配备相对 $g$ 的下降数据的 $S'$ 上拟凝聚模范畴的忠实全函子. 此外：

\begin{corollary}\label{cor:VIII.1.3}
(\emph{模的下降})\index[term]{modexiajiang@模的下降} 对 $S'$ 上的每个拟凝聚模 $F'$，$F'$ 上相对 $g$ 的每个下降数据有效，i.e. $F'$ 连同其下降数据同构于 $S$ 上某拟凝聚模经 $g$ 的逆像 (由 \ref{cor:VIII.1.2} 在唯一同构意义下确定).
\end{corollary}

换言之，前述忠实全函子甚至是等价. 实际上，这意味着给出 $S$ 上的拟凝聚模，与给出 $S'$ 上配备相对 $g$ 的下降数据的拟凝聚模，是同一回事.

\begin{proof}
先设 $T$ 是 $S$-预概形，它 $S$-同构于覆盖 $S$ 的诱导开集 $S_i$ 族 $S$ 的和. 则结构态射 $T \to S$ 显然是有效 $\mathcal{F}$-下降态射 (这精确意味着 $S$ 上拟凝聚模 $F$ 的给出等价于诸 $S_i$ 上拟凝聚模 $F_i$ 及满足众所周知的上链条件的粘合同构 $\varphi_{ji} \colon F_i|_{S_i \cap S_j} \to F_j|_{S_i \cap S_j}$ 的给出). 由 VII，8 得出，为验证 $g \colon S' \to S$ 是有效 $\mathcal{F}$-下降态射，只需对由 $g$ 经基变换 $T \to S$ 得到的态射 $g_T \colon T' = T \times_S S' \to T$ 验证. (注意对 $T \to S$ 的假设在任何基变换下保持，故 $T \to S$ 实际上是万有有效 $\mathcal{F}$-下降态射). 取覆盖 $S$ 的仿射开集作 $S_i$，于是归结到 $S$ 仿射的情形.

则 $S'$ 是有限个仿射开集的并，取这些开集的和 $S$-概形，得到仿射 $S$-概形 $S_1$ 和满射平坦的 $S$-态射 $S_1 \to S'$. 故 $S_1$ 在 $S$ 上也忠实平坦. 因此若证明忠实平坦且仿射的态射是有效 $\mathcal{F}$-下降态射，从而是万有严格 $\mathcal{F}$-下降态射 (该假设确实在基变换下保持)，则特别得出结论：结构态射 $S_1 \to S$ 是万有严格 $\mathcal{F}$-下降态射，且由于存在 $S$-态射 $S_1 \to S'$，由 \cite{exp08-D} 将得出 $g \colon S' \to S$ 是严格 $\mathcal{F}$-下降态射.

这归结到 $g$ 是仿射态射的情形，且如所见此时还可假设 $S$ 仿射，故可假设 $S$ 和 $S'$ 仿射. 此时 \ref{cor:VIII.1.2} 等价于
\end{proof}

\begin{lemma}\label{lem:VIII.1.4}
设 $A$ 是环，$A'$ 是忠实平坦的 $A$-代数，$M$ 和 $N$ 是两个 $A$-模，$M'$ 和 $N'$ 是由环变换 $A \to A'$ 得到的 $A'$-模，$M''$，$N''$ 是由环变换 $A \to A''$ 得到的 $A'' = A' \otimes_A A'$-模. 则集合映射序列
\[
\mathrm{Hom}_A(M, N) \longrightarrow \mathrm{Hom}_{A'}(M', N') \rightrightarrows \mathrm{Hom}_{A''}(M'', N'')
\]
正合.
\end{lemma}

由于 $N \to N'$ 单射 ($A'$ 在 $A$ 上忠实平坦)，看到第一个箭单射. 剩下证明：若 $A'$-同态 $u' \colon M' \to N'$ 与下降数据相容，则它来自 $A$-同态 $u \colon M \to N$. 这也简单意味着 $u'$ 把 $M'$ 的子集 $M$ 映入 $N'$ 的子集 $N$ (诱导映射 $u \colon M \to N$ 将自动 $A$-线性，因为 $u'$ 是 $A'$-线性的，同样看到 $u'$ 必然等于 $u \otimes_A A'$). 若 $x \in M$，则 $u'(x)$ 是映射对 $N' \rightrightarrows N''$ 的核中的元素. 为证 \ref{lem:VIII.1.4}，于是归结到下面的特殊情形 (对应于 $M = A$ 的情形)：

\begin{corollary}\label{cor:VIII.1.5}
设 $N$ 是 $A$-模，则集合映射序列
\[
N \longrightarrow N' \rightrightarrows N''
\]
正合.
\end{corollary}

设 $A_1$ 是忠实平坦的 $A$-代数. 为证所述序列正合，只需证经环变换 $A \to A_1$ 得到的序列正合. 后者显然正是 $A_1$-模 $N_1 = N \otimes_A A_1$ 和 $A_1$-代数 $A'_1 = A_1 \otimes_A A'$ 所对应的序列. 故只需找 $A$ 上忠实平坦的 $A_1$，使 $\mathrm{Spec}(A'_1) \to \mathrm{Spec}(A_1)$ 是严格 $\mathcal{F}$-下降态射. 取 $A_1 = A'$ 即可，因为此时前述态射有右逆态射，从而由 \cite{exp08-D} 它对 $\mathbf{Sch}$ 上的任何纤维范畴是有效下降态射.

最后剩下证明：若 $N'$ 是配备对 $A \to A'$ 的下降数据的 $A'$-模，i.e. 配备由 $N$ 经环变换 $A' \rightrightarrows A' \otimes_A A'$ 得到的两个模之间的同构
\[
\varphi \colon N'_1 \xrightarrow{\sim} N'_2,
\]
则 $N'$ 连同其下降数据同构于模 $N \otimes_A A'$. 考虑到 \ref{cor:VIII.1.5}，容易看到该断言等价于：

\begin{lemma}\label{lem:VIII.1.6}
设 $N'$ 是配备相对 $A \to A'$ (其中 $A'$ 是 $A$-代数) 的下降数据的 $A'$-模. 设 $N$ 是 $N'$ 的由满足 $\varphi(x \otimes_{A'} 1_{A'}) = 1_{A'} \otimes_{A'} x$ 的 $x$ 组成的子 $A$-模，考虑典范同态
\[
N \otimes_A A' \longrightarrow N',
\]
(它与下降数据相容). 若 $A'$ 在 $A$ 上忠实平坦，则该同态是同构.
\end{lemma}

\begin{proof}
证明该引理. 再设 $A_1$ 是忠实平坦的 $A$-代数，为证所述态射是同构，只需证它在环变换 $A \to A_1$ 后成为同构. 利用 $A_1$ 在 $A$ 上的平坦性，看到如此得到的同态正是按 $A'_1 = A' \otimes_A A_1$ 上的模 $N' \otimes_A A_1$ 项下、配备相对 $A_1 \to A'_1$ 的由 $N'$ 上给定数据经环变换典范得到的下降数据而直接得到的那一个. 故只需找 $A$ 上忠实平坦的 $A_1$，使 $\mathrm{Spec}(A'_1) \to \mathrm{Spec}(A_1)$ 是有效 $\mathcal{F}$-下降态射. 如上取 $A_1 = A'$. 这完成 \ref{lem:VIII.1.6} 的证明，从而完成 \ref{thm:VIII.1.1} 的证明.
\end{proof}

\begin{corollary}\label{cor:VIII.1.7}
(模的截面的下降)\index[term]{mojiemiandexiajiang@模截面的下降} 设 $g \colon S' \to S$ 是预概形态射，忠实平坦且拟紧. 对 $S$ 上的每个拟凝聚模 $G$，设 $G'$ 和 $G''$ 是它在 $S'$ 和 $S'' = S' \times_S S'$ 上的逆像，考虑 $S$ 上模同态图式
\[
G \longrightarrow g_{*}(G') \rightrightarrows h_{*}(G'')
\]
(其中 $h \colon S'' \to S$ 是结构态射). 该图式正合.
\end{corollary}

事实上，这意味着对 $S$ 中的每个开集 $U$，由 $U$ 上截面组成的对应图式正合. 显然可假设 $U = S$，此时所述正合性是 \ref{cor:VIII.1.2} 的特例，取 $F = \mathcal{O}_S$ 得到.

由于模的逆像函子右正合，由 \ref{thm:VIII.1.1} 形式地得出：

\begin{corollary}\label{cor:VIII.1.8}
(\emph{商模的下降})\index[term]{shangmodexiajiang@商模的下降} 用 \ref{cor:VIII.1.7} 的记号，设 $\mathrm{Quot}(F)$ 对每个预概形上的拟凝聚模 $F$ 表示 $F$ 的拟凝聚商模集. 在此约定下，集合映射图式
\[
\mathrm{Quot}(G) \longrightarrow \mathrm{Quot}(G') \rightrightarrows \mathrm{Quot}(G'')
\]
正合.
\end{corollary}

(用子模代替商模显然有同样的断言，因为两者一一对应). 特别取 $G = \mathcal{O}_S$，得到：

\begin{corollary}\label{cor:VIII.1.9}
(\emph{闭子预概形的下降})\index[term]{biziyugaixingdexiajiang@闭子预概形的下降} 对每个预概形 $X$，设 $H(X)$ 表示 $X$ 的闭子预概形集. 在此记号下，在 \ref{cor:VIII.1.7} 的条件下，下列集合映射图式
\[
H(S) \longrightarrow H(S') \rightrightarrows H(S'')
\]
正合.
\end{corollary}

有必要用下列结果补充\cref{thm:VIII.1.1}：

\begin{proposition}\label{prop:VIII.1.10}
(模的性质的下降)\index[term]{moxingzhidexiajiang@模性质的下降} 设 $g \colon S' \to S$ 是忠实平坦拟紧态射，$F$ 是 $S$ 上的拟凝聚模. 为使 $F$ 有限型 (resp. 有限表现，resp. 局部自由且有限型)，必要且充分条件是它在 $S'$ 上的逆像 $F'$ 如此.
\end{proposition}

只需证明"充分"部分. 显然可假设 $S$ 仿射，再把 $S'$ 换成覆盖 $S'$ 的仿射开集的和，归结到 $S'$ 也仿射的情形. 此时断言等价于：

\begin{corollary}\label{cor:VIII.1.11}
设 $A$ 是环，$A'$ 是忠实平坦的 $A$-代数，$M$ 是 $A$-模，$M'$ 是 $A'$-模 $M \otimes_A A'$. 为使 $M$ 有限型 (resp. 有限表现，resp. 有限型局部自由)，必要且充分条件是 $M'$ 如此.
\end{corollary}

事实上，有 $M = \varinjlim_i M_i$，其中 $M_i$ 是 $M$ 的有限型子模. 从而 $M' = \varinjlim_i M'_i$，若 $M'$ 有限型，则 $M'$ 等于某个 $M'_i$，由忠实平坦性 $M$ 等于 $M_i$，故 $M$ 有限型. 于是存在正合序列
\[
0 \longrightarrow R \longrightarrow L \longrightarrow M \longrightarrow 0,
\]
其中 $L$ 有限型自由，从而有正合序列
\[
0 \longrightarrow R' \longrightarrow L' \longrightarrow M' \longrightarrow 0,
\]
其中 $L'$ 有限型自由. 若 $M'$ 有限表现，则 $R'$ 有限型，由前面所述 $R$ 有限型，故 $M$ 有限表现. 最后，说 $M$ 局部自由且有限型，意味着它有限表现且平坦 (参见 IV，诺特情形；一般情形留给读者). 由于这些性质各自都随下降保持，其合取亦然，这完成证明.

\begin{remark}\label{rem:VIII.1.12}
\ref{thm:VIII.1.1} 与 \ref{prop:VIII.1.10} 的合取表明，把纤维范畴 $\mathcal{F}$ 换成由有限型 (resp. 有限表现，resp. 有限型局部自由，resp. 秩为给定 $n$ 的局部自由) 拟凝聚模组成的纤维子范畴时，断言 \ref{thm:VIII.1.1} 仍然成立.
\end{remark}

\section{仿射预概形在另一个之上的下降}

由于模的逆像函子与张量积及其他张量运算相容，\cref{thm:VIII.1.1} 蕴含各种变体，只需考虑(代替单个拟凝聚模)配备各种借助张量运算表达的附加结构的拟凝聚模或拟凝聚模组. 例如，$S$ 上三个拟凝聚模 $F$，$G$，$H$ 及配对
\[
F \otimes G \longrightarrow H
\]
的给出，等价于 $S'$ 上三个拟凝聚模 $F'$，$G'$，$H'$，配备相对 $g \colon S' \to S$ 的下降数据，并配备与这些下降数据"相容" (按此词的显然含义) 的配对
\[
F' \otimes G' \longrightarrow H'
\]
的给出. 例如，若 $F = G = H$，看到 $S$ 上配备代数律 (暂时不假设它满足结合律、交换律或单位截面的存在性公理) 的拟凝聚模 $F$ 的给出，等价于 $S'$ 上同样的给出，外加一个下降数据. 利用前一节的结果，立即看到为使 $F$ 满足刚才提到的通常公理之一，必要且充分条件是 $F'$ 如此. 例如，$S$ 上拟凝聚代数 (此后我们用此词总指：结合的、交换的、有单位截面的) $\mathcal{A}$ 的给出，等价于 $S'$ 上拟凝聚代数 $\mathcal{A}'$ 配备相对 $g \colon S' \to S$ 的下降数据的给出. 若回忆 $S$ 上拟凝聚代数的对偶范畴与 $S$ 上仿射 $S$-预概形范畴之间的等价 (EGA II §1)，立即得到：

\begin{theorem}\label{thm:VIII.2.1}
设 $\mathcal{F}'$ 是预概形仿射态射 $f \colon X \to S$ 的纤维范畴，把它看作预概形范畴 $\mathbf{Sch}$ 中箭的纤维范畴 (VI 11.a) 的纤维子范畴. 设 $g \colon S' \to S$ 是忠实平坦拟紧的预概形态射. 则 $g$ 是有效 $\mathcal{F}'$-下降态射.
\end{theorem}

\section{态射的集合性质与有限性性质的下降}

\begin{proposition}\label{prop:VIII.3.1}
设 $f \colon X \to Y$ 是 $S$-态射\footnote{关于第 3 和第 4 号处理的这类结果的其他结果，参见 EGA IV 2.3，2.6，2.7.}，$g \colon S' \to S$ 是满射态射，$f' \colon X' = X \times_S S' \to Y' = Y \times_S S'$ 是由 $f$ 借助 $g \colon S' \to S$ 作基变换得到的态射. 为使 $f$ 满射 (resp. 根式)，必要且充分条件是 $f'$ 如此.
\end{proposition}

注意 $f'$ 也可经基变换 $Y' \to Y$ 得到，后者也满射，因为它由满射的 $g \colon S' \to S$ 得到. 另一方面，对每个 $y \in Y$ 和 $y$ 上的每个 $y' \in Y'$，有同构
\[
X'_{y'} \xleftarrow{\sim} X_y \otimes_{k(y)} k(y'),
\]
其中 $X_y$ 表示 $X$ 在 $y$ 处的纤维，$X'_{y'}$ 表示 $X'$ 在 $y'$ 处的纤维. 由此得出，$X_y$ 非空 (resp. 至多有一个点，且该点对应根式剩余扩张) 当且仅当 $X'_{y'}$ 有同样性质. 这证明 3.1.

\begin{corollary}\label{cor:VIII.3.2}
在 \ref{prop:VIII.3.1} 的条件下，若 $f'$ 单射 (resp. 双射)，则 $f$ 也如此.
\end{corollary}

这是因为若 $X'_{y'}$ 至多有一个点 (resp. 恰好一个点)，则 $X_y$ 亦然；确实如此，因为态射 $X'_{y'} \to X_y$ 满射 (它由满射的 $\mathrm{Spec}(k(y')) \to \mathrm{Spec}(k(y))$ 得到).

\begin{proposition}\label{prop:VIII.3.3}
用 \ref{prop:VIII.3.1} 的记号. 假设 $g \colon S' \to S$ 满射且拟紧 (resp. 忠实平坦且拟紧). 为使 $f$ 拟紧 (resp. 有限型)，必要且充分条件是 $f'$ 如此.
\end{proposition}

只需证明"充分"部分. 显然可假设 $S = Y$，因为对 $g \colon S' \to S$ 的假设对 $Y' \to Y$ 保持. 还可假设 $Y$ 仿射. 则 $Y'$ 拟紧，故 $X'$ 拟紧 (因为 $f'$ 由假设拟紧). 设 $(X_i)_{i \in I}$ 是覆盖 $X$ 的仿射开集族，则诸 $X'_i$ 是覆盖 $X'$ 的开集，故有覆盖 $X'$ 的有限子族. 由于 $X' \to X$ 满射，得出对应的诸 $X_i$ 已经覆盖 $X$，故 $X$ 拟紧，i.e. $f$ 拟紧. 现在假设 $f'$ 有限型，在 $g$ 忠实平坦的假设下证明 $f$ 有限型. 把 $Y'$ 换成覆盖它的仿射开集族的和，可假设 $Y'$ 仿射. 最后，由前面所述 $X$ 被有限个仿射开集 $X_i$ 覆盖，在已知 $X'_i$ 在 $Y'$ 上有限型的条件下需证明它们在 $Y$ 上有限型. 这归结到：

\begin{corollary}\label{cor:VIII.3.4}
设 $B$ 是 $A$-代数，$A'$ 是忠实平坦的 $A$-代数，$B' = B \otimes_A A'$ 是由 $B$ 经环变换得到的 $A'$-代数. 为使 $B$ 有限型，必要且充分条件是 $B'$ 如此.
\end{corollary}

只需证明"充分"部分. 有 $B = \varinjlim_i B_i$，其中 $B_i$ 是 $B$ 的有限型子代数. 故有 $B' = \varinjlim_i B'_i$，若 $B'$ 在 $A'$ 上有限型，则 $B'$ 等于某个 $B'_i$，故 $B$ 等于 $B_i$，从而有限型.

\begin{corollary}\label{cor:VIII.3.5}
仍假设基变换态射 $g \colon S' \to S$ 忠实平坦且拟紧. 为使 $f$ 拟有限，必要且充分条件是 $f'$ 如此.
\end{corollary}

事实上，"拟有限"性质按定义是"有限型"与"纤维有限"的合取，二者各自都能随 $g$ 下降，第一个由 \ref{prop:VIII.3.3}，第二个由 \ref{prop:VIII.3.1} 的推理 (只用 $g$ 满射的事实).

\begin{remark}\label{rem:VIII.3.6}
设 $A$ 是环，$X$ 是 $A$-预概形. 容易看到下列条件等价：
\begin{enumerate}[i)]
\item 存在诺特环 $A_0$ (若愿意可假设它是 $A$ 的有限型子环)，有限型的 $A_0$-预概形 $X_0$，同态 $A_0 \to A$ 和 $A$-同构 $X \xrightarrow{\sim} X_0 \times_{A_0} A$.
\item 对角态射 $X \to X \times_{\mathrm{Spec}(A)} X$ 拟紧 (若 $X$ 在 $A$ 上分离则此条件为空的)，$X$ 是有限个仿射开集 $X_i$ 的并，其环 $B_i$ 是 $A$ 上有限表现的代数，i.e. 有限个未定元的多项式代数被有限型理想的商.
\end{enumerate}
若 $X$ 本身仿射，环为 $B$，这些条件简单意味着 $B$ 是 $A$ 上有限表现的代数.

态射 $f \colon X \to Y$ 称为\emph{有限表现态射}\index[term]{youxianbiaoxiantaishe@有限表现态射}，也说 $X$ 在 $Y$ 上有限表现，若 $Y$ 是仿射开集 $Y_i$ 的并，使 $X|_{Y_i}$ 作为 $Y_i$-预概形满足前面等价的条件. 则对 $Y$ 中每个仿射开集 $Y'$，$X|_{Y'}$ 亦然. 这是随基变换保持的性质，而且两个有限表现态射的复合有限表现.

这些概念确定后，由 (ii) 看到 (如 \ref{prop:VIII.1.10} 那样进行) 把其中的"有限型"换成"有限表现"，该断言仍然成立.
\end{remark}

\section{拓扑性质的下降}

\begin{theorem}\label{thm:VIII.4.1}
设 $g \colon Y' \to Y$ 是态射，$Z$ 是 $Y$ 的部分. 假设 $g$ 平坦，且存在拟紧态射 $f \colon X \to Y$ 使 $Z = f(X)$ (N.B. 若 $Y$ 诺特，这最后一个条件由"$Z$ 可构"蕴含). 则
\[
\overline{g^{-1}(Z)} = g^{-1}(\overline{Z}).
\]
\end{theorem}

可假设 $Y$ 仿射，再假设 $Y'$ 仿射. 由于 $Y$ 仿射，$X$ 是有限个仿射开集 $X_i$ 的并，把 $X$ 换成诸 $X_i$ 的和，可假设 $X$ 也仿射. 设 $A$，$A'$，$B$ 是 $Y$，$Y'$，$X$ 的环，$B' = B \otimes_A A'$ 是 $X' = X \times_Y Y'$ 的环，$I$ 是 $A \to B$ 的核，$I'$ 是 $A' \to B'$ 的核，故由这些理想定义的 $Y$ 和 $Y'$ 的闭部分分别是 $Z = f(X)$ 的闭包和 $Z' = f'(X') = g^{-1}(Z)$ 的闭包. 要证明后者等于 $g^{-1}(\overline{Z})$，这将由 $I' = IA'$ 得出，而后者又是 $A'$ 在 $A$ 上平坦的推论.

\begin{corollary}\label{cor:VIII.4.2}
设 $g \colon Y' \to Y$ 是平坦拟紧态射，$Z'$ 是 $Y'$ 的对于由 $g$ 定义的集合等价关系饱和的闭部分. 则 $Z' = g^{-1}(g(Z'))$.
\end{corollary}

确实有 $Z' = g^{-1}(Z)$，其中 $Z = g(Z')$. 可以应用 \ref{thm:VIII.4.1}，注意 \ref{thm:VIII.4.1} 中对 $Z$ 的条件确实满足，只需取 $X$ 为配备由 $Y'$ 诱导的既约结构的预概形 $Z'$. ($g$ 拟紧保证诱导态射 $f \colon Z' \to Y$ 拟紧).

断言 \ref{cor:VIII.4.2} 也意味着由 $Y$ 诱导的 $g(Y')$ 的拓扑是 $Y'$ 的拓扑的商. 特别：

\begin{corollary}\label{cor:VIII.4.3}
设 $g \colon Y' \to Y$ 是忠实平坦拟紧态射. 则 $g$ 使 $Y$ 成为 $Y'$ 的商拓扑空间，i.e. 对 $Y$ 的部分 $Z$，$Z$ 闭 (resp. 开) 当且仅当 $Z' = g^{-1}(Z)$ 如此.
\end{corollary}

回忆 $Y'$ 的两个元素 $a$，$b$ 在 $Y$ 中有相同的像，当且仅当它们形如 $p_1(c)$，$p_2(c)$，其中 $c$ 是 $Y'' = Y' \times_Y Y'$ 中适当的元素. 由此得出，若 $g$ 满射，有集合的精确图式
\[
P(Y) \longrightarrow P(Y') \rightrightarrows P(Y''),
\]
其中对每个集合 $E$，用 $P(E)$ 表示它的部分集. 这样，\ref{cor:VIII.4.3} 也可解释为：

\begin{corollary}\label{cor:VIII.4.4}
(开部分的下降 resp. 闭部分的下降) 设 $g \colon Y' \to Y$ 如 \ref{cor:VIII.4.3} 中. 对每个预概形 $X$，设 $\mathrm{Ouv}(X)$ (resp. $\mathrm{Fer}(X)$) 表示它的开部分集 (resp. 闭部分集). 则有集合映射的精确图式 (由 $g$ 和 $Y'' = Y' \times_Y Y'$ 的两个投影得到)：
\[
\mathrm{Ouv}(Y) \longrightarrow \mathrm{Ouv}(Y') \rightrightarrows \mathrm{Ouv}(Y''),
\]
\[
\mathrm{Fer}(Y) \longrightarrow \mathrm{Fer}(Y') \rightrightarrows \mathrm{Fer}(Y'').
\]
\end{corollary}

有下列对 \ref{cor:VIII.4.3} 的补充：

\begin{corollary}\label{cor:VIII.4.5}
设 $g \colon Y' \to Y$ 如 \ref{cor:VIII.4.3} 中，设 $Z$ 是 $Y$ 的以某个拟紧态射 $f \colon X \to Y$ 的像为 $Z$ 的部分 (例如 $Z$ 可构，$Y$ 诺特). 为使 $Z$ 是 $Y$ 的局部闭部分，必要且充分条件是 $Z' = g^{-1}(Z)$ 是 $Y'$ 的局部闭部分.
\end{corollary}

只需证明"充分"部分. 设 $Y_1$ 是 $Y$ 的闭子预概形，即 $Z$ 的闭包配备诱导既约结构，设 $Y'_1 = Y_1 \times_Y Y'$ 是 $Y'$ 的 $Y_1$ 的原像构成的闭子预概形. 它的底集是原像 $g^{-1}(Y_1) = g^{-1}(\overline{Z})$，由 \ref{thm:VIII.4.1} 等于 $Z'$. 由于 $Z'$ 在 $Y'$ 中局部闭，它在 $Z'$ 中开，故在 $Y'_1$ 中开. 而后者在 $Y_1$ 上忠实平坦且拟紧，由 \ref{cor:VIII.4.3} 得出 $Z$ 在 $Y_1$ 中开，i.e. 在 $\overline{Z}$ 中开，这意味着 $Z$ 局部闭.

\begin{corollary}\label{cor:VIII.4.6}
设 $g \colon S' \to S$ 是忠实平坦拟紧态射，$f \colon X \to Y$ 是 $S$-态射，$f' \colon X' \to Y'$ 是经基变换得到的 $S'$-态射. 假设 $f'$ 是开映射 (resp. 闭映射，resp. 拟紧且嵌入同胚，resp. 满同胚)；则 $f$ 有同样性质.
\end{corollary}

由于 $Y'$ 在 $Y$ 上忠实平坦且拟紧，可假设 $Y = S$. 设 $Q$ 是 $X$ 的部分，有 (用 $h$ 记投影态射 $X' \to X$)：
\[
g^{-1}(f(Q)) = f'(h^{-1}(Q)).
\]
若 $Q$ 开 (resp. 闭)，则 $h^{-1}(Q)$ 亦然，故假设 $f'$ 是开映射 (resp. 闭映射) 时 $f'(h^{-1}(Q))$ 亦然，由前面公式和 \ref{cor:VIII.4.3} 得 $f(Q)$ 亦然. 这证明 \ref{cor:VIII.4.6} 中前两个断言，剩下考察 $f'$ 是嵌入同胚的情形，证明 $f$ 是嵌入同胚. (满同胚的情形将由 \ref{prop:VIII.3.1} 得出). 由 \ref{cor:VIII.3.2} $f$ 单射，剩下证明映射 $X \to f(X)$ 开. 已知由 \ref{prop:VIII.3.3} $f$ 拟紧. 只需证明对 $X$ 的每个闭部分 $Z$，有 $Z = f^{-1}(\overline{f(Z)})$，这等价于满射映射 $h \colon X' \to X$ 的逆像的类似公式，i.e. 等价于
\[
Z' = f'^{-1}(g^{-1}(\overline{f(Z)})),
\]
其中令 $Z' = h^{-1}(Z)$. 而由 \ref{thm:VIII.4.1} 应用于 $Y$ 的部分 $f(Z)$，有 $g^{-1}(\overline{f(Z)}) = \overline{g^{-1}(f(Z))}$，待证公式等价于
\[
Z' = f'^{-1}(\overline{f'(Z')}),
\]
它由 $f'$ 是嵌入同胚的假设得出.

N.B. 在后一个推理中，已经假设 $f$ 拟紧，没有用到 $g$ 拟紧，只用 $g$ 忠实平坦. 故正是在这个假设下可以下降"嵌入同胚"或"满同胚"性质，或借助前面推理下降"$f'$ 拟紧且使 $f'(X')$ 成为 $X'$ 的商拓扑空间"的性质.

我们说预概形态射 $f \colon X \to Y$ 万有开 (resp. 万有闭，resp. 万有双连续，等等)，若对任何基变换 $Y' \to Y$，$f' \colon X' \to Y'$ 是开态射 (resp. 闭态射，resp. 到像空间的同胚). 由 \ref{cor:VIII.4.6} 得出：

\begin{corollary}\label{cor:VIII.4.7}
在 \ref{cor:VIII.4.6} 的条件下，为使 $f$ 万有开 (resp. 万有闭，resp. 万有嵌入同胚，resp. 万有同胚)，必要且充分条件是 $f'$ 如此.
\end{corollary}

\begin{corollary}\label{cor:VIII.4.8}
在 \ref{cor:VIII.4.6} 的条件下，为使 $f$ 分离 (resp. 紧合)，必要且充分条件是 $f'$ 如此.
\end{corollary}

说 $f$ 分离意味着对角态射 $X \to X \times_Y X$ 闭，或也万有闭，故 \ref{cor:VIII.4.8} 的第一个断言由 \ref{cor:VIII.4.7} 得出. 说 $f$ 紧合意味着 $f$ 满足条件 a) $f$ 有限型 b) $f$ 分离 c) $f$ 万有闭. 条件 a) 由 \ref{prop:VIII.3.3} 下降良好，b) 由刚才所述亦然，最后 c) 由 \ref{cor:VIII.4.7} 亦然.

\begin{remark}\label{rem:VIII.4.9}
回忆当 $g \colon Y' \to Y$ 是有限型平坦态射，$Y$ 局部诺特时，$g$ 是开态射 (\cref{prop:VI.6.6})，这是比 \ref{cor:VIII.4.3} 更精确的结果. 然而注意，若 $f$ 是诺特预概形的忠实平坦拟紧态射，$f$ 一般不一定是开态射. 例如设 $Y$ 是生成点 $y$ 不开的不可约概形 (例如代数曲线)，取 $Y'$ 为概形和 $Y \amalg \mathrm{Spec}(k(y))$，则开部分 $\mathrm{Spec}(k(y))$ 经结构态射 $Y' \to Y$ 的像不是 $Y$ 的开部分. 读者还会注意，本讲的各种断言若放弃所考虑的忠实平坦态射还拟紧的假设将变假，使断言失效的典型情形是把 $Y'$ 取为 $Y$ 诸点的局部环的谱的概形和. 例如，再取 $Y$ 为不可约代数曲线，$Z$ 为约化为生成点的部分，它在 $Y'$ 中的逆像开，而 $Z$ 不开.
\end{remark}

4.10. 本讲的各种断言中，把"$Y'$ 在 $Y$ 上平坦"的假设换成：存在 $Y'$ 上有限型、支集为 $Y'$、相对 $Y$ 平坦的模 $F$；忠实平坦假设则换成前述假设加上 $Y' \to Y$ 满射. 这适用于 \ref{prop:VIII.1.10} 中前两个断言，\ref{prop:VIII.3.3}，\ref{cor:VIII.3.5}，\ref{thm:VIII.4.1}，从而适用于本节所有结果.

\section{预概形态射的下降}

\begin{proposition}\label{prop:VIII.5.1}
设 $g \colon S' \to S$ 是预概形态射.
\begin{enumerate}[a)]
\item 假设 $g$ 满射，且同态
\[
g^{*} \colon \mathcal{O}_S \longrightarrow g_*(\mathcal{O}_{S'})
\]
单射，则 $g$ 是预概形范畴中的满态射，甚至在环化空间范畴中亦然.
\item 假设 $g$ 满射且使 $S$ 成为 $S'$ 的商拓扑空间. 设 $S'' = S' \times_S S'$，设 $h \colon S'' \to S$ 是结构态射，考虑典范同态图式
\[
\mathcal{O}_S \longrightarrow g_*(\mathcal{O}_{S'}) \rightrightarrows h_*(\mathcal{O}_{S''}),
\]
并假设该图式正合. 则 $g$ 是预概形范畴中的\emph{有效满态射} (也在环化空间范畴中)，i.e. 图式
\[
S \longleftarrow S' \leftleftarrows S''
\]
正合.
\end{enumerate}
\end{proposition}

\begin{proof}
a) 需证明环化空间态射 $f \colon S \to Z$ 在已知 $fg$ 时确定. 由于 $g$ 满射，已知 $f$ 的底集映射 $f_0$，剩下确定环层同态 $\mathcal{O}_Z \to \mathcal{O}_S$，或等价地由 $f$ 定义的环层同态
\[
u \colon f_0^{-1}(\mathcal{O}_Z) \longrightarrow \mathcal{O}_S.
\]
已知由 $fg$ 定义的同态
\[
(fg)_0^{-1}(\mathcal{O}_Z) = g_0^{-1}(f_0^{-1}(\mathcal{O}_Z)) \longrightarrow \mathcal{O}_{S'},
\]
或等价地，有一个同态
\[
f_0^{-1}(\mathcal{O}_Z) \longrightarrow g_{0*}(\mathcal{O}_{S'}) = g_*(\mathcal{O}_{S'}).
\]
立即看到后者正是 $g^{*} \colon \mathcal{O}_S \to g_*(\mathcal{O}_{S'})$ 与 $u$ 的复合，且由于 $g^{*}$ 单射，已知 $g^{*} \circ u$ 时 $u$ 确定. [N.B. 显然没有用到 $g \colon S' \to S$ 是预概形态射，断言对任何环化空间态射成立；同样的评注对 b) 成立，无论在环化空间范畴中，还是在以局部环为环的环化空间范畴中. 还注意若 $g$ 是预概形态射，未必满射，但 $g^{*} \colon \mathcal{O}_S \to g_*(\mathcal{O}_{S'})$ 单射，则对 $S$ 到概形 $Z$ 的两个态射 $f_1$，$f_2$，由 $f_1 g = f_2 g$ 得 $f_1 = f_2$；事实上，若 $I$ 是 $S$ 上定义 $f_1$，$f_2$ 的一致点子预概形的理想 (即 $Z \times Z$ 的对角子预概形经 $(f_1, f_2)$ 的原像)，看到 $I$ 包含在 $\mathrm{Ker}(g^{*})$ 中].

b) 需证明对每个环化空间 $Z$，下列映射图式
\[
\mathrm{Hom}(S, Z) \longrightarrow \mathrm{Hom}(S', Z) \rightrightarrows \mathrm{Hom}(S'', Z)
\]
正合，且当 $Z$ 是以局部环为环的环化空间、限于以局部环为环的环化空间同态时亦然. 由 a) 已知第一个映射单射，剩下看：若 $f' \colon S' \to Z$ 是环化空间同态且 $f' p_1 = f' p_2$，则 $f'$ 形如 $f g$，其中 $f \colon S \to Z$ 是环化空间同态. 由于 $g$ 满射，显然若 $f'$ 是局部环的环化空间态射，则 $f$ 亦然.

由对 $f'$ 的假设得出底集映射 $f'_0$ 在映射 $g_0$ 的诸纤维上为常值，故由于后者满射，$f'_0$ 唯一分解为 $f'_0 = f_0 g_0$，其中 $f_0 \colon S \to Z$ 是映射，它必然连续，因为 $g_0$ 把 $S$ 等同于 $S'$ 的商拓扑空间. 现在考虑由对应 $f'$ 的同态 $(f_0 g_0)^{-1}(\mathcal{O}_Z) \to \mathcal{O}_{S'}$ 得到的同态
\[
f_0^{-1}(\mathcal{O}_Z) \longrightarrow g_*(\mathcal{O}_{S'}).
\]
假设 $f' p_1 = f' p_2$ 解释为：前述同态与两个同态 $g_*(\mathcal{O}_{S'}) \rightrightarrows h_*(\mathcal{O}_{S''})$ 的复合相同，故由假设 b) 它经一个态射
\[
f_0^{-1}(\mathcal{O}_Z) \longrightarrow \mathcal{O}_S
\]
分解. 后者定义环化空间态射 $f \colon S \to Z$，即所求态射.
\end{proof}

\begin{theorem}\label{thm:VIII.5.2}
设 $\mathcal{F}$ 是预概形范畴 $\mathbf{Sch}$ 中箭的纤维范畴 (VI 11.a). 则任何忠实平坦拟紧态射 $g \colon S' \to S$ 是 $\mathcal{F}$-下降态射 (或如常所说的，$\mathbf{Sch}$ 中的下降态射).
\end{theorem}

这即意味着：设 $S'' = S' \times_S S'$，对 $S$ 上的两个预概形 $X$，$Y$，考虑它们在 $S'$ 上的逆像 $X'$，$Y'$ 及在 $S''$ 上的逆像 $X''$，$Y''$，得映射图式
\[
\mathrm{Hom}_S(X, Y) \longrightarrow \mathrm{Hom}_{S'}(X', Y') \rightrightarrows \mathrm{Hom}_{S''}(X'', Y'');
\]
记号确定后，\ref{thm:VIII.5.2} 意味着该图式正合. 注意一般并不成立 $g$ 是有效下降态射，i.e. 对 $S'$ 上的每个预概形 $X'$，$X'$ 上相对 $g \colon S' \to S$ 的每个下降数据有效. 有效性的问题 (常常微妙) 将在第 7 号考察.

已见 \cite{exp08-D} (考虑到 $\mathbf{Sch}$ 中纤维积存在) 断言 \ref{thm:VIII.5.2} 等价于：

\begin{corollary}\label{cor:VIII.5.3}
预概形的忠实平坦拟紧态射是万有有效满态射.
\end{corollary}

由于忠实平坦拟紧态射经任何基扩张后仍如此，归结到证明它是有效满态射\index[term]{youxiaomantaishe@有效满态射}. 应用判则 \ref{prop:VIII.5.1} b)，它考虑到 \ref{cor:VIII.4.3} 和 \ref{cor:VIII.1.7} 给出所求结果.

\begin{corollary}\label{cor:VIII.5.4}
设 $g \colon S' \to S$ 是忠实平坦拟紧态射，$f \colon X \to Y$ 是 $S$-态射，$f' \colon X' \to Y'$ 是经基变换 $S' \to S$ 得到的 $S'$-态射. 为使 $f$ 是同构，必要且充分条件是 $f'$ 如此.
\end{corollary}

事实上，若 $f'$ 是同构，它也是 $X'$，$Y'$ 上自然下降结构的同构，且由于函子 $X \mapsto X'$ (从 $\mathbf{Sch}_{/S}$ 到配备相对 $g$ 的下降数据的 $S'$ 上预概形范畴) 由 \ref{thm:VIII.5.2} 忠实全，得到所求结论.

\begin{corollary}\label{cor:VIII.5.5}
在 \ref{cor:VIII.5.4} 的条件下，为使 $f$ 是闭浸入 (resp. 开浸入，resp. 拟紧浸入)，必要且充分条件是 $f'$ 如此.
\end{corollary}

如常可假设 $Y = S$，只需证明"充分"部分. 注意 $X'/Y'$ 配备相对 $g \colon Y' \to Y$ 的下降数据，且结构态射 $f' \colon X' \to Y'$ 是浸入故是单态射，这蕴含 $Y''$ 的由 $S''$ 到 $S'$ 的两个投影之一另一个对 $X'/Y'$ 取逆像得到的两个子对象相同. 若 $f'$ 是闭浸入，由 \ref{cor:VIII.1.9} 得出存在 $Y$ 的闭子预概形 $X_1$，其经 $g \colon Y' \to Y$ 的逆像是 $X'$. 故由 $\mathcal{F}$-下降态射相对下降问题的解的唯一性，得出 $X_1$ 与 $X$ $Y$-同构，故 $f \colon X \to Y$ 是闭浸入. 对开浸入用 \ref{cor:VIII.4.4} 同样进行. 最后若 $f'$ 是拟紧浸入，由 \ref{prop:VIII.3.3} $f$ 拟紧，可以对 $Y$ 的部分 $f(X)$ 应用判则 \ref{cor:VIII.4.5}，它证明 $f(X)$ 局部闭，因为它在 $Y'$ 中的逆像 $f'(X')$ 如此. 再把 $Y$ 换成 $f(X)$ 在其中闭的开部分，归结到 $f'$ 是闭浸入的情形，故由前面所述 $f$ 如此.

\begin{corollary}\label{cor:VIII.5.6}
在 \ref{cor:VIII.5.4} 的条件下，为使 $f$ 仿射，必要且充分条件是 $f'$ 如此.
\end{corollary}

如 \ref{cor:VIII.5.5} 那样进行，用 \ref{thm:VIII.2.1} (也可用 Serre 的上同调判则 [EGA II 5.2]，它不借助下降技巧证明 \ref{cor:VIII.5.6}).

\begin{corollary}\label{cor:VIII.5.7}
在 \ref{cor:VIII.5.4} 的条件下，为使 $f$ 整 (resp. 有限，resp. 有限且局部自由)，必要且充分条件是 $f'$ 如此.
\end{corollary}

只需证明"充分"部分，如常可假设 $Y = S$，$Y$ 仿射，$Y'$ 仿射. 由于假设蕴含 $f'$ 仿射，由 \ref{cor:VIII.5.6} $f$ 也仿射，故 $X$ 从而 $X'$ 仿射. 设 $A$，$A'$，$B$，$B' = B \otimes_A A'$ 是 $Y$，$Y'$，$X$，$X'$ 的环. 有 $B = \varinjlim_i B_i$，其中 $B_i$ 遍历 $B$ 的在 $A$ 上有限型的子 $A$-代数，从而 $B' = \varinjlim_i B'_i$，其中 $B'_i$ 是 $A'$-代数 $B'$ 的有限型子代数. 若 $B'$ 在 $A'$ 上整，诸 $B'_i$ 是 $A'$ 上有限型模，故 $A'$ 在 $A$ 上忠实平坦时，诸 $B_i$ 是 $A$ 上有限型模，i.e. $B$ 在 $A$ 上整. 同样看到若 $B'$ 在 $A'$ 上有限，则 $B$ 在 $A$ 上有限. 对"有限型局部自由"同样结论，参见 1.11.

\begin{corollary}\label{cor:VIII.5.8}
在 \ref{cor:VIII.5.4} 的条件下，假设 $f$ 拟紧，设 $\mathcal{L}$ 是 $X$ 上的可逆模，$\mathcal{L}'$ 是它在 $X'$ 上的逆像. 为使 $\mathcal{L}$ 相对 $f$ 丰富 (resp. 非常丰富)，必要且充分条件是 $\mathcal{L}'$ 相对 $f'$ 丰富 (resp. 非常丰富).
\end{corollary}

只需证明"充分"部分. 对 $\mathcal{L}$ 的假设无论如何蕴含 $f'$ 分离，故由 \ref{cor:VIII.4.8} $f$ 分离，且由于 $f$ 拟紧，$g \colon Y' \to Y$ 平坦，经仿射覆盖计算直像表明有同构
\[
g^{*}(f_*(\mathcal{L}^{\otimes n})) \xrightarrow{\sim} f'_*(\mathcal{L}'^{\otimes n})
\]
对每个整数 $n$，故有同构
\[
g^{*}(\mathcal{S}) \xrightarrow{\sim} \mathcal{S}',
\]
其中 $\mathcal{S}$ (resp. $\mathcal{S}'$) 表示 $Y$ (resp. $Y'$) 上的拟凝聚分次代数，即诸 $f_*(\mathcal{L}^{\otimes n})$ (resp. 诸 $f'_*(\mathcal{L}'^{\otimes n})$) 对 $n \geqslant 0$ 的直和. 注意对每个 $n \geqslant 0$，典范同态 $f'^{*}(\mathcal{S}'_n) \to \mathcal{L}'^{\otimes n}$ 的余核是 $f^{*}(\mathcal{S}_n) \to \mathcal{L}^{\otimes n}$ 的余核经 $X' \to X$ 的原像，故其支集 $Z'_n$ 是支集 $Z_n$ 的原像. 若 $\mathcal{L}'$ 丰富，诸 $Z'_n$ 的交为空，故由于 $X' \to X$ 满射，诸 $Z_n$ 的交为空，i.e. 有典范态射
\[
j \colon X \longrightarrow \mathrm{Proj}(\mathcal{S})
\]
(EGA II 3). 而且，类似态射
\[
j' \colon X' \longrightarrow \mathrm{Proj}(\mathcal{S}')
\]
正是由前者经基变换 $Y' \to Y$ 得到的 (loc. cit.). 这样，说 $\mathcal{L}'$ 相对 $f'$ 丰富意味着 $j'$ 是浸入，而且必然拟紧，因为 $f'$ 拟紧. 故由 \ref{cor:VIII.5.5} $j$ 是浸入，i.e. $\mathcal{L}$ 相对 $f$ 丰富. — 对"非常丰富"的情形完全类似地进行，上面只需取 $n = 1$，并把 $\mathrm{Proj}(\mathcal{S})$ 的考虑换成由 $\mathcal{S}_1$ 对应的射影纤维丛 $P(\mathcal{S}_1)$ 的考虑.

回忆 (EGA II 5.1.1) 拟紧态射 $f$ 称为拟仿射\index[term]{nifangshetaishe@拟仿射态射}，若对 $Y$ 中每个仿射开集 $U$，$f^{-1}(U)$ 是同构于仿射概形的开子概形的预概形. 证明 (loc. cit.) 这等价于说 $\mathcal{O}_X$ 相对 $f$ 丰富 (或也说：非常丰富). 故 \ref{cor:VIII.5.8} 蕴含：

\begin{corollary}\label{cor:VIII.5.9}
在 \ref{cor:VIII.5.4} 的条件下，且假设 $f$ 拟紧，为使 $f$ 拟仿射，必要且充分条件是 $f'$ 如此.
\end{corollary}

\begin{remark}\label{rem:VIII.5.10}
Hironaka 的非射影簇例子表明，可以有非奇异代数簇 (其中 $Y$ 射影) 的紧合态射 $f \colon X \to Y$，使 $Y$ 是两个开集 $Y_i$ 的并，$X_i = X \times_Y Y_i$ 在 $Y_i$ 上射影，但 $f$ 不射影. 故令 $Y' = Y_1 \amalg Y_2$，$Y'$ 在 $Y$ 上忠实平坦且拟紧 (甚至拟有限)，$f' \colon X' \to Y'$ 射影，但 $f$ 不射影. 因此注意，为应用 \ref{cor:VIII.5.8}，从 $f'$ 射影推出 $f$ 的同样结论，需要 $X'$ 上已经有相对 $f'$ 丰富的可逆模 $\mathcal{L}'$，配备相对 $X' \to X$ 的下降数据 (这使得可以把 $\mathcal{L}'$ 看作 $X$ 上可逆模 $\mathcal{L}$ 的逆像，后者将由 \ref{cor:VIII.5.8} 相对 $f$ 丰富). 当 $g \colon S' \to S$ 有限且局部自由时，参见 7.7.
\end{remark}

\section{应用于有限与拟有限态射}

我们要证明下列两个定理：

\begin{theorem}\label{thm:VIII.6.1}
设 $f \colon X \to Y$ 是纤维有限紧合态射\footnote{参见 EGA IV 18.12，关于对未必局部诺特预概形的推广.}，其中 $Y$ 局部诺特. 则 $f$ 有限.
\end{theorem}

\begin{theorem}\label{thm:VIII.6.2}
设 $f \colon X \to Y$ 是拟有限分离态射，其中 $Y$ 局部诺特. 则 $f$ 拟仿射，更不用说拟射影.
\end{theorem}

\begin{remark}\label{rem:VIII.6.3}
\cref{thm:VIII.6.1} 众所周知，在代数簇情形归功于 Chevalley；也见 [EGA III 4] 中用"\emph{全纯函数}定理"的简单证明. 这里给出的证明不用后者，而用下降理论；我们把它作为赠品给读者，因为得到它与 \ref{thm:VIII.6.2} 的证明"不费事". 也回忆 ([EGA III 4] 或 \cite{exp08-1}) Zariski "主定理"的全局形式，由"全纯函数定理"推出，断言：若 $f \colon X \to Y$ 拟有限且拟射影，$Y$ 诺特，则 $X$ 与有限 $Y$-预概形 $Z$ 的开子预概形 $Y$-同构. "主定理"与 \ref{thm:VIII.6.2} 的合取于是表述为：
\end{remark}

\begin{corollary}\label{cor:VIII.6.4}
设 $f \colon X \to Y$ 是拟有限分离态射，$Y$ 诺特. 则 $X$ 与有限 $Y$-预概形 $Z$ 的开子预概形 $Y$-同构.
\end{corollary}

\ref{thm:VIII.6.2} 对下降理论的另一个有趣推论将随 \ref{cor:VIII.7.9} 给出.

\begin{proof}
[6.1 与 6.2 的证明] 我们承认下列事实，其证明容易\footnote{参见 EGA IV 8.}：
\end{proof}

\begin{lemma}\label{lem:VIII.6.5}
设 $X$ 是局部诺特 $Y$ 上有限型的预概形，$y \in Y$. 为使存在 $y$ 的开邻域 $U$ 使 $X|_U$ 在 $U$ 上有限 (resp. 拟仿射，resp. ...)，必要且充分条件是 $X \times_Y \mathrm{Spec}(\mathcal{O}_y)$ 在 $\mathrm{Spec}(\mathcal{O}_y)$ 上有限 (resp. 拟仿射，resp. ...).
\end{lemma}

另一方面，由于 $f \colon X \to Y$ 有限 (resp. 拟仿射) 的性质在 $Y$ 上局部，为证 \ref{thm:VIII.6.1} 和 \ref{thm:VIII.6.2} 归结到 $Y$ 是局部环的谱，从而尤其有限维的情形. (N.B. 预概形 $Y$ 的维数定义为它的局部环的 Krull 维数的上确界). 对
\[
n = \dim(Y)
\]
作递归，$n < 0$ 时断言平凡. 故可假设 $n \geqslant 0$，且对维数 $n' < n$ 已证. 又可假设 $Y$ 是维数 $n$ 的局部诺特环 $A$ 的谱. 注意 \ref{thm:VIII.6.1} 和 \ref{thm:VIII.6.2} 中的假设随基变换保持 (在开头的归结中已用到)，它们在基变换 $\mathrm{Spec}(\widehat{A}) \to \mathrm{Spec}(A)$ 后仍然成立. 由于后者忠实平坦且拟紧，断言 \ref{cor:VIII.5.7} 和 \ref{cor:VIII.5.9} 把我们归结到 $A$ 还完备的情形. 利用事实：每个在 $A$ 上拟有限的诺特局部环 $B$ 在 $A$ 上有限，以及 $X$ 在 $Y$ 上分离且 $y$ 的纤维由孤立点组成，得到分解
\[
X = X' \amalg X''
\]
其中 $X'$ 在 $Y$ 上有限，$X''$ 在 $y$ 处的纤维为空. 若 $X$ 在 $Y$ 上紧合，则 $X''$ 亦然，故它在 $Y$ 中的像闭，且由于它不含 $y$，它为空，故 $X'' = \varnothing$ 从而 $X = X'$，这表明 $X$ 在 $Y$ 上有限并证明 \ref{thm:VIII.6.1} (N.B. 这里递归假设无用). 若 $X$ 在 $Y$ 上拟有限，$X''$ 也如此，而 $X''$ 实际位于 $Y$ 的开集 $Y - \{y\}$ 上，其维数 $< n$. 由递归假设，$X''$ 在 $Y - \{y\}$ 上拟仿射，故在 $Y$ 上也拟仿射，$X'$ 显然亦然，故它们的和 $X$ 亦然，这证明 6.2.

\begin{remark}\label{rem:VIII.6.6}
\cref{thm:VIII.6.1,thm:VIII.6.2} 在不再假设 $Y$ 局部诺特时仍然成立，只要明确假设 $f$ 有限表现 (参见 \ref{rem:VIII.3.6}). 事实上，仍可假设 $Y$ 仿射，此时容易验证情形 $f \colon X \to Y$ 由基变换 $Y \to Y_0$ 自满足与 $f$ 同样假设、$Y_0$ 诺特的情形 $f_0 \colon X_0 \to Y_0$ 得到. 故由结果 \ref{thm:VIII.6.1} resp. \ref{thm:VIII.6.2}，$f_0$ 有限 resp. 拟仿射，故 $f$ 亦然. 这类推理常在应用中有用于摆脱诺特假设 (它在应用中总是碍事).
\end{remark}

\section{下降数据的有效性判据}

如常考虑预概形态射
\[
g \colon S' \longrightarrow S
\]
和 $S'$-预概形 $X'$. 按一般事实 VII，9\footnote{不存在；参见 Séminaire Bourbaki 的 212 号报告 (1962).}，$X'$ 上相对 $g$ 的下降数据的给出，等价于等价对 \cite{exp08-3}
\[
q_1, q_2 \colon X'' \rightrightarrows X'
\]
的给出，使结构态射 $X' \to S'$ 与该对及由 $g$ 定义的等价对
\[
p_1, p_2 \colon S'' = S' \times_S S' \rightrightarrows S'
\]
相容，且使对应图式
\[
\begin{tikzcd}[row sep=large, column sep=large]
X' \arrow[d] & X'' \arrow[l] \arrow[d] \\
S' & S'' \arrow[l]
\end{tikzcd}
\]
中提取的两个方形 (或两个之一，由对称性两者等价) (用 $p_1$，$q_1$ 或 $p_2$，$q_2$) 笛卡儿. 该下降数据提出的下降问题的解，i.e. $S$ 上配备与下降数据相容的同构 $X \times_S S' \xleftarrow{\sim} X'$ 的对象 $X$，等价于满足 $h q_1 = h q_2$ 的笛卡儿方形
\[
\begin{tikzcd}[row sep=large, column sep=large]
X \arrow[d] & X' \arrow[l,"h"] \arrow[d] \\
S & S' \arrow[l,"g"]
\end{tikzcd}
\]
的给出.

由于忠实平坦拟紧态射的集随基变换稳定，且忠实平坦拟紧态射由 \ref{cor:VIII.5.3} 是有效满态射，由一般理论 \cite{exp08-D} 得出：

\begin{proposition}\label{prop:VIII.7.1}
假设 $g \colon S' \to S$ 忠实平坦且拟紧. 为使 $X'$ 上相对 $g$ 的下降数据有效，必要且充分条件是它定义的等价关系 $R = (q_1, q_2)$ 有效 (i.e. 商 $X'/R$ 存在且 $X''$ 成为 $X'$ 在 $X'/R$ 上的纤维积方形)，且典范态射 $X' \to X'/R$ 忠实平坦且拟紧.
\end{proposition}

这样，下降数据的有效性问题是一个等价图的有效性问题的特殊情形，本节给出的各种有效性判则可由此得到. 尽管如此，在下降的语境中有\cref{thm:VIII.2.1}，它蕴含若 $X'$ 在 $S'$ 上仿射，则 $X'$ 上相对 $g$ 的每个下降数据有效，这个断言对由一般平坦等价图取商没有类似物. 我们这里给出的所有有效性判则也可看作由前者推出.

设 $U'$ 是 $X'$ 的子预概形 (或更一般地 $X'$ 在范畴 $\mathbf{Sch}$ 中的子对象)；称 $U'$ 被 $X'$ 上的下降数据稳定，若可以在 $U'$ 上找到相对 $g$ 的下降数据，使浸入 $U' \to X'$ 与下降数据相容. 这也意味着 $U'$ 在 $X''$ 中经 $q_1$ 和 $q_2$ 的逆像相同 (或如常所说的，$U'$ 被等价关系 $R$ 稳定)，当然 $U'$ 上的所求下降数据于是唯一，称为由 $X'$ 的下降数据诱导的下降数据. 这样：

\begin{proposition}\label{prop:VIII.7.2}
设 $(X'_i)$ 是 $X'$ 的被下降数据稳定的开集 $X'_i$ 的覆盖. 为使 $X'$ 上的下降数据有效，必要且充分条件是诸 $X'_i$ 中诱导的下降数据有效.
\end{proposition}

这是例如 \ref{prop:VIII.7.1} 的容易推论，证明细节留给读者.

\begin{corollary}\label{cor:VIII.7.3}
设 $(S_i)$ 是 $S$ 的开覆盖，对每个 $i$ 设 $S'_i$ 和 $X'_i$ 由 $S'$ 和 $X'$ 经基变换 $S_i \to S$ 得到. 为使 $X'$ 上的下降数据有效，必要且充分条件是：对每个 $i$，$X'_i$ 上相对 $g_i \colon S'_i \to S_i$ 的下降数据有效.
\end{corollary}

这个判则实际上总把我们归结到 $S$ 仿射的情形. 在 $S'$ 也仿射的情形 (应用中最常见)，有：

\begin{corollary}\label{cor:VIII.7.4}
假设 $S$ 和 $S'$ 仿射. 为使 $X'$ 上的下降数据有效，必要且充分条件是 $X'$ 是被下降数据稳定的仿射开集 $X'_i$ 的并.
\end{corollary}

充分性来自 \ref{prop:VIII.7.2} 及事实：若 $X'_i$ 仿射，它在 $S'$ 上仿射，可应用 2.1. 对必要性，注意若 $X'$ 来自 $X$，且 $X$ 被仿射开集 $X_i$ 覆盖，则诸 $X'_i = X_i \times_S S'$ 是被下降数据稳定且覆盖 $X'$ 的仿射开集.

\begin{corollary}\label{cor:VIII.7.5}
设 $g \colon S' \to S$ 是忠实平坦、拟紧且根式的态射. 则 $g$ 是有效下降态射，i.e. 对 $S'$ 上的每个 $X'$，$X'$ 上相对 $g \colon S' \to S$ 的每个下降数据有效.
\end{corollary}

事实上，由 \ref{cor:VIII.7.3} 可假设 $S$ 仿射，故由于 $S'$ 在 $S$ 上根式故分离，$S'$ 分离. 而且对每个 $x' \in X'$，由等价关系 $R$ 定义的集合等价关系的纤维 $R(x') = q_2(q_1^{-1}(x'))$ 约化为一点，因为 $g$ 根式时，由它经基变换 $S' \to S$ 得到的 $p_1$，$p_2$ 也根式，故由前者经基变换 $X' \to S''$ 得到的 $q_1$，$q_2$ 也根式. 故 $X'$ 的每个开集被下降数据稳定. 用仿射开集 $X'_i$ 覆盖 $X'$，由于 $S'$ 分离它们在 $S$ 上仿射，故由 \ref{thm:VIII.2.1} 诱导的下降数据有效. 再由 \ref{prop:VIII.7.2} 得出结论.

注意 \ref{cor:VIII.7.5} 给出预概形范畴中已知的有效下降态射的唯一情形，而且这很可能确实是唯一情形，即使限于诺特概形或域上有限型概形. 当假设 $S$ 局部诺特且 $S'$ 在 $S$ 上有限型时，断言 \ref{cor:VIII.7.5} 也是下列 (推广 Weil 的 Galois 下降和 Cartier 的不可分下降) 的特殊情形：

\begin{corollary}\label{cor:VIII.7.6}
设 $g \colon S' \to S$ 是有限局部自由 (i.e. 由 $S$ 上为有限型局部自由模的代数定义) 且满射 (故 $g$ 忠实平坦且拟紧，从而是下降态射) 的态射. 设 $X'$ 是配备下降数据的 $S'$-预概形. 为使该数据有效，必要且充分条件是：对每个 $x' \in X'$，纤维 $R(x') = q_2(q_1^{-1}(x'))$ 包含在某个仿射开集中 (若 $X'$ 在 $S'$ 上拟射影，该条件自动满足).
\end{corollary}

括号中的评注来自事实：若 $s$ 是 $x'$ 下方的 $S$ 的点，则 $R(x')$ 有限且包含在纤维 $X'_s$ 中，另一方面由于 $X'$ 在 $S'$ 上拟射影且 $S'$ 在 $S$ 上有限，$X'$ 在 $S$ 上拟射影，这意味着 $X'$ 在 $S$ 上的纤维包含在仿射开集中.

由于仿射概形的任何有限部分有仿射邻域基本系，看到限制到 $S$ 的某个仿射开集上方不损失假设，由 \ref{cor:VIII.7.3} 这归结到 $S$ 仿射的情形. 由 \ref{cor:VIII.7.4}，归结到证明 $x'$ 包含在被下降数据稳定的仿射开集中. 设 $U$ 是包含 $R(x')$ 的仿射开集，则饱和集
\[
R(X' - U) = q_2(q_1^{-1}(X' - U))
\]
与 $R(x')$ 不相交，另一方面由于 $q_2$ 有限 (因为 $g$ 故 $p_2$ 有限) 故闭，第二个成员是 $X'$ 的闭部分. 设 $U'$ 是它在 $X'$ 中的余集，它是饱和开集，且
\[
R(x') \subset U' \subset U
\]
其中 $U$ 仿射，但 $U'$ 事先未必仿射. 由于仿射概形 $U$ 中的有限部分 $R(x')$ 有 $U_f$ 形的仿射邻域基本系，把 $f$ 换成它在 $U'$ 上的限制，看到存在 $\mathcal{O}_U$ 的截面 $f$ 使：
\[
R(x') \subset U'_f, \quad U'_f \text{ 仿射}.
\]
设 $U'' = q_1^{-1}(U') = q_2^{-1}(U')$，仍用 $q_1$，$q_2$ 记诱导态射 $U'' \to U'$，考虑
\[
f' = N_{q_2}(q_1^{*}(f)),
\]
其中 $N_{q_2}$ 表示相对有限局部自由态射 $q_2 \colon U'' \to U'$ 的范数. 范数构造与基变换的相容性容易蕴含 $f'$ 是不变截面：
\[
q_1^{*}(f') = q_2^{*}(f')
\]
这蕴含 $U'_{f'}$ 是 $U'$ 的饱和开集. 而且更精确地，用 $Z(f')$ 记截面 $f'$ 的零点集，由范数的性质得到：
\[
Z(f') = q_2(Z(q_1^{*}(f))) = q_2(q_1^{*}(Z(f))) = R(U' - U'_f),
\]
这蕴含 $U'_{f'} = U' - Z(f')$ 饱和，包含 $R(x')$，且包含在 $U'_f$ 中. 由于后者仿射，得出 $U'_{f'}$ 也仿射 (因为它等于 $(U'_f)_{f''}$，其中 $f'' = f'|_{U'_{f'}}$). 故它是包含 $R(x')$ 从而 $x'$ 的仿射饱和开集，这完成证明.

注意这个推理适用于预概形 $X'$ 中有限局部自由的任何等价关系 (或甚至只是\emph{预等价}\index[term]{yudengjia@预等价}，参见 \cite{exp08-3})，而且 \ref{cor:VIII.7.6} 也是有限局部自由预等价的类似结果 (loc. cit.) 的特殊情形. 对下面的 \ref{cor:VIII.7.7} 同样评注. 也可以在得到包含 $x'$ 的仿射 (拟仿射) 饱和开集 $U'$ 的存在后，借助 \ref{cor:VIII.7.9} 和 \ref{prop:VIII.7.2}，这避免范数的使用. 还注意在 \ref{cor:VIII.7.6} 的条件下，若 $X'$ 上的下降数据有效，$X'$ 来自 $S$ 上的 $X$，则态射 $X' \to X$ 有限、局部自由且满射，因为它由 $g$ 经基变换 $X \to S$ 得到. 由此得出 (EGA II 6.6.4) 若 $X'$ 在 $S'$ 上故在 $S$ 上拟射影，则 $X$ 在 $S$ 上拟射影 (取 $X$ 上相对丰富可逆层为 $X'$ 上相对丰富 (于 $S$，或于 $S'$，两者相同) 可逆层的范数). 由此得到：

\begin{corollary}\label{cor:VIII.7.7}
有限局部自由且满射的态射 $g \colon S' \to S$ 对互相拟射影的预概形的纤维范畴是有效下降态射，i.e. 对每个在 $S'$ 上拟射影的 $X'$，$X'$ 上相对 $g$ 的每个下降数据有效，且下降得到的 $S$-预概形 $X$ 在 $S$ 上拟射影.
\end{corollary}

\begin{proposition}\label{prop:VIII.7.8}
设 $g \colon S' \to S$ 是忠实平坦拟紧态射. 则 $g$ 对预概形 $T$ 上配备相对 $T$ 丰富可逆层的拟紧预概形 $Z$ 的纤维范畴是有效下降态射. 特别，对 $S'$ 上的每个预概形 $X'$，配备相对 $g \colon S' \to S$ 的下降数据，以及对 $X'$ 上相对 $S'$ 丰富的每个可逆层 $\mathcal{L}'$，配备相对 $X'$ 上给定数据的下降数据，(i.e. 配备满足通常传递条件的 $q_1^{*}(\mathcal{L}')$ 与 $q_2^{*}(\mathcal{L}')$ 的同构)，$X'$ 上的下降数据有效，且由下降自 $\mathcal{L}'$ 得到的下降预概形 $X$ 上的可逆层 $\mathcal{L}$ 相对 $S$ 丰富.
\end{proposition}

证明与 \ref{cor:VIII.5.8} 完全类似，注意由 $\mathcal{L}'$ 定义的 $S'$ 上拟凝聚分次代数 $\mathcal{S}'$ 上有下降数据，使得可以借助 \ref{thm:VIII.1.1} 构造 $S$ 上拟凝聚分次代数 $\mathcal{S}$，从而得到 $S$ 上的 $P = \mathrm{Proj}(\mathcal{S})$，使 $P' = \mathrm{Proj}(\mathcal{S}')$ 连同其下降数据等同于 $P \times_S S'$. 由于由假设 $X'$ 等同于 $P'$ 的开集 (必然被 $P'$ 上的下降数据稳定)，$X'$ 上的下降数据也有效，得到的下降预概形是 $P$ 中的开集. 细节留给读者. — 特别取 $\mathcal{L}' = \mathcal{O}_{X'}$，得到：

\begin{corollary}\label{cor:VIII.7.9}
设 $g \colon S' \to S$ 是忠实平坦拟紧态射，设 $X'$ 是 $S'$ 上方的拟仿射预概形，则 $X'$ 上相对 $g$ 的每个下降数据有效，且下降得到的预概形 $X$ 在 $S$ 上拟仿射.
\end{corollary}

由 \ref{thm:VIII.6.2}，这个结果特别适用于 $S'$ 局部诺特且 $X'$ 在 $S'$ 上拟有限分离的情形，更一般地适用于 $S'$ 任意且 $X'$ 在 $S'$ 上有限表现、拟有限分离的情形 (参见 \ref{rem:VIII.6.6}).

\begin{remark}\label{rem:VIII.7.10}
本节给出的结果穷尽了目前已知的有效性判则，很可能甚至穷尽了所有有用的判则\footnote{这种看法部分被证明有误，参见 J.P. Murre，Séminaire Bourbaki 294 (附录)，1965年5月，以及一些特殊结果 (特别是 Néron 和 Raynaud 的). 关于群概形的下降，参见 M. Raynaud 的博士论文 (1968).}. 支持这一断言的注意下列反例：
\begin{enumerate}[i)]
\item 若 $S$ 是域的谱，$S'$ 是二次 伽罗瓦扩张的谱，可以找到 $S'$ 上紧合且光滑、维数 3 的 $X'$，配备不有效的下降数据 (Serre).
\item 可以找到 $S$ 为维数 3 的正则局部环的谱 (若愿意，可以取给定特征域上代数概形的局部环)，$T$ 是 $S$ 的群 $Z/2Z$ 的主覆盖，使若 $t$ 表示 $T$ 的位于 $S$ 的闭点 $s$ 上方的点之一，且 $S' = T - s$，则能找到 $S'$ 上射影且正则的 $X'$，配备相对 $g \colon S' \to S$ 的下降数据，该下降数据不有效.
\end{enumerate}
这些构造使用 Hironaka 的非射影簇例子. 对 (i)，只需用事实：可以在 $k$ 上方找到维数 3 的紧合光滑概形 $X_0$，$G = Z/2Z$ 无惯性地作用其上，其中有 $S$ 上两个有理点 $a$，$b$，在 $G$ 下相伴，不包含在仿射开集中. 令 $X' = X_0 \times_k k'$，借助 $G$ 在两个因子上的作用使 $G$ 作用在 $X'$ 上，这给出 $X'$ 上相对 $g \colon \mathrm{Spec}(k') \to \mathrm{Spec}(k)$ 的下降数据. 在 $a$ resp. $b$ 上方恰好有一个点 $a'$ resp. $b'$ (带二次剩余扩张)，且 $a'$ 和 $b'$ 在 $G$ 下相伴，因为 $X' \to X_0$ 与 $G$ 的作用相容. 则 $a'$ 和 $b'$ 不能包含在仿射开集 $U'$ 中，因为此时 $U = X_0 - \mathrm{Im}(X' - U')$ 将是包含 $(a, b)$ 且逆像在 $X'$ 中位于 $U'$ 内的开集，故拟仿射，从而 $U$ 拟仿射，于是 $(a, b)$ 在 $U$ 中有仿射邻域.

对 (ii)，用事实：在 Hironaka 的例子中，$X_0$ 作为 $k$-射影概形 $Y$ ($k$ 上光滑，态射 $f \colon X_0 \to Y$ 双有理，但无关紧要) 上方的紧合预概形得到，群 $G$ 也作用在 $Y$ 上，与它在 $X_0$ 上的作用相容，最后令 $S' = Y - f(b)$，$X' = X_0|_{S'}$，$X'$ 在 $S'$ 上射影. 则 $X_0$ 配备相对典范态射 $Y \to S = Y/G$ 的自然下降数据，借助 $G$ 在 $X_0$ 上与它在 $Y$ 上的作用相容的作用. 该下降数据不有效，因为 $(a, b)$ 不包含在仿射开集中. $X'$ 上相对 $g \colon S' \to S$ 的诱导下降数据于是不有效，容易验证.
\end{remark}

\section{文献}\printbibliography[heading=none]


\end{refsection}
\begin{refsection}% =====================================================================
% 第九讲 (Exposé IX) 平展态射的下降。基本群的应用
% (Descente des morphismes étales. Application au groupe fondamental)
% A. Grothendieck
% =====================================================================
\chapter{平展态射的下降。基本群的应用}

\section{平展态射的回顾}

我们在这里回顾平展态射的 (在 I 中发展的) 将要用到的性质，借机从理论中消除多余的诺特假设. 读者会注意，即使只关心诺特概形，下降技巧也导致引入非诺特概形 (例如 $\mathrm{Spec}(\widehat{A} \otimes_A \widehat{A})$，其中 $A$ 是诺特局部环)，而为能应用纤维范畴的语言，必须定义平展态射等概念时不引入诺特限制. 不愿验证或承认下列断言在无诺特假设下成立的读者，可以满足于在 I 的诺特假设下承认它们，只要在随后各节的断言中引入同样的诺特假设，并对推理中出现的非诺特概形使用下面的\cref{def:IX.1.1}.

\begin{definition}\label{def:IX.1.1}
设 $f \colon X \to S$ 是预概形态射，$x$ 是 $X$ 的点. 称 $f$ 在 $x$ 处平展，或 $X$ 在 $x$ 处平展于 $S$ 上，若存在 $s = f(x)$ 的仿射开邻域 $U$，$x$ 的位于 $U$ 上方的仿射开邻域 $V$，仿射诺特概形 $U_0$，平展 (I) 且仿射的 $U_0$-概形 $V_0$，态射 $U \to U_0$，和 $U$-同构
\[
V \xrightarrow{\sim} V_0 \times_{U_0} U.
\]
\end{definition}

注意当 $S$ 局部诺特时，这个术语与 loc. cit. 一致. 类似地称 $f$ 平展，或 $X$ 平展于 $S$ 上，若 $f$ 在 $X$ 的每个点 $x$ 处平展. 在这些定义下，下列命题毫无困难地归结到诺特情形，在 I 第 4，5，7 号中证明了. 细节读者可参见 EGA IV\footnote{精确地说，EGA IV 17，18.}.

\begin{remark}\label{rem:IX.1.2}
若 $f$ 在 $x$ 处平展，则 $f$ "在 $x$ 处有限表现" (\cref{cor:VIII.3.5})，$x$ 在纤维 $f^{-1}(s)$ 中的局部环是 $k(s)$ 的有限可分扩张，最后 $f$ 在 $x$ 处平坦. 可以证明反过来也成立，故\cref{def:IX.1.1} 与 $S$ 局部诺特的情形相同，只需把"在 $x$ 处有限型"换成"在 $x$ 处有限表现". 由于这个结果证明微妙，我们这里不想给出平展态射概念的这一定义，它不便直接证明下面要跟随的性质.
\end{remark}

先注意平凡地有：

\begin{proposition}\label{prop:IX.1.3}
若 $f \colon X \to S$ 平展，则由基变换 $S' \to S$ 得到的任何态射 $f' \colon X' \to S'$ 也平展.
\end{proposition}

故可以说平展态射构成 $\mathbf{Sch}$ 中箭的范畴的纤维子范畴 (参见 VI 11 a)). 本讲的对象是研究这个 $\mathbf{Sch}$ 上纤维范畴的正合性质.

\begin{proposition}\label{prop:IX.1.4}
设 $f \colon X \to S$ 是预概形态射. 为使它是开浸入，必要且充分条件是它平展且根式.
\end{proposition}

参见 I 5.1. 由此得出，若 $X$ 在 $S$ 上平展，$X$ 在 $S$ 上的每个截面是开浸入，再用 \ref{prop:IX.1.4}，得到：

\begin{corollary}\label{cor:IX.1.5}
设 $X$ 是平展 $S$-预概形. 则 $X$ 在 $S$ 上的截面集与 $X$ 的使结构态射诱导的态射 $\Gamma \to S$ 根式且满射的开部分 $\Gamma$ 集一一对应.
\end{corollary}

而且若 $X$ 在 $S$ 上分离，$\Gamma$ 将是 $X$ 的同时开且闭的部分，但无关紧要. — 作显然的基变换，可以把 \ref{cor:IX.1.5} 写成表面上更一般的形式：

\begin{corollary}\label{cor:IX.1.6}
设 $X$ 和 $Y$ 是两个 $S$-预概形，$Y$ 在 $S$ 上平展. 则映射 $f \mapsto \Gamma_f$ (把 $X$ 到 $Y$ 的每个 $S$-态射 $f$ 配给 $X \times_S Y$ 的 $f$ 的图的部分) 是从 $\mathrm{Hom}_S(X, Y)$ 到 $X \times_S Y$ 的使 $\mathrm{pr}_1$ 诱导的态射 $\Gamma \to X$ 根式且满射的开部分 $\Gamma$ 集的双射.
\end{corollary}

\begin{proposition}\label{prop:IX.1.7}
设 $S_0$ 是由拟凝聚幂零理想 (i.e. 使 $S_0$ 与 $S$ 有同样底集的理想) 定义的 $S$ 的子预概形. 则函子 $X \mapsto X \times_S S_0$ (从 $S$ 上平展预概形范畴到 $S_0$ 上平展预概形范畴) 是范畴等价.
\end{proposition}

该函子忠实全的事实是 \ref{cor:IX.1.6} 的直接推论. 本质满射的事实包含在 I 8.3 中. 注意在前面的等价中，$X$ 有限型 i.e. 在 $S$ 上拟有限，(resp. 有限 i.e. 是 $S$ 的平展覆盖)，当且仅当 $X_0$ 在 $S_0$ 上满足类似条件；对分离条件同样评注. 这些事实直接，也包含在下面的 \ref{prop:IX.2.4} 中.

\begin{corollary}\label{cor:IX.1.8}
设 $A$ 是剩余域 $k$ 的完备诺特局部环. 则函子 $B \mapsto B \otimes_A k$ 是 $A$ 上有限平展代数范畴与 $k$ 上有限平展代数范畴 (i.e. 由有限个 $k$ 的有限可分扩张组成的代数) 的等价.
\end{corollary}

\begin{proposition}\label{prop:IX.1.9}
为使 $X$ 是 $S$ 的平展覆盖 i.e. 在 $S$ 上有限且平展，必要且充分条件是 $X$ 与 $S$ 上某个代数 $\mathcal{A}$ (它是有限型局部自由模) 的谱 $S$-同构，且对每个 $s \in S$，$\mathcal{A}_s \otimes_{\mathcal{O}_s} k(s)$ 是 $k(s)$ 上的可分代数，从而在此情形是 $k(s)$ 的有限可分扩张的直积.
\end{proposition}

最后下列结果较不初等，它是 I 8.4 与代数几何中层存在定理 (EGA III 5；参见 \cite{exp09-1} 定理 3) 的合取.

\begin{theorem}\label{thm:IX.1.10}
设 $S$ 是完备诺特局部环的谱，$X$ 是 $S$-概形紧合，$X_0$ 是 $X$ 在 $S$ 的闭点处的纤维 (故 $X_0$ 是 $X$ 的闭子概形). 则限制函子 $X' \mapsto X' \times_X X_0$ 是 $X$ 的平展覆盖范畴与 $X_0$ 的平展覆盖范畴的等价.
\end{theorem}

\section{浸没态射与万有浸没态射}

\begin{definition}\label{def:IX.2.1}
预概形态射 $g \colon S' \to S$ 称为浸没\index[term]{jinmeitaishe@浸没态射}，若它满射，且使 $S$ 成为 $S'$ 的商拓扑空间 (i.e. $S$ 的部分 $U$ 使 $f^{-1}(U)$ 开，则 $U$ 开). 称 $f$ 万有浸没\index[term]{wanyoujinmeitaishe@万有浸没态射}，若对任何态射 $T \to S$，由 $f$ 经基变换得到的态射 $f' \colon T' = S' \times_S T \to T$ 浸没.
\end{definition}

两个浸没 (resp. 万有浸没) 态射的复合浸没 (resp. 万有浸没)，万有浸没态射的基变换给出万有浸没态射 (已为此做了所需的调整)，这是直接的. 若 $fg$ 浸没 (resp. 万有浸没)，则 $f$ 浸没.

\begin{example}\label{exa:IX.2.2}
a) 满射且开的或闭的态射浸没，故万有闭满射或万有开满射态射万有浸没. 例如紧合满射态射万有浸没. 另一方面，忠实平坦拟紧态射万有浸没 (\cref{cor:VIII.4.3}). 这将是对我们最重要的两个情形.
\end{example}

可以对浸没或万有浸没态射 $g \colon S' \to S$ 应用 VIII 4.3 的推理，特别得到：

\begin{proposition}\label{prop:IX.2.3}
假设 $g \colon S' \to S$ 浸没. 则下列映射图式正合：
\[
\mathrm{Ouv}(S) \longrightarrow \mathrm{Ouv}(S') \rightrightarrows \mathrm{Ouv}(S''),
\]
其中 $S'' = S' \times_S S'$，$\mathrm{Ouv}(X)$ 表示预概形 $X$ 的开部分集.
\end{proposition}

\begin{proposition}\label{prop:IX.2.4}
设 $g \colon S' \to S$ 是万有浸没态射，$f \colon X \to Y$ 是 $S$-态射，$f' \colon X' \to Y'$ 是经基变换得到的 $S'$-态射. 为使 $f$ 开 (resp. 闭)，只需 $f'$ 如此. 为使 $f$ 万有开 (resp. 万有闭，resp. 分离)，必要且充分条件是 $f'$ 如此. 若还假设 $g$ 拟紧，$f$ 局部有限型，为使 $f$ 紧合，必要且充分条件是 $f'$ 如此.
\end{proposition}

对最后一点，注意若 $f'$ 紧合故拟紧，则由 VIII 3.3 $f$ 拟紧，故有限型，因为它局部有限型. 另一方面由前面所述它分离且万有闭，故紧合.

\begin{proposition}\label{prop:IX.2.5}
设 $S'$ 是完备诺特局部环的谱 $S$ 上有限型的预概形，假设 $S$ 的闭点 $s$ 的纤维有限，故该纤维诸点 $s'$ 在 $S'$ 中的局部环在 $A = \mathcal{O}_s$ 上有限. 设 $S''$ 是所论诸 $\mathcal{O}_{S',s'}$ 的谱的和概形，看作有限 $S$-概形. 为使 $g \colon S' \to S$ 万有浸没，必要且充分条件是结构态射 $S'' \to S'$ 满射.
\end{proposition}

由于有自然 $S$-态射 $S'' \to S'$，且有限满射态射由 \ref{exa:IX.2.2} 万有浸没，所述条件充分. 证明若 $S'' \to S'$ 不满射，则 $g$ 不万有浸没. 事实上，设 $t$ 是不在 $S''$ 的像中的 $S$ 的点；则存在 $S$-概形 $T$，为离散赋值环的谱，它在 $S$ 中的像是 $\{s, t\}$. 注意 $S''$ 在 $S'$ 中的像开，因为态射 $S'' \to S'$ 局部同构，另一方面该像包含 $S'_s$，且与 $S'_t$ 不相交. 由此得出这最后一个在 $T' = S' \times_S T$ 中的原像开，且等于 $T$ 的闭点的原像. 这表明 $T' \to T$ 不浸没，故 $S' \to S$ 不万有浸没.

\begin{remark}\label{rem:IX.2.6}
用判则 IV 6.3 (诺特空间中可构部分开) 容易得到下列赋值判则：有限型态射 $g \colon S' \to S$ ($S$ 局部诺特) 万有浸没，必要且充分条件是：对每个为离散赋值环谱的 $S$-概形 $T$，令 $T' = S' \times_S T$，$T$ 的闭点在 $T'$ 中的原像不开.
\end{remark}

\section{平展预概形态射的下降}

\begin{proposition}\label{prop:IX.3.1}
设 $g \colon S' \to S$ 是预概形满射态射，$X$ 和 $Y$ 是 $S$ 上的两个预概形，$X'$，$Y'$ 是它们在 $S'$ 上的逆像. 若 $Y$ 在 $S$ 上非分歧，则典范映射
\[
\mathrm{Hom}_S(X, Y) \longrightarrow \mathrm{Hom}_{S'}(X', Y')
\]
单射.
\end{proposition}

事实上，由 \ref{cor:IX.1.6}，$S$-态射 $f \colon X \to Y$ 在已知它的图的底集 $\Gamma$ (它是 $Z = X \times_S Y$ 的部分) 时确定. 由于
\[
Z' = Z \times_S S' = X' \times_{S'} Y' \longrightarrow Z
\]
满射 (因为 $S' \to S$ 满射)，该部分 $\Gamma$ 在已知它在 $X' \times_{S'} Y'$ 中的原像时确定，后者正是 $f'$ 的图的底集. 由此得结论.

$Z$ 的部分 $\Gamma$ 是 $S$-态射 $f \colon X \to Y$ 的图，当且仅当它开，且 $\mathrm{pr}_1$ 诱导的态射 $\Gamma \to Y$ 根式且满射，参见 1. 当第一个性质满足时，第二个性质成立当且仅当 $\Gamma$ 在 $Z'$ 中的原像 $\Gamma'$ 满足同样条件 (\cref{prop:VIII.3.1}). 若还知道 $Z' \to Z$ 浸没 (特别当 $S' \to S$ 万有浸没时如此)，则 $\Gamma$ 开当且仅当 $\Gamma'$ 开. 这样，集合 $\mathrm{Hom}_S(X, Y)$ 于是与 $Z'$ 的下列开部分 $\Gamma'$ 的集一一对应：投影态射 $\mathrm{pr}_1 \colon Z' \to X'$ 根式且满射 (i.e. 对应 $S'$-态射 $f' \colon X' \to Y'$)，且对由 $Z' \to Z$ 定义的等价关系饱和，i.e. 它在 $Z'' = Z' \times_Z Z' = Z \times_S S''$ (其中 $S'' = S' \times_S S'$) 中经两个投影之一的另一个的原像相等. 而后者正是由 $f'$ 经基变换 ($S'' \to S'$ 的两个投影之一另一个) 得到的两个 $S''$-态射 $X'' \to Y''$ 的图. 于是得到：

\begin{proposition}\label{prop:IX.3.2}
设 $g \colon S' \to S$ 是预概形万有浸没态射，$S'' = S' \times_S S'$，$X$ 和 $Y$ 是两个 $S$-预概形，$X'$ 和 $Y'$ 是它们在 $S'$ 上的逆像，$X''$，$Y''$ 是它们在 $S''$ 上的逆像. 若 $Y$ 在 $S$ 上平展，则下列典范映射图式正合：
\[
\mathrm{Hom}_S(X, Y) \longrightarrow \mathrm{Hom}_{S'}(X', Y') \rightrightarrows \mathrm{Hom}_{S''}(X'', Y'').
\]
\end{proposition}

取 $X$ 和 $Y$ 在 $S$ 上平展，得到下列断言，它反过来给出 \ref{prop:IX.3.2} (甚至限于 $X = S$，此时总能由基变换 $X \to S$ 归结到 \ref{prop:IX.3.2})：

\begin{corollary}\label{cor:IX.3.3}
预概形的万有浸没态射对互相平展的预概形的纤维范畴是下降态射.
\end{corollary}

我也不知道它是否必然对所述纤维范畴是有效下降态射，即使再假设 $S$ 诺特且 $g$ 有限型，并限于平展覆盖. 我们仍将在下一节给出有用的有效性判据.

\begin{corollary}\label{cor:IX.3.4}
设 $g \colon S' \to S$ 是万有浸没态射，其诸纤维 $g^{-1}(s)$ "几何连通\index[term]{jiheliantong@几何连通}"，i.e. 对任何扩张 $K/k(s)$，$g^{-1}(s) \otimes_{k(s)} K$ 连通. 则 $S$ 连通时 $S'$ 连通. 由 $g$ 定义的从 $S$ 上平展预概形范畴到 $S'$ 上平展预概形范畴的函子忠实全.
\end{corollary}

$S'$ 的同时开且闭的部分对由 $g$ 定义的集合等价关系饱和，因为诸纤维连通，故它是 $S$ 的某个部分的原像，后者由于 $g$ 浸没必然开且闭. 故 $S$ 连通时 $S'$ 连通. 这也蕴含下列结果：两个万有连通纤维态射的复合 $fg$ ($f$ 万有浸没) 有万有连通纤维；若 $S$ 上的 $S'_1$ 和 $S'_2$ 有万有连通纤维，则 $S'_1 \times_S S'_2$ 亦然. 特别在 \ref{cor:IX.3.4} 的条件下，$S''$ 在 $S$ 上有万有连通纤维. 设 $X$ 和 $Y$ 在 $S$ 上平展，设 $u'$ 是 $X'$ 到 $Y'$ 的 $S'$-态射，证明它与下降数据相容 (这由 \ref{cor:IX.3.3} 给出所求结论). 设 $u''_1$ 和 $u''_2$ 是由 $u'$ 得到的两个 $S''$-态射 $X'' \to Y''$. $S''$ 的 $u''_1$ 和 $u''_2$ 的一致点子预概形是诱导开子预概形，纤维性地闭，它是 $Y''$ 在 $S''$ 上的对角预概形的原像\footnote{注意 $S'$ 在 $S$ 上的诸纤维分离!}. 故它是 $S$ 的某部分的原像. 由于它包含 $S''$ 中的对角线，它等于 $S''$，故 $u''_1 = u''_2$，证毕.

\section{平展预概形的下降：有效性判据}

\begin{proposition}\label{prop:IX.4.1}
设 $g \colon S' \to S$ 是忠实平坦拟紧态射. 则 $g$ 对互相平展、分离且有限型的预概形的纤维范畴是有效下降态射.
\end{proposition}

这确实是所述纤维范畴的下降态射，由 \ref{cor:IX.3.3} 或 VIII 5.2 任选其一. 剩下证明：若 $X'$ 在 $S'$ 上平展、分离且有限型，配备相对 $g \colon S' \to S$ 的下降数据，则该数据在所述纤维范畴中有效. 容易看到，若 $X$ 是 $S$ 上的预概形，则它在 $S$ 上平展当且仅当它在 $S'$ 上平展 (由\cref{def:IX.1.1} 和 loc. cit. 3.6). 故它在 $S$ 上平展、分离且有限型，当且仅当 $X'$ 在 $S'$ 上如此，参见例如 2.4. 故只需保证 $X$ 上下降数据对 $\mathbf{Sch}$ 的箭的纤维范畴有效. 而由 VIII 6.2 和 \ref{cor:IX.6.6}，$X'$ 在 $S'$ 上拟仿射. 于是可用 VIII 7.9 得出结论. 读者还会注意，若限于互相平展且有限的预概形，证明要求更少，因为可直接援引 VIII 2.1.

\begin{corollary}\label{cor:IX.4.2}
设 $g \colon S' \to S$ 是万有浸没态射，$X'$ 是 $S'$ 上配备相对 $g$ 的下降数据的平展分离有限型 $S'$-预概形，$S_1 \to S$ 是忠实平坦拟紧态射，$S'_1$ 和 $X'_1$ 由 $S'$ 和 $X'$ 经基变换得到，故 $S'_1 \to S_1$ 万有浸没，$X'_1$ 在 $S'_1$ 上平展分离且有限型，配备相对 $g_1 \colon S'_1 \to S_1$ 的下降数据. 为使 $X'$ 上的下降数据有效，必要且充分条件是 $X'_1$ 上的下降数据有效.
\end{corollary}

这由范畴中的下降理论 \cite{exp09-D} 得出，考虑到 \ref{prop:IX.4.1} 和 3.3.

类似地证明：

\begin{corollary}\label{cor:IX.4.3}
设 $g \colon S' \to S$ 是万有浸没态射，$X'$ 是 $S'$ 上配备相对 $g$ 的下降数据的平展 $S'$-预概形，$(S_i)$ 是 $S$ 的开覆盖. 为使下降数据有效，必要且充分条件是：对每个 $i$，相对态射 $g_i \colon S'_i = S' \times_S S_i \to S_i$ 在 $X'_i = X' \times_S S_i$ 上对应的下降数据有效.
\end{corollary}

最后一个结果导致给出局部有效性判据：

\begin{proposition}\label{prop:IX.4.4}
设 $g \colon S' \to S$ 是有限表现 (\cref{rem:VIII.3.6}) 且万有浸没的态射，$X'$ 是 $S'$ 上平展有限表现的预概形，配备相对 $g$ 的下降数据，最后 $a$ 是 $S$ 的点. 为使存在 $a$ 的开邻域 $U$，使相对态射
\[
g_U \colon S'_U = S' \times_S U \longrightarrow S_U = U
\]
在 $X'_U = X' \times_S U$ 上对应的下降数据有效，必要且充分条件是相对态射
\[
g_a \colon S'_a = S' \times_S \mathrm{Spec}(\mathcal{O}_a) \longrightarrow S_a = \mathrm{Spec}(\mathcal{O}_a)
\]
在 $X'_a = X' \times_S \mathrm{Spec}(\mathcal{O}_a)$ 上对应的下降数据有效.
\end{proposition}

必要性平凡，证明充分性. 于是有 $S_a$ 上有限型平展预概形 $X_a$ 和与下降数据相容的同构
\[
(\ast) \quad X'_a \xrightarrow{\sim} X_a \times_{S_a} S'_a.
\]
按关于环的归纳极限 (这里是环 $A_f$，其中 $A$ 是 $a$ 的仿射开邻域的环，$f$ 遍历不在对应 $a$ 的素理想中的元素) 上定义的预概形的容易的一般推论，可以找到 $a$ 的开邻域 $U$，$U = S_U$ 上有限型平展预概形 $X_U$，和 $S_a$-同构 $X_a \xrightarrow{\sim} X_U \times_{S_U} S_a$. 而且，取 $U$ 足够小，可假设同构 $(\ast)$ 来自同构：
\[
X'_U \xrightarrow{\sim} X_U \times_{S_U} S'_U;
\]
后者可能不与下降数据相容，但只要缩小 $U$，它将与下降数据相容. 这完成证明.

\begin{corollary}\label{cor:IX.4.5}
在 \ref{prop:IX.4.4} 的条件下，为使 $X'$ 上的下降数据有效，必要且充分条件是：对每个 $a \in S$，相对态射 $S'_a = S' \times_S \mathrm{Spec}(\mathcal{O}_a) \to \mathrm{Spec}(\mathcal{O}_a)$ 在 $X'_a$ 上对应的下降数据有效. 当 $S$ 局部诺特且 $X'$ 在 $S'$ 上分离时，在前面判则中还可把 $\mathcal{O}_a$ 换成它的完备化.
\end{corollary}

第一个断言由 \ref{prop:IX.4.4} 和 \ref{cor:IX.4.3} 得出，第二个断言则是 \ref{cor:IX.4.2} 的推论. 再用 \ref{cor:IX.4.2} 及事实：对每个诺特局部环 $A$，可以找到完备诺特局部环 $B$ 和局部同态 $A \to B$，使 $B$ 在 $A$ 上平坦且 $B/\mathfrak{m}B$ 是 $A$ 的剩余域 $k = A/\mathfrak{m}$ 的给定扩张，得到：

\begin{corollary}\label{cor:IX.4.6}
在 \ref{prop:IX.4.4} 的条件下，再假设 $X'$ 在 $S'$ 上分离，$S$ 局部诺特. 为使 $X'$ 上的下降数据有效，必要且充分条件是：对每个为剩余域代数闭的完备局部环谱的 $S$-概形 $S_1$，相对态射 $g_1 \colon S'_1 \to S_1$ 在 $X'_1 = X' \times_S S_1$ 上对应的下降数据有效.
\end{corollary}

\begin{theorem}\label{thm:IX.4.7}
设 $g \colon S' \to S$ 是有限满射有限表现态射 (若 $S$ 局部诺特，这最后一个假设是其他假设的推论)\footnote{可以证明实际上只需 $g$ 是整态射，用 EGA IV 8 风格的极限过程归结到正文的情形.}. 则 $g$ 对互相平展、分离、有限型的预概形的纤维范畴是有效下降态射.
\end{theorem}

需证明：若 $X'$ 在 $S'$ 上平展、分离、有限型，配备相对 $g$ 的下降数据，则该数据有效. 用 \ref{cor:IX.4.3} 容易归结到 $S$ 诺特的情形. 由 \ref{cor:IX.4.5}，可假设 $S$ 是诺特局部环的谱，从而尤其
\[
\dim S = n < +\infty.
\]
对 $\dim S = n$ 作递归，$n < 0$ 时断言平凡. 假设 $n \geqslant 0$ 且定理对维数 $n' < n$ 已证. 由 \ref{cor:IX.4.6} 归结到 $S$ 是完备局部环的谱的情形，故 $S'$ 是完备局部环的谱的有限并. 故有
\[
X' = X'_1 \sqcup X'_2
\]
其中 $X'_1$ 在 $S'$ 上有限，$X'_2$ 在 $S'$ 的任一个闭点上方都没有点. 考虑对应下降数据的态射
\[
q_1, q_2 \colon X'' \rightrightarrows X',
\]
它与 $p_1, p_2 \colon S'' \rightrightarrows S'$ 相容. 立即看到
\[
X'' = q_i^{-1}(X'_1) \sqcup q_i^{-1}(X'_2) \quad i = 1, 2
\]
是 $X''$ 在 $S''$ 上的对应典范分解，这蕴含 $q_1^{-1}(X'_1) = q_2^{-1}(X'_1)$，从而 $X'_1$ 和 $X'_2$ 配备诱导下降数据. 设 $T$ 是 $S$ 的闭点的余集的开集，故 $T' = S' \times_S T$ 是 $S'$ 的闭点集余集的部分，而整个位于 $T'$ 上方的 $X'_2$ 配备相对 $g$ 诱导的态射 $T' \to T$ 的下降数据. 由于后者有限满射，且 $\dim T < \dim S = n$，由递归假设该下降数据有效. 故只需证明 $X'_1$ 上的下降数据有效，现在可假设 $X'$ 在 $S'$ 上平展且有限. (N.B. 若在断言 \ref{thm:IX.4.7} 中限于平展覆盖，递归推理不必要). 设 $S_0$ 是 $A$ 的剩余域的谱，设 $S'_0 = S' \times_S S_0$，同样由 $S'$ 在 $S$ 上的纤维平方和立方 $S''$ 和 $S'''$ 定义 $S''_0$，$S'''_0$. 由 \ref{cor:IX.1.8}，态射 $S_0 \to S$，$S'_0 \to S'$ 等对平展覆盖范畴 (分别 $S$ 和 $S_0$ 的，$S'$ 和 $S'_0$ 的，等等) 诱导等价. 由范畴中下降理论的推论 \cite{exp09-D}，得出 $g \colon S' \to S$ 对平展覆盖纤维范畴是有效下降态射，当且仅当 $g_0 \colon S'_0 \to S_0$ 如此. 而确实如此，例如作为 \ref{prop:IX.4.1} 的特殊情形. 这完成 \ref{thm:IX.4.7} 的证明.

\begin{corollary}\label{cor:IX.4.8}
假设 $S$ 局部诺特，\ref{thm:IX.4.7} 的结论在只假设 $S' \to S$ 万有浸没、有限型且拟有限时仍成立.
\end{corollary}

由 \ref{cor:IX.4.6}，可假设 $S$ 是完备诺特局部环的谱. 则由 \ref{prop:IX.2.5}，存在有限满射态射 $S_1 \to S$ 和 $S$-态射 $S_1 \to S'$. 由于 $S_1 \to S$ 由 \ref{thm:IX.4.7} 对所述纤维范畴是万有严格下降态射，且 $S' \to S$ 对它万有下降，\ref{cor:IX.4.8} 由一般推论 \cite{exp09-D} 得出.

\begin{corollary}\label{cor:IX.4.9}
设 $g \colon S' \to S$ 是有限型、满射且万有开的态射，$S$ 局部诺特. 则 $g$ 对互相平展、分离且有限型的预概形的纤维范畴是有效下降态射.
\end{corollary}

如 \ref{thm:IX.4.7} 那样进行，归结到 $S$ 是诺特完备局部环 $A$ 的谱的情形. 设 $A_1$ 是 $A$ 上的有限代数，谱为 $S_1$，使 $S_1 \to S$ 有限满射，故由 \ref{thm:IX.4.7} 对所述纤维范畴是万有有效下降态射. 由一般定理 \cite{exp09-D} 得出，$g$ 对所述纤维范畴是有效下降态射，当且仅当对应态射 $g_1 \colon S'_1 = S' \times_S S_1 \to S_1$ 如此. 由于后者满足与 $g$ 同样的假设，归结到对 $S_1$ 代替 $S$ 证明 4.9. 先取 $A_1$ 为诸 $A/\mathfrak{p}_i$ 的直积 (其中 $\mathfrak{p}_i$ 遍历 $A$ 的极小素理想)，归结到 $A$ 整的情形. 此时证明\footnote{参见 EGA IV 14.3.13 和 14.5.4.}存在 $S'$ 的整子概形 $S_1$，在 $S$ 上拟有限且支配 $S$，经过 $S'$ 在 $S$ 的闭点 $y$ 处的纤维的某个点 (因为 $S'$ 在 $S$ (整局部诺特) 上万有开有限型，且 $S'_y \neq \varnothing$). 由于 $A$ 完备，$S_1$ 在 $S$ 上有限，且因它支配 $S$，态射 $S_1 \to S$ 满射. 再次把 $S$ 换成 $S_1$，归结到 $S'$ 在 $S$ 上有截面的情形，此时断言平凡.

\begin{theorem}\label{thm:IX.4.10}
设 $g \colon S' \to S$ 是有限根式满射有限表现态射 (若 $S$ 局部诺特，这最后一个条件多余\footnote{甚至只需 $g$ 整、根式且满射，如容易归结到正文情形所见，风格如 EGA IV 8，参见 SGA 4 VIII 1.1.}). 则逆像函子诱导 $S$ 上平展预概形范畴与 $S'$ 上平展预概形范畴的等价.
\end{theorem}

由于 $S'$ 在 $S' \times_S S'$ 和 $S' \times_S S' \times_S S'$ 中的对角态射是满射浸入，故由 \ref{prop:IX.1.9} 诱导 (分别 $S' \times_S S'$ 和 $S' \times_S S' \times_S S'$ 上平展预概形范畴与 $S'$ 上平展预概形范畴的) 等价，由下降推论 \cite{exp09-D} 得出每个在 $S'$ 上平展的 $X'$ 配备相对 $g \colon S' \to S$ 的下降数据，且只有一个. 故 \ref{cor:IX.3.3} 蕴含经 $g$ 的逆像函子 (从 $S$ 上平展预概形范畴到 $S'$ 上平展预概形范畴) 忠实全. 剩下证明它本质满射，i.e. 每个在 $S'$ 上平展的 $X'$ 同构于某个在 $S$ 上平展的 $X$ 的逆像. 问题显然在 $S$ 和 $X'$ 上局部，可假设 $S$，$S'$，$X'$ 仿射. 但此时 $X'$ 在 $S'$ 上分离且有限型，可应用有效性判据 4.7.

\begin{corollary}\label{cor:IX.4.11}
把对 $g$ 的假设换成：$g$ 忠实平坦、拟紧且根式，则 \ref{cor:IX.4.9} 的结论仍然成立.
\end{corollary}

同样的证明，用 \ref{prop:IX.4.1} 代替 4.7.

注意 \ref{thm:IX.4.7} 的证明"初等"，因为它不用紧合态射的有限性与比较定理 (EGA III 3，4，5). 对下列结果不再如此：

\begin{theorem}\label{thm:IX.4.12}
设 $g \colon S' \to S$ 是紧合、满射、有限表现 (若 $S$ 局部诺特，这最后一个假设是第一个的推论) 的态射. 则 $g$ 对预概形的平展覆盖纤维范畴是有效下降态射.
\end{theorem}

由 \ref{cor:IX.3.3} 和 \ref{exa:IX.2.2}，归结到证明：对 $S'$ 上的每个平展覆盖 $X'$，配备相对 $g \colon S' \to S$ 的下降数据，该下降数据有效. 用 \ref{cor:IX.4.3} 容易归结到 $S$ 诺特的情形，用 \ref{cor:IX.4.6} 可假设 $S$ 是完备诺特局部环 $A$ 的谱. 如常引入 $S''$ 和 $S'''$，设 $S_0$ 是 $A$ 的剩余域的谱，设 $S'_0$，$S''_0$，$S'''_0$ 由 $S'$，$S''$，$S'''$ 经基变换 $S_0 \to S$ 得到，i.e. 是 $S'$，$S''$，$S'''$ 在 $S$ 的闭点处的纤维. 由 \ref{thm:IX.1.10}，态射 $S_0 \to S$，$S'_0 \to S'$ 等对目标概形与源概形的平展覆盖范畴诱导等价. 从而 $g \colon S' \to S$ 对预概形的平展覆盖纤维范畴是严格下降态射，当且仅当 $g_0 \colon S'_0 \to S_0$ 如此，而由 \ref{prop:IX.4.1} 确实如此. 这完成 \ref{thm:IX.4.12} 的证明. (在这个推理中，只需要 \ref{thm:IX.1.10} 中所述函子忠实全的事实，它不用代数几何中凝聚层存在定理.)

\section{基本群术语的翻译}

设
\[
g \colon S' \longrightarrow S
\]
是对预概形的平展覆盖纤维范畴的有效下降态射，例如紧合、满射、有限表现态射 (\cref{thm:IX.4.12})，或忠实平坦拟紧态射. 如常引入 $S''$，$S'''$，并用 $\mathcal{C}$，$\mathcal{C}'$，$\mathcal{C}''$，$\mathcal{C}'''$ 表示 $S$，$S'$，$S''$，$S'''$ 的平展覆盖范畴，有范畴的 2-正合图式
\[
(\ast) \quad \mathcal{C} \xrightarrow{\;p^{*}\;} \mathcal{C}' \xrightarrow{\;p_1^{*}, p_2^{*}\;} \mathcal{C}'' \xrightarrow{\;p_{21}^{*}, p_{32}^{*}, p_{31}^{*}\;} \mathcal{C}'''
\]
(其中第二个和第三个箭分别表示两个箭 $p_1^{*}, p_2^{*}$ 和三个箭 $p_{21}^{*}, p_{32}^{*}, p_{31}^{*}$) 对应图式
\[
S \longleftarrow S' \xleftarrow{p_1, p_2} S'' \xleftarrow{p_{21}, p_{32}, p_{31}} S'''.
\]
假设预概形 $S$，$S'$，$S''$，$S'''$ 是连通预概形的无交和，特别当它们是局部连通预概形时如此，更不用说局部诺特 (例如 $S'$ 在局部诺特 $S$ 上有限型). 则 $(\ast)$ 中的范畴 $\mathcal{C}$，$\mathcal{C}'$ ... 是 V 9 的多 Galois 范畴，各自由一族拓扑紧全不连通的群描述，即预概形 $S$，$S'$，$S''$，$S'''$ 的连通分量的基本群. 为简单假设 $S$ 连通，将给出计算它的基本群的方法，用由 $\mathcal{C}'$，$\mathcal{C}''$，$\mathcal{C}'''$ 组成、借助表达这些范畴的基本群适当明确的纤维范畴. 读者会注意，勾勒的方法实际上在多 Galois 范畴的一般框架 (它们不必来自给定的预概形 $S$，$S'$，$S''$，$S'''$) 中有效. 它也是计算拓扑空间 $S$ (闭子空间 $S_i$ 的局部有限并，或开子空间 $S_i$ 的任意并) 基本群的熟知方法的类似物，用 $S_i$ 分量和 $S_i \cap S_j$ 分量的基本群. 当然，预概形框架中的类似情形确实落入下降的一般框架，只需引入 $S_i$ 的和预概形 $S'$ 和典范态射 $g \colon S' \to S$.

令
\[
E' = \pi_0(S'), \quad E'' = \pi_0(S''), \quad E''' = \pi_0(S'''),
\]
其中 $\pi_0$ 表示"连通分量集"函子. 由于 $S'$ 在 $S$ 上的纤维积构成 $\mathbf{Sch}$ 的单纯对象，它被函子 $\pi_0$ 变换成 $E'$，$E''$，$E'''$ 是维数 0，1，2 的分量的单纯集. 将用到单纯映射
\[
q_i = \pi_0(p_i), \quad (i = 1, 2), \quad \text{和} \quad q_{ij} = \pi_0(p_{ij}), \quad (i, j) = (2, 1), (3, 2), (3, 1),
\]
在图式中显示：
\[
(1) \quad E' \xleftarrow{\;q_1, q_2\;} E'' \xleftarrow{\;q_{21}, q_{32}, q_{31}\;} E'''.
\]
$E'$ 的对象用带一个撇的记号，如 $s'$，$E''$ resp. $E'''$ 的对象用带两个 resp. 三个撇的记号. $S$ 连通的事实翻译为 $\pi_0(K) = 0$，其中 $K$ 是由 $g \colon S' \to S$ 定义的单纯集，或也翻译为：映射对 $(q_1, q_2)$ 在 $E'$ 中生成的等价关系传递.

一次选定 $E'$ 的元素 $s'_0$，并对每个 $s' \in E'$ 选 $E''$ 的元素 $s'$ 使\footnote{注意隐式承认存在的元素 $s'$ 并非在所有情形存在. 故如所述陈述的\cref{thm:IX.5.1} 不适用于所有情形. 但参照所写的文本不难修改该定理的陈述，使它给出在所有情形适用的计算方法. 特别，该定理的诸推论原样成立.}
\[
q_1(s') = s'_0, \quad q_2(s') = s',
\]
从而显示 $S$ 的连通性. 对每个 $s' \in E'$，在 $S'$ 的连通分量 $s'$ 中选几何点 $s'$；该点实际上通过 $C'$ 上对应的纤维函子 $F'_{s'}$ 出现. 该函子的自同构群，i.e. $S'$ 在 $s'$ 处的基本群，记为 $\pi_{s'}$. 类似地选 $s''$ 和 $s'''$，从而选函子 $F''_{s''}$ 和 $F'''_{s'''}$，得基本群 $\pi_{s''}$ 和 $\pi_{s'''}$. 这样
\[
\pi_{s'} = \pi_1(S', s'), \quad \pi_{s''} = \pi_1(S'', s''), \quad \pi_{s'''} = \pi_1(S''', s''').
\]
对每个 $s'' \in E''$，$p_1(s'')$ 与 $q_1(s'')$ 在同一个连通分量中，故存在函子同构 $F''_{s''} \circ p_1^{*} \xrightarrow{\sim} F'_{q_1(s'')}$ (i.e. 从 $p_1(s'')$ 到 $q_1(s'')$ 的"道路类"). 对 $q_2$ 和诸 $q_{ij}$ 重复此评注. 选定所有这些道路类：
\[
F''_{s''} \circ p_i^{*} \xrightarrow{\sim} F'_{q_i(s'')}, \quad F'''_{s'''} \circ p_{ij}^{*} \xrightarrow{\sim} F''_{q_{ij}(s''')},
\]
(对 $i = 1, 2$ 和 $(i, j) = (2, 1), (3, 2), (3, 1)$). 特别由此得到群同态：
\[
(2) \quad q_i^{s''} \colon \pi_{s''} \longrightarrow \pi_{q_i(s'')}, \quad q_{ij}^{s'''} \colon \pi_{s'''} \longrightarrow \pi_{q_{ij}(s''')},
\]
($i$ 和 $(i, j)$ 取同样值). 最后回忆在纤维 $\mathcal{C}'$，$\mathcal{C}''$，$\mathcal{C}'''$ 的劈开范畴结构中还有函子同构：
\[
p_{21}^{*} p_1^{*} \xrightarrow{\sim} p_{31}^{*} p_1^{*}, \quad p_{21}^{*} p_2^{*} \xrightarrow{\sim} p_{32}^{*} p_1^{*}, \quad p_{31}^{*} p_2^{*} \xrightarrow{\sim} p_{32}^{*} p_2^{*},
\]
它们由两个成员分别与 $u_i^{*}$ ($i = 1, 2, 3$，其中 $u_i$ 是 $S'''$ 到 $S'$ 的三个投影) 的同构得到. 明确这些数据时，对每个 $s'''$ 找到完全确定的元素
\[
(3) \quad a_i^{s'''} \in \pi_{v_i(s''')},
\]
(其中 $v_i$，$i = 1, 2, 3$ 是由 $v_i = \pi_0(u_i)$ 定义的三个映射 $E''' \to E'$)，它们服从条件：
\[
q_1^{s''_1} q_{21}^{s'''} = \mathrm{int}(a_1^{s'''}) q_1^{s''_2} q_{31}^{s'''} \quad (s''_1 = q_{21}(s'''), s''_2 = q_{31}(s''')),
\]
以及两个类似条件，涉及 $a_2$ 和 $a_3$. 读者还会注意，数据 (1)，(2)，(3) 允许在纤维范畴等价的范围内重构所考虑的纤维 $\mathcal{C}'$，$\mathcal{C}''$，$\mathcal{C}'''$ 的纤维范畴. 故原则上应允许在等价范围内重构 $\mathcal{C}$，从而在同构范围内重构它的基本群. 我们实际上将在几何点 $p(s'_0)$ 处确定基本群，i.e. 函子 $F'_{s'_0} \circ p^{*}$ 的自同构群.

注意 $\mathcal{C}'$ 的对象 $X'$ 的给出本质上等价于有限集 $X'_{s'}$ ($s' \in E'$) 的给出，其中 $\pi_{s'}$ 连续作用. 这样的对象上的粘合数据于是等价于对每个 $s'' \in E''$ 给一个双射：
\[
\varphi_{s''} \colon X'_{q_1(s'')} \xrightarrow{\sim} X'_{q_2(s'')}
\]
与 $\pi_{s''}$ 的作用相容，后者借助同态 $q_i^{s''} \colon \pi_{s''} \to \pi_{q_i(s'')}$ 作用于两个成员之一另一个. 先取 $s''$ 为 $s'$ 形式的元素，看到这样的数据定义双射
\[
\psi_{s'} \colon X'_{s'_0} = F'_0(X') \xrightarrow{\sim} X'_{s'}
\]
这使得可以把诸 $X'_{s'}$ 等同于同一个集合 $F'_0(X') = X'_{s'_0}$，所有群 $\pi_{s'}$ 此后在其上作用. 这样，双射 $\varphi_{s''}$ 对应双射
\[
g_{s''} \colon F'_0(X') \xrightarrow{\sim} F'_0(X'),
\]
一方面服从与 $\pi_{s''}$ 的交换关系：
\[
\mathrm{a)} \quad g_{s''} q_1^{s''}(g'') = q_2^{s''}(g'') g_{s''} \quad (s'' \in E'', g'' \in \pi_{s''}),
\]
另一方面服从关系
\[
\mathrm{b)} \quad g_{s'} = g_{s'_0} \quad (s' \in E'),
\]
表达我们把诸 $X_{s'}$ 彼此等同的方式. 明确这样的粘合数据成为下降数据的条件时，找到关系：
\[
\mathrm{c)} \quad a_3^{s'''} g_{q_{31}(s''')} a_1^{s'''} = g_{q_{32}(s''')} a_2^{s'''} g_{q_{21}(s''')} \quad (s''' \in E''').
\]
这给出 $\mathcal{C}'$ 的配备下降数据的对象范畴与下列范畴的等价：有限集 ($\pi_{s'}$ 在其上连续作用)，再配备满足关系 a)，b)，c) 的双射 $g_{s''}$.

设 $G$ 是由群 $\pi_{s'}$ 和新生成元 $g_{s''}$ 服从关系 a)，b)，c) 生成的群，设 $\pi$ 是 $G$ 的使在诸群 $\pi_{s'}$ 中的原像为开子群的有限指数子群的商的投射极限. 也说 $\pi$ 是由诸 $\pi_{s'}$ 和诸 $g_{s''}$ 服从关系 a)，b)，c) 生成的 Galois 型群. 立即看到所考虑的范畴也等价于拓扑群 $\pi$ 在其上连续作用的有限集范畴. 这建立下列断言：

\begin{theorem}\label{thm:IX.5.1}
设 $g \colon S' \to S$ 是对预概形的平展覆盖纤维范畴的有效下降态射 (参见 \ref{cor:IX.4.9} 和 \ref{thm:IX.4.12}) 的预概形态射. 假设 $S$ 连通，且 $S'$，它的纤维平方 $S''$ 和纤维立方 $S'''$ 是连通预概形的和 (例如 $S'$ 在连通局部诺特 $S$ 上有限型时如此). 如上选择：$S'$，$S''$，$S'''$ 的每个连通分量中的几何点，某些道路类，$E'$ 中的元素 $s'_0$，以及对每个 $s' \in E'$ 满足在 $E'$ 中的两个像为 $s'_0$ 和 $s'$ 的 $E''$ 的元素 $s'$. ($E'$，$E''$，$E'''$ 分别表示 $S'$，$S''$，$S'''$ 的连通分量集). 则 $S$ 在 $s'_0$ 的像处的几何点的基本群典范同构于由诸 $\pi_{s'} = \pi_1(S', s')$ ($s' \in E'$) 和生成元 $g_{s''}$ ($s'' \in E''$) 服从上述关系 a)，b)，c) (涉及群 $\pi_{s''} = \pi_1(S'', s'')$ 的元素和前面引入的元素 $a_i^{s'''}$ ($i = 1, 2, 3$，$s''' \in E'''$)) 生成的 Galois 型群.
\end{theorem}

\begin{corollary}\label{cor:IX.5.2}
假设 $S'$ 和 $S''$ 只有有限个连通分量，且 $S'$ 的连通分量的基本群拓扑有限生成. 则 $S$ 的基本群拓扑有限生成.
\end{corollary}

这样，我们以后将证明代数闭域上射影\emph{正规概形}的基本群拓扑有限生成. 用 Chow 引理和代数概形的正规化，由此得出同样结果对代数闭域上任何紧合概形成立.

\begin{corollary}\label{cor:IX.5.3}
假设 $S'$，$S''$，$S'''$ 只有有限个连通分量，$S'$ 的连通分量的基本群拓扑有限表现，$S''$ 的连通分量的基本群拓扑有限生成. 则 $S$ 的基本群拓扑有限表现.
\end{corollary}

注意可以把 \ref{cor:IX.4.9} (限于平展覆盖) 表达为：诺特预概形的有限根式满射态射诱导基本群的同构；形象地说，基本群是预概形的拓扑不变量. 可以用 \ref{thm:IX.5.1} 更一般地明确对预概形操作 (例如预概形沿有限点集的"捏合"，它有简单拓扑意义) 对基本群的作用. 例如得到：

\begin{corollary}\label{cor:IX.5.4}
设 $g \colon S' \to S$ 是有限有限表现态射，$T$ 是 $S$ 的离散部分. 对每个 $s \in S$，设 $n(s)$ 是纤维 $g^{-1}(s)$ 中的"\emph{几何点数}" (它也可明确为 $g^{-1}(s)$ 在 $k(s)$ 上的可分次数，即它的剩余扩张的可分次数的和). 假设对 $s \in S - T$，有 $n(s) = 1$. 对每个 $s \in T$，设 $K_s$ 是 $k(s)$ 的代数闭扩张，$I_s$ 是 $S'$ 的在 $K_s$ 中取值的几何点集 (这是有 $n(s)$ 个元素的集合)，$I'_s$ 是 $I_s$ 去掉选定一个点的余集，最后 $I'$ 是诸 $I'_s$ 的并. 假设 $S'$ 连通. 则 $S$ 的基本群同构于由 $S'$ 的基本群和生成元 $g_i$ ($i \in I'$) 不受任何附加条件生成的 Galois 型群.
\end{corollary}

证明细节留给读者；所得断言只是下述事实在群论语言中的翻译：$S$ 的平展覆盖范畴 $\mathcal{C}$ 与 $S'$ 的平展覆盖 $X'$ 范畴 (对每个 $s \in T$，配备 $X'$ 在 $g^{-1}(s)$ 的诸点在 $K_s$ 中取值的 $n(s)$ 个纤维之间的传递双射系统) 等价. (当然在这个内蕴形式下，不再需要假设 $S'$ 连通).

\begin{example}\label{exa:IX.5.5}
容易证明代数闭域 $k$ 上的有理曲线 $\mathbf{P}^1_k$ 单连通\footnote{参见 XI 1.1.}. 故恰有一个二重点、有 $n$ 个解析分支的完备有理曲线的基本群是由 $n - 1$ 个生成元生成的自由 Galois 型群. 例如在普通二重点的情形，得到基本群 $\widehat{\mathbb{Z}}$，如 (I 11 a)) 中所宣告. 相反，尖点 (它是"几何单支"的点) 的存在对基本群没有影响.
\end{example}

\begin{corollary}\label{cor:IX.5.6}
设 $g \colon S' \to S$ 是纤维几何连通的万有浸没预概形态射，$S$ 连通. 则 $S'$ 连通，且选 $S'$ 中的几何点 $s'$、记 $s$ 为它在 $S$ 中的像，同态
\[
\pi_1(S', s') \longrightarrow \pi_1(S, s)
\]
满射. 若 $g$ 是对预概形的平展覆盖纤维范畴的有效下降态射 (参见 \ref{thm:IX.4.12})，引入 $S'' = S' \times_S S'$ 的几何点 $s'' = \mathrm{diag}(s')$，以及两个投影诱导的两个同态
\[
p_1^{*}, p_2^{*} \colon \pi_1(S'', s'') \longrightarrow \pi_1(S', s'),
\]
则 $\pi_1(S, s)$ 同构于在 Galois 型群范畴中这对态射的余核，i.e. 同构于 $\pi_1(S', s')$ 除以由形如 $p_1^{*}(g'') p_2^{*}(g'')^{-1}$ ($g'' \in \pi_1(S'', s'')$) 的元素生成的闭不变子群.
\end{corollary}

由 \ref{cor:IX.3.4} 确实知道函子 $X \mapsto X \times_S S'$ (从 $S$ 上平展覆盖到 $S'$ 上平展覆盖) 忠实全，这等价于基本群同态是满态射 (\cref{prop:V.6.9}). 最后一个断言是描述 \ref{thm:IX.5.1} 的直接推论.

\section{基本正合序列。具相对连通纤维态射的下降}

\begin{theorem}\label{thm:IX.6.1}
设 $S$ 是剩余域 $k$ 的 Artin 环 $A$ 的谱，$\overline{k}$ 是 $k$ 的代数闭包，$X$ 是 $S$-预概形，$X_0 = X \otimes_A \overline{k}$，$\overline{X}_0 = X \otimes_A \overline{k}$，$a$ 是 $X$ 的几何点，$\overline{a}$ 是它在 $X$ 中的像，$b$ 是它在 $S$ 中的像. 假设 $\overline{X}_0$ 在 $k$ 上拟紧且\emph{几何连通} (N.B. 若 $X$ 在 $S$ 上紧合，这意味着 $H^0(X_0, \mathcal{O}_{X_0})$ 是剩余域在 $k$ 上根式的 Artin 局部环). 则典范同态序列
\[
e \longrightarrow \pi_1(\overline{X}_0, \overline{a}) \longrightarrow \pi_1(X, a) \longrightarrow \pi_1(S, b) \longrightarrow e
\]
正合，且有
\[
\pi_1(S, b) \xleftarrow{\sim} \pi_1(\overline{k}, \overline{k}) = \overline{k} \text{ 在 } k \text{ 上的 伽罗瓦群}.
\]
\end{theorem}

由于基本群不随幂零元素去势而改变，可假设 $A = k$，这已使最后一个同构显然. 设 $k'$ 是 $k$ 在 $\overline{k}$ 中的可分闭包，考虑 $X' = X \otimes_k k'$ 和 $a$ 在 $X'$ 中的像 $a'$. 有典范同态序列
\[
e \longrightarrow \pi_1(\overline{X}_0, \overline{a}) \longrightarrow \pi_1(X', a') \longrightarrow \pi_1(S', b') \longrightarrow e
\]
(其中 $S' = \mathrm{Spec}(k')$). 最后，借助图式
\[
\begin{tikzcd}[row sep=large, column sep=large]
S & X \arrow[l] \\
S' \arrow[u] & X' \arrow[l] \arrow[u]
\end{tikzcd}
\]
有这个序列到 $X/k$ 的相应序列的典范同态. 另一方面看到，这个群序列的同态是同构，如 \ref{cor:IX.4.11} 所示. 故归结到证明第二个序列正合，i.e. 可假设 $k$ 完备. 设 $k_i$ 是 $k$ 在 $\overline{k}$ 中的有限 Galois 子扩张，令 $X_i = X \otimes_k k_i$，设 $a_i$ 是 $a$ 在 $X_i$ 中的像. 把验证自然同态
\[
\pi_1(\overline{X}_0, \overline{a}) \xrightarrow{\sim} \varprojlim_i \pi_1(X_i, a_i)
\]
是同构留给读者，它简单意味着 $X$ 的平展覆盖来自某个 $X_i$ 的平展覆盖，且后者在换成某个 $X_j$ ($j \geqslant i$) 的意义下本质唯一. 另一方面，设 $\pi_i$ 是 $k_i$ 在 $k$ 上的 伽罗瓦群，i.e. $S_i = \mathrm{Spec}(k_i)$ 的 $S$-自同构群的对偶群. 由于函子 $S' \mapsto X \times_S S'$ (从 $S$ 的平展覆盖到 $X$ 的平展覆盖) 忠实全 (\cref{cor:IX.3.4})，得出 $\pi_i$ 也同构于 $X$ 的连通主覆盖 $X_i$ 的 $X$-自同构群的对偶群. 故由 V 6.13 有正合序列
\[
e \longrightarrow \pi_1(X_i, a_i) \longrightarrow \pi_1(X, a) \longrightarrow \pi_i \longrightarrow e.
\]
对这些正合序列取 $i$ 的投射极限，得到正合序列 (因为在 Galois 型群范畴中)，它正是 \ref{thm:IX.6.1} 中考虑的序列. 这完成证明.

\ref{thm:IX.6.1} 中右正合性的几何翻译如下：

\begin{corollary}\label{cor:IX.6.2}
用前面记号，设 $X'$ 是 $X$ 的平展覆盖，$\overline{X}'_0$ 是 $\overline{X}_0$ 的对应平展覆盖. 下列条件等价：
\begin{enumerate}[i)]
\item 存在平展于 $S$ 上的 $S'$ 和 $X$-同构 $X' \xrightarrow{\sim} X \times_S S'$ ($S'$ 此时由 \ref{cor:IX.3.4} 在唯一同构意义下确定).
\item $\overline{X}'_0$ 在 $\overline{X}_0$ 上完全分解.
\end{enumerate}
若 $X'$ 连通，这些条件也等价于：
\begin{enumerate}[i)]
\item[(ii bis)] $\overline{X}'_0$ 在 $\overline{X}_0$ 上有截面.
\end{enumerate}
\end{corollary}

(N.B. 这最后一个补充本质；(i) 与 (ii) 的等价只意味着 $\pi_1(S, b)$ 是 $\pi_1(X, a)$ 除以由 $\pi_1(\overline{X}_0, \overline{a})$ 的像生成的闭不变子群所得的商群，而不是除以该像本身). 在上述条件下，称 $X'$ 是 $X$ 的\emph{几何平凡覆盖}\index[term]{jihepingfanfugai@几何平凡覆盖}.

\begin{remark}\label{rem:IX.6.3}
在断言 \ref{thm:IX.6.1} 中不能把 $k$ 换成 $k$ 的任意代数闭扩张，即使已假设 $k$ 代数闭. 换言之，一般不对：若 $X$ 是代数闭域 $k$ 上的连通代数概形，把 $k$ 换成代数闭扩张时它的基本群不变；例如在特征 $p > 0$ 中对 $k$ 上的仿射直线已经错了，因为无穷远点处有"高分歧"现象，蕴含基本群有"连续"结构. 我们将在下一讲看到，若 $X$ 在 $k$ 上紧合，这类现象不会发生. 还将在 $k$ 特征零时用超越方法证明同样结果.
\end{remark}

\begin{corollary}\label{cor:IX.6.4}
假设 $a$ 定位于 $X$ 中在 $k$ 上有理 (或更一般地剩余域在 $k$ 上根式) 的 $x$ 处. 则正合序列 \ref{thm:IX.6.1} 分裂.
\end{corollary}

可假设 $S = \mathrm{Spec}(k)$. 若 $x$ 在 $k$ 上有理，它对应 $X$ 在 $S$ 上的截面 $S \to X$，把 $b$ 变为 $a$，并定义同态 $\pi_1(S, b) \to \pi_1(X, a)$，即所求分裂. 若 $k(x)$ 在 $k$ 上根式，作基扩张 $\mathrm{Spec}(k(x)) \to \mathrm{Spec}(k)$ 归结到前一情形.

\begin{theorem}\label{thm:IX.6.5}
设 $f \colon X \to S$ 是纤维几何连通的紧合满射有限表现态射，$X'$ 是在 $X$ 上有限表现且紧合的预概形，$s$ 是 $S$ 的点，$F = X_s$ 是 $X$ 在 $s$ 处的纤维，$F'_1$ 是 $X'$ 在 $s$ 处的纤维 $F' = X'_s$ 的连通分量. 为使存在 $F'_1$ 在 $X'$ 中的开邻域 $X'_1$，平展 $S$-概形 $S'_1$ 和 $X$-同构 $X'_1 \xrightarrow{\sim} S'_1 \times_S X$，必要且充分条件是 $X'$ 在 $F'_1$ 的诸点处平展于 $X$ 上，且 $F'_1$ 是 $F$ 的几何平凡覆盖.
\end{theorem}

条件的必要性平凡，剩下证明充分性. 容易归结到 $S$ 诺特的情形. 考虑 $f$ 的 Stein 分解 $X \to T \to S$，其中 $T$ 是 $S$ 上代数 $f_*(\mathcal{O}_X)$ 的谱. 由于 $X$ 在 $S$ 上的诸纤维几何连通，且 $f$ 满射，态射 $T \to S$ 有限满射且根式，故 (\cref{thm:IX.4.10}) 每个平展于 $T$ 上的 $T'$ 来自某个平展于 $S$ 上的 $S'$ 的逆像. 这归结到把 $S$ 换成 $T$ 证明 \ref{thm:IX.6.5}，i.e. 在假设 $f_*(\mathcal{O}_X) = \mathcal{O}_S$ 的情形. 考虑紧合态射 $h \colon X' \to S$ 的 Stein 分解 $X' \to S' \to S$，其中 $S'$ 是代数 $h_*(\mathcal{O}_{X'})$ 的谱. 态射 $X' \to X$ 和 $X' \to S'$ 定义典范态射
\[
X' \longrightarrow X \times_S S',
\]
我们的断言包含在下列断言中：

\begin{corollary}\label{cor:IX.6.6}
设 $f \colon X \to S$ 是局部诺特预概形的紧合态射，使 $f_*(\mathcal{O}_X) = \mathcal{O}_S$，设 $X'$ 是在 $X$ 上紧合的预概形. 考虑 $X' \to S$ 的 Stein 分解 $X' \to S' \to S$ 和典范态射 $X' \to X \times_S S'$. 设 $s$ 是 $S$ 的点，$s'$ 是 $S'$ 的位于 $s$ 上方的点，对应 $X'$ 在 $s$ 处的纤维 $X'_s$ 的连通分量 $F'_1$. 为使态射 $X' \to X \times_S S'$ 在 $s'$ 的平展于 $S$ 上的开邻域 $U'$ 上方是同构，必要且充分条件是 $X'$ 在 $F'_1$ 的诸点处平展于 $X$ 上，且 $F'_1$ 是纤维 $F = X_s$ 的几何平凡覆盖.
\end{corollary}

必要性仍然平凡，剩下证明充分性. 结论也意味着 a) 由 $X' \to X \times_S S'$ 经基变换 $\mathrm{Spec}(\mathcal{O}_{b_{s'}}) \to S'$ 得到的态射是同构，b) $S'$ 在 $s'$ 处平展于 $S$ 上，i.e. $\mathcal{O}_{b_{s'}}$ 在 $\mathcal{O}_{b_s}$ 上平展. 在这种形式下，看到结论在基变换 $\mathrm{Spec}(\mathcal{O}_{b_s}) \to S$ 下不变. 假设也在此基变换下稳定，故可假设 $S$ 是完备诺特局部环的谱. 显然还可假设 $X'$ 连通，此时 $S' = \mathrm{Spec}(\mathcal{O}_{s'})$，$F' = F'_1$.

由于 $X'$ 在 $X$ 上平展的点的集开且包含纤维 $X'_{s'} = F'$，$X'$ 在 $S$ 上紧合，得出 $X'$ 在 $X$ 上平展. 由于它在 $F = X_s$ 上诱导的平展覆盖同构于 $F \otimes_{k(s)} L$，其中 $L$ 在 $k(s)$ 上平展，由 \ref{thm:IX.1.10} 得出它同构于形如 $X \times_S T$ 的覆盖，其中 $T$ 平展于 $S$ 上. (N.B. 这里仍只需用 \ref{thm:IX.1.10} 中函子忠实全的事实，它由 $X$ 上凝聚层的形式同构来自这些层的同构得出). 故若 $T$ 由 $A$ 上有限代数 $B$ 定义，$X'$ 等同于 $X$ 上代数 $\mathcal{O}_X \otimes_A B$ 的谱，由此立即得出 (因为 $f_*(\mathcal{O}_X) = \mathcal{O}_S$) $h_*(\mathcal{O}_{X'})$ 由 $B$ 定义，故典范态射 $X' \to X \times_S S'$ 正是所考虑的同构 $X' \xrightarrow{\sim} X \times_S T$. 这完成证明.

\begin{corollary}\label{cor:IX.6.7}
在 \ref{thm:IX.6.5} 的条件下，为使存在平展于 $S$ 上的预概形 $S'$ 和 $X$-同构 $X' \xrightarrow{\sim} X \times_S S'$，必要且充分条件是 $X'$ 在 $X$ 上平展，且对每个 $s \in S$，$X'_{s'}$ 是 $X_s$ 的几何平凡覆盖.
\end{corollary}

确实，若如此，$X'$ 是同构于平展于 $S$ 上的 $S'_i$ 的逆像的开集 $X'_i$ 的并. 容易看到这些 $S'_i$ 粘合成平展于 $S$ 上的 $S'$，并得到同构 $X' \xrightarrow{\sim} X \times_S S'$. 例如可以说诸 $X'_i$ 配备相对 $X \to S$ 的下降数据，它们必然粘合成整个 $X'$ 上相对 $X \to S$ 的下降数据. 由于后者在诸 $X'_i$ 上有效，容易得出 (借助第 4 号中遗漏的一个推论) 它有效. 也可以把 \ref{cor:IX.6.7} 表述成下列形式：

\begin{corollary}\label{cor:IX.6.8}
设 $f \colon X \to S$ 是纤维几何连通的紧合满射有限表现态射. 则 $f$ 对互相平展有限的预概形的纤维范畴是有效下降态射. 函子 $S' \mapsto X \times_S S'$ 诱导 $S$ 上平展有限预概形范畴与 $X$ 上对每个纤维 $X_s$ 诱导几何平凡覆盖的平展有限预概形范畴的等价.
\end{corollary}

\begin{remark}\label{rem:IX.6.9}
设 $f \colon X \to S$ 是紧合满射态射，$Y$ 局部诺特，则 $f$ 分解为满足 \ref{cor:IX.6.8} 假设的态射 $X \to S'$ 和服从 \ref{thm:IX.4.7} 的有限满射态射 $S' \to S$ 的复合，故为两个对互相平展有限的预概形的纤维范畴都是万有有效下降态射的态射的积. 可以由此得出 $f$ 本身对所述纤维范畴是万有有效下降态射. 于是用不同方法重新得到 4.12.
\end{remark}

\begin{remark}\label{rem:IX.6.10}
把 $f$ 紧合的假设换成：$X$ 在 $S$ 上有限型且在 $S$ 上有截面 (故 $f$ 万有浸没，且是对互相平展预概形的纤维范畴的下降态射)，\ref{cor:IX.6.7} 的结论不再成立，即使 $S$ 是离散赋值环的谱且 $X'$ 是 $X$ 的平展覆盖. 为看到这一点，从 $S$ 上紧合、一般纤维为非奇异有理曲线、特殊纤维 $Z_0$ 由两条共点直线组成的 $Z$ 出发. 例如，若 $t$ 是赋值环 $A$ 的幺模化元，取 $\mathbf{P}^2_A$ 的由齐次方程 $x^2 + y^2 + tz^2 = 0$ (齐次坐标 $x, y, z$) 定义的闭子概形 $Z$. 取 $X$ 为 $Z_0$ 的奇异点 $a$ 在并 $Z \cup \mathbf{P}^2_k$ 中的余集. $X$ 的诸纤维是 $\mathbf{P}^1_k$ 和 $\mathbf{P}^2_k - a$，故几何单连通 (i.e. 这样的纤维的每个平展覆盖几何平凡). 然而，如第 4 号那样进行，容易构造 $X$ 的平展覆盖，它不来自 $S$ 的平展覆盖，即把 $Z - a$ 和 $\mathbf{P}^2_k - a$ 的平凡覆盖粘合得到. 反之，把紧合性假设换成 $X$ 在 $S$ 上万有开有限表现，\ref{cor:IX.6.7} 的结论可能仍然成立\footnote{现在已证明，$g$ 只需万有开且满射；参见 SGA 4 XV 1.15.}. 至少假设 $X$ 在 $S$ 上的诸纤维几何不可约 (而不只是几何连通) 时如此. 只指出在这个问题中可以归结到 $S$ 是剩余域代数闭的完备离散赋值环的谱的情形.
\end{remark}

\ref{cor:IX.6.7} 在基本群术语中的解释如下：

\begin{corollary}\label{cor:IX.6.11}
设 $f \colon X \to S$ 是纤维几何连通的紧合满射有限表现态射. 假设 $X$ 故 $S$ 连通. 设 $a$ 是 $X$ 的几何点，$b$ 是它在 $S$ 中的像，对每个 $s \in S$，选 $k(s)$ 的代数闭包 $\overline{k(s)}$，在此扩张中取值的 $S_s$ 的几何点 $a_s$，以及从 $a_s$ 到 $a$ 的道路类，从而得到同态 $\pi_1(\overline{X}_s, a_s) \to \pi_1(X, a)$，其中 $\overline{X}_s = X_s \otimes_{k(s)} \overline{k(s)}$. 则同态 $\pi_1(X, a) \to \pi_1(S, b)$ 满射，且它的核是 $\pi_1(X, a)$ 的由诸 $\pi_1(\overline{X}_s, a_s)$ 的像生成的闭不变子群.
\end{corollary}

\begin{remark}\label{rem:IX.6.12}
在 \ref{cor:IX.6.7} 的条件下，假设 $S$ 诺特，容易看到使 $S'_{s'}$ 在 $X_s$ 上几何平凡的点 $s \in S$ 的集是可构集；若 $X'$ 在 $X$ 上紧合，它甚至开，如 \ref{cor:IX.6.6} 所见. 故若 $S$ 是 Jacobson 预概形 (例如域上有限型)，或 $X'$ 在 $X$ 上紧合，验证 \ref{cor:IX.6.7} 的条件时可限于 $S$ 的闭点 $s$. 同样，在 \ref{cor:IX.6.11} 中此时只需取 $S$ 的闭点的诸 $\pi_1(\overline{X}_s, a_s)$.
\end{remark}

\section{文献}\printbibliography[heading=none]


\end{refsection}
\begin{refsection}% =====================================================================
% 第十讲 (Exposé X) 基本群的特殊化理论
% (Théorie de la spécialisation du groupe fondamental)
% A. Grothendieck
% 原书页码: pp. 201-218
% =====================================================================
\chapter{基本群的特殊化理论}

在本讲中，我们限于研究紧合态射中几何纤维的基本群，i.e. 变动的紧合代数概形的基本群. 在以后的讲中，我们将把所用的技巧推广到"在无穷远处"\emph{温和分歧}的平展覆盖. 这将例如给出"三点问题"在阶与特征互素的 Galois 覆盖情形的解 (i.e. 确定 $\mathbf{P}^1_k$ 上至多在给定三点处分歧且在这些点温和分歧的 Galois 覆盖)，以及它的显然变体.

\section{紧合可分态射的同伦正合序列}

\begin{definition}\label{def:X.1.1}
域 $k$ 上的预概形 $X$ 称为可分\index[term]{kefentaishe@可分态射}，或在 $k$ 上可分，若对 $k$ 的任何扩张 $K$，$X \otimes_k K$ 既约. 若 $f \colon X \to Y$ 是预概形态射，称 $f$ 可分，或 $X$ 在 $Y$ 上可分，若 $X$ 在 $Y$ 上平坦，且对每个 $y \in Y$，纤维 $X \otimes_Y k(y)$ 在 $k(y)$ 上可分.
\end{definition}

若 $X$ 是域 $k$ 上的预概形，说它可分也意味着它既约，且诸 $k(x)$ (对 $X$ 的不可约分量的生成点 $x$) 是 $k$ 的可分扩张. 若 $k$ 完备，则说 $X$ 既约与之等价. 注意若 $X$ 在 $Y$ 上可分，则对任何基变换 $Y' \to Y$，$X' = X \times_Y Y'$ 在 $Y'$ 上可分. 在适当的有限性假设下还可证明两个\emph{可分态射}的复合可分. 我们只需下列形式：若 $X$ 在 $Y$ 上可分，且 $X'$ 在 $X$ 上平展，则 $X'$ 在 $Y$ 上可分. 这确实是定义和 I 9.2 的直接推论. 另一方面，"可分态射"假设将借助下列命题起作用：

\begin{proposition}\label{prop:X.1.2}
设 $f \colon X \to Y$ 是紧合可分态射，$Y$ 局部诺特，考虑它的 Stein 分解 $X \xrightarrow{f'} Y' \xrightarrow{f''} Y$ (其中 $f'_*(\mathcal{O}_X) = \mathcal{O}_{Y'}$，$Y'$ 在 $Y$ 上有限且同构于代数 $f_*(\mathcal{O}_X)$ 的谱). 则 $Y'$ 是 $Y$ 的平展覆盖.
\end{proposition}

这个命题将出现在 (EGA III 7)\footnote{参见 EGA III 7.8.10 (i).} 中. 指出证明的原理. 容易归结到 $Y$ 是完备局部环 $A$ 的谱的情形，再对后者作适当的有限平坦扩张 (对应适当的剩余扩张)，可以假设闭点 $y$ 的纤维的连通分量几何连通，这也意味着 $H^0(X_y, \mathcal{O}_{X_y})$ 分解成与 $k = k(y)$ 相同的体的乘积. 此时假设 $X$ 连通 (这是允许的)，将有 $H^0(X_y, \mathcal{O}_{X_y}) = k$，故同态 $A \to H^0(X_y, \mathcal{O}_{X_y})$ 满射. 由此用一般命题 (Künneth 型) 得出 $f_*(\mathcal{O}_X)$ 由 $A$ 上自由模 $B$ 定义，且 $B/\mathfrak{m}B \to H^0(X_y, \mathcal{O}_{X_y}) = k$ 双射. 故此时 $B$ 是 $A$ 上的平展代数，这完成证明.

\begin{theorem}\label{thm:X.1.3}
设 $f \colon X \to Y$ 是紧合可分态射，$Y$ 局部诺特且连通，假设 $f_*(\mathcal{O}_X) = \mathcal{O}_Y$ (这蕴含 $X$ 在 $Y$ 上的诸纤维几何连通，由 \ref{prop:X.1.2} 反之亦然). 设 $y$ 是 $Y$ 的点，$\overline{k(y)}$ 是 $k(y)$ 的代数闭包，$X_{\overline{y}} = X_y \otimes_{k(y)} \overline{k(y)}$. 最后设 $X'$ 是 $X$ 的连通平展覆盖，$\overline{X'}_y = X'_y \otimes_{k(y)} \overline{k(y)}$. 为使存在 $Y$ 的平展覆盖 $Y'$ 和 $X$-同构 $X' \xrightarrow{\sim} X \times_Y Y'$，必要且充分条件是 $\overline{X'}_y$ 在 $X_{\overline{y}}$ 上有截面.
\end{theorem}

令 $Y' = \mathrm{Spec}(h_*(\mathcal{O}_{X'}))$ (其中 $h \colon X' \to Y$ 是复合态射 $X' \to X \to Y$)，只需证明典范 $Y$-态射
\[
X' \longrightarrow X \times_Y Y'
\]
是同构，且 $Y'$ 在 $Y$ 上平展. 由 \ref{prop:X.1.2} 已知 $Y'$ 在 $Y$ 上平展，故 $X \times_Y Y'$ 在 $X$ 上平展，故态射 $X' \to X \times_Y Y'$ 也平展 (\cref{cor:I.4.8}). 而且，$Y'$ 作为连通集 $X'$ 的像连通，故由于 $X$ 在 $Y$ 上有连通纤维 (IX 3.4 和 V 6.9 (iii))，$X \times_Y Y'$ 连通. 故为证 $X' \to X \times_Y Y'$ 是同构，只需看它在 $X \times_Y Y'$ 的某点的投影次数等于 1. 这容易由 $\overline{X'}_y$ 在 $X_{\overline{y}}$ 上有截面的假设得出，或用 IX 6.6，或更简单地注意：只需在基变换 $\mathrm{Spec}(\overline{k(y)}) \to Y$ 后证明 $X \times_Y Y'$ 中存在这样的点，此时显然. 这完成 \ref{thm:X.1.3} 的证明.

考虑到 IX 3.4 和字典 V 6.9 与 V 6.11，可以把 \ref{thm:X.1.3} 写成下列等价形式

\begin{corollary}\label{cor:X.1.4}
用前面关于 $f \colon X \to Y$ 和 $X_{\overline{y}}$ 的记号，设 $\overline{a}$ 是 $X_{\overline{y}}$ 的几何点，$a$ 是它在 $X$ 中的像，$b$ 是它在 $Y$ 中的像. 则下列群同态序列正合：
\[
\pi_1(X_{\overline{y}}, \overline{a}) \longrightarrow \pi_1(X, a) \longrightarrow \pi_1(Y, b) \longrightarrow e.
\]
\end{corollary}

\begin{remark}\label{rem:X.1.5}
注意 \ref{thm:X.1.3} 的证明本质地用 \ref{prop:X.1.2}，从而用代数-形式几何的"第一比较定理". 相反，第 IX 讲的下降理论只通过 IX 3.4 出现，而在使 $f_*(\mathcal{O}_X) = \mathcal{O}_Y$ 的紧合态射 $f \colon X \to Y$ 的情形，\ref{cor:X.3.4} 的直接证明容易. 设 $Y'$ 平展于 $Y$ 上，假设 $X' = X \times_Y Y'$ 是两个非空开集的无交和，证明 $Y'$ 亦然. 确实有 $Y' = \mathrm{Spec}(\mathcal{A})$，故 $X' = \mathrm{Spec}(\mathcal{B})$，其中 $\mathcal{B} = \mathcal{A} \otimes_{\mathcal{O}_Y} \mathcal{O}_X$，而 $X'$ 的直和分解对应 $\mathcal{B}$ 分解成两个非零代数 $\mathcal{B}_1$ 和 $\mathcal{B}_2$ 的乘积. 由于 $f_*(\mathcal{O}_X) = \mathcal{O}_Y$，容易得出 $f_*(\mathcal{B}) = \mathcal{A}$，故 $\mathcal{A}$ 是两个代数 $f_*(\mathcal{B}_1)$ 和 $f_*(\mathcal{B}_2)$ 的和 (它们也非零，因为它们的单位截面非零)，证毕.
\end{remark}

\refstepcounter{theorem}\noindent\textbf{\thesection.\arabic{theorem}.}\quad 仍假设 $f$ 紧合且可分，但不再对 $f_*(\mathcal{O}_X)$ 作假设，它对应 $Y$ 的完全确定的平展覆盖 $Y'$，且在 $b$ 上方被 $a$ 的像 $b'$ 带点. 对典范态射 $X \to Y'$ 应用 \ref{cor:X.1.4}，并假设 $f$ 满射，正合序列 \ref{cor:X.1.4} 换成下列序列，它是代数拓扑中纤维空间的同伦正合序列的类似物：
\[
\pi_1(X_{\overline{y}}, \overline{a}) \longrightarrow \pi_1(X, a) \longrightarrow \pi_1(Y, b) \longrightarrow \pi_0(X_{\overline{y}}, \overline{a}) \longrightarrow \pi_0(X, a) \longrightarrow \pi_0(Y, b) \longrightarrow e.
\]
当然，在 \ref{cor:X.1.4} 中一般不能断言同态 $\pi_1(X_{\overline{y}}, \overline{a}) \to \pi_1(X, a)$ 单射；在代数拓扑中它的核是 $\pi_2(Y, b)$ 的像，在代数几何中也有必要引入各维数的同伦群，以及满足适当假设 (例如光滑态射) 的紧合态射的完整同伦正合序列. 目前除高同伦群的合理 (即使不是最终) 定义外，没有任何这方面的结果.

\begin{corollary}\label{cor:X.1.7}
设 $k$ 是代数闭域，$X$ 和 $Y$ 是 $k$ 上的两个连通预概形，假设 $X$ 在 $k$ 上紧合，$Y$ 局部诺特. 设 $a$ 是 $X$ 的几何点，$b$ 是 $Y$ 的在同一个 $k$ 的代数闭扩张 $K$ 中取值的几何点，考虑 $X \times_k Y$ 的几何点 $c = (a, b)$，以及由两个投影 $X \times_k Y \to X$ 和 $X \times_k Y \to Y$ 的基本群同态得到的同态 $\pi_1(X \times_k Y, c) \to \pi_1(X, a) \times \pi_1(Y, b)$. 前述同态是同构.
\end{corollary}

先假设 $K = k$. 令 $Z = X \times_k Y$，考虑投影 $f \colon Z \to Y$ 和 $Y$ 的几何点 $b$ 的局部化，对这种情况应用结果 1.4. 为此注意，换到 $X_{\mathrm{r\acute{e}d}}$ (不改变所考虑的基本群) 后可假设 $X$ 已既约故在 $k$ 上可分，故 $Z$ 在 $k$ 上可分，且显然有几何连通纤维 (因为 $X$ 连通). $Z$ 在 $b$ 处的几何纤维典范同构于 $X \otimes_k K = X$. 另一方面，由于复合态射 $X \to Z \to X$ 是恒等，得到 $\pi_1(X, a) \to \pi_1(Z, c)$ 单射，且 \ref{cor:X.1.4} 给出正合序列：
\[
e \longrightarrow \pi_1(X, a) \longrightarrow \pi_1(Z, c) \longrightarrow \pi_1(Y, b) \longrightarrow e.
\]
另一方面，有典范正合序列：
\[
e \longrightarrow \pi_1(X, a) \longrightarrow \pi_1(X, a) \times \pi_1(Y, b) \longrightarrow \pi_1(Y, b) \longrightarrow e,
\]
其中写的两个同态是典范单射和典范投影. 最后借助两端项上的恒等态射和中间项的典范同态 $\pi_1(Z, c) \to \pi_1(X, a) \times \pi_1(Y, b)$，有第一个正合序列到第二个的同态. 如此得到的图式的交换性平凡验证. 由于两端项上的同态是同构，中间项亦然，这在该情形证明 1.7.

当不再假设 $K = k$ 时，只得到同构
\[
\pi_1(Z, c) \longrightarrow \pi_1(X \otimes_k K, a) \times \pi_1(Y, b),
\]
而 \ref{cor:X.1.7} 等价于下列特殊情形：

\begin{corollary}\label{cor:X.1.8}
设 $X$ 是代数闭域 $k$ 上的紧合连通概形，$k'$ 是 $k$ 的代数闭扩张，$a'$ 是 $X \otimes_k k'$ 的几何点，$a$ 是它在 $X$ 中的像. 则典范同态 $\pi_1(X \otimes_k k', a') \to \pi_1(X, a)$ 是同构.
\end{corollary}

该同态满射的事实等价于：若 $X'$ 是 $X$ 的连通平展覆盖，则 $X' \otimes_k k'$ 也连通，它立即由 $k$ 代数闭的事实得出；它也是 IX 3.4 的特殊情形. 对 $X$ 的紧合性假设还没有用. 这样，说所考虑的同态单射也意味着：$X \otimes_k k'$ 的每个平展覆盖同构于 $X$ 的某个平展覆盖的逆像. 本质上是寻常的：可以找到 $k'$ 的有限型子 $k$-代数 $A$ 和 $X \otimes_k A$ 的平展覆盖，使它在 $X \otimes_k k'$ 上的逆像同构于给定覆盖. 设 $Y = \mathrm{Spec}(A)$，它是有限型的整 $k$-概形，故有 $k$ 上有理点. 对 $X \times Y$ 在 $k$ 上有理点 $(a, b)$ 处的基本群应用 \ref{cor:X.1.7}：得到 $X \times Y$ 的每个连通平展覆盖同构于覆盖 $X' \times Y'$ 的商，其中 $X'$，$Y'$ 分别是 $X$ 和 $Y$ 的群 $G$，$G'$ 的平展 Galois 覆盖，商按 $G \times G'$ 的子群 $H$. 这蕴含该 $X \times Y$ 覆盖在 $X \times Y'$ 上的逆像同构于形如 $X'_1 \times Y'$ 的覆盖，其中 $X'_1$ 是 $X$ 的平展覆盖. 故若 $L$ 是 $Y$ 的函数域，等于 $A$ 在 $k'$ 中的分式域，由给定 $X \times_k Y$ 覆盖诱导的 $X \otimes_k L$ 的平展覆盖，使存在 $L$ 的有限可分扩张 $L'$，使得所述覆盖在 $X \otimes_k L'$ 上的逆像同构于 $X'_1 \otimes_k L'$. 而 $k'$ 代数闭，可假设 $L$ 的扩张 $L'$ 包含在 $k'$ 中. 这证明给定的 $X \otimes_k k'$ 的平展覆盖同构于 $X'_1 \otimes_k k'$，证毕.

顺带指出的积 $X \times_k Y$ 的平展覆盖的显式形式立即蕴含下列结果：

\begin{corollary}\label{cor:X.1.9}
设 $k$ 是代数闭域，$X$ 和 $Y$ 是 $k$ 上两个局部诺特预概形，$Z = X \times_k Y$ 是它们的积，$Z'$ 是 $Z$ 的平展覆盖. 对每个在 $k$ 上有理的点 $y \in Y$，设 $i_y \colon \mathrm{Spec}(k) \to Y$ 是典范伴随的态射，$j_y = \mathrm{id}_X \times_k i_y$ 是对应的态射 $X \to Z$. 最后设 $X'_y$ 是 $Z'$ 经 $j_y$ 的逆像得到的 $X$ 的平展覆盖. 假设 $Y$ 连通，且 $X$ 或 $Y$ 在 $k$ 上紧合. 则 $X$ 的诸覆盖 $X'_y$ 全部同构.
\end{corollary}

形象地说，由连通预概形 $Y$ 参数化的 $X$ 的平展覆盖族是常的，若 $X$ 或参数预概形 $Y$ 在 $k$ 上紧合.

\begin{remark}\label{rem:X.1.10}
\cref{cor:X.1.7,cor:X.1.8,cor:X.1.9} 在正规代数概形的情形归功于 Lang-Serre \cite{exp10-2}. (他们的工作是这个讨论班发展的基本群理论的最初动机). 如这两位作者所指出的，放弃紧合性假设时这些结果变得不精确，至少在特征 $p > 0$ 时. 例如取 $X$ 为仿射直线 $X = \mathrm{Spec}(k[t])$，不难看到由方程
\[
x^p - x = st
\]
定义的、由仿射直线 $Y = \mathrm{Spec}(k[s])$ 参数化的 $X$ 的覆盖平展，且两两不同构. 这使 \ref{cor:X.1.9} 从而更不用说 \ref{cor:X.1.7} 失效，同样看到若 $s$ 看作 $k$ 的代数闭扩张 $K$ 中 $k$ 上的超越元，则找到 $X \otimes_k K$ 的平展覆盖 $X'$，它不来自 $X$ 的平展覆盖.
\end{remark}

\section{层存在定理的应用：紧合可分态射纤维基本群的半连续性定理}

\begin{theorem}\label{thm:X.2.1}
设 $Y$ 是剩余域 $k$ 的完备诺特局部环的谱，$X$ 是 $Y$-概形紧合，$X_0 = X \otimes_A k$，$a_0$ 是 $X_0$ 的几何点，$a$ 是 $X$ 的对应几何点. 则典范同态 $\pi_1(X_0, a_0) \to \pi_1(X, a)$ 是同构.
\end{theorem}

这只是在基本群语言中翻译 IX 1.10 中回顾的结果. 正是这里，代数-形式几何中的层存在定理本质地进入基本群理论.

现在引入剩余域 $k$ 的代数闭包 $\overline{k}$ 和几何纤维 $\overline{X}_0 = X_0 \otimes_k \overline{k}$. 有正合序列 (\cref{thm:IX.6.1})
\[
e \longrightarrow \pi_1(\overline{X}_0, \overline{a}) \longrightarrow \pi_1(X_0, a_0) \longrightarrow \pi_1(\overline{k}, \overline{k}) \longrightarrow e,
\]
另一方面有同构 \ref{thm:X.2.1} 和更初等的类似同构
\[
\pi_1(\overline{k}, \overline{k}) \longrightarrow \pi_1(Y, b),
\]
其中 $b$ 是 $a$ 在 $Y$ 中的像. 由此得到：

\begin{corollary}\label{cor:X.2.2}
用前面记号，假设 $\overline{X}_0 = X_0 \otimes_k \overline{k}$ 连通，设 $\overline{a}_0$ 是 $\overline{X}_0$ 的几何点，$a_0$ 是它在 $X$ 中的像，$b_0$ 是它在 $Y$ 中的像. 则下列典范同态序列
\[
e \longrightarrow \pi_1(\overline{X}_0, \overline{a}) \longrightarrow \pi_1(X, a_0) \longrightarrow \pi_1(Y, b_0) \longrightarrow e
\]
正合.
\end{corollary}

把这个序列与正合序列 \ref{cor:X.1.4} 比较，但注意 a) 对 $X \to Y$ 不必作平坦性或纤维可分性假设；b) 有重要补充：态射 $\pi_1(\overline{X}_0, \overline{a}_0) \to \pi_1(X, a_0)$ 单射.

这最后一个事实使我们能把 $X$ 在 $Y$ 上的其他几何纤维的基本群与 $\overline{X}_0$ 的比较. 设 $y_1$ 是 $Y$ 的任意点，$X_1$ 是它的纤维，$\overline{X}_1$ 是它的几何纤维 (相对 $k(y_1)$ 的某个代数闭扩张)，$a_1$ 是 $\overline{X}_1$ 的几何点，$a_1$ 是它在 $X$ 中的像，$b_1$ 是它在 $Y$ 中的像. 选择从 $a_1$ 到 $a_0$ 的"道路类"，从而有从 $b_1$ 到 $b_0$ 的道路类，从而有同态的交换图式：
\[
\begin{tikzcd}[row sep=large, column sep=large]
& \pi_1(\overline{X}_1, \overline{a}_1) \arrow[r] \arrow[d] & \pi_1(X, a_1) \arrow[r] \arrow[d] & \pi_1(Y, b_1) \arrow[r] & e \\
e \arrow[r] & \pi_1(\overline{X}_0, \overline{a}_0) \arrow[r] & \pi_1(X, a_0) \arrow[r] & \pi_1(Y, b_0) \arrow[r] & e,
\end{tikzcd}
\]
其中写的两个垂直箭头是同构. 由于第二行正合，找到典范同态，称之为基本群的特殊化同态\index[term]{jibenqunteshuhuatongtai@基本群特殊化同态} (只依赖所选的从 $a_1$ 到 $a_0$ 的道路类，故模 $\pi_1(X, a_0)$ 的内自同构确定)：
\[
\pi_1(\overline{X}_1, \overline{a}_1) \longrightarrow \pi_1(\overline{X}_0, \overline{a}_0).
\]
当上面第一行也正合时，立即得出特殊化同态满射. 因此考虑到 \ref{cor:X.1.4} 得到：

\begin{corollary}\label{cor:X.2.3}
在 \ref{thm:X.2.1} 的条件下，再假设态射 $f \colon X \to Y$ 可分 (\cref{def:X.1.1}) 且 $\overline{X}_0$ 连通 (故由 \ref{prop:X.1.2} 有 $f_*(\mathcal{O}_X) = \mathcal{O}_Y$). 则对 $X$ 在 $Y$ 上的每个几何纤维 $\overline{X}_1$ (配备几何点 $a_1$)，上面定义的特殊化同态是满同态.
\end{corollary}

这是基本群的半连续性结果，它在代数拓扑中似乎还没有类似物. 还可以在更一般的条件下陈述：

\begin{corollary}\label{cor:X.2.4}
设 $f \colon X \to Y$ 是纤维几何连通的紧合态射，$Y$ 局部诺特，$y_0$ 和 $y_1$ 是 $Y$ 的两个满足 $y_0 \in \overline{\{y_1\}}$ 的点，$\overline{X}_0$ 和 $\overline{X}_1$ 是对应 $k(y_0)$ 和 $k(y_1)$ 的给定代数闭扩张的几何纤维，$a_0$ resp. $a_1$ 是 $\overline{X}_0$ resp. $\overline{X}_1$ 的几何点. 则可以自然方式定义特殊化同态：
\[
\pi_1(\overline{X}_1, \overline{a}_1) \longrightarrow \pi_1(\overline{X}_0, \overline{a}_0),
\]
它在模内自同构意义下确定，且若 $f$ 是可分态射 (\cref{def:X.1.1}) 则是满同态.
\end{corollary}

事实上，首先由 \ref{cor:X.1.8} 得出 \ref{cor:X.2.4} 本质上与所选剩余域 $k(y_0)$ 和 $k(y_1)$ 的代数闭扩张无关. 这使我们能把 $Y$ 换成 $Y$ 上有位于 $y_0$ (resp. $y_1$) 上方的点 $y'_0$ (resp. $y'_1$) 的预概形 $Y'$. 取 $Y'$ 为 $Y$ 中 $y_0$ 的局部环的完备化的谱，应用 2.3.

\begin{remark}\label{rem:X.2.5}
\ref{cor:X.2.4} 关于特殊化同态满射的最终结论，更不用说作为其推论的结果 \ref{thm:X.1.3} 和 \ref{cor:X.1.4}，在不再假设 $f \colon X \to Y$ 是可分态射时变得不精确，即使在特征 0 的代数闭域上的射影概形情形. 我们后面将看到例子，既在 $f$ 平坦但 $f$ 有不可分纤维的情形 ($X$ 和 $Y$ 在 $k$ 上仍光滑)，也在 $f$ 的诸纤维确实可分但 $f$ 不平坦的情形 (例如 $f \colon X \to Y$ 是整正规概形的双有理态射)，参见 XI 3. 在这些例子中，一般几何纤维的基本群可能为零，而某个适当的特殊几何纤维的基本群非零. 另一方面，即使 $f \colon X \to Y$ 如 \ref{cor:X.2.4} 那样是紧合可分态射，特殊化态射也经常不是同构. 例如，容易给出 $\overline{X}_1$ 是非奇异椭圆曲线的例子 (故它的基本群交换，且对与特征互素的素数 $\ell$，它的 $\ell$-准素分量同构于 $\mathbb{Z}_\ell^2$，参见 XI)，而 $\overline{X}_0$ 由两条在两点相交的非奇异有理曲线组成，或由两条在一点相切的有理曲线组成，或最后由一条有尖点奇点的有理曲线组成. (关于退化椭圆曲线的完整分类，参见 Kodaira \cite{exp10-1} 和 Néron 的近期工作). 此时在这些情形中，$\overline{X}_0$ 的基本群分别是 $\widehat{\mathbb{Z}}$，$e$，$e$，故"严格小于"$\overline{X}_1$ 的. 我们后面将在 $f$ 是光滑态射时看到特殊化同态核的上界，它特别蕴含若 $k(y_0)$ 特征 0，则特殊化同态是同构. 但即使对光滑态射，若 $k(y_0)$ 的特征 $> 0$，特殊化同态也可能不是同构，例如 $X$ 是 $Y$ 上的阿贝尔概形 (若愿意，相对维数 1) 的情形，参见 XI 2. 基本群特殊化的满意理论必须考虑"真正的"基本群的"连续分量"，它对应无穷小群作为结构群的覆盖的分类；由此有理由期待：紧合光滑态射 $f \colon X \to Y$ 的几何纤维的"真正"基本群在 $X$ 上构成漂亮的局部系，它是 $X$ 上有限平坦群概形的投射极限\footnote{这个极具吸引力的猜想不幸被 M. Artin 的一个未发表例子否定，甚至在 $f$ 的诸纤维是给定亏格 $g \geqslant 2$ 的代数曲线时.}. 我们以后将回到这个观点，而当前的对象正相反，是尽量推进拓扑理论与基本群的概形理论共同的现象.
\end{remark}

现在设 $X_0$ 是代数闭域 $k$ 上亏格 $g$ 的紧合、光滑且连通的曲线. 若 $k$ 特征零，它的基本群可以用超越方式如下确定. 已知 $X_0$ 来自定义在素数域 $\mathbb{Q}$ 的有限超越次数的代数闭扩张上的曲线的基扩张，考虑到 \ref{cor:X.1.8}，可假设 $k$ 本身在 $\mathbb{Q}$ 上超越次数有限. 故可假设 $k$ 是复数域 $\mathbb{C}$ 的子域，再一次应用 \ref{cor:X.1.8} 可假设 $k = \mathbb{C}$. 此时不难验证 $X$ 的基本群同构于关联拓扑空间 $\tilde{X}$ (亏格 $g$ 的紧定向曲面) 的基本群关于有限指数子群定义的拓扑的完备化\footnote{这个推导在一个未整理的讲授中明确过.}. 另一方面，众所周知拓扑基本群由 $2g$ 个生成元 $s_i$，$t_i$ ($1 \leqslant i \leqslant g$) 生成，只服从一个关系：
\[
(s_1 t_1 s_1^{-1} t_1^{-1}) \cdots (s_g t_g s_g^{-1} t_g^{-1}) = 1.
\]
故 $X$ 的基本群有 $2g$ 个拓扑生成元 $s_i$，$t_i$ ($1 \leqslant i \leqslant g$)，只服从前面关系. 若 $k$ 的特征为 $p > 0$，用 $A$ 记由 $k$ 构造的 Witt 向量环，用 $K$ 记它的分式域的代数闭扩张. 在 III 7.4 中已见存在 $Y = \mathrm{Spec}(A)$ 上的概形 $X$，在 $Y$ 上紧合且光滑，约化为 $X_0$. 对它应用 \ref{cor:X.2.3}，找到满态射
\[
\pi_1(X_1) \longrightarrow \pi_1(X_0),
\]
其中 $X_1 = X \otimes_A K$. 立即\footnote{参见 EGA IV 12.2.} $X_1$ 在 $K$ 上光滑、连通 (\cref{prop:X.1.2})、维数 1，且它的亏格等于 $g$ (由 Euler-Poincaré 特征的不变性，参见 EGA III 7). 由于 $K$ 特征 0，可以应用前面结果. 于是用超越方式证明：

\begin{theorem}\label{thm:X.2.6}
设 $X_0$ 是代数闭域 $k$ 上的光滑紧合连通代数曲线，$g$ 是它的亏格. 则 $\pi_1(X_0)$ 有 $2g$ 个拓扑生成元的生成系，只服从上面写的关系. 当 $k$ 的特征为 0 时，$\pi_1(X_0)$ 甚至是前面生成元和关系的自由 Galois 型群.
\end{theorem}

\begin{remark}\label{rem:X.2.7}
据编写者所知，目前没有用纯代数方式证明前面结果的方法 (除亏格 0，1 外). 首先，看不出如何在 $\pi_1(X_0)$ 中区分 $2g$ 个元素，然后才能期望它们构成拓扑生成元系. 在这方面，去掉 $n$ 个点的有理直线及其中在这 $n$ 点温和分歧的覆盖研究的情形更友好，因为考虑在这 $n$ 点的分歧群提供待研究基本群的 $n$ 个元素，确实证明它们拓扑生成该基本群，如我们后面将看到的\footnote{参见第十二讲. 而且这些元素只在模共轭意义下真正确定，需要在它们的类中明智地同时选择.}. 但即使在这个特别具体的情形，似乎也不存在纯代数证明. 这样的证明显然极其有趣. 关于曲线基本群唯一能用纯代数方式证明的事实 (除下面 \ref{cor:X.2.12} 的弱有限性定理，Lang-Serre \cite{exp10-2} 用代数方式证明)，似乎是通过 Jacobi 簇确定化为阿贝尔群的基本群 (在 IX 5.8 末行指出).
\end{remark}

\refstepcounter{theorem}\noindent\textbf{\thesection.\arabic{theorem}.}\quad \ref{thm:X.2.6} 的最后一个断言在特征 $p > 0$ 时不再成立，如在椭圆曲线情形已经看到. 如已指出的，我们不知道 $X_0$ 的基本群是否拓扑有限表现；这似乎完全不可能.

\begin{theorem}\label{thm:X.2.9}
设 $k$ 是代数闭域，$X$ 是 $k$ 上紧合连通概形. 则 $X$ 的基本群拓扑有限生成.
\end{theorem}

对 $n = \dim X$ 作递归，$n \leqslant 0$ 时断言平凡. 假设 $n > 0$ 且定理对维数 $n' < n$ 已证. 由 Chow 引理 (EGA II 5.6.2) 存在 $k$ 上射影概形 $X'$ 和满射态射 $X' \to X$. 显然可假设 $X'$ 既约，再换到正规化后正规. 借助下降理论，只需证明 $X'$ 的诸连通分量的基本群拓扑有限生成 (\cref{cor:IX.5.2}). 这归结到 $X$ 射影且正规的情形. 若 $n = 1$，只需应用 2.6. 若 $n \geqslant 2$，考虑射影浸入 $X \to \mathbf{P}^r_k$，和满足 $Y \neq X$ (i.e. $H \not\supset X$) 的超平面截面 $Y = X \cdot H$ (配备诱导既约结构). 则 $\dim Y < n$，考虑到递归假设，只需证明 $\pi_1(Y) \to \pi_1(X)$ 满射. 更一般地：

\begin{lemma}\label{lem:X.2.10}
设 $X$ 是代数闭域 $k$ 上的紧合预概形，$g \colon X \to \mathbf{P}^r_k$ 是态射. 假设 $X$ 不可约且正规，$\dim g(X) \geqslant 2$. 设 $H$ 是 $\mathbf{P}^r_k$ 的超平面，$Y = X \times_{\mathbf{P}^r_k} H$. 则 $Y$ 连通，且同态 $\pi_1(Y) \to \pi_1(X)$ 满射.
\end{lemma}

这些断言确实由下列断言得出：

\begin{corollary}\label{cor:X.2.11}
在先前条件下，设 $X'$ 是 $X$ 的连通平展覆盖，$Y' = X' \times_X Y = X' \times_{\mathbf{P}^r_k} H$ 是 $Y$ 上诱导的覆盖. 则 $Y'$ 连通.
\end{corollary}

由于 $X$ 正规，$X'$ 正规，故连通时 $X'$ 不可约；且它在 $\mathbf{P}^r_k$ 中的像维数 $\geqslant 2$. 一个归功于 Zariski 的熟知引理 (称为"Bertini 定理") 蕴含：若 $H_1$ 是 $\mathbf{P}^r_k$ 中定义在 $k$ 的扩张 $K$ 上的一般超平面，则 $X' \times_{\mathbf{P}^r} H_1$ 在 $K$ 上万有不可约，故万有连通. Zariski 连通定理 (EGA III 4) 蕴含对每个超平面 $H$ (定义在 $k$ 的任意扩张上) $X' \times_{\mathbf{P}^r} H$ 几何连通. 这完成 \ref{cor:X.2.11} 的证明，从而 \ref{thm:X.2.9} 的证明.

\begin{corollary}\label{cor:X.2.12}
(Lang-Serre) 在 \ref{thm:X.2.9} 的条件下，对每个有限群 $G$，$X$ 的群 $G$ 的主覆盖的同构类的集有限.
\end{corollary}

\begin{remark}\label{rem:X.2.13}
在 \ref{lem:X.2.10} 的条件下，当 $\dim g(X) \geqslant 3$ (至少在 $g$ 是浸入且 $X$ 正则时) 我们将证明更精确的结果，在代数几何中以"Lefschetz 定理"之名著称：$\pi_1(Y) \to \pi_1(X)$ 是同构.\footnote{参见本讨论班的续篇 SGA 2 (1962)，定理 X 3.10.} 在经典情形对同调群和高同伦群有类似断言，它们迟早应被纳入抽象代数几何. 即使对 Hodge 上同调 $H^p(X, \Omega^q)$，问题似乎也还未被研究；对后者，Lefschetz 定理在特征 $p > 0$ 时原样成立也不太可能.

\begin{remark}[2003 年 (MR) 增补]
设 $R$ 是剩余域 $k$ 代数闭且特征 $p > 0$、分式域 $K$ 的完备离散赋值环，设 $Y$ 是 $R$ 上亏格 $g$ 的紧合光滑连通曲线. 有满的特殊化态射 $sp \colon \pi_1(Y_K) \to \pi_1(Y_k)$. 已注意到若 $K$ 特征 0，则只要 $g \geqslant 1$，$sp$ 就不是同构. 假设 $K$ 特征 $p$，故 $R$ 同构于形式幂级数环 $k[\![T]\!]$. 在等特征 $p > 0$ 时，基本群不由亏格 $g$ 确定，如在椭圆曲线 (可以是普通或超奇异的) 上已经看到. 引用 A. Tamagawa 的近期结果 (尚未发表). 若 $G$ 是 pro-有限群，记 $G^{\mathrm{r\acute{e}s}}$ 为 $G$ 的 pro-有限商群，即 $G$ 的有限可解拓扑商的投射极限.

\begin{theorem}[A. Tamagawa]
假设 $g \geqslant 2$，特殊纤维 $Y_k$ 在有限域上可定义，且由 $sp$ 取商得到的态射 $sp^{\mathrm{r\acute{e}s}} \colon \pi_1(Y_K)^{\mathrm{r\acute{e}s}} \to \pi_1(Y_k)^{\mathrm{r\acute{e}s}}$ 双射. 则曲线 $Y$ 在 $R$ 上为常.
\end{theorem}

注意 $\overline{K}/K$ 的 伽罗瓦群可解. 前述断言还可以翻译成：假设一般纤维 $Y_K$ 的每个有限平展覆盖 $X_K \to Y_K$ (可解 伽罗瓦群的 Galois 覆盖) 都有潜在好约化 (即在 $K$ 的某个有限扩张中用 $R$ 在其中的正规化替换 $R$ 后，延拓成 $Y$ 的有限平展覆盖). 则 $Y$ 在 $R$ 上为常.
\end{remark}
\end{remark}

\section{纯度定理的应用：紧合光滑态射纤维基本群的连续性定理}

不证地回忆\footnote{证明参见 SGA 2 X 3.4.}：

\begin{theorem}\label{thm:X.3.1}
[\emph{纯度定理} (Zariski-Nagata)]\index[term]{chundudingli@纯度定理} 设 $f \colon X \to Y$ 是整预概形的拟有限支配态射，$X$ 正规，$Y$ 正则局部诺特，设 $Z$ 是 $f$ 不平展 (i.e. $f$ 分歧，两者相同，I 9.5 (ii)) 的 $X$ 的点集. 若 $Z \neq X$，则 $Z$ 在它的每个点处余维数 1，i.e. 对 $Z$ 的每个不可约分量 $Z'$ (生成点 $z$)，$\mathcal{O}_{X,z}$ 的 Krull 维数等于 1.
\end{theorem}

回忆预概形称为正规 resp. 正则，若它的局部环正规 resp. 正则，且关系 $Z \neq X$ 也意味着有限扩张 $R(Z)/R(X)$ (其中 $R$ 表示有理函数域) 可分. 在 $Z$ 的分量 $Z'$ 的生成点 $z$ 处，并在 $z$ 下方的 $Y$ 的点 $y$ 处局部化，得到等价断言：

\begin{corollary}\label{cor:X.3.2}
设 $A$ 是诺特正则局部环，$A \to B$ 是单射局部同态，使 $B$ 正规，是 $A$ 上有限型代数的局部化，且在 $A$ 上拟有限；再假设 $\dim A (= \dim B) \geqslant 2$，且对 $B$ 的每个异于极大理想的素理想 $\mathfrak{p}$，$B$ 在 $\mathfrak{p}$ 处平展于 $A$ 上，i.e. $B_{\mathfrak{p}}$ 在 $A_{\mathfrak{q}}$ (其中 $\mathfrak{q} = A \cap \mathfrak{p}$) 上平展. 则 $B$ 在 $A$ 上平展.
\end{corollary}

而且不难把最后一个断言归结到 $A$ 是完备局部环、故 $B$ 在 $A$ 上有限的情形. Zariski \cite{exp10-5} 给出这个结果的简单证明，在等特征情形有效；一般情形归功于 Nagata \cite{exp10-3}，他依靠 Chow 的精细结果；后者没有被讨论班的任何参加者验证过，应成为以后讲授的对象. 这里只指出在 $\dim A = 2$ 的特殊情形下非常简单的证明，它足以用于我们本节的这个最重要应用. 由于 $B$ 正规，它是深度 (旧术语：上同调余维数) $\geqslant 2$ 的 $B$-模，故也是深度 $\geqslant 2$ 的 $A$-模，且 $A$ 维数 2 正则，由此得出 $B$ 是 $A$ 上的自由模\footnote{参见 EGA IV$_0$ 17.3.4.}. 于是由 I 4.10 得出 $B$ 在 $A$ 上分歧的素理想 $\mathfrak{q}$ 的集是 $\mathrm{Spec}(A)$ 的由主理想 (由 $B$ 在 $A$ 上的基的判别式生成) 定义的部分，故它包含在 $\mathrm{Spec}(A)$ 的闭点中时为空，这证明 $\dim A = 2$ 时的 3.2.

主要用 \ref{thm:X.3.1} 的等价形式：

\begin{corollary}\label{cor:X.3.3}
设 $X$ 是局部诺特正则预概形，$U$ 是 $X$ 的余维数 $\geqslant 2$ 的闭部分 $Z$ 的余集的开部分. 则函子 $X' \mapsto X' \times_X U$ (从 $X$ 的平展覆盖范畴到 $U$ 的平展覆盖范畴) 是范畴等价；特别，若 $a$ 是 $U$ 的几何点，典范同态 $\pi_1(U, a) \to \pi_1(X, a)$ 是同构.
\end{corollary}

最后一个断言显然是第一个的推论，对第一个显然可假设 $X$ 连通故不可约. 由 $X$ 的正规性已得函子 $X' \mapsto X' \times_X U$ (从 $X$ 的局部自由 (未必平展) 覆盖范畴到 $U$ 的覆盖范畴) 忠实全，因为函子 $E \mapsto E|_U$ (从 $X$ 上局部自由模范畴到 $U$ 上局部自由模范畴) 忠实全. 剩下证明对 $U$ 的每个平展覆盖 $U'$，存在 $X$ 的平展覆盖 $X'$ (由前面所述必然唯一)，使 $U'$ 同构于 $X' \times_X U$. 显然可假设 $U'$ 连通，故不可约，因为 ($U$ 正规时) $U'$ 正规. 设 $K$ 是 $X$ (或 $U$，相同) 上的有理函数域，$K'$ 是 $U'$ 的，则 $U'$ 等同于 $U$ 在 $K'$ 中的正规化 (\cref{cor:I.10.3}). 设 $X'$ 是 $X$ 在 $K'$ 中的正规化 (EGA II 6.3)，则 $X' \times_X U \xrightarrow{\sim} U'$，另一方面 $X'$ 正规、整，结构态射 $f \colon X' \to X$ 有限且支配 (因为 $X$ 正规且 $K'/K$ 是有限可分扩张). 它在 $f^{-1}(U) = X' - f^{-1}(Z)$ 上平展，且由于 $Z$ 在 $X$ 中余维数 $\geqslant 2$，$f^{-1}(Z)$ 在 $X'$ 中余维数 $\geqslant 2$. 于是由 \ref{thm:X.3.1} 得出 $X'$ 在 $X$ 上平展，这完成证明.

现在设 $f \colon X \to Y$ 是从局部诺特正则预概形 $X$ 到预概形 $Y$ 的有理映射，假设 $f$ 在余维数 $\geqslant 2$ 的闭部分的余集的开集 $U$ 上定义. 则由 \ref{cor:X.3.3} 推出函子 (在同构意义下确定) 从 $Y$ 的平展覆盖范畴到 $X$ 的平展覆盖范畴，从而对 $U$ 的每个像为 $b$ 的几何点 $a$，有典范同态
\[
\pi_1(X, a) \longrightarrow \pi_1(Y, b)
\]
(由典范同态 $\pi_1(U, a) \to \pi_1(Y, b)$ 借助同构 $\pi_1(U, a) \xrightarrow{\sim} \pi_1(X, a)$ 得到). 当 $f$ 是支配态射，$X$ 和 $Y$ 整、函数域为 $K$ 和 $L$ (故 $K$ 是 $L$ 的扩张)，$Y$ 正规时，这些对应借助域扩张明确：注意对 $L$ 的每个在 $Y$ 上非分歧的有限扩张 $L'$，代数 $K' = L' \otimes_L K$ 在 $X$ 上非分歧.

特别，这些思考表明：局部诺特连通正则预概形 (由余维数 $\leqslant 1$ 处局部化的几何点带点) 的基本群，在取该范畴中的态射为在余维数 $\geqslant 2$ 闭部分的余集上定义的支配有理映射时，是函子. 回忆例如域 $k$ 上正规概形到 $k$ 上紧合概形的有理映射在余维数 $\geqslant 2$ 的集的余集上定义，得到：

\begin{corollary}\label{cor:X.3.4}
(基本群的双有理不变性)\index[term]{jibenqunshuangyoulibubianxing@基本群双有理不变性} 设 $k$ 是域，$X$ 和 $Y$ 是 $k$ 上紧合的正则概形，$f \colon X \to Y$ 是从 $X$ 到 $Y$ 的双有理映射，$\Omega$ 是 $X$ 的函数域 $K$ 的代数闭扩张，用来定义 $X$ 的基本群和 $Y$ 的基本群. 则它们典范同构.
\end{corollary}

这也意味着：对 $K$ 的有限扩张 $K'$，若它在 $K$ 的某个紧合非奇异"模型" $X$ 上非分歧，则它在任何其他紧合非奇异模型上都非分歧.

\begin{remark}\label{rem:X.3.5}
关于纯度定理的其他应用，参见 Abhyankar 的工作，在 \cite{exp10-4} 中阐述，受 Zariski \cite[第八章]{exp10-6} 的结果启发，后者用拓扑方式证明. 它们还远未被"抽象"代数几何同化，值得新的努力.
\end{remark}

将需要分歧理论的一些初等事实. 设 $V$ 是分式域 $K$、剩余域 $k$ 的离散赋值环，$L$ 是 $K$ 的群 $G$ 的 伽罗瓦扩张，$V'$ 是 $V$ 在 $L$ 中的正规化 (它是 $V$ 上秩 $n = [L : K]$ 的自由模)，$\mathfrak{m}'$ 是 $V'$ 的极大理想，$G_d$ 是 $G$ 的保持 $\mathfrak{m}'$ 不变的元素的子群，故 $G_d$ 作用在 $k$ 的剩余扩张 $k' = V'/\mathfrak{m}'$ 上，$G_i$ 是 $G_d$ 的平凡作用的元素的子群 (回忆 $G_d$ 和 $G_i$ 分别称为 $G$ 的分解子群和惯性子群). 称 $L$ 在 $V$ 上温和分歧，若 $n_i = [G_i : e]$ 的阶与 $k$ 的特征互素 (若 $k$ 特征 0 该条件总满足). 众所周知此时 $G_i$ 典范嵌入 $k'^{*}$，故同构于 $k'$ 中 $n_i$ 次单位根群，这特别蕴含 $G_i$ 循环. 这种情况的典型情形是令 $L = K[t]/(t^n - u)$，其中 $u$ 是 $V$ 的幺模化元，$n$ 是与 $p$ 互素的整数：若 $K$ 含 $n$ 次单位根，$L$ 是 $K$ 的完全分歧 伽罗瓦扩张，伽罗瓦群 $G = G_i$ 同构于 $\mathbb{Z}/n\mathbb{Z}$.

\begin{lemma}\label{lem:X.3.6}
(Abhyankar 引理) 设 $V$ 是分式域 $K$ 的离散赋值环，$L$ 和 $K'$ 是 $K$ 的在 $V$ 上温和分歧的两个 伽罗瓦扩张，$n$ 和 $m$ 是对应惯性群的阶，$L'$ 是 $L$ 和 $K'$ 在 $K$ 上的复合扩张. 若 $m$ 是 $n$ 的倍数，则 $L'$ 在 $K'$ 的正规化 $V'$ 的诸局部化上非分歧.
\end{lemma}

设 $W'$ 是 $V'$ 在 $L'$ 中的正规化，$\mathfrak{m}'$ 是 $V'$ 的极大理想，$\mathfrak{n}'$ 是 $W'$ 的位于 $\mathfrak{m}'$ 上方的极大理想，$\mathfrak{n}$ 是它在 $V$ 在 $L$ 中的正规化 $W$ 上诱导的极大理想，$G$，$H$，$M$ 是 $L$，$K'$，$L'$ 在 $K$ 上的 伽罗瓦群，$G_i$，$H_i$，$L'_i$ 是对应所选极大理想的惯性群. 则 $M$ 嵌入乘积 $G \times H$，$M_i$ 嵌入乘积 $G_i \times H_i$，且投影 $M \to G$ 和 $M \to H$，$M_i \to G_i$ 和 $M_i \to H_i$ 满射 (中间域推论). 由此已经得出，由于由假设 $G_i$ 和 $H_i$ 是阶 $m$ 和 $n$ (与 $p$ 互素) 的循环群，$M_i$ 的阶与 $p$ 互素，故循环，且由于 $m$ 是 $n$ 的倍数故 $G_i \times H_i$ 的元素 $m$ 次幂为零，$M_i$ 的阶整除 $m$，故等于 $m$，因为 $M_i \to H_i$ 满射. 故这最后一个同态也单射. 而它的核是 $\mathfrak{n}'$ 在 $\mathfrak{m}'$ 上的惯性群，这证明 $L'$ 在 $\mathfrak{n}'$ 处非分歧于 $K'$ 上. 由此得引理.

现在置于 \ref{cor:X.2.4} 的条件下，有特殊化同态
\[
\pi_1(\overline{X}_1, \overline{a}_1) \longrightarrow \pi_1(\overline{X}_0, \overline{a}_0)
\]
它相对紧合可分态射 $f \colon X \to Y$ 满射. 想明确这个同态的核. 如 \ref{cor:X.2.4} 的证明那样进行，看到在这个问题中总可以假设 $Y$ 是离散赋值环 $V$ 的谱，完备且剩余域代数闭 (因为总可以找到这样的环和它的谱到 $Y$ 的、像为 $\{y_0, y_1\}$ 的态射). 则 $X_0 = \overline{X}_0$，$k(y_0) = k = V$ 的剩余域，$k(y_1) = K = V$ 的分式域. 设 $K_s$ 是 $K$ 的可分闭包，$\overline{K}$ 是它的代数闭包，对 $K$ 的每个含 $V$ 的子环 $W$，令 $X_W = X \otimes_V W$. 特别有
\[
X_V = X, \quad X_K = X_1, \quad X_{\overline{K}} = \overline{X}_1.
\]
而且典范态射 $\overline{X}_1 = X_{\overline{K}} \to X_{K_s}$ 诱导基本群同构 (\cref{cor:IX.4.11})，故考虑到同构 \ref{thm:X.2.1} $\pi_1(X_0) \to \pi_1(X)$，归结到研究与典范态射 $X_{K_s} \to X$ 关联的满同态
\[
\pi_1(X_{K_s}) \longrightarrow \pi_1(X).
\]
确定后者的核归结到下列问题的解：有 $X_{K_s}$ 的群 $G$ 的连通主覆盖 $Z_{K_s}$ (故对应 $\pi_1(X_{K_s})$ 到 $G$ 的同态)，确定在什么条件下它同构于 $X$ 的群 $G$ 的主覆盖 $Z$ 的原像.

先注意 $K_s$ 是它的有限子扩张 $K'$ 在 $K$ 上的递增滤过并，从而 $Z_{K_s}$ 同构于 $X_{K'}$ 的某个主覆盖 $Z_{K'}$ (对适当的 $K'$) 的逆像 (注意对固定的 $K'$，$Z_{K'}$ 不唯一确定). 说 $Z_{K_s}$ 同构于 $X$ 的主覆盖 $Z$ 的逆像，意味着存在 $K_s$ 的有限子扩张 $K'' \supset K'$，使 $Z_{K''} = Z_{K'} \otimes_{K'} K''$ 同构于 $Z \otimes_V K''$.

现在对 $K_s$ 的有限子扩张 $K'$，用 $V'$ 记 $V$ 在 $K'$ 中的正规化，它是离散赋值环，完备，剩余域 $k$. 故典范态射 $X_{V'} \to X_V$ 对 $Y = \mathrm{Spec}(V)$ 和 $Y' = \mathrm{Spec}(V')$ 的闭点上方的纤维诱导同构，于是由应用于 $X_V$ 和 $X_{V'}$ 的 \ref{thm:X.2.1} 得出基本群诱导同态 $\pi_1(X_{V'}) \to \pi_1(X_V)$ 是同构，或也即：$X_{V'}$ 的每个主覆盖是同构意义下确定的 $X_V$ 的主覆盖的逆像. 这蕴含

\begin{lemma}\label{lem:X.3.7}
设 $Z_{K'}$ 是 $X_{K'}$ 的群 $G$ 的连通主覆盖，$Z_{K_s}$ 是它在 $X_{K_s}$ 上的逆像. 为使后者同构于 $X$ 的主覆盖 $Z$ 的逆像，必要且充分条件是存在 $K$ 在 $K_s$ 中的有限扩张 $K'' \supset K'$，使 $X_{K''}$ 的主覆盖 $Z_{K''}$ 由 $X_{V''}$ 的主覆盖诱导.
\end{lemma}

特别假设诸 $X_{V''}$ 正规 (例如只需 $X_0$ 正规，更不用说 $X_0$ 简单，参见 I 9.1). 由于它们连通，故不可约. 设 $L$ 是 $X$ 和 $X_K$ 的有理函数域，$L'$ 是 $X_{V'}$ 和 $X_{K'}$ 的，$L''$ 是 $X_{V''}$ 和 $X_{K''}$ 的. 则在 \ref{lem:X.3.7} 的条件下，$Z_{K'}$ 定义 $L'$ 的有限可分扩张 $R'$，$Z_{K''}$ 定义扩张 $R'' = R' \otimes_{L'} L'' = R' \otimes_{K'} K''$. \ref{lem:X.3.7} 中考虑的条件于是也意味着存在 $K'$ 的有限可分扩张 $K''$，使 $R'' = R' \otimes_{K'} K''$ 在函数域 $L'' = L' \otimes_{K'} K''$ 的正规概形 $X_{V''}$ 上非分歧，而不只是在 $X_{V''}$ 的开部分 $X_{K''}$ 上非分歧.

此后假设 $f \colon X \to Y$ 是光滑态射，故诸态射 $X_{V'} \to \mathrm{Spec}(V')$ 光滑，故诸概形 $X_{V'}$ 正则. 注意 $\mathrm{Spec}(V')$ 的闭点在 $X_{V'}$ 中的纤维不可约且余维数 1. 设 $o'$ 是它的局部环，它是函数域 $L'$ 的离散赋值环，剩余域同构于 $X_0$ 的有理函数域，故与 $k$ 有同样特征. 类似地在 $L''$ 中定义 $o''$，它显然是 $o'$ 在 $L''$ 中的正规化. 于是由纯度\cref{thm:X.3.1} 或 \ref{cor:X.3.3} 得出：为使 $R''$ 在 $X_{V''}$ 上非分歧，必要且充分条件是 $R''$ 在 $o''$ ($o'$ 在 $L''$ 中的正规化) 上非分歧.

现在注意，若 $u'$ 是 $V'$ 的幺模化元，它也是 $o'$ 的幺模化元. 若 $n$ 是与 $k$ 的特征 $p$ 互素的整数，取 $K'' = K'[t]/(t^n - u)$，则 $K''$ 是 $K'$ 的有限 伽罗瓦扩张，且 $L''$ 同构于 $L'[t]/(t^n - u')$，故在 $o'$ 上温和分歧，惯性群阶 $n$. 假设 $G$ 的阶与 $p$ 互素，这蕴含 $R'$ 在 $o'$ 上温和分歧，取 $n$ 为 $R'$ 在 $o'$ 上的惯性群阶的与 $p$ 互素的倍数 (例如 $n = [G : e]$). 应用 Abhyankar \cref{lem:X.3.6}，看到 \ref{lem:X.3.7} 中考虑的条件满足.

这证明下列定理：

\begin{theorem}\label{thm:X.3.8}
设 $f \colon X \to Y$ 是纤维几何连通的紧合光滑态射，$Y$ 局部诺特，$y_0$ 和 $y_1$ 是满足 $y_0 \in \overline{\{y_1\}}$ 的 $Y$ 的两个点，$\overline{X}_0$ 和 $\overline{X}_1$ 是对应的几何纤维，考虑特殊化同态 (\cref{cor:X.2.4}) $\pi_1(\overline{X}_1) \to \pi_1(\overline{X}_0)$. 该同态满射，且 $\pi_1(\overline{X}_1)$ 到阶与 $k(y_0)$ 的特征 $p$ 互素的有限群 $G$ 的每个连续同态都来自 $\pi_1(\overline{X}_0)$ 到 $G$ 的同态.
\end{theorem}

换言之：

\begin{corollary}\label{cor:X.3.9}
若 $k(y_0)$ 特征零，则特殊化同态是同构. 若 $k(y_0)$ 特征 $p > 0$，则特殊化同态的核包含在 $\pi_1(\overline{X}_1)$ 到阶与 $p$ 互素的有限群的连续同态的核的交中 (或即：Galois 型群 $\pi_1(\overline{X}_1)$ 的 Sylow $p$-子群生成的闭不变子群)；故若 $\pi_1(\overline{X}_1)^{(p)}$ 表示 $\pi_1(\overline{X}_1)$ 除以前述闭子群的商群，且类似地定义 $\pi_1(\overline{X}_0)^{(p)}$，则特殊化同态诱导同构
\[
\pi_1(\overline{X}_1)^{(p)} \xrightarrow{\sim} \pi_1(\overline{X}_0)^{(p)}.
\]
\end{corollary}

注意 \ref{thm:X.3.8} 的证明纯代数. 如 \ref{thm:X.2.6} 那样进行，用超越方式由此得出：

\begin{corollary}\label{cor:X.3.10}
设 $X_0$ 是特征 $p$ 的代数闭域上亏格 $g$ 的紧合光滑连通曲线. 用 \ref{cor:X.3.9} 引入的记号，群 $\pi_1(X_0)^{(p)}$ 同构于 $\Gamma^{(p)}$，其中 $\Gamma$ 是由生成元 $s_i$，$t_i$ ($1 \leqslant i \leqslant g$) 服从关系
\[
(s_1 t_1 s_1^{-1} t_1^{-1}) \cdots (s_g t_g s_g^{-1} t_g^{-1}) = 1
\]
生成的 Galois 型群.
\end{corollary}

\begin{remark}\label{rem:X.3.11}
在 $k(y_0)$ 特征零的情形，结果 \ref{cor:X.3.9} 用超越方式众所周知. 注意 \ref{cor:X.3.10} 的证明在不等特征情形调用纯度定理，但只在维数 2 的环的情形，其中该定理的证明容易且在正文中已回忆.
\end{remark}

\section{文献}\printbibliography[heading=none]


\end{refsection}
\begin{refsection}% =====================================================================
% 第十一讲 (Exposé XI) 例与补充
% (Exemples et compléments)
% A. Grothendieck
% 原书页码: pp. 219-238
% =====================================================================
\chapter{例与补充}

\section{射影空间，单有理簇}

\begin{proposition}\label{prop:XI.1.1}
设 $k$ 是代数闭域，$X = \mathbf{P}^r_k$ 是 $k$ 上维数 $r$ 的射影空间. 则 $X$ 单连通，i.e. $\pi_1(X) = 0$.
\end{proposition}

对 $r = 0$ 平凡. 若 $r = 1$，需证明：若 $X'$ 是 $X = \mathbf{P}^1_k$ 的非空连通平展覆盖，则 $X' \xrightarrow{\sim} X$. 亏格公式这里给出，若 $g$ 和 $g'$ 是 $X$ 和 $X'$ 的亏格：
\[
1 - g' = d(1 - g),
\]
其中 $d$ 是 $X'$ 在 $X$ 上的次数. 由于 $g = 0$，有 $1 - g' = d$，这要求 $d = 1$，因为 $g' \geqslant 0$，故证明 $X' \xrightarrow{\sim} X$. 当 $r \geqslant 2$ 时对 $r$ 作递归，假设对 $r' < r$ 的 $\mathbf{P}^{r'}$ 单连通. 把它应用于 $\mathbf{P}^r$ 的超平面，用 (\cref{lem:X.2.10})，确实得出 $\mathbf{P}^r$ 单连通. 另一种证明：由 (\cref{cor:X.1.7}) 有 $\pi_1(\mathbf{P}^1 \times \cdots \times \mathbf{P}^1) = \pi_1(\mathbf{P}^1) \times \cdots \times \pi_1(\mathbf{P}^1)$，故 $(\mathbf{P}^1)^r$ 单连通，因为 $\mathbf{P}^1$ 如此，故由基本群的双有理不变性 (\cref{cor:X.3.4})，$\mathbf{P}^r$ 单连通. 这个证明更一般地说明：

\begin{corollary}\label{cor:XI.1.2}
设 $X$ 是代数闭域 $k$ 上紧合的正规概形；若 $X$ 是有理簇，i.e. 整且它的函数域是 $k$ 的纯超越扩张，则 $X$ 单连通.
\end{corollary}

这个结果特别适用于 Grassmann 簇，更一般地适用于 $G/H$，其中 $G$ 是 $k$ 上的连通线性群，$H$ 是含 $G$ 的 Borel 子群 $G$ 的代数子群.

回忆 $k$ 上的\emph{单有理簇}是 $k$ 上紧合且整、函数域 $K$ 包含在 $k$ 的纯超越扩张 $K'$ (在 $K$ 上有限，i.e. 在 $k$ 上有与 $K$ 同样的超越次数) 中的概形，i.e. 存在支配有理映射 $f \colon \mathbf{P}^r_k \to X$，其中 $r = \dim X$. 若 $X$ 正规，由 X 3.4 之前的思考看到，对 $X$ 的每个连通平展覆盖 $X'$ (函数域 $L/K$)，代数 $L \otimes_K K'$ 在 $K'$ 上的模型 $\mathbf{P}^r$ 上非分歧，故由 \ref{prop:XI.1.1} 完全分解，这说明 $L$ 是 $K$-同构于 $K'/K$ 的子扩张. 因此，考虑到 V 8.2，证明：

\begin{corollary}\label{cor:XI.1.3}
代数闭域上正规单有理簇的基本群有限\index[term]{danyoulicu@单有理簇}.
\end{corollary}

(N.B. 注意在"单有理"的定义中，不需要 $K'/K$ 有限).

\begin{remark}\label{rem:XI.1.4}
当然，本节的众所周知. 另一方面，J.-P. Serre 已证明 \cite{exp11-10}：特征零代数闭域上射影光滑单有理簇 $X$ 时，$X$ 单连通. 他的证明是超越的，因为它用 Hodge 对称定理.

\begin{remark}[2003 年 (MR) 增补]
限于特征 $p \geqslant 0$ 的代数闭域 $k$. 在特征 $p > 0$ 时，前面给出的 $k$-簇单有理的定义是弱单有理 $k$-簇的定义，相对的是强单有理 $k$-簇的定义，后者还假设 $K'$ 是 $K$ 的可分扩张.

在维数 2 的情形，存在射影光滑、弱单有理但基本群 (由 \ref{cor:XI.1.3} 有限) 非平凡故非有理的曲面 (T. Shioda，On unirationality of supersingular surfaces，Math. Ann. 225 (1977)，p. 155--159). 相反，强单有理曲面有理：这由 Castelnuovo 有理性判据得出，由 O. Zariski 推广到特征 $p > 0$ (参见 J.-P. Serre，S\'em. Bourbaki no 146，1957 与 \OE uvres--Collected papers，vol. 1，p. 467).

在复数域 $\mathbb{C}$ 上，已知有维数 $\geqslant 3$ 的射影光滑、单有理但非有理簇的例子 (参见 P. Deligne，Vari\'et\'es unirationnelles non rationnelles [d'apr\`es M. Artin et D. Mumford]，S\'em. Bourbaki no 402，1971-72，Lecture Notes vol. 317). 特别在射影空间 $\mathbf{P}^4$ 中的光滑三次超曲面就是这样的情形. 这些例子中有些推广到特征 $> 0$，给出强单有理但非有理的簇 (参见 J.P. Murre，Reduction of the proof of the non-rationality of a non singular cubic threefold to a result of Mumford，Compositio Math. 27 (1973)，p. 63--82).

设 $V$ 是射影正规整 $k$-簇. 称 $V$ 有理连通，若存在整的有限型 $k$-概形 $T$ 和 $k$-态射 $F \colon T \times \mathbf{P}^1 \to V$，使态射：
\begin{equation}
\begin{array}{ccc}
T \times \mathbf{P}^1 \times \mathbf{P}^1 & \longrightarrow & V \times V \\
(t, u, u') & \longmapsto & (F(t, u), F(t, u'))
\end{array}
\tag{1}
\label{eq:XI.1}
\end{equation}
支配. 特别，经过 $V$ 的足够一般的两个有理点有一条有理曲线. 这个术语的合理性在于：若 $V$ 有理连通，可以用有限条有理曲线的链连接两个有理点. 若 $k$ 的特征 $p > 0$，加强前面的定义，要求映射 (1) 一般光滑. 于是得到可分有理连通簇的概念. 这样单有理簇有理连通，且在特征 $p > 0$ 时强单有理簇可分有理连通. J. Koll\'ar 已证明可分有理连通簇有平凡代数基本群 (参见 O. Debarre，Vari\'et\'es rationnellement connexes [d'apr\`es T. Graber，J. Harris，J. Starr et A.J. de Jong]，S\'em. Bourbaki no 906，2001-2002).
\end{remark}
\end{remark}

\section{阿贝尔簇\index[term]{abeiercu@阿贝尔簇}}

设 $k$ 是代数闭域，$A$ 是 $k$ 上的\emph{阿贝尔簇}，i.e. $k$ 上紧合、光滑且连通的群概形，最后 $G$ 是 $k$ 上有限型交换群概形. 用 $\mathrm{Ext}(A, G)$ 记 $A$ 被 $G$ 的交换扩张的类群，用 $H^1(A, G)$ 记 $A$ 上群 $G$ 的主纤维丛的类群 (比较下面第 4 号)，考虑典范同态
\[
\mathrm{Ext}(A, G) \longrightarrow H^1(A, G).
\]
Serre \cite[第七章，定理 5]{exp11-5} 的推理说明它是单射同态，像为 $H^1(A, G)$ 的"本原元素"集，i.e. 满足：
\[
\pi^{*}(\xi) = \mathrm{pr}_1^{*}(\xi) + \mathrm{pr}_2^{*}(\xi)
\]
的元素 $\xi$ 的集，其中 $\mathrm{pr}_i$ 是 $A \times A$ 到 $A$ 的两个投影，$\pi \colon A \times A \to A$ 是 $A$ 的合成律 (N.B. Serre 只在 $G$ 线性且连通、当然在 $k$ 上光滑的情形陈述他的定理，但简化他的推理的第一部分，看到这些限制不必要：只需注意 $A$ 到 $k$ 上有限型群概形 $E$ 的、把单位变成单位的每个态射是群同态，并把它应用于 $A$ 被 $G$ 的扩张 $E$ 在 $A$ 上的截面).

把这个结果应用于 $G$ 是 $k$ 上有限可分群，i.e. 普通有限群，且假设交换的情形. 用 $\pi_1(A \times A) \xrightarrow{\sim} \pi_1(A) \times \pi_1(A)$ (\cref{cor:X.1.7})，并把每个代数概形 $X$ 的 $H^1(X, G)$ 解释为 $\mathrm{Hom}(\pi_1(X), G)$，特别对 $X = A$ 或 $X = A \times A$，看到 $H^1(A, G)$ 的每个类本原，故有同构
\[
\mathrm{Ext}(A, G) \xrightarrow{\sim} H^1(A, G),
\]
换言之，$A$ 的交换结构群 $G$ 的每个在 $A$ 的原点上方带点的主覆盖，唯一地配备把标记点作为原点、且使 $A' \to A$ 成为代数群同态的代数群结构. 特别，若 $A'$ 连通，它也是阿贝尔簇，且与 $A$ 同源.

另一方面，由于函子 $X \mapsto \pi_1(X)$ (从带点代数概形 $X$ 到群) 与乘积交换 (\cref{prop:IX.1.7})，它把第一种范畴中的群变换成群范畴中的群，i.e. 交换群. 故若 $A$ 是阿贝尔簇，$\pi_1(A)$ 是交换群. 故为知道 $\pi_1(A)$，只需知道 $G$ 在有限交换群中变动时的函子 $G \mapsto H^1(A, G) = \mathrm{Hom}(\pi_1(A), G)$. 最后回忆对每个整数 $n > 0$，$A$ 中乘以 $n$ 的同态：
\[
A \xrightarrow{n} A
\]
满射，故核有限，i.e. 是同源，且由此得出每个同源 $A' \to A$ 是前述型同源的商. 由这些及标准推理 (参见例如 \cite{exp11-6}) 得到：

\begin{theorem}\label{thm:XI.2.1}
(Serre-Lang) 设 $A$ 是代数闭域 $k$ 上的阿贝尔簇，对每个整数 $n > 0$ 考虑 $A$ 中乘以 $n$ 的核 ${}_n A$ 的底集普通有限群 $K_n$，最后对每个素数 $\ell$ 令：
\[
T_\ell(A) = \varprojlim_r K_{\ell^r}
\]
且
\[
T_\bullet(A) = \prod_\ell T_\ell(A) = \varprojlim_n K_n
\]
(其中 $m$ 是 $n$ 的倍数 $m = ns$ 时，用乘以 $s$ 把 $K_m$ 送到 $K_n$). 则群 $\pi_1(A)$ 典范同构于 $T_\bullet(A)$，故对每个素数 $\ell$，$\pi_1(A)$ 的 $\ell$-准素分量典范同构于 $T_\ell(A)$.
\end{theorem}

注意这些同构对变动的 $A$ 函子. 模 $T_\ell(A)$ 称为阿贝尔簇 $A$ 的 Tate $\ell$-进模. 它是 $A$ 的加性函子，特别给出 $A$ 的自同态环 $\mathrm{Hom}(A, A)$ 在 $T_\ell(A)$ 上的表示，称为 Weil 的 $\ell$-进表示，它在阿贝尔簇理论中起重要作用 (参见例如 \cite[第七章]{exp11-4}). \cref{thm:XI.2.1} 给出它在 $A$ 的 $\ell$-进同调群 $H_1(A, \mathbb{Z}_\ell) = \pi_1(A)_\ell$ 中的自然表示的解释，后者显然先验地更令人满意，特别从 Lefschetz 公式 \cite[第五章]{exp11-4} 的观点看. 这里回忆 Weil 关于 $T_\ell(A)$ 结构的结果：

a) 若 $n$ 与 $\mathrm{car}(k)$ 互素，则 $K_n$ 是 $\mathbb{Z}/n\mathbb{Z}$ 上秩等于 $2 \dim A$ 的自由模，故若 $\ell$ 是与 $\mathrm{car}(k)$ 不同的素数，$T_\ell(A)$ 是 $\ell$-进整数环 $\mathbb{Z}_\ell$ 上秩等于 $2 \dim A$ 的自由模；

b) 若 $n$ 是 $\mathrm{car}(k) = p$ 的幂，则 $K_n$ 是 $\mathbb{Z}/n\mathbb{Z}$ 上秩 $\nu \leqslant \dim A$ (与 $n$ 无关) 的自由模，故 $T_p(A)$ 是 $p$-进整数环 $\mathbb{Z}_p$ 上秩 $\nu \leqslant \dim A$ 的自由模.

这表明在这里发展的基本群理论中，变动的阿贝尔簇的基本群不随参数正则地变化，它的 $p$-准素分量在对应剩余特征 $p$ 的参数值 $t$ 处可能骤然减小；这个现象最熟知的例子是椭圆曲线. 然而注意，无论 $n$ (是否与特征互素)，$A$ 中乘以 $n$ 的真正核 ${}_n A$ 是 $k$ 上次数 $n^{2g}$ (其中 $g = \dim A$) 的有限群概形，若 $n$ 是 $p = \mathrm{car}(k)$ 的倍数则它在 $k$ 上不可分. 而且，当 $A$ 在阿贝尔簇族中变动，i.e. 在底概形 $S$ 上有阿贝尔概形 $A$ 时，更一般地证明 ${}_n A$ 是 $S$ 上次数 $n^{2g}$ 的有限平坦群概形，即只要考虑核 ${}_n A$ 的无穷小部分，它们对任何 $n$ 都正则地表现. 这提示阿贝尔簇 $A$ 的"真正"基本群是拟代数群 (有限群在 $k$ 上的形式投射极限) $\varprojlim_n {}_n A$，其中代数概形 $X$ 的"真正基本群"应理解为：分类 $X$ 的结构群为 $k$ 上任意有限群 $G$ (不必在 $k$ 上可分) 的主覆盖的拟群. 这样，例如通过 $\mathrm{Hom}(A, A)$ 在 $A$ 的真正基本群的 $p$-准素分量中的表示，比 Serre \cite{exp11-8} 的构造更自然地找回后者用 $\ell \neq p$ 定义的 Weil 特征多项式.

\section{投射锥，Zariski 例子}

为简单仍设 $k$ 是代数闭域，设 $V$ 是 $k$ 上连通的射影 $k$-概形，是 $\mathbf{P}^r_k$ 的闭子概形，若愿意可假设它非奇异. 设 $Y = \widehat{C}_V$ 是 $V$ 的射影\emph{投射锥}\index[term]{toushezhui@投射锥}，$y_0$ 是它的顶点，$X = \widehat{C}_V$ 是与 $\mathcal{O}_V(1)$ 关联的向量丛 $\mathcal{C}_V = \mathbb{V}(\mathcal{O}_V(1))$ 的通常射影闭包，最后
\[
f \colon X \longrightarrow Y
\]
是典范态射，把 $\mathcal{C}_V$ 在 $X$ 上的零截面 $X_0$ 收缩成一点 (EGA II 8.6.4). 由于 $X$ 是 $V$ 上的局部平凡丛，纤维为 $\mathbf{P}^1$ 故单连通，态射 $p \colon X \to V$ 由 (\cref{cor:X.1.4}) 诱导同构：
\[
\pi_1(X) \xrightarrow{\sim} \pi_1(V).
\]
由于 $p$ 诱导同构 $X_0 \to V$，得出 $X$ 的平展覆盖完全分解，当且仅当它在 $X_0$ 上的限制如此. 而 $Y$ 的每个平展覆盖 $Y'$ 的逆像 $X' = X \times_Y Y'$ 是 $X$ 的在纤维 $X_0$ 上完全分解的平展覆盖，故平凡. 由于同态 $\pi_1(X) \to \pi_1(Y)$ 满射 (\cref{cor:IX.3.4})，得出
\[
\pi_1(Y) = (e),
\]
换言之每个射影投射锥单连通. (N.B. 在特征 0 时，取 $Y$ 为仿射投射锥时同样结果成立).

现在假设 $V$ 正则 i.e. 在 $k$ 上光滑；则 $X$ 正则，对 $V$ 的适当射影浸入，找到正规的投射锥 $Y$. 若 $V$ 不单连通，故 $X$ 不单连通，设 $X'$ 是 $X$ 的非平凡连通平展覆盖. 由于 $X$ 在 $y \in Y$ ($y \neq y_0$) 处的纤维约化为一点，看到 $X'$ 在它的诸纤维 (特别在一般纤维) 上的限制平凡；然而 $X$ 不来自 $Y$ 的平展覆盖的逆像，因为 $Y$ 单连通且 $X'$ 将完全分解. 这说明 (X 1.3 和 \ref{rem:XI.1.4}) 在把 $f$ 可分的假设换成更弱的、它的诸纤维是 $k(s)$ 上的可分 (甚至光滑) 代数概形的假设时变假. 同样注意几何纤维 $\overline{X}_y$ ($y \neq y_0$) 的基本群显然约化为 $(e)$，因为这些纤维约化为一点，而 $\pi_1(X_0) \neq e$，故半连续性定理 (\cref{cor:X.2.4}) 对 $f$ 也失效.

最后指出 Zariski 指出的例子，它使这些定理在把 $f$ 可分的假设换成 $f$ 平坦的假设时失效. 设 $f \colon X \to Y$ 是从非奇异射影曲面到有理直线 $Y = \mathbf{P}^1$ 的态射，使 $K = k(X)$ 是 $k(f)$ 的"正则"扩张，i.e. 初等且可分，(i.e. 一般几何纤维连通且可分)，且使除子 $(f) = X_0 - X_\infty$ 是某个除子的 $n$ 倍 (其中 $n$ 是与特征互素的整数). 在任意特征中可以构造这样的例子. 设 $X'$ 是 $X$ 在 $K(f^{1/n})$ 中的正规化，其中 $K = k(X)$ 是 $X$ 的函数域. 由对 $(f)$ 的假设得出 $X'$ 在 $X$ 上平展. 设 $Y'$ 是 $Y$ 在 $k(t)(t^{1/n})$ 中的正规化，它恰在点 $t = 0$ 和 $t = \infty$ 处分歧于 $Y$ 上，且限制 $X'|_{f^{-1}(U)}$ 同构于 $Y'|_U$ 的逆像. 特别，$X'$ 在 $X$ 于 $Y$ 上的一般几何纤维上的限制完全分解. 然而，$X'$ 不同构于 $Y$ 的平展覆盖的逆像，因为立即看到后者必然为 $Y'$，这与 $Y'$ 在 $Y$ 上分歧\footnote{可以注意，从"平展拓扑"(SGA 4 VII) 的观点，在这个例子中 $R^1(f_{\mathrm{ét}})_*(\mathbb{Z}/n\mathbb{Z})$ 在 $S$ 上"不分离".} 矛盾.

下面 (按 Serre) 是实现这个例子条件的简单方式，借鉴 \cite[第 20 号]{exp11-5}：取 $n$ 为 $\geqslant 5$ 的素数，与特征不同，让 $G = \mathbb{Z}/n\mathbb{Z}$ 通过把坐标乘以 $G$ 的四个不同特征 (由于 $n \geqslant 5$ 可能) 作用在 $k^4$ 中\footnote{[2003 年增补：$k^4$ 表示 $k$ 上维数 4 的仿射空间.]}. 则 $G$ 作用在射影空间 $\mathbf{P}^3_k$ 上，$G$ 下唯一的不动点是坐标轴对应的四个点. 曲面 $X'$ 由方程 $x^n + y^n + z^n + t^n = 0$ 定义，在 $k$ 上光滑 (Jacobi 判则)，不含任何不动点，故由于 $G$ 是素阶群，"无不动点"地作用在 $X$ 上，i.e. $X$ 是 $X = X'/G$ 的群 $G$ 的主覆盖. 设 $g = x/y$ 是 $k(X') = K'$ 中的元素，若所选特征是 $\chi_i$，$i = 0, 1, 2, 3$ (其中 $\chi$ 是本原特征)，则它是 $K'$ 在 $K = k(X)$ 上的 Kummer 生成元，设 $f$ 是它的 $n$ 次幂，是 $K$ 的元素. 立即看到 $K'$ 是 $k(g)$ 的正则扩张，这由曲线 (在 $U, T, Z$ 中的齐次方程 $T^n + Z^n + (1 + g^n)U^n = 0$) 在 $k(g)$ 上光滑 (Jacobi 判则) 及已知每条平面曲线连通得出. 另一方面有 $k(f) = K \cap k(g)$，因为第二个成员是 $k(f)$ 的包含在 $k(g)$ (素阶扩张) 中且不同于 $k(g)$ (因为 $g \notin K$) 的扩张. 这蕴含 $K$ 是 $k(f)$ 的正则扩张. 最后 $f$ 在 $X$ 上的除子是某个除子的 $n$ 倍，因为它在 $X'$ 上的逆像是 $g^n$ 的除子，故是 $n$ 倍，且可以下降，因为 $X'$ 在 $X$ 上平展. 若有理映射 $f \colon X \to \mathbf{P}^1$ 是态射，即 $f$ 的零点除子与极点除子不相交，就完成了. 实际上，容易验证 (仍看 $X'$ 上) 这两个除子是 $k$ 上两条在一点 $a$ 横截相交的光滑曲线的 $n$ 倍. 现在把 $X$ 换成在 $a$ 处爆破得到的概形，前面条件 ($\mathrm{div}(f)$ 可被 $n$ 整除，且 $k(X_1) = k(X)$ 是 $k(X)$ 的正则扩张) 仍满足，而且 $f$ 是态射 $X_1 \to \mathbf{P}^1$，故满足所要条件.

\section{上同调正合序列}

设 $S$ 是预概形，故 $X$ 上预概形范畴 $\mathbf{Sch}_{/S}$ 确定，因而其中的群概念也确定，称为 $S$ 上的预概形群，或简称 $S$-群. 为简化叙述并明确思想，此后我们最常限于在 $S$ 上仿射且平坦\footnote{事实上，对下面要做的，用拟仿射假设代替仿射假设就够了，参见第 228 页脚注 (5).} 的群\footnote{参见 SGA 3 IV，特别 §4.}，这足以用于我们考虑的应用. (当然，遇到许多两个假设都不满足的情形). 设 $G$ 是这样的 $S$-群，$P$ 是 $G$ 在其上右作用的 $S$-预概形，这特别蕴含 $\mathbf{Sch}_{/S}$ 中的态射
\[
\pi \colon P \times_S G \longrightarrow P
\]
满足众所周知的公理. 称 $P$ 在 $G$ 下形式主齐性\index[term]{xingshizhuqixing@形式主齐性}，若态射
\[
P \times_S G \longrightarrow P \times_S P
\]
(分量为 $\mathrm{pr}_1$ 和 $\pi$) 是同构；这等价于说对 $\mathbf{Sch}_{/S}$ 的每个对象 $S'$，$P(S') = \mathrm{Hom}_S(S', P)$ 作为带算子群 $G(S') = \mathrm{Hom}_S(S', G)$ 的集，为空或主齐性 (i.e. 为空或同构于 $G(S')$，其上 $G(S')$ 经右平移作用). 称 $P$ 平凡，若 $P$ 同构于 $G$ (其上 $G$ 经右平移作用)，或等价地，若带算子集 $P(S')$ ($G(S')$ 下) 每个都平凡. 例如用著名的归结到集合情形的技巧验证，$P$ 平凡当且仅当它形式主齐性且在 $S$ 上有截面 (后者用范畴术语说 $P$ 在 $\mathbf{Sch}_{/S}$ 的终结对象 $e = S$ 上有截面，i.e. 存在 $e$ 到 $P$ 的态射). 定义比形式\emph{主齐性纤维丛}\index[term]{zhuqixingxianweicong@主齐性纤维丛}更强的 $G$ 下的主齐性纤维丛 $P$ 概念，需先在 $\mathbf{Sch}_{/S}$ 中明确一个用于"下降"的态射集，它起"局部化态射"的作用来"平凡化"纤维丛. 最适当的选择随上下文而变，没有任何一个包含其他所有. 这里方便采用下列定义：

\begin{definition}\label{def:XI.4.1}
设 $G$ 是 $S$-群. 称 $P$ 为 $G$ 下的 (右) 主齐性纤维丛，若 $P$ 是带右 $S$-群 $G$ 的 $S$-预概形，使存在 $S$ 的开集 $U_i$ 的覆盖，以及对每个 $i$ 忠实平坦且拟紧的基变换态射 $S'_i \to U_i$，使 $P' = P \times_S S'$ 是 $G' = G \times_S S'$ 下的平凡带算子预概形，其中 $S'$ 是诸 $S'_i$ 的无交和 $S$-预概形.
\end{definition}

(注意基变换函子 $X \mapsto X' = X \times_S S'$ 左正合，把群变成群，带算子群的对象变成带算子群的对象). 注意 \ref{def:XI.4.1} 随基变换稳定. 还注意：

\begin{proposition}\label{prop:XI.4.2}
设 $G$ 是在 $S$ 上平坦且拟紧的 $S$-群，$P$ 是 $G$ 在其上右作用的 $S$-预概形. 下列条件等价：
\begin{enumerate}[i)]
\item $P$ 是 $G$ 下的主齐性纤维丛；
\item $P$ 在 $G$ 下形式主齐性，且结构态射 $P \to S$ 忠实平坦且拟紧.
\end{enumerate}
\end{proposition}

若 $P$ 在 $G$ 下主齐性，则用 \ref{def:XI.4.1} 的记号，$P'$ 在 $S'$ 上忠实平坦且拟紧 (因为 $G'$ 如此，且 $P'$ 与它 $S'$-同构)，故 $P$ 在 $S$ 上方有同样性质 (对"满射"和"拟紧"，参见 VIII 3.1，对"平坦"是第 VIII 讲的推论中的疏漏). 反之，若 (ii) 满足，取基变换 $S' = P$，它在 $S$ 上确实忠实平坦且拟紧；则 $P'$ 在 $S'$ 上形式主齐性，因为 $P$ 在 $S$ 上如此且基变换函子左正合，另一方面 $P'$ 在 $S'$ 上有截面，即对角截面，故它是平凡主纤维丛，这完成证明.

\begin{corollary}\label{cor:XI.4.3}
若 $G$ 在 $S$ 上仿射且平坦，则 $G$ 下的每个主齐性纤维丛 $P$ 在 $S$ 上仿射且平坦.
\end{corollary}

确实，经基的忠实平坦拟紧扩张后它如此，应用 (\cref{cor:VIII.5.6}).

\cref{def:XI.4.1} 对 $S$ 上平坦仿射 $S$-群有用，因为 (\cref{thm:VIII.2.1})，i.e. \ref{def:XI.4.1} 中考虑的态射 $S' \to S$ 对互相仿射预概形的范畴是有效下降态射. 借助这个事实，下面概述的事实的验证基本上"范畴式"进行\footnote{参见第 226 页脚注 (3) 中的 loc. cit..}. 设 $E$ 是 $S$-群 $G$ 在其上左作用的 $S$-预概形，设 $P$ 是 $G$ 下的 (右) 主齐性纤维丛，想定义"局部"同构于 $E$ 的相伴纤维丛 $E^{(P)}$. 为此，让 $G$ 按律 $(x, y) \mapsto (xg, g^{-1}y)$ 右作用在 $P \times_S E$ 上，该律在集合情形描述此类作用，并经著名技巧推广到范畴. 令 (假定存在)：
\[
E^{(P)} = (P \times_S E)/G
\]
由此看到 $P \times_S E$ 是 $T = E^{(P)}$ 上带右算子群 $G_T = G \times_S T$ 的预概形；为方便，还希望 $P \times_S E$ 是 $T$ 上群 $G_T$ 的主齐性纤维丛. 为验证 $E^{(P)}$ 的存在和前述性质，取\cref{def:XI.4.1} 中的 $S'$，看初始情形在 $S'$ 上的逆像情形. 由于 $P'$ 平凡，i.e. 同构于 $G'_d$，立即看到 $E'^{(P')}$ 存在且有所要的正合性质. 事实上，$E' \times_{S'} P'$ 是 $G'$-同构于 $E' \times_{S'} G'$ 的，故 $E'^{(P')}$ 同构于 $E'$. 而且，带算子空间情形的"相伴纤维丛"的构造与任何基扩张交换，并取种种基扩张 $S'' \rightrightarrows S'$ 和 $S''' \rightrightarrows S''$，其中 $S''$ 和 $S'''$ 是 $S'$ 在 $S$ 上的二重和三重纤维积，看到 $E'^{(P')}$ 配备相对态射 $S' \to S$ 的下降数据，且 $E^{(P)}$ 以所需性质存在当且仅当该下降数据有效；当然 $E^{(P)}$ 正是下降对象. (用 $S' \to S$ 是 $S$-预概形范畴中的下降态射的事实，参见 VIII 5.2). 由此得出 $E$ 在 $S$ 上仿射时相伴纤维丛存在. 把这个构造应用于有 $S$-群同态 $G \to H$，取 $E$ 为 $S$-预概形 $H$ (配备由给定同态得到 $G$ 在 $H$ 上的左作用) 的情形；由于 $H$ 以与 $G$ 在 $H$ 上的作用交换的方式右作用于自身，且 (在 $S$ 上方存在的前提下) 相伴纤维丛的构造与基扩张交换，容易看到 $H$ 将右作用在 $P^{(H)}$ 上，后者按 \ref{def:XI.4.1} 的意义成为 $H$ 下的主齐性纤维丛，且精确地说被与 $P$ 同样的态射 $S' \to S$ 平凡化. 特别，对 $G$ 下的每个主齐性纤维丛 $P$ 和每个 $S$-群同态 $G \to H$ ($H$ 在 $S$ 上仿射)，以对 $(G \to H)$ 函子且与任何基变换 $T \to S$ 相容的方式，关联 $H$ 群的主齐性纤维丛.

\begin{definition}\label{def:XI.4.4}
设 $G$ 是 $S$-预概形. 用 $H^0(S, G)$ 记 $G$ 在 $S$ 上的截面集，当 $G$ 是 $S$-群时把它看作群. 此时用 $H^1(S, G)$ 记 $S$ 上群 $G$ 的主齐性纤维丛的同构类集，把 $H^1(S, G)$ 看作配备对应平凡纤维丛的"标记点"\footnote{这个记号只有在有下降有效性判据时才与一般上同调记号 (SGA 4 V) 相容，后者只在 $G$ 仿射 (或只拟仿射，参见 (\cref{cor:VIII.7.9})) 时几乎不保证.}.
\end{definition}

这样，$H^0(S, G)$ 是 $S$-预概形 $G$ 的、取值在集合范畴的函子. 该函子左正合，更不用说与有限积交换，这确实蕴含它把群变成群，交换群变成交换群. 类似地，借助相伴纤维丛的构造，$H^1(S, G)$ 是仿射 $S$-群 $G$ 的、取值在集合范畴的函子；容易看到该函子与有限积交换. 特别它把第一种范畴中的群，i.e. 交换仿射 $S$-群，变成第二种范畴中的群，甚至交换群 (因为第一种范畴中的群交换). 这样，若 $G$ 是交换仿射 $S$-群，$H^1(S, G)$ 是交换群，交换仿射 $S$-群的同态 $G \to H$ 给出群同态 $H^1(S, G) \to H^1(S, H)$.

为简单，此后限于交换仿射 $S$-群. 设
\[
0 \longrightarrow G' \xrightarrow{u} G \xrightarrow{v} G'' \longrightarrow 0
\]
是这样的群的态射序列，称该序列正合，若 $vu = 0$ (这可以把 $G$ 看作 $G''$ 上带右算子群 $G'_{G''}$ 的预概形) 且 $G$ 是 $G''$ 上群 $G'_{G''} = G' \times_S G''$ 的主齐性纤维丛. 这特别蕴含 $u \colon G' \to G$ 是 $v$ 的核，更不用说蕴含序列
\[
0 \longrightarrow H^0(X, G') \longrightarrow H^0(X, G) \longrightarrow H^0(X, G'')
\]
的正合性. 还蕴含可以定义映射
\[
\partial \colon H^0(X, G'') \longrightarrow H^1(X, G'),
\]
把 $G''$ 在 $S$ 上的每个截面，i.e. 每个 $S$-态射 $f \colon S \to G''$，配给 $S$ 上群 $G' \simeq f^{*}(G'_{G''})$ 的主齐性纤维丛 $P_f$，即 $G''$ 上主齐性纤维丛 $G$ 的逆像. 从 $S$-预概形观点，它正是经 $v \colon G \to G''$ 的、浸入 $f$ 的像 $S$ 的子预概形的逆像，且 $G'$ 在 $P_f$ 上的作用由 $G'$ 在 $G$ 上的右作用诱导. 下列命题的验证也没有除写法以外的困难，留给读者：

\begin{proposition}\label{prop:XI.4.5}
映射 $\partial \colon H^0(X, G'') \to H^1(X, G')$ 是群同态. 下列同态序列正合：
\[
0 \longrightarrow H^0(X, G') \xrightarrow{u_0} H^0(X, G) \xrightarrow{v_0} H^0(X, G'') \xrightarrow{\partial} H^1(X, G') \xrightarrow{u_1} H^1(X, G) \xrightarrow{v_1} H^1(X, G'')
\]
(其中除 $\partial$ 外的同态来自 $H^0$ resp. $H^1$ 的函子律).
\end{proposition}

\begin{remark}\label{rem:XI.4.6}
这里用于研究主齐性纤维丛的观点显然受 Serre \cite{exp11-7} 启发，读者会有兴趣查阅. 当想要一个也适用于在 $S$ 上拟射影的结构 $S$-群 (从而特别包含射影阿贝尔概形) 的形式体系时，有兴趣修改 \ref{def:XI.4.1}，要求 $S'$ 是在覆盖 $S$ 的 $S$ 的开集 $S_i$ 上有限且局部自由的预概形 $S'_i$ 的和. 前面的展开此时成立，特别包括 \ref{prop:XI.4.5}，只要处处把仿射假设换成拟射影假设，并相应解释 $S$-群正合序列的定义. 确实只需把 (\cref{thm:VIII.2.1}) 的引用换成 (\cref{cor:VIII.7.7})：所用态射 $S' \to S$ 对互相拟射影预概形的纤维范畴是有效下降态射. 然而注意这个第二种主齐性纤维丛概念比第一种 \ref{def:XI.4.1} 更严格.
\end{remark}

\refstepcounter{theorem}\noindent\textbf{\thesection.\arabic{theorem}.}\quad 要求 $S$ 被开集 $S_i$ 覆盖，使对每个 $i$，$P|_{S_i}$ 是 $G|_{S_i}$ 下的平凡带算子纤维丛，得到更严格的主齐性纤维丛概念：此时称 $P$ 是\emph{局部平凡主齐性纤维丛}\index[term]{jubupingfanzhuqixingxianweicong@局部平凡主齐性纤维丛}. 这些纤维丛 (对给定的 $G$) 的类构成 $H^1(X, G)$ 的部分，与 $H^1(X, \mathcal{O}_X(G))$ 一一对应，其中 $\mathcal{O}_X(G)$ 是 $G$ 在 $S$ 上的截面的 (通常意义) 层，参见 \cite{exp11-2}. 为使这些 $H^1$ 还给出上同调正合序列 \ref{prop:XI.4.5}，显然必须假设序列 $0 \to G' \to G \to G'' \to 0$ 对这个新上下文在合理意义下正合，i.e. $G$ 是 $G''$ 上群 $G'_{G''}$ 的局部平凡纤维丛；这也意味着 $u \colon G' \to G$ 是 $v \colon G \to G''$ 的核，且 $G$ 在 $G''$ 上局部有截面.

\refstepcounter{theorem}\noindent\textbf{\thesection.\arabic{theorem}.}\quad 引入高上同调群 $H^i(X, G)$ 延续正合序列 \ref{prop:XI.4.5} 显然非常可取. 这可以在"Weil 上同调"框架中做到：考虑 $S$ 上拟紧预概形范畴 $B$，配备忠实平坦拟紧态射集 $M$，称为局部化态射. $S$ 上 (或更好，在 $(B, M)$ 上) 的 Abel "Weil 层"是把和变成积、把
\[
T'' = T' \times_T T' \xrightarrow{\mathrm{pr}_1, \mathrm{pr}_2} T' \xrightarrow{f} T \quad (f \in M)
\]
的序列变成集合的精确图式
\[
F(T) \longrightarrow F(T') \rightrightarrows F(T'')
\]
的、从 $B$ 到 Abel 群范畴的反变函子 $F$. Weil 层构成 Abel 范畴，有正合归纳极限且有一个生成元，故有足够多内射对象 \cite{exp11-1}. 函子 $\Gamma(F) = F(S)$ 的右导出函子记为 $H^i(S, F)$. 另一方面，每个交换 $S$-群显然定义 Weil 层 (\cref{thm:VIII.5.2})，它的 $H^0$ 和 $H^1$ 正是 $H^0(S, G)$ 和 $H^1(S, G)$，从而可以合理方式定义其他 $H^i(S, G)$. 还证明 $S$-群的正合序列定义 Weil 层的正合序列，从而可以找回并延续正合序列 \ref{prop:XI.4.5}\footnote{关于这个观点的系统研究，参见 SGA 4 I 至 IX.}.

\refstepcounter{theorem}\noindent\textbf{\thesection.\arabic{theorem}.}\quad 如 \cite{exp11-2} 那样发展 \ref{prop:XI.4.5} 的非交换变体是适宜的. 关于本节概述的各种上同调概念的适当框架中的系统发展，参见 J. Giraud 的预备工作\footnote{参见 J. Giraud，Cohomologie non ab\'elienne，Springer-Verlag 1971.}.

\section{主纤维丛的特殊情形}

现在假设 $S$ 连通，配备几何点 $a$，从而有分类 $S$ 的平展覆盖的基本群 $\pi_1(S, a)$：$S$ 的平展覆盖范畴等价于 $\pi_1$ 连续作用的有限集范畴. 由此得出 $S$ 上有限平展群概形 $G$ 本质上由普通有限群 $\underline{G}$ (其上 $\pi_1$ 以群自同构连续作用) 确定. $G$ 在其上右作用的 $S$ 的平展覆盖 $P$ 本质上由有限集 $\underline{P}$ (其上 $\pi_1$ 连续地 (左) 作用) 确定，$\underline{G}$ 以与 $\pi_1$ 的作用相容的方式右作用于其上：
\[
s(p \cdot g) = (sp) \cdot (sg) \quad \text{对 } s \in \pi_1, p \in \underline{P}, g \in \underline{G}.
\]
验证 $P$ 按 \ref{def:XI.4.1} 的意义是主齐性纤维丛，当且仅当 $\underline{P}$ 是 $\underline{G}$ 下的主齐性集 (例如用判则 \ref{prop:XI.4.2}). 换言之，$S$ 上群 $G$ 的主齐性纤维丛范畴等价于 $\pi_1$ 连续作用的有限集范畴中群 $\underline{G}$ 的主齐性纤维丛范畴. 特别由此得到对 $G$ 函子的典范双射：
\[
(\ast) \quad H^1(S, G) \xrightarrow{\sim} H^1(\pi_1, \underline{G}),
\]
其中第二个成员表示 $\pi_1$ 作用的 (当然无需说明：连续) 有限集范畴中 $\underline{G}$ 下的主齐性纤维丛的同构类集，该集以熟知方式明确为 1-上链集 $Z^1(\pi_1, \underline{G})$ (满足 $\phi(1) = 1$，$\phi(st) = \phi(s)(s.\phi(t))$) 的商集，除以自然作用其上的群 $\underline{G}$.

一个重要的情形是 $\pi_1$ 平凡作用于 $\underline{G}$，i.e. $G$ 是 $S$ 的完全分解覆盖，同构于 $\underline{G}$ 个 $S$ 的和；此时也写 $H^1(S, \underline{G})$ 代替 $H^1(S, G)$，且该集与 $H^1(\pi_1, \underline{G}) = \mathrm{Hom}(\pi_1, \underline{G})/\underline{G}$ 的内自同构一一对应 ($\ast$). 还注意此时 $S$ 上群 $G$ 的主齐性纤维丛正是 $S$ 的群 $\underline{G}$ 的主覆盖 (\cref{cor:V.2.7})，前述一一对应正是从 $S$ 的群 $\underline{G}$、在 $a$ 上方带点的主覆盖与 $\pi_1(S, a)$ 到 $\underline{G}$ 的连续同态之间的对应 (V 第 5 号末) 推出的对应.

把平展覆盖理论与主纤维丛理论联系起来的兴趣 (已隐含在 A. Weil，G\'en\'eralisation des Fonctions Ab\'eliennes 中，并由 S. Lang 在他的类域论几何理论中明确，参见 Serre \cite{exp11-5})，来自每个在 $S$ 上有限且平展的 $S$-群可以浸入一个在 $S$ 上仿射且光滑、纤维连通、$G$ 交换时也交换的 $S$-群 $H$，使得由正合序列 \ref{prop:XI.4.5} (及其可能的非交换变体)，群 $G$ 的主覆盖的"离散"分类可以用群 $H$ 的主纤维丛的"连续"分类研究，反之亦然. 关于 $G$ 在 $H$ 中浸入的一般构造的想法 (看来在实践中很少用)，参见 \cite[VI 2.8]{exp11-5}. 我们只在下节发展两个重要的、当然经典的特殊情形. 需要一个辅助结果：

\begin{proposition}\label{prop:XI.5.1}
设 $S$ 是预概形，$G$ 是同构于 $\mathrm{GL}(n)_S$ (例如 $G_{m,S}$) 或 $G_{a,S}$ 的 $S$-群，则 $G$ 下的每个主齐性纤维丛局部平凡.
\end{proposition}

明确 $\mathrm{GL}(n)_S$ ($n$ 是 $\geqslant 0$ 的整数) 表示 $S$-预概形 $T$ 上反变函子 $T \mapsto \mathrm{GL}(n, \Gamma(T, \mathcal{O}_T))$ 的 $S$-群，特别 $G_{m,S}$ ("$S$ 上的乘法群") 表示反变函子 $T \mapsto \Gamma(T, \mathcal{O}_T^{*})$，故作为 $S$ 上的预概形同构于 $\mathrm{Spec} \mathcal{O}_S[t, t^{-1}]$，其中 $t$ 是未定元. 类似地 $G_{a,S}$ 表示反变函子 $T \mapsto \Gamma(T, \mathcal{O}_T)$，故作为 $S$-预概形同构于 $\mathrm{Spec}(\mathcal{O}_S[t])$，其中 $t$ 是未定元. 注意经分解论证，\ref{prop:XI.5.1} 重新给出 Rosenlicht 的局部平凡性结果，它关于 $G$ 有"合成序列"、相继因子为这里考虑类型的群的情形. (关于主齐性纤维丛局部平凡性问题的更精细研究，参见 \cite{exp11-7} 和 \cite{exp11-3}).

第一个断言由 $G(T) = \mathrm{Aut}(\mathcal{O}_T^n)$ 及 \ref{def:XI.4.1} 中涉及的态射 $S' \to S$ (i.e. 忠实平坦且拟紧的) 对局部同构于 $\mathcal{O}_T^n$ 的模，i.e. 秩 $n$ 局部自由模，的纤维范畴是\emph{有效下降态射} (\cref{rem:VIII.1.12}) 的事实证明. 第二个类似证明，注意此时 $G(T) = \mathrm{Aut}(E_T)$，其中 $E_T$ 是 $\mathcal{O}_T$ 被 $\mathcal{O}_T$ 的平凡扩张 (其中自同构当然必须保持扩张结构)，且 \ref{def:XI.4.1} 中涉及的态射 $S' \to S$ 对 $\mathcal{O}_T$ 被 $\mathcal{O}_T$ 的扩张的纤维范畴是有效下降态射 (容易由 VIII 1.1 得出)，且这样的扩张自动局部平凡.

\begin{remark}\label{rem:XI.5.2}
注意同样的证明类型适用于辛群 $\mathrm{Sp}(2n)_S$，考虑到：局部同构于 $\mathcal{O}_S^{2n}$ 的模上的交错型，若"非退化" i.e. 定义该模与其对偶的同构，则局部同构于标准型. 正交群的结果相反为假，即使在 $S$ 是域的谱时，因为域上可以有不与标准型同构的二次型. 而且基本上在 \cite{exp11-3} 中证明，群 $\mathrm{GL}$，$\mathrm{Sp}$，$G_a$ 及可分解成此类群的群，是这里考虑的局部平凡性结果几乎唯一的群.
\end{remark}

\begin{corollary}\label{cor:XI.5.3}
有典范双射
\[
H^1(S, \mathrm{GL}(n)_S) \xleftarrow{\sim} H^1(S, \mathrm{GL}(n, \mathcal{O}_S)),
\]
特别
\[
H^1(S, G_{m,S}) \xleftarrow{\sim} H^1(S, \mathcal{O}_S^{*}),
\]
和
\[
H^1(S, G_{a,S}) \xleftarrow{\sim} H^1(S, \mathcal{O}_S),
\]
其中第二个成员表示拓扑空间 $S$ 以通常层为系数的上同调.
\end{corollary}

特别，$H^1(S, \mathrm{GL}(n)_S)$ 等同于 $S$ 上秩 $n$ 局部自由模的同构类集，$H^1(S, G_{a,S})$ 等同于模 $\mathcal{O}_S$ 被它自己的扩张的类集.

\section{应用于主覆盖：Kummer 与 Artin-Schreier 理论}

\begin{proposition}\label{prop:XI.6.1}
设 $S$ 是预概形，$n$ 是 $> 0$ 的整数，$u_n \colon G_{m,S} \to G_{m,S}$ 是 $n$ 次幂同态，${}_n \mu_S$\index[notations]{${}_n\mu_S$} 是它的核. 则 ${}_n \mu_S$ 在 $S$ 上有限且秩 $n$ 局部自由，且当且仅当对每个 $s \in S$，$s$ 的特征与 $n$ 互素时它在 $S$ 上平展. 同态序列
\[
0 \longrightarrow {}_n \mu_S \longrightarrow G_{m,S} \xrightarrow{u_n} G_{m,S} \longrightarrow 0
\]
按第 4 号的意义正合. (称之为 $S$ 上相对整数 $n$ 的 \emph{Kummer 正合序列}\index[term]{Kummer zhenghexulie@Kummer 正合序列}).
\end{proposition}

有
\[
G_m = \mathrm{Spec} \mathcal{O}_S[t, t^{-1}],
\]
$u_n$ 对应仿射 $\mathcal{O}_S$-代数上的同态 $u_n$，由
\[
u_n(t) = t^n
\]
给出，另一方面 $G_{m,S}$ 的单位截面对应 $\mathcal{O}_S$-代数的增广同态，由
\[
\varepsilon(t) = 1
\]
给出，它的核是主理想 $(t - 1)$. 后者经 $u_n$ 的像是主理想 $(1 - t^n)$，得到：
\[
{}_n \mu_S = \mathrm{Spec} \mathcal{O}_S[t]/(1 - t^n),
\]
这特别说明 ${}_n \mu_S$ 在 $S$ 上有限，由 $S$ 上秩 $n$ 自由、以诸 $t^i$ ($0 \leqslant i \leqslant n - 1$) 为基的代数定义. 为使它平展于 $s \in S$，必要且充分条件是约化代数 $k[t]/(1 - t^n)$ (其中 $k = k(s)$，由向 $k$ 形式添加 $n$ 次单位根得到) 在 $k$ 上可分，i.e. $1 - t^n$ 在 $k$ 的代数闭包中的根全部不同，这等价于 $n$ 与特征互素. 最后，证明 \ref{prop:XI.6.1} 中同态序列正合归结到由判则 \ref{prop:XI.4.2} 证明 $u_n$ 忠实平坦. 显然可假设 $S$ 仿射，环为 $A$，故 $G_{m,S}$ 仿射，环为 $B = A[t, t^{-1}]$，只需验证 $u_n$ 使 $B$ 成为 $B$ 上秩 $n$ 的自由模，或等价地 $u_n$ 单射，且 $A[t, t^{-1}]$ 是 $A[t^n, t^{-n}]$ 上秩 $n$ 的自由模. 确实容易验证诸 $t^i$ ($0 \leqslant i \leqslant n - 1$) 构成其一在另一个上的基，这完成证明.

\begin{definition}\label{def:XI.6.2}
称 ${}_n \mu_S$ 为 $S$ 上秩 $n$ 的 Kummer 群，称 $S$ 上秩 $n$ 的 Kummer 主覆盖为 $S$ 上秩 $n$ 的 Kummer 群下的每个主齐性纤维丛.
\end{definition}

这些覆盖的集是群，记为 $H^1(S, {}_n \mu_S)$ 或简记 $H^1(S, {}_n \mu)$. 注意 $S$ 上秩 $n$ 的 Kummer 群的构造与基扩张相容，故 ${}_n \mu_S$ 由绝对 Kummer 群 ${}_n \mu$ 经基扩张自 $\mathrm{Spec}(\mathbb{Z})$ 得到.

用 $(\mathbb{Z}/n\mathbb{Z})_S$ 记由普通有限群 $\mathbb{Z}/n\mathbb{Z}$ 定义的 $S$-群. 若 $G$ 是任意 $S$-群，$(\mathbb{Z}/n\mathbb{Z})_S$ 到 $G$ 的 $S$-群同态 $u$ 与 $G$ 在 $S$ 上的 $n$ 次幂为单位截面的截面一一对应，且与基变换相容，对应 $u$ 的 $(\mathbb{Z}/n\mathbb{Z})_S$ 在 $S$ 上由 $\mathbb{Z}/n\mathbb{Z}$ 的生成元 $1 \mod n\mathbb{Z}$ 定义的截面的像. 这样：

\begin{corollary}\label{cor:XI.6.3}
若 ${}_n \mu_S$ 在 $S$ 上平展，则得到 $S$-群同构 $(\mathbb{Z}/n\mathbb{Z})_S \xrightarrow{\sim} {}_n \mu_S$ 与 $\mathcal{O}_S$ 的在 $S$ 的每个连通分量上精确阶为 $n$ 的截面 (这样的截面称为"$S$ 上的本原 $n$ 次单位根") 之间的一一对应. 故为使 ${}_n \mu_S$ 作为 $S$-群同构于 $(\mathbb{Z}/n\mathbb{Z})_S$，必要且充分条件是它平展于 $S$ 上，i.e. $S$ 的剩余特征与 $n$ 互素，且 $S$ 上存在本原 $n$ 次单位根.
\end{corollary}

这解释了在经典 Kummer 理论中假设基域 (充当 $S$ 的角色) 的特征与 $n$ 互素且含 $n$ 次单位根、并选择本原 $n$ 次单位根所起的作用. 一旦有了概形语言，就不必再纠缠这些假设，应当直接对 $\mu_n$ 而不是 $\mathbb{Z}/n\mathbb{Z}$ 推理. 这样，\ref{prop:XI.6.1}，\ref{prop:XI.4.5} 和 \ref{cor:XI.5.3} 的合取给出 Kummer 主覆盖理论与 Picard 群理论之间的下列一般关系：

\begin{proposition}\label{prop:XI.6.4}
设 $S$ 是预概形，有典范正合序列
\[
0 \to H^0(S, \mu_n) \to H^0(S, \mathcal{O}_S^{*}) \to H^0(S, \mathcal{O}_S^{*}) \to H^1(S, \mu_n) \to H^1(S, \mathcal{O}_S^{*}) \to H^1(S, \mathcal{O}_S^{*}),
\]
从而，令 $H^1(S, \mathcal{O}_S^{*}) = \mathrm{Pic}(S)$，并用 ${}_n A$ 和 $A^n$ 对每个 Abel 群 $A$ 记 $A$ 中乘以 $n$ 的核和余核，有正合序列：
\[
0 \longrightarrow H^0(S, \mathcal{O}_S)^{*n} \longrightarrow H^1(S, \mu_n) \longrightarrow {}_n \mathrm{Pic}(S) \longrightarrow 0.
\]
\end{proposition}

将明确两个重要情形，其中这个正合序列的一个或另一个端项为零：

\begin{corollary}\label{cor:XI.6.5}
假设 ${}_n \mathrm{Pic}(S) = 0$ (例如 $S$ 是局部环的谱，或因子环的谱)，设 $A$ 是环 $H^0(S, \mathcal{O}_S)$. 则有典范同构
\[
H^1(S, \mu_n) \xrightarrow{\sim} A^{*}/A^{*n}.
\]
\end{corollary}

这本质上是 $S$ 是域的谱时的 Kummer 理论经典陈述.

\begin{corollary}\label{cor:XI.6.6}
假设 $H^0(S, \mathcal{O}_S)$ 的每个元素都是 $n$ 次幂，例如 $H^0(S, \mathcal{O}_S)$ 是代数闭域的复合，或 $S$ 在代数闭域 $k$ 上既约且紧合. 则有典范同构
\[
H^1(S, \mu_n) \xrightarrow{\sim} {}_n \mathrm{Pic}(S).
\]
\end{corollary}

特别，当 $S$ 在代数闭域 $k$ 上紧合且连通时，这联系 $S$ 的基本群与 $S$ 在 $k$ 上的 Picard 概形 $P$ 的有限阶点；有同构
\[
\mathrm{Hom}(\pi_1(S), \mathbb{Z}/n\mathbb{Z}) \simeq {}_n P(k)
\]
对与特征互素的 $n$ 成立，它在代数几何中常用. 作为应用，当 $P$ 的连通分量 $P^0$ 是维数 $g$ 的完备群概形时，用第 2 号回顾的结果和 Néron-Severi 挠群的有限性，看到对每个与特征互素的素数 $\ell$，化为阿贝尔群的基本群 $\pi_1(S)$ 的 $\ell$-准素分量是 $\ell$-进整数环 $\mathbb{Z}_\ell$ 上有限型、秩 $2g$ 的模 (而且除至多有限个 $\ell$ 值外自由). 如 Serre 所注意的，这在某些条件下可以用来证明：$X$ 是连通 $S$ 上平坦射影概形时，$X$ 的诸纤维的 Picard 概形都有同样维数，只要应用半连续性定理 (\cref{cor:X.2.3})；Serre 的论证在 $X$ 在 $S$ 上的 Picard 概形存在、且 $X$ 在 $S$ 上的诸纤维的连通 Picard 群是紧合群概形时适用，例如 $X$ 在 $S$ 上的几何纤维正规时 (其中 $X$ 总在 $S$ 上平坦射影)，特别 $X$ 在 $S$ 上光滑射影时.

现在设 $p$ 是素数，假设 $S$ 是特征 $p$ 的预概形，i.e. $p \cdot \mathcal{O}_S = 0$. 则 $\mathcal{O}_S$ 中的 $p$ 次幂同态可加，对应态射 (把 $S$ 换成 $S$ 上变动的 $T$ 得到)：
\[
F \colon G_{a,S} \longrightarrow G_{a,S}
\]
因而是 $S$-群同态，称为 Frobenius 同态 (N.B. 这样的态射对每个由素数域 $\mathbb{Z}/p\mathbb{Z}$ 上的预概形 $G_0$ 经基扩张得到的 $S$-预概形 $G$ 定义，且 $G_0$ 是群概形时该态射是群同态). 令：
\[
\wp = \mathrm{id} - F \colon G_{a,S} \longrightarrow G_{a,S}.
\]
另一方面考虑由普通有限群 $\mathbb{Z}/p\mathbb{Z}$ 定义的 $S$-群 $(\mathbb{Z}/p\mathbb{Z})_S$，已说对每个 $S$-群 $G$，$(\mathbb{Z}/p\mathbb{Z})_S$ 到 $G$ 的 $S$-群同态与 $G$ 在 $S$ 上的 $p$ 次幂为单位截面的截面一一对应. 当 $G = G_{a,S}$ 时，它们对应 $G$ 在 $S$ 上的任意截面. 特别取 $G_{a,S}$ 在 $S$ 上对应环层 $\mathcal{O}_S$ 的单位截面的截面，找到 $S$-群同态
\[
i \colon (\mathbb{Z}/p\mathbb{Z})_S \longrightarrow G_{a,S}
\]

\begin{proposition}\label{prop:XI.6.7}
$S$-群同态序列
\[
0 \longrightarrow (\mathbb{Z}/p\mathbb{Z})_S \longrightarrow G_{a,S} \longrightarrow G_{a,S} \longrightarrow 0
\]
正合 (按第 4 号的意义). (称之为 $S$ 上的 \emph{Artin-Schreier 正合序列}\index[term]{Artin-Schreier zhenghexulie@Artin-Schreier 正合序列}).
\end{proposition}

只需在素数域 $k = \mathbb{Z}/p\mathbb{Z}$ 上证明. 只需注意同态 $\wp^{*} \colon k[t] \to k[t]$ (由 $\wp^{*}(t) = t - t^p$ 定义) 使 $k[t]$ 成为 $k[t]$ 上秩 $p$ 的自由模，确切地说 $k[t]$ 是 $k[s]$ (其中 $s = t - t^p$) 上以诸 $t^i$ ($0 \leqslant i \leqslant p - 1$) 为基的自由模.

由此用 \ref{prop:XI.4.5} 和 \ref{cor:XI.5.3} 得出：

\begin{proposition}\label{prop:XI.6.8}
有典范正合序列：
\[
0 \longrightarrow H^0(S, \mathbb{Z}/p\mathbb{Z}) \longrightarrow H^0(S, \mathcal{O}_S) \longrightarrow H^0(S, \mathcal{O}_S) \longrightarrow H^1(S, \mathbb{Z}/p\mathbb{Z}) \longrightarrow H^1(S, \mathcal{O}_S) \longrightarrow H^1(S, \mathcal{O}_S),
\]
从而有正合序列：
\[
0 \longrightarrow H^0(S, \mathcal{O}_S)/\wp H^0(S, \mathcal{O}_S) \longrightarrow H^1(S, \mathbb{Z}/p\mathbb{Z}) \longrightarrow H^1(S, \mathcal{O}_S)^F \longrightarrow 0,
\]
(其中最后一个项的指数 $F$ 表示端自同态 $F$ 的不变子群，等于 $\wp = \mathrm{id} - F$ 的核).
\end{proposition}

再明确两个极端情形：

\begin{corollary}\label{cor:XI.6.9}
假设 $H^1(S, \mathcal{O}_S)^F = 0$，例如 $S$ 是仿射概形. 则，令 $A = H^0(S, \mathcal{O}_S)$，有典范同构
\[
H^1(S, \mathbb{Z}/p\mathbb{Z}) \xrightarrow{\sim} A/\wp A.
\]
\end{corollary}

这是经典形式的 Artin-Schreier 理论，至少在 $A$ 是域的谱时.

\begin{corollary}\label{cor:XI.6.10}
假设 $\wp H^0(S, \mathcal{O}_S) = H^0(S, \mathcal{O}_S)$，例如 $H^0(S, \mathcal{O}_S)$ 是代数闭域的复合，或 $S$ 在代数闭域上紧合. 则有典范同构：
\[
H^1(S, \mathbb{Z}/p\mathbb{Z}) \xrightarrow{\sim} H^1(S, \mathcal{O}_S)^F
\]
\end{corollary}

\begin{remark}\label{rem:XI.6.11}
最后一个断言归功于 J.-P. Serre \cite{exp11-9}. 也可以对任意 $n$ 的 $n$ 结构群 $\mathbb{Z}/p^n\mathbb{Z}$ 发展类似理论，用 Witt 群概形 $W_n$ 代替 $G_a$，参见 loc. cit.. 注意在特征 $p > 0$ 时，Kummer 理论不再给出关于 $p$ 阶主覆盖的信息，因为 $\mu_p$ 此时是"无穷小"群，i.e. 在基上根式，与 $\mathbb{Z}/p\mathbb{Z}$ 没有直接关系；因此乍看，这些覆盖的理论 (当 $S$ 是代数闭域上的紧合概形以明确思想时) 不再服从 \ref{cor:XI.6.6} 那样的 Picard 概形理论. 然而，若回忆 $\mathrm{Pic}_{S/k}$ 在原点的 Zariski 切空间\footnote{关于 $\mathrm{Pic}_{S/K}$ 的定义，参见 A. Grothendieck，S\'em. Bourbaki N 232 (1962年2月).} 等同于 $H^1(S, \mathcal{O}_S)$，看到知道 $\mathrm{Pic}_{S/k}$ 中乘以 $p$ 的核 ${}_p \mathrm{Pic}_{S/k}$ 的群概形，蕴含知道 $H^1(S, \mathbb{Z}/p\mathbb{Z})$ 以及 $H^1(S, \mu_p)$；注意它也蕴含知道 $H^1(S, \alpha_p)$，其中 $\alpha_p$ 表示素数域上的无穷小群概形，即 $F \colon G_a \to G_a$ 的核 (也可描述为维数 1 的平凡受限 $p$-Lie 代数的受限包络代数的谱)：确实 \ref{prop:XI.4.5} 的正合序列这里给出：
\[
H^1(S, \alpha_p) \simeq \mathrm{Ker}(F \colon H^1(S, \mathcal{O}_S) \longrightarrow H^1(S, \mathcal{O}_S)),
\]
更一般地，用 $\alpha_{p^n}$ 记 $G_a$ 中 $F$ 的 $n$ 次迭代的核，有
\[
H^1(S, \alpha_{p^n}) \simeq \mathrm{Ker}(F^n \colon H^1(S, \mathcal{O}_S) \longrightarrow H^1(S, \mathcal{O}_S)).
\]
事实上，知道 ${}_p \mathrm{Pic}_{S/k}$ 等价于知道所有被 $p$ 零化的有限交换代数群 $G$ 的 $H^1(S, G)$，更一般地，知道 ${}_{p^n} \mathrm{Pic}_{S/k}$ 等价于知道所有被 $p^n$ 零化的有限交换代数群 $G$ 的 $H^1(S, G)$，由下列定理得出，它在所考虑的情形同时涵盖 Kummer 理论和 Artin-Schreier 理论：
设 $G$ 是 $k$ 上有限代数群，$D(G) = \mathrm{Hom}_{k\text{-群}}(G, G_m)$ 是它的 Cartier 对偶 (它的仿射代数由 $G$ 的仿射代数的对偶向量空间承载，i.e. 按 Dieudonn\'e-Cartier 的意义由 $G$ 的超代数承载)，则有典范同构：
\[
(\ast) \quad H^1(S, G) \simeq \mathrm{Hom}_{k\text{-群}}(D(G), \mathrm{Pic}_{S/k}).
\]
(N.B. $S$ 是 $k$ 代数闭上的紧合概形，使 $H^0(S, \mathcal{O}_S) = k$). 这个公式还可以表述为：$S$ 的"真正基本群" (第 2 号提到过) 化为阿贝尔群后同构于诸 $D(P_i)$ 的投射极限，其中 $P_i$ 遍历 $\mathrm{Pic}_{S/k}$ 的有限代数子群，记为 $T^{\bullet}(\mathrm{Pic}_{X/k})$. 当 $S$ 是阿贝尔簇时，在 \ref{thm:XI.2.1} 中已见这个群也同构于"真正"Tate 模 $T_\bullet(S) = \varprojlim_n {}_n S$，且同构 $(\ast)$ 此时写成更醒目的形式
\[
\mathrm{Ext}^1(A, G) \simeq \mathrm{Hom}(D(G), B),
\]
其中 $A$ 是阿贝尔簇，$B$ 是它的对偶，$G$ 是 $k$ 上有限代数群. 刚才指出的结果还可以推广到 $k$ 被任意底预概形代替、以及 $G$ 是有限群以外的其他系数群的情形.
\end{remark}

\section{文献}\printbibliography[heading=none]


\end{refsection}
\begin{refsection}% =====================================================================
% 第十二讲 (Exposé XII) 代数几何与解析几何
% (Géométrie algébrique et géométrie analytique)
% M. Raynaud
% 原书页码: pp. 239-257
% =====================================================================
\chapter{代数几何与解析几何}

\noindent Mme M. Raynaud\footnote{按 A. Grothendieck 的未发表笔记.}

如 \cite{exp12-10} 那样进行，给复数域 $\mathbb{C}$ 上每个局部有限型概形 $X$ 关联一个解析空间 $X^{\mathrm{an}}$\index[notations]{$X^{\mathrm{an}}$}，它的底集是 $X(\mathbb{C})$.

在本讲的第 2 和第 3 号中，给出 $X$ 与 $X^{\mathrm{an}}$ 的通常性质之间、以及态射 $f \colon X \to Y$ 与关联态射 $f^{\mathrm{an}} \colon X^{\mathrm{an}} \to Y^{\mathrm{an}}$ 的性质之间的"字典".

然后证明 $X$ 与 $X^{\mathrm{an}}$ 上凝聚层之间的比较定理 (在 \cite[第 12 号]{exp12-10} 中对射影簇建立) 在 $X$ 是紧合概形时仍然成立.

最后在第 5 号中证明 $X$ 的有限平展覆盖范畴与 $X^{\mathrm{an}}$ 的有限平展覆盖范畴等价. 作为给读者的赠品，给出 \emph{Grauert-Remmert 定理} \cite{exp12-6} 的新证明，用\emph{奇点消解} \cite{exp12-8}.

\section{概形的关联解析空间}

设 $X$ 是 $\mathbb{C}$ 上局部有限型概形. 设 $\Phi$ 是从解析空间范畴 \cite[第 9 号]{exp12-4} 到集合范畴的函子，它把解析空间 $X$ 配给 $\mathbb{C}$-代数环化空间态射集 $\mathrm{Hom}_{\mathbb{C}}(X, X)$. 有下列定理：

\begin{theorem}[定理与定义]\label{thm:XII.1.1}
函子 $\Phi$ 被解析空间 $X^{\mathrm{an}}$ 和态射 $\phi \colon X^{\mathrm{an}} \to X$ 表示. 称 $X^{\mathrm{an}}$ 为 $X$ 的\emph{关联解析空间}\index[term]{guanlianjiexikongjian@关联解析空间}.
\end{theorem}

若 $|X^{\mathrm{an}}|$ 是 $X^{\mathrm{an}}$ 的底集，$\phi$ 诱导从 $|X^{\mathrm{an}}|$ 到 $X$ 的在 $\mathbb{C}$ 中取值的点集 $X(\mathbb{C})$ 的双射. 而且，对 $X^{\mathrm{an}}$ 的每个点 $x$，态射
\[
\phi_x \colon \mathcal{O}_{X,\phi(x)} \longrightarrow \mathcal{O}_{X^{\mathrm{an}},x},
\]
必然局部，经完备化给出同构
\[
\widehat{\phi_x} \colon \widehat{\mathcal{O}}_{X,\phi(x)} \xrightarrow{\sim} \widehat{\mathcal{O}}_{X^{\mathrm{an}},x}.
\]
特别态射 $\phi$ 平坦.

注意 $\phi$ 诱导 $X^{\mathrm{an}}$ 到 $X(\mathbb{C})$ 的双射的事实由 $X^{\mathrm{an}}$ 的万有性质得出. 另一方面有下列断言：

a) 若定理对概形 $Y$ 成立，则对 $Y$ 的任何子概形 $X$ 也成立. 先假设 $X$ 是 $Y$ 的开子概形；若 $\psi \colon Y^{\mathrm{an}} \to Y$ 是典范态射，$\psi^{-1}(X)$ 是 $Y^{\mathrm{an}}$ 的开集，配备 $Y^{\mathrm{an}}$ 诱导的解析空间结构. 由于解析空间 $X$ 到 $X$ 的每个态射由后者的万有性质经 $Y^{\mathrm{an}}$ 分解，从而经 $X^{\mathrm{an}}$ (即纤维积 $Y^{\mathrm{an}} \times_Y X$) 分解，$X^{\mathrm{an}}$ 是 $X$ 的关联解析空间. 最后关于 $\phi_x$ 的断言显然.

现在只需考虑 $X$ 是 $Y$ 的闭子概形的情形. 设 $I$ 是定义 $X$ 的 $\mathcal{O}_Y$-凝聚理想；则 $I \cdot \mathcal{O}_{Y^{\mathrm{an}}}$ 是 $\mathcal{O}_{Y^{\mathrm{an}}}$ 上的凝聚理想层，它定义 $Y^{\mathrm{an}}$ 的闭解析子空间 $X^{\mathrm{an}}$；如开子概形的情形那样看到 $X^{\mathrm{an}}$ 是 $X$ 的关联解析空间. 设 $\phi \colon X^{\mathrm{an}} \to X$ 是典范态射. 对 $X^{\mathrm{an}}$ 的每个点 $x$，态射 $\phi_x$ 正是由 $\psi_x$ 诱导的态射
\[
\mathcal{O}_{Y,\psi(x)}/I_{\psi(x)} \longrightarrow \mathcal{O}_{Y^{\mathrm{an}},x}/I_{\psi(x)} \cdot \mathcal{O}_{Y^{\mathrm{an}},x}
\]
的完备化
\[
\widehat{\phi_x} \colon \widehat{\mathcal{O}}_{Y,\psi(x)}/I_{\psi(x)} \cdot \widehat{\mathcal{O}}_{Y,\psi(x)} \longrightarrow \widehat{\mathcal{O}}_{Y^{\mathrm{an}},x}/I_{\psi(x)} \cdot \widehat{\mathcal{O}}_{Y^{\mathrm{an}},x}
\]
是同构，因为 $\widehat{\psi_x}$ 是同构，这证明 a).

b) 若有两个 $\mathbb{C}$-概形 $X_1$，$X_2$，使 $X_1^{\mathrm{an}}$ 和 $X_2^{\mathrm{an}}$ 存在，则 $(X_1 \times X_2)^{\mathrm{an}}$ 也存在. 设 $\phi_1 \colon X_1^{\mathrm{an}} \to X_1$，$\phi_2 \colon X_2^{\mathrm{an}} \to X_2$ 是典范态射，$p_1$，$p_2$ 是 $X_1^{\mathrm{an}} \times X_2^{\mathrm{an}}$ 的两个投影. 由 EGA I 1.8.1 形式地得出 $X_1 \times X_2$ 是 $X_1$ 和 $X_2$ 在局部环环化空间范畴中的积；由此得出态射 $\phi_1 \cdot p_1$ 和 $\phi_2 \cdot p_2$ 定义态射 $\phi \colon X_1^{\mathrm{an}} \times X_2^{\mathrm{an}} \to X_1 \times X_2$，且对 $(X_1^{\mathrm{an}} \times X_2^{\mathrm{an}}, \phi)$ 表示函子 $X \mapsto \mathrm{Hom}_{\mathbb{C}}(X, X_1 \times X_2)$.

c) 若 $E^1$ 表示维数 1 的仿射空间，i.e. 配备全纯函数\index[term]{quanchunhanshu@全纯函数}层的拓扑空间 $\mathbb{C}$，函子\footnote{[2003 年增补：$E^1_{\mathbb{C}}$ 表示 $\mathbb{C}$ 上的 (代数) 仿射直线.]} $X \mapsto \mathrm{Hom}_{\mathbb{C}}(X, E^1_{\mathbb{C}})$ 被 $E^1$ 表示，典范态射 $\phi \colon E^1 \to E^1_{\mathbb{C}}$ 是显然态射. 确实，给出解析空间 $X$ 到 $E^1_{\mathbb{C}}$ 的态射等价于给出 $\Gamma(X, \mathcal{O}_X)$ 的元素，后者也等价于给出 $X$ 到 $E^1$ 的态射. 显然有双射 $|E^1| \simeq E^1_{\mathbb{C}}(\mathbb{C})$，且对每个点 $x \in E^1$，态射 $\widehat{\phi_x}$ 正是 $\mathbb{C}$ 上一元形式幂级数环的恒等态射.

由 b) 和 c) 得出定理对仿射空间 $E^n_{\mathbb{C}}$ ($n \geqslant 0$) 成立. 用 a)，看到对 $\mathbb{C}$ 上每个局部有限型仿射概形 $X$ 也成立. 若不再假设 $X$ 仿射，且 $(X_i)$ 是 $X$ 的仿射开集覆盖，由万有性质和 a) 得出诸 $X_i^{\mathrm{an}}$ 粘合，从而定义 $X$ 的关联解析空间 $X^{\mathrm{an}}$.

\refstepcounter{theorem}\noindent\textbf{\thesection.\arabic{theorem}.}\quad 设 $f \colon X \to Y$ 是 $\mathbb{C}$ 上局部有限型概形的态射. 若 $\phi \colon X^{\mathrm{an}} \to X$ 和 $\psi \colon Y^{\mathrm{an}} \to Y$ 是典范态射，由 $Y^{\mathrm{an}}$ 的万有性质得出存在唯一态射 $f^{\mathrm{an}} \colon X^{\mathrm{an}} \to Y^{\mathrm{an}}$，使图式
\[
\begin{tikzcd}[row sep=large, column sep=large]
X^{\mathrm{an}} \arrow[r] \arrow[d,"{f^{\mathrm{an}}}"] & X \arrow[d,"{f}"] \\
Y^{\mathrm{an}} \arrow[r] & Y
\end{tikzcd}
\]
交换. 于是定义了从 $\mathbb{C}$ 上局部有限型概形范畴到解析空间范畴的函子 $\Phi$.

函子 $\Phi$ 与有限投射极限交换. 只需看到 $\Phi$ 与纤维积交换. 若 $X$，$Y$，$Z$ 是 $\mathbb{C}$ 上局部有限型概形，由 $X \times_Z Y$ 是 $X$ 和 $Y$ 在 $Z$ 上方的局部环环化空间范畴中的纤维积的事实得出 $X^{\mathrm{an}} \times_{Z^{\mathrm{an}}} Y^{\mathrm{an}}$ 满足刻画 $(X \times_Z Y)^{\mathrm{an}}$ 的万有性质.

\refstepcounter{theorem}\noindent\textbf{\thesection.\arabic{theorem}.}\quad 设 $X$ 是 $\mathbb{C}$ 上局部有限型概形，$X^{\mathrm{an}}$ 是关联解析空间，$\phi \colon X^{\mathrm{an}} \to X$ 是典范态射. 若 $F$ 是 $\mathcal{O}_X$-模，逆像 $\phi^{*} F = F^{\mathrm{an}}$ 是 $\mathcal{O}_{X^{\mathrm{an}}}$ 上的模层. 于是定义了从 $\mathcal{O}_X$-模范畴到 $X^{\mathrm{an}}$ 上模范畴的函子. 该函子与归纳极限交换 (EGA IV$_0$ 4.3.2). 层 $\mathcal{O}_{X^{\mathrm{an}}}$ 凝聚 \cite[第 18 号 §2 定理 2]{exp12-4}，故它把凝聚层变成\emph{凝聚层} (EGA IV$_0$ 5.3.11). 还有：

\begingroup
\addtocounter{theorem}{-1}
\renewcommand{\thetheorem}{\thesection.3.1}
\begin{proposition}\label{prop:XII.1.3.1}
把 $\mathcal{O}_{X^{\mathrm{an}}}$-模 $F$ 配给它在 $X^{\mathrm{an}}$ 上的逆像 $F^{\mathrm{an}}$ 的函子正合、忠实、保守.
\end{proposition}
\endgroup

正合性由态射 $\phi \colon X^{\mathrm{an}} \to X$ 平坦 (\cref{thm:XII.1.1}) 得出. 证明函子 $F \mapsto F^{\mathrm{an}}$ 忠实. 考虑到正合性，只需证明 $F^{\mathrm{an}}$ 为零时 $F$ 为零. 对 $X^{\mathrm{an}}$ 的每个点 $x$，有 $F_{\phi(x)} \otimes_{\mathcal{O}_{X,\phi(x)}} \mathcal{O}_{X^{\mathrm{an}},x} = 0$. 态射 $\mathcal{O}_{X,\phi(x)} \to \mathcal{O}_{X^{\mathrm{an}},x}$ 忠实平坦，故对 $X$ 的每个闭点 $\phi(x)$ 有 $F_{\phi(x)} = 0$，且 $X$ 是 Jacobson 概形 (EGA IV 10.4.8)，这蕴含 $F$ 为零.

函子 $F \mapsto F^{\mathrm{an}}$ 保守的事实由正合性和忠实性形式地得出.

\section{概形与关联解析空间的性质比较}

\begin{proposition}\label{prop:XII.2.1}
设 $X$ 是 $\mathbb{C}$ 上局部有限型概形，$X^{\mathrm{an}}$ 是关联解析空间，$n$ 是整数. 考虑性质 $P$：
\begin{enumerate}[i)]
\item 非空
\item[i')] 离散
\item[ii)] Cohen-Macaulay
\item[iii)] $(S_n)$
\item[iv)] 正则
\item[v)] $(R_n)$
\item[vi)] 正规
\item[vii)] 既约
\item[viii)] 维数 $n$.
\end{enumerate}
则 $X$ 具有性质 $P$ 当且仅当 $X^{\mathrm{an}}$ 具有性质 $P$.
\end{proposition}

设 $\phi \colon X^{\mathrm{an}} \to X$ 是典范态射. (i) 由 $|X^{\mathrm{an}}| = X(\mathbb{C})$ (\cref{thm:XII.1.1}) 和 $X$ 是 Jacobson 概形 (EGA IV 10.4.8) 得出. 说 $X$ (resp. $X^{\mathrm{an}}$) 离散等价于说 $\dim X = 0$ (resp. 由 \cite[第 19 号 §4 推论 6]{exp12-4}，$\dim X^{\mathrm{an}} = 0$)；故 (i') 由 (viii) 得出.

设 $P$ 是 (ii) 至 (vii) 之一. 为使 $X$ 具有性质 $P$，必要且充分条件是 $P$ 在 $X$ 的每个闭点成立；确实，$X$ 是 excellent (EGA IV 7.8.6 (iii))，故 $X$ 验证 $P$ 的点集开 (loc. cit.)，且该开集含所有闭点时等于整个 $X$. 说 $X$ (resp. $X^{\mathrm{an}}$) 具有性质 $P$ 于是等价于：对 $X^{\mathrm{an}}$ 的每个点 $x$，局部环 $\mathcal{O}_{X,\phi(x)}$ (resp. $\mathcal{O}_{X^{\mathrm{an}},x}$) 具有性质 $P$. 由于 excellent 局部环具有性质 $P$ 的事实可在完备化后看到，命题由同构 $\widehat{\mathcal{O}}_{X,\phi(x)} \xrightarrow{\sim} \widehat{\mathcal{O}}_{X^{\mathrm{an}},x}$ 在情形 (ii) 至 (vii) 中得出. 情形 (viii) 也如此，考虑到关系
\[
\dim X = \sup_x \dim \mathcal{O}_{X,\phi(x)}, \quad \dim X^{\mathrm{an}} = \sup_x \dim \mathcal{O}_{X^{\mathrm{an}},x},
\]
其中 $x \in X^{\mathrm{an}}$. 这完成证明.

\begin{proposition}\label{prop:XII.2.2}
设 $X$ 是 $\mathbb{C}$ 上局部有限型概形，$\phi \colon X^{\mathrm{an}} \to X$ 是典范态射，$T$ 是 $X$ 的局部可构部分. 则有关系
\[
\overline{\phi^{-1}(T)} = \phi^{-1}(\overline{T}).
\]
\end{proposition}

可假设 $T$ 是 $X$ 的稠密开集. 设 $H$ 是 $X$ 的底集为 $X - T$ 的既约闭子概形；关联空间 $H^{\mathrm{an}}$ 是 $X^{\mathrm{an}}$ 的底集为 $X^{\mathrm{an}} - \phi^{-1}(T)$ 的闭解析子空间. 需证明 $H^{\mathrm{an}}$ 的每个点 $x$ 属于 $\phi^{-1}(\overline{T})$. 在这样的点 $x$，解析空间芽 $(X^{\mathrm{an}}, x)$ 含子芽 $(H^{\mathrm{an}}, x)$，后者由 $\mathcal{O}_{X^{\mathrm{an}},x}$ 的非幂零理想定义. 于是由 Nullstellensatz \cite[第 19 号 §4 推论 3]{exp12-4} 得出 $x$ 的每个开邻域含有 $X^{\mathrm{an}}$ 的不属于 $H^{\mathrm{an}}$ 的点，这确实证明 $x \in \phi^{-1}(\overline{T})$.

\begin{corollary}\label{cor:XII.2.3}
设 $X$ 是 $\mathbb{C}$ 上局部有限型概形，$\phi \colon X^{\mathrm{an}} \to X$ 是典范态射，$T$ 是 $X$ 的局部可构部分. 为使 $T$ 是开部分 (resp. 闭部分，resp. 稠密部分)，必要且充分条件是 $\phi^{-1}(T)$ 如此.
\end{corollary}

该推论由 \ref{prop:XII.2.2} 和 $X$ 是 Jacobson 概形 (EGA IV 10.4.8)、故 $X$ 的两个在非常稠密集 $X(\mathbb{C})$ 上有同样迹的局部可构部分相等的事实得出.

\begin{proposition}\label{prop:XII.2.4}
设 $X$ 是 $\mathbb{C}$ 上局部有限型概形. 为使 $X$ 连通 (resp. 不可约)，必要且充分条件是 $X^{\mathrm{an}}$ 如此.
\end{proposition}

假设 $X^{\mathrm{an}}$ 连通 (resp. 不可约). $X^{\mathrm{an}}$ 在 $X$ 中的像 $X(\mathbb{C})$ 此时连通 (resp. 不可约). 由此得出 $X$ 连通 (resp. 不可约)，因为 $X$ 与 $X(\mathbb{C})$ 的闭部分一一对应 (EGA IV 10.1.2).

反之假设 $X$ 连通 (resp. 不可约)，证明 $X^{\mathrm{an}}$ 亦然. 可限于 $X$ 不可约的情形. 假设 $X$ 连通. 给定 $X$ 的点 $x$，存在 $X$ 的不可约闭子概形 $X_1, \ldots, X_n$ 的有限序列 (其中 $x \in X_1$，$y \in X_n$，且对 $1 \leqslant i \leqslant n - 1$ 有 $X_i \cap X_{i+1} \neq \varnothing$) 的点 $y \in X$ 的集同时开且闭，故等于整个 $X$. 对前述那样的序列 $X_1, \ldots, X_n$，对 $1 \leqslant i \leqslant n - 1$ 也有 $X_i^{\mathrm{an}} \cap X_{i+1}^{\mathrm{an}} \neq \varnothing$；若已证诸 $X_i^{\mathrm{an}}$ 连通，则 $X^{\mathrm{an}}$ 也连通.

此后假设 $X$ 不可约. 还可假设 $X$ 仿射. 若 $(U_i)_{i \in I}$ 是 $X$ 的仿射开集覆盖，其中任意两个开集的交非空，则同样的性质对 $X^{\mathrm{an}}$ 的覆盖 $(U_i^{\mathrm{an}})_{i \in I}$ 成立；若已证诸 $U_i^{\mathrm{an}}$ 不可约，则 $X^{\mathrm{an}}$ 也如此.

还可假设 $X$ 正规. 设 $\widetilde{X}$ 是 $X$ 的正规化；由于态射 $\widetilde{X} \to X$ 满射，$\widetilde{X}^{\mathrm{an}} \to X^{\mathrm{an}}$ 也满射，这证明 $\widetilde{X}^{\mathrm{an}}$ 不可约时 $X^{\mathrm{an}}$ 也如此.

此后假设 $X$ 仿射正规. 由于 $X^{\mathrm{an}}$ 的局部环整，说 $X^{\mathrm{an}}$ 不可约与说它连通相同. 确实，若 $F$ 是 $X^{\mathrm{an}}$ 的闭解析部分，其中 $\mathrm{codim}_x(F, X^{\mathrm{an}}) = 0$ 的点 $x$ 的集是 $X^{\mathrm{an}}$ 的闭解析子集 \cite[第 20 号 A 推论 1]{exp12-4}，也开；若 $X^{\mathrm{an}}$ 连通，这证明 $F \neq X^{\mathrm{an}}$ 时 $F$ 稀疏，故 $X^{\mathrm{an}}$ 不可约. 于是归结到证明 $X^{\mathrm{an}}$ 连通.

设
\[
i \colon X \longrightarrow P
\]
是 $X$ 的紧化，其中 $P$ 是射影正规 $\mathbb{C}$-概形，$i$ 是支配开浸入. 由 \cite[第 12 号 定理 1]{exp12-10} 得出 $P^{\mathrm{an}}$ 连通. 由于 $X^{\mathrm{an}}$ 是从 $P^{\mathrm{an}}$ 去掉稀疏闭解析部分得到的，由下面 \ref{lem:XII.2.5} 得出 $X^{\mathrm{an}}$ 也连通.

\begin{lemma}\label{lem:XII.2.5}
设 $P$ 是连通正规解析空间，$Y$ 是稀疏闭解析部分，则 $X = P - Y$ 连通.
\end{lemma}

当 $Y$ 余维数 $\geqslant 2$ 时，命题由 \cite[第 3 号 命题 4]{exp12-11} 得出. 在一般情形，可以假设 (把 $P$ 中余维数 $\geqslant 2$ 的闭解析部分去掉) $P$ 和 $Y$ (看作 $P$ 的既约解析子空间) 正则. 由隐函数定理，$Y$ 的每个点 $y$ 有同构于仿射空间 $E^n$ 的球的邻域 $U$，使 $U \cap Y$ 由若干坐标函数的消没定义. 这证明 $U - U \cap Y$ 连通，故 $X$ 也连通.

\begin{corollary}\label{cor:XII.2.6}
设 $X$ 是 $\mathbb{C}$ 上局部有限型概形；由典范态射 $X^{\mathrm{an}} \to X$ 诱导的态射
\[
\pi_0(X^{\mathrm{an}}) \longrightarrow \pi_0(X)
\]
双射.
\end{corollary}

\section{态射性质的比较}

\begin{proposition}\label{prop:XII.3.1}
设 $f \colon X \to Y$ 是 $\mathbb{C}$ 上局部有限型概形的态射，$f^{\mathrm{an}} \colon X^{\mathrm{an}} \to Y^{\mathrm{an}}$ 是 $f$ 在关联解析空间上导出的态射. 设 $P$ 是下列性质：
\begin{enumerate}[i)]
\item 平坦
\item[ii)] 非分歧 (i.e. 非分歧)
\item[iii)] 平展
\item[iv)] 光滑
\item[v)] 正规
\item[vi)] 既约
\item[vii)] 单射
\item[viii)] 分离
\item[ix)] 同构
\item[x)] 单态射
\item[xi)] 开浸入.
\end{enumerate}
则 $f$ 具有性质 $P$ 当且仅当 $f^{\mathrm{an}}$\index[notations]{$f^{\mathrm{an}}$} 具有性质 $P$.
\end{proposition}

记 $\phi \colon X^{\mathrm{an}} \to X$ 和 $\psi \colon Y^{\mathrm{an}} \to Y$ 为典范态射. 设 $x$ 是 $X^{\mathrm{an}}$ 的点，$y = f^{\mathrm{an}}(x)$. 由 $f$ 和 $f^{\mathrm{an}}$ 导出的态射 $\mathcal{O}_{Y^{\mathrm{an}},y} \to \mathcal{O}_{X^{\mathrm{an}},x}$ 和 $\mathcal{O}_{Y,\psi(y)} \to \mathcal{O}_{X,\phi(x)}$ 经完备化给出同样态射 (\cref{thm:XII.1.1}). 由 \cite[第 3 章 §5 命题 4]{exp12-2} (resp. EGA IV 17.4.4)，说 $f^{\mathrm{an}}$ 验证性质 (i) (resp. (ii)) 与说 $f$ 在 $X$ 的每个闭点验证 (i) (resp. (ii)) 相同. 由于 $X$ 中验证 (i) (resp. (ii)) 的点集开 (EGA IV 11.1.1 和 I 3.3)，这证明 (i) 和 (ii)，从而 (iii).

设 $P$ 是性质 (iv) (resp. (v)，resp. (vi)). 考虑到 \ref{prop:XII.2.1} ((v)，(vi)，(vii))，说 $f^{\mathrm{an}}$ 在 $Y^{\mathrm{an}}$ 的各点 $y$ 处的几何纤维正则 (resp. 正规，resp. 既约) 与说 $f$ 在 $Y$ 的各闭点 $\psi(y)$ 处的几何纤维如此相同. 情形 (iv) (resp. (v)，resp. (vi)) 于是由 (i) 和 $Y$ 中 $f$ 的几何纤维正则的点集开 (EGA IV 12.1.7) 的事实得出.

(vii). 若 $f$ 单射，$f^{\mathrm{an}}$ 也单射. 反之假设 $f^{\mathrm{an}}$ 单射，证明 $f$ 也单射. 可假设 $f$ 有限型. 由于 $f^{\mathrm{an}}$ 单射，$f$ 在 $Y$ 的闭点处的纤维根式；由于 $Y$ 中纤维根式的点集局部可构 (EGA IV 9.6.1)，且 $Y$ 是 Jacobson 概形，$f$ 的所有纤维根式故单射.

(viii). 设 $\Delta \colon X \to X \times_Y X$ 和 $\Delta^{\mathrm{an}} \colon X^{\mathrm{an}} \to X^{\mathrm{an}} \times_{Y^{\mathrm{an}}} X^{\mathrm{an}}$ 是对角浸入，$\Theta \colon X^{\mathrm{an}} \times_{Y^{\mathrm{an}}} X^{\mathrm{an}} \to X \times_Y X$ 是典范态射. 由 \ref{cor:XII.2.3}，说 $\Delta(X)$ 在 $X \times_Y X$ 中闭与说 $\Delta^{\mathrm{an}}(X^{\mathrm{an}})$ 在 $X^{\mathrm{an}} \times_{Y^{\mathrm{an}}} X^{\mathrm{an}}$ 中闭相同.

由于开浸入正是平展单射态射 (EGA IV 17.9.1 和 \cite[第 13 号 §1]{exp12-4})，(xi) 由 (iii) 和 (vii) 得出. 同构正是满射开浸入，故 (ix) 由 (xi) 和下面 (3.2 (i)) 得出. 说 $f$ 是单态射等价于说对角态射 $\Delta \colon X \to X \times_Y X$ 是同构，故 (x) 由 (ix) 得出.

\begin{proposition}\label{prop:XII.3.2}
设 $X$ 和 $Y$ 是 $\mathbb{C}$ 上两个局部有限型概形，$f \colon X \to Y$ 是有限型态射，$f^{\mathrm{an}} \colon X^{\mathrm{an}} \to Y^{\mathrm{an}}$ 是 $f$ 在关联解析空间上导出的态射. 设 $P$ 是下列性质：
\begin{enumerate}[i)]
\item 满射
\item[ii)] 支配
\item[iii)] 闭浸入
\item[iv)] 浸入
\item[v)] 紧合\footnote{称解析空间态射紧合，若按 \cite[第 1 章 §10 第 1 号]{exp12-1} 的意义紧合且分离.}
\item[vi)] 有限.
\end{enumerate}
则 $f$ 具有性质 $P$ 当且仅当 $f^{\mathrm{an}}$ 具有性质 $P$.
\end{proposition}

设 $\phi \colon X^{\mathrm{an}} \to X$ 和 $\psi \colon Y^{\mathrm{an}} \to Y$ 是典范态射.

(i). 若 $f$ 满射，对 $Y^{\mathrm{an}}$ 的每个点 $y$，$f^{-1}(\psi(y))$ 是 $X$ 的非空闭部分；它含至少一个闭点，这证明 $f^{\mathrm{an}}$ 满射. 反之，若 $f^{\mathrm{an}}$ 满射，$f(X)$ 是 $Y$ 的局部\emph{可构部分} (EGA IV 1.8.4)，含 $Y$ 的所有闭点；故有 $f(X) = Y$.

(ii) 由 \ref{prop:XII.2.2} 得出.

(iii). 若 $f$ 是闭浸入，由 \ref{thm:XII.1.1} a) $f^{\mathrm{an}}$ 也如此. 反之，若 $f^{\mathrm{an}}$ 是闭浸入，由 \ref{prop:XII.3.1} (x) 和 \ref{prop:XII.3.2} (v) $f$ 也如此，因为说 $f$ 是闭浸入等价于说 $f$ 是紧合单态射 (EGA IV 8.11.5).

(iv). 显然 $f$ 是浸入时 $f^{\mathrm{an}}$ 也是. 反之假设 $f^{\mathrm{an}}$ 是浸入，设 $T$ 是 $X$ 在 $Y$ 中的像，$\overline{T}$ 是 $f$ 的概形闭包. 有 $f$ 的分解
\[
X \xrightarrow{i} \overline{T} \xrightarrow{j} Y,
\]
其中 $j$ 是闭浸入，$i$ 是典范态射，由此得到 $f^{\mathrm{an}}$ 的分解
\[
X^{\mathrm{an}} \xrightarrow{i^{\mathrm{an}}} \overline{T}^{\mathrm{an}} \xrightarrow{j^{\mathrm{an}}} Y^{\mathrm{an}}.
\]
由于 $T = f(X)$ 是 $Y$ 的局部可构部分 (EGA IV 1.8.4)，由 \ref{prop:XII.2.2} 有 $\overline{T}^{\mathrm{an}} = \overline{f^{\mathrm{an}}(X^{\mathrm{an}})}$. 由此得出 $i^{\mathrm{an}}(X^{\mathrm{an}})$ 是 $\overline{T}^{\mathrm{an}}$ 的开集，故 $i(X)$ 是 $\overline{T}$ 的开集. 考虑 $i$ 的典范分解
\[
X \xrightarrow{i_1} i(X) \xrightarrow{i_2} \overline{T}.
\]
态射 $i_1^{\mathrm{an}}$ 是紧合单态射，故由 \ref{prop:XII.3.2} (v) 和 \ref{prop:XII.3.1} (x) $i_1$ 也如此；这证明 $i_1$ 从而 $f$ 是浸入.

(v). 假设 $f$ 紧合，证明 $f^{\mathrm{an}}$ 也紧合. $f^{\mathrm{an}}$ 紧合的事实局部于 $Y^{\mathrm{an}}$，可假设 $Y$ 仿射. 由 Chow 引理 (EGA II 5.6.1)，可以找到 $Y$-射影概形 $X'$ 和射影满射态射
\[
g \colon X' \longrightarrow X.
\]
态射 $(fg)^{\mathrm{an}} = f^{\mathrm{an}} g^{\mathrm{an}}$ 射影故紧合，$g^{\mathrm{an}}$ 满射，由 \cite[第 1 章 §10]{exp12-1} 得出 $f^{\mathrm{an}}$ 紧合.

反之假设 $f^{\mathrm{an}}$ 紧合，证明 $f$ 也紧合. 由 \ref{prop:XII.3.1} (viii) $f$ 分离. 剩下证明 $f$ 万有闭，甚至只需证明 $f$ 闭；确实，对每个局部有限型 $Y$-概形 $Y'$，态射
\[
f_{(Y')} = h \colon X \times_Y Y' \longrightarrow Y'
\]
也闭，因为 $h^{\mathrm{an}}$ 紧合. 设 $T$ 是 $X$ 的闭部分；$f(T)$ 是局部可构集，且有
\[
f^{\mathrm{an}}(\phi^{-1}(T)) = \psi^{-1}(f(T)).
\]
由于 $f^{\mathrm{an}}$ 紧合，$\psi^{-1}(f(T))$ 是 $Y^{\mathrm{an}}$ 的闭部分，由 \ref{prop:XII.2.2} 得出
\[
\overline{\psi^{-1}(f(T))} = \psi^{-1}(\overline{f(T)}).
\]
这蕴含 $f(T) = \overline{f(T)}$，i.e. $f$ 闭故紧合.

(vi). 说态射有限与说它纤维有限紧合相同 (EGA III 4.4.2 和 \cite[第 19 号 §5]{exp12-4}). 由于 $f$ 的纤维有限的点集局部可构 (EGA IV 9.7.9)，$f$ 的诸纤维有限当且仅当 $f^{\mathrm{an}}$ 的诸纤维如此；故 (vi) 由 (v) 得出.

\begin{remark}\label{rem:XII.3.3}
a) 设 $f \colon X \to Y$ 是 $\mathbb{C}$ 上局部有限型概形的态射. $f^{\mathrm{an}}$ 是局部同构的事实不蕴含 $f$ 也如此. 确实，若 $f$ 平展，$f^{\mathrm{an}}$ 平展故是局部同构 \cite[第 13 号 §1]{exp12-4}，但 $f$ 未必如此.

b) 不假设 $f$ 有限型时断言 \ref{prop:XII.3.2} 不成立. 例如证明 $f^{\mathrm{an}}$ 可以是闭浸入而 $f$ 不是. 只需取 $X$ 为 $\mathbb{Z}$ 个 $\mathrm{Spec} \mathbb{C}$ 的和，$Y$ 为仿射直线，$f$ 为把 $X$ 的点送到 $Y$ 的构成离散部分的互异点的态射.
\end{remark}

\section{上同调比较定理与存在定理}

本号的对象是重新证明 \cite[第 2 号 定理 5 和 定理 6]{exp12-3} 的结果；后者把 $X$ 射影时在 \cite[第 12 号]{exp12-10} 中建立的定理推广到紧合概形的情形，并推广到相对情形. 关于解析空间上的相对紧合概形的更一般结果在 \cite[第八章 第 3 号]{exp12-7} 中证明.

回忆 \cite[第 12 号]{exp12-10} 中用的 Čech 上同调在代数情形与解析情形都等于通常上同调 (EGA III 1.4.1 和 \cite[II 5.10]{exp12-5}).

\refstepcounter{theorem}\noindent\textbf{\thesection.\arabic{theorem}.}\quad 设 $f \colon X \to Y$ 是 $\mathbb{C}$ 上局部有限型概形的态射，考虑交换图式
\[
\begin{tikzcd}[row sep=large, column sep=large]
X^{\mathrm{an}} \arrow[r,"{\phi}"] \arrow[d,"{f^{\mathrm{an}}}"] & X \arrow[d,"{f}"] \\
Y^{\mathrm{an}} \arrow[r,"{\psi}"] & Y.
\end{tikzcd}
\]
若 $F$ 是 $\mathcal{O}_X$-模，对每个整数 $p \geqslant 0$ 有态射
\[
R^p f_* F \xrightarrow{i} R^p f_*(\phi_* F^{\mathrm{an}}) \xrightarrow{j} R^p (f \cdot \phi)_* F^{\mathrm{an}} \xrightarrow{k} \psi_* (R^p f_*^{\mathrm{an}} F^{\mathrm{an}}),
\]
其中 $i$ 由典范态射 $F \to \phi_* F^{\mathrm{an}}$ 得出，$j$，$k$ 是 Leray 谱序列的"边缘同态". 与复合 $k \cdot j \cdot i$ 关联典范态射
\[
(4.1.1) \quad \theta_p \colon (R^p f_* F)^{\mathrm{an}} \longrightarrow R^p f_*^{\mathrm{an}} (F^{\mathrm{an}})
\]

\begin{theorem}\label{thm:XII.4.2}
设 $f \colon X \to Y$ 是 $\mathbb{C}$ 上局部有限型概形的紧合态射，$F$ 是凝聚 $\mathcal{O}_X$-模. 则对每个整数 $p \geqslant 0$，态射 (4.1.1)
\[
\theta_p \colon (R^p f_* F)^{\mathrm{an}} \longrightarrow R^p f_*^{\mathrm{an}} (F^{\mathrm{an}})
\]
是同构.
\end{theorem}

1) $f$ 射影的情形. 证明与 \cite[第 13 号]{exp12-10} 类似. 简要回忆. 归结到 $X$ 是 $Y$ 上方射影空间型 $\mathbf{P}^r_Y$ 的情形. 设 $\mathcal{Y} = Y^{\mathrm{an}}$，$\mathcal{P} = \mathbf{P}^r_\mathcal{Y}$；先证明
\[
f_*^{\mathrm{an}} \mathcal{O}_{\mathcal{P}} = \mathcal{O}_{\mathcal{Y}}, \quad R^p f_*^{\mathrm{an}}(\mathcal{O}_{\mathcal{P}}) = 0 \quad \text{对 } p > 0.
\]
为验证前面关系，可归结到 $\mathcal{Y}$ 是仿射空间 $E^n$ 的球 $B$ 的情形. 考虑 $\mathcal{P}$ 的由 $r + 1$ 个同构于 $B \times E^r$ 的开集组成的"标准覆盖" $\{U_i\}$. 由于这些开集是 Stein 的，对每个整数 $p \geqslant 0$ 有同构
\[
H^p(\{U_i\}, \mathcal{O}_{\mathcal{P}}) \xrightarrow{\sim} H^p(\mathcal{P}, \mathcal{O}_{\mathcal{P}})
\]
于是可以把结构层 $\mathcal{O}_{\mathcal{P}}$ 在开集 $U_i$ 及其交上的截面用 Laurent 级数表达；容易的计算证明
\[
H^0(\mathcal{P}, \mathcal{O}_{\mathcal{P}}) \xrightarrow{\sim} H^0(\mathcal{Y}, \mathcal{O}_{\mathcal{Y}}), \quad H^p(\mathcal{P}, \mathcal{O}_{\mathcal{P}}) = 0 \quad \text{对 } p > 0.
\]
证明于是照抄 \cite[第 12 号 引理 5]{exp12-10} 完成，上同调群换成上同调层.

2) $f$ 紧合的情形. 用 EGA III 3.1.2 归结到射影情形. 设 $K$ 是使 $\theta_p$ 对每个 $p \geqslant 0$ 是同构的凝聚 $\mathcal{O}_X$-模范畴. 只需证明：对每个正合序列 $0 \to F' \to F \to F'' \to 0$ (其中两项属于 $K$)，第三项也属于 $K$，$K$ 的对象的直和因子属于 $K$，且对 $X$ 的每个点 $x$，可以找到 $K$ 的对象 $F$ 使 $F_x \neq 0$.

第一个条件由五引理应用于下列行正合的交换图式得出：
\[
\begin{tikzcd}[row sep=large, column sep=large]
(R^p f_* F')^{\mathrm{an}} \arrow[r] \arrow[d] & (R^p f_* F)^{\mathrm{an}} \arrow[r] \arrow[d] & (R^p f_* F'')^{\mathrm{an}} \arrow[r] \arrow[d] & (R^{p+1} f_* F')^{\mathrm{an}} \\
R^p f_*^{\mathrm{an}} F'^{\mathrm{an}} \arrow[r] & R^p f_*^{\mathrm{an}} F^{\mathrm{an}} \arrow[r] & R^p f_*^{\mathrm{an}} F''^{\mathrm{an}} \arrow[r] & R^{p+1} f_*^{\mathrm{an}} F'^{\mathrm{an}},
\end{tikzcd}
\]
类似验证第二个条件.

为验证第三个条件，可限于 $X$ 是不可约概形、生成点为 $x$ 的情形. 可从一开始假设 $Y$ 诺特. 由 Chow 引理 (EGA II 5.6.1)，可以找到 $Y$-射影概形 $X'$ 和射影满射态射 $g \colon X' \to X$. 另一方面存在整数 $n$，使对每个 $p > 0$ 有 $R^p g_*(\mathcal{O}_{X'}(n)) = 0$，且典范态射 $g^{*} g_*(\mathcal{O}_{X'}(n)) \to \mathcal{O}_{X'}(n)$ 满射 (EGA III 2.2.1). 令 $F = g_*(\mathcal{O}_{X'}(n))$，层 $F$ 回答这个问题. 确实有 $F_x \neq 0$；而且 Leray 谱序列
\[
R^p f_*(R^q g_*(\mathcal{O}_{X'}(n))) \Longrightarrow R^{p+q} (f \cdot g)_* (\mathcal{O}_{X'}(n))
\]
退化，有同构
\[
R^p f_* F \xrightarrow{\sim} R^p (f \cdot g)_* (\mathcal{O}_{X'}(n)).
\]
如同代数情形有典范同构
\[
R^p f_*^{\mathrm{an}} F^{\mathrm{an}} \xrightarrow{\sim} R^p (f \cdot g)^{\mathrm{an}}_* (\mathcal{O}_{X'}(n)^{\mathrm{an}}),
\]
且图式
\[
\begin{tikzcd}[row sep=large, column sep=large]
(R^p f_* F)^{\mathrm{an}} \arrow[r,"{\sim}"] \arrow[d,"{\theta_p}"] & (R^p (f \cdot g)_* (\mathcal{O}_{X'}(n)))^{\mathrm{an}} \arrow[d,"{\psi_p}"] \\
R^p f_*^{\mathrm{an}} F^{\mathrm{an}} \arrow[r,"{\sim}"] & R^p (f \cdot g)^{\mathrm{an}}_* (\mathcal{O}_{X'}(n)^{\mathrm{an}})
\end{tikzcd}
\]
交换. 由 1) $\psi_p$ 是同构；故 $\theta_p$ 也同构，这完成证明.

\begin{corollary}\label{cor:XII.4.3}
设 $X$ 是 $\mathbb{C}$ 上紧合概形，$F$ 是凝聚 $\mathcal{O}_X$-模. 则对每个整数 $p \geqslant 0$，典范态射
\[
H^p(X, F) \longrightarrow H^p(X^{\mathrm{an}}, F^{\mathrm{an}})
\]
是同构.
\end{corollary}

\begin{theorem}\label{thm:XII.4.4}
设 $X$ 是 $\mathbb{C}$ 上紧合概形. 把凝聚 $\mathcal{O}_X$-模 $F$ 配给它在 $X^{\mathrm{an}}$ 上的逆像 $F^{\mathrm{an}}$ 的函子是范畴等价.
\end{theorem}

1) 该函子忠实全. 设 $F$ 和 $G$ 是两个凝聚 $\mathcal{O}_X$-模. 典范态射
\[
\mathrm{Hom}_{\mathcal{O}_X}(F, G) \longrightarrow \mathrm{Hom}_{\mathcal{O}_{X^{\mathrm{an}}}}(F^{\mathrm{an}}, G^{\mathrm{an}})
\]
等同于典范态射
\[
H^0(X, \mathcal{H}om_{\mathcal{O}_X}(F, G)) \longrightarrow H^0(X^{\mathrm{an}}, \mathcal{H}om_{\mathcal{O}_X}(F, G))
\]
(EGA III$_0$ 6.7.6). 由于 $\mathcal{H}om_{\mathcal{O}_X}(F, G)$ 凝聚，由 \ref{cor:XII.4.3} 该态射双射.

2) 该函子本质满射. 当 $X$ 射影时断言由 \cite[第 12 号 定理 3]{exp12-10} 得出. 一般情形用 Chow 引理 (EGA II 5.6.1) 归结到前述. 设 $X'$ 是 $\mathbb{C}$-射影概形，$f \colon X' \to X$ 是射影满射态射，$U$ 是 $X$ 的稠密开集，使 $f$ 诱导同构 $f^{-1}(U) \simeq U$. 对 $X$ 作诺特递归；可假设：对 $X^{\mathrm{an}}$ 上每个可以找到 $X$ 的与 $X$ 不同的闭部分 $Y$ (满足 $Y^{\mathrm{an}} \supset \mathrm{Supp} \mathcal{G}$) 的凝聚层 $\mathcal{G}$，存在 $X$ 上的凝聚层 $G$ 使有同构 $G^{\mathrm{an}} \simeq \mathcal{G}$.

设 $F$ 是 $\mathcal{O}_{X^{\mathrm{an}}}$ 上的凝聚模层，$K$ 和 $L$ 是由序列
\[
0 \longrightarrow K \longrightarrow F \longrightarrow f_*^{\mathrm{an}} f^{\mathrm{an}*} F \longrightarrow L \longrightarrow 0
\]
正合定义的凝聚层. 由于 $X'$ 射影，存在凝聚 $\mathcal{O}_{X'}$-模 $F'$ 使 $F'^{\mathrm{an}} \simeq f^{\mathrm{an}*} F$；由 \ref{thm:XII.4.2} 得出同构 $(f_* F')^{\mathrm{an}} \simeq f_*^{\mathrm{an}} f^{\mathrm{an}*} F$. 由于 $K|_{U^{\mathrm{an}}}$ 和 $L|_{U^{\mathrm{an}}}$ 为零，存在凝聚 $\mathcal{O}_X$-模 $K$ 和 $L$ 使有同构 $K^{\mathrm{an}} \simeq K$，$L^{\mathrm{an}} \simeq L$. 由 1) 态射 $f_*^{\mathrm{an}} f^{\mathrm{an}*} F \to L$ 来自唯一态射 $f_* F' \to L$；设 $I = \mathrm{Ker}(f_* F' \to L)$. 层 $F$ 是 $I^{\mathrm{an}}$ 被 $K^{\mathrm{an}}$ 的扩张，只需看这个扩张来自 $I$ 被 $K$ 的扩张的逆像. 故只需证明典范态射
\[
(\ast) \quad \mathrm{Ext}^q_{\mathcal{O}_X}(I, K)^{\mathrm{an}} \longrightarrow \mathrm{Ext}^q_{\mathcal{O}_{X^{\mathrm{an}}}}(I^{\mathrm{an}}, K^{\mathrm{an}}) \quad q \neq 1
\]
双射. 有同构 $\mathrm{Ext}^q_{\mathcal{O}_X}(I, K)^{\mathrm{an}} \xrightarrow{\sim} \mathrm{Ext}^q_{\mathcal{O}_{X^{\mathrm{an}}}}(I^{\mathrm{an}}, K^{\mathrm{an}})$ 对每个整数 $q \geqslant 0$ 成立 (EGA III$_0$ 12.3.5)，以及谱序列的态射
\[
\begin{tikzcd}[row sep=large, column sep=large]
H^p(X, \mathcal{E}xt^q_{\mathcal{O}_X}(I, K)) \arrow[r, Rightarrow] \arrow[d] & \mathrm{Ext}^{p+q}_{\mathcal{O}_X}(I, K) \arrow[d] \\
H^p(X^{\mathrm{an}}, \mathcal{E}xt^q_{\mathcal{O}_{X^{\mathrm{an}}}}(I^{\mathrm{an}}, K^{\mathrm{an}})) \arrow[r, Rightarrow] & \mathrm{Ext}^{p+q}_{\mathcal{O}_{X^{\mathrm{an}}}}(I^{\mathrm{an}}, K^{\mathrm{an}}).
\end{tikzcd}
\]
该态射是同构，因为由 \ref{cor:XII.4.3} 在 $E^{pq}_2$ 项上如此，这证明 $(\ast)$ 双射.

\begin{corollary}\label{cor:XII.4.5}
把紧合 $\mathbb{C}$-概形 $X$ 配给 $X^{\mathrm{an}}$ 的函子忠实全.
\end{corollary}

需证明：若 $X$ 和 $Y$ 是两个紧合 $\mathbb{C}$-概形，典范映射
\[
\mathrm{Hom}_{\mathbb{C}}(X, Y) \longrightarrow \mathrm{Hom}(X^{\mathrm{an}}, Y^{\mathrm{an}})
\]
双射. 给出 $X$ 到 $Y$ 的态射 (resp. $X^{\mathrm{an}}$ 到 $Y^{\mathrm{an}}$ 的态射) 等价于给出它的图，i.e. $X \times Y$ 的闭子概形 $Z$ (resp. $X^{\mathrm{an}} \times Y^{\mathrm{an}}$ 的闭解析子空间 $Z$)，使第一个投影 $X \times Y \to X$ (resp. $X^{\mathrm{an}} \times Y^{\mathrm{an}} \to X^{\mathrm{an}}$) 在 $Z$ 上的限制是同构. 由于给出 $X \times Y$ 的闭子概形 (resp. $X^{\mathrm{an}} \times Y^{\mathrm{an}}$ 的闭解析子空间) 等价于给出 $\mathcal{O}_{X \times Y}$ (resp. $\mathcal{O}_{X^{\mathrm{an}} \times Y^{\mathrm{an}}}$) 上的凝聚理想层，该推论由 \ref{thm:XII.4.4} 得出.

\begin{corollary}\label{cor:XII.4.6}
设 $X$ 是 $\mathbb{C}$ 上紧合概形. 把 $X$ 上方的有限概形 (resp. 有限平展概形) $X'$ 配给 $X'^{\mathrm{an}}$ 的函子，是 $X$ 上方有限概形 (resp. 有限平展概形) 范畴到 $X^{\mathrm{an}}$ 上方有限 (resp. 有限平展) 解析空间范畴的等价.
\end{corollary}

确实，给出有限态射 $X' \to X$ (resp. $X'^{\mathrm{an}} \to X^{\mathrm{an}}$) 等价于给出 $\mathcal{O}_X$ (resp. $\mathcal{O}_{X^{\mathrm{an}}}$) 上的凝聚代数层 \cite[第 19 号 §5 定理 2]{exp12-4}. 该推论在非 resp. 情形由 \ref{thm:XII.4.4} 得出，resp. 情形考虑到 \ref{prop:XII.3.1} (iii) 由此推出.

\section{平展覆盖的比较定理}

\refstepcounter{theorem}\addtocounter{theorem}{-1}\noindent\textbf{\thesection.0.}\quad 明确\emph{解析空间的有限覆盖}概念. 若 $X$ 是解析空间，称 $X$ 上方有限的解析空间 $X'$ 为 $X$ 的有限覆盖\index[term]{jiexikongjiandeyouxianfugai@解析空间的有限覆盖}，若 $X'$ 的每个不可约分量支配 $X$ 的某个不可约分量.

\begin{theorem}\label{thm:XII.5.1}
[Riemann 存在定理]\index[term]{Riemann cunzaidingli@Riemann 存在定理} 设 $X$ 是 $\mathbb{C}$ 上局部有限型概形，$X^{\mathrm{an}}$ 是 $X$ 的关联解析空间. 把 $X$ 的每个有限平展覆盖 $X'$ 配给 $X'^{\mathrm{an}}$ 的函子 $\Psi$，是 $X$ 的有限平展覆盖范畴到 $X^{\mathrm{an}}$ 的有限平展覆盖范畴的等价.
\end{theorem}

1) 函子 $\Psi$ 忠实全. 设 $X'$ 和 $X''$ 是 $X$ 的两个有限平展覆盖，证明典范映射
\[
(\ast) \quad \mathrm{Hom}_X(X', X'') \longrightarrow \mathrm{Hom}_{X^{\mathrm{an}}}(X'^{\mathrm{an}}, X''^{\mathrm{an}})
\]
双射. 可假设 $X'$ 连通. 给出 $X'$ 到 $X''$ 的 $X$-态射等价于给出 $X' \times_X X''$ 的连通分量 $X_i$，使第一个投影诱导的态射 $X_i \to X'$ 是同构. 由于 $X' \times_X X''$ 的连通分量与 $X'^{\mathrm{an}} \times_{X^{\mathrm{an}}} X''^{\mathrm{an}}$ 的连通分量一一对应 (\cref{cor:XII.2.6})，且态射 $X_i \to X'$ 是同构当且仅当 $X_i^{\mathrm{an}} \to X'^{\mathrm{an}}$ 如此，这证明 $(\ast)$ 双射.

2) 函子 $\Psi$ 本质满射. 设 $X'$ 是 $X^{\mathrm{an}}$ 的有限平展覆盖，证明存在 $X$ 的平展覆盖 $X'$，使有同构 $X'^{\mathrm{an}} \xrightarrow{\sim} X'$. 考虑到 1) 问题局部于 $X$，可假设 $X$ 仿射.

a) 归结到 $X$ 正规的情形. 可假设 $X$ 既约. 假设定理对 $X_{\mathrm{r\'ed}}$ 已证. 把 $X$ 的有限平展覆盖 $X'$ 配给 $X_{\mathrm{r\'ed}}$ 的有限平展覆盖 $X_{\mathrm{r\'ed}}'^{\mathrm{an}}$ 的函子此时是等价. 由于它由 $\Psi$ 与把 $X^{\mathrm{an}}$ 的有限平展覆盖配给它在 $X_{\mathrm{r\'ed}}^{\mathrm{an}}$ 上的逆像的函子 $\Theta$ 复合得到，且 $\Theta$ 忠实全，这证明 $\Psi$ 是范畴等价.

可假设 $X$ 正规. 设 $\widetilde{X}$ 是 $X$ 的正规化，$p \colon \widetilde{X} \to X$ 是典范态射. 由于 $p$ 有限，$p$ 对平展覆盖范畴是有效下降态射 (\cref{thm:IX.4.7}). 假设定理对 $\widetilde{X}$ 已证，若令 $\widetilde{X}' = X' \times_{X^{\mathrm{an}}} \widetilde{X}^{\mathrm{an}}$，存在 $\widetilde{X}$ 的平展覆盖 $\widetilde{X}'$ 和同构 $\widetilde{X}'^{\mathrm{an}} \simeq \widetilde{X}'$. 由 1) 得出 $\widetilde{X}'$ 上的自然下降数据提升为相对 $\widetilde{X} \to X$ 的 $\widetilde{X}'$ 上的下降数据；这证明存在 $X$ 的平展覆盖 $X'$，使有同构 $i \colon X'^{\mathrm{an}} \times_{X^{\mathrm{an}}} \widetilde{X}^{\mathrm{an}} \simeq \widetilde{X}'$，它在 $\widetilde{X}^{\mathrm{an}} \times_{X^{\mathrm{an}}} \widetilde{X}^{\mathrm{an}}$ 的两个投影下的逆像相同. 由 IX 3.2 (其证明在解析情形有效)，态射 $\widetilde{X}^{\mathrm{an}} \to X^{\mathrm{an}}$ 对平展覆盖范畴是下降态射，从而 $i$ 来自同构 $X'^{\mathrm{an}} \simeq X'$.

b) 归结到 $X$ 正则的情形. 设 $U$ 是 $X$ 的正则点的开集，$i \colon U \to X$，$i^{\mathrm{an}} \colon U^{\mathrm{an}} \to X^{\mathrm{an}}$ 是典范态射；由于 $X$ 正规，有 $\mathrm{codim}(X - U, X) \geqslant 2$. 假设存在 $U$ 的平展覆盖 $U'$ 使 $U'^{\mathrm{an}} \simeq X'|_{U^{\mathrm{an}}}$，证明此时 $U'$ 延拓成 $X$ 的平展覆盖 $X'$，使 $X'^{\mathrm{an}} \simeq X'$. 只需看 $U'$ 延拓成 $X$ 的平展覆盖 $X'$；确实将有同构 $X'^{\mathrm{an}}|_{U^{\mathrm{an}}} \simeq X'|_{U^{\mathrm{an}}}$；但若 $F$ 和 $G$ 是 $\mathcal{O}_{X^{\mathrm{an}}}$ 上分别定义 $X'$ 和 $X'^{\mathrm{an}}$ 的凝聚代数层，$X$ 正规且 $\mathrm{codim}(X - U, X) \geqslant 2$ 的事实蕴含典范态射
\[
F \longrightarrow i_*^{\mathrm{an}}(F|_{U^{\mathrm{an}}}), \quad G \longrightarrow i_*^{\mathrm{an}}(G|_{U^{\mathrm{an}}})
\]
是同构 \cite[第 3 号 命题 4]{exp12-11}. 由此得出 $F$ 和 $G$ 从而 $X'^{\mathrm{an}}$ 和 $X'$ 同构.

设 $\phi \colon X^{\mathrm{an}} \to X$ 是典范态射. 由于把 $U'$ 延拓到 $X$ 的问题局部于 $X$，只需证明：对 $X^{\mathrm{an}} - U^{\mathrm{an}}$ 的每个点 $y$，$U_{\phi(y)} = U \times_X \mathrm{Spec} \mathcal{O}_{X,\phi(y)}$ 的平展覆盖 $U'_{\phi(y)} = U' \times_X \mathrm{Spec} \mathcal{O}_{X,\phi(y)}$ 延拓到 $\mathrm{Spec} \mathcal{O}_{X,\phi(y)}$. 设 $H$ 是定义 $U'$ 的 $\mathcal{O}_U$-凝聚代数. 典范态射
\[
\alpha \colon (i_* H)^{\mathrm{an}} \longrightarrow i_*^{\mathrm{an}}(H^{\mathrm{an}}) = F
\]
定义 $\mathrm{Spec} \mathcal{O}_{X^{\mathrm{an}},y}$ 上模层的态射：
\[
\alpha_y \colon (i_* H)^{\mathrm{an}}_y \longrightarrow F_y,
\]
它在 $U_y = U_{\phi(y)} \times_{\mathrm{Spec} \mathcal{O}_{X,\phi(y)}} \mathrm{Spec} \mathcal{O}_{X^{\mathrm{an}},y}$ 上的限制是同构. 但这证明 $H|_{U_y}$ 平凡，故 $U'_{\phi(y)}$ 延拓到 $\mathrm{Spec} \mathcal{O}_{X,\phi(y)}$.

c) $X$ 仿射正则的情形. 设
\[
j \colon X \longrightarrow P
\]
是 $X$ 的紧化，其中 $P$ 是射影 $\mathbb{C}$-概形，$j$ 是支配开浸入. 由奇点消解定理\index[term]{qidianxiaojie@奇点消解} \cite{exp12-8}，可以找到正则概形 $R$、射影态射 $r \colon R \to P$，使 $r$ 诱导同构 $r^{-1}(X) \simeq X$，且 $r^{-1}(X)$ 是 $R$ 中正规交叉除子的余集. 设
\[
k \colon X \longrightarrow R
\]
是典范浸入. 将证明存在 $R^{\mathrm{an}}$ 的延拓平展覆盖 $X'^{\mathrm{an}}$ 的有限正规覆盖 (5.0) $R'$. 由下面\cref{prop:XII.5.3}，这样的覆盖唯一；延拓 $X'^{\mathrm{an}}$ 的问题于是局部于 $R^{\mathrm{an}}$ 在 $R^{\mathrm{an}} - X^{\mathrm{an}}$ 附近的邻域. 而 $R^{\mathrm{an}} - X^{\mathrm{an}}$ 的每个点有同构于仿射空间 $E^n$ 的球的开邻域 $V$，使 $V - V \cap X^{\mathrm{an}}$ 由前 $p$ 个坐标函数 $z_1, \ldots, z_p$ (其中 $0 \leqslant p \leqslant n$) 的消没定义. $U = V \cap X^{\mathrm{an}}$ 的基本群同构于 $\mathbb{Z}^p$，$U$ 的每个平展覆盖是形如
\[
U'' = U[T_1, \ldots, T_p]/(T_1^{n_1} - z_1, \ldots, T_p^{n_p} - z_p)
\]
的覆盖除以 $U''$ 的 伽罗瓦群 $\mathbb{Z}/n_1\mathbb{Z} \times \cdots \times \mathbb{Z}/n_p\mathbb{Z}$ 的子群 $H$ 的商，其中诸 $n_i$ 是 $> 0$ 的整数. 而 $U''$ 延拓成正则覆盖
\[
V'' = V[T_1, \ldots, T_p]/(T_1^{n_1} - z_1, \ldots, T_p^{n_p} - z_p),
\]
$H$ 作用在其上，$V''$ 除以 $H$ 的商是所求延拓.

证明于是由 \ref{cor:XII.4.6} 完成. 覆盖 $R'$ 来自 $R$ 的有限覆盖 $R'$；$R'$ 在 $X$ 上的限制是 $X$ 的覆盖 $X'$，使 $X'^{\mathrm{an}} \simeq X'$，且由 \ref{prop:XII.3.1} (iii) $X'$ 是 $X$ 的平展覆盖.

\begin{corollary}\label{cor:XII.5.2}
设 $X$ 是 $\mathbb{C}$ 上连通局部有限型概形，$\phi \colon X^{\mathrm{an}} \to X$ 是典范态射，$x$ 是 $X^{\mathrm{an}}$ 的点. 设 $\pi_1(X^{\mathrm{an}}, x)$ 是拓扑空间 $X^{\mathrm{an}}$ 在点 $x$ 的基本群，$\pi_1(X, \phi(x))$ 是概形 $X$ 在点 $\phi(x)$ 的基本群 (V 7). 则 $\pi_1(X, \phi(x))$ 典范同构于 $\pi_1(X^{\mathrm{an}}, x)$ 关于有限指数子群拓扑的完备化.
\end{corollary}

设 $\mathcal{C}$ 是 $X^{\mathrm{an}}$ 的有限平展覆盖范畴，$F$ 是从 $\mathcal{C}$ 到 $\mathrm{Ens}$ 的、把 $X^{\mathrm{an}}$ 的每个有限平展覆盖 $X'$ 配给 $X'$ 的位于 $x$ 上方的点集的函子，设 $\widehat{\pi_1}(X^{\mathrm{an}}, x)$ 是按 V 4 所述与 $\mathcal{C}$ 和 $F$ 关联的 pro-有限群. 由于 $X^{\mathrm{an}}$ 的每个有限平展覆盖是万有覆盖除以有限指数子群的商，$\widehat{\pi_1}(X^{\mathrm{an}}, x)$ 正是 $\pi_1(X^{\mathrm{an}}, x)$ 关于有限指数子群拓扑的完备化. 该推论于是由 \ref{thm:XII.5.1} 和 V 6.10 得出.

\begin{proposition}\label{prop:XII.5.3}
设 $X$ 是正规解析空间，$Y$ 是 $U = X - Y$ 在 $X$ 中稠密的闭解析子集. 则把 $X$ 的每个有限正规覆盖 (5.0) $X'$ 配给它在 $U$ 上的限制的函子忠实全.
\end{proposition}

设 $X'$ 和 $X''$ 是 $X$ 的两个有限正规覆盖. 需证明典范映射
\[
\mathrm{Hom}_X(X', X'') \longrightarrow \mathrm{Hom}_U(X'|_U, X''|_U)
\]
双射. 设 $u$，$v$ 是 $X'$ 到 $X''$ 的两个在 $U$ 上的限制相同的 $X$-态射，证明 $u = v$. 态射 $u$ 和 $v$ 在稠密开集 $U \times_X X'$ 上重合，故在底拓扑空间上重合. 由 \cite[第 19 号 §4 推论 5]{exp12-4} 这证明 $u = v$.

现在设 $u$ 是 $X'|_U$ 到 $X''|_U$ 的 $U$-态射，证明它延拓到整个 $X'$. 可假设 $X'$ 正则. 确实，$X'$ 正规，可以找到余集为余维数 $\geqslant 2$ 的解析部分的 $X$ 的开集 $V$，使 $X' \times_X V = V'$ 正则. 设 $V'' = X'' \times_X V$，假设命题对 $V$ 已证. 考虑交换图式
\[
\begin{tikzcd}[row sep=large, column sep=large]
V' \arrow[r,"{i'}"] & X' \arrow[d,"{f'}"] \\
V'' \arrow[r,"{i''}"'] & X'' \arrow[d,"{f''}"] \\
V \arrow[r,"{i}"'] & X
\end{tikzcd}
\]
与 $u$ 关联 $\mathcal{O}_V$-代数的态射 $g_*'' \mathcal{O}_{V''} \to g_*' \mathcal{O}_{V'}$，由此得到态射
\[
i_* g_*'' \mathcal{O}_{V''} \longrightarrow i_* g_*' \mathcal{O}_{V'}.
\]
考虑到同构 $i'_* \mathcal{O}_{V'} \simeq \mathcal{O}_{X'}$，$i''_* \mathcal{O}_{V''} \simeq \mathcal{O}_{X''}$ \cite[第 3 号 命题 4]{exp12-11}，由此得到 $\mathcal{O}_X$-代数的态射
\[
f_*'' \mathcal{O}_{X''} \longrightarrow f_*' \mathcal{O}_{X'},
\]
从而得到所求态射 $X' \to X''$.

此后假设 $X'$ 正则. 设 $U' = U \times_X X'$，$Y' = X' - U'$. 把 $Y'$ 看作 $X'$ 的既约解析子空间；若 $Y'_1$ 是 $Y'$ 的奇异闭集，有 $\dim Y'_1 < \dim Y'$ \cite[第 20 号 D. 定理 3]{exp12-4}. 对 $\dim Y'$ 作递归，可假设 $Y'$ 光滑. 由于只需把 $u$ 延拓到 $Y'$ 每个点的开邻域，由隐函数定理可假设 $X'$ 是仿射空间 $E^n$ 的球，$Y'$ 由前 $p$ 个坐标函数 $z_1, \ldots, z_p$ (其中 $0 \leqslant p \leqslant n$) 的消没定义的闭集.

把 $u$ 关联 $p \colon X' \times_X X'' \to X'$ 在 $U'$ 上方的截面 $s$. 可以限制 $X'$，使 $A = p_*(\mathcal{O}_{X' \times_X X''})$ 由 $\Gamma(X', A)$ 的元素 $x_1, \ldots, x_q$ 生成. 设 $u_1, \ldots, u_q \in \Gamma(U', \mathcal{O}_{X'})$ 是 $s$ 对 $x_1|_{U'}, \ldots, x_q|_{U'}$ 的像. 说 $s$ 延拓到 $X'$ 等价于说诸 $u_1, \ldots, u_q$ 延拓成 $\Gamma(X', \mathcal{O}_{X'})$ 的截面. 但 $f$ 有限，每个 $u_i$ 是 $z_1, \ldots, z_p$ 的 Laurent 级数 (系数是 $z_{p+1}, \ldots, z_n$ 的幂级数)，满足整依赖关系. 由此得出 $u_i$ 有界故是 $z_1, \ldots, z_n$ 的幂级数，从而延拓到 $X'$.

可以问 \ref{prop:XII.5.3} 引入的函子是否是范畴等价. 借助 Grauert-Remmert 定理 \cite{exp12-6} 可以回答这个问题，下面用奇点消解给出它的证明. 在 \cite{exp12-8} 之前也可以用 Grauert-Remmert 定理证明 \ref{thm:XII.5.1}；正是这样做的.

\begin{theorem}\label{thm:XII.5.4}
[Grauert-Remmert 定理]\index[term]{Grauert-Remmert dingli@Grauert-Remmert 定理} 设 $X$ 是正规解析空间，$Y$ 是 $U = X - Y$ 在 $X$ 中稠密的闭解析子集. 设 $U'$ 是 $U$ 的有限正规覆盖；假设存在 $X$ 的稀疏闭解析部分 $S$，使 $U'$ 在 $U - U \cap S$ 上的限制平展. 则存在 $X$ 的延拓 $U'$ 的有限正规覆盖 $X'$，且 $X'$ 在同构意义下唯一.
\end{theorem}

唯一性由 \ref{prop:XII.5.3} 得出. 延拓 $U'$ 的问题于是局部于 $X$. 可假设 $U$ 正则且 $U'$ 在 $U$ 上平展. 确实 $U$ 的正则点集是 $X$ 中稠密、余集为解析部分的开集 $V$ \cite[第 20 号 D 定理 2]{exp12-4}，只需把 $U$ 换成开集 $V - V \cap S$.

设 $y$ 是 $X - U$ 的点，证明可以把 $U'$ 延拓到 $y$ 的邻域. 把 $X$ 限制到 $y$ 的开邻域，由奇点消解定理 \cite{exp12-8} 可以找到正则解析空间 $X_1$、射影态射 $f \colon X_1 \to X$ (在 $U$ 上的限制诱导同构 $U_1 = f^{-1}(U) \simeq U$)，使 $U_1$ 是 $X_1$ 中正规交叉除子的余集. 证明 $U'$ 延拓成 $X_1$ 的有限正规覆盖. 由于问题局部于 $X_1$，可假设 $X_1$ 是仿射空间 $E^n$ 的球，$X_1 - U_1$ 由前 $p$ 个坐标函数 $z_1, \ldots, z_p$ (其中 $0 \leqslant p \leqslant n$) 的消没定义. $U_1$ 的平展覆盖 $U'$ 是形如
\[
U_2 = U_1[T_1, \ldots, T_p]/(T_1^{n_1} - z_1, \ldots, T_p^{n_p} - z_p)
\]
的覆盖除以 $U_2$ 的 伽罗瓦群的子群 $H$ 的商. 覆盖 $U_2$ 延拓成覆盖
\[
X_2 = X_1[T_1, \ldots, T_p]/(T_1^{n_1} - z_1, \ldots, T_p^{n_p} - z_p)
\]
$H$ 作用在其上，$X_2/H$ 把 $U'$ 延拓到 $X_1$.

记 $X'_1$ 为 $X_1$ 的延拓 $U'$ 的有限正规覆盖，$F_1$ 为 $F_1$ 定义的 $\mathcal{O}_{X_1}$-凝聚代数. 由 Grauert-Remmert 有限性定理 \cite[第 15 号 \cref{thm:XII.1.1}]{exp12-4}，$f_* F_1$ 是 $\mathcal{O}_X$-凝聚代数. 它对应 $X$ 的有限覆盖 $X'$，而且 $X'$ 正规，因为 $X'_1$ 正规，且 $X'$ 是所求 $U'$ 的延拓.

\begin{remark}\label{rem:XII.5.5}
在断言 \ref{thm:XII.5.4} 中，不能删除关于 $U' \to U$ 不平展的点集的条件. 例如取 $X$ 为复平面的单位圆盘，$U$ 为 $X$ 中原点的余集，$U' = U[T]/(T^2 - \sin 1/z)$，其中 $z$ 是 $X$ 上的坐标函数. 则 $U'$ 是 $U$ 的不延拓到 $X$ 的有限正规覆盖. 确实假设 $U'$ 延拓成 $X$ 的有限覆盖 $X'$；则 $X' \to X$ 不平展的点集是含所有满足 $\sin 1/z = 0$ 的点 $z$ 的闭解析集，这荒谬.

然而，当 $\mathrm{codim}(X - U, X) \geqslant 2$ 时可以删除关于 $U' \to U$ 的奇异点集的条件. 可假设 $U$ 正则. $U' \to U$ 不平展的点集是 $U$ 的除子，由 Remmert-Stein 定理 \cite[定理 3]{exp12-9} 得出它是 $X$ 的除子在 $U$ 上的迹. 此时，若 $A$ 是使 $U' = \mathrm{Spec}^{\mathrm{an}}(A)$ 的 $\mathcal{O}_U$-凝聚代数，$i \colon U \to X$ 是典范态射，只需取 $X' = \mathrm{Spec}^{\mathrm{an}}(i_* A)$；确实知道 $i_* A$ 凝聚 \cite[第 1 号 定理 1]{exp12-11}.

\begin{remark}[2003 年 (MR) 增补]
存在有限表现的非平凡群 $G$，没有与 $G$ 不同的有限指数子群，例如 Higman 群 (参见 J.-P. Serre，Arbres et amalgames，Ast\'erisque N\textdegree 46，命题 6，第一章 §1). 于是，若这样的群实现为 $\mathbb{C}$ 上概形 $V$ (说光滑且射影) 的拓扑基本群，拓扑空间 $V_{\mathrm{top}}$ (即 $V$ 的底拓扑空间) 不单连通，但 $V$ 代数单连通. 目前不知道这样的 $V$. 但指出 D. Toledo 构造了 $\mathbb{C}$ 上光滑射影概形 $V$，它的拓扑基本群不关于有限指数子群拓扑分离 (自然态射 $\pi_1(V_{\mathrm{top}}) \to \pi_1(V_{\mathrm{alg}})$ 不单射). [D. Toledo，Projective varieties with non-residually finite fundamental group，Publ. Math. IHES 77 (1993)，p. 103--119.]
\end{remark}
\end{remark}

\section{文献}\printbibliography[heading=none]


\end{refsection}
\begin{refsection}% =====================================================================
% 第十三讲 (Exposé XIII) 集合层与非交换群层的上同调真性
% (Propreté cohomologique des faisceaux d'ensembles et des faisceaux
%  de groupes non commutatifs)
% M. Raynaud
% 原书页码: pp. 259-322
% =====================================================================
\chapter{集合层与非交换群层的上同调真性}

\noindent Mme M. Raynaud\footnote{按 A. Grothendieck 的未发表笔记.}

本讲旨在用平展上同调推广 IX 和 X 的某些结果. 还说明如何把 SGA 5 II 中对这类层仍有意义的结果推广到未必交换的群层. 假定已知 SGA 4 中阐述的平展上同调概念.

主要结果 (\cref{thm:XIII.2.4}) 给出非紧合态射 $f \colon U \to S$ 的重要例子，它在"维数 $\leqslant 1$ 时上同调真"，i.e. 对 $U$ 上某些 (平展拓扑意义的) 群层 $F$，$f_* F$ 和 $R^1 f_* F$ 的构造与任何基变换 $S' \to S$ 交换. 这个性质确实被 $S$ 上紧合概形 $X$ 的、相对 $S$ 为正规交叉除子 $D$ 的余集的开集 $U$ 满足，至少当对 $F$ 要求它常有限、阶与 $S$ 的剩余特征互素时. 若不再假设 $F$ 的阶与 $S$ 的剩余特征互素，有类似结果，只要把 $R^1 f_* F$ 换成子层 $R^1_t f_* F$ (只考虑 $F$ 下"相对 $S$ 在 $X$ 上温和分歧"的挠子). 特别这可以用来证明可分闭域上紧合光滑代数曲线去掉有限个闭点后所得曲线的\emph{温和分歧基本群}拓扑有限型 (\cref{cor:XIII.2.12}).

第 4 号致力于同伦正合序列和 Künneth 公式.

最后附录给出 X.3.6 中证明的 Abhyankar 引理的有用变体.

\setcounter{section}{-1}

\section{叠理论的回顾}

下面将用 \cite{exp13-1} 和 \cite{exp13-2} 中阐述的叠理论. 限于概形的平展位点. 给定概形 $X$，用 $X_{\mathrm{ét}}$ 记 $X$ 的平展位点. 回忆 $X$ 上的叠 $F$ 是 $X_{\mathrm{ét}}$ 上的纤维范畴\index[term]{die@叠}，使对每个平展于 $X$ 上的概形 $X'$ 和纤维 $F_{X'}$ 的每对对象 $x$，$y$，预层 $\mathrm{Hom}_{X'}(x, y)$ 是层，且使对每个平展满射态射 $X'' \to X'$，$F_{X''}$ 的配备相对 $X'' \to X'$ 的下降数据的每个对象是 $F_{X'}$ 的某个对象的逆像.

用 $F(X')$ 记 $F/X'$ 的笛卡儿截面的范畴. 更一般地，若 $\mathbf{Sch}_X$ 是配备平展拓扑的 $X$ 上方概形范畴，叠 $F$ 可以延拓成 $\mathbf{Sch}_X$ 上的叠 $F$，对每个态射 $f \colon X' \to X$，仍用 $F(X')$ 记该叠 $F$ 在 $X'$ 上方的笛卡儿截面范畴.

\emph{叠丛}是满足下列条件的叠\index[term]{diecong@叠丛}：对每个平展于 $X$ 上的概形 $X'$ 和纤维 $F_{X'}$ 的每对对象 $x$，$y$，$x$ 到 $y$ 的每个态射是同构，$x$ 和 $y$ 局部同构，且 $X_{\mathrm{ét}}$ 中使 $F_{X'}$ 非空的对象 $X'$ 的集是 $X_{\mathrm{ét}}$ 的加细. 例如群层下的挠子叠是叠丛，而且它有笛卡儿截面. 反之有截面的叠丛，i.e. 存在 $F_X$ 的对象 $x$，等价于自同构群层 $\mathrm{Aut}_X(x)$ 下的挠子叠.

有叠 $F$ 的子叠丛和\emph{极大子叠丛}的显然概念. 给定 $F(X)$ 的笛卡儿截面 $x$，存在唯一的 $F$ 的极大子叠丛\index[term]{jidazidiecong@极大子叠丛} $G_x$，使 $x$ 经 $G_x$ 分解. 称 $G_x$ 为 $x$ 生成的子叠丛；按定义它是\emph{平凡叠丛}\index[term]{pingfandiecong@平凡叠丛}. 由
\[
S_F(X') = \{F|_{X'}\index[notations]{$S_F$} \text{ 的极大子叠丛}\}
\]
定义的预层 $S_F$ 是层，称为 $F$ 的\emph{极大子叠丛层}\index[term]{jidazidiecongceng@极大子叠丛层}. 设 $\mathcal{O}$ 是由
\[
\mathcal{O}(X') = \{F_{X'} \text{ 的对象模同构的类}\}
\]
定义的预层. 把 $F_{X'}$ 的每个对象 $x$ 配给 $x$ 生成的 $F|_{X'}$ 的极大子叠丛，得到态射
\[
\mathcal{O} \longrightarrow S_F;
\]
由 \cite[III 2.1.4]{exp13-2}，该态射使 $S_F$ 成为与 $\mathcal{O}$ 关联的层.

称叠 $F$ 可构\index[term]{kegoudie@可构叠} (resp. ind-$L$-有限\index[term]{ind-L-youxiandie@ind-$L$-有限叠}，$L$ 是素数集)，若对每个平展于 $X$ 上的概形 $X'$ 和 $F_{X'}$ 的每个对象 $x$，层 $\mathrm{Aut}_{X'}(x)$ 如此 \cite[VII 2.2.1]{exp13-2}. 称叠 1-可构\index[term]{1-kegoudie@1-可构叠}，若它可构且极大子叠丛层可构.

\section{上同调真性}

\refstepcounter{theorem}\noindent\textbf{\thesection.\arabic{theorem}.}\quad 设 $S$ 是概形，$f \colon X \to Y$ 是 $S$-概形的态射. 若 $S'$ 是 $S$-概形，考虑下列所有方形笛卡儿的图式：
\begin{equation}
\begin{tikzcd}[row sep=large, column sep=large]
X \arrow[d,"{f}"] & X' \arrow[l,"h"] \arrow[d,"{f'}"] \\
Y \arrow[d] & Y' \arrow[l,"g"] \arrow[d] \\
S & S' \arrow[l]
\end{tikzcd}
\tag{1.0.1}
\label{eq:XIII.1.0.1}
\end{equation}
若 $Y_1$ 是 $Y$ 上方的平展概形，令 $X_1 = X \times_Y Y_1$，$Y'_1 = Y' \times_Y Y_1$，考虑笛卡儿方形
\begin{equation}
\begin{tikzcd}[row sep=large, column sep=large]
X_1 \arrow[d,"{f_1}"] & X'_1 \arrow[l,"{h_1}"] \arrow[d,"{f'_1}"] \\
Y_1 & Y'_1 \arrow[l,"{g_1}"]
\end{tikzcd}
\tag{1.0.2}
\label{eq:XIII.1.0.2}
\end{equation}

\begin{definition}\label{def:XIII.1.1}
设 $F$ 是 $X$ 上的叠. 称 $(F, f)$ 相对 $S$ 在维数 $\leqslant -1$\index[term]{shangtongdiaozhen@上同调真} (resp. $\leqslant 0$，resp. $\leqslant 1$) 时上同调真，若对每个 $S$-概形 $S'$，典范函子 (由叠的逆像的万有性质以显然方式定义)：
\[
g^{*} f_* F \longrightarrow f'_* h^{*} F \quad \text{(参见 1.0.1)}
\]
忠实 (resp. 忠实全，resp. 范畴等价).
\end{definition}

若在 $S$ 上无混淆可能，特别当 $S = Y$ 时，说上同调真而不再说相对 $S$ 上同调真.

\refstepcounter{theorem}\noindent\textbf{\thesection.\arabic{theorem}.}\quad 设 $F$ 是 $X$ 上的集合层；设 $\Phi$ 是与 $F$ 关联的离散范畴叠，i.e. 纤维在平展于 $X$ 上的每个概形 $X_1$ 上方是对象集为 $F(X_1)$ 的离散范畴的叠. 称 $(F, f)$ 相对 $S$ 在维数 $\leqslant -1$ (resp. $\leqslant 0$) 时上同调真，若 $(\Phi, f)$ 相对 $S$ 在维数 $\leqslant 0$ (resp. $\leqslant 1$) 时上同调真.

典范态射
\[
(1.2.1) \quad g^{*} f_* F \longrightarrow f'_* h^{*} F
\]
经离散范畴叠的过渡给出典范态射
\[
g^{*} f_* \Phi \longrightarrow f'_* h^{*} \Phi.
\]
从而说 $(F, f)$ 相对 $S$ 在维数 $\leqslant -1$ (resp. $\leqslant 0$) 时上同调真，等价于说对每个 $S$-概形 $S'$，态射 (1.2.1) 单射 (resp. 双射).

\refstepcounter{theorem}\noindent\textbf{\thesection.\arabic{theorem}.}\quad 设 $F$ 是 $X$ 上的群层，$\Phi$ 是 $X$ 上群 $F$ 的挠子叠 \cite[II 2.3.2]{exp13-1}. 称 $(F, f)$ 相对 $S$ 在维数 $\leqslant -1$ (resp. $\leqslant 0$，resp. $\leqslant 1$) 时上同调真，若 $(\Phi, f)$ 相对 $S$ 在维数 $\leqslant -1$ (resp. $\leqslant 0$，resp. $\leqslant 1$) 时上同调真. 上同调真性条件可以如下明确.

\begin{proposition}\label{prop:XIII.1.3.1}
记号同 \eqref{eq:XIII.1.0.1} 和 \eqref{eq:XIII.1.0.2}. 设 $F$ 是 $X$ 上的群层. 用 $F'$ (resp. $F_1$，resp. $F'_1$，等等) 记 $F$ 在 $X'$ (resp. $X_1$，resp. $X'_1$，等等) 上的逆像. 则下列条件等价：
\begin{enumerate}[i)]
\item $(F, f)$ 相对 $S$ 在维数 $\leqslant -1$ (resp. $\leqslant 0$，resp. $\leqslant 1$) 时上同调真.
\item 对每个态射 $S' \to S$、每个平展于 $Y$ 上方的概形 $Y_1$ 和 $X_1$ 上群 $F_1$ 的每个挠子 $P$，若 $P F_1$ 表示 $F_1$ 被 $P$ 扭转的群 \cite[II 4.1.2.3]{exp13-1}，典范态射
\[
a_0 \colon g_{1*}(f_{1*}(P F_1)) \longrightarrow f'_{1*}(P' F'_1)
\]
单射 (resp. $a_0$ 双射且典范态射
\[
a_1 \colon g^{*}(R^1 f_* F) \longrightarrow R^1 f'_* F'
\]
单射，resp. $a_0$ 和 $a_1$ 双射).
\item[(ii bis)] 对每个态射 $S' \to S$、每个平展于 $Y$ 上方的概形 $Y_1$、$X_1$ 上群 $F_1$ 的每个挠子 $P$ 和 $P F_1$ 下的每个挠子 $R$，典范态射
\[
\alpha_0 \colon g_{1*}(f_{1*} R) \longrightarrow f'_{1*} R'
\]
单射 (resp. $\alpha_0$ 双射，resp. 态射 $\alpha_0$ 和
\[
\alpha_1 \colon g_{1*}(R^1 f_{1*}(P F_1)) \longrightarrow R^1 f'_{1*}(P' F'_1)
\]
双射).
\end{enumerate}
\end{proposition}

\begin{proof}
(i)$\Rightarrow$(ii bis). 由 \cite[II 4.2.5]{exp13-1}，群 $P F_1$ 的每个挠子 $R$ 形如 $R = Q \wedge^{F_1} P^{\circ}$，其中 $Q$ 是群 $F_1$ 的挠子，$P^{\circ}$ 是 $P$ 的逆挠子. 有 $R' \simeq Q' \wedge^{F'_1} P'^{\circ}$. 设 $\Phi$ 是 $F$ 下的挠子叠，设 $x$，$y$ (resp. $x'$，$y'$) 是纤维范畴 $(g^{*} f_* \Phi)_{Y'_1}$ (resp. $(f'_* \Phi')_{Y'_1}$) 的对应 $P$，$Q$ (resp. $P'$，$Q'$) 的对象. 有关系
\[
Q \wedge^{F_1} P^{\circ} \simeq \mathrm{Hom}_{F_1}(P, Q),
\]
由此得出典范同构
\[
\mathrm{Hom}_{Y'_1}(x, y) \simeq g_{1*} f_{1*}(Q \wedge^{F_1} P^{\circ}), \quad \mathrm{Hom}_{Y'_1}(x', y') \simeq f'_{1*}(Q' \wedge^{F'_1} P'^{\circ}).
\]
从而态射 $\alpha_0$ 等同于态射
\[
\mathrm{Hom}_{Y'_1}(x, y) \longrightarrow \mathrm{Hom}_{Y'_1}(x', y'),
\]
由此得出：若 $(F, f)$ 相对 $S$ 在维数 $\leqslant -1$ 时上同调真，则 $\alpha_0$ 单射，若 $(F, f)$ 在维数 $\leqslant 0$ 时上同调真，则 $\alpha_0$ 双射.

现在假设 $(F, f)$ 相对 $S$ 在维数 $\leqslant 1$ 时上同调真，i.e. 典范态射
\[
\phi \colon g^{*} f_* \Phi \longrightarrow f'_* \Phi'
\]
是等价. 设 $G$ 是叠 $f_* \Phi$ 的极大子叠丛层 \cite[III 2.1.8]{exp13-1}；有同构 $G \simeq R^1 f_* F$. 由于 $g^{*} G$ 是 $g^{*} f_* \Phi$ 的极大子叠丛层 \cite[III 2.1.5.5]{exp13-2}，态射 $\alpha_1$ 由 $\phi|_{Y'_1}$ 取极大子叠丛层得到，故是同构.

(ii bis)$\Rightarrow$(ii). 只需证明 $\alpha_0$ 双射时 $a_1$ 单射. 设 $Y'_1$ 是平展于 $Y'$ 上方的概形，$s$ 和 $t$ 是 $g^{*}(R^1 f_* F)(Y'_1)$ 的在 $R^1 f'_* F'(Y'_1)$ 中有同样像的两个元素，证明 $s = t$. 该断言局部于 $Y'_1$ 的平展拓扑，考虑到 $g^{*}(R^1 f_* F)$ 的逆像定义，可假设 $Y'_1$ 是平展于 $Y$ 上方的概形 $Y_1$ 的逆像，且 $s$ 和 $t$ 来自 $X_1$ 上的挠子 $P$ 和 $Q$. 对 $s$ 和 $t$ 的假设于是意味着 $P$ 和 $Q$ 在 $X'_1$ 上的逆像 $P'$ 和 $Q'$ 在 $Y'_1$ 的平展拓扑下局部同构. 令 $R = \mathrm{Hom}_{F_1}(P, Q)$，态射
\[
g_{1*} f_{1*} R \longrightarrow f'_{1*} R'
\]
双射的事实证明 $P$ 和 $Q$ 在 $Y_1$ 的平展拓扑下局部同构，故 $s = t$.

(ii)$\Rightarrow$(i). 为证明 $\phi$ 忠实 (resp. 忠实全)，只需证明：若 $Y_1$ 是平展于 $Y$ 上的概形，$P$，$Q$ 是 $X_1$ 上群 $F_1$ 的两个挠子，$x$，$y$ (resp. $x'$，$y'$) 是 $(g^{*} f_* \Phi)_{Y'_1}$ (resp. $(f'_* \Phi')_{Y'_1}$) 的对应 $P$，$Q$ (resp. $P'$，$Q'$) 的对象，则态射
\[
a \colon \mathrm{Hom}(x, y) \longrightarrow \mathrm{Hom}(x', y')
\]
单射 (resp. 双射). $a$ 等同于典范态射
\[
H^0(Y'_1, g_{1*} f_{1*}(Q \wedge^{F_1} P^{\circ})) \longrightarrow H^0(Y'_1, f'_{1*}(Q' \wedge^{F'_1} P'^{\circ})).
\]
若 $\mathrm{Hom}(x, y) \neq \varnothing$，则 $Q \wedge^{F_1} P^{\circ}$ 是 $P F_1$ 下在 $Y_1$ 上局部平凡的挠子；由此得出 $f_{1*}(Q \wedge^{F_1} P^{\circ})$ 是 $f_{1*}(P F_1)$ 下的挠子，且 $g_{1*} f_{1*}(Q \wedge^{F_1} P^{\circ})$ 是平凡挠子. 态射 $a$ 于是等同于典范态射
\[
H^0(Y_1, g_{1*} f_{1*}(P F_1)) \longrightarrow H^0(Y'_1, f'_{1*}(P' F'_1)).
\]
若 $\mathrm{Hom}(x', y') \neq \varnothing$ 且 $a_1$ 单射也如此，因为此时 $Q' \wedge^{F'_1} P'^{\circ}$ 平凡，且由 $a_1$ 的单射得出 $P$ 和 $Q$ 在 $Y_1$ 上局部同构. 由此得出：$a_0$ 单射 (resp. $a_0$ 双射且 $a_1$ 单射) 时 $\phi$ 忠实 (resp. 忠实全).

剩下证明 $a_0$ 和 $a_1$ 双射时函子 $\phi$ 本质满射. 设 $Y''$ 是平展于 $Y'$ 上方的概形，$X'' = X' \times_{Y'} Y''$，设 $P''$ 是 $X''$ 上群 $F'' = F'|_{X''}$ 的挠子. 将证明存在 $(g^{*} f_* \Phi)_{Y''}$ 的元素 $x$，使它在 $(f'_* \Phi')_{Y''}$ 中的像同构于 $P''$. 设 $p''$ 是 $P''$ 的类. 由 $a_1$ 满射得出可以找到平展满射态射 $Y''_1 \to Y''$、平展态射 $Y_1 \to Y$ (使有态射 $Y''_1 \to Y'_1$) 和 $X_1$ 上群 $F_1$ 的挠子 $P_1$，使 $P_1$ 在 $X''_1$ 上的逆像 $P''_1$ 同构于 $P''$ 的逆像. 用 $\phi$ 忠实全的事实，看到 $(g^{*} f_* \Phi)_{Y''_1}$ 中对应 $P''_1$ 的对象 $x_1$ 配备相对 $Y''_1 \to Y''$ 的下降数据，故来自 $(g^{*} f_* \Phi)_{Y''}$ 的元素 $x$. 由于 $x$ 在 $(f'_* \Phi')_{Y''}$ 中的像是 $P''$，这证明 $\phi$ 本质满射，完成证明.
\end{proof}

\begin{example}\label{exa:XIII.1.4}
设 $f \colon X \to Y$ 是紧合态射. 由 \cite[VII 2.2.2]{exp13-2} 得出，对 $X$ 上每个 ind-有限叠 $F$，对 $(F, f)$ 相对 $Y$ 在维数 $\leqslant 1$ 时上同调真. 特别，对 $X$ 上每个集合层 (resp. 每个群层，resp. 每个 ind-有限群层) $F$，$(F, f)$ 在维数 $\leqslant 0$ (resp. $\leqslant 0$，resp. $\leqslant 1$) 时上同调真.
\end{example}

\begin{remark}\label{rem:XIII.1.5}
a) 设 $F$ 是 $X$ 上使 $(F, f)$ 相对 $S$ 在维数 $\leqslant -1$ (resp. $\leqslant 0$) 时上同调真的群层. 若把 $F$ 看作集合层，$(F, f)$ 相对 $S$ 在维数 $\leqslant -1$ (resp. $\leqslant 0$) 时上同调真，但反之不真.

例如取 $Y$ 为严格局部离散赋值环的谱，闭点为 $t$，生成点为 $s$，$f \colon X \to Y$ 为闭纤维为空的非空 $Y$-概形，$F$ 为 $X$ 上的非平凡常群层，$P$ 为使 $H^0(X_s, P F|_{X_s}) = 1$ 的 $F$ 下挠子. 则把 $P F$ 看作集合层时，$(P F, f)$ 相对 $Y$ 在维数 $\leqslant -1$ 时上同调真. 把 $P F$ 看作群层时，有同构 $(P F)^{\circ P} \simeq F$；由于典范态射
\[
H^0(X, F) \longrightarrow H^0(X_t, F|_{X_t}) = 1
\]
不单射，这证明 $(P F, f)$ 相对 $Y$ 在维数 $\leqslant -1$ 时不真.

b) 假设 $f$ 相容 (i.e. 拟紧且拟分离). 设 $F$ 是 $X$ 上的叠. 对 $Y'$ 的每个几何点 $y$，用 $Y$ (resp. $Y'$) 记 $Y$ (resp. $Y'$) 在 $y$ 处的严格局部化，令 $X = X \times_Y Y$，$X' = X' \times_{Y'} Y'$，等等. 为使 $(F, f)$ 相对 $S$ 在维数 $\leqslant -1$ (resp. $\leqslant 0$，resp. $\leqslant 1$) 时上同调真，必要且充分条件是：对每个 $S$-概形 $S'$ 和 $Y'$ 的每个几何点 $y$，典范函子
\[
F(X) \longrightarrow F'(X')
\]
忠实 (resp. 忠实全，resp. 范畴等价).

确实，若 $S'$ 是 $S$-概形，为使函子
\[
g^{*} f_* F \longrightarrow f'_* F'
\]
忠实 (resp. 忠实全，resp. 等价)，必要且充分条件是它诱导的 $Y'$ 的各几何点 $y'$ 上的纤维函子如此 \cite[III 2.1.5.9]{exp13-2}. 该断言于是由相容态射的叠的直像的几何纤维计算 \cite[VII 2.1.5]{exp13-2} 得出.

c) 设 $F$ 是 $X$ 上的叠. $(F, f)$ 相对 $S$ 在维数 $\leqslant -1$ (resp. $\leqslant 0$，resp. $\leqslant 1$) 时上同调真的事实局部于 $Y$ 的平展拓扑.

设 $S'$ 是 $S$-概形，$F'$ 是 $F$ 在 $X'$ 上的逆像 (参见 \eqref{eq:XIII.1.0.1}). 若 $(F, f)$ 相对 $S$ 在维数 $\leqslant 1$ 时上同调真，则 $(F', f')$ 也如此. 但 $(F, f)$ 相对 $S$ 在维数 $\leqslant -1$ (resp. $\leqslant 0$) 时上同调真不蕴含 $(F', f')$ 如此.

例如设 $S'$ 是离散赋值环，$f' \colon E_{S'} \to S'$ 是 $S'$ 上方的仿射空间，$x$ 是 $E_{S'}$ 的位于 $S'$ 的生成点上方的一个闭点，$F'$ 是 $E_{S'}$ 上的集合层，它在 $E_{S'} - \{x\}$ 上的限制是单元常层，在 $x$ 上方的几何点处的纤维有两个元素. 则 $(F', f')$ 相对 $S'$ 在维数 $\leqslant -1$ 时不真. 设 $S = S'[Z]$，$f \colon E_S \to S$ 是 $S$ 上的仿射空间，$T$ 是 $X = E_S$ 的不与闭集 $Z = 0$ 相交、且使 $f(T)$ 含 $S$ 的生成点的闭部分. 设 $G$ 是 $F'$ 在 $X$ 上的逆像，$F$ 是把 $G|_{X - T}$ 用空集延拓得到的 $X$ 上集合层. 则 $(F, f)$ 相对 $S$ 在维数 $\leqslant -1$ 时上同调真，但经由 $Z = 0$ 定义的基变换 $S' \to S$ 后不再如此.

d) 设 $F$ 是 $X$ 上使 $(F, f)$ 相对 $Y$ 在维数 $\leqslant -1$ (resp. $\leqslant 0$，resp. $\leqslant 1$) 时上同调真的叠. 则对 $Y$ 的每个几何点 $y$，典范函子
\[
(f_* F)_y \longrightarrow F(X_y)
\]
忠实 (resp. 忠实全，resp. 范畴等价).
\end{remark}

\begin{proposition}\label{prop:XIII.1.6}
设 $f \colon X \to Y$ 和 $g \colon Y \to Z$ 是两个 $S$-态射，$\Phi$ 是 $X$ 上的叠.
\begin{enumerate}[1)]
\item 假设 $(\Phi, f)$ 和 $(f_* \Phi, g)$ 相对 $S$ 在维数 $\leqslant -1$ (resp. $\leqslant 0$，resp. $\leqslant 1$) 时上同调真. 则 $(\Phi, gf)$ 也如此.
\item 假设 $(\Phi, gf)$ 相对 $S$ 在维数 $\leqslant -1$ 时\emph{上同调真} (resp. $(\Phi, gf)$ 相对 $S$ 在维数 $\leqslant 0$ 时上同调真且 $(\Phi, f)$ 相对 $S$ 在维数 $\leqslant -1$ 时上同调真，resp. $(\Phi, gf)$ 相对 $S$ 在维数 $\leqslant 1$ 时上同调真且 $(\Phi, f)$ 相对 $S$ 在维数 $\leqslant 0$ 时上同调真). 则 $(f_* \Phi, g)$ 相对 $S$ 在维数 $\leqslant -1$ (resp. $\leqslant 0$，resp. $\leqslant 1$) 时上同调真.
\end{enumerate}
\end{proposition}

对每个 $S$-概形 $S'$，考虑下列所有方形笛卡儿的图式：
\begin{equation}
\begin{tikzcd}[row sep=large, column sep=large]
X' \arrow[r,"{f'}"] \arrow[d,"{h}"] & Y' \arrow[r,"{g'}"] \arrow[d,"{k}"] & Z' \arrow[r] \arrow[d,"{m}"] & S' \arrow[d] \\
X \arrow[r,"{f}"] & Y \arrow[r,"{g}"] & Z \arrow[r] & S.
\end{tikzcd}
\tag{1.6.1}
\label{eq:XIII.1.6.1}
\end{equation}
典范态射
\[
m^{*}(g_* f_* \Phi) \longrightarrow g'_* f'_*(h^{*} \Phi)
\]
等同于典范态射
\[
i \quad j
\]
的复合：
\[
m^{*}(g_* f_* \Phi) \xrightarrow{i} g'_*(k^{*} f_* \Phi) \xrightarrow{j} g'_* f'_*(h^{*} \Phi).
\]
1) 假设蕴含 $i$ 和 $j$ 忠实 (resp. 忠实全，resp. 等价)；故 $ji$ 如此.

2) 假设蕴含 $ji$ 忠实 (resp. $ji$ 忠实全且 $j$ 忠实，resp. $ji$ 等价且 $j$ 忠实全)；由此得出 $i$ 忠实 (resp. 忠实全，resp. 等价).

\begin{corollary}\label{cor:XIII.1.7}
设 $f \colon X \to Y$ 和 $g \colon Y \to Z$ 是两个 $S$-态射，$F$ 是 $X$ 上的群层. 假设 $(F, gf)$ 相对 $S$ 在维数 $\leqslant -1$ 时上同调真 (resp. $(F, gf)$ 相对 $S$ 在维数 $\leqslant 0$ 时上同调真且 $(F, f)$ 相对 $S$ 在维数 $\leqslant -1$ 时上同调真). 则 $(f_* F, g)$ 相对 $S$ 在维数 $\leqslant -1$ (resp. $\leqslant 0$) 时上同调真.
\end{corollary}

沿用 \eqref{eq:XIII.1.6.1} 的记号，对每个平展于 $Y$ 上方的概形 $Y_1$，用 $f_1$，$F_1$ 记 $f$，$F$ 经态射 $Y_1 \to Y$ 的逆像. 设 $\Phi$ 是 $F$ 下的挠子叠，$\Psi$ 是 $f_* F$ 下的挠子叠. 有典范函子
\[
\phi \colon \Psi \longrightarrow f_* \Phi,
\]
它把平展于 $Y$ 上的每个概形 $Y_1$ 和 $Y_1$ 上群 $f_{1*} F_1$ 的每个挠子 $P$ 配给由 $f_1^{*} P$ 经结构群扩张 $f_1^{*} f_{1*} F_1 \to F_1$ 得到的 $X_1$ 上的挠子 $\widetilde{P}$. 函子 $\phi$ 忠实全. 确实，若 $P$ 和 $Q$ 是 $Y_1$ 上群 $f_{1*} F_1$ 的两个挠子，有典范态射
\[
\mathrm{Isom}_{f_{1*} F_1}(P, Q) \longrightarrow f_{1*}(\mathrm{Isom}_{F_1}(\widetilde{P}, \widetilde{Q}))
\]
它局部如此故是同构. 由此得出典范态射
\[
\mathrm{Isom}_{f_{1*} F_1}(P, Q) \longrightarrow \mathrm{Isom}_{F_1}(\widetilde{P}, \widetilde{Q})
\]
是同构，故 $\phi$ 忠实全.

有交换图式
\[
\begin{tikzcd}[row sep=large, column sep=large]
g'_* k^{*} \Psi \arrow[r,"{i}"] \arrow[d,"{g'_* k^{*}(\phi)}"] & m^{*}(g_* \Psi) \arrow[d,"{m^{*} g_*(\phi)}"] \\
g'_* k^{*}(f_* \Phi) \arrow[r,"{j}"] & m^{*}(g_* f_* \Phi),
\end{tikzcd}
\]
其中 $i$ 和 $j$ 是基变换态射. 由 \ref{prop:XIII.1.6} 2) 得出 $j$ 忠实 (resp. 忠实全). 由于 $g'_* k^{*}(\phi)$ 和 $m^{*} g_*(\phi)$ 忠实全，由上面图式得出 $i$ 忠实 (resp. 忠实全).

\begin{corollary}\label{cor:XIII.1.8}
设 $f \colon X \to Y$ 是相容 $S$-态射，$g \colon Y \to Z$ 是紧合 $S$-态射，$\Phi$ 是 $X$ 上的 ind-有限叠 \cite[VII 2.2.1]{exp13-2}. 假设 $(\Phi, f)$ 相对 $S$ 在维数 $\leqslant -1$ (resp. $\leqslant 0$，resp. $\leqslant 1$) 时上同调真，则 $(\Phi, gf)$ 也如此.
\end{corollary}

由于 $f$ 相容，$f_* \Phi$ 是 ind-有限\emph{叠} (SGA 4 IX 1.6 (ii)). 该推论于是由 \ref{prop:XIII.1.6} 1) 和 \ref{exa:XIII.1.4} 得出.

\begin{corollary}\label{cor:XIII.1.9}
设 $f \colon X \to Y$ 是整 $S$-态射，$g \colon Y \to Z$ 是 $S$-态射. 若 $F$ 是 $X$ 上的集合层，为使 $(f_* F, g)$ 相对 $S$ 在维数 $\leqslant -1$ (resp. $\leqslant 0$) 时上同调真，必要且充分条件是 $(F, gf)$ 如此. 若 $F$ 是 $X$ 上的群层，为使 $(f_* F, g)$ 相对 $S$ 在维数 $\leqslant -1$ (resp. $\leqslant 0$，resp. $\leqslant 1$) 时上同调真，必要且充分条件是 $(F, gf)$ 如此.
\end{corollary}

集合层情形的断言由 \ref{prop:XIII.1.6} 和 $(F, f)$ 相对 $S$ 在维数 $\leqslant 0$ 时上同调真的事实得出. 设 $F$ 是 $X$ 上的群层，$\Phi$ 是 $F$ 下的挠子叠. 由 SGA 4 VIII 5.8，$F$ 下的每个挠子局部平凡于 $Y$ 上. 由此得出叠 $f_* \Phi$ 等价于 $f_* F$ 下的挠子叠，等价由把平展于 $Y$ 上的每个概形 $Y_1$ 和 $X_1 = X \times_Y Y_1$ 上群 $F|_{X_1}$ 的每个挠子 $P$ 配给群 $f_* F|_{Y_1}$ 的挠子 $f_* P$ 得到. 由于 $(F, f)$ 相对 $S$ 在维数 $\leqslant 1$ 时上同调真，该推论于是由 \ref{prop:XIII.1.6} 得出.

\begin{definition}\label{def:XIII.1.10}
1.10.1. 设 $E$ 是范畴，考虑图式
\[
\Phi \xrightarrow{p} \Phi_1 \xrightarrow[p_2]{p_1} \Phi_2,
\]
其中 $\Phi$，$\Phi_1$，$\Phi_2$ 是 $E$ 上方的纤维范畴，箭是纤维范畴态射，设
\[
a \colon p_1 p \xrightarrow{\sim} p_2 p
\]
是函子同构.

称上图式正合，若满足条件 a)：对 $\Phi$ 的每对对象 $x$，$y$ 和满足 $p_1(f) = p_2(f)$ (借 $a$ 等同 $p_1 p$ 和 $p_2 p$) 的态射 $f \colon p(x) \to p(y)$，存在唯一态射 $g \colon x \to y$ 使 $p(g) = f$.

1.10.2. 考虑图式
\[
\Phi \xrightarrow{p} \Phi_1 \xrightarrow[p_1, p_2]{} \Phi_2 \xrightarrow[p_{12}, p_{23}, p_{31}]{} \Phi_3,
\]
其中 $\Phi$，$\Phi_i$ ($1 \leqslant i \leqslant 3$) 是 $E$ 上的纤维范畴，箭是纤维范畴态射. 假设给定函子同构
\[
a \colon p_1 p \xrightarrow{\sim} p_2 p
\]
\[
a_1 \colon p_{31} p_2 \xrightarrow{\sim} p_{12} p_1, \quad a_2 \colon p_{12} p_2 \xrightarrow{\sim} p_{23} p_1, \quad a_3 \colon p_{23} p_2 \xrightarrow{\sim} p_{31} p_1,
\]
使下列图式交换：
\[
\begin{tikzcd}[row sep=large, column sep=large]
p_{23} p_1 p \arrow[r,"{\mathrm{id}.a}"] \arrow[d,"{a_2 . \mathrm{id}}"] & p_{23} p_2 p \arrow[r,"{a_3 . \mathrm{id}}"] & p_{31} p_1 p \arrow[d,"{\mathrm{id}.a}"] \\
p_{12} p_2 p & p_{12} p_1 p \arrow[l,"{\mathrm{id}.a}"'] & p_{31} p_2 p \arrow[l,"{a_1 . \mathrm{id}}"']
\end{tikzcd}
\]
等同 $p_1 p$ 和 $p_2 p$，$p_{31} p_2$ 和 $p_{12} p_1$，等等.

称上图式正合，若满足条件：
a) 与 1.10.1 的条件 a) 类似.
b) 对 $\Phi_1$ 的每个对象 $x_1$ 和满足
\[
(1.10.2.1) \quad p_{23}(u) p_{31}(u) = p_{12}(u)^{-1}
\]
的同构 $u \colon p_1(x_1) \xrightarrow{\sim} p_2(x_1)$，存在 $\Phi$ 的对象 $x$ 和同构 $i \colon p(x) \xrightarrow{\sim} x_1$，使图式
\begin{equation}
\begin{tikzcd}[row sep=large, column sep=large]
p_1 p(x) \arrow[r,"{p_1(i)}"] \arrow[d] & p_1(x_1) \arrow[d,"{u}"] \\
p_2 p(x) \arrow[r,"{\sim}","{p_2(i)}"'] & p_2(x_1)
\end{tikzcd}
\tag{1.10.2.2}
\label{eq:XIII.1.10.2.2}
\end{equation}
交换.
\end{definition}

1.10.3. 以显然方式定义范畴 $E$ 上方纤维范畴的正合图式态射概念.

1.10.4. 特别将用 $E$ 是位点、$\Phi$，$\Phi_i$ ($1 \leqslant i \leqslant 3$) 是 $E$ 上的叠的情形中的正合图式概念.

设 $f \colon E \to E'$ 是位点态射，
\[
(1.10.4.1) \quad \Phi \xrightarrow{p} \Phi_1 \xrightarrow[p_2]{p_1} \Phi_2 \xrightarrow[p_{12}, p_{23}, p_{31}]{} \Phi_3
\]
是 $E$ 上叠的正合图式. 经直像得到显然正合的图式
\[
f_* \Phi \longrightarrow f_* \Phi_1 \rightrightarrows f_* \Phi_2 \Rrightarrow f_* \Phi_3.
\]
若 $h \colon E'' \to E$ 是位点态射，同样有正合图式
\[
h^{*} \Phi \xrightarrow{p''} h^{*} \Phi_1 \xrightarrow[p''_1, p''_2]{} h^{*} \Phi_2 \xrightarrow[p''_{12}, p''_{23}, p''_{31}]{} h^{*} \Phi_3.
\]
先验证 1.10.2 的条件 a). 设 $F''$ 是 $E''$ 的对象，$x''$ 和 $y''$ 是 $(h^{*} \Phi)_{F''}$ 的两个对象，$x''_1$ 和 $y''_1$ 是它们在 $h^{*} \Phi_1$ 中的像，$x''_2$，$y''_2$ 是它们在 $h^{*} \Phi_2$ 中的像. 设 $u''_1 \colon x''_1 \to y''_1$ 是满足 $p''_1(u''_1) = p''_2(u''_1)$ 的态射，证明 $u''_1$ 来自唯一态射 $u'' \colon x'' \to y''$. 问题局部于 $F''$，可假设有 $E$ 的对象 $F_1$，$F''$ 到 $F_1$ 经 $h$ 的逆像 $F''_1$ 的态射，$\Phi_{F_1}$ 的在 $F''$ 上的逆像是 $x''$ 和 $y''$ 的对象 $x$，$y$. 设 $x_1$，$y_1$ (resp. $x_2$，$y_2$) 是 $x$，$y$ 在 $\Phi_1$ (resp. $\Phi_2$) 中的像. 可假设 $u''_1$ 来自满足 $p_1(u_1) = p_2(u_1)$ 的态射 $u_1 \colon x_1 \to y_1$. 由 (1.10.4.1) 的正合性，得到唯一态射 $u \colon x \to y$，它的 $h$ 逆像是所求态射 $u''$.

1.10.2 的条件 b) 类似验证. 设 $x''_1$ 是 $(h^{*} \Phi_1)_{F''}$ 的对象，$u'' \colon p''_1(x''_1) \to p''_2(x''_1)$ 是满足关系
\[
p''_{23}(u'') p''_{31}(u'') = p''_{12}(u'')^{-1}
\]
的态射，证明存在 $(h^{*} \Phi)_{F''}$ 的对象 $x''$ 和同构 $i'' \colon p''(x'') \simeq x''_1$，使类似 \eqref{eq:XIII.1.10.2.2} 的图式交换. 由于问题局部于 $F''$，可假设有对象 $F_1$、如上 $F'' \to F''_1$ 的态射和 $(\Phi_1)_{F_1}$ 的在 $(h^{*} \Phi_1)_{F''}$ 中的逆像是 $x''_1$ 的对象 $x_1$. 同样可假设 $u''$ 来自满足 \eqref{eq:XIII.1.10.2.2} 的态射 $u \colon p_1(x_1) \to p_2(x_1)$. $\Phi_{F_1}$ 的对象 $x$ (它的 $h$ 逆像是回答问题的元素 $x''$) 的存在于是由 (1.10.4.1) 的正合性得出.

\begin{example}\label{exa:XIII.1.11}
1) 设 $f \colon X_1 \to X$ 是变动的概形上平展层范畴的下降态射 (例如\emph{万有浸没态射} (SGA 4 VIII 9.3)). 设 $X_2 = X_1 \times_X X_1$，$g \colon X_2 \to X$ 是典范投影，$F$ 是 $X$ 上的集合层. 由 loc. cit. 得出有集合层的正合序列
\[
(1.11.1) \quad F \longrightarrow f_* f^{*} F \rightrightarrows g_* g^{*} F.
\]
若 $\Phi$ 是与 $F$ 关联的离散范畴叠，$\Phi_3$ 是 $X$ 上的终叠，i.e. 所有纤维约化为一个元素 (唯一态射是恒等态射) 的叠，说序列 (1.11.1) 正合等价于说叠图式
\[
\Phi \longrightarrow f_* f^{*} \Phi \rightrightarrows g_* g^{*} \Phi \Rrightarrow \Phi_3
\]
如此.

2) 设 $f \colon X_1 \to X$ 是变动的概形上平展层范畴的有效下降态射 (例如紧合满射态射，或整满射态射，或局部有限表现忠实平坦态射 (SGA 4 VIII 9.4)). 设 $X_2 = X_1 \times_X X_1$，$g \colon X_2 \to X$ 是典范投影，$X_3 = X_1 \times_X X_1 \times_X X_1$，$h \colon X_3 \to X$ 是典范态射. 设 $\Phi$ 是 $X$ 上的叠，$\Phi_1 = f_* f^{*} \Phi$，$\Phi_2 = g_* g^{*} \Phi$，$\Phi_3 = h_* h^{*} \Phi$. 则有正合图式
\[
\Phi \longrightarrow \Phi_1 \rightrightarrows \Phi_2 \Rrightarrow \Phi_3,
\]
其中箭是投影关联的典范态射.

确实把 $\Phi$ 看作配备平展拓扑的 $X$ 上方概形范畴 $\mathbf{Sch}_X$ 上的叠. 由 \cite[VII 2.2.8]{exp13-2}，$\Phi$ 对 $\mathbf{Sch}_X$ 上使覆盖态射为平展层范畴的有效下降态射的最粗拓扑也是叠. 上图式的正合性由此立即得出.
\end{example}

\begin{proposition}\label{prop:XIII.1.12}
设 $S$ 是概形，$f \colon X \to Y$ 是 $S$-态射.
\begin{enumerate}[1)]
\item 设
\[
\Phi \longrightarrow \Phi_1 \rightrightarrows \Phi_2
\]
是 $X$ 上叠的正合图式. 假设 $(\Phi_1, f)$ 相对 $S$ 在维数 $\leqslant 0$ 时上同调真，$(\Phi_2, f)$ 相对 $S$ 在维数 $\leqslant -1$ 时上同调真. 则 $(\Phi, f)$ 相对 $S$ 在维数 $\leqslant 0$ 时上同调真.
\item 设
\[
\Phi \longrightarrow \Phi_1 \rightrightarrows \Phi_2 \Rrightarrow \Phi_3
\]
是 $X$ 上叠的正合图式. 假设 $(\Phi, f)$ 相对 $S$ 在维数 $\leqslant 1$ 时上同调真，$(\Phi_2, f)$ 相对 $S$ 在维数 $\leqslant 0$ 时上同调真，$(\Phi_3, f)$ 相对 $S$ 在维数 $\leqslant -1$ 时上同调真. 则 $(\Phi, f)$ 相对 $S$ 在维数 $\leqslant 1$ 时上同调真.
\end{enumerate}
\end{proposition}

对每个 $S$-概形 $S'$，考虑下列所有方形笛卡儿的交换图式：
\[
\begin{tikzcd}[row sep=large, column sep=large]
X' \arrow[r,"{f'}"] \arrow[d,"{h}"] & Y' \arrow[r] \arrow[d,"{g}"] & S' \arrow[d] \\
X \arrow[r,"{f}"] & Y \arrow[r] & S.
\end{tikzcd}
\]
证明 2)，1) 的证明类似. 由于直像和逆像函子把叠的正合图式变成正合图式 (1.10.4)，有叠的正合图式态射：
\[
\begin{tikzcd}[row sep=large, column sep=large]
g^{*} f_* \Phi \arrow[r,"{\pi}"] \arrow[d,"{\phi}"] & g^{*} f_* \Phi_1 \arrow[r,"{\pi_1}","{\pi_2}"'] \arrow[d,"{\phi_1}"] & g^{*} f_* \Phi_2 \arrow[r, Rightarrow] \arrow[d,"{\phi_2}"] & g^{*} f_* \Phi_3 \arrow[d,"{\phi_3}"] \\
f'_* h^{*} \Phi \arrow[r] & f'_* h^{*} \Phi_1 \arrow[r, Rightarrow] & f'_* h^{*} \Phi_2 \arrow[r, Rightarrow] & f'_* h^{*} \Phi_3.
\end{tikzcd}
\]
由假设 $\phi_1$ 是范畴等价，$\phi_2$ 忠实全，$\phi_3$ 忠实. 由前面图式得出 $\phi$ 是等价.

\begin{proposition}\label{prop:XIII.1.13}
设 $f \colon X \to Y$ 是 $S$-态射.
\begin{enumerate}[1)]
\item 设
\[
F \longrightarrow G \rightrightarrows H
\]
是 $X$ 上集合层的正合图式. 假设 $(G, f)$ 相对 $S$ 在维数 $\leqslant 0$ 时上同调真，$(H, f)$ 相对 $S$ 在维数 $\leqslant -1$ 时上同调真. 则 $(F, f)$ 相对 $S$ 在维数 $\leqslant 0$ 时上同调真.
\item 设 $F \to G$ 是 $X$ 上群层的单态射. 若 $Y_1$ 是平展于 $Y$ 上的概形，令 $X_1 = Y_1 \times_Y X$，用 $f_1$ (resp. $F_1$，resp. $G_1$) 记 $f$ (resp. $F$，resp. $G$) 在 $Y_1$ 上的逆像 (参见 1.0.2). 假设 $(G, f)$ 相对 $S$ 在维数 $\leqslant 0$ (resp. $\leqslant 1$) 时上同调真，且对每个平展于 $Y$ 上的概形 $Y_1$ 和 $G_1$ 下的每个挠子 $Q$，$(Q/F_1, f_1)$ 相对 $S$ 在维数 $\leqslant -1$ (resp. $\leqslant 0$) 时上同调真. 则 $(F, f)$ 相对 $S$ 在维数 $\leqslant 0$ (resp. $\leqslant 1$) 时上同调真.
\item 设 $F \to G$ 是 $X$ 上群层的单态射. 假设 $(F, f)$ 相对 $S$ 在维数 $\leqslant 1$ 时上同调真，$(G, f)$ 相对 $S$ 在维数 $\leqslant 0$ 时上同调真. 则对 $G$ 下的每个挠子 $Q$，$(Q/F, f)$ 相对 $S$ 在维数 $\leqslant 0$ 时上同调真.
\end{enumerate}
\end{proposition}

\begin{proof}
1) 设 $\Phi$ (resp. $\Phi_1$，resp. $\Phi_2$) 是与 $F$ (resp. $G$，resp. $H$) 关联的离散范畴叠，$\Phi_3$ 是 $X$ 上的终叠. 有正合图式
\[
\Phi \longrightarrow \Phi_1 \rightrightarrows \Phi_2 \Rrightarrow \Phi_3.
\]
由假设 $(\Phi_1, f)$ 相对 $S$ 在维数 $\leqslant 1$ 时上同调真，$(\Phi_2, f)$ 相对 $S$ 在维数 $\leqslant 0$ 时上同调真 (1.2). 由于 $(\Phi_3, f)$ 显然相对 $S$ 在维数 $\leqslant -1$ 时上同调真，由 \ref{prop:XIII.1.12} 得出 $(\Phi, f)$ 相对 $S$ 在维数 $\leqslant 1$ 时上同调真，i.e. $(F, f)$ 相对 $S$ 在维数 $\leqslant 0$ 时上同调真.

2) 先证明：若 $(G, f)$ 相对 $S$ 在维数 $\leqslant 0$ 时上同调真，且诸 $(Q/F_1, f_1)$ 相对 $S$ 在维数 $\leqslant -1$ 时上同调真，则 $(F, f)$ 相对 $S$ 在维数 $\leqslant 0$ 时上同调真. 由 \ref{prop:XIII.1.3.1} 只需证明：对每个平展于 $Y$ 上方的概形 $Y_1$ 和 $X_1$ 上群 $F_1$ 的每个挠子 $P$，把 $P F_1$ 看作集合层时 $(P F_1, f_1)$ 相对 $S$ 在维数 $\leqslant 0$ 时上同调真，且典范态射
\[
d \colon g^{*}(R^1 f_* F) \longrightarrow R^1 f'_* F'
\]
单射. 第一个断言由 1) 立即得出，因为若 $Q$ 表示由 $P$ 经结构群扩张 $F_1 \to G_1$ 得到的挠子，有同构
\[
Q^{G_1}/P F_1 \xrightarrow{\sim} Q/F_1.
\]
证明 $d$ 单射. 只需证明：若 $Y_1$ 是平展于 $Y$ 上方的概形，$P$ 和 $\widetilde{P}$ 是 $F_1$ 下的两个在 $X'_1$ 上的逆像 $P'$ 和 $\widetilde{P}'$ 同构的挠子，则 (必要时作 $Y_1$ 的平展满射扩张) $P$ 和 $\widetilde{P}$ 变成同构. 选同构 $p' \colon P' \xrightarrow{\sim} \widetilde{P}'$. 设 $Q$ (resp. $\widetilde{Q}$) 是由 $P$ (resp. $\widetilde{P}$) 经结构群扩张 $F_1 \to G_1$ 得到的挠子. $Q$ (resp. $\widetilde{Q}$) 在 $X'_1$ 上的逆像 $Q'$ (resp. $\widetilde{Q}'$) 由 $P'$ (resp. $\widetilde{P}'$) 经结构群扩张 $F'_1 \to G'_1$ 得到；设 $q' \colon Q' \xrightarrow{\sim} \widetilde{Q}'$ 是由 $p'$ 同样得到的同构. 由于 $(G, f)$ 相对 $S$ 在维数 $\leqslant 0$ 时上同调真，可假设 (必要时作 $Y_1$ 的平展满射扩张) $q'$ 是同构 $q \colon Q \xrightarrow{\sim} \widetilde{Q}$ 的像. 与挠子 $P$ (resp. $\widetilde{P}$) 关联 $Q/F_1$ 的截面 $x$ (resp. $\widetilde{Q}/\widetilde{F}_1$ 的截面 $\widetilde{x}$)，且为使 $P$ 和 $\widetilde{P}$ 同构，必要且充分条件是在同构 $Q \xrightarrow{\sim} \widetilde{Q}$ 使导出的同构
\[
e \colon H^0(X_1, Q/F_1) \longrightarrow H^0(X_1, \widetilde{Q}/\widetilde{F}_1)
\]
把 $x$ 变成 $\widetilde{x}$. 取同构 $q$. 截面 $e(x)$ 和 $\widetilde{x}$ 在 $H^0(X'_1, \widetilde{Q}'/\widetilde{F}'_1)$ 中有同样的像. 由于 $(\widetilde{Q}/\widetilde{F}_1, f_1)$ 相对 $S$ 在维数 $\leqslant -1$ 时上同调真，必要时作 $Y_1$ 的平展满射扩张，有 $e(x) = \widetilde{x}$，这证明 $d$ 单射.

为完成证明，剩下证明：若 $(G, f)$ 相对 $S$ 在维数 $\leqslant 1$ 时上同调真，且对每个平展于 $Y$ 上方的概形 $Y_1$ 和 $X_1$ 上群 $F_1$ 的每个挠子 $Q$，$(Q/F_1, f_1)$ 相对 $S$ 在维数 $\leqslant 0$ 时上同调真，则态射 $d$ 满射. 设 $P'$ 是 $X'_1$ 上群 $F'_1$ 的挠子，$Q'$ 是由 $P'$ 经结构群扩张得到的 $G'_1$ 下挠子. $P'$ 的给出等价于 $Q'$ 和 $H^0(X'_1, Q'/F'_1)$ 的截面 $x'$ 的给出. 由态射
\[
g^{*}(R^1 f_* G) \longrightarrow R^1 f'_* G'
\]
的满射得出，必要时作 $Y_1$ 的平展满射扩张，存在 $G_1$ 下的挠子 $Q$，使它在 $X'_1$ 上的逆像同构于 $Q'$. 用 $(Q/F_1, f_1)$ 相对 $S$ 在维数 $\leqslant 0$ 时上同调真的事实，同样可假设存在 $H^0(X_1, Q/F_1)$ 的元素 $x$，使它在 $H^0(X'_1, Q'/F'_1)$ 中的像是 $x'$. $Q$ 和 $x$ 的给出确定 $F_1$ 下的挠子 $P$，使它在 $X_1$ 上的逆像同构于 $P'$，这证明 $d$ 满射.

3) 证明 $(Q/F, f)$ 相对 $S$ 在维数 $\leqslant -1$ 时上同调真，i.e. 对每个 $S$-概形 $S'$ 和每个平展于 $Y$ 上方的概形 $Y_1$，若 $x$，$\widetilde{x}$ 是 $H^0(X_1, Q_1/F_1)$ 的在 $H^0(X'_1, Q'_1/F'_1)$ 中的像 $x'$，$\widetilde{x}'$ 相等的两个元素，则 (作 $Y_1$ 的满射扩张后) $x = \widetilde{x}$. 与 $x$ (resp. $\widetilde{x}$) 关联 $F_1$ 下的挠子 $P$ (resp. $\widetilde{P}$)，使 $Q_1$ 由 $P$ (resp. $\widetilde{P}$) 经结构群扩张 $F_1 \to G_1$ 得到. 由关系 $x' = \widetilde{x}'$ 得出同构 $u' \colon P' \xrightarrow{\sim} \widetilde{P}'$，使它在 $Q'_1$ 上诱导的同构是恒等. 由于 $(F, f)$ 相对 $S$ 在维数 $\leqslant 0$ 时上同调真，由此得出 (作 $Y_1$ 的平展满射扩张后) 有提升 $u'$ 的同构 $u \colon P \to \widetilde{P}$；$(G, f)$ 相对 $S$ 在维数 $\leqslant -1$ 时上同调真的事实蕴含 $x = \widetilde{x}$.

证明 $(Q/F, f)$ 相对 $S$ 在维数 $\leqslant 0$ 时上同调真. 设 $Y''$ 是平展于 $Y'$ 上的概形，$x''$ 是 $H^0(X'', Q''/F'')$ 的元素. 与 $x''$ 关联 $X''$ 上群 $F''$ 的挠子 $P''$. 由于 $(F, f)$ 相对 $S$ 在维数 $\leqslant 1$ 时上同调真，可以找到平展满射态射 $Y''_1 \to Y''$ 和 $Y_1 \to Y$，使有态射 $Y''_1 \to Y'_1$，以及 $X_1$ 上群 $F_1$ 的挠子 $P$，使 $P$ 在 $X''_1$ 上的逆像同构于 $P''$ 的逆像. 由 $(G, f)$ 相对 $S$ 在维数 $\leqslant 0$ 时上同调真的事实得出，甚至可以选 $Y''_1$ 和 $Y_1$，使由 $P$ 经结构群扩张 $F_1 \to G_1$ 得到的挠子同构于 $Q_1$；与 $P$ 对应 $H^0(X_1, Q_1/F_1)$ 的元素 $x$，使它在 $H^0(X''_1, Q''_1/F''_1)$ 中的像同构于 $x''$ 的逆像，这完成证明.
\end{proof}

\begin{proposition}\label{prop:XIII.1.14}
设 $f \colon X \to S$ 是 $S$-概形，$F$ 是 $X$ 上的集合层或群层 (resp. ind-$L$-群层，$L$ 是素数集). 假设 $F$ 局部常，$(F, f)$ 在维数 $\leqslant 0$ (resp. $\leqslant 1$) 时上同调真，且 $f$ 局部 0-无环 (resp. 对 $L$ 局部 1-非球) (SGA 4 XV 1.11). 则对 $S$ 的几何点的每个特殊化 $s_1 \to s_2$，特殊化态射 (SGA 4 VIII 7.1)
\[
a_0 \colon (f_* F)_{s_2} \longrightarrow (f_* F)_{s_1}
\]
是同构，且若 $F$ 是群层，态射
\[
a_1 \colon (R^1 f_* F)_{s_2} \longrightarrow (R^1 f_* F)_{s_1}
\]
单射 (resp. 态射 $a_0$ 和 $a_1$ 是同构).
\end{proposition}

证明照抄 SGA 4 XVI 2.3 的证明，把"紧合"换成"上同调真"即可.

\begin{corollary}\label{cor:XIII.1.15}
设 $f \colon X \to S$ 是态射，$\Phi$ 是 $X$ 上的叠，$L$ 是素数集. 假设对每个平展于 $X$ 上的概形 $X_1$ 和 $\Phi_{X_1}$ 的每对对象 $x$，$y$，层 $\mathrm{Hom}_{X_1}(x, y)$ 局部常，层 $\mathrm{Aut}_{X_1}(x)$ 是局部常 ind-$L$-群，$\Phi$ 的极大子叠丛层 $S_\Phi$ \cite[III 2.1.7]{exp13-1} 局部常. 假设 $(\Phi, f)$ 在维数 $\leqslant 1$ 时上同调真，$f$ 对 $L$ 局部 1-非球. 则对 $S$ 的几何点的每个特殊化 $s_1 \to s_2$，特殊化态射
\[
a \colon (f_* \Phi)_{s_2} \longrightarrow (f_* \Phi)_{s_1}
\]
是范畴等价.
\end{corollary}

设 $S_1$ (resp. $S_2$) 是 $S$ 在 $s_1$ (resp. $s_2$) 处的严格局部化，$X_2$，$\Phi_2$ (resp. $X_1$，$\Phi_1$) 是 $X_2$，$\Phi_2$ 在 $S_2$ 上 (resp. $X_1$，$\Phi_1$ 在 $S_1$ 上) 的逆像，考虑笛卡儿方形
\[
\begin{tikzcd}[row sep=large, column sep=large]
X_1 \arrow[d,"{f_1}"] & X_2 \arrow[l,"h"] \arrow[d,"{f_2}"] \\
S_1 & S_2 \arrow[l,"g"]
\end{tikzcd}
\]
需证明函子
\[
\phi \colon \Phi_2(X_2) \longrightarrow \Phi_1(X_1)
\]
是等价. 函子 $\phi$ 忠实全；设 $x$，$y$ 是 $(\Phi_2)_{X_2}$ 的两个对象；典范态射
\[
\mathrm{Hom}_{X_2}(x, y) \longrightarrow \mathrm{Hom}_{X_1}(\phi(x), \phi(y))
\]
等同于典范态射
\[
H^0(X_2, \mathrm{Hom}_{X_2}(x, y)) \longrightarrow H^0(X_1, h^{*}(\mathrm{Hom}_{X_2}(x, y))).
\]
由 \ref{prop:XIII.1.14} 该态射是同构.

证明 $\phi$ 是等价. 设 $x_1$ 是 $\Phi_1(X_1)$ 的对象，$G_1$ 是 $x_1$ 生成的 $\Phi_1$ 的极大子叠丛. 态射
\[
H^0(X_2, S_{\Phi_2}) \longrightarrow H^0(X_1, h^{*}(S_{\Phi_2})) = H^0(X_1, S_{\Phi_1})
\]
双射，故存在 $\Phi_2$ 的极大子叠丛 $G_2$，使 $h^{*} G_2$ 同构于 $G_1$. 于是只需证明函子
\[
G_2 \longrightarrow h_* h^{*} G_2
\]
是等价. 但这种形式的问题局部于 $X_2$ 的平展拓扑，可假设 $G_2$ 是 $G_2$ 的某个对象的自同构群下的挠子叠丛，此时断言由 \ref{prop:XIII.1.14} 得出.

\begin{corollary}\label{cor:XIII.1.16}
假设同 1.14. 再假设 $F$ 是集合层 (resp. ind-$L$-群层)，且 $f_* F$ (resp. $R^1 f_* F$) 可构，则 $f_* F$ (resp. $R^1 f_* F$) 局部常.
\end{corollary}

该推论由 \ref{prop:XIII.1.14} 借助 SGA 4 IX 2.11 得出.

\begin{remark}\label{rem:XIII.1.17}
回忆 $f$ 局部 0-无环的条件在 $f$ 平坦、纤维可分、$X$ 和 $Y$ 局部诺特时满足 (SGA 4 XV 4.1)，$f$ 对 $L$ 局部 1-非球的条件在 $f$ 光滑、$L$ 为与 $S$ 的剩余特征不同的素数集时满足 (SGA 4 XV 2.1).
\end{remark}

\section{上同调真性的特殊情形：相对正规交叉除子}

2.0. 设 $R$ 是分式域 $K$ 的离散赋值环，$L$ 是 $K$-平展代数；则 $L$ 是有限个域 $L_i$ 的直积，其中 $L_i$ 是 $K$ 的平展扩张. 若 $L'_i$ 表示 $L_i$ 在 $L_i$ 的代数闭包中生成的 伽罗瓦扩张，称 $L$ 在 $R$ 上温和分歧，若诸 $L'_i$ 是 X 3 意义下的温和分歧扩张，i.e. $L'_i|K$ 的某个惯性群 $I_i$ 的阶与 $R$ 的剩余特征 $p$ 互素.

已知 $I_i$ 在一切情形都是与 $p$ 互素的循环群被 $p$-群扩张得到的. (当 $R$ 的剩余扩张可分时，这由 \cite[第四章 命题 7 推论 4]{exp13-5} 得出. loc. cit. 中给出的证明按下列方式推广到一般情形. 沿用 loc. cit. 的假设和记号但不假设剩余扩张可分. 设 $H_i$ 是惯性群 $G_0$ 的由使 $A_L$ 的每个幺模化元 $\pi$ 满足 $s\pi/\pi \in U^i$ 的元素 $s$ 组成的子群. 验证 $G_0/H_1$ 是与 $p$ 互素的阶群，且对 $i \geqslant 1$，$H^i/H^{i+1}$ 是 $p$-群，由此得到所述结果.)

2.0.1. 当 $R$ 严格局部时，有下列简单刻画：$K$-代数 $L$ 在 $R$ 上温和分歧，当且仅当诸 $[L_i : K]$ 与 $p$ 互素. 而且，若 $L$ 在 $R$ 上温和分歧，诸 $L_i$ 是 $K$ 的循环扩张. 确实，当 $R$ 严格局部时，$I_i$ 等于 $L'_i$ 在 $K$ 上的 伽罗瓦群. 正如刚才所回忆，$I_i$ 是与 $p$ 互素的循环群被 $p$-群扩张得到的. 若假设 $L'_i$ 在 $R$ 上温和分歧，$I_i$ 是与 $p$ 互素的阶循环群. 由此得出 $[L_i : K]$ 与 $p$ 互素且 $L_i = L'_i$. 反之，若 $[L_i : K]$ 与 $p$ 互素，$I_i$ 不能含非平凡正规 $p$-子群；故 $I_i$ 是与 $p$ 互素的阶循环群，这证明 $L_i$ 在 $R$ 上温和分歧.

2.0.2. 设 $R$ 是分式域 $K$ 的离散赋值环，$L$ 是 $K$-平展代数，$R$ 是 $R$ 的严格局部化，$K$ 是它的分式域，$L = L \otimes_K K$. 则 $L$ 在 $R$ 上温和分歧\index[term]{wenhefenqidaishu@温和分歧代数}当且仅当 $L$ 在 $R$ 上温和分歧. 归结到 $L$ 是域的情形. 设 $L = \prod_i L_i$，其中诸 $L_i$ 是 $K$ 的扩张域；若 $L'$ 是 $L$ 生成的 伽罗瓦扩张，且 $L' = L' \otimes_K K$，同样有 $L'$ 分解成域之积 $L' = \prod_j L'_j$，且每个 $L_i$ 是某个 $L'_j$ 的子扩张. 由于 $L'$ 是 $K$ 的 伽罗瓦扩张，诸 $L'_j$ 是 $K$ 的 伽罗瓦扩张. 假设 $L$ 在 $R$ 上温和分歧；由于 $L'_j|K$ 的 伽罗瓦群同构于 $L'|K$ 的惯性群，诸 $L'_j$ 在 $R$ 上温和分歧，故诸 $L_i$ 也如此. 反之，假设 $L$ 在 $R$ 上温和分歧. 对每个 $j$，设 $v_j$ 是 $L'_j$ 的延拓 $K$ 的赋值的离散赋值，仍用 $v_j$ 记它在 $L'$ 上诱导的赋值. $j$ 变动时，$v_j$ 遍历 $L'$ 的延拓 $K$ 的赋值的诸赋值. 设 $G = \mathrm{Gal}(L'|K)$，$H = \mathrm{Gal}(L'|L)$，$I_j$ 是 $L'|K$ 在 $v_j$ 处的惯性群，$J_j$ 是 $L'|L$ 在 $v_j$ 处的惯性群. 群 $I_j$ 是与 $p$ 互素的循环群被 $p$-群 $P_j$ 扩张得到的. 由于诸 $L_i$ 在 $R$ 上温和分歧，$I_j/J_j$ 与 $p$ 互素，故有 $P_j \subset J_j$. 从而 $H$ 含所有 $P_j$，故也含 $j$ 变动时诸 $P_j$ 生成的群 $P$. 但 $P$ 在 $G$ 中不变，因为 $G$ 的内自同构把诸 $I_j$ 彼此变换，故也把诸 $P_j$ 彼此变换. 由此得出 $P$ 是 $H$ 的在 $G$ 中正规的子群，故由于 $L'$ 是 $L$ 生成的 伽罗瓦扩张，有 $P = 1$，这证明 $L$ 在 $R$ 上温和分歧.

更一般地设 $R \to R'$ 是离散赋值环的态射，使 $R$ 的幺模化元 $\pi$ 的像是 $R'$ 的幺模化元 $\pi'$，且剩余扩张 $k(R')$ 是 $k(R)$ 的可分扩张. 设 $K$ 是 $R$ 的分式域，$K'$ 是 $R'$ 的分式域，$L$ 是 $K$-平展代数，$L' = L \otimes_K K'$. 则 $L$ 在 $R$ 上温和分歧当且仅当 $L'$ 在 $R'$ 上温和分歧. 可假设 $R$ 和 $R'$ 严格局部. 由 2.0.1 只需证明 $L$ 是域时 $L'$ 也是. 设 $\widetilde{R}$ 是 $R$ 在 $L$ 中的正规化，$\widetilde{\pi}$ 是 $\widetilde{R}$ 的幺模化元，$\widetilde{R}' = \widetilde{R} \otimes_R R'$. 由于扩张 $k(\widetilde{R})|k(R)$ 根式，且扩张 $k(R')|k(R)$ 平展，$k(\widetilde{R}) \otimes_{k(R)} k(R')$ 是域 [EGA IV 4.3.2 和 4.3.5]. 这证明 $\widetilde{R}'$ 是局部环，且由于 $\pi$ 在 $R'$ 中的像是 $\pi'$，有 $k(\widetilde{R}') = \widetilde{R}'/(\widetilde{\pi})$；从而 $\widetilde{R}'$ 是离散赋值环 \cite[第一章 §2 命题 2]{exp13-5}，故 $L'$ 是域.

2.0.3. 经归结到严格局部情形，看到\emph{温和分歧代数}的子代数和代数，两个温和分歧代数的张量积温和分歧，温和分歧代数经离散赋值环扩张后仍温和分歧，经温和分歧扩张后变成温和分歧的代数和代数温和分歧.

\refstepcounter{theorem}\noindent\textbf{\thesection.\arabic{theorem}.}\quad 设 $X$ 是 $S$-概形，$D$ 是 $X$ 上 $\geqslant 0$ 的除子. 回忆 (SGA 5 II 4.2) 称 $D$ 相对 $S$ 严格正规交叉\index[term]{yangezhengguijiaochachuzi@严格正规交叉除子}，若存在 $\Gamma(X, \mathcal{O}_X)$ 的元素的有限族 $(f_i)_{i \in I}$，使 $D = \sum_{i \in I} \mathrm{div}(f_i)$，且满足下列条件：

2.1.0. 对 $\mathrm{Supp}\, D$ 的每个点 $x$，$X$ 在 $x$ 处平展于 $S$ 上，且若用 $I(x)$ 记使 $f_i(x) = 0$ 的 $i \in I$ 的集，子概形 $V((f_i)_{i \in I(x)})$ 在 $S$ 上光滑且在 $X$ 中余维数 $\mathrm{card}\, I(x)$.

称除子 $D$ 相对 $S$ 正规交叉\index[term]{zhengguijiaochachuzi@正规交叉除子}，若它在 $X$ 上对平展拓扑局部严格正规交叉.

设 $D$ 是相对 $S$ 正规交叉除子. 令 $Y = \mathrm{Supp}\, D$，$U = X - Y$，用 $i \colon U \to X$ 记典范浸入. 对 $S$ 的每个几何点 $s$ 和几何纤维 $Y_s$ 的每个极大点 $y$，环 $R = \mathcal{O}_{X_s, y}$ 是离散赋值环.

在本号余下部分，将用下列技术定义：

\begingroup
\addtocounter{theorem}{-1}
\renewcommand{\thetheorem}{\thesection.1.1}
\begin{definition}\label{def:XIII.2.1.1}
设 $F$ 是 $U$ 上的集合层. 称 $F$ 在 $X$ 上 (沿 $D$) 相对 $S$ 温和分歧\index[term]{wenhefenqiceng@温和分歧层}，若对 $S$ 的每个几何点 $s$ 满足下列条件：
对 $Y_s$ 的每个极大点 $y$，$F$ 在 $\mathcal{O}_{X_s}$ 的分式域 $K$ 上的限制可由 $K$-平展代数 $L$ 的谱表示，后者在 $\mathcal{O}_{X_s, y}$ 上温和分歧.
\end{definition}
\endgroup

最常无混淆可能时，在术语中省略 $D$ 的提及.

\begingroup
\addtocounter{theorem}{-1}
\renewcommand{\thetheorem}{\thesection.1.2}
\begin{definition}\label{def:XIII.2.1.2}
若 $F$ 是 $U$ 上在 $X$ 上相对 $S$ 温和分歧的群层，用
\[
H^1_t(U, F)
\]
记 $H^1(U, F)$ 的由在 $X$ 上相对 $S$ 温和分歧的 $F$ 下挠子的类组成的子集\index[notations]{$H_t^1(U, F)$}.
\end{definition}
\endgroup

设
\[
\begin{tikzcd}[row sep=large, column sep=large]
U \arrow[r,"i"] \arrow[dr,"g"] & X \arrow[d,"f"] \\
& T
\end{tikzcd}
\]
是 $S$-概形的交换图式，$i$ 如 2.1 中；用
\[
R^1_t g_* F
\]
记与预层 $T' \mapsto H^1_t(U', F)$ 关联的\index[notations]{$R_t^1 g_* F$} $T$ 上层，其中 $T'$ 遍历 $T$ 上方的平展概形，$U' = U \times_T T'$；$R^1_t g_* F$ 是 $R^1 g_* F$ 的子层.

注意 $g$ 相容、$t$ 是 $T$ 的几何点、$T$ 是 $T$ 在 $t$ 处的严格局部化、$U = U \times_T T$ 时，有同构
\[
(2.1.2.1) \quad (R^1_t g_* F)_t \simeq H^1_t(U, F)
\]

2.1.3. 设 $\mathcal{C}_t((U, X)/S)$ 或简记 $\mathcal{C}_t$\index[notations]{$\mathcal{C}_t$} 是 $U$ 的在 $X$ 上相对 $S$ 温和分歧的平展覆盖范畴. 假设 $U$ 连通，$a$ 是 $U$ 的几何点. 设 $\Gamma_t$ 是把 $U$ 的在 $X$ 上相对 $S$ 温和分歧的平展覆盖 $U'$ 配给 $U'$ 的位于 $a$ 上方的几何点集的函子. 由 2.0 得出对 $(\mathcal{C}_t, \Gamma_t)$ 满足 V 4 的公理 $(G_1)$ 至 $(G_6)$. 从而 $\Gamma_t$ 可被 pro-对象表示，称之为 $(U, X)$ 在 $a$ 处带点、相对 $S$ 的\emph{温和分歧万有覆盖}\index[term]{wenhefenqiwanyoufugai@温和分歧万有覆盖}. 温和分歧万有覆盖的 $U$-自同构群的逆群称为温和分歧基本群\index[term]{wenhefenqijibenqun@温和分歧基本群}，记为
\[
\pi^t_1((U, X)/S, a)\index[notations]{$\pi_1^t((U, X)/S, a)$} \quad \text{或简记} \quad \pi^t_1(U, a) \quad \text{甚至} \quad \pi^t_1(U).
\]
它显然是基本群 $\pi_1(U, a)$ 的商 (\cref{prop:V.6.9}).

2.1.4. 设 $F$ 是 $U$ 上的群层，$P$ 是群 $F$ 的右挠子，$Q$ 是群 $F$ 的左挠子，假设 $P$ 和 $Q$ 在 $X$ 上相对 $S$ 温和分歧. 则 $P \wedge^F Q$ 也如此. 归结到证明：若 $R$ 是分式域 $K$ 的离散赋值环，$F$ 是 $K$ 上有限平展群概形，$P$ 和 $Q$ 是在 $R$ 上温和分歧的 $F$ 下两个挠子，则 $P \wedge^F Q$ 也如此. 而 $T = P \wedge^F Q$ 是 $P \times_K Q$ 的商. 若 $L$，$M$，$N$ 分别表示表示 $T$，$P$，$Q$ 的 $K$-代数，则 $L$ 是 $M \otimes_K N$ 的子代数，由 2.0.3 得出 $L$ 在 $R$ 上温和分歧.

由前面得出：若 $F$ 是 $U$ 上的群层，且存在 $F$ 下的在 $X$ 上相对 $S$ 温和分歧的挠子，则 $F$ 在 $X$ 上相对 $S$ 温和分歧. 确实 $P$ 的逆挠子 $P^{\circ}$ 在 $X$ 上相对 $S$ 温和分歧，因为作为集合层它同构于 $P$. 若 $P F$ 是 $F$ 被 $P$ 扭转的群，有同构
\[
F \simeq P^{\circ} \wedge^F P,
\]
从而 $F$ 在 $X$ 上相对 $S$ 温和分歧.

如前看到，若 $F \to F'$ 是 $U$ 上两个在 $X$ 上相对 $S$ 温和分歧的群层的态射，$P$ 是在 $X$ 上相对 $S$ 温和分歧的 $F$ 下挠子，则由 $P$ 经结构群扩张 $F \to F'$ 得到的挠子 $P'$ 在 $X$ 上相对 $S$ 温和分歧.

特别，典范态射
\[
H^1(U, F) \longrightarrow H^1(U, F')
\]
经限制到 $H^1_t(U, F)$ 给出典范态射
\[
H^1_t(U, F) \longrightarrow H^1_t(U, F')
\]

2.1.5. 设 $S' \to S$ 是态射，用 $U'$ (resp. $X'$，等等) 记 $U$ (resp. $X$，等等) 在 $S'$ 上的逆像. 若 $F$ 是 $U$ 上在 $X$ 上相对 $S$ 温和分歧的集合层，由\cref{def:XIII.2.1.1} 和 2.0.3 得出 $F'$ 在 $X'$ 上相对 $S'$ 温和分歧.

若 $F$ 是 $U$ 上的群层，$F$ 下在 $X$ 上相对 $S$ 温和分歧的挠子在 $S'$ 上的逆像是 $F'$ 下在 $X'$ 上相对 $S'$ 温和分歧的挠子. 特别有典范函子
\[
(2.1.5.1) \quad \mathcal{C}_t((U, X)/S) \longrightarrow \mathcal{C}_t((U', X')/S').
\]
假设 $U$ 和 $U'$ 连通，$a$ 是 $U$ 的几何点，$a'$ 是 $U'$ 的位于 $a$ 上方的几何点；由前面得出典范态射
\[
(2.1.5.2) \quad \pi^t_1(U', a') \longrightarrow \pi^t_1(U, a).
\]
若 $S' \to S$ 是态射，$h \colon T' \to T$ 是典范投影，态射
\[
h^{*}(R^1 g_* F) \longrightarrow R^1 g'_* F'
\]
经限制给出典范态射
\[
(2.1.5.3) \quad h^{*}(R^1_t g_* F) \longrightarrow R^1_t g'_* F'
\]

2.1.6. 设 $F$ 是 $U$ 上在 $X$ 上相对 $S$ 温和分歧的群层. 记号同 \ref{def:XIII.2.1.2}，有典范正合序列：
\[
(2.1.6.1) \quad 1 \longrightarrow H^1(X, i_* F) \longrightarrow H^1_t(U, F) \longrightarrow H^0(X, R^1_t i_* F)
\]
\[
1 \longrightarrow R^1 f_*(i_* F) \longrightarrow R^1_t g_* F \longrightarrow f_*(R^1_t i_* F).
\]
第一个由正合序列 (SGA 4 III 3.2) 得出：
\[
1 \longrightarrow H^1(X, i_* F) \longrightarrow H^1(U, F) \longrightarrow H^0(X, R^1 i_* F).
\]
只需证明 $H^1(X, i_* F)$ 在 $H^1(U, F)$ 中的像实际包含在 $H^1_t(U, F)$ 中，且 $H^1_t(U, F)$ 在 $H^0(X, R^1 i_* F)$ 中的像实际包含在 $H^0(X, R^1_t i_* F)$ 中. $i_* F$ 下的挠子在 $U$ 上的逆像是 $i^{*} i_* F$ 下的挠子，后者显然在 $X$ 上相对 $S$ 温和分歧，故经结构群扩张 $i^{*} i_* F \to F$ 后仍如此，这证明箭 $H^1(X, i_* F) \to H^1_t(U, F)$ 存在. $H^1_t(U, F)$ 在 $H^0(X, R^1 i_* F)$ 中的像包含在 $H^0(X, R^1_t i_* F)$ 中的事实由 $R^1_t i_* F$ 的定义立即得出. 这证明第一个正合序列存在，第二个由第一个经局部化得出.

\refstepcounter{theorem}\noindent\textbf{\thesection.\arabic{theorem}.}\quad 沿用 2.1 的记号. 当 $U$ 上的叠 $\Phi$ 局部于 $X$ 上和 $S$ 上 (对平展拓扑) 是 $S$ 上叠 $\Psi$ 的逆像时，定义它的温和分歧对象概念.

先设 $G$ 是 $U$ 上的叠丛，假设给定平展满射态射 $S_1 \to S$、平展满射态射 $X_2 \to X \times_S S_1$、$S_1$ 上的平凡叠丛 $H$ 和同构
\[
G|_{U_2} \longrightarrow H|_{U_2},
\]
其中 $U_1 = U \times_X X_1$，$U_2 = U \times_X X_2$. 选 $H|_{X_2}$ 的平凡化时，上面同构把 $G|_{U_2}$ 等同于群层 $F$ 下的挠子叠. 称 $G_U$ 的元素 $x$ 在 $X$ 上相对 $S$ 温和分歧，若 $x$ 在 $U_2$ 上的限制是在 $X$ 上相对 $S$ 温和分歧的挠子. 由 2.1.4 这个概念不依赖 $H|_{X_2}$ 的平凡化方式.

现在设 $\Phi$ 是 $U$ 上的叠，假设给定平展满射态射 $S_1 \to S$、平展满射态射 $X_2 \to X \times_S S_1$、$S_1$ 上的叠 $\Psi$ 和同构
\[
i \colon \Phi|_{U_2} \longrightarrow \Psi|_{U_2}.
\]
设 $x$ 是 $\Phi_U$ 的元素，$G_x$ 是 $x$ 生成的 $\Phi$ 的极大子叠丛 \cite[III 2.1.7]{exp13-1}，$S_\Phi$ 是 $\Phi$ 的极大子叠丛层. 同构 $i$ 诱导同构
\[
S_\Phi|_{U_2} \longrightarrow S_\Psi|_{U_2}.
\]
由 \ref{cor:XIII.5.7} 得出，必要时把 $S_1$ 换成平展满射扩张，有 $\Psi$ 的唯一极大子叠丛 $H$ (可假设平凡)，使 $i$ 定义同构
\[
G_x|_{U_2} \longrightarrow H|_{U_2}.
\]
称元素 $x$ 在 $X$ 上相对 $S$ 温和分歧\index[term]{diedewenhefenqiyuansu@叠的温和分歧元素}，若它作为配备上述同构的 $G_x$ 的元素如此.

2.2.1. 设 $\Phi$ 是 $U$ 上的叠，局部于 $X$ 上和 $S$ 上是 $S$ 上叠的逆像，设
\[
\begin{tikzcd}[row sep=large, column sep=large]
U \arrow[r,"i"] \arrow[dr,"g"] & X \arrow[d,"f"] \\
& T
\end{tikzcd}
\]
是如 \ref{def:XIII.2.1.2} 的图式. 对每个平展于 $T$ 上的概形 $T'$，若 $U' = U \times_T T'$，考虑 $(g_* \Phi)_{T'} = \Phi_{U'}$ 的由 $\Phi_{U'}$ 的在 $X$ 上相对 $S$ 温和分歧的元素组成的子集 $(g_{*t} \Phi)_{T'}$\index[notations]{$(g_{*t}\Phi)_{T'}$}. 称 $\Phi$ 经 $g$ 的温和分歧直像\index[term]{diedewenhefenqizhixiang@叠的温和分歧直像}，记为
\[
g_{*t} \Phi
\]
是 $g_* \Phi$ 的由平展于 $T$ 上方的概形 $T'$ 上方的对象为 $(g_{*t} \Phi)_{T'}$ 的元素组成的全子范畴. $g_{*t} \Phi$ 显然是 $g_* \Phi$ 的子叠.

2.2.2. 经归结到挠子叠的情形，看到若 $\Phi$ 是 $U$ 上局部于 $S$ 和 $X$ 的平展拓扑是 $S$ 上叠的逆像的叠，典范态射
\[
h^{*}(g_* \Phi) \longrightarrow g'_* \Phi'
\]
经限制给出典范态射
\[
h^{*}(g_{*t} \Phi) \longrightarrow g'_{*t} \Phi'
\]

\begin{remark}\label{rem:XIII.2.3}
a) 若 $F$ 是 $U$ 上局部常可构集合层，为使 $F$ 在 $X$ 上相对 $S$ 温和分歧，只需 \ref{def:XIII.2.1.1} 的条件对 $S$ 的位于 $S$ 的极大点上方的几何点满足. 可假设除子 $D$ 严格正规交叉. 层 $F$ 可由 $U$ 的平展覆盖 $V$ 表示. 若 $s$ 是 $S$ 的几何点，$y$ 是 $Y_s$ 的极大点，记 $S$ 为 $S$ 在 $s$ 处的严格局部化，$X$ 为 $X$ 在 $y$ 处的严格局部化，$U = U \times_X X$，$V = V \times_X X$. 若 \ref{def:XIII.2.1.1} 的条件对位于 $S$ 的极大点上方的几何点满足，由下面 \ref{prop:XIII.5.5} 得出 $V$ 是在 $X$ 上相对 $S$ 温和分歧的 $U$ 的覆盖；从而 $V$ 是在 $X$ 上相对 $S$ 温和分歧的 $U$ 的平展覆盖.

b) 设 $F$ 是 $U$ 上在 $X$ 上相对 $S$ 温和分歧的群层. 若 $s$ 是 $S$ 的几何点，$y$ 是 $Y_s$ 的极大点，记 $K$ 为 $\mathcal{O}_{X_s, y}$ 的分式域. 假设对每个点 $s$ 和每个点 $y$，谱表示 $F|_K$ 的 $K$-代数 $L$ 的秩与 $\mathcal{O}_{X_s}$ 的剩余特征 $p$ 互素. 有时滥用语言说 $F$ 与 $S$ 的剩余特征互素\index[term]{yushengyutezhenghusu@与剩余特征互素}. 此时，$F$ 下的每个挠子 $P$ 在 $X$ 上相对 $S$ 温和分歧.

设 $R$ 是 $\mathcal{O}_{X_s, y}$ 在 $y$ 处的严格局部化，$K$ 是它的分式域，$F$ 是 $F$ 在 $K$ 上的逆像. 证明可假设 $F$ 常. 由于 $F$ 在 $X$ 上相对 $S$ 温和分歧，$F$ 可由 $K$-代数 $L = \prod L_i$ 的谱表示，其中诸 $L_i$ 是 $K$ 的与 $p$ 互素的阶扩张. 可以找到 $K$ 的与 $p$ 互素的阶扩张 $K'$，使 $F|_{K'}$ 是常层. 由 2.0.3，为证明 $P|_K$ 在 $R$ 上温和分歧，只需看 $P|_{K'}$ 在 $R$ 在 $K'$ 中的整闭包上温和分歧，归结到 $F$ 常的情形. 此后假设 $F$ 常. 表示 $P|_K$ 的 $K$-代数 $H$ 是彼此同构的 $K$ 的扩张 $H_i$ 的积. 由于 $H$ 的秩与 $p$ 互素，$[H_1 : K]$ 也如此，这证明 $H$ 在 $X$ 上相对 $S$ 温和分歧.

c) 设 $X$ 是正则概形，$D$ 是 $X$ 的\emph{正规交叉除子} (SGA 5 I 3.1.5)，$U = X - \mathrm{Supp}\, D$，$F$ 是 $U$ 上的集合层. 若 $y$ 是 $\mathrm{Supp}\, D$ 的极大点，记 $K$ 为 $\mathcal{O}_{X, y}$ 的分式域. 称 $F$ 相对 $D$ 温和分歧\index[term]{wenhefenqi@温和分歧}，若对 $\mathrm{Supp}\, D$ 的每个极大点 $y$，$F|_K$ 可由在 $\mathcal{O}_{X, y}$ 上温和分歧的 $K$-代数表示.
\end{remark}

\begin{theorem}\label{thm:XIII.2.4}
设 $f \colon X \to S$ 是 $S$-概形，$D$ 是 $X$ 上相对 $S$ 正规交叉 (2.1) 的除子，$Y = \mathrm{Supp}\, D$，$U = X - Y$，$i \colon U \to X$ 是典范浸入. 设 $F$ 是 $U$ 上满足下列条件之一的集合层 (resp. 群层)：
\begin{enumerate}[a)]
\item $F$ 局部于 $X$ 和 $S$ 的平展拓扑是 $S$ 上集合层 (resp. 可构群层) 的逆像.
\item $F$ 在 $U$ 上局部常可构，且在 $X$ 上相对 $S$ 温和分歧.
\end{enumerate}
则有下列结论：
\begin{enumerate}[1)]
\item $(F, i)$ 相对 $S$ 在维数 $\leqslant 0$ 时上同调真 (resp. 对每个态射 $h \colon S' \to S$，若 $i' \colon U' \to X'$ 是 $i$ 在 $S'$ 上的逆像，$F' = F|_{U'}$，$k = h_{(X)}$，典范态射
\[
\Psi \colon k^{*}(R^1_t i_* F) \longrightarrow R^1_t i'_* F'
\]
是同构).
若 $F$ 是与 $S$ 的剩余特征互素 (2.3.b)) (在 $X$ 上相对 $S$ 温和分歧) 的群层，则 $(F, i)$ 相对 $S$ 在维数 $\leqslant 1$ 时上同调真.
\item 若 $F$ 是可构集合层 (resp. 群层)，$i_* F$ (resp. $R^1_t i_* F$) 可构.
\end{enumerate}
\end{theorem}

\begin{proof}
对每个 $S$-概形 $S'$，考虑下列所有方形笛卡儿的图式：
\[
\begin{tikzcd}[row sep=large, column sep=large]
U \arrow[d,"{i}"] & U' \arrow[l,"g"] \arrow[d,"{i'}"] \\
X \arrow[d,"{f}"] & X' \arrow[l,"k"] \arrow[d,"{f'}"] \\
S & S' \arrow[l,"h"]
\end{tikzcd}
\]
由于问题局部于 $X$ 的平展拓扑，可假设 $D$ 是相对 $S$ \emph{严格正规交叉除子} (2.1)；而且，把 $X$ 限制到 $Y$ 的邻域后可假设 $X$ 在 $S$ 上光滑.

**\ref{thm:XIII.2.4} 1) 的证明**

2.4.1. 满足 a) 的集合层的情形. 可假设 $F = g^{*} G$，其中 $G$ 是 $S$ 上的层. 由 (SGA 4 XVI 3.2) 得出典范态射
\[
(2.4.1.1) \quad f^{*} G \longrightarrow i_* F
\]
是同构. 对每个 $S$-概形 $h \colon S' \to S$，同样有同构 $f'^{*} G' \to i'_* F'$；从而典范态射
\[
\phi \colon k^{*}(i_* F) \longrightarrow i'_* F'
\]
等同于自然同构
\[
k^{*} f^{*} G \simeq f'^{*} G'.
\]

2.4.2. 满足 b) 的集合层的情形. 需证明 $\phi$ 是同构，只需看它在 $X'$ 的每个几何点 $x'$ 处如此. 设 $S$ (resp. $X$，resp. $S'$，resp. $X'$) 是 $S$ (resp. $X$，resp. $S'$，resp. $X'$) 在 $x'$ 处的严格局部化，令 $U = U_{(X)}$，$U' = U_{(X')}$，等等. 态射 $\phi_x$ 等同于典范态射
\[
\phi \colon H^0(U, F) \longrightarrow H^0(U', F')
\]
可以找到 $U$ 的 \ref{prop:XIII.5.4} 中出现类型的主覆盖 $V$，使 $F$ 在 $V$ 上的逆像是取值 $C$ 的常层. 若 $\Pi$ 是 $V$ 在 $U$ 上的 伽罗瓦群，$\Pi$ 作用在 $F|_V$ 上，有
\[
(2.4.2.1) \quad H^0(U, F) \simeq H^0(V, C_V)^{\Pi}
\]
其中第二个成员表示 $H^0(V, C_V)$ 的在 $\Pi$ 下不变的元素的集. 由于 $V' = V \times_U U'$ 是 $U'$ 的群 $\Pi' \simeq \Pi$ 的主覆盖，看到态射 $\phi$ 由典范态射
\[
H^0(V, C_V) \longrightarrow H^0(V', C_{V'})
\]
取 $\Pi$ 的不变元得到. 由于 $V$ 和 $V'$ 连通 (\cref{prop:XIII.5.4})，该态射从而 $\phi$ 是同构.

注意若 $F$ 还是群层，$P$ 是 $U$ 上在 $X$ 上相对 $D$ 温和分歧的群 $F$ 的挠子，由前面证明和 2.2 得出 $(P F, i)$ 相对 $S$ 在维数 $\leqslant 0$ 时上同调真.

2.4.3. 群层的情形. 为证明 $\Psi$ 是同构，只需证明对 $Y'$ 的每个几何点 $y'$，态射
\[
\Psi_{y'} \colon (k^{*}(R^1_t i_* F))_{y'} \longrightarrow (R^1_t i'_* F')_{y'}
\]
是同构. 由 \ref{def:XIII.2.1.2}，$\Psi_{y'}$ 等同于典范态射
\[
\Psi \colon H^1_t(U, F) \longrightarrow H^1_t(U', F').
\]
设 $\widetilde{U}$ 是 $U$ 的温和分歧万有覆盖 (2.1.3)，$\widetilde{F}$ 是 $F$ 在 $\widetilde{U}$ 上的逆像. 由情形 a) 的 \ref{cor:XIII.5.7} 和情形 b) 的 \ref{prop:XIII.5.5} 得出，$H^1_t(U, F)$ 等同于 $H^1(U, F)$ 的在 $\widetilde{U}$ 上的逆像平凡的 $F$-挠子类组成的子集. 另一方面，经典推理 (参见 IX 5，第 231 页) 表明在 $\widetilde{U}$ 上的逆像平凡的 $H^1(U, F)$ 的元素集等同于 $H^1(\pi^t_1(U), H^0(\widetilde{U}, \widetilde{F}))$. 由此得到典范同构
\[
(2.4.3.1) \quad H^1_t(U, F) \xrightarrow{\sim} H^1(\pi^t_1(U), H^0(\widetilde{U}, \widetilde{F})).
\]
从而态射 $\Psi$ 等同于典范态射
\[
H^1(\pi^t_1(U), H^0(\widetilde{U}, \widetilde{F})) \longrightarrow H^1(\pi^t_1(U'), H^0(\widetilde{U}', \widetilde{F}')).
\]
证明该态射是同构. 态射 $\pi^t_1(U') \to \pi^t_1(U)$ 由 \ref{cor:XIII.5.6} 是同构，态射 $H^0(\widetilde{U}, \widetilde{F}) \to H^0(\widetilde{U}', \widetilde{F}')$ 也如此. 确实，在情形 b) 中显然，因为 $F$ 常且 $\widetilde{U}$ 和 $\widetilde{U}'$ 连通. 在情形 a) 中，设 $G$ 是 $S$ 上使 $F = g^{*} G$ 的可构群层；由于态射 $\widetilde{U} \to S$ 和 $\widetilde{U}' \to S'$ 0-无环 (\cref{cor:XIII.5.7})，有
\[
H^0(\widetilde{U}, \widetilde{F}) \simeq H^0(S, G) \simeq H^0(S', G') \simeq H^0(\widetilde{U}', \widetilde{F}'),
\]
这蕴含 $\Psi$ 是同构. \ref{thm:XIII.2.4} 1) 的最后一个断言由前面，考虑到 \ref{rem:XIII.2.3} b) 得出.

**\ref{thm:XIII.2.4} 2) 的证明** — 满足 a) 的可构集合层情形由 (2.4.1.1) 立即得出. 设 $F = g^{*} G$ 是满足 a) 的群层，其中 $G$ 可构；可假设 $S$ 仿射；设 $(S_j)_{j \in J}$ 是 $S$ 的既约闭子概形的有限族，其并覆盖 $S$，使 $G$ 在 $S_j$ 上的逆像是局部常层. 考虑到 \ref{thm:XIII.2.4} 1)，为建立 $R^1_t i_* F$ 可构，只需看它在基变换 $S_j \to S$ ($j \in J$) 后如此. 归结到 $F$ 局部常的情形 b).

此后假设 $F$ 是满足 b) 的集合层或群层. 由于问题局部于 $X$ 的平展拓扑，可假设 $X$ 在 $S$ 上有限表现，经极限过渡可假设 $X$ 和 $S$ 诺特.

设 $D = \sum_{1 \leqslant i \leqslant r} \mathrm{div}\, f_i$，其中对 $\mathrm{Supp}\, D$ 的每个点 $x$，若 $I(x)$ 是使 $f_i(x) = 0$ 的 $i$ 的集，子概形 $V((f_i)_{i \in I(x)})$ 在 $S$ 上光滑且在 $X$ 中余维数 $\mathrm{card}\, I(x)$. 设 $P$ 是 $[1, r]$ 的部分集，对每个 $I \in P$，令
\[
X_I = \Big(\bigcap_{i \in I} V(f_i)\Big) \cap \Big(\bigcap_{i \notin I} X_{f_i}\Big).
\]
设 $z$ 是 $X_I$ 的点. 先限制到 $z$ 的平展邻域，可以找到 $X$ 的含 $z$ 的开集 $W$ 和 $U \cap W$ 的在 $W$ 上相对 $S$ 温和分歧、\ref{rem:XIII.5.6.1} 中考虑类型的主覆盖 $V$，使 $F$ 在 $V$ 上的逆像是取值 $C$ 的常层. 设 $\pi$ 是该覆盖 $V$ 的 伽罗瓦群. 对 $X_I$ 的每个几何点 $x$，由 (2.4.2.1) 有
\[
H^0(U, F) \simeq H^0(V, C_V)^{\pi}.
\]
由此得出 $i_* F|_{X_I \cap W}$ 局部常 (SGA 4 IX 2.13)，从而 $i_* F$ 可构.

最后证明 $F$ 是局部常群层时 $R^1_t i_* F|_{X_I}$ 可构. 若 $x$ 是 $X$ 的几何点，在 (2.4.3.1) 中得到表达式
\[
H^1_t(U, F) \xrightarrow{\sim} H^1(\pi^t_1(U), H^0(\widetilde{U}, \widetilde{F})).
\]
若 $p$ 是 $X$ 的剩余特征，由 \ref{cor:XIII.5.6} 有 $\pi^t_1(U) = \prod_{\ell \neq p} \mathbb{Z}_\ell(1)^{\mathrm{card}\, I}$. 设 $L$ 是整除有限群 $H^0(\widetilde{U}, \widetilde{F})$ 的阶的素数集，令 $K = \prod_{\ell \in L - \{p\} \cap L} \mathbb{Z}_\ell(1)^{\mathrm{card}\, I}$. 由 \cite[第一章 §5 例 2]{exp13-4} 有
\[
H^1_t(U, F) \simeq H^1(K, H^0(\widetilde{U}, \widetilde{F})).
\]
由于 $K$ 拓扑有限型且 $H^0(\widetilde{U}, \widetilde{F})$ 有限，首先得出层 $R^1_t i_* F|_{X_I}$ 的纤维有限. 另一方面，$L$ 不依赖点 $x$. 对每个 $q \in L$，设 $X_{I,q}$ 是 $X_I$ 的方程 $q = 0$ 的闭集，$X_{I'}$ 是 $X_I$ 的诸 $X_{I,q}$ 的并的余集的开集. 则 $R^1_t i_* F|_{X_{I,q}}$ 和 $R^1_t i_* F|_{X_{I'}}$ 局部常；确实，$X_{I,q}$ (resp. $X_{I'}$) 的几何点的特殊化箭对群 $K$ (\cref{rem:XIII.5.6.1}) 从而对集合 $H^1_t(U, F)$ 诱导同构，且可以应用 SGA 4 IX 2.13.
\end{proof}

\begin{corollary}\label{cor:XIII.2.5}
设 $f \colon X \to S$ 是态射，$D$ 是 $X$ 上相对 $S$ 正规交叉 (2.1) 的除子，$Y = \mathrm{Supp}\, D$，$U = X - Y$，$i \colon U \to X$ 是典范浸入. 设 $\Phi$ 是 $U$ 上的叠，假设给定平展满射态射 $S_1 \to S$ 和 $X_2 \to X \times_S S_1$，$S_1$ 上的叠 $\Psi$ 和同构 $\Phi|_{U_2} \simeq \Psi|_{U_2}$ (参见 2.2).
则对每个态射 $h \colon S' \to S$，若 $k = h_{(X)}$，典范函子
\[
\phi \colon k^{*} i_* \Phi \longrightarrow i'_* \Phi'
\]
忠实全. 若 $\Psi$ 可构，典范函子
\[
\psi \colon k^{*} i_{*t} \Phi \longrightarrow i'_{*t} \Phi'
\]
是范畴等价.
而且，若叠 $\Psi$ 可构 (resp. $\Psi$ 1-可构 (0) 且 $S$ 局部诺特)，$i_* \Phi$ 可构 (resp. $i_{*t} \Phi$ 1-可构).
\end{corollary}

证明 $\phi$ 忠实全. 只需看对 $X'$ 的每个几何点 $x'$，函子
\[
\phi \colon \Phi(U) \longrightarrow \Phi'(U')
\]
如此 (沿用 2.4.2 的记号). 设 $a$，$b$ 是 $\Phi(U)$ 的两个元素，$a'$，$b'$ 是它们在 $\phi$ 下的像. 由于态射 $U \to S$ 局部 0-无环 (\cref{cor:XIII.5.7})，有同构
\[
H^0(U, S_\Phi) \xrightarrow{\sim} H^0(S, S_\Psi).
\]
从而 $a$ 和 $b$ 由 $\Psi$ 的元素经逆像得到，故 $F = \mathrm{Hom}_U(a, b)$ 也如此. 由于层 $F' = \mathrm{Hom}_{U'}(a', b')$ 是 $F$ 在 $U'$ 上的逆像，由 2.4.1 得出典范态射
\[
H^0(U, F) \longrightarrow H^0(U', F')
\]
是同构，这证明 $\phi$ 忠实全.

证明对 $X'$ 的每个几何点 $x'$，函子
\[
\psi \colon i_{*t} \Phi(X') \longrightarrow i'_{*t} \Phi'(X')
\]
是等价. 由前面 $\psi$ 忠实全. 证明 $\psi$ 本质满射. 设 $a'$ 是 $\Phi'(U')$ 的在 $X'$ 上相对 $S'$ 温和分歧 (2.2) 的元素，证明它是 $\Phi(U)$ 的温和分歧元素的像. 由 \ref{thm:XIII.2.4} 1) 得出典范态射
\[
H^0(U, S_\Phi) \longrightarrow H^0(U', S_{\Phi'})
\]
是同构. 设 $G'$ 是 $a'$ 生成的 $\Phi'$ 的极大子叠丛；则存在 $\Phi$ 的极大子叠丛 $G$ (某个 $S$ 上叠的逆像)，使
\[
m^{*} G \simeq G',
\]
其中 $m$ 是态射 $U' \to U$. 典范函子
\[
(\ast) \quad k^{*} i_{*t} G \longrightarrow i'_{*t} G'
\]
是等价，因为 $G$ 等同于来自 $S$ 的可构群层下的挠子叠丛，且可以应用 \ref{thm:XIII.2.4} 1). 由 $(\ast)$ 得出存在 $G(U)$ 的在 $X$ 上相对 $S$ 温和分歧的元素 $a$，使它在 $U'$ 上的逆像是 $a'$，这证明 $\psi$ 是等价.

$\Psi$ 可构时 $i_* \Phi$ 也可构；确实 $i_* \Phi$ 的对象 $x$ 局部于 $S$ 的平展拓扑是 $\Psi$ 的对象 $y$ 的逆像；由 2.4.1.1 得出 $\mathrm{Aut}(x)$ 是 $\mathrm{Aut}(y)$ 的逆像，故可构. 最后，$\Psi$ 1-可构时由下面 \ref{cor:XIII.6.3} 得出 $i_{*t} \Phi$ 也 1-可构.

\begin{corollary}\label{cor:XIII.2.6}
记号同 2.4. 假设 $S$ 在使 $Y_s \neq \varnothing$ 的每个点 $s$ 处特征零. 则若 $F$ 是 $U$ 上局部常可构群层 (resp. $U$ 上局部于 $X$ 和 $S$ 是 $S$ 上\emph{可构叠}的逆像的可构叠)，$(F, i)$ 相对 $S$ 在维数 $\leqslant 1$ 时上同调真.
\end{corollary}

由于 $U$ 上每个可构集合层在 $X$ 上相对 $S$ 温和分歧，该推论由 \ref{thm:XIII.2.4} (resp. \ref{cor:XIII.2.5}) 得出.

\begin{corollary}\label{cor:XIII.2.7}
记号同 \ref{thm:XIII.2.4}，但再给定 $S$-概形 $T$、紧合态射 $p \colon X \to T$，并假设 $X$ 和 $T$ 在 $S$ 上有限表现；设 $q = pi$. 设 $F$ 是 $U$ 上满足 \ref{thm:XIII.2.4} 的条件 a) 或 b) 之一的可构集合层 (resp. 满足 \ref{thm:XIII.2.4} 的条件 a) 或 b) 之一的群层，resp. $U$ 上局部于 $X$ 和 $S$ 是 $S$ 上可构叠 $G$ 的逆像的叠). 则有下列结论：
\begin{enumerate}[1)]
\item $(F, q)$ 相对 $S$ 在维数 $\leqslant 0$ 时上同调真 (resp. 对每个态射 $h \colon S' \to S$，若 $m = h_{(T)}$，典范态射：
\[
\Theta \colon m^{*}(R^1_t q_* F) \longrightarrow R^1_t q'_* F'
\]
是同构，resp. 对每个态射 $S' \to S$，典范态射
\[
\xi \colon m^{*}(q_{*t} F) \longrightarrow q'_{*t} F'
\]
是等价).
\item 层 $q_* F$ (resp. 层 $R^1_t q_* F$，resp. 叠 $q_{*t} F$) 可构. 在最后情形，若假设 $S$ 局部诺特且 $G$ 1-可构，$q_{*t} F$ 也 1-可构.
\end{enumerate}
\end{corollary}

第一部分由 \ref{thm:XIII.2.4}，\ref{cor:XIII.2.5} 和 \ref{cor:XIII.1.8} 的证明立即得出. 证明 2). 若 $F$ 是 $U$ 上满足 \ref{thm:XIII.2.4} a) 或 \ref{thm:XIII.2.4} b) 的可构集合层，由 \ref{thm:XIII.2.4} 2) 得出 $i_* F$ 可构；故 $q_* F = p_*(i_* F)$ 也可构 (SGA 4 XIV 1.1).

设 $F$ 是 $U$ 上满足 \ref{thm:XIII.2.4} a) 或 \ref{thm:XIII.2.4} b) 的可构群层，证明 $R^1_t q_* F$ 可构. 经极限过渡 (EGA IV 8.10.5 和 17.7.8) 并用 1)，可假设 $S$ 诺特. 设 $\Phi$ 是 $X$ 上的叠，它在每个平展于 $X$ 上的概形 $X'$ 处的纤维由 $U' = U \times_X X'$ 上在 $X'$ 上相对 $S$ 温和分歧的群 $F|_U$ 的挠子组成；故有
\[
S(i_{*t} \Phi) \simeq R^1_t i_* F
\]
且由 \ref{thm:XIII.2.4} 2) 该层可构. 由下面 \ref{cor:XIII.6.3} 得出 $S(p_* \Phi)$ 可构，i.e. $R^1_t q_* F$ 可构.

最后，若 $F$ 是 $U$ 上局部于 $X$ 和 $S$ 是 $S$ 上可构叠的逆像的叠，$i_{*t} F$ 可构，故 $q_{*t} F = p_*(i_{*t} F)$ 也可构. 若还假设 $S$ 局部诺特且 $S_G$ 可构，由 \ref{cor:XIII.6.3} 得出 $S(i_{*t} F)$ 可构；故 $S(q_* i_{*t} F)$，i.e. $S(q_{*t} F)$ 也可构 (\cref{cor:XIII.6.2}).

\begin{corollary}\label{cor:XIII.2.8}
设
\[
\begin{tikzcd}[row sep=large, column sep=large]
U \arrow[r,"i"] \arrow[dr,"g"] & X \arrow[d,"f"] \\
& S
\end{tikzcd}
\]
是概形的交换图式，其中 $U$ 是 $X$ 中相对 $S$ 正规交叉除子的余集的开集，$f$ 是有限表现紧合态射. 设 $L$ 是素数集. 假设 $g$ 局部 0-无环 (resp. 对 $L$ 局部 1-非球). 则若 $F$ 是 $U$ 上局部常可构、在 $X$ 上相对 $S$ 温和分歧的集合层 (resp. $L$-群层)，$f_* F$ (resp. $R^1_t f_* F$) 局部常可构，且 $(F, f)$ 相对 $S$ 在维数 $\leqslant 0$ 时上同调真 (resp. $R^1_t f_* F$ 的构造与任何基变换 $S' \to S$ 交换). 在非 resp. 情形，若 $F$ 是群层，对 $S$ 的几何点的每个特殊化 $s_1 \to s_2$，特殊化态射
\[
(R^1_t f_* F)_{s_2} \longrightarrow (R^1_t f_* F)_{s_1}
\]
单射.
\end{corollary}

该推论由 \ref{cor:XIII.2.6} 和 \ref{prop:XIII.1.14} (resp. $R^1_t f_* F$ 的 \ref{prop:XIII.1.14} 类似物，按 loc. cit. 证明) 立即得出.

\begin{corollary}\label{cor:XIII.2.9}
设
\[
\begin{tikzcd}[row sep=large, column sep=large]
U \arrow[r,"i"] \arrow[dr,"g"] & X \arrow[d,"f"] \\
& S
\end{tikzcd}
\]
是概形的交换图式，其中 $U$ 是 $X$ 中相对 $S$ 正规交叉除子的余集的开集，$f$ 是有限表现紧合光滑态射. 设 $L$ 是与 $S$ 的剩余特征不同的素数集. 设 $F$ 是 $U$ 上局部常可构、在 $X$ 上相对 $S$ 温和分歧的 $L$-群层. 则 $R^1 f_* F$ 局部常可构，且 $(F, f)$ 相对 $S$ 在维数 $\leqslant 1$ 时上同调真.
\end{corollary}

该推论由 \ref{cor:XIII.2.8} 和 $R^1_t f_* F = R^1 f_* F$ (2.3 b)) 的事实得出.

\refstepcounter{theorem}\noindent\textbf{\thesection.\arabic{theorem}.}\quad 若 $U$ 是连通概形，$a$ 是 $U$ 的几何点，$L$ 是素数集，记
\[
(2.10.0) \quad \pi^L_1(U, a)\index[notations]{$\pi_1^L(U, a)$}
\]
为 $\pi_1(U, a)$ 的阶的所有素因子都在 $L$ 中的有限商的投射极限.

将定义基本群的特殊化态射，推广 X.2.

设 $g \colon U \to S$ 是纤维几何连通的相容态射 (resp. $g = f i$ 形式的态射，其中 $f \colon X \to S$ 是有限表现紧合态射，$i \colon U \to X$ 是使 $U$ 为 $X$ 中相对 $S$ 正规交叉除子的余集的开浸入 (参见 \ref{cor:XIII.2.8})). 设 $L$ 是素数集，在非 resp. 情形假设对每个有限常 $L$-群 $C$，$(C_U, g)$ 相对 $S$ 在维数 $\leqslant 1$ 时上同调真. 设 $s_1 \to s_2$ 是 $S$ 的几何点的特殊化态射，$S$ 是 $S$ 在 $s_2$ 处的严格局部化，$U = U \times_S S$. 有交换图式
\[
\begin{tikzcd}[row sep=large, column sep=large]
U_{s_1} \arrow[d] & U \arrow[l,"{h_1}"] \arrow[r,"{h_2}"] \arrow[d,"{g}"] & U_{s_2} \arrow[d] \\
s_1 & S \arrow[l] \arrow[r] & s_2.
\end{tikzcd}
\]
若 $a_1$ 是 $U_{s_1}$ 的几何点，$a_2$ 是 $U_{s_2}$ 的几何点，态射 $h_1$ 和 $h_2$ 定义典范态射
\[
\pi_1 \colon \pi^L_1(U_{s_1}, a_1) \longrightarrow \pi^L_1(U, a_1), \quad \pi_2 \colon \pi^L_1(U_{s_2}, a_2) \longrightarrow \pi^L_1(U, a_2)
\]
(resp. $\pi_1 \colon \pi^t_1(U_{s_1}, a_1) \longrightarrow \pi^t_1(U, a_1)$，$\pi_2 \colon \pi^t_1(U_{s_2}, a_2) \longrightarrow \pi^t_1(U, a_2)$) (V 7 和 2.1.5.2). 上同调真性假设 (resp. \ref{cor:XIII.2.8}) 证明 $\pi_2$ 是同构. 若选 $a_1$ 到 $a_2$ 的道路类，得到同构
\[
\pi_{12} \colon \pi^L_1(U, a_1) \xrightarrow{\sim} \pi^L_1(U, a_2) \quad \text{(resp. } \pi_{12} \colon \pi^t_1(U, a_1) \xrightarrow{\sim} \pi^t_1(U, a_2)\text{)}
\]
从而有态射 $\pi = \pi_2^{-1} \pi_{12} \pi_1$
\[
\pi \colon \pi^L_1(U_{s_1}, a_1) \longrightarrow \pi^L_1(U_{s_2}, a_2) \quad \text{(resp. } \pi \colon \pi^t_1(U_{s_1}, a_1) \longrightarrow \pi^t_1(U_{s_2}, a_2)\text{)}.
\]
换 $a_1$ 到 $a_2$ 的道路类，等价于用 $\pi^L_1(X_{s_2}, a_2)$ (resp. $\pi^t_1(X_{s_2}, a_2)$) 的内自同构修改 $\pi$. 称上述定义的一个态射为与特殊化态射 $s_1 \to s_2$ 关联的基本群特殊化态射，简记为
\[
\pi \colon \pi^L_1(X_{s_1}) \longrightarrow \pi^L_1(X_{s_2}) \quad \text{(resp. } \pi \colon \pi^t_1(X_{s_1}) \longrightarrow \pi^t_1(X_{s_2})\text{)}.
\]

\begin{lemma}\label{lem:XIII.2.11}
设 $f \colon X \to S$ 是有限表现紧合态射，$D$ 是 $X$ 上相对 $S$ 正规交叉除子，$Y = \mathrm{Supp}\, D$，$U = X - Y$，$i \colon U \to X$ 是典范态射，$s_1 \to s_2$ 是 $S$ 的几何点的特殊化态射，$y_1$ 是 $Y_{s_1}$ 的几何点，$y_2$ 是 $Y_{s_2}$ 的几何点，使 $y_1$ 在 $X$ 上的投影 $z_1$ 是 $y_2$ 的投影 $z_2$ 的一般化. 设 $I_{y_1}$ 是 $\pi^t_1(U_{s_1})$ 在 $y_1$ 处的惯性子群. 则 $I_{y_1}$ 经特殊化态射
\[
\pi \colon \pi^t_1(U_{s_1}) \longrightarrow \pi^t_1(U_{s_2})
\]
的像是 $\pi^t_1(U_{s_2})$ 在 $y_2$ 处的惯性子群.
\end{lemma}

设 $X$ (resp. $\widetilde{X}$) 是 $X$ 在 $y_2$ (resp. $y_1$) 处的严格局部化，$U = U \times_X X$ (resp. $\widetilde{U} = U \times_X \widetilde{X}$). 有典范态射 $\widetilde{U} \to U$，由 \ref{def:XIII.1.10} 得出有交换图式
\[
\begin{tikzcd}[row sep=large, column sep=large]
\pi^t_1(\widetilde{U}_{s_1}) \arrow[r,"{\pi'}"] \arrow[d] & \pi^t_1(U_{s_2}) \arrow[d] \\
\pi^t_1(U_{s_1}) \arrow[r,"{\pi}"] & \pi^t_1(U_{s_2}),
\end{tikzcd}
\]
其中 $\pi'$ 是典范态射 $\pi^t_1(\widetilde{U}_{s_1}) \to \pi^t_1(U_{s_1})$ 与特殊化态射的复合. 由于 $\pi^t_1(\widetilde{U}_{s_1})$ (resp. $\pi^t_1(U_{s_1})$) 是 $\pi^t_1(U_{s_1})$ 在 $y_1$ 处 (resp. $\pi^t_1(U_{s_2})$ 在 $y_2$ 处) 的惯性群，只需证明 $\pi'$ 满射. 但这由 \ref{cor:XIII.5.6} 中得到的表达式得出.

\begin{corollary}\label{cor:XIII.2.12}
设 $X$ 是特征 $p \geqslant 0$ 的可分闭域 $k$ 上亏格 $g$ 的紧合光滑连通曲线. 设 $U$ 是从 $X$ 去掉 $n$ 个互异闭点 $a_1, \ldots, a_n$ 得到的开集. 则温和分歧基本群 $\pi^t_1(U)$ (2.1.3) 可由 $2g + n$ 个元素 $x_i$，$y_i$，$\sigma_j$ ($1 \leqslant i \leqslant g$，$1 \leqslant j \leqslant n$) 生成，使 $\sigma_j$ 是对应 $a_j$ 的惯性群的生成元，且满足关系
\[
(\ast) \quad \prod_{1 \leqslant i \leqslant g} (x_i y_i x_i^{-1} y_i^{-1}) \cdot \prod_{1 \leqslant j \leqslant n} \sigma_j = 1.
\]
对每个阶与 $p$ 互素、由满足关系 $(\ast)$ 的元素 $\overline{x_i}$，$\overline{y_i}$，$\overline{\sigma_j}$ 生成的有限群 $G$，存在 $U$ 的群 $G$ 的平展覆盖，对应把 $x_i$，$y_i$，$\sigma_j$ 分别送到 $\overline{x_i}$，$\overline{y_i}$，$\overline{\sigma_j}$ 的同态 $\pi^t_1(U) \to G$. 换言之，若 $p'$ 表示与 $p$ 不同的素数集，$\pi^{p'}_1(U)$ 是由生成元 $x_i$，$y_i$，$\sigma_j$ 只服从关系 $(\ast)$ 生成的 pro-$p'$-群.
\end{corollary}

\begin{proof}
可假设 $k$ 代数闭. 先假设 $k$ 特征零. 则存在 $k$ 的在 $\mathbb{Q}$ 上超越次数有限的代数闭子扩张 $k'$，使 $X$ 由定义在 $k'$ 上的紧合光滑曲线 $X'$ 经标量扩张得到，且可假设点 $a_1, \ldots, a_n$ 来自 $X'$ 的有理点 $a'_1, \ldots, a'_n$. 由于 $k'$ 在 $\mathbb{Q}$ 上超越次数有限，可以找到 $k'$ 到复数域 $\mathbb{C}$ 的嵌入；设 $\widetilde{U} = U' \times_{k'} \mathbb{C}$. 设 $k''$ 是 $k$ 的代数闭扩张，使有 $k$ 和 $\mathbb{C}$ 到 $k''$ 的 $k'$-态射. 若 $g' \colon U' \to k'$ 是结构态射，$F$ 是 $U''$ 上有限常群层，由 \ref{cor:XIII.2.9} 得出特殊化态射
\[
(R^1 g'_* F')_{\mathbb{C}} \longrightarrow (R^1 g'_* F')_{k''} \longleftarrow (R^1 g'_* F')_k
\]
是同构. 用基本群的语言，这证明有同构 (模内自同构确定)
\[
\pi_1(U) \xrightarrow{\sim} \pi_1(\widetilde{U}),
\]
显然该同构把 $X' - U'$ 的某点处的惯性群变成同一点处的惯性群. 可假设 $k = \mathbb{C}$. 此时由 \emph{Riemann 存在定理} (\cref{cor:XII.5.2}) 得出基本群 $\pi_1(U)$ 正是 $\pi_1(U)$ 的关联解析空间的拓扑基本群关于有限指数子群拓扑的完备化. 后者可以用超越方式计算 \cite[第 7 章 §47]{exp13-3}；它可由 $2g + n$ 个元素 $x_i$，$y_i$，$\sigma_j$ 生成，其中 $\sigma_j$ 是中心在 $a_j$ 的小圆盘的局部基本群 $\pi_1(D_j)$ 的生成元的像，i.e. 对应点 $a_j$ 的惯性群的生成元，这些元素只满足关系 $(\ast)$.

若 $k$ 的特征 $p > 0$，可以找到剩余域 $k$、分式域 $K$ 特征零的完备离散赋值环 $A$，和 $S = \mathrm{Spec}\, A$ 上连通、紧合且光滑的概形 $X_1$，使 $X_1 \times_S \mathrm{Spec}\, k \simeq X$ (\cref{cor:III.7.4}). 诸点 $a_j$ 提升为 $X_1$ 在 $S$ 上方的截面 $s_j$；设 $Y_{1j}$ 是底集为 $s_j(S)$ 的既约闭子概形，$Y_1$ 是诸 $Y_{1j}$ 的并，$U_1 = X_1 - Y_1$，$g_1 \colon U_1 \to S$ 是结构态射. 设 $K$ 是 $K$ 的代数闭扩张，$U = U_1 \times_S K$. 若 $C$ 是有限常群，由 \ref{cor:XIII.2.8} 得出特殊化态射
\[
(R^1 g_{1*} C_{U_1})_k \longrightarrow (R^1 g_{1*} C_{U_1})_K
\]
单射，且 $C$ 的阶与 $p$ 互素时甚至双射. 用基本群的语言，这意味着特殊化态射 (\cref{def:XIII.1.10})
\[
\pi \colon \pi_1(U) \longrightarrow \pi^t_1(U)
\]
满射，且特殊化态射
\[
\pi^{p'}_1(U) \longrightarrow \pi^{p'}_1(U)
\]
双射. 最后，若 $x_i$，$y_i$，$\sigma_j$ 是 $\pi_1(U)$ 的生成元，其中 $\sigma_j$ 是对应点 $b_j = Y_{1j}(K)$ 的惯性群的生成元，则由 \ref{exa:XIII.1.11}，$\pi(\sigma_j)$ 是对应 $a_j$ 的惯性群的生成元，这完成证明.
\end{proof}

\begin{remark}[2003 年 (MR) 增补]
设 $k$ 是特征 $p > 0$ 的代数闭域，$X$ 是 $k$ 上亏格 $g$ 的紧合光滑连通代数曲线，$U$ 是 $X$ 的仿射开集，即 $X$ 的 $r \geqslant 1$ 个有理点的余集. 有基本群 $\pi_1(U)$，它的商 $\pi^t_1(U)$ 分类 $U$ 的在 $X - U$ 的点处温和分歧的有限平展覆盖，以及 $\pi^t_1(U)$ 的商 $\pi^{p'}_1(U)$ 分类 $U$ 的 伽罗瓦群阶与 $p$ 互素的 Galois 平展覆盖.

在 (Coverings of algebraic curves，Amer. J. Math. 79 (1957)，p. 825--856) 中，S. Abhyankar 对 $U$ 的有限平展覆盖的 伽罗瓦群的有限群结构提出若干猜想.

关于阶与 $p$ 互素、且是 $U$ 的连通平展覆盖的 伽罗瓦群的有限群的猜想，已在\cref{cor:XIII.2.12} 中证明并精确化. 在探讨 $\pi_1(U)$ 的有限商之前，先给出这个群的"大小"的若干指示. 设 $X = \mathrm{Spec}(A)$. 已知 $\mathrm{Hom}_{\mathrm{cont}}(\pi_1(U), \mathbb{Z}/p\mathbb{Z})$ 由 Artin-Schreier 理论描述 (\cref{cor:XI.6.9}). 该群同构于 $A/\wp(A)$，其中 $\wp$ 是把 $a$ 送到 $a^p - a$ 的 $A$ 到 $A$ 的映射. 先假设 $U$ 是仿射直线 $\mathbf{A}^1$，环为 $k[T]$. 设 $E$ 是 $k[T]$ 的形如 $\sum_i a_i T^i$ (其中 $p$ 整除 $i$ 时 $a_i = 0$) 的元素集. 立即看到复合映射 $E \to k[T] \to k[T]/\wp k[T]$ 双射. 于是系数 $a_i$ ($(i, p) = 1$) 表现得像参数化仿射直线的 $p$ 阶循环覆盖的空间的坐标. 特别由此得出 $\pi_1(\mathbf{A}^1)$ 不拓扑有限型，且若作代数闭域的扩张 $k \to k'$，自然态射 $\pi_1(\mathbf{A}^1 \times_k k') \to \pi_1(\mathbf{A}^1)$ (满射) 不双射，与紧合情形相反. 在一般情形，曲线 $U$ 实现为仿射直线上的有限概形，得出同样的现象对 $\pi_1(U)$ 发生 (而且更一般地对维数 $\geqslant 1$ 的任何连通有限型仿射 $k$-概形 $V$ 的 $\pi_1(V)$ 发生).

于是，若 $G$ 是有限群，记 $G^{(p')}$ 为 $G$ 的阶与 $p$ 互素的最大商群. 为使有限群 $G$ 是 $\pi_1(U)$ 的拓扑商，必须 $G^{(p')}$ 是 $\pi^{p'}_1(U)$ 的拓扑商，而原则上可以借助\cref{cor:XIII.2.12} 回答这个条件. 在前述文章中，Abhyankar 猜想这个必要条件也充分. 该猜想已被 M. Raynaud 在仿射直线的情形证明，被 D. Harbater 在一般情形证明 (Invent. Math. 116 (1994)，p. 425--462 和 117，p. 1--25). 例如，在仿射直线的情形 $\pi^{p'}_1(\mathbf{A}^1) = 1$，得出有限群 $G$ 是 $\mathbf{A}^1$ 的连通平展覆盖的 伽罗瓦群，当且仅当 $G^{(p')} = 1$，i.e. 当且仅当 $G$ 由它的 Sylow $p$-子群生成. 从而任何阶为 $p$ 的倍数的有限单群都合适.
\end{remark}

\section{上同调真性与一般局部无环性}

\begin{theorem}\label{thm:XIII.3.1}
设 $S$ 是生成点为 $s$ 的不可约概形，$X$ 和 $Y$ 是两个有限表现 $S$-概形，$f \colon X \to Y$ 是 $S$-态射. 对每个 $S$-概形 $S'$，用 $Y'$，$X'$，等等记 $Y$，$X$，等等经态射 $S' \to S$ 的逆像. 有下列性质：
\begin{enumerate}[1)]
\item a) 可以找到 $S$ 的非空开集 $S'$，使对 $X'$ 上每个有限常集合层 $F'$，$f'_* F'$ 可构，且 $(F', f')$ 相对 $S'$ 在维数 $\leqslant 0$ 时上同调真.
\item b) 设 $F$ 是 $X$ 上的可构集合层. 则可以找到 $S$ 的非空开集 $S'$ (依赖 $F$)，使 $f'_* F'$ 可构，且 $(F', f')$ 相对 $S'$ 在维数 $\leqslant 0$ 时上同调真.
\item 假设 $k(s)$ 的代数闭包 $k$ 上维数 $\leqslant \dim X_s$ 的有限型概形强可消解奇点 (SGA 5 I 3.1.5). 则还有下列性质：
\item a) 可以找到 $S$ 的非空开集 $S'$，使对 $X'$ 上每个与 $S$ 的剩余特征互素的阶有限常群层 $F'$，若 $\Phi'$ 是 $F'$ 下的挠子叠，$f'_* \Phi'$ 1-可构，且 $(F', f')$ 相对 $S'$ 在维数 $\leqslant 1$ 时上同调真.
\item b) 设 $L$ 是与 $S$ 的剩余特征不同的素数集，$\Phi$ 是 $X$ 上 (0) 的 1-可构 ind-$L$-叠，使对每个平展于 $X$ 上的概形 $X_1$ 和 $\Phi_{X_1}$ 的每对对象 $x$，$x_1$，层 $\mathrm{Hom}_{X_1}(x, x_1)$ 可构. 再假设 $S$ 局部诺特. 则可以找到 $S$ 的非空开集 $S'$，使 $f'_* \Phi'$ 1-可构，对 $(f'_* \Phi')_{Y_1}$ 的每对对象 $y$，$y_1$，$\mathrm{Hom}_{Y_1}(y, y_1)$ 可构，且 $(\Phi', f')$ 相对 $S'$ 在维数 $\leqslant 1$ 时上同调真.
\end{enumerate}
\end{theorem}

\begin{proof}
可假设 $S$ 仿射；由 SGA 4 VIII 1.1，可假设 $S$ 整；最后经极限过渡，可假设 $S$ 是 $\mathbb{Z}$ 上有限型代数的谱；特别 $S$ 此时诺特. 由于问题局部于 $Y$，可假设 $Y$ 仿射. 而且只需在有限扩张 $S' \to S$ 后证明定理，其中 $S'$ 是整概形，$S' \to S$ 由平展态射和有限根式满射态射复合而成.

**1) 常集合层的情形**

1) 1. 归结到 $X$ 在 $S$ 上正规的情形. 设 $X_{1s}$ 是 $(X_s)_{\mathrm{r\acute{e}d}}$ 的正规化；把 $S$ 限制到非空开集并作 $S$ 的根式扩张后，可假设 $X_{1s}$ 来自 $S$ 上正规的概形 $X_1$，且态射 $X_{1s} \to X_s$ 来自有限满射态射 $p \colon X_1 \to X$ (EGA IV 8.8.2 和 9.6.1). 假设定理对 $fp$ 已证. 把 $S$ 限制到开集后，可假设对 $X$ 上每个常集合层 $F$，$(p_* F, fp)$ 相对 $S$ 在维数 $\leqslant 0$ 时上同调真，且 $f_* p_*(p^{*} F)$ 可构. 由 \ref{cor:XIII.1.9}，$(p^{*} p_* F, f)$ 相对 $S$ 在维数 $\leqslant 0$ 时上同调真. 态射
\[
F \longrightarrow p_* p^{*} F = G
\]
是单态射. 已经由此得出，$f_* F$ 是 $f_* G$ 的子层，故 $f_* F$ 可构 (SGA 4 IX 2.9 (ii))，且 $(F, f)$ 相对 $S$ 在维数 $\leqslant -1$ 时上同调真.

设 $X_2 = X_1 \times_X X_1$，$q \colon X_2 \to X$ 是典范态射. 由 \ref{exa:XIII.1.11} 1)，有正合序列
\[
F \longrightarrow G \rightrightarrows q_* q^{*} F.
\]
由刚才对 $fq$ 代替 $f$ 证明的，可假设 (把 $S$ 限制到非空开集后) 对 $X$ 上每个常集合层 $F$，$(q^{*} F, fq)$ 相对 $S$ 在维数 $\leqslant -1$ 时上同调真，故 $(q_* q^{*} F, f)$ 相对 $S$ 在维数 $\leqslant -1$ 时上同调真. 由 \ref{prop:XIII.1.13} 1) 得出 $(F, f)$ 相对 $S$ 在维数 $\leqslant 0$ 时上同调真.

1) 2. 归结到 $X$ 在 $S$ 上正规仿射的情形.

设 $U_s$ 是 $X_s$ 的稠密仿射开集. 把 $S$ 限制到非空开集后，可假设 $U_s \to X_s$ 提升为相对 $S$ 概形支配的开浸入 $i \colon U \to X$ (EGA IV 8.9.1). 由于态射 $X \to S$ 正规，由 SGA 2 XIV 1.18 有：
\[
\mathrm{prof}^{\mathrm{ét}}_{S-U}(X) \geqslant 2;
\]
从而对 $X$ 上每个常层 $F$，典范态射
\[
F \longrightarrow i_* i^{*} F
\]
是同构. 由此得出，若假设定理对 $fi$ 和 $i$ 已证，则把 $S$ 限制到非空开集后，$(i_* F, i)$ 和 $(i_* F, fi)$ 相对 $S$ 在维数 $\leqslant 0$ 时上同调真. 故 $(F, f)$ 也如此 (1.6 2)). 又由于 $f_* F = (fi)_*(i_* F)$ 可构，这完成归结.

1) 3. 证明结束.

可假设 $S$ 正规 (EGA IV 7.8.3). 可以找到 $X_s$ 的紧化：
\[
\begin{tikzcd}[row sep=large, column sep=large]
X_s \arrow[r,"{j_s}"] \arrow[dr,"{f_s}"] & P_s \arrow[d,"{g_s}"] \\
& Y_s
\end{tikzcd}
\]
其中 $j_s$ 是支配开浸入，$g_s$ 是紧合态射；作 $k(s)$ 的根式扩张并把 $P_s$ 换成它的正规化 (不改变 $X_s$) 后，可假设 $P_s$ 几何正规. 把 $S$ 限制到非空开集并作根式满射扩张后，可假设上图式来自图式
\[
\begin{tikzcd}[row sep=large, column sep=large]
X \arrow[r,"{j}"] \arrow[dr,"{f}"] & P \arrow[d,"{g}"] \\
& Y
\end{tikzcd}
\]
其中 $P$ 是 $S$ 上正规概形，$j$ 是相对 $S$ 概形支配的开浸入，$g$ 是紧合态射 (EGA IV 6.9.1，9.9.4 和 9.6.1). 对 $X$ 上取值 $C$ 的每个有限常集合层 $F$，$j_* F$ 是取值 $C$ 的常层 (SGA 4 2.14.1)，经任何基变换 $S' \to S$ 后也如此；由此得出 $(F, j)$ 相对 $S$ 在维数 $\leqslant 0$ 时上同调真. 故 $(F, f)$ 也如此，因为 $g$ 紧合 (\cref{cor:XIII.1.8}). 由于 $g$ 紧合，$f_* F = g_* C_P$ 可构，这完成 1) a) 的证明.

**2) 可构集合层的情形**

设 $F$ 是 $X$ 上的可构集合层. 由 SGA 4 IX 2.14 (ii)，可以找到有限族态射 $p_i \colon Z_i \to X$，以及在每个 $Z_i$ 上取值 $C_i$ 的有限常集合层 $C_i$，使有单态射
\[
j \colon F \longrightarrow \prod_i p_{i*} C_i = G.
\]
由 1) a)，可假设 (把 $S$ 限制到非空开集后) 诸 $(C_i, fp_i)$ 相对 $S$ 在维数 $\leqslant 0$ 时上同调真，诸 $f_* p_{i*} C_i$ 可构. 已经得出 $(G, f)$ 相对 $S$ 在维数 $\leqslant 0$ 时上同调真 (\cref{cor:XIII.1.9})，故 $(F, f)$ 相对 $S$ 在维数 $\leqslant -1$ 时上同调真，且 $f_* F$ 可构. 设 $K$ 是余直和 $K = G \amalg_F G$；由于 $F$ 和 $G$ 可构，$K$ 也可构. 由前面得出，把 $S$ 限制到非空开集后可假设 $(K, f)$ 相对 $S$ 在维数 $\leqslant -1$ 时上同调真. 由 \ref{prop:XIII.1.13} 1) 得出 $(F, f)$ 相对 $S$ 在维数 $\leqslant 0$ 时上同调真.

**3) 常群层的情形**

若 $F$ 是 $X$ 上的常群层，记 $\Phi$ 为 $F$ 下的挠子叠.

3) 1. 先证明，把 $S$ 限制到非空开集后，对 $X$ 上每个与 $S$ 的剩余特征互素的阶有限常群层 $F$，$(F, f)$ 相对 $S$ 在维数 $\leqslant 0$ 时上同调真，且 $f_* \Phi$ 可构.

为此归结到 $X$ 在 $S$ 上光滑的情形. 作 $k(s)$ 的有限扩张 (可把它看作平展扩张与根式扩张的复合) 后，可以找到紧合满射态射 $p_s \colon X_{1s} \to X_s$，其中 $X_{1s}$ 是与 $X_s$ 同维数、在 $S$ 上光滑的概形，且把 $S$ 限制到非空开集后可假设 $p_s$ 来自紧合满射态射 $p \colon X_1 \to X$，其中 $X_1$ 在 $S$ 上光滑 (EGA IV 9.6.1 和 12.1.6). 设 $X_2 = X_1 \times_X X_1$，$q \colon X_2 \to X$ 是典范态射. 有 $X$ 上叠的正合图式
\[
\Phi \longrightarrow p_* p^{*} \Phi \rightrightarrows q_* q^{*} \Phi
\]
(1.11 2)). 假设定理在光滑情形已证，首先看到可假设 $f_* p_* p^{*} \Phi$ 可构；故 $f_* \Phi$ 也可构 (下面 \ref{lem:XIII.3.1.1}). 而且由 \ref{prop:XIII.1.6} 2)，可假设 $(p_* p^{*} \Phi, f)$ 相对 $S$ 在维数 $\leqslant 0$ 时上同调真；由此得出 $(\Phi, f)$ 相对 $S$ 在维数 $\leqslant -1$ 时上同调真. 可假设 $(q_* q^{*} \Phi, f)$ 相对 $S$ 在维数 $\leqslant -1$ 时上同调真，这蕴含 $(\Phi, f)$ 相对 $S$ 在维数 $\leqslant 0$ 时上同调真 (1.12 1)).

随后如 1) 2 归结到 $X$ 在 $S$ 上光滑仿射的情形. 设
\[
\begin{tikzcd}[row sep=large, column sep=large]
X \arrow[r,"{i}"] \arrow[dr,"{f}"] & P \arrow[d,"{q}"] \\
& Y
\end{tikzcd}
\]
是 $X$ 的紧化，其中 $i$ 是支配开浸入，$q$ 是紧合态射. 由于 $\dim P_s = \dim X_s$，可以对 $P_s$ 应用奇点消解假设. 作 $S$ 的平展扩张和根式扩张后，可以找到紧合态射 $r \colon Z \to P$，其中 $Z$ 在 $S$ 上光滑，$r^{-1}(X) \simeq X$，且 $r^{-1}(X)$ 是 $Z$ 中相对 $S$ 正规交叉除子的余集. 由 \ref{rem:XIII.2.3} b)，$F$ 下的每个挠子相对 $S$ 在 $Z$ 上温和分歧. 由 \ref{cor:XIII.2.7} 得出 $(F, f)$ 相对 $S$ 在维数 $\leqslant 0$ 时上同调真，这证明我们的断言.

3) 2. 归结到 $X$ 在 $S$ 上光滑的情形.

作 $k(s)$ 的有限扩张后，可以找到紧合满射态射 $p_s \colon X_{1s} \to X$，其中 $X_{1s}$ 在 $s$ 上光滑，且可假设 $p_s$ 来自紧合满射态射 $p \colon X_1 \to X$，其中 $X_1$ 在 $S$ 上光滑. 假设定理对 $fp$ 已证，对 $f$ 证明. 设 $F$ 是 $X$ 上与 $S$ 的剩余特征互素的阶有限常群层，$\Phi$ 是 $F$ 下的挠子叠. 设 $X_2 = X_1 \times_X X_1$，$X_3 = X_1 \times_X X_1 \times_X X_1$，$q \colon X_2 \to X$，$r \colon X_3 \to X$ 是典范态射. 由 \ref{exa:XIII.1.11} 2)，有叠的正合图式
\[
\Phi \longrightarrow p_* p^{*} \Phi \rightrightarrows q_* q^{*} \Phi \Rrightarrow r_* r^{*} \Phi.
\]
为证明 $(\Phi, f)$ 相对 $S$ 在维数 $\leqslant 1$ 时上同调真，只需证明 $(p_* p^{*} \Phi, f)$ 如此，$(q_* q^{*} \Phi, f)$ 相对 $S$ 在维数 $\leqslant 0$ 时上同调真，$(r_* r^{*} \Phi, f)$ 相对 $S$ 在维数 $\leqslant -1$ 时上同调真 (1.12 2)). 由上面 3) 1，可假设对每个有限常群层 $F$，$(q^{*} \Phi, fq)$，$(q^{*} \Phi, q)$，$(r^{*} \Phi, fr)$，$(r^{*} \Phi, r)$ 相对 $S$ 在维数 $\leqslant 0$ 时上同调真. 由 \ref{prop:XIII.1.6} 2) 得出 $(q_* q^{*} \Phi, f)$ 和 $(r_* r^{*} \Phi, f)$ 相对 $S$ 在维数 $\leqslant 0$ 时上同调真. 假设定理在光滑情形已证，$(p_* \Phi, fp)$ 和 $(p_* \Phi, p)$ 从而 $(p_* p^{*} \Phi, f)$ (1.6 2)) 相对 $S$ 在维数 $\leqslant 1$ 时上同调真. 这确实证明 $(\Phi, f)$ 相对 $S$ 在维数 $\leqslant 1$ 时上同调真.

而且由假设 $f_* p_* p^{*} \Phi$ 1-可构；由 \ref{thm:XIII.3.1} 可假设 $f_* q_* q^{*} \Phi$ 可构；由下面 \ref{lem:XIII.3.1.1} 得出 $f_* \Phi$ 1-可构.

3) 3. 归结到 $X$ 在 $S$ 上光滑仿射的情形.

由 3) 2，可假设 $X$ 在 $S$ 上光滑. 设 $(U_i)_{i \in I}$ 是 $X$ 的仿射开集的有限覆盖，$X_1$ 是诸 $U_i$ 的直和，$p \colon X_1 \to X$ 是典范态射. 由于 $p$ 是变动的概形上有限型平展层范畴的有效下降态射，如 3) 2 看到，若假设定理对诸 $U_i$，i.e. 对 $X_1$ 已证，则对 $X$ 也真.

3) 4. $X$ 在 $S$ 上光滑仿射的情形.

如 3) 1 看到，把 $S$ 限制到非空开集并作 $S$ 的平展满射扩张和根式满射扩张后，可以找到交换图式
\[
\begin{tikzcd}[row sep=large, column sep=large]
X \arrow[r,"{i}"] \arrow[dr,"{f}"] & P \arrow[d,"{q}"] \\
& Y
\end{tikzcd}
\]
其中 $P$ 在 $S$ 上光滑，$X$ 是 $P$ 中相对 $S$ 正规交叉除子的余集，$g$ 是紧合态射. 若 $F$ 是与 $S$ 的剩余特征互素的阶常层，$F$ 下的每个挠子相对 $S$ 在 $P$ 上温和分歧 (2.3 b)). $(F, f)$ 相对 $S$ 在维数 $\leqslant 1$ 时上同调真且 $f_* \Phi$ 可构的事实于是由 \ref{cor:XIII.2.7} 得出.

**4) 2) b) 的证明**

4) 1. $\Phi$ 是叠丛的情形.

可以找到有限型平展满射态射 $p \colon X_1 \to X$，使 $p^{*} \Phi$ 是平凡叠丛. 经下降，如 3) 2，看到只需对 $X_1$，$X_1 \times_X X_1$ 和 $X_1 \times_X X_1 \times_X X_1$ 证明定理. 归结到 $\Phi$ 是可构群层 $F$ (其纤维阶与 $S$ 的剩余特征互素) 下的挠子叠丛的情形.

由 SGA 4 IX 2.14，可以找到有限态射 $p_i \colon Z_i \to X$ 的有限族，以及对每个 $i$，与 $S$ 的剩余特征互素的阶有限常群层 $C_i$，使有单态射
\[
j \colon F \longrightarrow \prod_i p_{i*} C_i = G.
\]
设 $\Phi_i$ 是 $C_i$ 下的挠子叠，$\Psi$ 是 $G$ 下的挠子叠. 由 2) a)，把 $S$ 限制到非空开集后可假设诸 $(C_i, fp_i)$ 相对 $S$ 在维数 $\leqslant 1$ 时上同调真，诸叠 $f_* p_{i*} \Phi_i$ 1-可构. 由 \ref{cor:XIII.1.9} 得出诸 $(p_{i*} C_i, f)$ 相对 $S$ 在维数 $\leqslant 1$ 时上同调真；故 $(G, f)$ 也如此. 而且，由于 $p_{i*} \Phi_i$ 等价于群 $p_{i*} C_i$ 下的挠子叠 (SGA 4 VIII 5.8)，看到 $f_* \Psi$ 1-可构.

由于 $R^1 f_* G$ 可构，可以找到由有限型平展 $Y$-概形 $T$ 表示的层和满态射
\[
a \colon T \longrightarrow R^1 f_* G
\]
(SGA 4 IX 2.7)；且可假设 $T(T)$ 的恒等截面的像由 $X \times_Y T = X_T$ 上群 $G|_{X_T}$ 的挠子 $Q$ 定义. 设 $f_T \colon X_T \to T$ 是典范态射，$F_T = F|_{X_T}$，等等. 由 1) b) 可假设 (把 $S$ 限制到非空开集后) $(Q/F_T, f_T)$ 相对 $S$ 在维数 $\leqslant 0$ 时上同调真，且 $f_{T*}(Q/F_T)$ 可构. 由 \ref{lem:XIII.3.1.2} 得出 $f_* \Psi$ 可构.

证明 $(F, f)$ 相对 $S$ 在维数 $\leqslant 1$ 时上同调真. 由 \ref{prop:XIII.1.13} 2) 只需证明：对每个平展于 $Y$ 上的概形 $Y_1$ 和 $X_1 = X \times_Y Y_1$ 上的每个挠子 $Q_1$，若 $f_1 \colon X_1 \to Y_1$ 是典范态射，则 $(Q_1/F_1, f_1)$ 相对 $S$ 在维数 $\leqslant 0$ 时上同调真. 由 $T$ 的定义，$Q_1$ 局部于 $Y_1$ 的平展拓扑是 $Q$ 的逆像，这证明归结.

4) 2. 一般情形.

用\cref{lem:XIII.6.1.1}、4) 1 和 1) a) 看到，把 $S$ 限制到非空开集后可假设 $S(f_* \Phi)$ 可构，$(S_\Phi, f)$ 相对 $S$ 在维数 $\leqslant 0$ 时上同调真. 则可以找到由有限型平展 $Y$-概形 $T$ 表示的层和满态射
\[
a \colon T \longrightarrow S(f_* \Phi).
\]
沿用 4) 1 的记号，再令 $Z = T \times_Y T$，$X_Z = X \times_Y Z$，用 $f_Z \colon X_Z \to Z$ 记典范态射. 可假设 $T$ 选得使 $T(T)$ 的恒等截面经 $a$ 的像 $q$ 由 $(f_* \Phi)_T = \Phi_{X_T}$ 的对象 $p$ 定义. 设 $p_1$ 和 $p_2$ (resp. $q_1$ 和 $q_2$) 是 $p$ (resp. $q$) 经 $Z$ 到 $T$ 的两个投影的逆像. 可假设 (把 $S$ 限制到非空开集后) $(\mathrm{Aut}_{X_T}(p), f_T)$ 相对 $S$ 在维数 $\leqslant 1$ 时上同调真，$(\mathrm{Hom}_{X_Z}(p_1, p_2), f_Z)$ 相对 $S$ 在维数 $\leqslant 0$ 时上同调真，且 $f_{T*}(\mathrm{Aut}_{X_T}(p))$，$f_{Z*}(\mathrm{Hom}_{X_Z}(p_1, p_2))$ 可构 (a) 1) 和 4) 1)).

由此首先得出，对每个平展于 $Y$ 上的概形 $Y_1$ 和 $(f_* \Phi)_{Y_1}$ 的每对对象 $y$，$y_1$，层 $\mathrm{Aut}_{Y_1}(y)$ (resp. $\mathrm{Hom}_{Y_1}(y, y_1)$) 可构，这样的层局部于 $Y_1$ 的平展拓扑是 $f_{T*}(\mathrm{Aut}_{X_T}(p))$ (resp. $f_{Z*}(\mathrm{Hom}_{X_Z}(p_1, p_2))$) 的逆像.

剩下证明 $(\Phi, f)$ 相对 $S$ 在维数 $\leqslant 1$ 时上同调真. 只需证明：对每个 $S'$-概形 $S$ 和 $Y$ 的每个几何点 $y'$，若记 $Y$ 为 $Y$ 在 $y'$ 处的严格局部化，$X = X \times_Y Y$，等等，则典范函子
\[
\phi \colon \Phi(X) \longrightarrow \Phi'(X')
\]
是范畴等价.

证明 $\phi$ 忠实全. 设 $x$，$y \in \Phi(X)$，$x'$，$y'$ 是它们在 $\Phi'(X')$ 中的像，证明典范态射
\[
(\ast) \quad \mathrm{Hom}_X(x, y) \longrightarrow \mathrm{Hom}_{X'}(x', y')
\]
双射. 由 $T$ 的定义，存在两个 $Y$ 到 $T$ 的态射，使 $x$ 和 $y$ 是 $p$ 经这两个态射的逆像. 这等价于有态射 $Y \to Z$，从而有态射 $h \colon X \to X_Z$，使
\[
h^{*}(p_1) = x \quad h^{*}(p_2) = y.
\]
从而有典范同构
\[
\mathrm{Hom}_X(x, y) = h^{*}(\mathrm{Hom}_{X_Z}(p_1, p_2)).
\]
但考虑到 $(\mathrm{Hom}_{X_Z}(p_1, p_2), f_Z)$ 相对 $S$ 在维数 $\leqslant 0$ 时上同调真，看到 $(\mathrm{Hom}_X(x, y), f)$ 也如此，这证明态射 $(\ast)$ 双射.

证明 $\phi$ 本质满射. 设 $x' \in \Phi'(X')$. 由于 $(S_\Phi, f)$ 相对 $S$ 在维数 $\leqslant 0$ 时上同调真，典范态射
\[
H^0(X, S_\Phi) \longrightarrow H^0(X', S_{\Phi'})
\]
双射. 设 $G'$ 是 $x'$ 生成的 $\Phi'$ 的极大子叠丛；则存在 $\Phi$ 的极大子叠丛 $G$，使 $G$ 在 $X'$ 上的逆像是 $G'$. 由 4) 1，$(G, f)$ 相对 $S$ 在维数 $\leqslant 1$ 时上同调真；从而典范函子
\[
G(X) \longrightarrow G'(X')
\]
是范畴等价，这证明存在 $\Phi(X)$ 的元素 $x$，使 $x$ 在 $\Phi'(X')$ 中的像同构于 $x'$，完成定理的证明.
\end{proof}

\begin{lemma}\label{lem:XIII.3.1.1}
设 $S$ 是局部诺特概形，
\[
\Phi \xrightarrow{p} \Phi_1 \xrightarrow[p_2]{p_1} \Phi_2
\]
是 $S$ 上叠的正合图式 (1.10.1). 若 $\Phi_1$ 可构，$\Phi$ 也可构. 若对每个平展于 $S$ 上的概形 $S'$ 和 $(\Phi_1)_{S'}$ 的每对对象 $x_1$，$y_1$，层 $\mathrm{Hom}_{S'}(x_1, y_1)$ 可构，则对 $\Phi_{S'}$ 的每对对象 $x$，$y$，$\mathrm{Hom}_{S'}(x, y)$ 也可构. 假设 $\Phi_1$ 1-可构 (0) 且 $\Phi_2$ 可构，则 $\Phi$ 1-可构.
\end{lemma}

对每个平展于 $S$ 上的概形 $S'$ 和 $\Phi_{S'}$ 的每个对象 $x$，有单态射
\[
\mathrm{Aut}_{S'}(x) \longrightarrow \mathrm{Aut}_{S'}(p(x)).
\]
由 SGA 4 IX 2.9 得出，$\Phi_1$ 可构时 $\Phi$ 也可构，引理的第二个断言同样证明.

现在假设 $\Phi_1$ 1-可构，$\Phi_2$ 可构. 态射 $p$ 在极大子叠丛层上诱导态射
\[
\phi \colon S_\Phi \longrightarrow S_{\Phi_1}.
\]
设 $G$ 是 $S_\Phi$ 经 $\phi$ 的像；由 SGA 4 IX 2.9，$G$ 是可构层. 可以找到由有限型平展 $S$-概形 $T$ 表示的层和满态射
\[
a \colon T \longrightarrow G
\]
(SGA 4 IX 2.7). 且可以选 $T$，使 $T(T)$ 的恒等截面的像 $y$ 由 $(\Phi_1)_T$ 的形如 $x_1 = p(x)$ (其中 $x \in \Phi_T$) 的对象 $x_1$ 定义.

只需证明对 $S$ 的每个点 $s$，存在 $\overline{\{s\}}$ 的非空开集 $U$，使 $S_\Phi|_U$ 局部常可构. 设 $s \in S$，$s$ 是 $s$ 上方的几何点，$y_1, \ldots, y_n$ 是 $G_s$ 的元素. 由 $T$ 的定义存在态射 $h_i \colon s \to T$，使 $h^{*}_i(y) = y_i$. 设 $S'$ 是 $n$ 个同构于 $T$ 的概形在 $S$ 上的纤维积，$h \colon s \to S'$ 是诸 $h_i$ 的纤维积，$y_i$ (resp. $x_i$) 是 $y$ (resp. $x$) 经 $S'$ 到 $T$ 的第 $i$ 个投影的逆像. 设 $F_i$ 是 $S_\Phi|_{S'}$ 的 $y_i$ 的逆像子层，证明诸 $F_i$ 可构. 层 $F_i$ 是层 $F'_i$ 的商，后者对每个平展于 $S'$ 上方的概形 $S''$ 满足
\[
F'_i(S'') = \{\Phi_{S''} \text{ 的配备同构 } i \colon p(z) \simeq p(x_i|_{S''}) \text{ 的对象 } z \text{ 的模同构的类}\}.
\]
只需证明诸 $F'_i$ 可构. 若令 $z_i = p_1 p(x_i)$，$z'_i = p_2 p(x_i)$，有单态射
\[
\Psi_i \colon F'_i \longrightarrow \mathrm{Isom}_{S'}(z_i, z'_i),
\]
它把每个平展于 $S'$ 上的概形 $S''$ 和使有同构 $i \colon p(z) \simeq p(x_i|_{S''})$ 的 $\Phi_{S''}$ 的对象 $z$ 配给由下列图式交换条件定义的 $z_i$ 到 $z'_i$ 的同构
\[
\begin{tikzcd}[row sep=large, column sep=large]
p_1 p(z) \arrow[r,"{p_1(i)}"] \arrow[d,"{j}"] & z_i \arrow[d] \\
p_2 p(z) \arrow[r,"{p_2(i)}"'] & z'_i
\end{tikzcd}
\]
($j$ 是与正合图式关联的典范态射).

态射 $\Psi_i$ 单射，因为说两个使有同构 $i \colon p(z) \xrightarrow{\sim} p(x_i|_{S''})$，$i' \colon p(z') \xrightarrow{\sim} p(x_i|_{S''})$ 的对象 $z$，$z'$ 定义 $\mathrm{Isom}_{S''}(z_i, z'_i)$ 的同一个元素，等价于说 $p_1(i'^{-1} i) = p_2(i'^{-1} i)$，i.e. $i'^{-1} i$ 来自同构 $z \to z'$. 层 $F'_i$ 是 $\mathrm{Isom}_{S'}(z_i, z'_i)$ 的子层，故可构.

由于可以找到 $\overline{\{s\}}$ 的非空开集 $U$，使 $y_1|_U, \ldots, y_n|_U$ 生成 $G$，由下面\cref{lem:XIII.6.1.2} 得出 $S_\Phi|_U$ 可构，故把 $U$ 缩小后 $S_\Phi|_U$ 局部常可构.

\begin{lemma}\label{lem:XIII.3.1.2}
设 $S$ 是局部诺特概形，$f \colon X \to S$ 是态射，$F \to G$ 是 $X$ 上群层的单态射，$\Psi$ (resp. $\Psi_1$) 是 $F$ (resp. $G$) 下的挠子叠，$\Phi = f_* \Psi$，$\Phi_1 = f_* \Psi_1$. 假设给定 $S$ 上由有限型平展 $S$-概形 $T$ 表示的层和满态射
\[
a \colon T \longrightarrow S_{\Phi_1} \simeq R^1 f_* G,
\]
使存在 $X_T = X \times_S T$ 上群 $G|_{X_T}$ 的挠子 $Q$，在 $R^1 f_* G(T)$ 中定义 $T(T)$ 的恒等截面经 $a$ 的像. 设 $f_T \colon X_T \to T$ 是典范态射，令 $F_T = F|_{X_T}$. 假设 $\Phi_1$ 1-可构且 $f_{T*}(Q/F_T)$ 可构. 则 $\Phi$ 1-可构.
\end{lemma}

只需照抄 \ref{lem:XIII.3.1.1} 的证明，把存在态射 $\Psi_i$ 换成存在同构
\[
F'_i \xrightarrow{\sim} f_*(Q/F_T)|_{S'}.
\]

\begin{remark}\label{rem:XIII.3.1.3}
假设 $k = k(s)$ 特征零；则 $k$ 上维数 $\leqslant \dim X$ 的有限型概形强可消解奇点，\ref{thm:XIII.3.1} 的证明可以证明下列结果：
a) 存在 $S$ 的非空开集 $S_1$，使对 $S_1$ 上方的每个极大点特征零的概形 $S'$ 和 $X' = X \times_S S'$ 上每个局部常可构集合层 $F$，$(F, f')$ 相对 $S'$ 在维数 $\leqslant 0$ 时上同调真.
b) 若 $S$ 的所有剩余特征为零，存在 $S$ 的非空开集 $S_1$，使对 $S_1$ 上方的每个概形 $S'$ 和 $X'$ 上每个局部常可构群层 $F$，$(F, f')$ 相对 $S'$ 在维数 $\leqslant 1$ 时上同调真.

只需照抄 \ref{thm:XIII.3.1} 2) a) 的证明. 3) 4 中用的命题 \ref{cor:XIII.2.7} 适用于局部常层 $F$ 的情形，因为沿用 3) 4 的记号，由于 $S$ 的所有剩余特征为零，$F$ 下的每个挠子相对 $S$ 在 $P$ 上温和分歧.
\end{remark}

\begin{corollary}\label{cor:XIII.3.2}
设 $k$ 是特征 $p \geqslant 0$ 的域，$p'$ 是与 $p$ 不同的素数集，$f \colon X \to k$ 是相容态射.
\begin{enumerate}[1)]
\item 对每个集合层 $F$，$(F, f)$ 在维数 $\leqslant 0$ 时上同调真.
\item 假设满足下列条件之一：
\begin{enumerate}[a)]
\item $f$ 有限型，且 $k$ 的代数闭包上维数 $\leqslant \dim X$ 的有限型概形强可消解奇点.
\item $k$ 的代数闭包上有限型概形强可消解奇点.
\end{enumerate}
则对每个 ind-$p'$-群层 $F$，$(F, f)$ 在维数 $\leqslant 1$ 时上同调真.
\end{enumerate}
\end{corollary}

设 $F$ 是集合层 (resp. ind-$p'$-群层). 由 SGA 4 IX 2.7.2 可以把 $F$ 写成滤过归纳极限
\[
F = \varinjlim_i F_i,
\]
其中诸 $F_i$ 是可构集合层 (resp. ind-$p'$-群层). 由于 $f$ 相容，$f_*$ (resp. $R^1 f_*$) 与归纳极限交换 (SGA 4 VII 3.3). 若知道诸 $(F_i, f)$ 在维数 $\leqslant 0$ (resp. 每个集合层在维数 $\leqslant 0$ 且诸 $(F_i, f)$ 在维数 $\leqslant 1$) 时上同调真，则 $(F, f)$ 也如此. 可假设 $F$ 可构.

若假设 $f$ 有限型 (resp. 满足 a))，命题由 \ref{thm:XIII.3.1} 1) b) (resp. \ref{thm:XIII.3.1} 2) b)) 得出. 现在证明不再假设 $f$ 有限型时的 3.2. 对每个 $k$ 上方的概形 $S'$ 和 $S'$ 的每个几何点 $s$，记 $k$ (resp. $S'$) 为 $k$ 在 $s$ 处 (resp. $S'$ 在 $s$ 处) 的严格局部化，$X$ 为 $X$ 在 $k$ 上的逆像，考虑笛卡儿方形
\[
\begin{tikzcd}[row sep=large, column sep=large]
X \arrow[d,"{f}"] & X' \arrow[l,"g"] \arrow[d,"{f'}"] \\
k & S' \arrow[l]
\end{tikzcd}
\]
只需证明对每个 $S'$ 和每个 $s$，典范态射
\[
H^0(X, F) \longrightarrow H^0(X', F') \quad \text{(resp. } H^1(X, F) \longrightarrow H^1(X', F')\text{)}
\]
是同构. 只需证明有关系
\[
(\ast) \quad F \simeq g_* g^{*} F \quad \text{(resp. } R^1 g_*(g^{*} F) = 0\text{)}.
\]
这种形式的问题局部于 $X$ 的平展拓扑，可假设 $X$ 仿射；经极限过渡可假设 $X$ 在 $k$ 上有限型. 已知 $(F, f)$ 在维数 $\leqslant 0$ (resp. $\leqslant 1$) 时上同调真，把 $X$ 换成平展于 $X$ 上的有限型概形时也如此，这证明 $(\ast)$.

\begin{theorem}\label{thm:XIII.3.3}
设 $S$ 是生成点为 $s$ 的不可约概形，$f \colon X \to S$ 是有限表现态射. 假设 $k$ 的代数闭包 $k$ 上维数 $\leqslant \dim X_s$ 的有限型概形可消解奇点 (EGA IV 7.9.1). 则若 $L$ 表示与 $S$ 的剩余特征不同的素数集，可以找到 $S$ 的非空开集 $S_1$，使态射 $f|_{S_1}$ 对 $L$ 万有局部 1-非球.
\end{theorem}

可假设 $S$ 整且 $X$ 既约 (SGA 4 VIII 1.1). 经极限过渡可假设 $S$ 诺特. 而且只需在有限扩张 $S_1 \to S$ 后证明定理，其中 $S_1$ 是整概形，$S_1 \to S$ 由平展扩张和根式满射扩张复合而成.

先证明，把 $S$ 限制到非空开集后 $f$ 万有局部 0-无环. 作 $k(s)$ 的根式扩张后，可假设态射 $(X_s)_{\mathrm{r\acute{e}d}} \to s$ 可分；故可假设 (把 $S$ 限制到非空开集并作 $S$ 的根式满射扩张后) 态射 $f$ 平坦、几何纤维可分 (EGA IV 12.1.1)，这蕴含 $f$ 万有 0-无环 (SGA 4 XV 4.1).

证明，把 $S$ 限制到非空开集后 $f$ 对 $L$ 万有局部 1-非球. 作 $k(s)$ 的有限扩张 (可把它看作平展扩张与根式扩张的复合) 后，可以找到紧合满射态射 $p_s \colon Y_s \to X_s$，其中 $Y_s$ 是与 $X_s$ 同维数、在 $S$ 上光滑的概形，且可假设 $p_s$ 来自紧合满射态射 $p \colon Y \to X$，其中 $Y$ 在 $S$ 上光滑 (EGA IV 9.6.1 和 12.1.6). 只需证明，把 $S$ 限制到非空开集后，对每个方形笛卡儿的图式
\[
\begin{tikzcd}[row sep=large, column sep=large]
S'' \arrow[d,"{i}"] & X'' \arrow[l,"{f''}"] \arrow[d,"{j}"] \\
S' \arrow[d] & X' \arrow[l,"{f'}"] \arrow[d] \\
S & X \arrow[l,"{f}"]
\end{tikzcd}
\]
其中 $i$ 有限表现平展，和 $S''$ 上每个 ind-$L$-群层 $F$，若 $\Phi$ 是 $F$ 下的挠子叠，典范态射
\[
f'^{*} i_* \Phi \longrightarrow j_* f''^{*} \Phi
\]
是等价. 设 $Z = Y \times_X Y$，$T = Y \times_X Y \times_X Y$. 自然地有交换图式
\[
\begin{tikzcd}[row sep=large, column sep=large]
S'' \arrow[d,"{i}"] & X'' \arrow[l,"{f''}"] \arrow[d,"{j}"] & Y'' \arrow[l,"{p''}"] \arrow[d,"{p'}"] & Z'' \arrow[l] \arrow[d] & T'' \arrow[l] \arrow[d] \\
S' \arrow[d] & X' \arrow[l,"{f'}"] \arrow[d,"{f}"] & Y' \arrow[l,"{p'}"] \arrow[d,"{p}"] & Z' \arrow[l] \arrow[d] & T' \arrow[l] \arrow[d] \\
S & X \arrow[l] & Y \arrow[l] & Z \arrow[l] & T
\end{tikzcd}
\]
设 $q \colon Z \to X$ 和 $r \colon T \to X$ 是典范态射，$q'$，$r'$，$q''$，$r''$ 的记号显然. 由 \ref{exa:XIII.1.11} 2)，有行正合的本质上交换图式：
\[
\begin{tikzcd}[row sep=large, column sep=large]
f'^{*} i_* \Phi \arrow[r] \arrow[d,"{a}"] & p'_* p'^{*}(f'^{*} i_* \Phi) \arrow[r, Rightarrow] \arrow[d,"{b}"] & q'_* q'^{*}(f'^{*} i_* \Phi) \arrow[r, Rightarrow] \arrow[d,"{c}"] & r'_* r'^{*}(f'^{*} i_* \Phi) \arrow[d,"{d}"] \\
j_* f''^{*} \Phi \arrow[r] & j_* p''_* p''^{*}(f''^{*} \Phi) \arrow[r, Rightarrow] & j_* q''_* q''^{*}(f''^{*} \Phi) \arrow[r, Rightarrow] & j_* r''_* r''^{*}(f''^{*} \Phi)
\end{tikzcd}
\]
由于 $Y$ 在 $S$ 上光滑，态射 $fp$ 对 $L$ 万有局部 1-非球 (SGA 4 XV 2.1)，由 \cite[VII 2.1.7]{exp13-2} 得出 $b$ 是范畴等价. 另一方面，把 $S$ 限制到非空开集后可假设态射 $Z \to S$ 和 $T \to S$ 万有局部 0-无环. 由此得出函子 $c$ 和 $d$ 忠实全，上面图式于是证明 $a$ 是等价，这完成证明.

\begin{corollary}\label{cor:XIII.3.4}
设 $k$ 是特征 $p \geqslant 0$ 的域，$p'$ 是与 $p$ 不同的素数集，$f \colon X \to k$ 是相容态射. 假设满足下列条件之一：
a) $f$ 有限型，且 $k$ 的代数闭包上维数 $\leqslant \dim X$ 的有限型概形可消解奇点.
b) $k$ 的代数闭包上有限型概形可消解奇点.
则 $f$ 对 $p'$ 万有局部 1-非球.
\end{corollary}

情形 a) 由 \ref{thm:XIII.3.3} 得出. 在情形 b)，问题局部于 $X$，可假设 $X$ 仿射；经极限过渡 (SGA 4 XV 1.3) 归结到 $X$ 在 $k$ 上有限型的情形.

\begin{corollary}\label{cor:XIII.3.5}
设 $S$ 是生成点为 $s$ 的不可约概形，$f \colon X \to S$ 是有限表现态射. 假设 $k(s)$ 的代数闭包 $k$ 上维数 $\leqslant \dim X_s$ 的有限型概形强可消解奇点 (SGA 5 I 3.1.5). 若 $L$ 表示与 $S$ 的剩余特征不同的素数集，可以找到 $S$ 的非空开集 $S_1$，使对 $S_1$ 的几何点的每个特殊化 $s_1 \to s_2$，特殊化态射 (2.10)
\[
\pi^L_1(X_{s_1}) \longrightarrow \pi^L_1(X_{s_2})
\]
双射.
\end{corollary}

由 \ref{thm:XIII.3.1} 和 \ref{thm:XIII.3.3}，可 (把 $S$ 限制到非空开集后) 假设 $f$ 对 $L$ 局部 1-非球，且对 $X$ 上每个有限常 $L$-群层 $F$，$(F, f)$ 在维数 $\leqslant 1$ 时上同调真. 由 \ref{prop:XIII.1.14} 得出，对 $S_1$ 的几何点的每个特殊化 $s_1 \to s_2$，特殊化态射
\[
(R^1 f_* F)_{s_2} \longrightarrow (R^1 f_* F)_{s_1}
\]
双射. 该推论正是前面结果用基本群语言的翻译.

\section{同伦正合序列}

\refstepcounter{theorem}\addtocounter{theorem}{-1}\noindent\textbf{\thesection.0.}\quad 设 $X$ 和 $S$ 是两个连通概形，$f \colon X \to S$ 是态射，$a$ 是 $X$ 的几何点，$L$ 是素数集. 设 $K$ 是典范同态 $\pi_1(X, a) \to \pi_1(S, a)$ 的核，$N$ 是使 $K/N$ 是 pro-$L$-群 $K^L$ 的 $K$ 的最小 pro-正规子群. 则 $N$ 在 $\pi_1(X, a)$ 中正规，记
\[
\pi'_1(X, a)
\]
为 $\pi_1(X, a)$ 除以 $N$ 的商. 若 $a$ 是几何纤维 $X_s$ 的几何点，典范态射
\[
\pi_1(X_s, a) \longrightarrow \pi_1(X, a) \longrightarrow \pi_1(S, a)
\]
可以定义典范态射
\[
\pi^L_1(X_s, a) \xrightarrow{u} \pi'_1(X, a)\index[notations]{$\pi_1'(X, a)$} \xrightarrow{v} \pi_1(S, a)
\]
有 $vu = 0$.

\begin{proposition}\label{prop:XIII.4.1}
设 $S$ 是连通概形，$f \colon X \to S$ 是局部 0-无环 (SGA 4 XV 1.11) 的态射；再假设 $f$ 0-无环 (当 $f$ 相容时，这等价于 $f$ 的几何纤维连通 (SGA 4 XV 1.16)). 设 $L$ 是素数集. 若 $S'$ 是平展于 $S$ 上的概形，记 $X'$，$f'$ 为 $X$，$f$ 在 $S'$ 上的逆像. 假设对 $S$ 的每个平展覆盖 $S'$ 和 $X'$ 的每个平展覆盖 $E$ (它是群为 $L$-群的 Galois 覆盖的商)，$(E, f')$ 相对 $S'$ 在维数 $\leqslant 0$ 时上同调真，且 $f'_* E$ 可构. 则若 $s$ 是 $S$ 的几何点，$a$ 是纤维 $X_s$ 的几何点，群同态序列
\[
(4.1.1) \quad \pi^L_1(X_s, a) \xrightarrow{u} \pi'_1(X, a) \xrightarrow{v} \pi_1(S, a) \longrightarrow 1
\]
正合.
\end{proposition}

这个断言推广 X.1.4，将照抄它的证明.

先证明 $v$ 满射. 只需证明对 $S$ 的每个连通平展覆盖 $S'$，$X'$ 也连通 (\cref{prop:V.6.9}). 设 $C$ 是至少有两个元素的集合. 由 $f$ 0-无环的事实得出典范态射
\[
H^0(S', C_{S'}) \longrightarrow H^0(X', C_{X'})
\]
双射，故 $X'$ 连通，从而 $v$ 满射.

由 $K^L$ 的定义 (4.0)，有正合序列
\[
1 \longrightarrow K^L \longrightarrow \pi'_1(X, a) \longrightarrow \pi_1(S, a) \longrightarrow 1
\]
设 $\widetilde{S}$ 是 $S$ 的万有覆盖，$\widetilde{X} = \widetilde{S} \times_S X$；群 $K^L$ 分类群为 $L$-群的 Galois 覆盖 $P$，使存在 $S$ 的平展覆盖 $S'$ 和 $X' = X \times_S S'$ 的 Galois 覆盖 $Q$，使有同构 $P \simeq Q \times_{X'} \widetilde{X}$. 为使序列 (4.1.1) 正合，必要且充分条件是典范态射
\[
\pi^L_1(X_s, a) \longrightarrow K^L
\]
满射. 由 $K^L$ 的解释，这等价于：对 $S$ 的每个平展覆盖 $S'$ 和 $X'$ 的每个群为 $L$-群的 Galois 覆盖 $Q$，若 $P = Q \times_{X'} \widetilde{X}$ 连通，则 $Q|_{X_s}$ 连通. 证明这最后一个条件满足. 设 $S'$ 是 $S$ 的平展覆盖，$Q$ 是 $X'$ 的群为 $L$-群 $F$ 的 Galois 覆盖，使 $Q|_{X_s}$ 不连通，证明 (把 $S'$ 换成平展覆盖后) $Q$ 变成不连通. 存在 $F$ 的与 $F$ 不同的子群 $G$ 和 $G|_{X_s}$ 下的挠子 $R$，使 $Q|_{X_s}$ 由 $R$ 经结构群扩张 $G \to F$ 得到. $X'$ 的平展覆盖 $E = Q/G$ 使 $E|_{X_s}$ 有截面. 由 \ref{cor:XIII.1.16}，$f_* E$ 局部常可构，把 $S'$ 换成平展覆盖后甚至可假设 $f_* E$ 常. 由于 $(E, f')$ 相对 $S'$ 在维数 $\leqslant 0$ 时上同调真，且 $H^0(X_s, E|_{X_s})$ 非空，看到 $E$ 有截面. 但这证明 $Q$ 不连通，完成证明.

由 \ref{prop:XIII.4.1} 推出下列将在 \ref{prop:XIII.4.6} 使用的引理.

\begin{lemma}\label{lem:XIII.4.2}
设 $S$ 是连通概形，$f \colon X \to S$ 是局部 0-无环且 0-无环的态射，$L$ 是素数集. 假设对 $X$ 上每个有限常 $L$-群层 $F$，$(F, f)$ 相对 $S$ 在维数 $\leqslant 1$ 时上同调真，且对 $S$ 的每个平展覆盖 $S'$ 和 $X'$ 的每个平展覆盖 $E$ (它是群为 $L$-群的 Galois 覆盖的商)，$f'_* E$ 可构. 则若 $s$ 是 $S$ 的几何点，$a$ 是纤维 $X_s$ 的几何点，群同态序列
\[
\pi^L_1(X_s, a) \longrightarrow \pi'_1(X, a) \longrightarrow \pi_1(S, a) \longrightarrow 1
\]
正合.
\end{lemma}

\ref{prop:XIII.4.1} 的假设满足. 确实由 \ref{prop:XIII.1.13} 3) 得出，对每个平展于 $S$ 上的概形 $S'$ 和 $X'$ 的每个平展覆盖 $E$ (它是群为 $L$-群的 Galois 覆盖的商)，$(E, f')$ 相对 $S'$ 在维数 $\leqslant 0$ 时上同调真.

\begin{proposition}\label{prop:XIII.4.3}
设 $S$ 是连通概形，$L$ 是素数集，$f \colon X \to S$ 是 0-无环、对 $L$ 局部 1-非球 (SGA 4 XV 1.11) 的态射，$g \colon S \to X$ 是 $f$ 的截面. 设 $s$ 是 $S$ 的几何点，$a$ 是纤维 $X_s$ 的几何点. 假设对每个有限常 $L$-群层 $F$，$(F, f)$ 在维数 $\leqslant 1$ 时上同调真，$F$ 下的挠子叠经 $f$ 的直像是 \emph{1-可构叠} (0)，且对 $X'$ 的每个平展覆盖 $E$ (它是群为 $L$-群的 Galois 覆盖的商)，$f'_* E$ 可构. 则群同态序列
\[
(4.3.1) \quad 1 \longrightarrow \pi^L_1(X_s, a) \xrightarrow{u} \pi'_1(X, a) \xrightarrow{v} \pi_1(S, a) \longrightarrow 1
\]
正合.
\end{proposition}

考虑到 \ref{lem:XIII.4.2}，只需证明态射 $u$ 单射，i.e. 证明：对 $X_s$ 的每个群为 $L$-群 $C$ 的主覆盖 $Z$，存在 $X$ 的平展覆盖 $\widetilde{Z}$ 和 $Z|_{X_s}$ 的某连通分量到 $Z$ 的态射 (\cref{cor:V.6.8}). 设 $Z$ 是 $X_s$ 的群为 $L$-群 $C$ 的主覆盖，$z$ 是它在 $H^1(X_s, C_{X_s})$ 中的类. 由 \ref{rem:XIII.1.5} d)，有典范同构
\[
R^1 f_* C_X(s) \xrightarrow{\sim} H^1(X_s, C_{X_s}),
\]
且由 \ref{cor:XIII.1.16}，$R^1 f_* C_X$ 是局部常可构层. 可以找到 $S$ 的平展覆盖 $S'$，使 $R^1 f_* C_X|_{S'}$ 常. 若 $s \to S'$ 是位于几何点 $s \to S$ 上方的几何点，存在 $H^0(S', R^1 f_* C_X)$ 的元素 $z$，使它在 $H^1(X_s, C_{X_s})$ 中的像是 $z$. 由下面\cref{lem:XIII.4.3.1}，可以找到 $S'$ 的平展覆盖 $S'_1$ 和 $X'_1$ 上群 $C$ 的挠子 $P$，使它在 $H^0(S'_1, R^1 f_* C)$ 中的像等于 $z$ 的限制. 挠子 $P$ 可由 $X'_1 = X \times_S S'_1$ 的平展覆盖 $Z$ 表示，使 $Z \times_{X'_1} X_s$ 同构于 $Z$. 把 $Z$ 看作 $X$ 的平展覆盖，则有 $Z \times_X X_s$ 到 $Z$ 的态射，完成证明.

\begin{lemma}\label{lem:XIII.4.3.1}
设 $f \colon X \to S$ 是 0-无环且局部 0-无环的态射，$g$ 是 $f$ 的截面. 设 $C$ 是使 $(C_X, f)$ 在维数 $\leqslant 0$ 时上同调真、$C_X$ 下的挠子叠经 $f$ 的直像可构的有限常群. 则对 $H^0(S, R^1 f_* C_X)$ 的每个截面 $z$，可以找到 $S$ 的平展覆盖 $S_1$，且若 $X_1 = X \times_S S_1$，存在 $H^1(X_1, C_{X_1})$ 的元素，使它在典范态射
\[
H^1(X_1, C_{X_1}) \longrightarrow H^0(S_1, R^1 f_* C_{X_1})
\]
下的像等于 $z$ 在 $H^0(S_1, R^1 f_* C_{X_1})$ 中的限制.
\end{lemma}

对每个平展于 $S$ 上的概形 $S'$，令 $X' = X \times_S S'$，用 $g'$ (resp. $F'$，等等) 记 $g$ (resp. $F$，等等) 经态射 $S' \to S$ 的逆像. 由
\[
G(S') = \{X' \text{ 上群 } C_{X'} \text{ 的配备同构 } g'^{*} P \xrightarrow{\sim} C_{S'} \text{ 的挠子 } P \text{ 的模同构的类}\}
\]
定义的 $S$ 上预层 $G$ 是层. 这由下降得出，因为 $X'$ 上挠子 $P$ 的同构由它在 $g'(S')$ 上的限制确定. 而且有满态射
\[
G \longrightarrow R^1 f_* F.
\]
设 $z$ 是 $H^0(S, R^1 f_* C_X)$ 的元素，$H$ 是 $G$ 的 $z$ 的逆像子层. 只需证明 $H$ 是局部常可构层. 该性质局部于 $S$，可假设 $z$ 来自 $H^1(X, C_X)$ 的由使 $g^{*} P$ 同构于 $C_S$ 的挠子 $P$ 表示的元素. 给出同构 $i \colon g^{*} P \xrightarrow{\sim} C_S$ 等价于给出 $\mathrm{Aut}_{C_S}(g^{*} P)$ 的全局截面，且两个同构 $i$ 和 $i'$ 定义 $G(X)$ 的同一个元素，当且仅当 $ii'^{-1}$ 是 $\mathrm{Aut}_{C_X}(P)$ 的某个元素的像. 若考虑典范单射
\[
f_* \mathrm{Aut}_{C_X}(P) \longrightarrow \mathrm{Aut}_{C_S}(g^{*} P) \simeq C,
\]
$H$ 等同于 $\mathrm{Aut}_{C_S}(g^{*} P)$ 除以 $f_* \mathrm{Aut}_{C_X}(P)$ 的商. 由 \ref{cor:XIII.1.16}，$f_* \mathrm{Aut}_{C_X}(P)$ 局部常；故 $H$ 也局部常，完成证明.

\begin{example}\label{exa:XIII.4.4}
注意，若 $S$ 是连通概形，态射 $f$ 紧合平坦有限表现、几何纤维可分连通时 \ref{prop:XIII.4.1} 的假设满足，$L$ 任意 (参见 X 1.3). 若 $f$ 还光滑且有截面，$L$ 为与 $S$ 的剩余特征不同的素数集 (SGA 4 XV 2.1 和 XVI 5.2) 时，\ref{prop:XIII.4.3} 的假设满足.

若 $S$ 连通，有紧合有限表现概形 $Z$ 在 $S$ 上平坦、几何纤维可分连通，$X$ 是 $Z$ 中相对 $S$ 正规交叉除子的余集，$L$ 为与 $S$ 的剩余特征不同的素数集 (\cref{cor:XIII.2.9}) 时，\ref{prop:XIII.4.1} 的假设也满足. 若 $f$ 还光滑且有截面，\ref{prop:XIII.4.3} 的假设满足.
\end{example}

\refstepcounter{theorem}\noindent\textbf{\thesection.\arabic{theorem}.}\quad 沿用 \ref{prop:XIII.4.3} 的记号和假设. 若 $s$ 是 $S$ 的几何点，$a = g(s)$，截面 $g$ 可以定义态射
\[
w \colon \pi_1(S, a) \longrightarrow \pi'_1(X, a),
\]
使 $\pi'_1(X, a)$ 等同于 $\pi_1(S, a)$ 被 $\pi_1(X_s, a)$ 的半直积. pro-有限群 $\pi_1(S, a)$ 于是作用在 $\pi_1(X_s, a)$ 上. 由于 $\pi^L_1(X_s, a)$ 是 $\pi_1(S, a)$ 的作用不变群的严格投射极限，配备该作用的 $\pi^L_1(X_s, a)$ 的给出等价于 $S$ 上有限平展群概形的严格投射系统的给出，记为
\[
\pi^L_1(X/S, g, s) \quad \text{或简记}\index[notations]{$\pi_1^L(X/S, g, s)$} \quad \pi^L_1(X/S, g).
\]
有下列性质：

4.5.1. 对 $S$ 上每个纤维为 $L$-群的有限平展群概形 $G$，配备同构 $g^{*} P \xrightarrow{\sim} G$ 的 $G$ 在 $X$ 上的逆像 $G_X$ 下的挠子 $P$ 的类的集 $E$，典范同构于集
\[
\mathrm{Hom}_S(\pi^L_1(X/S, g, s), G) \quad \text{模 } G \text{ 的内自同构}.
\]

4.5.2. 对 $S$ 上每个纤维为 $L$-群的有限平展群概形 $G$，层 $R^1 f_* G_X$ 典范同构于与预层
\[
S' \longmapsto \mathrm{Hom}_{S'}(\pi^L_1(X/S, g, s), G) \text{ 模 } G \text{ 的内自同构}
\]
关联的层. ($S'$ 表示平展于 $S$ 上的概形).

4.5.3. 设 $S'$ 是连通 $S$-概形，$s$ 是 $S'$ 的几何点，$X'$，$g'$ 是 $X$，$g$ 在 $S'$ 上的逆像. 则 $\pi^L_1(X'/S', g', s)$ 典范同构于 $\pi^L_1(X/S, g, s)$ 在 $S'$ 上的逆像. 对 $S$ 的每个几何点 $\xi$，纤维 $\pi^L_1(X/S, g, s)_\xi$ 同构于 $\pi^L_1(X_\xi)$.

$G$ 的给出确实等价于配备 $\pi_1(S, a)$ 的作用的抽象 $L$-群 $G$ 的给出，从而有 $\pi'_1(X, a)$ 在 $G$ 上的作用. 4.5.1 中定义的同构由典范态射
\[
H^1(\pi'_1(X, a), G) \longrightarrow H^1(\pi^L_1(X_s, a), G) = \mathrm{Hom}(\pi^L_1(X_s, a), G)/\mathrm{aut.\,int.}\, G
\]
限制到子集 $E$ 得到，$E$ 双射到与 $\pi_1(S, a)$ 的作用相容的 $\pi^L_1(X_s, a)$ 到 $G$ 的态射子集. 断言 4.5.3 由 $\pi^L_1(X/S, g, s)$ 的定义，考虑到同伦正合序列 (\cref{lem:XIII.4.3.1}) 得出，4.5.2 由 4.5.1 和 4.5.3 推出.

\begin{proposition}\label{prop:XIII.4.6}
[Künneth 公式] 设 $k$ 是特征 $p \geqslant 0$ 的可分闭域，$X$ 和 $Y$ 是两个连通 $k$-概形，$a$ 是 $X$ 的几何点，$b$ 是 $Y$ 的几何点，$c$ 是 $X \times_k Y$ 的位于 $a$ 和 $b$ 上方的几何点. 假设满足下列条件之一：
a) $X$ 在 $k$ 上有限型，且 $k$ 的代数闭包 $k$ 上维数 $\leqslant \dim X$ 的有限型概形强可消解奇点 (SGA 5 I 3.1.5).
b) $X$ 拟紧拟分离，且 $k$ 上每个有限型概形强可消解奇点.
则若 $p'$ 是与 $p$ 不同的素数集，态射
\[
(4.6.0) \quad \pi^{p'}_1(X \times_k Y, c) \longrightarrow \pi^{p'}_1(X, a) \times \pi^{p'}_1(Y, b),
\]
由与投影
\[
X \times_k Y \longrightarrow X \quad \text{和} \quad X \times_k Y \longrightarrow Y
\]
关联的基本群同态得出，是同构.
\end{proposition}

可假设 $k$ 代数闭且 $X$ 既约 (SGA 4 VIII 1.1). 设 $Z = X \times_k Y$，$g \colon X \to k$ 和 $f \colon Z \to Y$ 是典范态射. 由 \ref{cor:XIII.3.5}，态射 $g$ 从而 $f$ 对 $p'$ 万有局部 1-非球. 由于 $X$ 连通，$f$ 0-无环 (SGA 4 XV 1.16). 另一方面由 \ref{cor:XIII.3.2} 得出，对每个 $p'$-有限群 $C$，$(C_X, g)$ 相对 $k$ 在维数 $\leqslant 1$ 时上同调真. 由此得出 $(C_Z, f)$ 相对 $Y$ 在维数 $\leqslant 1$ 时上同调真 (1.5 c))，且 $f_* C_Z$ 和 $R^1 f_* C_Z$ 是常层. 从而 $g$ 满足 \ref{lem:XIII.4.2} 的所有假设. 有正合序列
\[
\pi^{p'}_1(X_b, c) \longrightarrow \pi^{p'}_1(Z, c) \longrightarrow \pi^{p'}_1(Y, b) \longrightarrow 1.
\]
而且复合态射
\[
\pi^{p'}_1(X_b, c) \longrightarrow \pi^{p'}_1(Z, c) \longrightarrow \pi^{p'}_1(Y, b)
\]
是同构，故有正合序列
\[
1 \longrightarrow \pi^{p'}_1(X, a) \longrightarrow \pi^{p'}_1(Z, c) \longrightarrow \pi^{p'}_1(Y, b) \longrightarrow 1.
\]
另一方面态射 (4.6.0) 可以定义该正合序列到正合序列
\[
1 \longrightarrow \pi^{p'}_1(X, a) \longrightarrow \pi^{p'}_1(X, a) \times \pi^{p'}_1(Y, b) \longrightarrow \pi^{p'}_1(Y, b) \longrightarrow 1
\]
的态射，由此得出态射 (4.6.0) 是同构.

\refstepcounter{theorem}\noindent\textbf{\thesection.\arabic{theorem}.}\quad 设
\[
\begin{tikzcd}[row sep=large, column sep=large]
X \arrow[d,"{f}"] & X_U \arrow[l] \arrow[d,"{f_U}"] & X_s \arrow[l] \arrow[d] \\
S & U \arrow[l] & s \arrow[l]
\end{tikzcd}
\]
是所有方形笛卡儿的图式，其中 $S$ 是弧连通概形 (SGA 4 IX 2.12)，$U$ 是 $S$ 的连通开集，$s$ 是 $U$ 的几何点. 设 $a$ 是 $X_s$ 的几何点，$L$ 是素数集. 设 $g$ 是 $f$ 的截面，假设满足下列条件：
a) 态射 $f$ 0-无环、局部 0-无环，且对 $S$ 的每个平展覆盖 $S'$ 和 $X \times_S S'$ 的每个平展覆盖 $E$ (它是群为 $L$-群的 Galois 覆盖的商)，$(E, f_{(S')})$ 相对 $S'$ 在维数 $\leqslant 0$ 时上同调真.
b) 态射 $f_U$ 对 $L$ 局部 1-非球，且对 $X_U$ 上每个有限常 $L$-群层 $F$，$(F, f_U)$ 在维数 $\leqslant 1$ 时上同调真，$R^1 f_{U*} F$ 的纤维有限.
由 \ref{prop:XIII.4.1} 和 \ref{prop:XIII.4.3} 得出下列行正合的交换图式：
\begin{equation}
\begin{tikzcd}[row sep=large, column sep=large]
1 \arrow[r] & \pi^L_1(X_s, a) \arrow[r] \arrow[d] & \pi'_1(X_U, a) \arrow[r] \arrow[d] & \pi_1(U, a) \arrow[r] \arrow[d] & 1 \\
& \pi^L_1(X_s, a) \arrow[r] & \pi'_1(X, a) \arrow[r] & \pi_1(S, a) \arrow[r] & 1
\end{tikzcd}
\tag{4.7.0}
\label{eq:XIII.4.7.0}
\end{equation}
借助截面 $g$，有态射
\[
\pi_1(U, a) \longrightarrow \pi'_1(X_U, a), \quad \pi_1(S, a) \longrightarrow \pi'_1(X, a)
\]
由此得到 $\pi_1(S, a)$ 和 $\pi'_1(X_U, a)$ 在 $\pi_1(U, a)$ 上方的余直和到 $\pi'_1(X, a)$ 的态射：
\[
(4.7.1) \quad \phi \colon \pi = \pi_1(S, a) \amalg_{\pi_1(U, a)} \pi'_1(X_U, a) \longrightarrow \pi'_1(X, a).
\]
假设满足下列条件：
c) 若 $T = S - U$，有 $\mathrm{prof}^{\mathrm{ét}}_T(S) \geqslant 2$ (SGA 2 XIV 1.1).
则把 $S$ 的平展覆盖配给它在 $U$ 上的限制的函子忠实全 (SGA 2 XVI 1.4). 由此得出态射 $\pi_1(U, a) \to \pi_1(S, a)$ 满射 (\cref{prop:V.6.9})，由 \eqref{eq:XIII.4.7.0} 图式得出态射 $\pi'_1(X_U, a) \to \pi'_1(X, a)$ 也满射；更不用说 $\phi$ 是满态射. 设
\[
(4.7.2) \quad K = \mathrm{Ker}(\pi_1(U, a) \longrightarrow \pi_1(S, a)).
\]
(4.7.1) 中的群 $\pi$ 等同于 $\pi'_1(X_U, a)$ 除以 $K$ 的像 $L$ 在 $\pi'_1(X_U, a)$ 中生成的闭不变子群的商. 把 $\pi'_1(X_U, a)$ 看作 $\pi_1(U, a)$ 被 $\pi^L_1(X_s, a)$ 的半直积. 群 $K$ 于是经内自同构作用在 $\pi^L_1(X_s, a)$ 上，商 $\pi = \pi'_1(X_U, a)/L$ 等同于 $\pi_1(U, a)/K = \pi_1(S, a)$ 被 $\pi^L_1(X_s, a)$ 在 $K$ 下的余不变量群 $\pi^L_1(X_s, a)_K$ 的半直积. 最终有满态射
\[
(4.7.3) \quad \phi \colon \pi = \pi^L_1(X_s, a)_K\index[notations]{$\pi_1^L(X_s, a)_K$} \cdot \pi_1(S, a) \longrightarrow \pi'_1(X, a).
\]

下列命题给出 $\phi$ 是同构的条件.

\begin{proposition}\label{prop:XIII.4.7.4}
记号同 4.7. 假设除条件 a)，b)，c) 外还满足下列条件：
d) 对 $T = S - U$ 的每个点 $t$，态射 $f$ 在 $g(t)$ 处对 $L$ 局部 1-非球.
e) 对每个点 $t \in T$，纤维 $X_t$ 的每个不可约分量含 $g(t)$，且对 $X_t - \{g(t)\}$ 的每个非极大点 $x$，有
\[
\mathrm{prof}^{\mathrm{hop}}_x(X) \geqslant 3 \quad \text{(SGA 2 XIV 1.2)}
\]
且环 $\mathcal{O}_{X, x}$ 诺特.
则态射 (4.7.3) 是同构.
\end{proposition}

如前所说，群 $\pi$ 等同于 $\pi'_1(X_U, a)$ 除以 $K$ (4.7.2) 的像在 $\pi'_1(X_U, a)$ 中生成的闭不变子群 $L$ 的商. 这等价于说 $\pi$ 分类 $X_U$ 的主覆盖 $Z$ (使 $g^{*}_U(Z)$ 延拓成 $S$ 的平展覆盖，且 $Z$ 在 $X_s$ 上诱导的覆盖由群为 $L$-群的主覆盖经结构群扩张得到). 为证明 $\phi$ 是同构，只需证明这样的覆盖 $Z$ 延拓到整个 $X$.

先证明 $Z$ 延拓成含 $X_U$ 和 $g(S)$ 的开集的平展覆盖. 设 $W$ 是平展于 $X$ 上的概形，它的像含 $X_U$ 和 $g(S)$，令 $W_U = W \times_S U$. 由态射 $W \to S$ 0-无环和关系 $\mathrm{prof}^{\mathrm{ét}}_T S \geqslant 2$ 得出
\[
\mathrm{prof}^{\mathrm{ét}}_{(W - W_U)} W \geqslant 2
\]
(SGA 2 XIV 1.13)；从而若 $Z|_{W_U}$ 延拓成 $W$ 的平展覆盖，该延拓在唯一同构意义下唯一. 由此得出把 $Z$ 延拓到 $g(S) \cup X_U$ 的邻域的问题局部于 $g(T)$ 的点处的平展拓扑. 若 $t$ 是 $T$ 的点，令 $x = g(t)$，记 $X$ (resp. $S$) 为 $X$ 在 $x$ 处 (resp. $S$ 在 $t$ 处) 的严格局部化，$U = U \times_S S$，$X_U = X \times_X X_U$，$g \colon S \to X$ 为由 $g$ 导出的态射. 只需证明对 $T$ 的每个点 $t$，$Z$ 在 $X_U$ 上的逆像 $Z$ 延拓到 $X$ (或等价地平凡). 由 $Z$ 的定义，$Z$ 在 $U$ 上的逆像平凡. 为证明 $Z$ 平凡，只需证明它形如 $f^{*}_U E$，其中 $E$ 是 $U$ 的主覆盖；将有 $E \simeq g^{*}_U f^{*}_U E \simeq g^{*}_U Z$，从而由 $g^{*}_U Z$ 平凡得出结果. 由于态射 $f_U$ 0-无环且局部 0-无环，为证明 $Z$ 来自 $U$，只需证明对位于 $U$ 的某点上方 (可假设是点 $s$) 的每个几何点，$Z|_{X_s}$ 平凡 (SGA 4 XV 1.15). 但 $Z|_{X_s}$ 由群为 $L$-群的主覆盖经结构群扩张得到，由态射 $f$ 对 $L$ 1-非球得出结果.

于是证明存在 $g(S) \cup X_U$ 的开邻域 $V$，使 $Z$ 延拓成 $V$ 的平展覆盖 $Z_V$. 证明 $Z_V$ 延拓到整个 $X$. 只需看对 $X - V$ 的每个点 $x$，有
\[
\mathrm{prof}^{\mathrm{hop}}_x X \geqslant 3.
\]
这由假设 e) 和 $X - V$ 的点 $x$ 不能是它的纤维 $X_t$ 的极大点的事实得出，因为 $X_t$ 的每个不可约分量含 $g(t)$，$X_t$ 的每个极大点属于 $V$.

\begin{corollary}\label{cor:XIII.4.8}
假设同 \ref{prop:XIII.4.7.4}，但再假设 $\pi_1(S, a) = 1$. 则有同构
\[
(\ast) \quad \pi^L_1(X, a) \xrightarrow{\sim} \pi^L_1(X_s, a)_K.
\]
特别，若 $\pi^L_1(X_s, a)$ 拓扑有限表现，且 $K$ 经由有限型群作用在 $\pi^L_1(X_s, a)$ 上，则 $\pi^L_1(X, a)$ 拓扑有限表现.
\end{corollary}

同构 $(\ast)$ 已在 \ref{prop:XIII.4.7.4} 中证明. 假设 $\pi^L_1(X_s, a)$ 是 $n$ 个生成元的自由 pro-$L$-群 $L(x_1, \ldots, x_n)$ 除以由 $L(x_1, \ldots, x_n)$ 的元素 $y_1, \ldots, y_p$ 生成的闭不变子群的商，且 $K$ 经由由元素 $k_1, \ldots, k_q$ 生成的群作用. 若对每个 $i \in [1, n]$，$j \in [1, q]$，记 $z_{ij}$ 为 $L(x_1, \ldots, x_n)$ 的提升元素 $(k_j \cdot x_i)x^{-1}_i$，则 $\pi^L_1(X_s, a)_K$ 是 $L(x_1, \ldots, x_n)$ 除以由元素 $(y_i)_{i \in [1,p]}$，$(z_{ij})_{i \in [1,n], j \in [1,q]}$ 生成的闭不变子群的商.

\begin{remark}\label{rem:XIII.4.9}
a) 当 $S$ 是连通正规概形，$U$ 是 $S$ 的稠密拟紧开集，$f$ 是有限表现紧合态射、在 $T$ 的每个点 $t$ 处几何纤维连通且不可约，$f$ 还在 $X_U \cup g(T)$ 的点处可分光滑，$L$ 为与 $S$ 的剩余特征不同的素数集，$X$ 在 $X_t$ 的每个点处正则时，4.7 的条件 a) 至 e) 满足. 条件 a) 确实由 SGA 4 XV 4.1 和 \ref{exa:XIII.1.4} 得出；条件 b) 和 d) 由 \ref{exa:XIII.1.4} 和 SGA 4 XV 2.1 及 XVI 5.2 得出. 最后 e) 由 SGA 2 XIV 1.11 得出.

b) 推论 4.5 可以用来计算特征 $p$ 的可分闭域 $k$ 上紧合光滑曲面 $X$ 的基本群 $\pi^{p'}_1(X)$ ($p'$ 表示与 $p$ 不同的素数集). 该方法由 J.P. Murre 告诉我们；它归结到 (爆破 $X$ 后) 有纤维化 $X \to \mathbf{P}^1_k$ 和满足 4.7 的假设的 $\mathbf{P}^1_k$ 的开集 $U$ 的情形 (参见 SGA 7 的细节). 同样方法可以更一般地 (loc. cit.) 用来证明：若 $X$ 是连通有限型 $k$-概形，且 $k$ 的代数闭包上维数 $\leqslant \dim X$ 的有限型概形强可消解奇点 (SGA 5 I 3.1.5)，则 $\pi^{p'}_1(X)$ 拓扑有限表现.
\end{remark}

\section{附录 I：Abhyankar 引理的变体}

本附录含 Abhyankar 引理的不同变体.

\begin{proposition}\label{prop:XIII.5.1}
设 $X = \mathrm{Spec}\, A$ 是正则局部概形，$D = \sum_{1 \leqslant i \leqslant r} \mathrm{div}\, f_i$ 是正规交叉除子，其中诸 $f_i$ 是 $A$ 的极大理想的属于正则参数系的元素. 设诸 $n_i$ ($1 \leqslant i \leqslant r$) 是 $\geqslant 0$ 的整数，令
\[
X' = X[T_1, \ldots, T_r]/(T_1^{n_1} - f_1, \ldots, T_r^{n_r} - f_r),
\]
$U' = U \times_X X'$. 则 $X'$ 正则，且 $U'$ 是 $X'$ 中正规交叉除子 $\sum_{1 \leqslant i \leqslant r} \mathrm{div}\, T_i$ 的余集. 若诸整数 $n_i$ 与 $X$ 的剩余特征 $p$ 互素，$U'$ 是 $U$ 的相对 $D$ 温和分歧 (2.3 c)) 的连通平展覆盖.
\end{proposition}

确实 $X'$ 是极大理想由 $T_1, \ldots, T_r$ 生成的局部环 $A'$ 的谱. (可假设 $r = \dim(A)$.) 由于 $A'$ 在 $A$ 上有限且平坦，故维数 $r$，$A'$ 正则 (EGA IV$_0$ 17.1.1)，诸 $T_i$ 构成 $A'$ 的正则参数系. 假设诸 $n_i$ 与 $p$ 互素. 由于诸 $f_i$ 在 $U$ 上可逆，$U'$ 在 $U$ 上平展的事实由 I 7.4 得出. 而且 $U'$ 相对 $D$ 温和分歧；设 $x_i$ 是 $V(f_i)$ 的生成点，$R$ 是 $R = \mathcal{O}_{X, x_i}$ 的严格局部化，$K$ 是 $R$ 的分式域. 则表示 $U'|_K$ 的 $K$-代数由域 $K[T_i]/(T_i^{n_i} - f_i)$ 经非分歧扩张得到；故它在 $R$ 上温和分歧.

\begin{proposition}\label{prop:XIII.5.2}
[\emph{绝对 Abhyankar 引理}]\index[term]{juedui Abhyankar yinli@绝对 Abhyankar 引理} 设 $X$ 是正则局部概形，
\[
D = \sum_{1 \leqslant i \leqslant r} \mathrm{div}\, f_i
\]
是如 \ref{prop:XIII.5.1} 的正规交叉除子，$Y = \mathrm{Supp}\, D$，$U = X - Y$. 设 $V$ 是 $U$ 的相对 $D$ 温和分歧的平展覆盖. 若 $x_i$ 是闭集 $V(f_i)$ 的生成点，$\mathcal{O}_{X, x_i}$ 是分式域 $K_i$ 的离散赋值环，且 $V|_{K_i} = \mathrm{Spec}(\prod_{j \in J_i} L_j)$，其中诸 $L_j$ 是 $K_i$ 的有限可分扩张；记 $n_j$ 为 $L_j$ 生成的 伽罗瓦扩张的惯性群的阶，$n_i$ 为 $j$ 遍历 $J_i$ 时诸 $n_j$ 的最小公倍数. 若令
\[
X' = X[T_1, \ldots, T_r]/(T_1^{n_1} - f_1, \ldots, T_r^{n_r} - f_r),
\]
$U' = U_{(X')}$，$V' = V_{(X')}$，等等，$U'$ 的平展覆盖 $V'$ 唯一地 (模唯一同构) 延拓成 $X'$ 的平展覆盖，且诸 $n_i$ 与 $X$ 的剩余特征 $p$ 互素.
\end{proposition}

唯一性由 $X'$ 正规 (\cref{prop:XIII.5.1}) 得出；确实，延拓 $V'$ 的 $X'$ 的平展覆盖同构于 $X'$ 在 $V'$ 在 $X'$ 的生成点处的纤维中的正规化 (\cref{cor:I.10.2}). 若 $x'$ 是 $Y'$ 的几何点，记 $X'$ 为 $X'$ 在 $x'$ 处的严格局部化，$V' = V'_{(X')}$，等等. 经下降，考虑到唯一性，只需证明对 $Y'$ 的每个几何点 $x'$，$U'$ 的平展覆盖 $V'$ 延拓到 $X'$. 给定正则概形 $X'$ 的含所有满足 $\dim \mathcal{O}_{X', x'} \leqslant 1$ 的点 $x'$ 的开集的平展覆盖延拓到整个 $X'$ (SGA 2 XIV 1.11)，可限于 $x'$ 在 $Y'$ 的极大点上的投影. 在这样的点 $x'$，$V'$ 延拓成 $X'$ 的平展覆盖的事实由 3.6 得出.

证明诸 $n_i$ 与 $p$ 互素. 若不然，有例如 $p|n_1$. 把 $X$ 换成
\[
X[T_1, \ldots, T_r]/(T_1^{n_1/p} - f_1, T_2^{n_2} - f_2, \ldots, T_r^{n_r} - f_r),
\]
归结到 $X' = X[T_1]/(T_1^p - f_1)$ 的情形. 只需证明 $V$ 延拓成 $X$ 的平展覆盖，因为此时将有 $n_1 = 1$，与假设矛盾. 为此可假设 $X$ 严格局部. 设 $Z$ 是 $X$ 的方程 $p = 0$ 的闭子概形，$Z_1 = Z \cap X_{f_1}$；$Z_1$ 是 $Z$ 的非空开集. 由前面，$U'$ 的平展覆盖 $V'$ 延拓成 $X'$ 的平展覆盖 $W'$. 设 $W''_1$ 和 $W''_2$ 是 $W'$ 经两个投影 $X'' = X' \times_X X' \rightrightarrows X'$ 的逆像，证明下降同构 $u \colon W''_1|_{U''} \to W''_2|_{U''}$ 延拓成 $X''$-态射 $W''_1 \to W''_2$，它必然是 $W'$ 上相对 $X' \to X$ 的下降数据. 设 $Z''$ (resp. $Z''_1$) 是 $Z$ (resp. $Z_1$) 在 $X''$ 中的逆像. 由于态射 $Z'' \to Z$ 根式，存在延拓同构 $u|_{Z''_1}$ 的同构 $v \colon W''_1|_{Z''} \to W''_2|_{Z''}$. 但 $X$ Hensel 局部，有双射
\[
\mathrm{Hom}_{X''}(W''_1, W''_2) \simeq \mathrm{Hom}_{Z''}(W''_1|_{Z''}, W''_2|_{Z''}),
\]
从而有提升 $v$ 的态射 $w \colon W''_1 \to W''_2$. $X''$ 的 $u$ 和 $w$ 重合的同时开且闭子概形含 $Z''_1$，故等于整个 $X''$，由此得出 $V$ 延拓到 $X$.

5.3.0. 沿用 \ref{prop:XIII.5.2} 的假设和记号，再假设 $X$ 严格局部. 由 loc. cit. 得出 $U$ 的每个相对 $D$ 温和分歧的连通平展覆盖是形如
\[
U' = U[T_1, \ldots, T_r]/(T_1^{n_1} - f_1, \ldots, T_r^{n_r} - f_r)
\]
的温和分歧覆盖的商，其中诸 $n_i$ 是与 $p$ 互素的整数. 设 $\mu_n$ 是 $U$ 的 $n$ 次单位根群. $U'$ 的 $U$-自同构群正是 $\mu_{n_1} \times \cdots \times \mu_{n_r}$，其中 $n_i$ 次单位根 $\xi_i$ 把 $T_i$ 变成 $\xi_i T_i$ 地作用在 $U'$ 上. 故有下列断言：

\begin{corollary}\label{cor:XIII.5.3}
设 $X$ 是剩余特征 $p \geqslant 0$ 的严格局部正则概形，$D = \sum_{1 \leqslant i \leqslant n} \mathrm{div}\, f_i$ 是 $X$ 上的正规交叉除子，$U = X - \mathrm{Supp}\, D$. 令
\[
\widetilde{U} = \varprojlim_{(n_i)} U[T_1, \ldots, T_r]/(T_1^{n_1} - f_1, \ldots, T_r^{n_r} - f_r),
\]
投射极限按与 $p$ 互素的整数族 $n_i > 0$ (对整除关系) 的滤过有序集取. 则 $\widetilde{U}$ 是 $U$ 的温和分歧万有覆盖. 从而 $U$ 的温和分歧基本群为
\[
\pi^t_1(U) \simeq \prod_{\ell \neq p} \mathbb{Z}_\ell[1]^r \quad \text{(典范同构)}
\]
其中令 $\mathbb{Z}_\ell[1] = \varprojlim_{n > 0} \mu_{\ell^n}$\index[notations]{$\mathbb{Z}_\ell[1]$}. 群 $\pi^t_1(U)$ 非典范同构于 $\prod_{\ell \neq p} \mathbb{Z}_\ell^r$.
\end{corollary}

\begin{proposition}\label{prop:XIII.5.4}
设 $f \colon X \to S$ 是概形态射，$D = \sum_{1 \leqslant i \leqslant r} \mathrm{div}\, f_i$ 是相对 $S$ 正规交叉 (2.1) 的除子，其中对 $Y = \mathrm{Supp}\, D$ 的每个点 $x$，若 $I(x) \subset [1, r]$ 是使 $f_i(x) = 0$ 的 $i$ 的集，子概形 $V((f_i)_{i \in I(x)})$ 在 $S$ 上光滑且在 $X$ 中余维数 $\mathrm{card}\, I(x)$. 设 $U = X - Y$. 设 $x$ 是 $Y$ 的点，$X_1 = \mathrm{Spec}\, \mathcal{O}_{X, x}$，$U_1 = U \times_X X_1$，诸 $n_i$ ($i \in I(x)$) 是整数，
\[
X' = X[T_i]_{i \in I(x)}/(T_i^{n_i} - f_i).
\]
则若 $x'$ 是 $X'$ 的位于 $x$ 上方的点，$X'$ 在 $x'$ 处平展于 $S$ 上. 若诸整数 $n_i$ 与 $k(x)$ 的特征 $p$ 互素，$U'_1 = U_1 \times_X X'$ 是 $U_1$ 的相对 $S_1$ 在 $X_1$ 上温和分歧 (\cref{def:XIII.2.1.1}) 的连通平展覆盖.
\end{proposition}

若 $s = f(x)$，几何纤维 $X'_s$ 在点 $x'$ 处正则 (\cref{prop:XIII.5.1})；由于 $X'$ 在 $Y$ 的邻域内在 $S$ 上平坦，这证明 $X'$ 在 $x'$ 处平展于 $S$ 上 (EGA IV 12.1.6). 若诸整数 $n_i$ 与 $p$ 互素，$U'_1$ 是 $U_1$ 的平展覆盖 (\cref{cor:I.7.4})；它相对 $S$ 在 $X_1$ 上温和分歧，因为在 $S$ 的每个点处的几何纤维上如此 (\cref{prop:XIII.5.1}). 最后 $U'_1$ 连通的事实由 (SGA 4 XVI 3.2) 得出.

\begin{proposition}\label{prop:XIII.5.5}
[\emph{相对 Abhyankar 引理}]\index[term]{xiangdui Abhyankar yinli@相对 Abhyankar 引理} 设 $X$ 是 $S$-概形，$D$ 是如 (5.4) 的相对 $S$ 正规交叉除子. 设 $Y = X - \mathrm{Supp}\, D$，$U = X - Y$，$x$ 是 $Y$ 的点，$X_1$ 是 $X$ 在 $x$ 上方的几何点处的严格局部化，$U_1 = U \times_X X_1$，$V_1$ 是 $U_1$ 的平展覆盖. 假设对 $S$ 的每个极大点 $s$，$V_{1s}$ 相对 $s$ 在 $X_{1s}$ 上温和分歧. 则可以找到与 $k(x)$ 的特征 $p$ 互素、$i \in I(x)$ 的整数 $n_i$，使若令
\[
X'_1 = X_1[T_i]_{i \in I(x)}/(T_i^{n_i} - f_i),
\]
$U'_1 = U_1 \times_{X_1} X'_1$，等等，$U'_1$ 的平展覆盖 $V'_1$ 唯一地 (模唯一同构) 延拓成 $X'_1$ 的平展覆盖. 特别 $V_1$ 相对 $S$ 在 $X_1$ 上温和分歧.
\end{proposition}

可假设 $S$ 是闭点 $f(x)$ 的局部诺特局部环的谱. 对 $S$ 的每个极大点 $s$ 和每个 $i \in I(x)$，设 $x_i$ 是纤维 $X_{1s}$ 的闭集 $V(f_i)$ 的生成点. 局部环 $(\mathcal{O}_{X_1, x_i})_{\mathrm{r\acute{e}d}}$ 是分式域 $K_i$ 的离散赋值环，且 $V_1|_{K_i} = \mathrm{Spec}(\prod_{j \in I(x_i)} L_j)$，其中 $L_j$ 是 $K$ 的有限可分扩张；记 $n_j$ 为 $L_j$ 生成的 伽罗瓦扩张的惯性群的阶，$n_i$ 为 $s$ 遍历 $S$ 的极大点、$j \in I(x_i)$ 时诸 $n_j$ 的最小公倍数.

$n_i$ 这样选定后，证明 $V'_1$ 唯一地延拓成 $X'_1$ 的平展覆盖. 唯一性由 $X'$ 在 $Y$ 的点处平展于 $S$ 上 (故 $\mathrm{prof}^{\mathrm{ét}}_{Y'_1}(X'_1) \geqslant 2$ (SGA 4 XVI 3.2 或 SGA 2 XIV 1.19)) 得出. 设 $x'_1$ 是 $Y'_1$ 的点，$x'_1$ 是 $x'_1$ 上方的几何点，记 $X'_1$ 为 $X'_1$ 在 $x'_1$ 处的严格局部化，$U'_1 = U'_{1(X'_1)}$，等等. 经下降，考虑到唯一性，只需证明 $V'_1$ 延拓到 $X'_1$. 而且可以限于 $x'_1$ 为 $Y'_1$ 的极大点；确实此时在 $X'_1$ 的含 $Y'_1$ 的极大点的开集 $W'_1$ 上有 $V'_1$ 的延拓；若 $Z'_1 = X'_1 - W'_1$，当 $s$ 是 $S$ 的极大点时 $\mathrm{codim}(Z'_{1s}, X'_{1s}) \geqslant 2$，当 $s$ 是非极大的 $S$ 的点时 $\mathrm{codim}(Z'_{1s}, X'_{1s}) \geqslant 1$，$\mathrm{prof}^{\mathrm{ét}}_s(S) \geqslant 1$；$V'_1$ 延拓到整个 $X'_1$ 的事实由 (SGA 2 XIV 1.20) 得出. 而在位于 $Y'_1$ 的极大点上方的几何点 $x'_1$ 处，$X'_{1\,\mathrm{r\acute{e}d}}$ 是离散赋值环的谱，$V'_1$ 延拓到 $X'_1$ 的事实由 X 3.6 得出.

证明诸 $n_i$ 与 $p$ 互素. 若不然，有使 $p$ 整除 $n_{i_0}$ 的指标 $i_0 \in I(x)$. 把 $X$ 换成
\[
X_1[T_{i_0}, T_i]_{i \in I(x)}/(T_{i_0}^{n_{i_0}/p} - f_{i_0}, T_i^{n_i} - f_i),
\]
归结到 $X'_1 = X_1[T]/T^p - f_{i_0}$ 的情形. 由前面，$U'_1$ 的平展覆盖 $V'_1$ 延拓成 $X'_1$ 的平展覆盖 $E'_1$. 设 $\eta$ 是 $S$ 的闭点；由于态射 $X'_{1\eta} \to X_{1\eta}$ 根式，$V_{1\eta}$ 延拓成 $X_{1\eta}$ 的平展覆盖 $E_{1\eta}$. 由此如 \ref{prop:XIII.5.2} 得出，$E'_1$ 配备相对态射 $X'_1 \to X_1$ 的下降数据，它延拓 $E'_1|_{U'_1}$ 上的自然下降数据. 由此得出 $V_1$ 延拓到 $X_1$；但这蕴含 $n_{i_0} = 1$，与假设 $n_{i_0} = p$ 矛盾.

\begin{corollary}\label{cor:XIII.5.6}
设 $X$ 是 $S$-概形，$D = \sum_{1 \leqslant i \leqslant r} \mathrm{div}\, f_i$ 是如 \ref{prop:XIII.5.4} 的相对 $S$ 正规交叉除子. 设 $x$ 是 $X$ 的几何点，$X$ 是 $X$ 在 $x$ 处的严格局部化，$Y = Y_{(X)}$，$U = X - Y$，
\[
\widetilde{U} = \varprojlim_{(n_i)} U[T_i]_{i \in I(x)}/(T_i^{n_i} - f_i),
\]
投射极限按与 $k(x)$ 的特征 $p$ 互素的整数族 $n_i > 0$ 的滤过集取. 则 $\widetilde{U}$ 是 $U$ 的相对 $S$ 温和分歧万有覆盖. 从而 $U$ 的温和分歧基本群为
\[
\pi^t_1(U) \simeq \prod_{\ell \neq p} \mathbb{Z}_\ell[1]^{I(x)} \quad \text{(典范同构)}.
\]
群 $\pi^t_1(U)$ 非典范同构于 $\prod_{\ell \neq p} \mathbb{Z}_\ell^{I(x)}$.
\end{corollary}

\begin{remark}\label{rem:XIII.5.6.1}
设 $X$ 是 $S$-概形，$D = \sum_{1 \leqslant i \leqslant r} \mathrm{div}\, f_i$ 是如 \ref{prop:XIII.5.4} 的相对 $S$ 正规交叉除子，$U = X - \mathrm{Supp}\, D$. 对每个部分 $I \subset [1, r]$，设
\[
X_I = \Big(\bigcap_{i \in I} V(f_i)\Big) \cap \Big(\bigcap_{i \notin I} X_{f_i}\Big).
\]
设 $p$ 是素数或零，$Z$ 是 $X_I$ 的诸点特征都为 $p$ 的子集. 设
\[
\widetilde{U}_I = \varprojlim_{(n_i)} U[T_i]_{i \in I}/(T_i^{n_i} - f_i),
\]
其中投射极限按与 $p$ 互素的整数族 $n_i > 0$ 的滤过集取. 则对 $Z$ 的每个几何点 $x$，$\widetilde{U}_I$ 在 $U$ 上的逆像等同于 $U$ 的温和分歧万有覆盖.
\end{remark}

\begin{corollary}\label{cor:XIII.5.7}
记号同 5.6. 设 $S$ 是 $S$ 在 $x$ 处的严格局部化，
\[
g \colon U \longrightarrow S \quad \text{和} \quad \widetilde{g} \colon \widetilde{U} \longrightarrow \widetilde{S}
\]
是典范态射. 则态射 $g$ 和 $\widetilde{g}$ 0-无环 (SGA 4 XV 1.3). 设 $G$ 是 $S$ 上的可构群层，$F = g^{*} G$，$P$ 是 $F$ 下的挠子. 则 $P$ 相对 $S$ 在 $X$ 上温和分歧，当且仅当它在 $\widetilde{U}$ 上的逆像 $\widetilde{P}$ 平凡.
\end{corollary}

确实，对每个概形 $X' = X[T_i]_{i \in I(x)}/(T_i^{n_i} - f_i)$ (其中诸 $n_i$ 是与 $p$ 互素的整数 $> 0$)，态射 $f' \colon X' \to S$ 0-无环. $f'$ 在 $S$ 的各点处的几何纤维因此连通甚至不可约. 态射 $g' \colon U' \to S$ 的几何纤维也如此，这证明诸 $g'$，从而 $\widetilde{g}$，0-无环 (SGA 4 XV 1.16).

$U$ 上群 $F$ 的在 $\widetilde{U}$ 上的逆像平凡的挠子 $P$ 显然相对 $S$ 在 $X$ 上温和分歧. 证明反之，若 $P$ 相对 $S$ 在 $X$ 上温和分歧，它在 $\widetilde{U}$ 上的逆像平凡.

由 (SGA 4 IX 2.14 (ii)) 得出可以找到有限态射 $n \colon S_1 \to S$、$S_1$ 上的常群层 $C$ 和单态射 $G \to n_* C$. 考虑下列所有方形笛卡儿的交换图式：
\[
\begin{tikzcd}[row sep=large, column sep=large]
\widetilde{U}_1 \arrow[r] \arrow[d,"{r}"] & U_1 \arrow[r,"{g_1}"] \arrow[d,"{q}"] & S_1 \arrow[d,"{n}"] \\
\widetilde{U} \arrow[r] & U \arrow[r,"{g}"] & S.
\end{tikzcd}
\]
设 $C_1$ (resp. $\widetilde{C}_1$) 是 $C$ 在 $U_1$ (resp. $\widetilde{U}_1$) 上的逆像. 有 $i$ 和 $j$ 为同构 (SGA 4 VIII 5.8) 的交换图式：
\begin{equation}
\begin{tikzcd}[row sep=large, column sep=large]
H^1(U, q_* C_1) \arrow[r,"{i}"] \arrow[d] & H^1(U_1, C_1) \arrow[d,"{j}"] \\
H^1(\widetilde{U}, r_* \widetilde{C}_1) \arrow[r] & H^1(\widetilde{U}_1, \widetilde{C}_1).
\end{tikzcd}
\tag{$\ast$}
\end{equation}
设 $Q$ 是由 $P$ 经结构群扩张 $F \to q_* C_1$ 得到的 $q_* C_1$ 下挠子. 由 2.1.4，$Q$ 相对 $S$ 在 $X$ 上温和分歧. 与挠子 $Q$ 对应 (借助 $i$) $C_1$ 下的挠子 $Q_1$，且显然 $Q_1$ 相对 $S_1$ 在 $X_1 = X \times_S S_1$ 上温和分歧. 由 \ref{cor:XIII.5.6} 得出 $Q_1$ 在 $\widetilde{U}_1$ 上的逆像 $\widetilde{Q}_1$ 平凡，图式 ($\ast$) 于是表明 $Q$ 在 $\widetilde{U}$ 上的逆像 $\widetilde{Q}$ 平凡.

考虑下列第二行正合 (SGA 4 XII 3.1) 的交换图式：
\begin{equation}
\begin{tikzcd}[row sep=large, column sep=large]
H^0(S, n_* C/G) \arrow[r] \arrow[d,"{k}"] & H^1(S, G) \arrow[r, equal] & 1 \arrow[d] \\
H^0(\widetilde{U}, r_* \widetilde{C}_1/\widetilde{F}) \arrow[r] & H^1(\widetilde{U}, \widetilde{F}) \arrow[r] & H^1(\widetilde{U}, r_* \widetilde{C}_1).
\end{tikzcd}
\tag{$\ast\ast$}
\end{equation}
由于态射 $\widetilde{U} \to S$ 0-无环，$k$ 是同构. $\widetilde{P}$ 平凡的事实于是由 ($\ast\ast$) 得出.

\section{附录 II：叠的直像的有限性定理}

\begin{proposition}\label{prop:XIII.6.1}
设 $S$ 是局部诺特概形，$f \colon X \to S$ 是态射. 若 $S'$ 是 $S$-概形，记 $X'$ (resp. $f'$，等等) 为 $X$ (resp. $f$，等等) 经态射 $S' \to S$ 的逆像. 假设对每个平展于 $S$ 上的概形 $S'$ 和 $X'$ 上每个可构集合层 $F$，$f'_* F$ 可构，且对 $X'$ 上每个可构群层 $F$，$R^1 f'_* F$ 可构. 设 $\Phi$ 是 $X$ 上 (0) 的 1-可构叠. 则 $f_* \Phi$ 1-可构.
\end{proposition}

对每个平展于 $S$ 上的概形 $S'$ 和 $(f_* \Phi)_{S'}$ 的每个对象 $x$，有同构
\[
\mathrm{Aut}_{S'}(x) \simeq f'_* \mathrm{Aut}_{X'}(x),
\]
其中第二个成员中 $x$ 看作 $\Phi_{X'}$ 的对象. 所做假设于是蕴含 $f_* \Phi$ 可构. 设 $S_\Phi$ 是 $\Phi$ \cite[III 2.1.7]{exp13-1} 的极大子叠丛层. 由于 $f_*(S_\Phi)$ 可构，可以给它应用 (SGA 4 IX 2.7)，$f_* \Phi$ 1-可构的事实于是由下列引理得出.

\begin{lemma}\label{lem:XIII.6.1.1}
设 $S$ 是局部诺特概形，$f \colon X \to S$ 是态射，$\Phi$ 是 $X$ 上的叠. 假设给定 $S$ 上由有限型平展 $S$-概形 $T$ 表示的层、满态射
\[
a \colon T \longrightarrow f_*(S_\Phi)
\]
和纤维 $\Phi_{X_T}$ (其中 $X_T = X \times_S T$) 的对象 $p$，它在 $f_*(S_\Phi)(T) = S_\Phi(X_T)$ 中定义等于 $T(T)$ 的恒等截面经 $a$ 的像 $q$ 的元素. 设 $f_T \colon X_T \to T$ 是典范态射，假设层 $R^1 f_{T*}(\mathrm{Aut}_{X_T}(p))$ 可构；则 $S(f_* \Phi)$ 也可构.
\end{lemma}

典范态射 $f^{*} f_* \Phi \to \Phi$ 给出态射
\[
S(f^{*} f_* \Phi) \simeq f^{*}(S(f_* \Phi)) \longrightarrow S_\Phi,
\]
从而有典范态射
\[
\phi \colon S(f_* \Phi) \longrightarrow f_*(S_\Phi).
\]
设 $F = S(f_* \Phi)$，$G$ 是 $F$ 经 $\phi$ 的像；由 (SGA 4 IX 2.9)，$G$ 是可构层.

只需证明对 $S$ 的每个点 $s$，存在 $s$ 的非空开集 $U$，使 $F|_U$ 局部常可构. 设 $s \in S$，$s$ 是 $s$ 上方的几何点，$q_1, \ldots, q_n$ 是 $G_s$ 的元素. 由 $T$ 的定义，存在 $S$-态射 $h_i \colon s \to S'$，使 $q_i = h^{*}_i(q)$. 设 $S'$ 是 $n$ 个同构于 $T$ 的概形在 $S$ 上的纤维积，$s \to S'$ 是诸 $h_i$ 的纤维积，$X' = X \times_S S'$，$q_i$ (resp. $p_i$) 是 $q$ (resp. $p$) 经 $S'$ 到 $T$ 的第 $i$ 个投影的逆像. 若 $\Psi_i$ 是 $p_i$ 生成的 $\Phi|_{X'}$ 的极大子叠丛，层 $F_i = R^1 f'_*(\mathrm{Aut}_{X'}(p_i))$ 正是 $f'_* \Psi_i$ 的极大子叠丛层 $S(f'_* \Psi_i)$. 特别典范单射 $\Psi_i \to \Phi|_{X'}$ 给出态射
\[
\alpha_i \colon F_i \longrightarrow F|_{S'}.
\]
证明 $\alpha_i$ 是 $F_i$ 到 $q_i$ 在 $F|_{S'}$ 中的逆像的双射. 对每个平展于 $S'$ 上的概形 $S''$，$F_i(S'')$ 的每个截面 $y$ 在 $F(S'')$ 中的像是 $q_i|_{S''}$，因为局部于 $S''$ 的平展拓扑，$y$ 由同构于 $p_i|_{X''}$ 的 $\Phi_{X''}$ 的对象 $x$ 定义. 反之，若 $y \in F(S'')$ 在 $F(S'')$ 中的像是 $q_i|_{S''}$，局部于 $S''$ 的平展拓扑，$y$ 由同构于 $p_i$ 的 $\Phi_{X''}$ 的对象 $x$ 定义；从而 $x$ 是 $\Psi_{iX''}$ 的对象，故 $y \in F_i(S'')$.

用下面 \ref{lem:XIII.6.1.2} 完成证明. 可以找到 $s$ 的开邻域 $U'$，使 $q_1|_{U'}, \ldots, q_n|_{U'}$ 是 $G(U')$ 的截面并生成该层. 由于诸 $F_i|_{U'}$ 和 $G|_{U'}$ 可构，由 \ref{lem:XIII.6.1.2} 得出 $F|_{U'}$ 也可构；把 $U$ 换成更小的开集后 $F|_U$ 局部常，完成证明.

\begin{lemma}\label{lem:XIII.6.1.2}
设 $S$ 是局部诺特概形，$F \to G$ 是 $S$ 上群层的满态射. 设 $q_i$ 是 $X$ 上生成 $G$ 的 $G$ 截面的有限族，对每个 $i$，设 $F_i$ 是 $F$ 的 $q_i$ 的逆像子层. 则若 $G$ 和诸 $F_i$ 可构，$F$ 也可构.
\end{lemma}

为证明 $F$ 可构，只需证明对 $S$ 的每个点 $s$，存在 $s$ 的开邻域 $U$，使 $F|_U$ 局部常可构. 设 $s$ 是 $S$ 的点. 由于层 $F_i$ 和 $G$ 可构，可以找到 $s$ 的开邻域 $U$，使 $F_i|_U$ 和 $G|_U$ 局部常. 证明 $F|_U$ 局部常. 由 (SGA 4 IX 2.13 (i))，只需看：若 $s$ 是 $s$ 上方的几何点，$\widetilde{s}$ 是 $U$ 的几何点，$s \to \widetilde{s}$ 是特殊化态射，典范态射
\[
F_{\widetilde{s}} \longrightarrow F_s
\]
双射.

考虑交换图式
\[
\begin{tikzcd}[row sep=large, column sep=large]
(F_i)_{\widetilde{s}} \arrow[r,"{\widetilde{a}_i}"] \arrow[d,"{\simeq}"] & F_{\widetilde{s}} \arrow[r,"{\widetilde{a}}"] \arrow[d,"{b}"] & G_{\widetilde{s}} \arrow[d,"{\simeq}"] \\
(F_i)_s \arrow[r,"{a_i}"'] & F_s \arrow[r,"{a}"'] & G_s
\end{tikzcd}
\]
设 $q_i$ (resp. $\widetilde{q}_i$) 是 $q_i$ 在 $G_s$ (resp. $G_{\widetilde{s}}$) 中的逆像；态射 $a$ 和 $\widetilde{a}$ 满射，态射 $a_i$ (resp. $\widetilde{a}_i$) 诱导 $(F_i)_s$ (resp. $(F_i)_{\widetilde{s}}$) 到 $a^{-1}(q_i)$ (resp. $\widetilde{a}^{-1}(q_i)$) 的双射. 由上面图式得出 $b$ 是同构.

\begin{corollary}\label{cor:XIII.6.2}
设 $S$ 是局部诺特概形，$f \colon X \to S$ 是紧合态射. 设 $\Phi$ 是 $X$ 上 1-可构叠，则 $f_* \Phi$ 是 1-可构叠.
\end{corollary}

\ref{prop:XIII.6.1} 的证明也证明下列结果，考虑到 \ref{thm:XIII.2.4} 2)：

\begin{corollary}\label{cor:XIII.6.3}
设 $S$ 是局部诺特概形，$f \colon X \to S$ 是态射，$D$ 是 $X$ 上相对 $S$ 正规交叉 (2.1) 的除子，$Y = \mathrm{Supp}\, D$，$U = X - Y$，$i \colon U \to X$ 是典范浸入. 设 $\Phi$ 是 $U$ 上的叠，局部于 $X$ 和 $S$ 的平展拓扑是 $S$ 上 1-可构叠 $\Psi$ 的逆像. 则叠 $i_{*t} \Phi$ 1-可构.
\end{corollary}

\section{文献}\printbibliography[heading=none]


\end{refsection}

\backmatter
\printindex[term]
\printindex[notations]

\end{document}
